\documentclass[a4paper,12pt]{article}

% --- Pakete ---
\usepackage[utf8]{inputenc}   % Umlaute direkt eingeben
\usepackage[T1]{fontenc}
\usepackage[english]{babel}
\usepackage{lmodern}          % Schöne Schrift
\usepackage{geometry}         % Seitenränder anpassen
\usepackage{booktabs}         % Für schönere Tabellen
\usepackage{array}
\usepackage{setspace}         % Zeilenabstand
\usepackage{graphicx}
\usepackage{caption}
\usepackage{url}
\usepackage{todonotes}
\usepackage{float}
\usepackage{pgfplots}
\usepackage{xfp}      % provides \fpeval for calculations
\usepackage{multirow}
\usepackage[hidelinks, bookmarks=true]{hyperref}
\pgfplotsset{compat=1.18}

% --- Einstellungen ---
\geometry{a4paper, margin=2.5cm}
\setstretch{1.25}

% --- Dokument ---
\begin{document}
	
	% Titel
	\begin{center}
		{\LARGE \textbf{Praktikum Hochfrequenz-Schaltungstechnik WS25/26}} \\[1em]
		{\large  \textbf{Bericht}} \\[0.5em]
	\end{center}
	
\begin{center}
	\begin{tabular}{|l|p{4cm}|l|p{4cm}|}
		\hline
		\textbf{Name:} &  Sebastian Grigorevski & \textbf{Matrikelnr.:} &  35690104 \\
		\hline
		\textbf{Versuchsnr.:} & 1 & \textbf{Datum:} & 11.11.2025 \\
		\hline
	\end{tabular}
\end{center}
	
	\vspace{1cm}
	\setcounter{section}{1}  % setzt den Section Zähler auf 1
	\setcounter{subsection}{5}  % setzt den Subsection Zähler auf 5
	\setcounter{subsubsection}{0}  % setzt den Subsubsection Zähler auf 0

	\subsubsection{}
	Figure \ref{fig:resistors} shows the measured resistance $R_s$ and impedance $|Z|$ of a carbon resistor and a coaxial termination over a frequency range between 1 MHz and 30 MHz. % \cite[]{ltspice}. 
	The coaxial termination exhibits behavior very close to the ideal expectation. Both the $R_s$ and the $|Z|$ remain nearly constant at $\sim 50 \Omega$ over the entire frequency range. Thus it it well impedance matched.
	
	In contrast, the carbon resistor exhibits a clear frequency dependency. The impedance magnitude increases with frequency and exceeds 51~$\Omega$ at higher frequencies. This behavior is caused by parasitic inductive effects inherent to the resistor’s internal structure and material.
	
	
	\begin{figure}[H]
		\centering
		%% Creator: Matplotlib, PGF backend
%%
%% To include the figure in your LaTeX document, write
%%   \input{<filename>.pgf}
%%
%% Make sure the required packages are loaded in your preamble
%%   \usepackage{pgf}
%%
%% Also ensure that all the required font packages are loaded; for instance,
%% the lmodern package is sometimes necessary when using math font.
%%   \usepackage{lmodern}
%%
%% Figures using additional raster images can only be included by \input if
%% they are in the same directory as the main LaTeX file. For loading figures
%% from other directories you can use the `import` package
%%   \usepackage{import}
%%
%% and then include the figures with
%%   \import{<path to file>}{<filename>.pgf}
%%
%% Matplotlib used the following preamble
%%
\begingroup%
\makeatletter%
\begin{pgfpicture}%
\pgfpathrectangle{\pgfpointorigin}{\pgfqpoint{6.300000in}{3.500000in}}%
\pgfusepath{use as bounding box, clip}%
\begin{pgfscope}%
\pgfsetbuttcap%
\pgfsetmiterjoin%
\definecolor{currentfill}{rgb}{1.000000,1.000000,1.000000}%
\pgfsetfillcolor{currentfill}%
\pgfsetlinewidth{0.000000pt}%
\definecolor{currentstroke}{rgb}{1.000000,1.000000,1.000000}%
\pgfsetstrokecolor{currentstroke}%
\pgfsetdash{}{0pt}%
\pgfpathmoveto{\pgfqpoint{0.000000in}{0.000000in}}%
\pgfpathlineto{\pgfqpoint{6.300000in}{0.000000in}}%
\pgfpathlineto{\pgfqpoint{6.300000in}{3.500000in}}%
\pgfpathlineto{\pgfqpoint{0.000000in}{3.500000in}}%
\pgfpathlineto{\pgfqpoint{0.000000in}{0.000000in}}%
\pgfpathclose%
\pgfusepath{fill}%
\end{pgfscope}%
\begin{pgfscope}%
\pgfsetbuttcap%
\pgfsetmiterjoin%
\definecolor{currentfill}{rgb}{1.000000,1.000000,1.000000}%
\pgfsetfillcolor{currentfill}%
\pgfsetlinewidth{0.000000pt}%
\definecolor{currentstroke}{rgb}{0.000000,0.000000,0.000000}%
\pgfsetstrokecolor{currentstroke}%
\pgfsetstrokeopacity{0.000000}%
\pgfsetdash{}{0pt}%
\pgfpathmoveto{\pgfqpoint{0.688581in}{0.565123in}}%
\pgfpathlineto{\pgfqpoint{6.061265in}{0.565123in}}%
\pgfpathlineto{\pgfqpoint{6.061265in}{3.301775in}}%
\pgfpathlineto{\pgfqpoint{0.688581in}{3.301775in}}%
\pgfpathlineto{\pgfqpoint{0.688581in}{0.565123in}}%
\pgfpathclose%
\pgfusepath{fill}%
\end{pgfscope}%
\begin{pgfscope}%
\pgfpathrectangle{\pgfqpoint{0.688581in}{0.565123in}}{\pgfqpoint{5.372684in}{2.736651in}}%
\pgfusepath{clip}%
\pgfsetbuttcap%
\pgfsetroundjoin%
\pgfsetlinewidth{0.803000pt}%
\definecolor{currentstroke}{rgb}{0.690196,0.690196,0.690196}%
\pgfsetstrokecolor{currentstroke}%
\pgfsetdash{{0.800000pt}{1.320000pt}}{0.000000pt}%
\pgfpathmoveto{\pgfqpoint{1.429641in}{0.565123in}}%
\pgfpathlineto{\pgfqpoint{1.429641in}{3.301775in}}%
\pgfusepath{stroke}%
\end{pgfscope}%
\begin{pgfscope}%
\pgfsetbuttcap%
\pgfsetroundjoin%
\definecolor{currentfill}{rgb}{0.000000,0.000000,0.000000}%
\pgfsetfillcolor{currentfill}%
\pgfsetlinewidth{0.803000pt}%
\definecolor{currentstroke}{rgb}{0.000000,0.000000,0.000000}%
\pgfsetstrokecolor{currentstroke}%
\pgfsetdash{}{0pt}%
\pgfsys@defobject{currentmarker}{\pgfqpoint{0.000000in}{-0.048611in}}{\pgfqpoint{0.000000in}{0.000000in}}{%
\pgfpathmoveto{\pgfqpoint{0.000000in}{0.000000in}}%
\pgfpathlineto{\pgfqpoint{0.000000in}{-0.048611in}}%
\pgfusepath{stroke,fill}%
}%
\begin{pgfscope}%
\pgfsys@transformshift{1.429641in}{0.565123in}%
\pgfsys@useobject{currentmarker}{}%
\end{pgfscope}%
\end{pgfscope}%
\begin{pgfscope}%
\definecolor{textcolor}{rgb}{0.000000,0.000000,0.000000}%
\pgfsetstrokecolor{textcolor}%
\pgfsetfillcolor{textcolor}%
\pgftext[x=1.429641in,y=0.467901in,,top]{\color{textcolor}\rmfamily\fontsize{10.000000}{12.000000}\selectfont \(\displaystyle {0.5}\)}%
\end{pgfscope}%
\begin{pgfscope}%
\pgfpathrectangle{\pgfqpoint{0.688581in}{0.565123in}}{\pgfqpoint{5.372684in}{2.736651in}}%
\pgfusepath{clip}%
\pgfsetbuttcap%
\pgfsetroundjoin%
\pgfsetlinewidth{0.803000pt}%
\definecolor{currentstroke}{rgb}{0.690196,0.690196,0.690196}%
\pgfsetstrokecolor{currentstroke}%
\pgfsetdash{{0.800000pt}{1.320000pt}}{0.000000pt}%
\pgfpathmoveto{\pgfqpoint{2.355966in}{0.565123in}}%
\pgfpathlineto{\pgfqpoint{2.355966in}{3.301775in}}%
\pgfusepath{stroke}%
\end{pgfscope}%
\begin{pgfscope}%
\pgfsetbuttcap%
\pgfsetroundjoin%
\definecolor{currentfill}{rgb}{0.000000,0.000000,0.000000}%
\pgfsetfillcolor{currentfill}%
\pgfsetlinewidth{0.803000pt}%
\definecolor{currentstroke}{rgb}{0.000000,0.000000,0.000000}%
\pgfsetstrokecolor{currentstroke}%
\pgfsetdash{}{0pt}%
\pgfsys@defobject{currentmarker}{\pgfqpoint{0.000000in}{-0.048611in}}{\pgfqpoint{0.000000in}{0.000000in}}{%
\pgfpathmoveto{\pgfqpoint{0.000000in}{0.000000in}}%
\pgfpathlineto{\pgfqpoint{0.000000in}{-0.048611in}}%
\pgfusepath{stroke,fill}%
}%
\begin{pgfscope}%
\pgfsys@transformshift{2.355966in}{0.565123in}%
\pgfsys@useobject{currentmarker}{}%
\end{pgfscope}%
\end{pgfscope}%
\begin{pgfscope}%
\definecolor{textcolor}{rgb}{0.000000,0.000000,0.000000}%
\pgfsetstrokecolor{textcolor}%
\pgfsetfillcolor{textcolor}%
\pgftext[x=2.355966in,y=0.467901in,,top]{\color{textcolor}\rmfamily\fontsize{10.000000}{12.000000}\selectfont \(\displaystyle {1.0}\)}%
\end{pgfscope}%
\begin{pgfscope}%
\pgfpathrectangle{\pgfqpoint{0.688581in}{0.565123in}}{\pgfqpoint{5.372684in}{2.736651in}}%
\pgfusepath{clip}%
\pgfsetbuttcap%
\pgfsetroundjoin%
\pgfsetlinewidth{0.803000pt}%
\definecolor{currentstroke}{rgb}{0.690196,0.690196,0.690196}%
\pgfsetstrokecolor{currentstroke}%
\pgfsetdash{{0.800000pt}{1.320000pt}}{0.000000pt}%
\pgfpathmoveto{\pgfqpoint{3.282291in}{0.565123in}}%
\pgfpathlineto{\pgfqpoint{3.282291in}{3.301775in}}%
\pgfusepath{stroke}%
\end{pgfscope}%
\begin{pgfscope}%
\pgfsetbuttcap%
\pgfsetroundjoin%
\definecolor{currentfill}{rgb}{0.000000,0.000000,0.000000}%
\pgfsetfillcolor{currentfill}%
\pgfsetlinewidth{0.803000pt}%
\definecolor{currentstroke}{rgb}{0.000000,0.000000,0.000000}%
\pgfsetstrokecolor{currentstroke}%
\pgfsetdash{}{0pt}%
\pgfsys@defobject{currentmarker}{\pgfqpoint{0.000000in}{-0.048611in}}{\pgfqpoint{0.000000in}{0.000000in}}{%
\pgfpathmoveto{\pgfqpoint{0.000000in}{0.000000in}}%
\pgfpathlineto{\pgfqpoint{0.000000in}{-0.048611in}}%
\pgfusepath{stroke,fill}%
}%
\begin{pgfscope}%
\pgfsys@transformshift{3.282291in}{0.565123in}%
\pgfsys@useobject{currentmarker}{}%
\end{pgfscope}%
\end{pgfscope}%
\begin{pgfscope}%
\definecolor{textcolor}{rgb}{0.000000,0.000000,0.000000}%
\pgfsetstrokecolor{textcolor}%
\pgfsetfillcolor{textcolor}%
\pgftext[x=3.282291in,y=0.467901in,,top]{\color{textcolor}\rmfamily\fontsize{10.000000}{12.000000}\selectfont \(\displaystyle {1.5}\)}%
\end{pgfscope}%
\begin{pgfscope}%
\pgfpathrectangle{\pgfqpoint{0.688581in}{0.565123in}}{\pgfqpoint{5.372684in}{2.736651in}}%
\pgfusepath{clip}%
\pgfsetbuttcap%
\pgfsetroundjoin%
\pgfsetlinewidth{0.803000pt}%
\definecolor{currentstroke}{rgb}{0.690196,0.690196,0.690196}%
\pgfsetstrokecolor{currentstroke}%
\pgfsetdash{{0.800000pt}{1.320000pt}}{0.000000pt}%
\pgfpathmoveto{\pgfqpoint{4.208615in}{0.565123in}}%
\pgfpathlineto{\pgfqpoint{4.208615in}{3.301775in}}%
\pgfusepath{stroke}%
\end{pgfscope}%
\begin{pgfscope}%
\pgfsetbuttcap%
\pgfsetroundjoin%
\definecolor{currentfill}{rgb}{0.000000,0.000000,0.000000}%
\pgfsetfillcolor{currentfill}%
\pgfsetlinewidth{0.803000pt}%
\definecolor{currentstroke}{rgb}{0.000000,0.000000,0.000000}%
\pgfsetstrokecolor{currentstroke}%
\pgfsetdash{}{0pt}%
\pgfsys@defobject{currentmarker}{\pgfqpoint{0.000000in}{-0.048611in}}{\pgfqpoint{0.000000in}{0.000000in}}{%
\pgfpathmoveto{\pgfqpoint{0.000000in}{0.000000in}}%
\pgfpathlineto{\pgfqpoint{0.000000in}{-0.048611in}}%
\pgfusepath{stroke,fill}%
}%
\begin{pgfscope}%
\pgfsys@transformshift{4.208615in}{0.565123in}%
\pgfsys@useobject{currentmarker}{}%
\end{pgfscope}%
\end{pgfscope}%
\begin{pgfscope}%
\definecolor{textcolor}{rgb}{0.000000,0.000000,0.000000}%
\pgfsetstrokecolor{textcolor}%
\pgfsetfillcolor{textcolor}%
\pgftext[x=4.208615in,y=0.467901in,,top]{\color{textcolor}\rmfamily\fontsize{10.000000}{12.000000}\selectfont \(\displaystyle {2.0}\)}%
\end{pgfscope}%
\begin{pgfscope}%
\pgfpathrectangle{\pgfqpoint{0.688581in}{0.565123in}}{\pgfqpoint{5.372684in}{2.736651in}}%
\pgfusepath{clip}%
\pgfsetbuttcap%
\pgfsetroundjoin%
\pgfsetlinewidth{0.803000pt}%
\definecolor{currentstroke}{rgb}{0.690196,0.690196,0.690196}%
\pgfsetstrokecolor{currentstroke}%
\pgfsetdash{{0.800000pt}{1.320000pt}}{0.000000pt}%
\pgfpathmoveto{\pgfqpoint{5.134940in}{0.565123in}}%
\pgfpathlineto{\pgfqpoint{5.134940in}{3.301775in}}%
\pgfusepath{stroke}%
\end{pgfscope}%
\begin{pgfscope}%
\pgfsetbuttcap%
\pgfsetroundjoin%
\definecolor{currentfill}{rgb}{0.000000,0.000000,0.000000}%
\pgfsetfillcolor{currentfill}%
\pgfsetlinewidth{0.803000pt}%
\definecolor{currentstroke}{rgb}{0.000000,0.000000,0.000000}%
\pgfsetstrokecolor{currentstroke}%
\pgfsetdash{}{0pt}%
\pgfsys@defobject{currentmarker}{\pgfqpoint{0.000000in}{-0.048611in}}{\pgfqpoint{0.000000in}{0.000000in}}{%
\pgfpathmoveto{\pgfqpoint{0.000000in}{0.000000in}}%
\pgfpathlineto{\pgfqpoint{0.000000in}{-0.048611in}}%
\pgfusepath{stroke,fill}%
}%
\begin{pgfscope}%
\pgfsys@transformshift{5.134940in}{0.565123in}%
\pgfsys@useobject{currentmarker}{}%
\end{pgfscope}%
\end{pgfscope}%
\begin{pgfscope}%
\definecolor{textcolor}{rgb}{0.000000,0.000000,0.000000}%
\pgfsetstrokecolor{textcolor}%
\pgfsetfillcolor{textcolor}%
\pgftext[x=5.134940in,y=0.467901in,,top]{\color{textcolor}\rmfamily\fontsize{10.000000}{12.000000}\selectfont \(\displaystyle {2.5}\)}%
\end{pgfscope}%
\begin{pgfscope}%
\pgfpathrectangle{\pgfqpoint{0.688581in}{0.565123in}}{\pgfqpoint{5.372684in}{2.736651in}}%
\pgfusepath{clip}%
\pgfsetbuttcap%
\pgfsetroundjoin%
\pgfsetlinewidth{0.803000pt}%
\definecolor{currentstroke}{rgb}{0.690196,0.690196,0.690196}%
\pgfsetstrokecolor{currentstroke}%
\pgfsetdash{{0.800000pt}{1.320000pt}}{0.000000pt}%
\pgfpathmoveto{\pgfqpoint{6.061265in}{0.565123in}}%
\pgfpathlineto{\pgfqpoint{6.061265in}{3.301775in}}%
\pgfusepath{stroke}%
\end{pgfscope}%
\begin{pgfscope}%
\pgfsetbuttcap%
\pgfsetroundjoin%
\definecolor{currentfill}{rgb}{0.000000,0.000000,0.000000}%
\pgfsetfillcolor{currentfill}%
\pgfsetlinewidth{0.803000pt}%
\definecolor{currentstroke}{rgb}{0.000000,0.000000,0.000000}%
\pgfsetstrokecolor{currentstroke}%
\pgfsetdash{}{0pt}%
\pgfsys@defobject{currentmarker}{\pgfqpoint{0.000000in}{-0.048611in}}{\pgfqpoint{0.000000in}{0.000000in}}{%
\pgfpathmoveto{\pgfqpoint{0.000000in}{0.000000in}}%
\pgfpathlineto{\pgfqpoint{0.000000in}{-0.048611in}}%
\pgfusepath{stroke,fill}%
}%
\begin{pgfscope}%
\pgfsys@transformshift{6.061265in}{0.565123in}%
\pgfsys@useobject{currentmarker}{}%
\end{pgfscope}%
\end{pgfscope}%
\begin{pgfscope}%
\definecolor{textcolor}{rgb}{0.000000,0.000000,0.000000}%
\pgfsetstrokecolor{textcolor}%
\pgfsetfillcolor{textcolor}%
\pgftext[x=6.061265in,y=0.467901in,,top]{\color{textcolor}\rmfamily\fontsize{10.000000}{12.000000}\selectfont \(\displaystyle {3.0}\)}%
\end{pgfscope}%
\begin{pgfscope}%
\definecolor{textcolor}{rgb}{0.000000,0.000000,0.000000}%
\pgfsetstrokecolor{textcolor}%
\pgfsetfillcolor{textcolor}%
\pgftext[x=3.374923in,y=0.288889in,,top]{\color{textcolor}\rmfamily\fontsize{10.000000}{12.000000}\selectfont Frequency (Hz)}%
\end{pgfscope}%
\begin{pgfscope}%
\definecolor{textcolor}{rgb}{0.000000,0.000000,0.000000}%
\pgfsetstrokecolor{textcolor}%
\pgfsetfillcolor{textcolor}%
\pgftext[x=6.061265in,y=0.302778in,right,top]{\color{textcolor}\rmfamily\fontsize{10.000000}{12.000000}\selectfont \(\displaystyle \times{10^{7}}{}\)}%
\end{pgfscope}%
\begin{pgfscope}%
\pgfpathrectangle{\pgfqpoint{0.688581in}{0.565123in}}{\pgfqpoint{5.372684in}{2.736651in}}%
\pgfusepath{clip}%
\pgfsetbuttcap%
\pgfsetroundjoin%
\pgfsetlinewidth{0.803000pt}%
\definecolor{currentstroke}{rgb}{0.690196,0.690196,0.690196}%
\pgfsetstrokecolor{currentstroke}%
\pgfsetdash{{0.800000pt}{1.320000pt}}{0.000000pt}%
\pgfpathmoveto{\pgfqpoint{0.688581in}{0.565123in}}%
\pgfpathlineto{\pgfqpoint{6.061265in}{0.565123in}}%
\pgfusepath{stroke}%
\end{pgfscope}%
\begin{pgfscope}%
\pgfsetbuttcap%
\pgfsetroundjoin%
\definecolor{currentfill}{rgb}{0.000000,0.000000,0.000000}%
\pgfsetfillcolor{currentfill}%
\pgfsetlinewidth{0.803000pt}%
\definecolor{currentstroke}{rgb}{0.000000,0.000000,0.000000}%
\pgfsetstrokecolor{currentstroke}%
\pgfsetdash{}{0pt}%
\pgfsys@defobject{currentmarker}{\pgfqpoint{-0.048611in}{0.000000in}}{\pgfqpoint{-0.000000in}{0.000000in}}{%
\pgfpathmoveto{\pgfqpoint{-0.000000in}{0.000000in}}%
\pgfpathlineto{\pgfqpoint{-0.048611in}{0.000000in}}%
\pgfusepath{stroke,fill}%
}%
\begin{pgfscope}%
\pgfsys@transformshift{0.688581in}{0.565123in}%
\pgfsys@useobject{currentmarker}{}%
\end{pgfscope}%
\end{pgfscope}%
\begin{pgfscope}%
\definecolor{textcolor}{rgb}{0.000000,0.000000,0.000000}%
\pgfsetstrokecolor{textcolor}%
\pgfsetfillcolor{textcolor}%
\pgftext[x=0.344444in, y=0.516898in, left, base]{\color{textcolor}\rmfamily\fontsize{10.000000}{12.000000}\selectfont \(\displaystyle {48.0}\)}%
\end{pgfscope}%
\begin{pgfscope}%
\pgfpathrectangle{\pgfqpoint{0.688581in}{0.565123in}}{\pgfqpoint{5.372684in}{2.736651in}}%
\pgfusepath{clip}%
\pgfsetbuttcap%
\pgfsetroundjoin%
\pgfsetlinewidth{0.803000pt}%
\definecolor{currentstroke}{rgb}{0.690196,0.690196,0.690196}%
\pgfsetstrokecolor{currentstroke}%
\pgfsetdash{{0.800000pt}{1.320000pt}}{0.000000pt}%
\pgfpathmoveto{\pgfqpoint{0.688581in}{0.907205in}}%
\pgfpathlineto{\pgfqpoint{6.061265in}{0.907205in}}%
\pgfusepath{stroke}%
\end{pgfscope}%
\begin{pgfscope}%
\pgfsetbuttcap%
\pgfsetroundjoin%
\definecolor{currentfill}{rgb}{0.000000,0.000000,0.000000}%
\pgfsetfillcolor{currentfill}%
\pgfsetlinewidth{0.803000pt}%
\definecolor{currentstroke}{rgb}{0.000000,0.000000,0.000000}%
\pgfsetstrokecolor{currentstroke}%
\pgfsetdash{}{0pt}%
\pgfsys@defobject{currentmarker}{\pgfqpoint{-0.048611in}{0.000000in}}{\pgfqpoint{-0.000000in}{0.000000in}}{%
\pgfpathmoveto{\pgfqpoint{-0.000000in}{0.000000in}}%
\pgfpathlineto{\pgfqpoint{-0.048611in}{0.000000in}}%
\pgfusepath{stroke,fill}%
}%
\begin{pgfscope}%
\pgfsys@transformshift{0.688581in}{0.907205in}%
\pgfsys@useobject{currentmarker}{}%
\end{pgfscope}%
\end{pgfscope}%
\begin{pgfscope}%
\definecolor{textcolor}{rgb}{0.000000,0.000000,0.000000}%
\pgfsetstrokecolor{textcolor}%
\pgfsetfillcolor{textcolor}%
\pgftext[x=0.344444in, y=0.858979in, left, base]{\color{textcolor}\rmfamily\fontsize{10.000000}{12.000000}\selectfont \(\displaystyle {48.5}\)}%
\end{pgfscope}%
\begin{pgfscope}%
\pgfpathrectangle{\pgfqpoint{0.688581in}{0.565123in}}{\pgfqpoint{5.372684in}{2.736651in}}%
\pgfusepath{clip}%
\pgfsetbuttcap%
\pgfsetroundjoin%
\pgfsetlinewidth{0.803000pt}%
\definecolor{currentstroke}{rgb}{0.690196,0.690196,0.690196}%
\pgfsetstrokecolor{currentstroke}%
\pgfsetdash{{0.800000pt}{1.320000pt}}{0.000000pt}%
\pgfpathmoveto{\pgfqpoint{0.688581in}{1.249286in}}%
\pgfpathlineto{\pgfqpoint{6.061265in}{1.249286in}}%
\pgfusepath{stroke}%
\end{pgfscope}%
\begin{pgfscope}%
\pgfsetbuttcap%
\pgfsetroundjoin%
\definecolor{currentfill}{rgb}{0.000000,0.000000,0.000000}%
\pgfsetfillcolor{currentfill}%
\pgfsetlinewidth{0.803000pt}%
\definecolor{currentstroke}{rgb}{0.000000,0.000000,0.000000}%
\pgfsetstrokecolor{currentstroke}%
\pgfsetdash{}{0pt}%
\pgfsys@defobject{currentmarker}{\pgfqpoint{-0.048611in}{0.000000in}}{\pgfqpoint{-0.000000in}{0.000000in}}{%
\pgfpathmoveto{\pgfqpoint{-0.000000in}{0.000000in}}%
\pgfpathlineto{\pgfqpoint{-0.048611in}{0.000000in}}%
\pgfusepath{stroke,fill}%
}%
\begin{pgfscope}%
\pgfsys@transformshift{0.688581in}{1.249286in}%
\pgfsys@useobject{currentmarker}{}%
\end{pgfscope}%
\end{pgfscope}%
\begin{pgfscope}%
\definecolor{textcolor}{rgb}{0.000000,0.000000,0.000000}%
\pgfsetstrokecolor{textcolor}%
\pgfsetfillcolor{textcolor}%
\pgftext[x=0.344444in, y=1.201061in, left, base]{\color{textcolor}\rmfamily\fontsize{10.000000}{12.000000}\selectfont \(\displaystyle {49.0}\)}%
\end{pgfscope}%
\begin{pgfscope}%
\pgfpathrectangle{\pgfqpoint{0.688581in}{0.565123in}}{\pgfqpoint{5.372684in}{2.736651in}}%
\pgfusepath{clip}%
\pgfsetbuttcap%
\pgfsetroundjoin%
\pgfsetlinewidth{0.803000pt}%
\definecolor{currentstroke}{rgb}{0.690196,0.690196,0.690196}%
\pgfsetstrokecolor{currentstroke}%
\pgfsetdash{{0.800000pt}{1.320000pt}}{0.000000pt}%
\pgfpathmoveto{\pgfqpoint{0.688581in}{1.591368in}}%
\pgfpathlineto{\pgfqpoint{6.061265in}{1.591368in}}%
\pgfusepath{stroke}%
\end{pgfscope}%
\begin{pgfscope}%
\pgfsetbuttcap%
\pgfsetroundjoin%
\definecolor{currentfill}{rgb}{0.000000,0.000000,0.000000}%
\pgfsetfillcolor{currentfill}%
\pgfsetlinewidth{0.803000pt}%
\definecolor{currentstroke}{rgb}{0.000000,0.000000,0.000000}%
\pgfsetstrokecolor{currentstroke}%
\pgfsetdash{}{0pt}%
\pgfsys@defobject{currentmarker}{\pgfqpoint{-0.048611in}{0.000000in}}{\pgfqpoint{-0.000000in}{0.000000in}}{%
\pgfpathmoveto{\pgfqpoint{-0.000000in}{0.000000in}}%
\pgfpathlineto{\pgfqpoint{-0.048611in}{0.000000in}}%
\pgfusepath{stroke,fill}%
}%
\begin{pgfscope}%
\pgfsys@transformshift{0.688581in}{1.591368in}%
\pgfsys@useobject{currentmarker}{}%
\end{pgfscope}%
\end{pgfscope}%
\begin{pgfscope}%
\definecolor{textcolor}{rgb}{0.000000,0.000000,0.000000}%
\pgfsetstrokecolor{textcolor}%
\pgfsetfillcolor{textcolor}%
\pgftext[x=0.344444in, y=1.543142in, left, base]{\color{textcolor}\rmfamily\fontsize{10.000000}{12.000000}\selectfont \(\displaystyle {49.5}\)}%
\end{pgfscope}%
\begin{pgfscope}%
\pgfpathrectangle{\pgfqpoint{0.688581in}{0.565123in}}{\pgfqpoint{5.372684in}{2.736651in}}%
\pgfusepath{clip}%
\pgfsetbuttcap%
\pgfsetroundjoin%
\pgfsetlinewidth{0.803000pt}%
\definecolor{currentstroke}{rgb}{0.690196,0.690196,0.690196}%
\pgfsetstrokecolor{currentstroke}%
\pgfsetdash{{0.800000pt}{1.320000pt}}{0.000000pt}%
\pgfpathmoveto{\pgfqpoint{0.688581in}{1.933449in}}%
\pgfpathlineto{\pgfqpoint{6.061265in}{1.933449in}}%
\pgfusepath{stroke}%
\end{pgfscope}%
\begin{pgfscope}%
\pgfsetbuttcap%
\pgfsetroundjoin%
\definecolor{currentfill}{rgb}{0.000000,0.000000,0.000000}%
\pgfsetfillcolor{currentfill}%
\pgfsetlinewidth{0.803000pt}%
\definecolor{currentstroke}{rgb}{0.000000,0.000000,0.000000}%
\pgfsetstrokecolor{currentstroke}%
\pgfsetdash{}{0pt}%
\pgfsys@defobject{currentmarker}{\pgfqpoint{-0.048611in}{0.000000in}}{\pgfqpoint{-0.000000in}{0.000000in}}{%
\pgfpathmoveto{\pgfqpoint{-0.000000in}{0.000000in}}%
\pgfpathlineto{\pgfqpoint{-0.048611in}{0.000000in}}%
\pgfusepath{stroke,fill}%
}%
\begin{pgfscope}%
\pgfsys@transformshift{0.688581in}{1.933449in}%
\pgfsys@useobject{currentmarker}{}%
\end{pgfscope}%
\end{pgfscope}%
\begin{pgfscope}%
\definecolor{textcolor}{rgb}{0.000000,0.000000,0.000000}%
\pgfsetstrokecolor{textcolor}%
\pgfsetfillcolor{textcolor}%
\pgftext[x=0.344444in, y=1.885224in, left, base]{\color{textcolor}\rmfamily\fontsize{10.000000}{12.000000}\selectfont \(\displaystyle {50.0}\)}%
\end{pgfscope}%
\begin{pgfscope}%
\pgfpathrectangle{\pgfqpoint{0.688581in}{0.565123in}}{\pgfqpoint{5.372684in}{2.736651in}}%
\pgfusepath{clip}%
\pgfsetbuttcap%
\pgfsetroundjoin%
\pgfsetlinewidth{0.803000pt}%
\definecolor{currentstroke}{rgb}{0.690196,0.690196,0.690196}%
\pgfsetstrokecolor{currentstroke}%
\pgfsetdash{{0.800000pt}{1.320000pt}}{0.000000pt}%
\pgfpathmoveto{\pgfqpoint{0.688581in}{2.275530in}}%
\pgfpathlineto{\pgfqpoint{6.061265in}{2.275530in}}%
\pgfusepath{stroke}%
\end{pgfscope}%
\begin{pgfscope}%
\pgfsetbuttcap%
\pgfsetroundjoin%
\definecolor{currentfill}{rgb}{0.000000,0.000000,0.000000}%
\pgfsetfillcolor{currentfill}%
\pgfsetlinewidth{0.803000pt}%
\definecolor{currentstroke}{rgb}{0.000000,0.000000,0.000000}%
\pgfsetstrokecolor{currentstroke}%
\pgfsetdash{}{0pt}%
\pgfsys@defobject{currentmarker}{\pgfqpoint{-0.048611in}{0.000000in}}{\pgfqpoint{-0.000000in}{0.000000in}}{%
\pgfpathmoveto{\pgfqpoint{-0.000000in}{0.000000in}}%
\pgfpathlineto{\pgfqpoint{-0.048611in}{0.000000in}}%
\pgfusepath{stroke,fill}%
}%
\begin{pgfscope}%
\pgfsys@transformshift{0.688581in}{2.275530in}%
\pgfsys@useobject{currentmarker}{}%
\end{pgfscope}%
\end{pgfscope}%
\begin{pgfscope}%
\definecolor{textcolor}{rgb}{0.000000,0.000000,0.000000}%
\pgfsetstrokecolor{textcolor}%
\pgfsetfillcolor{textcolor}%
\pgftext[x=0.344444in, y=2.227305in, left, base]{\color{textcolor}\rmfamily\fontsize{10.000000}{12.000000}\selectfont \(\displaystyle {50.5}\)}%
\end{pgfscope}%
\begin{pgfscope}%
\pgfpathrectangle{\pgfqpoint{0.688581in}{0.565123in}}{\pgfqpoint{5.372684in}{2.736651in}}%
\pgfusepath{clip}%
\pgfsetbuttcap%
\pgfsetroundjoin%
\pgfsetlinewidth{0.803000pt}%
\definecolor{currentstroke}{rgb}{0.690196,0.690196,0.690196}%
\pgfsetstrokecolor{currentstroke}%
\pgfsetdash{{0.800000pt}{1.320000pt}}{0.000000pt}%
\pgfpathmoveto{\pgfqpoint{0.688581in}{2.617612in}}%
\pgfpathlineto{\pgfqpoint{6.061265in}{2.617612in}}%
\pgfusepath{stroke}%
\end{pgfscope}%
\begin{pgfscope}%
\pgfsetbuttcap%
\pgfsetroundjoin%
\definecolor{currentfill}{rgb}{0.000000,0.000000,0.000000}%
\pgfsetfillcolor{currentfill}%
\pgfsetlinewidth{0.803000pt}%
\definecolor{currentstroke}{rgb}{0.000000,0.000000,0.000000}%
\pgfsetstrokecolor{currentstroke}%
\pgfsetdash{}{0pt}%
\pgfsys@defobject{currentmarker}{\pgfqpoint{-0.048611in}{0.000000in}}{\pgfqpoint{-0.000000in}{0.000000in}}{%
\pgfpathmoveto{\pgfqpoint{-0.000000in}{0.000000in}}%
\pgfpathlineto{\pgfqpoint{-0.048611in}{0.000000in}}%
\pgfusepath{stroke,fill}%
}%
\begin{pgfscope}%
\pgfsys@transformshift{0.688581in}{2.617612in}%
\pgfsys@useobject{currentmarker}{}%
\end{pgfscope}%
\end{pgfscope}%
\begin{pgfscope}%
\definecolor{textcolor}{rgb}{0.000000,0.000000,0.000000}%
\pgfsetstrokecolor{textcolor}%
\pgfsetfillcolor{textcolor}%
\pgftext[x=0.344444in, y=2.569387in, left, base]{\color{textcolor}\rmfamily\fontsize{10.000000}{12.000000}\selectfont \(\displaystyle {51.0}\)}%
\end{pgfscope}%
\begin{pgfscope}%
\pgfpathrectangle{\pgfqpoint{0.688581in}{0.565123in}}{\pgfqpoint{5.372684in}{2.736651in}}%
\pgfusepath{clip}%
\pgfsetbuttcap%
\pgfsetroundjoin%
\pgfsetlinewidth{0.803000pt}%
\definecolor{currentstroke}{rgb}{0.690196,0.690196,0.690196}%
\pgfsetstrokecolor{currentstroke}%
\pgfsetdash{{0.800000pt}{1.320000pt}}{0.000000pt}%
\pgfpathmoveto{\pgfqpoint{0.688581in}{2.959693in}}%
\pgfpathlineto{\pgfqpoint{6.061265in}{2.959693in}}%
\pgfusepath{stroke}%
\end{pgfscope}%
\begin{pgfscope}%
\pgfsetbuttcap%
\pgfsetroundjoin%
\definecolor{currentfill}{rgb}{0.000000,0.000000,0.000000}%
\pgfsetfillcolor{currentfill}%
\pgfsetlinewidth{0.803000pt}%
\definecolor{currentstroke}{rgb}{0.000000,0.000000,0.000000}%
\pgfsetstrokecolor{currentstroke}%
\pgfsetdash{}{0pt}%
\pgfsys@defobject{currentmarker}{\pgfqpoint{-0.048611in}{0.000000in}}{\pgfqpoint{-0.000000in}{0.000000in}}{%
\pgfpathmoveto{\pgfqpoint{-0.000000in}{0.000000in}}%
\pgfpathlineto{\pgfqpoint{-0.048611in}{0.000000in}}%
\pgfusepath{stroke,fill}%
}%
\begin{pgfscope}%
\pgfsys@transformshift{0.688581in}{2.959693in}%
\pgfsys@useobject{currentmarker}{}%
\end{pgfscope}%
\end{pgfscope}%
\begin{pgfscope}%
\definecolor{textcolor}{rgb}{0.000000,0.000000,0.000000}%
\pgfsetstrokecolor{textcolor}%
\pgfsetfillcolor{textcolor}%
\pgftext[x=0.344444in, y=2.911468in, left, base]{\color{textcolor}\rmfamily\fontsize{10.000000}{12.000000}\selectfont \(\displaystyle {51.5}\)}%
\end{pgfscope}%
\begin{pgfscope}%
\pgfpathrectangle{\pgfqpoint{0.688581in}{0.565123in}}{\pgfqpoint{5.372684in}{2.736651in}}%
\pgfusepath{clip}%
\pgfsetbuttcap%
\pgfsetroundjoin%
\pgfsetlinewidth{0.803000pt}%
\definecolor{currentstroke}{rgb}{0.690196,0.690196,0.690196}%
\pgfsetstrokecolor{currentstroke}%
\pgfsetdash{{0.800000pt}{1.320000pt}}{0.000000pt}%
\pgfpathmoveto{\pgfqpoint{0.688581in}{3.301775in}}%
\pgfpathlineto{\pgfqpoint{6.061265in}{3.301775in}}%
\pgfusepath{stroke}%
\end{pgfscope}%
\begin{pgfscope}%
\pgfsetbuttcap%
\pgfsetroundjoin%
\definecolor{currentfill}{rgb}{0.000000,0.000000,0.000000}%
\pgfsetfillcolor{currentfill}%
\pgfsetlinewidth{0.803000pt}%
\definecolor{currentstroke}{rgb}{0.000000,0.000000,0.000000}%
\pgfsetstrokecolor{currentstroke}%
\pgfsetdash{}{0pt}%
\pgfsys@defobject{currentmarker}{\pgfqpoint{-0.048611in}{0.000000in}}{\pgfqpoint{-0.000000in}{0.000000in}}{%
\pgfpathmoveto{\pgfqpoint{-0.000000in}{0.000000in}}%
\pgfpathlineto{\pgfqpoint{-0.048611in}{0.000000in}}%
\pgfusepath{stroke,fill}%
}%
\begin{pgfscope}%
\pgfsys@transformshift{0.688581in}{3.301775in}%
\pgfsys@useobject{currentmarker}{}%
\end{pgfscope}%
\end{pgfscope}%
\begin{pgfscope}%
\definecolor{textcolor}{rgb}{0.000000,0.000000,0.000000}%
\pgfsetstrokecolor{textcolor}%
\pgfsetfillcolor{textcolor}%
\pgftext[x=0.344444in, y=3.253549in, left, base]{\color{textcolor}\rmfamily\fontsize{10.000000}{12.000000}\selectfont \(\displaystyle {52.0}\)}%
\end{pgfscope}%
\begin{pgfscope}%
\definecolor{textcolor}{rgb}{0.000000,0.000000,0.000000}%
\pgfsetstrokecolor{textcolor}%
\pgfsetfillcolor{textcolor}%
\pgftext[x=0.288889in,y=1.933449in,,bottom,rotate=90.000000]{\color{textcolor}\rmfamily\fontsize{10.000000}{12.000000}\selectfont Resistance / Impedance (\(\displaystyle \Omega\))}%
\end{pgfscope}%
\begin{pgfscope}%
\pgfpathrectangle{\pgfqpoint{0.688581in}{0.565123in}}{\pgfqpoint{5.372684in}{2.736651in}}%
\pgfusepath{clip}%
\pgfsetbuttcap%
\pgfsetroundjoin%
\pgfsetlinewidth{2.007500pt}%
\definecolor{currentstroke}{rgb}{0.121569,0.466667,0.705882}%
\pgfsetstrokecolor{currentstroke}%
\pgfsetdash{{7.400000pt}{3.200000pt}}{0.000000pt}%
\pgfpathmoveto{\pgfqpoint{0.678581in}{0.701956in}}%
\pgfpathlineto{\pgfqpoint{0.916907in}{0.701956in}}%
\pgfpathlineto{\pgfqpoint{0.952821in}{0.770372in}}%
\pgfpathlineto{\pgfqpoint{1.096481in}{0.770372in}}%
\pgfpathlineto{\pgfqpoint{1.132396in}{0.838788in}}%
\pgfpathlineto{\pgfqpoint{1.347885in}{0.838788in}}%
\pgfpathlineto{\pgfqpoint{1.383800in}{0.907205in}}%
\pgfpathlineto{\pgfqpoint{1.563375in}{0.907205in}}%
\pgfpathlineto{\pgfqpoint{1.599290in}{0.975621in}}%
\pgfpathlineto{\pgfqpoint{1.814779in}{0.975621in}}%
\pgfpathlineto{\pgfqpoint{1.850694in}{1.044037in}}%
\pgfpathlineto{\pgfqpoint{2.066184in}{1.044037in}}%
\pgfpathlineto{\pgfqpoint{2.102099in}{1.112454in}}%
\pgfpathlineto{\pgfqpoint{2.317588in}{1.112454in}}%
\pgfpathlineto{\pgfqpoint{2.353503in}{1.180870in}}%
\pgfpathlineto{\pgfqpoint{2.568992in}{1.180870in}}%
\pgfpathlineto{\pgfqpoint{2.604907in}{1.249286in}}%
\pgfpathlineto{\pgfqpoint{2.784482in}{1.249286in}}%
\pgfpathlineto{\pgfqpoint{2.820397in}{1.317702in}}%
\pgfpathlineto{\pgfqpoint{2.999971in}{1.317702in}}%
\pgfpathlineto{\pgfqpoint{3.035886in}{1.386119in}}%
\pgfpathlineto{\pgfqpoint{3.215461in}{1.386119in}}%
\pgfpathlineto{\pgfqpoint{3.251376in}{1.454535in}}%
\pgfpathlineto{\pgfqpoint{3.395035in}{1.454535in}}%
\pgfpathlineto{\pgfqpoint{3.430950in}{1.522951in}}%
\pgfpathlineto{\pgfqpoint{3.646440in}{1.522951in}}%
\pgfpathlineto{\pgfqpoint{3.682355in}{1.591368in}}%
\pgfpathlineto{\pgfqpoint{3.826014in}{1.591368in}}%
\pgfpathlineto{\pgfqpoint{3.861929in}{1.659784in}}%
\pgfpathlineto{\pgfqpoint{4.005589in}{1.659784in}}%
\pgfpathlineto{\pgfqpoint{4.041504in}{1.728200in}}%
\pgfpathlineto{\pgfqpoint{4.149249in}{1.728200in}}%
\pgfpathlineto{\pgfqpoint{4.185163in}{1.796616in}}%
\pgfpathlineto{\pgfqpoint{4.328823in}{1.796616in}}%
\pgfpathlineto{\pgfqpoint{4.364738in}{1.865033in}}%
\pgfpathlineto{\pgfqpoint{4.508398in}{1.865033in}}%
\pgfpathlineto{\pgfqpoint{4.544313in}{1.933449in}}%
\pgfpathlineto{\pgfqpoint{4.652057in}{1.933449in}}%
\pgfpathlineto{\pgfqpoint{4.687972in}{2.001865in}}%
\pgfpathlineto{\pgfqpoint{4.795717in}{2.001865in}}%
\pgfpathlineto{\pgfqpoint{4.831632in}{2.070282in}}%
\pgfpathlineto{\pgfqpoint{4.975292in}{2.070282in}}%
\pgfpathlineto{\pgfqpoint{5.011206in}{2.138698in}}%
\pgfpathlineto{\pgfqpoint{5.154866in}{2.138698in}}%
\pgfpathlineto{\pgfqpoint{5.190781in}{2.207114in}}%
\pgfpathlineto{\pgfqpoint{5.298526in}{2.207114in}}%
\pgfpathlineto{\pgfqpoint{5.334441in}{2.275530in}}%
\pgfpathlineto{\pgfqpoint{5.442185in}{2.275530in}}%
\pgfpathlineto{\pgfqpoint{5.478100in}{2.343947in}}%
\pgfpathlineto{\pgfqpoint{5.549930in}{2.343947in}}%
\pgfpathlineto{\pgfqpoint{5.585845in}{2.412363in}}%
\pgfpathlineto{\pgfqpoint{5.693590in}{2.412363in}}%
\pgfpathlineto{\pgfqpoint{5.729505in}{2.480779in}}%
\pgfpathlineto{\pgfqpoint{5.837249in}{2.480779in}}%
\pgfpathlineto{\pgfqpoint{5.873164in}{2.549196in}}%
\pgfpathlineto{\pgfqpoint{5.944994in}{2.549196in}}%
\pgfpathlineto{\pgfqpoint{5.980909in}{2.617612in}}%
\pgfpathlineto{\pgfqpoint{6.052739in}{2.617612in}}%
\pgfpathlineto{\pgfqpoint{6.071265in}{2.652904in}}%
\pgfpathlineto{\pgfqpoint{6.071265in}{2.652904in}}%
\pgfusepath{stroke}%
\end{pgfscope}%
\begin{pgfscope}%
\pgfpathrectangle{\pgfqpoint{0.688581in}{0.565123in}}{\pgfqpoint{5.372684in}{2.736651in}}%
\pgfusepath{clip}%
\pgfsetrectcap%
\pgfsetroundjoin%
\pgfsetlinewidth{2.007500pt}%
\definecolor{currentstroke}{rgb}{1.000000,0.498039,0.054902}%
\pgfsetstrokecolor{currentstroke}%
\pgfsetdash{}{0pt}%
\pgfpathmoveto{\pgfqpoint{0.678581in}{0.701956in}}%
\pgfpathlineto{\pgfqpoint{0.916907in}{0.701956in}}%
\pgfpathlineto{\pgfqpoint{0.952821in}{0.770372in}}%
\pgfpathlineto{\pgfqpoint{1.096481in}{0.770372in}}%
\pgfpathlineto{\pgfqpoint{1.132396in}{0.838788in}}%
\pgfpathlineto{\pgfqpoint{1.276056in}{0.838788in}}%
\pgfpathlineto{\pgfqpoint{1.311971in}{0.907205in}}%
\pgfpathlineto{\pgfqpoint{1.491545in}{0.907205in}}%
\pgfpathlineto{\pgfqpoint{1.527460in}{0.975621in}}%
\pgfpathlineto{\pgfqpoint{1.707035in}{0.975621in}}%
\pgfpathlineto{\pgfqpoint{1.742949in}{1.044037in}}%
\pgfpathlineto{\pgfqpoint{1.886609in}{1.044037in}}%
\pgfpathlineto{\pgfqpoint{1.922524in}{1.112454in}}%
\pgfpathlineto{\pgfqpoint{2.102099in}{1.112454in}}%
\pgfpathlineto{\pgfqpoint{2.138014in}{1.180870in}}%
\pgfpathlineto{\pgfqpoint{2.281673in}{1.180870in}}%
\pgfpathlineto{\pgfqpoint{2.317588in}{1.249286in}}%
\pgfpathlineto{\pgfqpoint{2.497163in}{1.249286in}}%
\pgfpathlineto{\pgfqpoint{2.533078in}{1.317702in}}%
\pgfpathlineto{\pgfqpoint{2.640822in}{1.317702in}}%
\pgfpathlineto{\pgfqpoint{2.676737in}{1.386119in}}%
\pgfpathlineto{\pgfqpoint{2.820397in}{1.386119in}}%
\pgfpathlineto{\pgfqpoint{2.856312in}{1.454535in}}%
\pgfpathlineto{\pgfqpoint{2.999971in}{1.454535in}}%
\pgfpathlineto{\pgfqpoint{3.035886in}{1.522951in}}%
\pgfpathlineto{\pgfqpoint{3.143631in}{1.522951in}}%
\pgfpathlineto{\pgfqpoint{3.179546in}{1.591368in}}%
\pgfpathlineto{\pgfqpoint{3.287291in}{1.591368in}}%
\pgfpathlineto{\pgfqpoint{3.323206in}{1.659784in}}%
\pgfpathlineto{\pgfqpoint{3.466865in}{1.659784in}}%
\pgfpathlineto{\pgfqpoint{3.502780in}{1.728200in}}%
\pgfpathlineto{\pgfqpoint{3.574610in}{1.728200in}}%
\pgfpathlineto{\pgfqpoint{3.610525in}{1.796616in}}%
\pgfpathlineto{\pgfqpoint{3.754185in}{1.796616in}}%
\pgfpathlineto{\pgfqpoint{3.790099in}{1.865033in}}%
\pgfpathlineto{\pgfqpoint{3.861929in}{1.865033in}}%
\pgfpathlineto{\pgfqpoint{3.897844in}{1.933449in}}%
\pgfpathlineto{\pgfqpoint{4.005589in}{1.933449in}}%
\pgfpathlineto{\pgfqpoint{4.041504in}{2.001865in}}%
\pgfpathlineto{\pgfqpoint{4.149249in}{2.001865in}}%
\pgfpathlineto{\pgfqpoint{4.185163in}{2.070282in}}%
\pgfpathlineto{\pgfqpoint{4.256993in}{2.070282in}}%
\pgfpathlineto{\pgfqpoint{4.292908in}{2.138698in}}%
\pgfpathlineto{\pgfqpoint{4.400653in}{2.138698in}}%
\pgfpathlineto{\pgfqpoint{4.436568in}{2.207114in}}%
\pgfpathlineto{\pgfqpoint{4.508398in}{2.207114in}}%
\pgfpathlineto{\pgfqpoint{4.544313in}{2.275530in}}%
\pgfpathlineto{\pgfqpoint{4.616142in}{2.275530in}}%
\pgfpathlineto{\pgfqpoint{4.652057in}{2.343947in}}%
\pgfpathlineto{\pgfqpoint{4.687972in}{2.343947in}}%
\pgfpathlineto{\pgfqpoint{4.723887in}{2.412363in}}%
\pgfpathlineto{\pgfqpoint{4.831632in}{2.412363in}}%
\pgfpathlineto{\pgfqpoint{4.867547in}{2.480779in}}%
\pgfpathlineto{\pgfqpoint{4.975292in}{2.480779in}}%
\pgfpathlineto{\pgfqpoint{5.011206in}{2.549196in}}%
\pgfpathlineto{\pgfqpoint{5.083036in}{2.549196in}}%
\pgfpathlineto{\pgfqpoint{5.118951in}{2.617612in}}%
\pgfpathlineto{\pgfqpoint{5.226696in}{2.617612in}}%
\pgfpathlineto{\pgfqpoint{5.262611in}{2.686028in}}%
\pgfpathlineto{\pgfqpoint{5.298526in}{2.686028in}}%
\pgfpathlineto{\pgfqpoint{5.334441in}{2.754444in}}%
\pgfpathlineto{\pgfqpoint{5.406270in}{2.754444in}}%
\pgfpathlineto{\pgfqpoint{5.442185in}{2.822861in}}%
\pgfpathlineto{\pgfqpoint{5.514015in}{2.822861in}}%
\pgfpathlineto{\pgfqpoint{5.549930in}{2.891277in}}%
\pgfpathlineto{\pgfqpoint{5.585845in}{2.891277in}}%
\pgfpathlineto{\pgfqpoint{5.621760in}{2.959693in}}%
\pgfpathlineto{\pgfqpoint{5.693590in}{2.959693in}}%
\pgfpathlineto{\pgfqpoint{5.729505in}{3.028110in}}%
\pgfpathlineto{\pgfqpoint{5.765420in}{3.028110in}}%
\pgfpathlineto{\pgfqpoint{5.801334in}{3.096526in}}%
\pgfpathlineto{\pgfqpoint{5.873164in}{3.096526in}}%
\pgfpathlineto{\pgfqpoint{5.909079in}{3.164942in}}%
\pgfpathlineto{\pgfqpoint{5.944994in}{3.164942in}}%
\pgfpathlineto{\pgfqpoint{5.980909in}{3.233358in}}%
\pgfpathlineto{\pgfqpoint{6.052739in}{3.233358in}}%
\pgfpathlineto{\pgfqpoint{6.071265in}{3.268650in}}%
\pgfpathlineto{\pgfqpoint{6.071265in}{3.268650in}}%
\pgfusepath{stroke}%
\end{pgfscope}%
\begin{pgfscope}%
\pgfpathrectangle{\pgfqpoint{0.688581in}{0.565123in}}{\pgfqpoint{5.372684in}{2.736651in}}%
\pgfusepath{clip}%
\pgfsetbuttcap%
\pgfsetroundjoin%
\pgfsetlinewidth{2.007500pt}%
\definecolor{currentstroke}{rgb}{0.172549,0.627451,0.172549}%
\pgfsetstrokecolor{currentstroke}%
\pgfsetdash{{7.400000pt}{3.200000pt}}{0.000000pt}%
\pgfpathmoveto{\pgfqpoint{0.678581in}{2.480779in}}%
\pgfpathlineto{\pgfqpoint{1.204226in}{2.480779in}}%
\pgfpathlineto{\pgfqpoint{1.240141in}{2.549196in}}%
\pgfpathlineto{\pgfqpoint{1.276056in}{2.480779in}}%
\pgfpathlineto{\pgfqpoint{1.311971in}{2.480779in}}%
\pgfpathlineto{\pgfqpoint{1.347885in}{2.549196in}}%
\pgfpathlineto{\pgfqpoint{1.383800in}{2.480779in}}%
\pgfpathlineto{\pgfqpoint{1.419715in}{2.480779in}}%
\pgfpathlineto{\pgfqpoint{1.455630in}{2.549196in}}%
\pgfpathlineto{\pgfqpoint{1.491545in}{2.480779in}}%
\pgfpathlineto{\pgfqpoint{1.527460in}{2.549196in}}%
\pgfpathlineto{\pgfqpoint{1.563375in}{2.549196in}}%
\pgfpathlineto{\pgfqpoint{1.599290in}{2.480779in}}%
\pgfpathlineto{\pgfqpoint{1.671120in}{2.480779in}}%
\pgfpathlineto{\pgfqpoint{1.707035in}{2.549196in}}%
\pgfpathlineto{\pgfqpoint{1.742949in}{2.549196in}}%
\pgfpathlineto{\pgfqpoint{1.778864in}{2.480779in}}%
\pgfpathlineto{\pgfqpoint{1.814779in}{2.549196in}}%
\pgfpathlineto{\pgfqpoint{1.850694in}{2.480779in}}%
\pgfpathlineto{\pgfqpoint{1.886609in}{2.480779in}}%
\pgfpathlineto{\pgfqpoint{1.922524in}{2.549196in}}%
\pgfpathlineto{\pgfqpoint{2.389418in}{2.549196in}}%
\pgfpathlineto{\pgfqpoint{2.425333in}{2.480779in}}%
\pgfpathlineto{\pgfqpoint{2.461248in}{2.549196in}}%
\pgfpathlineto{\pgfqpoint{2.604907in}{2.549196in}}%
\pgfpathlineto{\pgfqpoint{2.640822in}{2.480779in}}%
\pgfpathlineto{\pgfqpoint{2.676737in}{2.480779in}}%
\pgfpathlineto{\pgfqpoint{2.712652in}{2.549196in}}%
\pgfpathlineto{\pgfqpoint{2.892227in}{2.549196in}}%
\pgfpathlineto{\pgfqpoint{2.928142in}{2.480779in}}%
\pgfpathlineto{\pgfqpoint{2.964056in}{2.549196in}}%
\pgfpathlineto{\pgfqpoint{2.999971in}{2.480779in}}%
\pgfpathlineto{\pgfqpoint{3.071801in}{2.480779in}}%
\pgfpathlineto{\pgfqpoint{3.107716in}{2.549196in}}%
\pgfpathlineto{\pgfqpoint{3.215461in}{2.549196in}}%
\pgfpathlineto{\pgfqpoint{3.251376in}{2.480779in}}%
\pgfpathlineto{\pgfqpoint{3.287291in}{2.549196in}}%
\pgfpathlineto{\pgfqpoint{3.395035in}{2.549196in}}%
\pgfpathlineto{\pgfqpoint{3.430950in}{2.480779in}}%
\pgfpathlineto{\pgfqpoint{3.466865in}{2.549196in}}%
\pgfpathlineto{\pgfqpoint{3.826014in}{2.549196in}}%
\pgfpathlineto{\pgfqpoint{3.861929in}{2.480779in}}%
\pgfpathlineto{\pgfqpoint{3.897844in}{2.549196in}}%
\pgfpathlineto{\pgfqpoint{3.933759in}{2.480779in}}%
\pgfpathlineto{\pgfqpoint{4.005589in}{2.480779in}}%
\pgfpathlineto{\pgfqpoint{4.041504in}{2.549196in}}%
\pgfpathlineto{\pgfqpoint{4.256993in}{2.549196in}}%
\pgfpathlineto{\pgfqpoint{4.292908in}{2.480779in}}%
\pgfpathlineto{\pgfqpoint{4.328823in}{2.480779in}}%
\pgfpathlineto{\pgfqpoint{4.364738in}{2.549196in}}%
\pgfpathlineto{\pgfqpoint{4.400653in}{2.480779in}}%
\pgfpathlineto{\pgfqpoint{4.508398in}{2.480779in}}%
\pgfpathlineto{\pgfqpoint{4.544313in}{2.549196in}}%
\pgfpathlineto{\pgfqpoint{4.867547in}{2.549196in}}%
\pgfpathlineto{\pgfqpoint{4.903462in}{2.480779in}}%
\pgfpathlineto{\pgfqpoint{4.939377in}{2.480779in}}%
\pgfpathlineto{\pgfqpoint{4.975292in}{2.549196in}}%
\pgfpathlineto{\pgfqpoint{5.190781in}{2.549196in}}%
\pgfpathlineto{\pgfqpoint{5.226696in}{2.480779in}}%
\pgfpathlineto{\pgfqpoint{5.262611in}{2.549196in}}%
\pgfpathlineto{\pgfqpoint{5.298526in}{2.549196in}}%
\pgfpathlineto{\pgfqpoint{5.334441in}{2.480779in}}%
\pgfpathlineto{\pgfqpoint{5.406270in}{2.480779in}}%
\pgfpathlineto{\pgfqpoint{5.442185in}{2.549196in}}%
\pgfpathlineto{\pgfqpoint{5.478100in}{2.549196in}}%
\pgfpathlineto{\pgfqpoint{5.514015in}{2.480779in}}%
\pgfpathlineto{\pgfqpoint{5.729505in}{2.480779in}}%
\pgfpathlineto{\pgfqpoint{5.765420in}{2.549196in}}%
\pgfpathlineto{\pgfqpoint{5.873164in}{2.549196in}}%
\pgfpathlineto{\pgfqpoint{5.909079in}{2.480779in}}%
\pgfpathlineto{\pgfqpoint{5.980909in}{2.480779in}}%
\pgfpathlineto{\pgfqpoint{6.016824in}{2.549196in}}%
\pgfpathlineto{\pgfqpoint{6.071265in}{2.549196in}}%
\pgfpathlineto{\pgfqpoint{6.071265in}{2.549196in}}%
\pgfusepath{stroke}%
\end{pgfscope}%
\begin{pgfscope}%
\pgfpathrectangle{\pgfqpoint{0.688581in}{0.565123in}}{\pgfqpoint{5.372684in}{2.736651in}}%
\pgfusepath{clip}%
\pgfsetrectcap%
\pgfsetroundjoin%
\pgfsetlinewidth{2.007500pt}%
\definecolor{currentstroke}{rgb}{0.839216,0.152941,0.156863}%
\pgfsetstrokecolor{currentstroke}%
\pgfsetdash{}{0pt}%
\pgfpathmoveto{\pgfqpoint{0.678581in}{2.480779in}}%
\pgfpathlineto{\pgfqpoint{1.024651in}{2.480779in}}%
\pgfpathlineto{\pgfqpoint{1.060566in}{2.549196in}}%
\pgfpathlineto{\pgfqpoint{1.096481in}{2.549196in}}%
\pgfpathlineto{\pgfqpoint{1.132396in}{2.480779in}}%
\pgfpathlineto{\pgfqpoint{1.168311in}{2.549196in}}%
\pgfpathlineto{\pgfqpoint{1.563375in}{2.549196in}}%
\pgfpathlineto{\pgfqpoint{1.599290in}{2.480779in}}%
\pgfpathlineto{\pgfqpoint{1.635205in}{2.480779in}}%
\pgfpathlineto{\pgfqpoint{1.671120in}{2.549196in}}%
\pgfpathlineto{\pgfqpoint{2.640822in}{2.549196in}}%
\pgfpathlineto{\pgfqpoint{2.676737in}{2.480779in}}%
\pgfpathlineto{\pgfqpoint{2.712652in}{2.549196in}}%
\pgfpathlineto{\pgfqpoint{2.999971in}{2.549196in}}%
\pgfpathlineto{\pgfqpoint{3.035886in}{2.480779in}}%
\pgfpathlineto{\pgfqpoint{3.071801in}{2.480779in}}%
\pgfpathlineto{\pgfqpoint{3.107716in}{2.549196in}}%
\pgfpathlineto{\pgfqpoint{3.826014in}{2.549196in}}%
\pgfpathlineto{\pgfqpoint{3.861929in}{2.480779in}}%
\pgfpathlineto{\pgfqpoint{3.897844in}{2.549196in}}%
\pgfpathlineto{\pgfqpoint{3.933759in}{2.480779in}}%
\pgfpathlineto{\pgfqpoint{4.005589in}{2.480779in}}%
\pgfpathlineto{\pgfqpoint{4.041504in}{2.549196in}}%
\pgfpathlineto{\pgfqpoint{4.364738in}{2.549196in}}%
\pgfpathlineto{\pgfqpoint{4.400653in}{2.480779in}}%
\pgfpathlineto{\pgfqpoint{4.436568in}{2.549196in}}%
\pgfpathlineto{\pgfqpoint{5.190781in}{2.549196in}}%
\pgfpathlineto{\pgfqpoint{5.226696in}{2.480779in}}%
\pgfpathlineto{\pgfqpoint{5.262611in}{2.549196in}}%
\pgfpathlineto{\pgfqpoint{5.334441in}{2.549196in}}%
\pgfpathlineto{\pgfqpoint{5.370356in}{2.480779in}}%
\pgfpathlineto{\pgfqpoint{5.406270in}{2.480779in}}%
\pgfpathlineto{\pgfqpoint{5.442185in}{2.549196in}}%
\pgfpathlineto{\pgfqpoint{5.478100in}{2.549196in}}%
\pgfpathlineto{\pgfqpoint{5.514015in}{2.480779in}}%
\pgfpathlineto{\pgfqpoint{5.549930in}{2.480779in}}%
\pgfpathlineto{\pgfqpoint{5.585845in}{2.549196in}}%
\pgfpathlineto{\pgfqpoint{5.621760in}{2.549196in}}%
\pgfpathlineto{\pgfqpoint{5.657675in}{2.480779in}}%
\pgfpathlineto{\pgfqpoint{5.693590in}{2.549196in}}%
\pgfpathlineto{\pgfqpoint{5.729505in}{2.480779in}}%
\pgfpathlineto{\pgfqpoint{5.765420in}{2.549196in}}%
\pgfpathlineto{\pgfqpoint{5.909079in}{2.549196in}}%
\pgfpathlineto{\pgfqpoint{5.944994in}{2.480779in}}%
\pgfpathlineto{\pgfqpoint{5.980909in}{2.549196in}}%
\pgfpathlineto{\pgfqpoint{6.071265in}{2.549196in}}%
\pgfpathlineto{\pgfqpoint{6.071265in}{2.549196in}}%
\pgfusepath{stroke}%
\end{pgfscope}%
\begin{pgfscope}%
\pgfsetrectcap%
\pgfsetmiterjoin%
\pgfsetlinewidth{0.803000pt}%
\definecolor{currentstroke}{rgb}{0.000000,0.000000,0.000000}%
\pgfsetstrokecolor{currentstroke}%
\pgfsetdash{}{0pt}%
\pgfpathmoveto{\pgfqpoint{0.688581in}{0.565123in}}%
\pgfpathlineto{\pgfqpoint{0.688581in}{3.301775in}}%
\pgfusepath{stroke}%
\end{pgfscope}%
\begin{pgfscope}%
\pgfsetrectcap%
\pgfsetmiterjoin%
\pgfsetlinewidth{0.803000pt}%
\definecolor{currentstroke}{rgb}{0.000000,0.000000,0.000000}%
\pgfsetstrokecolor{currentstroke}%
\pgfsetdash{}{0pt}%
\pgfpathmoveto{\pgfqpoint{6.061265in}{0.565123in}}%
\pgfpathlineto{\pgfqpoint{6.061265in}{3.301775in}}%
\pgfusepath{stroke}%
\end{pgfscope}%
\begin{pgfscope}%
\pgfsetrectcap%
\pgfsetmiterjoin%
\pgfsetlinewidth{0.803000pt}%
\definecolor{currentstroke}{rgb}{0.000000,0.000000,0.000000}%
\pgfsetstrokecolor{currentstroke}%
\pgfsetdash{}{0pt}%
\pgfpathmoveto{\pgfqpoint{0.688581in}{0.565123in}}%
\pgfpathlineto{\pgfqpoint{6.061265in}{0.565123in}}%
\pgfusepath{stroke}%
\end{pgfscope}%
\begin{pgfscope}%
\pgfsetrectcap%
\pgfsetmiterjoin%
\pgfsetlinewidth{0.803000pt}%
\definecolor{currentstroke}{rgb}{0.000000,0.000000,0.000000}%
\pgfsetstrokecolor{currentstroke}%
\pgfsetdash{}{0pt}%
\pgfpathmoveto{\pgfqpoint{0.688581in}{3.301775in}}%
\pgfpathlineto{\pgfqpoint{6.061265in}{3.301775in}}%
\pgfusepath{stroke}%
\end{pgfscope}%
\begin{pgfscope}%
\pgfsetbuttcap%
\pgfsetmiterjoin%
\definecolor{currentfill}{rgb}{1.000000,1.000000,1.000000}%
\pgfsetfillcolor{currentfill}%
\pgfsetfillopacity{0.800000}%
\pgfsetlinewidth{1.003750pt}%
\definecolor{currentstroke}{rgb}{0.800000,0.800000,0.800000}%
\pgfsetstrokecolor{currentstroke}%
\pgfsetstrokeopacity{0.800000}%
\pgfsetdash{}{0pt}%
\pgfpathmoveto{\pgfqpoint{3.779240in}{0.634568in}}%
\pgfpathlineto{\pgfqpoint{5.964043in}{0.634568in}}%
\pgfpathquadraticcurveto{\pgfqpoint{5.991821in}{0.634568in}}{\pgfqpoint{5.991821in}{0.662346in}}%
\pgfpathlineto{\pgfqpoint{5.991821in}{1.423148in}}%
\pgfpathquadraticcurveto{\pgfqpoint{5.991821in}{1.450926in}}{\pgfqpoint{5.964043in}{1.450926in}}%
\pgfpathlineto{\pgfqpoint{3.779240in}{1.450926in}}%
\pgfpathquadraticcurveto{\pgfqpoint{3.751462in}{1.450926in}}{\pgfqpoint{3.751462in}{1.423148in}}%
\pgfpathlineto{\pgfqpoint{3.751462in}{0.662346in}}%
\pgfpathquadraticcurveto{\pgfqpoint{3.751462in}{0.634568in}}{\pgfqpoint{3.779240in}{0.634568in}}%
\pgfpathlineto{\pgfqpoint{3.779240in}{0.634568in}}%
\pgfpathclose%
\pgfusepath{stroke,fill}%
\end{pgfscope}%
\begin{pgfscope}%
\pgfsetbuttcap%
\pgfsetroundjoin%
\pgfsetlinewidth{2.007500pt}%
\definecolor{currentstroke}{rgb}{0.121569,0.466667,0.705882}%
\pgfsetstrokecolor{currentstroke}%
\pgfsetdash{{7.400000pt}{3.200000pt}}{0.000000pt}%
\pgfpathmoveto{\pgfqpoint{3.807018in}{1.346759in}}%
\pgfpathlineto{\pgfqpoint{3.945907in}{1.346759in}}%
\pgfpathlineto{\pgfqpoint{4.084795in}{1.346759in}}%
\pgfusepath{stroke}%
\end{pgfscope}%
\begin{pgfscope}%
\definecolor{textcolor}{rgb}{0.000000,0.000000,0.000000}%
\pgfsetstrokecolor{textcolor}%
\pgfsetfillcolor{textcolor}%
\pgftext[x=4.195907in,y=1.298148in,left,base]{\color{textcolor}\rmfamily\fontsize{10.000000}{12.000000}\selectfont Rs – Carbon resistor}%
\end{pgfscope}%
\begin{pgfscope}%
\pgfsetrectcap%
\pgfsetroundjoin%
\pgfsetlinewidth{2.007500pt}%
\definecolor{currentstroke}{rgb}{1.000000,0.498039,0.054902}%
\pgfsetstrokecolor{currentstroke}%
\pgfsetdash{}{0pt}%
\pgfpathmoveto{\pgfqpoint{3.807018in}{1.153086in}}%
\pgfpathlineto{\pgfqpoint{3.945907in}{1.153086in}}%
\pgfpathlineto{\pgfqpoint{4.084795in}{1.153086in}}%
\pgfusepath{stroke}%
\end{pgfscope}%
\begin{pgfscope}%
\definecolor{textcolor}{rgb}{0.000000,0.000000,0.000000}%
\pgfsetstrokecolor{textcolor}%
\pgfsetfillcolor{textcolor}%
\pgftext[x=4.195907in,y=1.104475in,left,base]{\color{textcolor}\rmfamily\fontsize{10.000000}{12.000000}\selectfont |Z| – Carbon resistor}%
\end{pgfscope}%
\begin{pgfscope}%
\pgfsetbuttcap%
\pgfsetroundjoin%
\pgfsetlinewidth{2.007500pt}%
\definecolor{currentstroke}{rgb}{0.172549,0.627451,0.172549}%
\pgfsetstrokecolor{currentstroke}%
\pgfsetdash{{7.400000pt}{3.200000pt}}{0.000000pt}%
\pgfpathmoveto{\pgfqpoint{3.807018in}{0.959413in}}%
\pgfpathlineto{\pgfqpoint{3.945907in}{0.959413in}}%
\pgfpathlineto{\pgfqpoint{4.084795in}{0.959413in}}%
\pgfusepath{stroke}%
\end{pgfscope}%
\begin{pgfscope}%
\definecolor{textcolor}{rgb}{0.000000,0.000000,0.000000}%
\pgfsetstrokecolor{textcolor}%
\pgfsetfillcolor{textcolor}%
\pgftext[x=4.195907in,y=0.910802in,left,base]{\color{textcolor}\rmfamily\fontsize{10.000000}{12.000000}\selectfont Rs – Coaxial termination}%
\end{pgfscope}%
\begin{pgfscope}%
\pgfsetrectcap%
\pgfsetroundjoin%
\pgfsetlinewidth{2.007500pt}%
\definecolor{currentstroke}{rgb}{0.839216,0.152941,0.156863}%
\pgfsetstrokecolor{currentstroke}%
\pgfsetdash{}{0pt}%
\pgfpathmoveto{\pgfqpoint{3.807018in}{0.765741in}}%
\pgfpathlineto{\pgfqpoint{3.945907in}{0.765741in}}%
\pgfpathlineto{\pgfqpoint{4.084795in}{0.765741in}}%
\pgfusepath{stroke}%
\end{pgfscope}%
\begin{pgfscope}%
\definecolor{textcolor}{rgb}{0.000000,0.000000,0.000000}%
\pgfsetstrokecolor{textcolor}%
\pgfsetfillcolor{textcolor}%
\pgftext[x=4.195907in,y=0.717129in,left,base]{\color{textcolor}\rmfamily\fontsize{10.000000}{12.000000}\selectfont |Z| – Coaxial termination}%
\end{pgfscope}%
\end{pgfpicture}%
\makeatother%
\endgroup%

		\caption{Comparison of resistance and impedance of carbon resistor and coaxial termination}
		\label{fig:resistors}
	\end{figure}


	\subsubsection{}
	
	Figure \ref{fig:352} shows the measured impedance magnitude $|Z|$ and phase of the inductor as a
	function of frequency.
	At $f = 3.26\,\mathrm{MHz}$, the inductance is measured at $\sim50 \Omega$. The theoretical value is obtained from
	$|Z| \approx X_L = 2\pi f L \approx 45 \Omega$ with L=2.2 $\mu$H and is close to the nominal value of the inductor.
	Small deviations are caused by series resistance and measurement uncertainty.
	The impedance reaches $|Z| = 50\,\Omega$ at low frequency where the phase is still
	strongly positive ($+88.86^\circ$), confirming inductive behavior.
	A pronounced maximum of $|Z|$ occurs at approximately $30\,\mathrm{MHz}$, indicating
	the self-resonant frequency of the inductor due to parasitic capacitance.
	Therefore, the maximum operating frequency of the inductor is below resonance.
	
	\begin{figure}[H]
		\centering
		%% Creator: Matplotlib, PGF backend
%%
%% To include the figure in your LaTeX document, write
%%   \input{<filename>.pgf}
%%
%% Make sure the required packages are loaded in your preamble
%%   \usepackage{pgf}
%%
%% Also ensure that all the required font packages are loaded; for instance,
%% the lmodern package is sometimes necessary when using math font.
%%   \usepackage{lmodern}
%%
%% Figures using additional raster images can only be included by \input if
%% they are in the same directory as the main LaTeX file. For loading figures
%% from other directories you can use the `import` package
%%   \usepackage{import}
%%
%% and then include the figures with
%%   \import{<path to file>}{<filename>.pgf}
%%
%% Matplotlib used the following preamble
%%   \usepackage{fontspec}
%%   \setmainfont{DejaVuSerif.ttf}[Path=\detokenize{C:/Users/wenze/AppData/Local/Programs/Python/Python310/Lib/site-packages/matplotlib/mpl-data/fonts/ttf/}]
%%   \setsansfont{DejaVuSans.ttf}[Path=\detokenize{C:/Users/wenze/AppData/Local/Programs/Python/Python310/Lib/site-packages/matplotlib/mpl-data/fonts/ttf/}]
%%   \setmonofont{DejaVuSansMono.ttf}[Path=\detokenize{C:/Users/wenze/AppData/Local/Programs/Python/Python310/Lib/site-packages/matplotlib/mpl-data/fonts/ttf/}]
%%
\begingroup%
\makeatletter%
\begin{pgfpicture}%
\pgfpathrectangle{\pgfpointorigin}{\pgfqpoint{6.300000in}{3.500000in}}%
\pgfusepath{use as bounding box, clip}%
\begin{pgfscope}%
\pgfsetbuttcap%
\pgfsetmiterjoin%
\definecolor{currentfill}{rgb}{1.000000,1.000000,1.000000}%
\pgfsetfillcolor{currentfill}%
\pgfsetlinewidth{0.000000pt}%
\definecolor{currentstroke}{rgb}{1.000000,1.000000,1.000000}%
\pgfsetstrokecolor{currentstroke}%
\pgfsetdash{}{0pt}%
\pgfpathmoveto{\pgfqpoint{0.000000in}{0.000000in}}%
\pgfpathlineto{\pgfqpoint{6.300000in}{0.000000in}}%
\pgfpathlineto{\pgfqpoint{6.300000in}{3.500000in}}%
\pgfpathlineto{\pgfqpoint{0.000000in}{3.500000in}}%
\pgfpathlineto{\pgfqpoint{0.000000in}{0.000000in}}%
\pgfpathclose%
\pgfusepath{fill}%
\end{pgfscope}%
\begin{pgfscope}%
\pgfsetbuttcap%
\pgfsetmiterjoin%
\definecolor{currentfill}{rgb}{1.000000,1.000000,1.000000}%
\pgfsetfillcolor{currentfill}%
\pgfsetlinewidth{0.000000pt}%
\definecolor{currentstroke}{rgb}{0.000000,0.000000,0.000000}%
\pgfsetstrokecolor{currentstroke}%
\pgfsetstrokeopacity{0.000000}%
\pgfsetdash{}{0pt}%
\pgfpathmoveto{\pgfqpoint{0.834028in}{0.582778in}}%
\pgfpathlineto{\pgfqpoint{5.593472in}{0.582778in}}%
\pgfpathlineto{\pgfqpoint{5.593472in}{3.020000in}}%
\pgfpathlineto{\pgfqpoint{0.834028in}{3.020000in}}%
\pgfpathlineto{\pgfqpoint{0.834028in}{0.582778in}}%
\pgfpathclose%
\pgfusepath{fill}%
\end{pgfscope}%
\begin{pgfscope}%
\pgfpathrectangle{\pgfqpoint{0.834028in}{0.582778in}}{\pgfqpoint{4.759444in}{2.437222in}}%
\pgfusepath{clip}%
\pgfsetbuttcap%
\pgfsetroundjoin%
\pgfsetlinewidth{0.803000pt}%
\definecolor{currentstroke}{rgb}{0.690196,0.690196,0.690196}%
\pgfsetstrokecolor{currentstroke}%
\pgfsetdash{{0.800000pt}{1.320000pt}}{0.000000pt}%
\pgfpathmoveto{\pgfqpoint{1.360503in}{0.582778in}}%
\pgfpathlineto{\pgfqpoint{1.360503in}{3.020000in}}%
\pgfusepath{stroke}%
\end{pgfscope}%
\begin{pgfscope}%
\pgfsetbuttcap%
\pgfsetroundjoin%
\definecolor{currentfill}{rgb}{0.000000,0.000000,0.000000}%
\pgfsetfillcolor{currentfill}%
\pgfsetlinewidth{0.803000pt}%
\definecolor{currentstroke}{rgb}{0.000000,0.000000,0.000000}%
\pgfsetstrokecolor{currentstroke}%
\pgfsetdash{}{0pt}%
\pgfsys@defobject{currentmarker}{\pgfqpoint{0.000000in}{-0.048611in}}{\pgfqpoint{0.000000in}{0.000000in}}{%
\pgfpathmoveto{\pgfqpoint{0.000000in}{0.000000in}}%
\pgfpathlineto{\pgfqpoint{0.000000in}{-0.048611in}}%
\pgfusepath{stroke,fill}%
}%
\begin{pgfscope}%
\pgfsys@transformshift{1.360503in}{0.582778in}%
\pgfsys@useobject{currentmarker}{}%
\end{pgfscope}%
\end{pgfscope}%
\begin{pgfscope}%
\definecolor{textcolor}{rgb}{0.000000,0.000000,0.000000}%
\pgfsetstrokecolor{textcolor}%
\pgfsetfillcolor{textcolor}%
\pgftext[x=1.360503in,y=0.485556in,,top]{\color{textcolor}\sffamily\fontsize{10.000000}{12.000000}\selectfont 0.2}%
\end{pgfscope}%
\begin{pgfscope}%
\pgfpathrectangle{\pgfqpoint{0.834028in}{0.582778in}}{\pgfqpoint{4.759444in}{2.437222in}}%
\pgfusepath{clip}%
\pgfsetbuttcap%
\pgfsetroundjoin%
\pgfsetlinewidth{0.803000pt}%
\definecolor{currentstroke}{rgb}{0.690196,0.690196,0.690196}%
\pgfsetstrokecolor{currentstroke}%
\pgfsetdash{{0.800000pt}{1.320000pt}}{0.000000pt}%
\pgfpathmoveto{\pgfqpoint{1.889624in}{0.582778in}}%
\pgfpathlineto{\pgfqpoint{1.889624in}{3.020000in}}%
\pgfusepath{stroke}%
\end{pgfscope}%
\begin{pgfscope}%
\pgfsetbuttcap%
\pgfsetroundjoin%
\definecolor{currentfill}{rgb}{0.000000,0.000000,0.000000}%
\pgfsetfillcolor{currentfill}%
\pgfsetlinewidth{0.803000pt}%
\definecolor{currentstroke}{rgb}{0.000000,0.000000,0.000000}%
\pgfsetstrokecolor{currentstroke}%
\pgfsetdash{}{0pt}%
\pgfsys@defobject{currentmarker}{\pgfqpoint{0.000000in}{-0.048611in}}{\pgfqpoint{0.000000in}{0.000000in}}{%
\pgfpathmoveto{\pgfqpoint{0.000000in}{0.000000in}}%
\pgfpathlineto{\pgfqpoint{0.000000in}{-0.048611in}}%
\pgfusepath{stroke,fill}%
}%
\begin{pgfscope}%
\pgfsys@transformshift{1.889624in}{0.582778in}%
\pgfsys@useobject{currentmarker}{}%
\end{pgfscope}%
\end{pgfscope}%
\begin{pgfscope}%
\definecolor{textcolor}{rgb}{0.000000,0.000000,0.000000}%
\pgfsetstrokecolor{textcolor}%
\pgfsetfillcolor{textcolor}%
\pgftext[x=1.889624in,y=0.485556in,,top]{\color{textcolor}\sffamily\fontsize{10.000000}{12.000000}\selectfont 0.4}%
\end{pgfscope}%
\begin{pgfscope}%
\pgfpathrectangle{\pgfqpoint{0.834028in}{0.582778in}}{\pgfqpoint{4.759444in}{2.437222in}}%
\pgfusepath{clip}%
\pgfsetbuttcap%
\pgfsetroundjoin%
\pgfsetlinewidth{0.803000pt}%
\definecolor{currentstroke}{rgb}{0.690196,0.690196,0.690196}%
\pgfsetstrokecolor{currentstroke}%
\pgfsetdash{{0.800000pt}{1.320000pt}}{0.000000pt}%
\pgfpathmoveto{\pgfqpoint{2.418746in}{0.582778in}}%
\pgfpathlineto{\pgfqpoint{2.418746in}{3.020000in}}%
\pgfusepath{stroke}%
\end{pgfscope}%
\begin{pgfscope}%
\pgfsetbuttcap%
\pgfsetroundjoin%
\definecolor{currentfill}{rgb}{0.000000,0.000000,0.000000}%
\pgfsetfillcolor{currentfill}%
\pgfsetlinewidth{0.803000pt}%
\definecolor{currentstroke}{rgb}{0.000000,0.000000,0.000000}%
\pgfsetstrokecolor{currentstroke}%
\pgfsetdash{}{0pt}%
\pgfsys@defobject{currentmarker}{\pgfqpoint{0.000000in}{-0.048611in}}{\pgfqpoint{0.000000in}{0.000000in}}{%
\pgfpathmoveto{\pgfqpoint{0.000000in}{0.000000in}}%
\pgfpathlineto{\pgfqpoint{0.000000in}{-0.048611in}}%
\pgfusepath{stroke,fill}%
}%
\begin{pgfscope}%
\pgfsys@transformshift{2.418746in}{0.582778in}%
\pgfsys@useobject{currentmarker}{}%
\end{pgfscope}%
\end{pgfscope}%
\begin{pgfscope}%
\definecolor{textcolor}{rgb}{0.000000,0.000000,0.000000}%
\pgfsetstrokecolor{textcolor}%
\pgfsetfillcolor{textcolor}%
\pgftext[x=2.418746in,y=0.485556in,,top]{\color{textcolor}\sffamily\fontsize{10.000000}{12.000000}\selectfont 0.6}%
\end{pgfscope}%
\begin{pgfscope}%
\pgfpathrectangle{\pgfqpoint{0.834028in}{0.582778in}}{\pgfqpoint{4.759444in}{2.437222in}}%
\pgfusepath{clip}%
\pgfsetbuttcap%
\pgfsetroundjoin%
\pgfsetlinewidth{0.803000pt}%
\definecolor{currentstroke}{rgb}{0.690196,0.690196,0.690196}%
\pgfsetstrokecolor{currentstroke}%
\pgfsetdash{{0.800000pt}{1.320000pt}}{0.000000pt}%
\pgfpathmoveto{\pgfqpoint{2.947867in}{0.582778in}}%
\pgfpathlineto{\pgfqpoint{2.947867in}{3.020000in}}%
\pgfusepath{stroke}%
\end{pgfscope}%
\begin{pgfscope}%
\pgfsetbuttcap%
\pgfsetroundjoin%
\definecolor{currentfill}{rgb}{0.000000,0.000000,0.000000}%
\pgfsetfillcolor{currentfill}%
\pgfsetlinewidth{0.803000pt}%
\definecolor{currentstroke}{rgb}{0.000000,0.000000,0.000000}%
\pgfsetstrokecolor{currentstroke}%
\pgfsetdash{}{0pt}%
\pgfsys@defobject{currentmarker}{\pgfqpoint{0.000000in}{-0.048611in}}{\pgfqpoint{0.000000in}{0.000000in}}{%
\pgfpathmoveto{\pgfqpoint{0.000000in}{0.000000in}}%
\pgfpathlineto{\pgfqpoint{0.000000in}{-0.048611in}}%
\pgfusepath{stroke,fill}%
}%
\begin{pgfscope}%
\pgfsys@transformshift{2.947867in}{0.582778in}%
\pgfsys@useobject{currentmarker}{}%
\end{pgfscope}%
\end{pgfscope}%
\begin{pgfscope}%
\definecolor{textcolor}{rgb}{0.000000,0.000000,0.000000}%
\pgfsetstrokecolor{textcolor}%
\pgfsetfillcolor{textcolor}%
\pgftext[x=2.947867in,y=0.485556in,,top]{\color{textcolor}\sffamily\fontsize{10.000000}{12.000000}\selectfont 0.8}%
\end{pgfscope}%
\begin{pgfscope}%
\pgfpathrectangle{\pgfqpoint{0.834028in}{0.582778in}}{\pgfqpoint{4.759444in}{2.437222in}}%
\pgfusepath{clip}%
\pgfsetbuttcap%
\pgfsetroundjoin%
\pgfsetlinewidth{0.803000pt}%
\definecolor{currentstroke}{rgb}{0.690196,0.690196,0.690196}%
\pgfsetstrokecolor{currentstroke}%
\pgfsetdash{{0.800000pt}{1.320000pt}}{0.000000pt}%
\pgfpathmoveto{\pgfqpoint{3.476988in}{0.582778in}}%
\pgfpathlineto{\pgfqpoint{3.476988in}{3.020000in}}%
\pgfusepath{stroke}%
\end{pgfscope}%
\begin{pgfscope}%
\pgfsetbuttcap%
\pgfsetroundjoin%
\definecolor{currentfill}{rgb}{0.000000,0.000000,0.000000}%
\pgfsetfillcolor{currentfill}%
\pgfsetlinewidth{0.803000pt}%
\definecolor{currentstroke}{rgb}{0.000000,0.000000,0.000000}%
\pgfsetstrokecolor{currentstroke}%
\pgfsetdash{}{0pt}%
\pgfsys@defobject{currentmarker}{\pgfqpoint{0.000000in}{-0.048611in}}{\pgfqpoint{0.000000in}{0.000000in}}{%
\pgfpathmoveto{\pgfqpoint{0.000000in}{0.000000in}}%
\pgfpathlineto{\pgfqpoint{0.000000in}{-0.048611in}}%
\pgfusepath{stroke,fill}%
}%
\begin{pgfscope}%
\pgfsys@transformshift{3.476988in}{0.582778in}%
\pgfsys@useobject{currentmarker}{}%
\end{pgfscope}%
\end{pgfscope}%
\begin{pgfscope}%
\definecolor{textcolor}{rgb}{0.000000,0.000000,0.000000}%
\pgfsetstrokecolor{textcolor}%
\pgfsetfillcolor{textcolor}%
\pgftext[x=3.476988in,y=0.485556in,,top]{\color{textcolor}\sffamily\fontsize{10.000000}{12.000000}\selectfont 1.0}%
\end{pgfscope}%
\begin{pgfscope}%
\pgfpathrectangle{\pgfqpoint{0.834028in}{0.582778in}}{\pgfqpoint{4.759444in}{2.437222in}}%
\pgfusepath{clip}%
\pgfsetbuttcap%
\pgfsetroundjoin%
\pgfsetlinewidth{0.803000pt}%
\definecolor{currentstroke}{rgb}{0.690196,0.690196,0.690196}%
\pgfsetstrokecolor{currentstroke}%
\pgfsetdash{{0.800000pt}{1.320000pt}}{0.000000pt}%
\pgfpathmoveto{\pgfqpoint{4.006109in}{0.582778in}}%
\pgfpathlineto{\pgfqpoint{4.006109in}{3.020000in}}%
\pgfusepath{stroke}%
\end{pgfscope}%
\begin{pgfscope}%
\pgfsetbuttcap%
\pgfsetroundjoin%
\definecolor{currentfill}{rgb}{0.000000,0.000000,0.000000}%
\pgfsetfillcolor{currentfill}%
\pgfsetlinewidth{0.803000pt}%
\definecolor{currentstroke}{rgb}{0.000000,0.000000,0.000000}%
\pgfsetstrokecolor{currentstroke}%
\pgfsetdash{}{0pt}%
\pgfsys@defobject{currentmarker}{\pgfqpoint{0.000000in}{-0.048611in}}{\pgfqpoint{0.000000in}{0.000000in}}{%
\pgfpathmoveto{\pgfqpoint{0.000000in}{0.000000in}}%
\pgfpathlineto{\pgfqpoint{0.000000in}{-0.048611in}}%
\pgfusepath{stroke,fill}%
}%
\begin{pgfscope}%
\pgfsys@transformshift{4.006109in}{0.582778in}%
\pgfsys@useobject{currentmarker}{}%
\end{pgfscope}%
\end{pgfscope}%
\begin{pgfscope}%
\definecolor{textcolor}{rgb}{0.000000,0.000000,0.000000}%
\pgfsetstrokecolor{textcolor}%
\pgfsetfillcolor{textcolor}%
\pgftext[x=4.006109in,y=0.485556in,,top]{\color{textcolor}\sffamily\fontsize{10.000000}{12.000000}\selectfont 1.2}%
\end{pgfscope}%
\begin{pgfscope}%
\pgfpathrectangle{\pgfqpoint{0.834028in}{0.582778in}}{\pgfqpoint{4.759444in}{2.437222in}}%
\pgfusepath{clip}%
\pgfsetbuttcap%
\pgfsetroundjoin%
\pgfsetlinewidth{0.803000pt}%
\definecolor{currentstroke}{rgb}{0.690196,0.690196,0.690196}%
\pgfsetstrokecolor{currentstroke}%
\pgfsetdash{{0.800000pt}{1.320000pt}}{0.000000pt}%
\pgfpathmoveto{\pgfqpoint{4.535230in}{0.582778in}}%
\pgfpathlineto{\pgfqpoint{4.535230in}{3.020000in}}%
\pgfusepath{stroke}%
\end{pgfscope}%
\begin{pgfscope}%
\pgfsetbuttcap%
\pgfsetroundjoin%
\definecolor{currentfill}{rgb}{0.000000,0.000000,0.000000}%
\pgfsetfillcolor{currentfill}%
\pgfsetlinewidth{0.803000pt}%
\definecolor{currentstroke}{rgb}{0.000000,0.000000,0.000000}%
\pgfsetstrokecolor{currentstroke}%
\pgfsetdash{}{0pt}%
\pgfsys@defobject{currentmarker}{\pgfqpoint{0.000000in}{-0.048611in}}{\pgfqpoint{0.000000in}{0.000000in}}{%
\pgfpathmoveto{\pgfqpoint{0.000000in}{0.000000in}}%
\pgfpathlineto{\pgfqpoint{0.000000in}{-0.048611in}}%
\pgfusepath{stroke,fill}%
}%
\begin{pgfscope}%
\pgfsys@transformshift{4.535230in}{0.582778in}%
\pgfsys@useobject{currentmarker}{}%
\end{pgfscope}%
\end{pgfscope}%
\begin{pgfscope}%
\definecolor{textcolor}{rgb}{0.000000,0.000000,0.000000}%
\pgfsetstrokecolor{textcolor}%
\pgfsetfillcolor{textcolor}%
\pgftext[x=4.535230in,y=0.485556in,,top]{\color{textcolor}\sffamily\fontsize{10.000000}{12.000000}\selectfont 1.4}%
\end{pgfscope}%
\begin{pgfscope}%
\pgfpathrectangle{\pgfqpoint{0.834028in}{0.582778in}}{\pgfqpoint{4.759444in}{2.437222in}}%
\pgfusepath{clip}%
\pgfsetbuttcap%
\pgfsetroundjoin%
\pgfsetlinewidth{0.803000pt}%
\definecolor{currentstroke}{rgb}{0.690196,0.690196,0.690196}%
\pgfsetstrokecolor{currentstroke}%
\pgfsetdash{{0.800000pt}{1.320000pt}}{0.000000pt}%
\pgfpathmoveto{\pgfqpoint{5.064351in}{0.582778in}}%
\pgfpathlineto{\pgfqpoint{5.064351in}{3.020000in}}%
\pgfusepath{stroke}%
\end{pgfscope}%
\begin{pgfscope}%
\pgfsetbuttcap%
\pgfsetroundjoin%
\definecolor{currentfill}{rgb}{0.000000,0.000000,0.000000}%
\pgfsetfillcolor{currentfill}%
\pgfsetlinewidth{0.803000pt}%
\definecolor{currentstroke}{rgb}{0.000000,0.000000,0.000000}%
\pgfsetstrokecolor{currentstroke}%
\pgfsetdash{}{0pt}%
\pgfsys@defobject{currentmarker}{\pgfqpoint{0.000000in}{-0.048611in}}{\pgfqpoint{0.000000in}{0.000000in}}{%
\pgfpathmoveto{\pgfqpoint{0.000000in}{0.000000in}}%
\pgfpathlineto{\pgfqpoint{0.000000in}{-0.048611in}}%
\pgfusepath{stroke,fill}%
}%
\begin{pgfscope}%
\pgfsys@transformshift{5.064351in}{0.582778in}%
\pgfsys@useobject{currentmarker}{}%
\end{pgfscope}%
\end{pgfscope}%
\begin{pgfscope}%
\definecolor{textcolor}{rgb}{0.000000,0.000000,0.000000}%
\pgfsetstrokecolor{textcolor}%
\pgfsetfillcolor{textcolor}%
\pgftext[x=5.064351in,y=0.485556in,,top]{\color{textcolor}\sffamily\fontsize{10.000000}{12.000000}\selectfont 1.6}%
\end{pgfscope}%
\begin{pgfscope}%
\pgfpathrectangle{\pgfqpoint{0.834028in}{0.582778in}}{\pgfqpoint{4.759444in}{2.437222in}}%
\pgfusepath{clip}%
\pgfsetbuttcap%
\pgfsetroundjoin%
\pgfsetlinewidth{0.803000pt}%
\definecolor{currentstroke}{rgb}{0.690196,0.690196,0.690196}%
\pgfsetstrokecolor{currentstroke}%
\pgfsetdash{{0.800000pt}{1.320000pt}}{0.000000pt}%
\pgfpathmoveto{\pgfqpoint{5.593472in}{0.582778in}}%
\pgfpathlineto{\pgfqpoint{5.593472in}{3.020000in}}%
\pgfusepath{stroke}%
\end{pgfscope}%
\begin{pgfscope}%
\pgfsetbuttcap%
\pgfsetroundjoin%
\definecolor{currentfill}{rgb}{0.000000,0.000000,0.000000}%
\pgfsetfillcolor{currentfill}%
\pgfsetlinewidth{0.803000pt}%
\definecolor{currentstroke}{rgb}{0.000000,0.000000,0.000000}%
\pgfsetstrokecolor{currentstroke}%
\pgfsetdash{}{0pt}%
\pgfsys@defobject{currentmarker}{\pgfqpoint{0.000000in}{-0.048611in}}{\pgfqpoint{0.000000in}{0.000000in}}{%
\pgfpathmoveto{\pgfqpoint{0.000000in}{0.000000in}}%
\pgfpathlineto{\pgfqpoint{0.000000in}{-0.048611in}}%
\pgfusepath{stroke,fill}%
}%
\begin{pgfscope}%
\pgfsys@transformshift{5.593472in}{0.582778in}%
\pgfsys@useobject{currentmarker}{}%
\end{pgfscope}%
\end{pgfscope}%
\begin{pgfscope}%
\definecolor{textcolor}{rgb}{0.000000,0.000000,0.000000}%
\pgfsetstrokecolor{textcolor}%
\pgfsetfillcolor{textcolor}%
\pgftext[x=5.593472in,y=0.485556in,,top]{\color{textcolor}\sffamily\fontsize{10.000000}{12.000000}\selectfont 1.8}%
\end{pgfscope}%
\begin{pgfscope}%
\definecolor{textcolor}{rgb}{0.000000,0.000000,0.000000}%
\pgfsetstrokecolor{textcolor}%
\pgfsetfillcolor{textcolor}%
\pgftext[x=3.213750in,y=0.295587in,,top]{\color{textcolor}\sffamily\fontsize{10.000000}{12.000000}\selectfont Frequency (Hz)}%
\end{pgfscope}%
\begin{pgfscope}%
\definecolor{textcolor}{rgb}{0.000000,0.000000,0.000000}%
\pgfsetstrokecolor{textcolor}%
\pgfsetfillcolor{textcolor}%
\pgftext[x=5.593472in,y=0.309476in,right,top]{\color{textcolor}\sffamily\fontsize{10.000000}{12.000000}\selectfont 1e8}%
\end{pgfscope}%
\begin{pgfscope}%
\pgfpathrectangle{\pgfqpoint{0.834028in}{0.582778in}}{\pgfqpoint{4.759444in}{2.437222in}}%
\pgfusepath{clip}%
\pgfsetbuttcap%
\pgfsetroundjoin%
\pgfsetlinewidth{0.803000pt}%
\definecolor{currentstroke}{rgb}{0.690196,0.690196,0.690196}%
\pgfsetstrokecolor{currentstroke}%
\pgfsetdash{{0.800000pt}{1.320000pt}}{0.000000pt}%
\pgfpathmoveto{\pgfqpoint{0.834028in}{0.684437in}}%
\pgfpathlineto{\pgfqpoint{5.593472in}{0.684437in}}%
\pgfusepath{stroke}%
\end{pgfscope}%
\begin{pgfscope}%
\pgfsetbuttcap%
\pgfsetroundjoin%
\definecolor{currentfill}{rgb}{0.000000,0.000000,0.000000}%
\pgfsetfillcolor{currentfill}%
\pgfsetlinewidth{0.803000pt}%
\definecolor{currentstroke}{rgb}{0.000000,0.000000,0.000000}%
\pgfsetstrokecolor{currentstroke}%
\pgfsetdash{}{0pt}%
\pgfsys@defobject{currentmarker}{\pgfqpoint{-0.048611in}{0.000000in}}{\pgfqpoint{-0.000000in}{0.000000in}}{%
\pgfpathmoveto{\pgfqpoint{-0.000000in}{0.000000in}}%
\pgfpathlineto{\pgfqpoint{-0.048611in}{0.000000in}}%
\pgfusepath{stroke,fill}%
}%
\begin{pgfscope}%
\pgfsys@transformshift{0.834028in}{0.684437in}%
\pgfsys@useobject{currentmarker}{}%
\end{pgfscope}%
\end{pgfscope}%
\begin{pgfscope}%
\definecolor{textcolor}{rgb}{0.000000,0.000000,0.000000}%
\pgfsetstrokecolor{textcolor}%
\pgfsetfillcolor{textcolor}%
\pgftext[x=0.648440in, y=0.631675in, left, base]{\color{textcolor}\sffamily\fontsize{10.000000}{12.000000}\selectfont 0}%
\end{pgfscope}%
\begin{pgfscope}%
\pgfpathrectangle{\pgfqpoint{0.834028in}{0.582778in}}{\pgfqpoint{4.759444in}{2.437222in}}%
\pgfusepath{clip}%
\pgfsetbuttcap%
\pgfsetroundjoin%
\pgfsetlinewidth{0.803000pt}%
\definecolor{currentstroke}{rgb}{0.690196,0.690196,0.690196}%
\pgfsetstrokecolor{currentstroke}%
\pgfsetdash{{0.800000pt}{1.320000pt}}{0.000000pt}%
\pgfpathmoveto{\pgfqpoint{0.834028in}{1.185750in}}%
\pgfpathlineto{\pgfqpoint{5.593472in}{1.185750in}}%
\pgfusepath{stroke}%
\end{pgfscope}%
\begin{pgfscope}%
\pgfsetbuttcap%
\pgfsetroundjoin%
\definecolor{currentfill}{rgb}{0.000000,0.000000,0.000000}%
\pgfsetfillcolor{currentfill}%
\pgfsetlinewidth{0.803000pt}%
\definecolor{currentstroke}{rgb}{0.000000,0.000000,0.000000}%
\pgfsetstrokecolor{currentstroke}%
\pgfsetdash{}{0pt}%
\pgfsys@defobject{currentmarker}{\pgfqpoint{-0.048611in}{0.000000in}}{\pgfqpoint{-0.000000in}{0.000000in}}{%
\pgfpathmoveto{\pgfqpoint{-0.000000in}{0.000000in}}%
\pgfpathlineto{\pgfqpoint{-0.048611in}{0.000000in}}%
\pgfusepath{stroke,fill}%
}%
\begin{pgfscope}%
\pgfsys@transformshift{0.834028in}{1.185750in}%
\pgfsys@useobject{currentmarker}{}%
\end{pgfscope}%
\end{pgfscope}%
\begin{pgfscope}%
\definecolor{textcolor}{rgb}{0.000000,0.000000,0.000000}%
\pgfsetstrokecolor{textcolor}%
\pgfsetfillcolor{textcolor}%
\pgftext[x=0.383344in, y=1.132989in, left, base]{\color{textcolor}\sffamily\fontsize{10.000000}{12.000000}\selectfont 1000}%
\end{pgfscope}%
\begin{pgfscope}%
\pgfpathrectangle{\pgfqpoint{0.834028in}{0.582778in}}{\pgfqpoint{4.759444in}{2.437222in}}%
\pgfusepath{clip}%
\pgfsetbuttcap%
\pgfsetroundjoin%
\pgfsetlinewidth{0.803000pt}%
\definecolor{currentstroke}{rgb}{0.690196,0.690196,0.690196}%
\pgfsetstrokecolor{currentstroke}%
\pgfsetdash{{0.800000pt}{1.320000pt}}{0.000000pt}%
\pgfpathmoveto{\pgfqpoint{0.834028in}{1.687064in}}%
\pgfpathlineto{\pgfqpoint{5.593472in}{1.687064in}}%
\pgfusepath{stroke}%
\end{pgfscope}%
\begin{pgfscope}%
\pgfsetbuttcap%
\pgfsetroundjoin%
\definecolor{currentfill}{rgb}{0.000000,0.000000,0.000000}%
\pgfsetfillcolor{currentfill}%
\pgfsetlinewidth{0.803000pt}%
\definecolor{currentstroke}{rgb}{0.000000,0.000000,0.000000}%
\pgfsetstrokecolor{currentstroke}%
\pgfsetdash{}{0pt}%
\pgfsys@defobject{currentmarker}{\pgfqpoint{-0.048611in}{0.000000in}}{\pgfqpoint{-0.000000in}{0.000000in}}{%
\pgfpathmoveto{\pgfqpoint{-0.000000in}{0.000000in}}%
\pgfpathlineto{\pgfqpoint{-0.048611in}{0.000000in}}%
\pgfusepath{stroke,fill}%
}%
\begin{pgfscope}%
\pgfsys@transformshift{0.834028in}{1.687064in}%
\pgfsys@useobject{currentmarker}{}%
\end{pgfscope}%
\end{pgfscope}%
\begin{pgfscope}%
\definecolor{textcolor}{rgb}{0.000000,0.000000,0.000000}%
\pgfsetstrokecolor{textcolor}%
\pgfsetfillcolor{textcolor}%
\pgftext[x=0.383344in, y=1.634303in, left, base]{\color{textcolor}\sffamily\fontsize{10.000000}{12.000000}\selectfont 2000}%
\end{pgfscope}%
\begin{pgfscope}%
\pgfpathrectangle{\pgfqpoint{0.834028in}{0.582778in}}{\pgfqpoint{4.759444in}{2.437222in}}%
\pgfusepath{clip}%
\pgfsetbuttcap%
\pgfsetroundjoin%
\pgfsetlinewidth{0.803000pt}%
\definecolor{currentstroke}{rgb}{0.690196,0.690196,0.690196}%
\pgfsetstrokecolor{currentstroke}%
\pgfsetdash{{0.800000pt}{1.320000pt}}{0.000000pt}%
\pgfpathmoveto{\pgfqpoint{0.834028in}{2.188378in}}%
\pgfpathlineto{\pgfqpoint{5.593472in}{2.188378in}}%
\pgfusepath{stroke}%
\end{pgfscope}%
\begin{pgfscope}%
\pgfsetbuttcap%
\pgfsetroundjoin%
\definecolor{currentfill}{rgb}{0.000000,0.000000,0.000000}%
\pgfsetfillcolor{currentfill}%
\pgfsetlinewidth{0.803000pt}%
\definecolor{currentstroke}{rgb}{0.000000,0.000000,0.000000}%
\pgfsetstrokecolor{currentstroke}%
\pgfsetdash{}{0pt}%
\pgfsys@defobject{currentmarker}{\pgfqpoint{-0.048611in}{0.000000in}}{\pgfqpoint{-0.000000in}{0.000000in}}{%
\pgfpathmoveto{\pgfqpoint{-0.000000in}{0.000000in}}%
\pgfpathlineto{\pgfqpoint{-0.048611in}{0.000000in}}%
\pgfusepath{stroke,fill}%
}%
\begin{pgfscope}%
\pgfsys@transformshift{0.834028in}{2.188378in}%
\pgfsys@useobject{currentmarker}{}%
\end{pgfscope}%
\end{pgfscope}%
\begin{pgfscope}%
\definecolor{textcolor}{rgb}{0.000000,0.000000,0.000000}%
\pgfsetstrokecolor{textcolor}%
\pgfsetfillcolor{textcolor}%
\pgftext[x=0.383344in, y=2.135617in, left, base]{\color{textcolor}\sffamily\fontsize{10.000000}{12.000000}\selectfont 3000}%
\end{pgfscope}%
\begin{pgfscope}%
\pgfpathrectangle{\pgfqpoint{0.834028in}{0.582778in}}{\pgfqpoint{4.759444in}{2.437222in}}%
\pgfusepath{clip}%
\pgfsetbuttcap%
\pgfsetroundjoin%
\pgfsetlinewidth{0.803000pt}%
\definecolor{currentstroke}{rgb}{0.690196,0.690196,0.690196}%
\pgfsetstrokecolor{currentstroke}%
\pgfsetdash{{0.800000pt}{1.320000pt}}{0.000000pt}%
\pgfpathmoveto{\pgfqpoint{0.834028in}{2.689692in}}%
\pgfpathlineto{\pgfqpoint{5.593472in}{2.689692in}}%
\pgfusepath{stroke}%
\end{pgfscope}%
\begin{pgfscope}%
\pgfsetbuttcap%
\pgfsetroundjoin%
\definecolor{currentfill}{rgb}{0.000000,0.000000,0.000000}%
\pgfsetfillcolor{currentfill}%
\pgfsetlinewidth{0.803000pt}%
\definecolor{currentstroke}{rgb}{0.000000,0.000000,0.000000}%
\pgfsetstrokecolor{currentstroke}%
\pgfsetdash{}{0pt}%
\pgfsys@defobject{currentmarker}{\pgfqpoint{-0.048611in}{0.000000in}}{\pgfqpoint{-0.000000in}{0.000000in}}{%
\pgfpathmoveto{\pgfqpoint{-0.000000in}{0.000000in}}%
\pgfpathlineto{\pgfqpoint{-0.048611in}{0.000000in}}%
\pgfusepath{stroke,fill}%
}%
\begin{pgfscope}%
\pgfsys@transformshift{0.834028in}{2.689692in}%
\pgfsys@useobject{currentmarker}{}%
\end{pgfscope}%
\end{pgfscope}%
\begin{pgfscope}%
\definecolor{textcolor}{rgb}{0.000000,0.000000,0.000000}%
\pgfsetstrokecolor{textcolor}%
\pgfsetfillcolor{textcolor}%
\pgftext[x=0.383344in, y=2.636930in, left, base]{\color{textcolor}\sffamily\fontsize{10.000000}{12.000000}\selectfont 4000}%
\end{pgfscope}%
\begin{pgfscope}%
\definecolor{textcolor}{rgb}{0.000000,0.000000,0.000000}%
\pgfsetstrokecolor{textcolor}%
\pgfsetfillcolor{textcolor}%
\pgftext[x=0.327789in,y=1.801389in,,bottom,rotate=90.000000]{\color{textcolor}\sffamily\fontsize{10.000000}{12.000000}\selectfont Impedance |Z| (\(\displaystyle \Omega\))}%
\end{pgfscope}%
\begin{pgfscope}%
\pgfpathrectangle{\pgfqpoint{0.834028in}{0.582778in}}{\pgfqpoint{4.759444in}{2.437222in}}%
\pgfusepath{clip}%
\pgfsetrectcap%
\pgfsetroundjoin%
\pgfsetlinewidth{2.007500pt}%
\definecolor{currentstroke}{rgb}{0.000000,0.000000,0.000000}%
\pgfsetstrokecolor{currentstroke}%
\pgfsetdash{}{0pt}%
\pgfpathmoveto{\pgfqpoint{0.857838in}{0.702484in}}%
\pgfpathlineto{\pgfqpoint{0.862941in}{0.693561in}}%
\pgfpathlineto{\pgfqpoint{0.919075in}{0.710004in}}%
\pgfpathlineto{\pgfqpoint{0.990517in}{0.733415in}}%
\pgfpathlineto{\pgfqpoint{1.051753in}{0.757227in}}%
\pgfpathlineto{\pgfqpoint{1.097681in}{0.777681in}}%
\pgfpathlineto{\pgfqpoint{1.118093in}{0.788459in}}%
\pgfpathlineto{\pgfqpoint{1.128299in}{0.793573in}}%
\pgfpathlineto{\pgfqpoint{1.189535in}{0.832775in}}%
\pgfpathlineto{\pgfqpoint{1.215050in}{0.853229in}}%
\pgfpathlineto{\pgfqpoint{1.220153in}{0.856889in}}%
\pgfpathlineto{\pgfqpoint{1.235462in}{0.873131in}}%
\pgfpathlineto{\pgfqpoint{1.255875in}{0.892733in}}%
\pgfpathlineto{\pgfqpoint{1.260978in}{0.898247in}}%
\pgfpathlineto{\pgfqpoint{1.266081in}{0.907120in}}%
\pgfpathlineto{\pgfqpoint{1.271184in}{0.913437in}}%
\pgfpathlineto{\pgfqpoint{1.276287in}{0.916746in}}%
\pgfpathlineto{\pgfqpoint{1.291596in}{0.938703in}}%
\pgfpathlineto{\pgfqpoint{1.301802in}{0.955748in}}%
\pgfpathlineto{\pgfqpoint{1.306905in}{0.960360in}}%
\pgfpathlineto{\pgfqpoint{1.317111in}{0.980513in}}%
\pgfpathlineto{\pgfqpoint{1.322214in}{0.985776in}}%
\pgfpathlineto{\pgfqpoint{1.327317in}{0.997407in}}%
\pgfpathlineto{\pgfqpoint{1.332420in}{1.003573in}}%
\pgfpathlineto{\pgfqpoint{1.342626in}{1.030794in}}%
\pgfpathlineto{\pgfqpoint{1.352832in}{1.046185in}}%
\pgfpathlineto{\pgfqpoint{1.357935in}{1.063029in}}%
\pgfpathlineto{\pgfqpoint{1.363038in}{1.072053in}}%
\pgfpathlineto{\pgfqpoint{1.368141in}{1.091403in}}%
\pgfpathlineto{\pgfqpoint{1.373244in}{1.101831in}}%
\pgfpathlineto{\pgfqpoint{1.378347in}{1.124390in}}%
\pgfpathlineto{\pgfqpoint{1.388553in}{1.150408in}}%
\pgfpathlineto{\pgfqpoint{1.393656in}{1.178682in}}%
\pgfpathlineto{\pgfqpoint{1.408965in}{1.228212in}}%
\pgfpathlineto{\pgfqpoint{1.419171in}{1.267114in}}%
\pgfpathlineto{\pgfqpoint{1.424274in}{1.311981in}}%
\pgfpathlineto{\pgfqpoint{1.439584in}{1.390688in}}%
\pgfpathlineto{\pgfqpoint{1.449790in}{1.457312in}}%
\pgfpathlineto{\pgfqpoint{1.459996in}{1.537673in}}%
\pgfpathlineto{\pgfqpoint{1.465099in}{1.537673in}}%
\pgfpathlineto{\pgfqpoint{1.475305in}{1.636532in}}%
\pgfpathlineto{\pgfqpoint{1.480408in}{1.636532in}}%
\pgfpathlineto{\pgfqpoint{1.490614in}{1.761008in}}%
\pgfpathlineto{\pgfqpoint{1.495717in}{1.761008in}}%
\pgfpathlineto{\pgfqpoint{1.500820in}{1.836155in}}%
\pgfpathlineto{\pgfqpoint{1.505923in}{1.836155in}}%
\pgfpathlineto{\pgfqpoint{1.511026in}{1.268768in}}%
\pgfpathlineto{\pgfqpoint{1.516129in}{1.922431in}}%
\pgfpathlineto{\pgfqpoint{1.521232in}{2.022393in}}%
\pgfpathlineto{\pgfqpoint{1.526335in}{2.022393in}}%
\pgfpathlineto{\pgfqpoint{1.531438in}{2.139550in}}%
\pgfpathlineto{\pgfqpoint{1.541644in}{2.139550in}}%
\pgfpathlineto{\pgfqpoint{1.546747in}{2.278665in}}%
\pgfpathlineto{\pgfqpoint{1.551850in}{2.062298in}}%
\pgfpathlineto{\pgfqpoint{1.556953in}{2.278665in}}%
\pgfpathlineto{\pgfqpoint{1.562056in}{2.446304in}}%
\pgfpathlineto{\pgfqpoint{1.582468in}{2.446304in}}%
\pgfpathlineto{\pgfqpoint{1.587571in}{2.651943in}}%
\pgfpathlineto{\pgfqpoint{1.592674in}{2.724884in}}%
\pgfpathlineto{\pgfqpoint{1.597777in}{2.651943in}}%
\pgfpathlineto{\pgfqpoint{1.602880in}{2.724884in}}%
\pgfpathlineto{\pgfqpoint{1.607984in}{2.472573in}}%
\pgfpathlineto{\pgfqpoint{1.613087in}{2.724884in}}%
\pgfpathlineto{\pgfqpoint{1.618190in}{2.651943in}}%
\pgfpathlineto{\pgfqpoint{1.633499in}{2.651943in}}%
\pgfpathlineto{\pgfqpoint{1.638602in}{2.909217in}}%
\pgfpathlineto{\pgfqpoint{1.766177in}{2.909217in}}%
\pgfpathlineto{\pgfqpoint{1.771280in}{2.472573in}}%
\pgfpathlineto{\pgfqpoint{1.776383in}{2.651943in}}%
\pgfpathlineto{\pgfqpoint{1.827414in}{2.651943in}}%
\pgfpathlineto{\pgfqpoint{1.832517in}{2.314007in}}%
\pgfpathlineto{\pgfqpoint{1.837620in}{2.446304in}}%
\pgfpathlineto{\pgfqpoint{1.842723in}{2.314007in}}%
\pgfpathlineto{\pgfqpoint{1.868238in}{2.314007in}}%
\pgfpathlineto{\pgfqpoint{1.873341in}{2.178653in}}%
\pgfpathlineto{\pgfqpoint{1.883547in}{2.178653in}}%
\pgfpathlineto{\pgfqpoint{1.888650in}{2.278665in}}%
\pgfpathlineto{\pgfqpoint{1.893753in}{2.139550in}}%
\pgfpathlineto{\pgfqpoint{1.898856in}{2.062298in}}%
\pgfpathlineto{\pgfqpoint{1.903959in}{2.139550in}}%
\pgfpathlineto{\pgfqpoint{1.909062in}{2.062298in}}%
\pgfpathlineto{\pgfqpoint{1.914165in}{2.139550in}}%
\pgfpathlineto{\pgfqpoint{1.919268in}{2.022393in}}%
\pgfpathlineto{\pgfqpoint{1.934577in}{2.022393in}}%
\pgfpathlineto{\pgfqpoint{1.939680in}{1.922431in}}%
\pgfpathlineto{\pgfqpoint{1.949886in}{1.922431in}}%
\pgfpathlineto{\pgfqpoint{1.954990in}{1.836155in}}%
\pgfpathlineto{\pgfqpoint{1.965196in}{1.836155in}}%
\pgfpathlineto{\pgfqpoint{1.970299in}{1.761008in}}%
\pgfpathlineto{\pgfqpoint{1.980505in}{1.761008in}}%
\pgfpathlineto{\pgfqpoint{1.990711in}{1.627508in}}%
\pgfpathlineto{\pgfqpoint{1.995814in}{1.694985in}}%
\pgfpathlineto{\pgfqpoint{2.000917in}{1.636532in}}%
\pgfpathlineto{\pgfqpoint{2.006020in}{1.636532in}}%
\pgfpathlineto{\pgfqpoint{2.011123in}{1.584445in}}%
\pgfpathlineto{\pgfqpoint{2.021329in}{1.584445in}}%
\pgfpathlineto{\pgfqpoint{2.026432in}{1.537673in}}%
\pgfpathlineto{\pgfqpoint{2.031535in}{1.537673in}}%
\pgfpathlineto{\pgfqpoint{2.036638in}{1.495512in}}%
\pgfpathlineto{\pgfqpoint{2.046844in}{1.495512in}}%
\pgfpathlineto{\pgfqpoint{2.051947in}{1.457312in}}%
\pgfpathlineto{\pgfqpoint{2.057050in}{1.457312in}}%
\pgfpathlineto{\pgfqpoint{2.062153in}{1.422471in}}%
\pgfpathlineto{\pgfqpoint{2.067256in}{1.422471in}}%
\pgfpathlineto{\pgfqpoint{2.072359in}{1.390688in}}%
\pgfpathlineto{\pgfqpoint{2.077462in}{1.390688in}}%
\pgfpathlineto{\pgfqpoint{2.082565in}{1.361461in}}%
\pgfpathlineto{\pgfqpoint{2.087668in}{1.361461in}}%
\pgfpathlineto{\pgfqpoint{2.092771in}{1.334591in}}%
\pgfpathlineto{\pgfqpoint{2.097874in}{1.334591in}}%
\pgfpathlineto{\pgfqpoint{2.102977in}{1.309725in}}%
\pgfpathlineto{\pgfqpoint{2.113183in}{1.309725in}}%
\pgfpathlineto{\pgfqpoint{2.118286in}{1.286715in}}%
\pgfpathlineto{\pgfqpoint{2.123389in}{1.286715in}}%
\pgfpathlineto{\pgfqpoint{2.128493in}{1.265309in}}%
\pgfpathlineto{\pgfqpoint{2.133596in}{1.265309in}}%
\pgfpathlineto{\pgfqpoint{2.138699in}{1.245357in}}%
\pgfpathlineto{\pgfqpoint{2.143802in}{1.245357in}}%
\pgfpathlineto{\pgfqpoint{2.148905in}{1.222447in}}%
\pgfpathlineto{\pgfqpoint{2.154008in}{1.222447in}}%
\pgfpathlineto{\pgfqpoint{2.159111in}{1.205402in}}%
\pgfpathlineto{\pgfqpoint{2.164214in}{1.209312in}}%
\pgfpathlineto{\pgfqpoint{2.169317in}{1.192969in}}%
\pgfpathlineto{\pgfqpoint{2.174420in}{1.189410in}}%
\pgfpathlineto{\pgfqpoint{2.179523in}{1.174371in}}%
\pgfpathlineto{\pgfqpoint{2.184626in}{1.174371in}}%
\pgfpathlineto{\pgfqpoint{2.189729in}{1.160183in}}%
\pgfpathlineto{\pgfqpoint{2.194832in}{1.160183in}}%
\pgfpathlineto{\pgfqpoint{2.199935in}{1.149505in}}%
\pgfpathlineto{\pgfqpoint{2.205038in}{1.146748in}}%
\pgfpathlineto{\pgfqpoint{2.210141in}{1.136371in}}%
\pgfpathlineto{\pgfqpoint{2.215244in}{1.134115in}}%
\pgfpathlineto{\pgfqpoint{2.220347in}{1.134115in}}%
\pgfpathlineto{\pgfqpoint{2.230553in}{1.110704in}}%
\pgfpathlineto{\pgfqpoint{2.235656in}{1.110704in}}%
\pgfpathlineto{\pgfqpoint{2.240759in}{1.099875in}}%
\pgfpathlineto{\pgfqpoint{2.245862in}{1.099875in}}%
\pgfpathlineto{\pgfqpoint{2.250965in}{1.089599in}}%
\pgfpathlineto{\pgfqpoint{2.261171in}{1.089599in}}%
\pgfpathlineto{\pgfqpoint{2.266274in}{1.079823in}}%
\pgfpathlineto{\pgfqpoint{2.271377in}{1.079823in}}%
\pgfpathlineto{\pgfqpoint{2.281583in}{1.061575in}}%
\pgfpathlineto{\pgfqpoint{2.286686in}{1.063029in}}%
\pgfpathlineto{\pgfqpoint{2.291789in}{1.061575in}}%
\pgfpathlineto{\pgfqpoint{2.296892in}{1.053053in}}%
\pgfpathlineto{\pgfqpoint{2.301996in}{1.054406in}}%
\pgfpathlineto{\pgfqpoint{2.307099in}{1.044931in}}%
\pgfpathlineto{\pgfqpoint{2.312202in}{1.044931in}}%
\pgfpathlineto{\pgfqpoint{2.322408in}{1.030794in}}%
\pgfpathlineto{\pgfqpoint{2.327511in}{1.029641in}}%
\pgfpathlineto{\pgfqpoint{2.332614in}{1.029641in}}%
\pgfpathlineto{\pgfqpoint{2.337717in}{1.022523in}}%
\pgfpathlineto{\pgfqpoint{2.342820in}{1.023575in}}%
\pgfpathlineto{\pgfqpoint{2.347923in}{1.016607in}}%
\pgfpathlineto{\pgfqpoint{2.353026in}{1.016607in}}%
\pgfpathlineto{\pgfqpoint{2.358129in}{1.009990in}}%
\pgfpathlineto{\pgfqpoint{2.363232in}{1.009990in}}%
\pgfpathlineto{\pgfqpoint{2.368335in}{1.003573in}}%
\pgfpathlineto{\pgfqpoint{2.373438in}{1.003573in}}%
\pgfpathlineto{\pgfqpoint{2.378541in}{0.997407in}}%
\pgfpathlineto{\pgfqpoint{2.383644in}{0.997407in}}%
\pgfpathlineto{\pgfqpoint{2.388747in}{0.989436in}}%
\pgfpathlineto{\pgfqpoint{2.393850in}{0.974346in}}%
\pgfpathlineto{\pgfqpoint{2.398953in}{0.985776in}}%
\pgfpathlineto{\pgfqpoint{2.404056in}{0.985776in}}%
\pgfpathlineto{\pgfqpoint{2.409159in}{0.980312in}}%
\pgfpathlineto{\pgfqpoint{2.414262in}{0.980312in}}%
\pgfpathlineto{\pgfqpoint{2.419365in}{0.974998in}}%
\pgfpathlineto{\pgfqpoint{2.424468in}{0.974998in}}%
\pgfpathlineto{\pgfqpoint{2.429571in}{0.968230in}}%
\pgfpathlineto{\pgfqpoint{2.434674in}{0.969885in}}%
\pgfpathlineto{\pgfqpoint{2.439777in}{0.964922in}}%
\pgfpathlineto{\pgfqpoint{2.444880in}{0.936397in}}%
\pgfpathlineto{\pgfqpoint{2.449983in}{0.960159in}}%
\pgfpathlineto{\pgfqpoint{2.460189in}{0.960159in}}%
\pgfpathlineto{\pgfqpoint{2.465292in}{0.955547in}}%
\pgfpathlineto{\pgfqpoint{2.470395in}{0.955547in}}%
\pgfpathlineto{\pgfqpoint{2.475498in}{0.948529in}}%
\pgfpathlineto{\pgfqpoint{2.480602in}{0.951086in}}%
\pgfpathlineto{\pgfqpoint{2.485705in}{0.946774in}}%
\pgfpathlineto{\pgfqpoint{2.490808in}{0.946774in}}%
\pgfpathlineto{\pgfqpoint{2.495911in}{0.942563in}}%
\pgfpathlineto{\pgfqpoint{2.501014in}{0.942563in}}%
\pgfpathlineto{\pgfqpoint{2.506117in}{0.938553in}}%
\pgfpathlineto{\pgfqpoint{2.511220in}{0.938553in}}%
\pgfpathlineto{\pgfqpoint{2.516323in}{0.934191in}}%
\pgfpathlineto{\pgfqpoint{2.521426in}{0.934191in}}%
\pgfpathlineto{\pgfqpoint{2.526529in}{0.930381in}}%
\pgfpathlineto{\pgfqpoint{2.531632in}{0.930381in}}%
\pgfpathlineto{\pgfqpoint{2.536735in}{0.926722in}}%
\pgfpathlineto{\pgfqpoint{2.546941in}{0.926722in}}%
\pgfpathlineto{\pgfqpoint{2.552044in}{0.923112in}}%
\pgfpathlineto{\pgfqpoint{2.557147in}{0.923112in}}%
\pgfpathlineto{\pgfqpoint{2.562250in}{0.919653in}}%
\pgfpathlineto{\pgfqpoint{2.572456in}{0.919653in}}%
\pgfpathlineto{\pgfqpoint{2.577559in}{0.916294in}}%
\pgfpathlineto{\pgfqpoint{2.587765in}{0.916294in}}%
\pgfpathlineto{\pgfqpoint{2.592868in}{0.913036in}}%
\pgfpathlineto{\pgfqpoint{2.603074in}{0.913036in}}%
\pgfpathlineto{\pgfqpoint{2.608177in}{0.909827in}}%
\pgfpathlineto{\pgfqpoint{2.618383in}{0.909827in}}%
\pgfpathlineto{\pgfqpoint{2.623486in}{0.906719in}}%
\pgfpathlineto{\pgfqpoint{2.628589in}{0.906719in}}%
\pgfpathlineto{\pgfqpoint{2.633692in}{0.903711in}}%
\pgfpathlineto{\pgfqpoint{2.638795in}{0.903711in}}%
\pgfpathlineto{\pgfqpoint{2.643898in}{0.900754in}}%
\pgfpathlineto{\pgfqpoint{2.649001in}{0.900754in}}%
\pgfpathlineto{\pgfqpoint{2.654105in}{0.897896in}}%
\pgfpathlineto{\pgfqpoint{2.659208in}{0.897896in}}%
\pgfpathlineto{\pgfqpoint{2.664311in}{0.895089in}}%
\pgfpathlineto{\pgfqpoint{2.674517in}{0.895089in}}%
\pgfpathlineto{\pgfqpoint{2.684723in}{0.889725in}}%
\pgfpathlineto{\pgfqpoint{2.694929in}{0.889725in}}%
\pgfpathlineto{\pgfqpoint{2.700032in}{0.887168in}}%
\pgfpathlineto{\pgfqpoint{2.705135in}{0.887168in}}%
\pgfpathlineto{\pgfqpoint{2.710238in}{0.884661in}}%
\pgfpathlineto{\pgfqpoint{2.715341in}{0.884661in}}%
\pgfpathlineto{\pgfqpoint{2.720444in}{0.882155in}}%
\pgfpathlineto{\pgfqpoint{2.725547in}{0.882155in}}%
\pgfpathlineto{\pgfqpoint{2.730650in}{0.879799in}}%
\pgfpathlineto{\pgfqpoint{2.735753in}{0.879799in}}%
\pgfpathlineto{\pgfqpoint{2.740856in}{0.877443in}}%
\pgfpathlineto{\pgfqpoint{2.745959in}{0.877443in}}%
\pgfpathlineto{\pgfqpoint{2.751062in}{0.875136in}}%
\pgfpathlineto{\pgfqpoint{2.756165in}{0.875136in}}%
\pgfpathlineto{\pgfqpoint{2.761268in}{0.872931in}}%
\pgfpathlineto{\pgfqpoint{2.771474in}{0.872931in}}%
\pgfpathlineto{\pgfqpoint{2.776577in}{0.870725in}}%
\pgfpathlineto{\pgfqpoint{2.781680in}{0.870725in}}%
\pgfpathlineto{\pgfqpoint{2.786783in}{0.868569in}}%
\pgfpathlineto{\pgfqpoint{2.791886in}{0.868569in}}%
\pgfpathlineto{\pgfqpoint{2.796989in}{0.866514in}}%
\pgfpathlineto{\pgfqpoint{2.807195in}{0.866514in}}%
\pgfpathlineto{\pgfqpoint{2.812298in}{0.864458in}}%
\pgfpathlineto{\pgfqpoint{2.817401in}{0.864458in}}%
\pgfpathlineto{\pgfqpoint{2.822504in}{0.862453in}}%
\pgfpathlineto{\pgfqpoint{2.827608in}{0.862453in}}%
\pgfpathlineto{\pgfqpoint{2.832711in}{0.860498in}}%
\pgfpathlineto{\pgfqpoint{2.842917in}{0.860498in}}%
\pgfpathlineto{\pgfqpoint{2.848020in}{0.858593in}}%
\pgfpathlineto{\pgfqpoint{2.853123in}{0.858593in}}%
\pgfpathlineto{\pgfqpoint{2.858226in}{0.856688in}}%
\pgfpathlineto{\pgfqpoint{2.868432in}{0.856688in}}%
\pgfpathlineto{\pgfqpoint{2.873535in}{0.854883in}}%
\pgfpathlineto{\pgfqpoint{2.878638in}{0.854984in}}%
\pgfpathlineto{\pgfqpoint{2.883741in}{0.853229in}}%
\pgfpathlineto{\pgfqpoint{2.888844in}{0.853079in}}%
\pgfpathlineto{\pgfqpoint{2.893947in}{0.851324in}}%
\pgfpathlineto{\pgfqpoint{2.904153in}{0.851324in}}%
\pgfpathlineto{\pgfqpoint{2.909256in}{0.849720in}}%
\pgfpathlineto{\pgfqpoint{2.914359in}{0.849720in}}%
\pgfpathlineto{\pgfqpoint{2.919462in}{0.848015in}}%
\pgfpathlineto{\pgfqpoint{2.924565in}{0.848015in}}%
\pgfpathlineto{\pgfqpoint{2.929668in}{0.846361in}}%
\pgfpathlineto{\pgfqpoint{2.934771in}{0.846361in}}%
\pgfpathlineto{\pgfqpoint{2.939874in}{0.844707in}}%
\pgfpathlineto{\pgfqpoint{2.950080in}{0.844707in}}%
\pgfpathlineto{\pgfqpoint{2.955183in}{0.843153in}}%
\pgfpathlineto{\pgfqpoint{2.960286in}{0.843153in}}%
\pgfpathlineto{\pgfqpoint{2.965389in}{0.841548in}}%
\pgfpathlineto{\pgfqpoint{2.970492in}{0.841548in}}%
\pgfpathlineto{\pgfqpoint{2.975595in}{0.840044in}}%
\pgfpathlineto{\pgfqpoint{2.980698in}{0.839944in}}%
\pgfpathlineto{\pgfqpoint{2.985801in}{0.838541in}}%
\pgfpathlineto{\pgfqpoint{2.990904in}{0.838541in}}%
\pgfpathlineto{\pgfqpoint{2.996007in}{0.837037in}}%
\pgfpathlineto{\pgfqpoint{3.006214in}{0.837037in}}%
\pgfpathlineto{\pgfqpoint{3.011317in}{0.835583in}}%
\pgfpathlineto{\pgfqpoint{3.016420in}{0.835583in}}%
\pgfpathlineto{\pgfqpoint{3.021523in}{0.834179in}}%
\pgfpathlineto{\pgfqpoint{3.026626in}{0.834179in}}%
\pgfpathlineto{\pgfqpoint{3.031729in}{0.832775in}}%
\pgfpathlineto{\pgfqpoint{3.036832in}{0.832775in}}%
\pgfpathlineto{\pgfqpoint{3.041935in}{0.831422in}}%
\pgfpathlineto{\pgfqpoint{3.052141in}{0.831422in}}%
\pgfpathlineto{\pgfqpoint{3.057244in}{0.830068in}}%
\pgfpathlineto{\pgfqpoint{3.062347in}{0.830068in}}%
\pgfpathlineto{\pgfqpoint{3.067450in}{0.828715in}}%
\pgfpathlineto{\pgfqpoint{3.072553in}{0.828715in}}%
\pgfpathlineto{\pgfqpoint{3.077656in}{0.827411in}}%
\pgfpathlineto{\pgfqpoint{3.087862in}{0.827411in}}%
\pgfpathlineto{\pgfqpoint{3.092965in}{0.826158in}}%
\pgfpathlineto{\pgfqpoint{3.098068in}{0.826158in}}%
\pgfpathlineto{\pgfqpoint{3.103171in}{0.824905in}}%
\pgfpathlineto{\pgfqpoint{3.113377in}{0.824905in}}%
\pgfpathlineto{\pgfqpoint{3.118480in}{0.823652in}}%
\pgfpathlineto{\pgfqpoint{3.123583in}{0.823652in}}%
\pgfpathlineto{\pgfqpoint{3.133789in}{0.822348in}}%
\pgfpathlineto{\pgfqpoint{3.143995in}{0.821245in}}%
\pgfpathlineto{\pgfqpoint{3.149098in}{0.821245in}}%
\pgfpathlineto{\pgfqpoint{3.154201in}{0.820042in}}%
\pgfpathlineto{\pgfqpoint{3.159304in}{0.820042in}}%
\pgfpathlineto{\pgfqpoint{3.164407in}{0.818889in}}%
\pgfpathlineto{\pgfqpoint{3.174614in}{0.818839in}}%
\pgfpathlineto{\pgfqpoint{3.179717in}{0.817686in}}%
\pgfpathlineto{\pgfqpoint{3.184820in}{0.817786in}}%
\pgfpathlineto{\pgfqpoint{3.195026in}{0.816583in}}%
\pgfpathlineto{\pgfqpoint{3.200129in}{0.815480in}}%
\pgfpathlineto{\pgfqpoint{3.205232in}{0.815530in}}%
\pgfpathlineto{\pgfqpoint{3.215438in}{0.814427in}}%
\pgfpathlineto{\pgfqpoint{3.230747in}{0.813324in}}%
\pgfpathlineto{\pgfqpoint{3.246056in}{0.811319in}}%
\pgfpathlineto{\pgfqpoint{3.271571in}{0.810016in}}%
\pgfpathlineto{\pgfqpoint{3.286880in}{0.808261in}}%
\pgfpathlineto{\pgfqpoint{3.302189in}{0.807309in}}%
\pgfpathlineto{\pgfqpoint{3.312395in}{0.806356in}}%
\pgfpathlineto{\pgfqpoint{3.332807in}{0.805404in}}%
\pgfpathlineto{\pgfqpoint{3.343013in}{0.804451in}}%
\pgfpathlineto{\pgfqpoint{3.363426in}{0.803549in}}%
\pgfpathlineto{\pgfqpoint{3.373632in}{0.802646in}}%
\pgfpathlineto{\pgfqpoint{3.394044in}{0.801794in}}%
\pgfpathlineto{\pgfqpoint{3.409353in}{0.800040in}}%
\pgfpathlineto{\pgfqpoint{3.434868in}{0.799187in}}%
\pgfpathlineto{\pgfqpoint{3.450177in}{0.797533in}}%
\pgfpathlineto{\pgfqpoint{3.470589in}{0.796731in}}%
\pgfpathlineto{\pgfqpoint{3.485898in}{0.795127in}}%
\pgfpathlineto{\pgfqpoint{3.506310in}{0.794375in}}%
\pgfpathlineto{\pgfqpoint{3.521619in}{0.792821in}}%
\pgfpathlineto{\pgfqpoint{3.542032in}{0.792069in}}%
\pgfpathlineto{\pgfqpoint{3.557341in}{0.790565in}}%
\pgfpathlineto{\pgfqpoint{3.577753in}{0.789863in}}%
\pgfpathlineto{\pgfqpoint{3.593062in}{0.788459in}}%
\pgfpathlineto{\pgfqpoint{3.613474in}{0.787757in}}%
\pgfpathlineto{\pgfqpoint{3.628783in}{0.786404in}}%
\pgfpathlineto{\pgfqpoint{3.649195in}{0.785702in}}%
\pgfpathlineto{\pgfqpoint{3.659401in}{0.785050in}}%
\pgfpathlineto{\pgfqpoint{3.674710in}{0.784399in}}%
\pgfpathlineto{\pgfqpoint{3.684916in}{0.783747in}}%
\pgfpathlineto{\pgfqpoint{3.705329in}{0.783095in}}%
\pgfpathlineto{\pgfqpoint{3.720638in}{0.781842in}}%
\pgfpathlineto{\pgfqpoint{3.741050in}{0.781240in}}%
\pgfpathlineto{\pgfqpoint{3.756359in}{0.780037in}}%
\pgfpathlineto{\pgfqpoint{3.776771in}{0.779436in}}%
\pgfpathlineto{\pgfqpoint{3.786977in}{0.778834in}}%
\pgfpathlineto{\pgfqpoint{3.802286in}{0.778283in}}%
\pgfpathlineto{\pgfqpoint{3.817595in}{0.777130in}}%
\pgfpathlineto{\pgfqpoint{3.853316in}{0.776027in}}%
\pgfpathlineto{\pgfqpoint{3.878832in}{0.774372in}}%
\pgfpathlineto{\pgfqpoint{3.909450in}{0.773320in}}%
\pgfpathlineto{\pgfqpoint{3.929862in}{0.772317in}}%
\pgfpathlineto{\pgfqpoint{3.955377in}{0.771264in}}%
\pgfpathlineto{\pgfqpoint{3.975789in}{0.770312in}}%
\pgfpathlineto{\pgfqpoint{4.006407in}{0.769309in}}%
\pgfpathlineto{\pgfqpoint{4.026819in}{0.768357in}}%
\pgfpathlineto{\pgfqpoint{4.057438in}{0.767454in}}%
\pgfpathlineto{\pgfqpoint{4.082953in}{0.766051in}}%
\pgfpathlineto{\pgfqpoint{4.108468in}{0.765198in}}%
\pgfpathlineto{\pgfqpoint{4.139086in}{0.764346in}}%
\pgfpathlineto{\pgfqpoint{4.159498in}{0.763494in}}%
\pgfpathlineto{\pgfqpoint{4.190116in}{0.762642in}}%
\pgfpathlineto{\pgfqpoint{4.215631in}{0.761438in}}%
\pgfpathlineto{\pgfqpoint{4.246250in}{0.760636in}}%
\pgfpathlineto{\pgfqpoint{4.266662in}{0.759834in}}%
\pgfpathlineto{\pgfqpoint{4.297280in}{0.759082in}}%
\pgfpathlineto{\pgfqpoint{4.322795in}{0.757929in}}%
\pgfpathlineto{\pgfqpoint{4.368722in}{0.756877in}}%
\pgfpathlineto{\pgfqpoint{4.394237in}{0.755774in}}%
\pgfpathlineto{\pgfqpoint{4.440165in}{0.754721in}}%
\pgfpathlineto{\pgfqpoint{4.480989in}{0.753417in}}%
\pgfpathlineto{\pgfqpoint{4.521813in}{0.752415in}}%
\pgfpathlineto{\pgfqpoint{4.552431in}{0.751462in}}%
\pgfpathlineto{\pgfqpoint{4.598359in}{0.750560in}}%
\pgfpathlineto{\pgfqpoint{4.628977in}{0.749658in}}%
\pgfpathlineto{\pgfqpoint{4.669801in}{0.748755in}}%
\pgfpathlineto{\pgfqpoint{4.700419in}{0.747903in}}%
\pgfpathlineto{\pgfqpoint{4.741243in}{0.747051in}}%
\pgfpathlineto{\pgfqpoint{4.771862in}{0.746199in}}%
\pgfpathlineto{\pgfqpoint{4.827995in}{0.745146in}}%
\pgfpathlineto{\pgfqpoint{4.868819in}{0.744093in}}%
\pgfpathlineto{\pgfqpoint{4.919850in}{0.743090in}}%
\pgfpathlineto{\pgfqpoint{4.960674in}{0.742088in}}%
\pgfpathlineto{\pgfqpoint{5.016807in}{0.741135in}}%
\pgfpathlineto{\pgfqpoint{5.072940in}{0.739982in}}%
\pgfpathlineto{\pgfqpoint{5.144383in}{0.738879in}}%
\pgfpathlineto{\pgfqpoint{5.185207in}{0.738228in}}%
\pgfpathlineto{\pgfqpoint{5.256649in}{0.737175in}}%
\pgfpathlineto{\pgfqpoint{5.302577in}{0.736323in}}%
\pgfpathlineto{\pgfqpoint{5.379122in}{0.735320in}}%
\pgfpathlineto{\pgfqpoint{5.425049in}{0.734568in}}%
\pgfpathlineto{\pgfqpoint{5.491389in}{0.733565in}}%
\pgfpathlineto{\pgfqpoint{5.522007in}{0.733014in}}%
\pgfpathlineto{\pgfqpoint{5.588346in}{0.732112in}}%
\pgfpathlineto{\pgfqpoint{5.588346in}{0.732112in}}%
\pgfusepath{stroke}%
\end{pgfscope}%
\begin{pgfscope}%
\pgfsetrectcap%
\pgfsetmiterjoin%
\pgfsetlinewidth{0.803000pt}%
\definecolor{currentstroke}{rgb}{0.000000,0.000000,0.000000}%
\pgfsetstrokecolor{currentstroke}%
\pgfsetdash{}{0pt}%
\pgfpathmoveto{\pgfqpoint{0.834028in}{0.582778in}}%
\pgfpathlineto{\pgfqpoint{0.834028in}{3.020000in}}%
\pgfusepath{stroke}%
\end{pgfscope}%
\begin{pgfscope}%
\pgfsetrectcap%
\pgfsetmiterjoin%
\pgfsetlinewidth{0.803000pt}%
\definecolor{currentstroke}{rgb}{0.000000,0.000000,0.000000}%
\pgfsetstrokecolor{currentstroke}%
\pgfsetdash{}{0pt}%
\pgfpathmoveto{\pgfqpoint{5.593472in}{0.582778in}}%
\pgfpathlineto{\pgfqpoint{5.593472in}{3.020000in}}%
\pgfusepath{stroke}%
\end{pgfscope}%
\begin{pgfscope}%
\pgfsetrectcap%
\pgfsetmiterjoin%
\pgfsetlinewidth{0.803000pt}%
\definecolor{currentstroke}{rgb}{0.000000,0.000000,0.000000}%
\pgfsetstrokecolor{currentstroke}%
\pgfsetdash{}{0pt}%
\pgfpathmoveto{\pgfqpoint{0.834028in}{0.582778in}}%
\pgfpathlineto{\pgfqpoint{5.593472in}{0.582778in}}%
\pgfusepath{stroke}%
\end{pgfscope}%
\begin{pgfscope}%
\pgfsetrectcap%
\pgfsetmiterjoin%
\pgfsetlinewidth{0.803000pt}%
\definecolor{currentstroke}{rgb}{0.000000,0.000000,0.000000}%
\pgfsetstrokecolor{currentstroke}%
\pgfsetdash{}{0pt}%
\pgfpathmoveto{\pgfqpoint{0.834028in}{3.020000in}}%
\pgfpathlineto{\pgfqpoint{5.593472in}{3.020000in}}%
\pgfusepath{stroke}%
\end{pgfscope}%
\begin{pgfscope}%
\pgfsetbuttcap%
\pgfsetmiterjoin%
\definecolor{currentfill}{rgb}{1.000000,1.000000,1.000000}%
\pgfsetfillcolor{currentfill}%
\pgfsetfillopacity{0.800000}%
\pgfsetlinewidth{1.003750pt}%
\definecolor{currentstroke}{rgb}{0.800000,0.800000,0.800000}%
\pgfsetstrokecolor{currentstroke}%
\pgfsetstrokeopacity{0.800000}%
\pgfsetdash{}{0pt}%
\pgfpathmoveto{\pgfqpoint{4.637106in}{2.496699in}}%
\pgfpathlineto{\pgfqpoint{5.496250in}{2.496699in}}%
\pgfpathquadraticcurveto{\pgfqpoint{5.524028in}{2.496699in}}{\pgfqpoint{5.524028in}{2.524476in}}%
\pgfpathlineto{\pgfqpoint{5.524028in}{2.922778in}}%
\pgfpathquadraticcurveto{\pgfqpoint{5.524028in}{2.950556in}}{\pgfqpoint{5.496250in}{2.950556in}}%
\pgfpathlineto{\pgfqpoint{4.637106in}{2.950556in}}%
\pgfpathquadraticcurveto{\pgfqpoint{4.609328in}{2.950556in}}{\pgfqpoint{4.609328in}{2.922778in}}%
\pgfpathlineto{\pgfqpoint{4.609328in}{2.524476in}}%
\pgfpathquadraticcurveto{\pgfqpoint{4.609328in}{2.496699in}}{\pgfqpoint{4.637106in}{2.496699in}}%
\pgfpathlineto{\pgfqpoint{4.637106in}{2.496699in}}%
\pgfpathclose%
\pgfusepath{stroke,fill}%
\end{pgfscope}%
\begin{pgfscope}%
\pgfsetrectcap%
\pgfsetroundjoin%
\pgfsetlinewidth{2.007500pt}%
\definecolor{currentstroke}{rgb}{0.000000,0.000000,0.000000}%
\pgfsetstrokecolor{currentstroke}%
\pgfsetdash{}{0pt}%
\pgfpathmoveto{\pgfqpoint{4.664883in}{2.837478in}}%
\pgfpathlineto{\pgfqpoint{4.803772in}{2.837478in}}%
\pgfpathlineto{\pgfqpoint{4.942661in}{2.837478in}}%
\pgfusepath{stroke}%
\end{pgfscope}%
\begin{pgfscope}%
\definecolor{textcolor}{rgb}{0.000000,0.000000,0.000000}%
\pgfsetstrokecolor{textcolor}%
\pgfsetfillcolor{textcolor}%
\pgftext[x=5.053772in,y=2.788867in,left,base]{\color{textcolor}\sffamily\fontsize{10.000000}{12.000000}\selectfont |Z|}%
\end{pgfscope}%
\begin{pgfscope}%
\pgfsetbuttcap%
\pgfsetroundjoin%
\pgfsetlinewidth{2.007500pt}%
\definecolor{currentstroke}{rgb}{0.000000,0.000000,0.000000}%
\pgfsetstrokecolor{currentstroke}%
\pgfsetdash{{7.400000pt}{3.200000pt}}{0.000000pt}%
\pgfpathmoveto{\pgfqpoint{4.664883in}{2.629755in}}%
\pgfpathlineto{\pgfqpoint{4.803772in}{2.629755in}}%
\pgfpathlineto{\pgfqpoint{4.942661in}{2.629755in}}%
\pgfusepath{stroke}%
\end{pgfscope}%
\begin{pgfscope}%
\definecolor{textcolor}{rgb}{0.000000,0.000000,0.000000}%
\pgfsetstrokecolor{textcolor}%
\pgfsetfillcolor{textcolor}%
\pgftext[x=5.053772in,y=2.581144in,left,base]{\color{textcolor}\sffamily\fontsize{10.000000}{12.000000}\selectfont Phase}%
\end{pgfscope}%
\begin{pgfscope}%
\pgfsetbuttcap%
\pgfsetroundjoin%
\definecolor{currentfill}{rgb}{0.000000,0.000000,0.000000}%
\pgfsetfillcolor{currentfill}%
\pgfsetlinewidth{0.803000pt}%
\definecolor{currentstroke}{rgb}{0.000000,0.000000,0.000000}%
\pgfsetstrokecolor{currentstroke}%
\pgfsetdash{}{0pt}%
\pgfsys@defobject{currentmarker}{\pgfqpoint{0.000000in}{0.000000in}}{\pgfqpoint{0.048611in}{0.000000in}}{%
\pgfpathmoveto{\pgfqpoint{0.000000in}{0.000000in}}%
\pgfpathlineto{\pgfqpoint{0.048611in}{0.000000in}}%
\pgfusepath{stroke,fill}%
}%
\begin{pgfscope}%
\pgfsys@transformshift{5.593472in}{0.673954in}%
\pgfsys@useobject{currentmarker}{}%
\end{pgfscope}%
\end{pgfscope}%
\begin{pgfscope}%
\definecolor{textcolor}{rgb}{0.000000,0.000000,0.000000}%
\pgfsetstrokecolor{textcolor}%
\pgfsetfillcolor{textcolor}%
\pgftext[x=5.690694in, y=0.621193in, left, base]{\color{textcolor}\sffamily\fontsize{10.000000}{12.000000}\selectfont 0}%
\end{pgfscope}%
\begin{pgfscope}%
\pgfsetbuttcap%
\pgfsetroundjoin%
\definecolor{currentfill}{rgb}{0.000000,0.000000,0.000000}%
\pgfsetfillcolor{currentfill}%
\pgfsetlinewidth{0.803000pt}%
\definecolor{currentstroke}{rgb}{0.000000,0.000000,0.000000}%
\pgfsetstrokecolor{currentstroke}%
\pgfsetdash{}{0pt}%
\pgfsys@defobject{currentmarker}{\pgfqpoint{0.000000in}{0.000000in}}{\pgfqpoint{0.048611in}{0.000000in}}{%
\pgfpathmoveto{\pgfqpoint{0.000000in}{0.000000in}}%
\pgfpathlineto{\pgfqpoint{0.048611in}{0.000000in}}%
\pgfusepath{stroke,fill}%
}%
\begin{pgfscope}%
\pgfsys@transformshift{5.593472in}{0.992754in}%
\pgfsys@useobject{currentmarker}{}%
\end{pgfscope}%
\end{pgfscope}%
\begin{pgfscope}%
\definecolor{textcolor}{rgb}{0.000000,0.000000,0.000000}%
\pgfsetstrokecolor{textcolor}%
\pgfsetfillcolor{textcolor}%
\pgftext[x=5.690694in, y=0.939992in, left, base]{\color{textcolor}\sffamily\fontsize{10.000000}{12.000000}\selectfont 20}%
\end{pgfscope}%
\begin{pgfscope}%
\pgfsetbuttcap%
\pgfsetroundjoin%
\definecolor{currentfill}{rgb}{0.000000,0.000000,0.000000}%
\pgfsetfillcolor{currentfill}%
\pgfsetlinewidth{0.803000pt}%
\definecolor{currentstroke}{rgb}{0.000000,0.000000,0.000000}%
\pgfsetstrokecolor{currentstroke}%
\pgfsetdash{}{0pt}%
\pgfsys@defobject{currentmarker}{\pgfqpoint{0.000000in}{0.000000in}}{\pgfqpoint{0.048611in}{0.000000in}}{%
\pgfpathmoveto{\pgfqpoint{0.000000in}{0.000000in}}%
\pgfpathlineto{\pgfqpoint{0.048611in}{0.000000in}}%
\pgfusepath{stroke,fill}%
}%
\begin{pgfscope}%
\pgfsys@transformshift{5.593472in}{1.311553in}%
\pgfsys@useobject{currentmarker}{}%
\end{pgfscope}%
\end{pgfscope}%
\begin{pgfscope}%
\definecolor{textcolor}{rgb}{0.000000,0.000000,0.000000}%
\pgfsetstrokecolor{textcolor}%
\pgfsetfillcolor{textcolor}%
\pgftext[x=5.690694in, y=1.258792in, left, base]{\color{textcolor}\sffamily\fontsize{10.000000}{12.000000}\selectfont 40}%
\end{pgfscope}%
\begin{pgfscope}%
\pgfsetbuttcap%
\pgfsetroundjoin%
\definecolor{currentfill}{rgb}{0.000000,0.000000,0.000000}%
\pgfsetfillcolor{currentfill}%
\pgfsetlinewidth{0.803000pt}%
\definecolor{currentstroke}{rgb}{0.000000,0.000000,0.000000}%
\pgfsetstrokecolor{currentstroke}%
\pgfsetdash{}{0pt}%
\pgfsys@defobject{currentmarker}{\pgfqpoint{0.000000in}{0.000000in}}{\pgfqpoint{0.048611in}{0.000000in}}{%
\pgfpathmoveto{\pgfqpoint{0.000000in}{0.000000in}}%
\pgfpathlineto{\pgfqpoint{0.048611in}{0.000000in}}%
\pgfusepath{stroke,fill}%
}%
\begin{pgfscope}%
\pgfsys@transformshift{5.593472in}{1.630353in}%
\pgfsys@useobject{currentmarker}{}%
\end{pgfscope}%
\end{pgfscope}%
\begin{pgfscope}%
\definecolor{textcolor}{rgb}{0.000000,0.000000,0.000000}%
\pgfsetstrokecolor{textcolor}%
\pgfsetfillcolor{textcolor}%
\pgftext[x=5.690694in, y=1.577591in, left, base]{\color{textcolor}\sffamily\fontsize{10.000000}{12.000000}\selectfont 60}%
\end{pgfscope}%
\begin{pgfscope}%
\pgfsetbuttcap%
\pgfsetroundjoin%
\definecolor{currentfill}{rgb}{0.000000,0.000000,0.000000}%
\pgfsetfillcolor{currentfill}%
\pgfsetlinewidth{0.803000pt}%
\definecolor{currentstroke}{rgb}{0.000000,0.000000,0.000000}%
\pgfsetstrokecolor{currentstroke}%
\pgfsetdash{}{0pt}%
\pgfsys@defobject{currentmarker}{\pgfqpoint{0.000000in}{0.000000in}}{\pgfqpoint{0.048611in}{0.000000in}}{%
\pgfpathmoveto{\pgfqpoint{0.000000in}{0.000000in}}%
\pgfpathlineto{\pgfqpoint{0.048611in}{0.000000in}}%
\pgfusepath{stroke,fill}%
}%
\begin{pgfscope}%
\pgfsys@transformshift{5.593472in}{1.949152in}%
\pgfsys@useobject{currentmarker}{}%
\end{pgfscope}%
\end{pgfscope}%
\begin{pgfscope}%
\definecolor{textcolor}{rgb}{0.000000,0.000000,0.000000}%
\pgfsetstrokecolor{textcolor}%
\pgfsetfillcolor{textcolor}%
\pgftext[x=5.690694in, y=1.896391in, left, base]{\color{textcolor}\sffamily\fontsize{10.000000}{12.000000}\selectfont 80}%
\end{pgfscope}%
\begin{pgfscope}%
\pgfsetbuttcap%
\pgfsetroundjoin%
\definecolor{currentfill}{rgb}{0.000000,0.000000,0.000000}%
\pgfsetfillcolor{currentfill}%
\pgfsetlinewidth{0.803000pt}%
\definecolor{currentstroke}{rgb}{0.000000,0.000000,0.000000}%
\pgfsetstrokecolor{currentstroke}%
\pgfsetdash{}{0pt}%
\pgfsys@defobject{currentmarker}{\pgfqpoint{0.000000in}{0.000000in}}{\pgfqpoint{0.048611in}{0.000000in}}{%
\pgfpathmoveto{\pgfqpoint{0.000000in}{0.000000in}}%
\pgfpathlineto{\pgfqpoint{0.048611in}{0.000000in}}%
\pgfusepath{stroke,fill}%
}%
\begin{pgfscope}%
\pgfsys@transformshift{5.593472in}{2.267952in}%
\pgfsys@useobject{currentmarker}{}%
\end{pgfscope}%
\end{pgfscope}%
\begin{pgfscope}%
\definecolor{textcolor}{rgb}{0.000000,0.000000,0.000000}%
\pgfsetstrokecolor{textcolor}%
\pgfsetfillcolor{textcolor}%
\pgftext[x=5.690694in, y=2.215190in, left, base]{\color{textcolor}\sffamily\fontsize{10.000000}{12.000000}\selectfont 100}%
\end{pgfscope}%
\begin{pgfscope}%
\pgfsetbuttcap%
\pgfsetroundjoin%
\definecolor{currentfill}{rgb}{0.000000,0.000000,0.000000}%
\pgfsetfillcolor{currentfill}%
\pgfsetlinewidth{0.803000pt}%
\definecolor{currentstroke}{rgb}{0.000000,0.000000,0.000000}%
\pgfsetstrokecolor{currentstroke}%
\pgfsetdash{}{0pt}%
\pgfsys@defobject{currentmarker}{\pgfqpoint{0.000000in}{0.000000in}}{\pgfqpoint{0.048611in}{0.000000in}}{%
\pgfpathmoveto{\pgfqpoint{0.000000in}{0.000000in}}%
\pgfpathlineto{\pgfqpoint{0.048611in}{0.000000in}}%
\pgfusepath{stroke,fill}%
}%
\begin{pgfscope}%
\pgfsys@transformshift{5.593472in}{2.586751in}%
\pgfsys@useobject{currentmarker}{}%
\end{pgfscope}%
\end{pgfscope}%
\begin{pgfscope}%
\definecolor{textcolor}{rgb}{0.000000,0.000000,0.000000}%
\pgfsetstrokecolor{textcolor}%
\pgfsetfillcolor{textcolor}%
\pgftext[x=5.690694in, y=2.533990in, left, base]{\color{textcolor}\sffamily\fontsize{10.000000}{12.000000}\selectfont 120}%
\end{pgfscope}%
\begin{pgfscope}%
\pgfsetbuttcap%
\pgfsetroundjoin%
\definecolor{currentfill}{rgb}{0.000000,0.000000,0.000000}%
\pgfsetfillcolor{currentfill}%
\pgfsetlinewidth{0.803000pt}%
\definecolor{currentstroke}{rgb}{0.000000,0.000000,0.000000}%
\pgfsetstrokecolor{currentstroke}%
\pgfsetdash{}{0pt}%
\pgfsys@defobject{currentmarker}{\pgfqpoint{0.000000in}{0.000000in}}{\pgfqpoint{0.048611in}{0.000000in}}{%
\pgfpathmoveto{\pgfqpoint{0.000000in}{0.000000in}}%
\pgfpathlineto{\pgfqpoint{0.048611in}{0.000000in}}%
\pgfusepath{stroke,fill}%
}%
\begin{pgfscope}%
\pgfsys@transformshift{5.593472in}{2.905551in}%
\pgfsys@useobject{currentmarker}{}%
\end{pgfscope}%
\end{pgfscope}%
\begin{pgfscope}%
\definecolor{textcolor}{rgb}{0.000000,0.000000,0.000000}%
\pgfsetstrokecolor{textcolor}%
\pgfsetfillcolor{textcolor}%
\pgftext[x=5.690694in, y=2.852789in, left, base]{\color{textcolor}\sffamily\fontsize{10.000000}{12.000000}\selectfont 140}%
\end{pgfscope}%
\begin{pgfscope}%
\definecolor{textcolor}{rgb}{0.000000,0.000000,0.000000}%
\pgfsetstrokecolor{textcolor}%
\pgfsetfillcolor{textcolor}%
\pgftext[x=6.011346in,y=1.801389in,,top,rotate=90.000000]{\color{textcolor}\sffamily\fontsize{10.000000}{12.000000}\selectfont Phase (deg)}%
\end{pgfscope}%
\begin{pgfscope}%
\pgfpathrectangle{\pgfqpoint{0.834028in}{0.582778in}}{\pgfqpoint{4.759444in}{2.437222in}}%
\pgfusepath{clip}%
\pgfsetbuttcap%
\pgfsetroundjoin%
\pgfsetlinewidth{2.007500pt}%
\definecolor{currentstroke}{rgb}{0.000000,0.000000,0.000000}%
\pgfsetstrokecolor{currentstroke}%
\pgfsetdash{{7.400000pt}{3.200000pt}}{0.000000pt}%
\pgfpathmoveto{\pgfqpoint{0.857838in}{2.404398in}}%
\pgfpathlineto{\pgfqpoint{0.862941in}{2.909217in}}%
\pgfpathlineto{\pgfqpoint{0.883353in}{2.561566in}}%
\pgfpathlineto{\pgfqpoint{0.898662in}{2.337131in}}%
\pgfpathlineto{\pgfqpoint{0.919075in}{2.090381in}}%
\pgfpathlineto{\pgfqpoint{0.929281in}{1.983742in}}%
\pgfpathlineto{\pgfqpoint{0.944590in}{1.846340in}}%
\pgfpathlineto{\pgfqpoint{0.959899in}{1.728543in}}%
\pgfpathlineto{\pgfqpoint{0.975208in}{1.624774in}}%
\pgfpathlineto{\pgfqpoint{0.990517in}{1.537742in}}%
\pgfpathlineto{\pgfqpoint{1.005826in}{1.459317in}}%
\pgfpathlineto{\pgfqpoint{1.026238in}{1.372285in}}%
\pgfpathlineto{\pgfqpoint{1.046650in}{1.296570in}}%
\pgfpathlineto{\pgfqpoint{1.067062in}{1.232013in}}%
\pgfpathlineto{\pgfqpoint{1.092578in}{1.164746in}}%
\pgfpathlineto{\pgfqpoint{1.097681in}{1.153588in}}%
\pgfpathlineto{\pgfqpoint{1.107887in}{1.128244in}}%
\pgfpathlineto{\pgfqpoint{1.118093in}{1.105928in}}%
\pgfpathlineto{\pgfqpoint{1.123196in}{1.097480in}}%
\pgfpathlineto{\pgfqpoint{1.143608in}{1.055398in}}%
\pgfpathlineto{\pgfqpoint{1.153814in}{1.038502in}}%
\pgfpathlineto{\pgfqpoint{1.158917in}{1.027344in}}%
\pgfpathlineto{\pgfqpoint{1.164020in}{1.021765in}}%
\pgfpathlineto{\pgfqpoint{1.194638in}{0.971235in}}%
\pgfpathlineto{\pgfqpoint{1.199741in}{0.965656in}}%
\pgfpathlineto{\pgfqpoint{1.209947in}{0.948760in}}%
\pgfpathlineto{\pgfqpoint{1.220153in}{0.937602in}}%
\pgfpathlineto{\pgfqpoint{1.225256in}{0.929153in}}%
\pgfpathlineto{\pgfqpoint{1.230359in}{0.923574in}}%
\pgfpathlineto{\pgfqpoint{1.235462in}{0.915126in}}%
\pgfpathlineto{\pgfqpoint{1.260978in}{0.887072in}}%
\pgfpathlineto{\pgfqpoint{1.266081in}{0.878624in}}%
\pgfpathlineto{\pgfqpoint{1.271184in}{0.873045in}}%
\pgfpathlineto{\pgfqpoint{1.276287in}{0.870335in}}%
\pgfpathlineto{\pgfqpoint{1.301802in}{0.842281in}}%
\pgfpathlineto{\pgfqpoint{1.306905in}{0.839411in}}%
\pgfpathlineto{\pgfqpoint{1.317111in}{0.828253in}}%
\pgfpathlineto{\pgfqpoint{1.322214in}{0.825384in}}%
\pgfpathlineto{\pgfqpoint{1.327317in}{0.819805in}}%
\pgfpathlineto{\pgfqpoint{1.332420in}{0.816936in}}%
\pgfpathlineto{\pgfqpoint{1.342626in}{0.805778in}}%
\pgfpathlineto{\pgfqpoint{1.352832in}{0.800199in}}%
\pgfpathlineto{\pgfqpoint{1.357935in}{0.794620in}}%
\pgfpathlineto{\pgfqpoint{1.363038in}{0.791751in}}%
\pgfpathlineto{\pgfqpoint{1.368141in}{0.786172in}}%
\pgfpathlineto{\pgfqpoint{1.373244in}{0.783303in}}%
\pgfpathlineto{\pgfqpoint{1.378347in}{0.777724in}}%
\pgfpathlineto{\pgfqpoint{1.388553in}{0.772145in}}%
\pgfpathlineto{\pgfqpoint{1.393656in}{0.766566in}}%
\pgfpathlineto{\pgfqpoint{1.419171in}{0.752539in}}%
\pgfpathlineto{\pgfqpoint{1.424274in}{0.746800in}}%
\pgfpathlineto{\pgfqpoint{1.459996in}{0.727194in}}%
\pgfpathlineto{\pgfqpoint{1.465099in}{0.727194in}}%
\pgfpathlineto{\pgfqpoint{1.475305in}{0.721615in}}%
\pgfpathlineto{\pgfqpoint{1.480408in}{0.721615in}}%
\pgfpathlineto{\pgfqpoint{1.490614in}{0.716036in}}%
\pgfpathlineto{\pgfqpoint{1.495717in}{0.716036in}}%
\pgfpathlineto{\pgfqpoint{1.500820in}{0.713167in}}%
\pgfpathlineto{\pgfqpoint{1.511026in}{0.713167in}}%
\pgfpathlineto{\pgfqpoint{1.521232in}{0.707588in}}%
\pgfpathlineto{\pgfqpoint{1.526335in}{0.707588in}}%
\pgfpathlineto{\pgfqpoint{1.531438in}{0.704878in}}%
\pgfpathlineto{\pgfqpoint{1.541644in}{0.704878in}}%
\pgfpathlineto{\pgfqpoint{1.546747in}{0.702009in}}%
\pgfpathlineto{\pgfqpoint{1.551850in}{0.704878in}}%
\pgfpathlineto{\pgfqpoint{1.562056in}{0.699140in}}%
\pgfpathlineto{\pgfqpoint{1.582468in}{0.699140in}}%
\pgfpathlineto{\pgfqpoint{1.587571in}{0.696430in}}%
\pgfpathlineto{\pgfqpoint{1.633499in}{0.696430in}}%
\pgfpathlineto{\pgfqpoint{1.638602in}{0.693561in}}%
\pgfpathlineto{\pgfqpoint{1.766177in}{0.693561in}}%
\pgfpathlineto{\pgfqpoint{1.771280in}{0.696430in}}%
\pgfpathlineto{\pgfqpoint{1.827414in}{0.696430in}}%
\pgfpathlineto{\pgfqpoint{1.832517in}{0.699140in}}%
\pgfpathlineto{\pgfqpoint{1.868238in}{0.699140in}}%
\pgfpathlineto{\pgfqpoint{1.873341in}{0.702009in}}%
\pgfpathlineto{\pgfqpoint{1.888650in}{0.702009in}}%
\pgfpathlineto{\pgfqpoint{1.893753in}{0.704878in}}%
\pgfpathlineto{\pgfqpoint{1.914165in}{0.704878in}}%
\pgfpathlineto{\pgfqpoint{1.919268in}{0.707588in}}%
\pgfpathlineto{\pgfqpoint{1.934577in}{0.707588in}}%
\pgfpathlineto{\pgfqpoint{1.939680in}{0.710457in}}%
\pgfpathlineto{\pgfqpoint{1.949886in}{0.710457in}}%
\pgfpathlineto{\pgfqpoint{1.954990in}{0.713167in}}%
\pgfpathlineto{\pgfqpoint{1.965196in}{0.713167in}}%
\pgfpathlineto{\pgfqpoint{1.970299in}{0.716036in}}%
\pgfpathlineto{\pgfqpoint{1.980505in}{0.716036in}}%
\pgfpathlineto{\pgfqpoint{1.985608in}{0.718905in}}%
\pgfpathlineto{\pgfqpoint{1.995814in}{0.718905in}}%
\pgfpathlineto{\pgfqpoint{2.000917in}{0.721615in}}%
\pgfpathlineto{\pgfqpoint{2.006020in}{0.721615in}}%
\pgfpathlineto{\pgfqpoint{2.011123in}{0.724484in}}%
\pgfpathlineto{\pgfqpoint{2.021329in}{0.724484in}}%
\pgfpathlineto{\pgfqpoint{2.026432in}{0.727194in}}%
\pgfpathlineto{\pgfqpoint{2.031535in}{0.727194in}}%
\pgfpathlineto{\pgfqpoint{2.036638in}{0.730063in}}%
\pgfpathlineto{\pgfqpoint{2.046844in}{0.730063in}}%
\pgfpathlineto{\pgfqpoint{2.051947in}{0.732932in}}%
\pgfpathlineto{\pgfqpoint{2.057050in}{0.732932in}}%
\pgfpathlineto{\pgfqpoint{2.062153in}{0.735642in}}%
\pgfpathlineto{\pgfqpoint{2.067256in}{0.735642in}}%
\pgfpathlineto{\pgfqpoint{2.072359in}{0.738511in}}%
\pgfpathlineto{\pgfqpoint{2.077462in}{0.738511in}}%
\pgfpathlineto{\pgfqpoint{2.082565in}{0.741221in}}%
\pgfpathlineto{\pgfqpoint{2.087668in}{0.741221in}}%
\pgfpathlineto{\pgfqpoint{2.092771in}{0.744090in}}%
\pgfpathlineto{\pgfqpoint{2.097874in}{0.744090in}}%
\pgfpathlineto{\pgfqpoint{2.102977in}{0.746800in}}%
\pgfpathlineto{\pgfqpoint{2.113183in}{0.746800in}}%
\pgfpathlineto{\pgfqpoint{2.118286in}{0.749669in}}%
\pgfpathlineto{\pgfqpoint{2.123389in}{0.749669in}}%
\pgfpathlineto{\pgfqpoint{2.128493in}{0.752539in}}%
\pgfpathlineto{\pgfqpoint{2.133596in}{0.752539in}}%
\pgfpathlineto{\pgfqpoint{2.138699in}{0.755248in}}%
\pgfpathlineto{\pgfqpoint{2.143802in}{0.755248in}}%
\pgfpathlineto{\pgfqpoint{2.148905in}{0.758118in}}%
\pgfpathlineto{\pgfqpoint{2.154008in}{0.758118in}}%
\pgfpathlineto{\pgfqpoint{2.159111in}{0.760827in}}%
\pgfpathlineto{\pgfqpoint{2.164214in}{0.760827in}}%
\pgfpathlineto{\pgfqpoint{2.169317in}{0.763696in}}%
\pgfpathlineto{\pgfqpoint{2.174420in}{0.763696in}}%
\pgfpathlineto{\pgfqpoint{2.179523in}{0.766566in}}%
\pgfpathlineto{\pgfqpoint{2.184626in}{0.766566in}}%
\pgfpathlineto{\pgfqpoint{2.189729in}{0.769275in}}%
\pgfpathlineto{\pgfqpoint{2.194832in}{0.769275in}}%
\pgfpathlineto{\pgfqpoint{2.199935in}{0.772145in}}%
\pgfpathlineto{\pgfqpoint{2.210141in}{0.772145in}}%
\pgfpathlineto{\pgfqpoint{2.215244in}{0.774854in}}%
\pgfpathlineto{\pgfqpoint{2.220347in}{0.774854in}}%
\pgfpathlineto{\pgfqpoint{2.230553in}{0.780593in}}%
\pgfpathlineto{\pgfqpoint{2.235656in}{0.780593in}}%
\pgfpathlineto{\pgfqpoint{2.240759in}{0.783303in}}%
\pgfpathlineto{\pgfqpoint{2.245862in}{0.783303in}}%
\pgfpathlineto{\pgfqpoint{2.250965in}{0.786172in}}%
\pgfpathlineto{\pgfqpoint{2.261171in}{0.786172in}}%
\pgfpathlineto{\pgfqpoint{2.266274in}{0.788882in}}%
\pgfpathlineto{\pgfqpoint{2.271377in}{0.788882in}}%
\pgfpathlineto{\pgfqpoint{2.281583in}{0.794620in}}%
\pgfpathlineto{\pgfqpoint{2.291789in}{0.794620in}}%
\pgfpathlineto{\pgfqpoint{2.296892in}{0.797330in}}%
\pgfpathlineto{\pgfqpoint{2.301996in}{0.797330in}}%
\pgfpathlineto{\pgfqpoint{2.307099in}{0.800199in}}%
\pgfpathlineto{\pgfqpoint{2.312202in}{0.800199in}}%
\pgfpathlineto{\pgfqpoint{2.322408in}{0.805778in}}%
\pgfpathlineto{\pgfqpoint{2.332614in}{0.805778in}}%
\pgfpathlineto{\pgfqpoint{2.337717in}{0.808647in}}%
\pgfpathlineto{\pgfqpoint{2.342820in}{0.808647in}}%
\pgfpathlineto{\pgfqpoint{2.347923in}{0.811357in}}%
\pgfpathlineto{\pgfqpoint{2.353026in}{0.811357in}}%
\pgfpathlineto{\pgfqpoint{2.358129in}{0.814226in}}%
\pgfpathlineto{\pgfqpoint{2.363232in}{0.814226in}}%
\pgfpathlineto{\pgfqpoint{2.368335in}{0.816936in}}%
\pgfpathlineto{\pgfqpoint{2.373438in}{0.816936in}}%
\pgfpathlineto{\pgfqpoint{2.378541in}{0.819805in}}%
\pgfpathlineto{\pgfqpoint{2.383644in}{0.819805in}}%
\pgfpathlineto{\pgfqpoint{2.388747in}{0.822674in}}%
\pgfpathlineto{\pgfqpoint{2.393850in}{0.830963in}}%
\pgfpathlineto{\pgfqpoint{2.398953in}{0.825384in}}%
\pgfpathlineto{\pgfqpoint{2.404056in}{0.825384in}}%
\pgfpathlineto{\pgfqpoint{2.409159in}{0.828253in}}%
\pgfpathlineto{\pgfqpoint{2.414262in}{0.828253in}}%
\pgfpathlineto{\pgfqpoint{2.419365in}{0.830963in}}%
\pgfpathlineto{\pgfqpoint{2.424468in}{0.830963in}}%
\pgfpathlineto{\pgfqpoint{2.429571in}{0.833832in}}%
\pgfpathlineto{\pgfqpoint{2.434674in}{0.833832in}}%
\pgfpathlineto{\pgfqpoint{2.439777in}{0.836702in}}%
\pgfpathlineto{\pgfqpoint{2.444880in}{0.836702in}}%
\pgfpathlineto{\pgfqpoint{2.449983in}{0.839411in}}%
\pgfpathlineto{\pgfqpoint{2.460189in}{0.839411in}}%
\pgfpathlineto{\pgfqpoint{2.465292in}{0.842281in}}%
\pgfpathlineto{\pgfqpoint{2.470395in}{0.842281in}}%
\pgfpathlineto{\pgfqpoint{2.475498in}{0.844990in}}%
\pgfpathlineto{\pgfqpoint{2.480602in}{0.844990in}}%
\pgfpathlineto{\pgfqpoint{2.485705in}{0.847860in}}%
\pgfpathlineto{\pgfqpoint{2.490808in}{0.847860in}}%
\pgfpathlineto{\pgfqpoint{2.495911in}{0.850729in}}%
\pgfpathlineto{\pgfqpoint{2.501014in}{0.850729in}}%
\pgfpathlineto{\pgfqpoint{2.506117in}{0.853439in}}%
\pgfpathlineto{\pgfqpoint{2.511220in}{0.853439in}}%
\pgfpathlineto{\pgfqpoint{2.516323in}{0.856308in}}%
\pgfpathlineto{\pgfqpoint{2.521426in}{0.856308in}}%
\pgfpathlineto{\pgfqpoint{2.526529in}{0.859018in}}%
\pgfpathlineto{\pgfqpoint{2.531632in}{0.859018in}}%
\pgfpathlineto{\pgfqpoint{2.536735in}{0.861887in}}%
\pgfpathlineto{\pgfqpoint{2.546941in}{0.861887in}}%
\pgfpathlineto{\pgfqpoint{2.552044in}{0.864597in}}%
\pgfpathlineto{\pgfqpoint{2.557147in}{0.864597in}}%
\pgfpathlineto{\pgfqpoint{2.562250in}{0.867466in}}%
\pgfpathlineto{\pgfqpoint{2.572456in}{0.867466in}}%
\pgfpathlineto{\pgfqpoint{2.577559in}{0.870335in}}%
\pgfpathlineto{\pgfqpoint{2.587765in}{0.870335in}}%
\pgfpathlineto{\pgfqpoint{2.592868in}{0.873045in}}%
\pgfpathlineto{\pgfqpoint{2.603074in}{0.873045in}}%
\pgfpathlineto{\pgfqpoint{2.608177in}{0.875914in}}%
\pgfpathlineto{\pgfqpoint{2.618383in}{0.875914in}}%
\pgfpathlineto{\pgfqpoint{2.623486in}{0.878624in}}%
\pgfpathlineto{\pgfqpoint{2.628589in}{0.878624in}}%
\pgfpathlineto{\pgfqpoint{2.633692in}{0.881493in}}%
\pgfpathlineto{\pgfqpoint{2.638795in}{0.881493in}}%
\pgfpathlineto{\pgfqpoint{2.643898in}{0.884362in}}%
\pgfpathlineto{\pgfqpoint{2.649001in}{0.884362in}}%
\pgfpathlineto{\pgfqpoint{2.654105in}{0.887072in}}%
\pgfpathlineto{\pgfqpoint{2.659208in}{0.887072in}}%
\pgfpathlineto{\pgfqpoint{2.664311in}{0.889941in}}%
\pgfpathlineto{\pgfqpoint{2.674517in}{0.889941in}}%
\pgfpathlineto{\pgfqpoint{2.684723in}{0.895520in}}%
\pgfpathlineto{\pgfqpoint{2.694929in}{0.895520in}}%
\pgfpathlineto{\pgfqpoint{2.700032in}{0.898389in}}%
\pgfpathlineto{\pgfqpoint{2.705135in}{0.898389in}}%
\pgfpathlineto{\pgfqpoint{2.710238in}{0.901099in}}%
\pgfpathlineto{\pgfqpoint{2.715341in}{0.901099in}}%
\pgfpathlineto{\pgfqpoint{2.720444in}{0.903968in}}%
\pgfpathlineto{\pgfqpoint{2.725547in}{0.903968in}}%
\pgfpathlineto{\pgfqpoint{2.730650in}{0.906678in}}%
\pgfpathlineto{\pgfqpoint{2.735753in}{0.906678in}}%
\pgfpathlineto{\pgfqpoint{2.740856in}{0.909547in}}%
\pgfpathlineto{\pgfqpoint{2.745959in}{0.909547in}}%
\pgfpathlineto{\pgfqpoint{2.751062in}{0.912416in}}%
\pgfpathlineto{\pgfqpoint{2.756165in}{0.912416in}}%
\pgfpathlineto{\pgfqpoint{2.761268in}{0.915126in}}%
\pgfpathlineto{\pgfqpoint{2.771474in}{0.915126in}}%
\pgfpathlineto{\pgfqpoint{2.776577in}{0.917995in}}%
\pgfpathlineto{\pgfqpoint{2.781680in}{0.917995in}}%
\pgfpathlineto{\pgfqpoint{2.786783in}{0.920705in}}%
\pgfpathlineto{\pgfqpoint{2.791886in}{0.920705in}}%
\pgfpathlineto{\pgfqpoint{2.796989in}{0.923574in}}%
\pgfpathlineto{\pgfqpoint{2.807195in}{0.923574in}}%
\pgfpathlineto{\pgfqpoint{2.812298in}{0.926444in}}%
\pgfpathlineto{\pgfqpoint{2.817401in}{0.926444in}}%
\pgfpathlineto{\pgfqpoint{2.822504in}{0.929153in}}%
\pgfpathlineto{\pgfqpoint{2.827608in}{0.929153in}}%
\pgfpathlineto{\pgfqpoint{2.832711in}{0.932023in}}%
\pgfpathlineto{\pgfqpoint{2.842917in}{0.932023in}}%
\pgfpathlineto{\pgfqpoint{2.848020in}{0.934732in}}%
\pgfpathlineto{\pgfqpoint{2.853123in}{0.934732in}}%
\pgfpathlineto{\pgfqpoint{2.858226in}{0.937602in}}%
\pgfpathlineto{\pgfqpoint{2.868432in}{0.937602in}}%
\pgfpathlineto{\pgfqpoint{2.873535in}{0.940471in}}%
\pgfpathlineto{\pgfqpoint{2.878638in}{0.940471in}}%
\pgfpathlineto{\pgfqpoint{2.883741in}{0.943181in}}%
\pgfpathlineto{\pgfqpoint{2.888844in}{0.943181in}}%
\pgfpathlineto{\pgfqpoint{2.893947in}{0.946050in}}%
\pgfpathlineto{\pgfqpoint{2.904153in}{0.946050in}}%
\pgfpathlineto{\pgfqpoint{2.909256in}{0.948760in}}%
\pgfpathlineto{\pgfqpoint{2.914359in}{0.948760in}}%
\pgfpathlineto{\pgfqpoint{2.919462in}{0.951629in}}%
\pgfpathlineto{\pgfqpoint{2.924565in}{0.951629in}}%
\pgfpathlineto{\pgfqpoint{2.929668in}{0.954498in}}%
\pgfpathlineto{\pgfqpoint{2.934771in}{0.954498in}}%
\pgfpathlineto{\pgfqpoint{2.939874in}{0.957208in}}%
\pgfpathlineto{\pgfqpoint{2.950080in}{0.957208in}}%
\pgfpathlineto{\pgfqpoint{2.955183in}{0.960077in}}%
\pgfpathlineto{\pgfqpoint{2.960286in}{0.960077in}}%
\pgfpathlineto{\pgfqpoint{2.965389in}{0.962787in}}%
\pgfpathlineto{\pgfqpoint{2.970492in}{0.962787in}}%
\pgfpathlineto{\pgfqpoint{2.975595in}{0.965656in}}%
\pgfpathlineto{\pgfqpoint{2.980698in}{0.965656in}}%
\pgfpathlineto{\pgfqpoint{2.985801in}{0.968525in}}%
\pgfpathlineto{\pgfqpoint{2.990904in}{0.968525in}}%
\pgfpathlineto{\pgfqpoint{2.996007in}{0.971235in}}%
\pgfpathlineto{\pgfqpoint{3.006214in}{0.971235in}}%
\pgfpathlineto{\pgfqpoint{3.011317in}{0.974104in}}%
\pgfpathlineto{\pgfqpoint{3.016420in}{0.974104in}}%
\pgfpathlineto{\pgfqpoint{3.021523in}{0.976814in}}%
\pgfpathlineto{\pgfqpoint{3.026626in}{0.976814in}}%
\pgfpathlineto{\pgfqpoint{3.031729in}{0.979683in}}%
\pgfpathlineto{\pgfqpoint{3.036832in}{0.979683in}}%
\pgfpathlineto{\pgfqpoint{3.041935in}{0.982393in}}%
\pgfpathlineto{\pgfqpoint{3.052141in}{0.982393in}}%
\pgfpathlineto{\pgfqpoint{3.057244in}{0.985262in}}%
\pgfpathlineto{\pgfqpoint{3.062347in}{0.985262in}}%
\pgfpathlineto{\pgfqpoint{3.067450in}{0.988131in}}%
\pgfpathlineto{\pgfqpoint{3.072553in}{0.988131in}}%
\pgfpathlineto{\pgfqpoint{3.077656in}{0.990841in}}%
\pgfpathlineto{\pgfqpoint{3.087862in}{0.990841in}}%
\pgfpathlineto{\pgfqpoint{3.092965in}{0.993710in}}%
\pgfpathlineto{\pgfqpoint{3.098068in}{0.993710in}}%
\pgfpathlineto{\pgfqpoint{3.103171in}{0.996420in}}%
\pgfpathlineto{\pgfqpoint{3.113377in}{0.996420in}}%
\pgfpathlineto{\pgfqpoint{3.118480in}{0.999289in}}%
\pgfpathlineto{\pgfqpoint{3.123583in}{0.999289in}}%
\pgfpathlineto{\pgfqpoint{3.128686in}{1.002159in}}%
\pgfpathlineto{\pgfqpoint{3.133789in}{1.002159in}}%
\pgfpathlineto{\pgfqpoint{3.138892in}{1.004868in}}%
\pgfpathlineto{\pgfqpoint{3.149098in}{1.004868in}}%
\pgfpathlineto{\pgfqpoint{3.154201in}{1.007738in}}%
\pgfpathlineto{\pgfqpoint{3.159304in}{1.007738in}}%
\pgfpathlineto{\pgfqpoint{3.164407in}{1.010447in}}%
\pgfpathlineto{\pgfqpoint{3.174614in}{1.010447in}}%
\pgfpathlineto{\pgfqpoint{3.179717in}{1.013317in}}%
\pgfpathlineto{\pgfqpoint{3.184820in}{1.013317in}}%
\pgfpathlineto{\pgfqpoint{3.189923in}{1.016186in}}%
\pgfpathlineto{\pgfqpoint{3.195026in}{1.016186in}}%
\pgfpathlineto{\pgfqpoint{3.200129in}{1.018896in}}%
\pgfpathlineto{\pgfqpoint{3.205232in}{1.018896in}}%
\pgfpathlineto{\pgfqpoint{3.210335in}{1.021765in}}%
\pgfpathlineto{\pgfqpoint{3.220541in}{1.021765in}}%
\pgfpathlineto{\pgfqpoint{3.225644in}{1.024474in}}%
\pgfpathlineto{\pgfqpoint{3.230747in}{1.024474in}}%
\pgfpathlineto{\pgfqpoint{3.235850in}{1.027344in}}%
\pgfpathlineto{\pgfqpoint{3.240953in}{1.027344in}}%
\pgfpathlineto{\pgfqpoint{3.246056in}{1.030213in}}%
\pgfpathlineto{\pgfqpoint{3.256262in}{1.030213in}}%
\pgfpathlineto{\pgfqpoint{3.261365in}{1.032923in}}%
\pgfpathlineto{\pgfqpoint{3.271571in}{1.032923in}}%
\pgfpathlineto{\pgfqpoint{3.276674in}{1.035792in}}%
\pgfpathlineto{\pgfqpoint{3.281777in}{1.035792in}}%
\pgfpathlineto{\pgfqpoint{3.286880in}{1.038502in}}%
\pgfpathlineto{\pgfqpoint{3.291983in}{1.038502in}}%
\pgfpathlineto{\pgfqpoint{3.297086in}{1.041371in}}%
\pgfpathlineto{\pgfqpoint{3.302189in}{1.041371in}}%
\pgfpathlineto{\pgfqpoint{3.307292in}{1.044240in}}%
\pgfpathlineto{\pgfqpoint{3.317498in}{1.044240in}}%
\pgfpathlineto{\pgfqpoint{3.322601in}{1.046950in}}%
\pgfpathlineto{\pgfqpoint{3.332807in}{1.046950in}}%
\pgfpathlineto{\pgfqpoint{3.337910in}{1.049819in}}%
\pgfpathlineto{\pgfqpoint{3.348116in}{1.049819in}}%
\pgfpathlineto{\pgfqpoint{3.353220in}{1.052529in}}%
\pgfpathlineto{\pgfqpoint{3.363426in}{1.052529in}}%
\pgfpathlineto{\pgfqpoint{3.368529in}{1.055398in}}%
\pgfpathlineto{\pgfqpoint{3.378735in}{1.055398in}}%
\pgfpathlineto{\pgfqpoint{3.383838in}{1.058267in}}%
\pgfpathlineto{\pgfqpoint{3.394044in}{1.058267in}}%
\pgfpathlineto{\pgfqpoint{3.399147in}{1.060977in}}%
\pgfpathlineto{\pgfqpoint{3.404250in}{1.060977in}}%
\pgfpathlineto{\pgfqpoint{3.409353in}{1.063846in}}%
\pgfpathlineto{\pgfqpoint{3.419559in}{1.063846in}}%
\pgfpathlineto{\pgfqpoint{3.424662in}{1.066556in}}%
\pgfpathlineto{\pgfqpoint{3.434868in}{1.066556in}}%
\pgfpathlineto{\pgfqpoint{3.439971in}{1.069425in}}%
\pgfpathlineto{\pgfqpoint{3.445074in}{1.069425in}}%
\pgfpathlineto{\pgfqpoint{3.450177in}{1.072294in}}%
\pgfpathlineto{\pgfqpoint{3.460383in}{1.072294in}}%
\pgfpathlineto{\pgfqpoint{3.465486in}{1.075004in}}%
\pgfpathlineto{\pgfqpoint{3.470589in}{1.075004in}}%
\pgfpathlineto{\pgfqpoint{3.475692in}{1.077873in}}%
\pgfpathlineto{\pgfqpoint{3.480795in}{1.077873in}}%
\pgfpathlineto{\pgfqpoint{3.485898in}{1.080583in}}%
\pgfpathlineto{\pgfqpoint{3.496104in}{1.080583in}}%
\pgfpathlineto{\pgfqpoint{3.501207in}{1.083452in}}%
\pgfpathlineto{\pgfqpoint{3.506310in}{1.083452in}}%
\pgfpathlineto{\pgfqpoint{3.511413in}{1.086322in}}%
\pgfpathlineto{\pgfqpoint{3.516516in}{1.086322in}}%
\pgfpathlineto{\pgfqpoint{3.521619in}{1.089031in}}%
\pgfpathlineto{\pgfqpoint{3.531826in}{1.089031in}}%
\pgfpathlineto{\pgfqpoint{3.536929in}{1.091901in}}%
\pgfpathlineto{\pgfqpoint{3.542032in}{1.091901in}}%
\pgfpathlineto{\pgfqpoint{3.547135in}{1.094610in}}%
\pgfpathlineto{\pgfqpoint{3.552238in}{1.094610in}}%
\pgfpathlineto{\pgfqpoint{3.557341in}{1.097480in}}%
\pgfpathlineto{\pgfqpoint{3.562444in}{1.097480in}}%
\pgfpathlineto{\pgfqpoint{3.567547in}{1.100189in}}%
\pgfpathlineto{\pgfqpoint{3.577753in}{1.100189in}}%
\pgfpathlineto{\pgfqpoint{3.582856in}{1.103059in}}%
\pgfpathlineto{\pgfqpoint{3.587959in}{1.103059in}}%
\pgfpathlineto{\pgfqpoint{3.593062in}{1.105928in}}%
\pgfpathlineto{\pgfqpoint{3.603268in}{1.105928in}}%
\pgfpathlineto{\pgfqpoint{3.608371in}{1.108638in}}%
\pgfpathlineto{\pgfqpoint{3.613474in}{1.108638in}}%
\pgfpathlineto{\pgfqpoint{3.618577in}{1.111507in}}%
\pgfpathlineto{\pgfqpoint{3.623680in}{1.111507in}}%
\pgfpathlineto{\pgfqpoint{3.628783in}{1.114217in}}%
\pgfpathlineto{\pgfqpoint{3.638989in}{1.114217in}}%
\pgfpathlineto{\pgfqpoint{3.644092in}{1.117086in}}%
\pgfpathlineto{\pgfqpoint{3.649195in}{1.117086in}}%
\pgfpathlineto{\pgfqpoint{3.654298in}{1.119955in}}%
\pgfpathlineto{\pgfqpoint{3.664504in}{1.119955in}}%
\pgfpathlineto{\pgfqpoint{3.669607in}{1.122665in}}%
\pgfpathlineto{\pgfqpoint{3.674710in}{1.122665in}}%
\pgfpathlineto{\pgfqpoint{3.679813in}{1.125534in}}%
\pgfpathlineto{\pgfqpoint{3.690019in}{1.125534in}}%
\pgfpathlineto{\pgfqpoint{3.695122in}{1.128244in}}%
\pgfpathlineto{\pgfqpoint{3.705329in}{1.128244in}}%
\pgfpathlineto{\pgfqpoint{3.710432in}{1.131113in}}%
\pgfpathlineto{\pgfqpoint{3.715535in}{1.131113in}}%
\pgfpathlineto{\pgfqpoint{3.720638in}{1.133982in}}%
\pgfpathlineto{\pgfqpoint{3.730844in}{1.133982in}}%
\pgfpathlineto{\pgfqpoint{3.735947in}{1.136692in}}%
\pgfpathlineto{\pgfqpoint{3.741050in}{1.136692in}}%
\pgfpathlineto{\pgfqpoint{3.746153in}{1.139561in}}%
\pgfpathlineto{\pgfqpoint{3.751256in}{1.139561in}}%
\pgfpathlineto{\pgfqpoint{3.756359in}{1.142271in}}%
\pgfpathlineto{\pgfqpoint{3.766565in}{1.142271in}}%
\pgfpathlineto{\pgfqpoint{3.771668in}{1.145140in}}%
\pgfpathlineto{\pgfqpoint{3.776771in}{1.145140in}}%
\pgfpathlineto{\pgfqpoint{3.781874in}{1.148009in}}%
\pgfpathlineto{\pgfqpoint{3.792080in}{1.148009in}}%
\pgfpathlineto{\pgfqpoint{3.797183in}{1.150719in}}%
\pgfpathlineto{\pgfqpoint{3.802286in}{1.150719in}}%
\pgfpathlineto{\pgfqpoint{3.807389in}{1.153588in}}%
\pgfpathlineto{\pgfqpoint{3.812492in}{1.153588in}}%
\pgfpathlineto{\pgfqpoint{3.817595in}{1.156298in}}%
\pgfpathlineto{\pgfqpoint{3.827801in}{1.156298in}}%
\pgfpathlineto{\pgfqpoint{3.832904in}{1.159167in}}%
\pgfpathlineto{\pgfqpoint{3.838007in}{1.159167in}}%
\pgfpathlineto{\pgfqpoint{3.843110in}{1.162036in}}%
\pgfpathlineto{\pgfqpoint{3.853316in}{1.162036in}}%
\pgfpathlineto{\pgfqpoint{3.858419in}{1.164746in}}%
\pgfpathlineto{\pgfqpoint{3.863522in}{1.164746in}}%
\pgfpathlineto{\pgfqpoint{3.868625in}{1.167615in}}%
\pgfpathlineto{\pgfqpoint{3.873729in}{1.167615in}}%
\pgfpathlineto{\pgfqpoint{3.878832in}{1.170325in}}%
\pgfpathlineto{\pgfqpoint{3.889038in}{1.170325in}}%
\pgfpathlineto{\pgfqpoint{3.894141in}{1.173194in}}%
\pgfpathlineto{\pgfqpoint{3.899244in}{1.173194in}}%
\pgfpathlineto{\pgfqpoint{3.904347in}{1.176064in}}%
\pgfpathlineto{\pgfqpoint{3.909450in}{1.176064in}}%
\pgfpathlineto{\pgfqpoint{3.914553in}{1.178773in}}%
\pgfpathlineto{\pgfqpoint{3.924759in}{1.178773in}}%
\pgfpathlineto{\pgfqpoint{3.929862in}{1.181643in}}%
\pgfpathlineto{\pgfqpoint{3.934965in}{1.181643in}}%
\pgfpathlineto{\pgfqpoint{3.940068in}{1.184352in}}%
\pgfpathlineto{\pgfqpoint{3.945171in}{1.184352in}}%
\pgfpathlineto{\pgfqpoint{3.950274in}{1.187222in}}%
\pgfpathlineto{\pgfqpoint{3.955377in}{1.187222in}}%
\pgfpathlineto{\pgfqpoint{3.960480in}{1.190091in}}%
\pgfpathlineto{\pgfqpoint{3.970686in}{1.190091in}}%
\pgfpathlineto{\pgfqpoint{3.975789in}{1.192801in}}%
\pgfpathlineto{\pgfqpoint{3.980892in}{1.192801in}}%
\pgfpathlineto{\pgfqpoint{3.985995in}{1.195670in}}%
\pgfpathlineto{\pgfqpoint{3.996201in}{1.195670in}}%
\pgfpathlineto{\pgfqpoint{4.001304in}{1.198380in}}%
\pgfpathlineto{\pgfqpoint{4.006407in}{1.198380in}}%
\pgfpathlineto{\pgfqpoint{4.011510in}{1.201249in}}%
\pgfpathlineto{\pgfqpoint{4.021716in}{1.201249in}}%
\pgfpathlineto{\pgfqpoint{4.026819in}{1.204118in}}%
\pgfpathlineto{\pgfqpoint{4.031922in}{1.204118in}}%
\pgfpathlineto{\pgfqpoint{4.037025in}{1.206828in}}%
\pgfpathlineto{\pgfqpoint{4.042128in}{1.206828in}}%
\pgfpathlineto{\pgfqpoint{4.047232in}{1.209697in}}%
\pgfpathlineto{\pgfqpoint{4.057438in}{1.209697in}}%
\pgfpathlineto{\pgfqpoint{4.062541in}{1.212407in}}%
\pgfpathlineto{\pgfqpoint{4.067644in}{1.212407in}}%
\pgfpathlineto{\pgfqpoint{4.072747in}{1.215276in}}%
\pgfpathlineto{\pgfqpoint{4.077850in}{1.215276in}}%
\pgfpathlineto{\pgfqpoint{4.082953in}{1.217986in}}%
\pgfpathlineto{\pgfqpoint{4.093159in}{1.217986in}}%
\pgfpathlineto{\pgfqpoint{4.098262in}{1.220855in}}%
\pgfpathlineto{\pgfqpoint{4.103365in}{1.220855in}}%
\pgfpathlineto{\pgfqpoint{4.108468in}{1.223724in}}%
\pgfpathlineto{\pgfqpoint{4.118674in}{1.223724in}}%
\pgfpathlineto{\pgfqpoint{4.123777in}{1.226434in}}%
\pgfpathlineto{\pgfqpoint{4.128880in}{1.226434in}}%
\pgfpathlineto{\pgfqpoint{4.133983in}{1.229303in}}%
\pgfpathlineto{\pgfqpoint{4.139086in}{1.229303in}}%
\pgfpathlineto{\pgfqpoint{4.144189in}{1.232013in}}%
\pgfpathlineto{\pgfqpoint{4.154395in}{1.232013in}}%
\pgfpathlineto{\pgfqpoint{4.159498in}{1.234882in}}%
\pgfpathlineto{\pgfqpoint{4.164601in}{1.234882in}}%
\pgfpathlineto{\pgfqpoint{4.169704in}{1.237751in}}%
\pgfpathlineto{\pgfqpoint{4.174807in}{1.237751in}}%
\pgfpathlineto{\pgfqpoint{4.179910in}{1.240461in}}%
\pgfpathlineto{\pgfqpoint{4.190116in}{1.240461in}}%
\pgfpathlineto{\pgfqpoint{4.195219in}{1.243330in}}%
\pgfpathlineto{\pgfqpoint{4.200322in}{1.243330in}}%
\pgfpathlineto{\pgfqpoint{4.205425in}{1.246040in}}%
\pgfpathlineto{\pgfqpoint{4.210528in}{1.246040in}}%
\pgfpathlineto{\pgfqpoint{4.215631in}{1.248909in}}%
\pgfpathlineto{\pgfqpoint{4.220734in}{1.248909in}}%
\pgfpathlineto{\pgfqpoint{4.225838in}{1.251779in}}%
\pgfpathlineto{\pgfqpoint{4.236044in}{1.251779in}}%
\pgfpathlineto{\pgfqpoint{4.241147in}{1.254488in}}%
\pgfpathlineto{\pgfqpoint{4.246250in}{1.254488in}}%
\pgfpathlineto{\pgfqpoint{4.251353in}{1.257358in}}%
\pgfpathlineto{\pgfqpoint{4.261559in}{1.257358in}}%
\pgfpathlineto{\pgfqpoint{4.266662in}{1.260067in}}%
\pgfpathlineto{\pgfqpoint{4.271765in}{1.260067in}}%
\pgfpathlineto{\pgfqpoint{4.276868in}{1.262937in}}%
\pgfpathlineto{\pgfqpoint{4.281971in}{1.262937in}}%
\pgfpathlineto{\pgfqpoint{4.287074in}{1.265806in}}%
\pgfpathlineto{\pgfqpoint{4.297280in}{1.265806in}}%
\pgfpathlineto{\pgfqpoint{4.302383in}{1.268516in}}%
\pgfpathlineto{\pgfqpoint{4.307486in}{1.268516in}}%
\pgfpathlineto{\pgfqpoint{4.312589in}{1.271385in}}%
\pgfpathlineto{\pgfqpoint{4.317692in}{1.271385in}}%
\pgfpathlineto{\pgfqpoint{4.322795in}{1.274095in}}%
\pgfpathlineto{\pgfqpoint{4.333001in}{1.274095in}}%
\pgfpathlineto{\pgfqpoint{4.338104in}{1.276964in}}%
\pgfpathlineto{\pgfqpoint{4.343207in}{1.276964in}}%
\pgfpathlineto{\pgfqpoint{4.348310in}{1.279833in}}%
\pgfpathlineto{\pgfqpoint{4.353413in}{1.279833in}}%
\pgfpathlineto{\pgfqpoint{4.358516in}{1.282543in}}%
\pgfpathlineto{\pgfqpoint{4.368722in}{1.282543in}}%
\pgfpathlineto{\pgfqpoint{4.373825in}{1.285412in}}%
\pgfpathlineto{\pgfqpoint{4.378928in}{1.285412in}}%
\pgfpathlineto{\pgfqpoint{4.384031in}{1.288122in}}%
\pgfpathlineto{\pgfqpoint{4.389134in}{1.288122in}}%
\pgfpathlineto{\pgfqpoint{4.394237in}{1.290991in}}%
\pgfpathlineto{\pgfqpoint{4.399341in}{1.290991in}}%
\pgfpathlineto{\pgfqpoint{4.404444in}{1.293860in}}%
\pgfpathlineto{\pgfqpoint{4.414650in}{1.293860in}}%
\pgfpathlineto{\pgfqpoint{4.419753in}{1.296570in}}%
\pgfpathlineto{\pgfqpoint{4.424856in}{1.296570in}}%
\pgfpathlineto{\pgfqpoint{4.429959in}{1.299439in}}%
\pgfpathlineto{\pgfqpoint{4.440165in}{1.299439in}}%
\pgfpathlineto{\pgfqpoint{4.445268in}{1.302149in}}%
\pgfpathlineto{\pgfqpoint{4.450371in}{1.302149in}}%
\pgfpathlineto{\pgfqpoint{4.455474in}{1.305018in}}%
\pgfpathlineto{\pgfqpoint{4.465680in}{1.305018in}}%
\pgfpathlineto{\pgfqpoint{4.470783in}{1.307887in}}%
\pgfpathlineto{\pgfqpoint{4.475886in}{1.307887in}}%
\pgfpathlineto{\pgfqpoint{4.480989in}{1.310597in}}%
\pgfpathlineto{\pgfqpoint{4.486092in}{1.310597in}}%
\pgfpathlineto{\pgfqpoint{4.491195in}{1.313466in}}%
\pgfpathlineto{\pgfqpoint{4.501401in}{1.313466in}}%
\pgfpathlineto{\pgfqpoint{4.506504in}{1.316176in}}%
\pgfpathlineto{\pgfqpoint{4.511607in}{1.316176in}}%
\pgfpathlineto{\pgfqpoint{4.516710in}{1.319045in}}%
\pgfpathlineto{\pgfqpoint{4.521813in}{1.319045in}}%
\pgfpathlineto{\pgfqpoint{4.526916in}{1.321914in}}%
\pgfpathlineto{\pgfqpoint{4.537122in}{1.321914in}}%
\pgfpathlineto{\pgfqpoint{4.542225in}{1.324624in}}%
\pgfpathlineto{\pgfqpoint{4.547328in}{1.324624in}}%
\pgfpathlineto{\pgfqpoint{4.552431in}{1.327493in}}%
\pgfpathlineto{\pgfqpoint{4.562637in}{1.327493in}}%
\pgfpathlineto{\pgfqpoint{4.567740in}{1.330203in}}%
\pgfpathlineto{\pgfqpoint{4.572844in}{1.330203in}}%
\pgfpathlineto{\pgfqpoint{4.577947in}{1.333072in}}%
\pgfpathlineto{\pgfqpoint{4.583050in}{1.333072in}}%
\pgfpathlineto{\pgfqpoint{4.588153in}{1.335782in}}%
\pgfpathlineto{\pgfqpoint{4.598359in}{1.335782in}}%
\pgfpathlineto{\pgfqpoint{4.603462in}{1.338651in}}%
\pgfpathlineto{\pgfqpoint{4.608565in}{1.338651in}}%
\pgfpathlineto{\pgfqpoint{4.613668in}{1.341521in}}%
\pgfpathlineto{\pgfqpoint{4.618771in}{1.341521in}}%
\pgfpathlineto{\pgfqpoint{4.623874in}{1.344230in}}%
\pgfpathlineto{\pgfqpoint{4.634080in}{1.344230in}}%
\pgfpathlineto{\pgfqpoint{4.639183in}{1.347100in}}%
\pgfpathlineto{\pgfqpoint{4.644286in}{1.344230in}}%
\pgfpathlineto{\pgfqpoint{4.649389in}{1.349809in}}%
\pgfpathlineto{\pgfqpoint{4.659595in}{1.349809in}}%
\pgfpathlineto{\pgfqpoint{4.664698in}{1.352679in}}%
\pgfpathlineto{\pgfqpoint{4.669801in}{1.352679in}}%
\pgfpathlineto{\pgfqpoint{4.674904in}{1.355548in}}%
\pgfpathlineto{\pgfqpoint{4.680007in}{1.355548in}}%
\pgfpathlineto{\pgfqpoint{4.685110in}{1.358258in}}%
\pgfpathlineto{\pgfqpoint{4.695316in}{1.358258in}}%
\pgfpathlineto{\pgfqpoint{4.700419in}{1.361127in}}%
\pgfpathlineto{\pgfqpoint{4.705522in}{1.361127in}}%
\pgfpathlineto{\pgfqpoint{4.710625in}{1.363837in}}%
\pgfpathlineto{\pgfqpoint{4.715728in}{1.363837in}}%
\pgfpathlineto{\pgfqpoint{4.720831in}{1.366706in}}%
\pgfpathlineto{\pgfqpoint{4.731037in}{1.366706in}}%
\pgfpathlineto{\pgfqpoint{4.736140in}{1.369575in}}%
\pgfpathlineto{\pgfqpoint{4.741243in}{1.369575in}}%
\pgfpathlineto{\pgfqpoint{4.746347in}{1.372285in}}%
\pgfpathlineto{\pgfqpoint{4.756553in}{1.372285in}}%
\pgfpathlineto{\pgfqpoint{4.761656in}{1.375154in}}%
\pgfpathlineto{\pgfqpoint{4.766759in}{1.375154in}}%
\pgfpathlineto{\pgfqpoint{4.771862in}{1.377864in}}%
\pgfpathlineto{\pgfqpoint{4.782068in}{1.377864in}}%
\pgfpathlineto{\pgfqpoint{4.787171in}{1.380733in}}%
\pgfpathlineto{\pgfqpoint{4.792274in}{1.380733in}}%
\pgfpathlineto{\pgfqpoint{4.797377in}{1.383602in}}%
\pgfpathlineto{\pgfqpoint{4.802480in}{1.383602in}}%
\pgfpathlineto{\pgfqpoint{4.807583in}{1.386312in}}%
\pgfpathlineto{\pgfqpoint{4.817789in}{1.386312in}}%
\pgfpathlineto{\pgfqpoint{4.822892in}{1.389181in}}%
\pgfpathlineto{\pgfqpoint{4.827995in}{1.389181in}}%
\pgfpathlineto{\pgfqpoint{4.833098in}{1.391891in}}%
\pgfpathlineto{\pgfqpoint{4.843304in}{1.391891in}}%
\pgfpathlineto{\pgfqpoint{4.848407in}{1.394760in}}%
\pgfpathlineto{\pgfqpoint{4.853510in}{1.394760in}}%
\pgfpathlineto{\pgfqpoint{4.858613in}{1.397629in}}%
\pgfpathlineto{\pgfqpoint{4.863716in}{1.397629in}}%
\pgfpathlineto{\pgfqpoint{4.868819in}{1.400339in}}%
\pgfpathlineto{\pgfqpoint{4.873922in}{1.400339in}}%
\pgfpathlineto{\pgfqpoint{4.879025in}{1.403208in}}%
\pgfpathlineto{\pgfqpoint{4.889231in}{1.403208in}}%
\pgfpathlineto{\pgfqpoint{4.894334in}{1.405918in}}%
\pgfpathlineto{\pgfqpoint{4.899437in}{1.405918in}}%
\pgfpathlineto{\pgfqpoint{4.904540in}{1.408787in}}%
\pgfpathlineto{\pgfqpoint{4.909643in}{1.408787in}}%
\pgfpathlineto{\pgfqpoint{4.914746in}{1.411656in}}%
\pgfpathlineto{\pgfqpoint{4.919850in}{1.411656in}}%
\pgfpathlineto{\pgfqpoint{4.924953in}{1.414366in}}%
\pgfpathlineto{\pgfqpoint{4.935159in}{1.414366in}}%
\pgfpathlineto{\pgfqpoint{4.940262in}{1.417235in}}%
\pgfpathlineto{\pgfqpoint{4.945365in}{1.417235in}}%
\pgfpathlineto{\pgfqpoint{4.950468in}{1.419945in}}%
\pgfpathlineto{\pgfqpoint{4.955571in}{1.419945in}}%
\pgfpathlineto{\pgfqpoint{4.960674in}{1.422814in}}%
\pgfpathlineto{\pgfqpoint{4.965777in}{1.422814in}}%
\pgfpathlineto{\pgfqpoint{4.970880in}{1.425684in}}%
\pgfpathlineto{\pgfqpoint{4.981086in}{1.425684in}}%
\pgfpathlineto{\pgfqpoint{4.986189in}{1.428393in}}%
\pgfpathlineto{\pgfqpoint{4.991292in}{1.428393in}}%
\pgfpathlineto{\pgfqpoint{4.996395in}{1.431263in}}%
\pgfpathlineto{\pgfqpoint{5.001498in}{1.431263in}}%
\pgfpathlineto{\pgfqpoint{5.006601in}{1.433972in}}%
\pgfpathlineto{\pgfqpoint{5.016807in}{1.433972in}}%
\pgfpathlineto{\pgfqpoint{5.021910in}{1.436842in}}%
\pgfpathlineto{\pgfqpoint{5.027013in}{1.436842in}}%
\pgfpathlineto{\pgfqpoint{5.032116in}{1.439711in}}%
\pgfpathlineto{\pgfqpoint{5.042322in}{1.439711in}}%
\pgfpathlineto{\pgfqpoint{5.047425in}{1.442421in}}%
\pgfpathlineto{\pgfqpoint{5.052528in}{1.442421in}}%
\pgfpathlineto{\pgfqpoint{5.057631in}{1.445290in}}%
\pgfpathlineto{\pgfqpoint{5.062734in}{1.445290in}}%
\pgfpathlineto{\pgfqpoint{5.067837in}{1.448000in}}%
\pgfpathlineto{\pgfqpoint{5.078043in}{1.448000in}}%
\pgfpathlineto{\pgfqpoint{5.083146in}{1.450869in}}%
\pgfpathlineto{\pgfqpoint{5.088249in}{1.450869in}}%
\pgfpathlineto{\pgfqpoint{5.093352in}{1.453579in}}%
\pgfpathlineto{\pgfqpoint{5.103559in}{1.453579in}}%
\pgfpathlineto{\pgfqpoint{5.108662in}{1.456448in}}%
\pgfpathlineto{\pgfqpoint{5.118868in}{1.456448in}}%
\pgfpathlineto{\pgfqpoint{5.123971in}{1.459317in}}%
\pgfpathlineto{\pgfqpoint{5.129074in}{1.459317in}}%
\pgfpathlineto{\pgfqpoint{5.134177in}{1.462027in}}%
\pgfpathlineto{\pgfqpoint{5.144383in}{1.462027in}}%
\pgfpathlineto{\pgfqpoint{5.149486in}{1.464896in}}%
\pgfpathlineto{\pgfqpoint{5.159692in}{1.464896in}}%
\pgfpathlineto{\pgfqpoint{5.164795in}{1.467606in}}%
\pgfpathlineto{\pgfqpoint{5.175001in}{1.467606in}}%
\pgfpathlineto{\pgfqpoint{5.180104in}{1.470475in}}%
\pgfpathlineto{\pgfqpoint{5.190310in}{1.470475in}}%
\pgfpathlineto{\pgfqpoint{5.195413in}{1.473344in}}%
\pgfpathlineto{\pgfqpoint{5.205619in}{1.473344in}}%
\pgfpathlineto{\pgfqpoint{5.210722in}{1.476054in}}%
\pgfpathlineto{\pgfqpoint{5.215825in}{1.476054in}}%
\pgfpathlineto{\pgfqpoint{5.220928in}{1.478923in}}%
\pgfpathlineto{\pgfqpoint{5.231134in}{1.478923in}}%
\pgfpathlineto{\pgfqpoint{5.236237in}{1.481633in}}%
\pgfpathlineto{\pgfqpoint{5.246443in}{1.481633in}}%
\pgfpathlineto{\pgfqpoint{5.251546in}{1.484502in}}%
\pgfpathlineto{\pgfqpoint{5.256649in}{1.484502in}}%
\pgfpathlineto{\pgfqpoint{5.261752in}{1.487371in}}%
\pgfpathlineto{\pgfqpoint{5.271959in}{1.487371in}}%
\pgfpathlineto{\pgfqpoint{5.277062in}{1.490081in}}%
\pgfpathlineto{\pgfqpoint{5.287268in}{1.490081in}}%
\pgfpathlineto{\pgfqpoint{5.292371in}{1.492950in}}%
\pgfpathlineto{\pgfqpoint{5.297474in}{1.492950in}}%
\pgfpathlineto{\pgfqpoint{5.302577in}{1.495660in}}%
\pgfpathlineto{\pgfqpoint{5.312783in}{1.495660in}}%
\pgfpathlineto{\pgfqpoint{5.317886in}{1.498529in}}%
\pgfpathlineto{\pgfqpoint{5.328092in}{1.498529in}}%
\pgfpathlineto{\pgfqpoint{5.333195in}{1.501399in}}%
\pgfpathlineto{\pgfqpoint{5.338298in}{1.501399in}}%
\pgfpathlineto{\pgfqpoint{5.343401in}{1.504108in}}%
\pgfpathlineto{\pgfqpoint{5.353607in}{1.504108in}}%
\pgfpathlineto{\pgfqpoint{5.358710in}{1.506978in}}%
\pgfpathlineto{\pgfqpoint{5.363813in}{1.506978in}}%
\pgfpathlineto{\pgfqpoint{5.368916in}{1.509687in}}%
\pgfpathlineto{\pgfqpoint{5.379122in}{1.509687in}}%
\pgfpathlineto{\pgfqpoint{5.384225in}{1.512557in}}%
\pgfpathlineto{\pgfqpoint{5.394431in}{1.512557in}}%
\pgfpathlineto{\pgfqpoint{5.399534in}{1.515426in}}%
\pgfpathlineto{\pgfqpoint{5.404637in}{1.515426in}}%
\pgfpathlineto{\pgfqpoint{5.409740in}{1.518136in}}%
\pgfpathlineto{\pgfqpoint{5.419946in}{1.518136in}}%
\pgfpathlineto{\pgfqpoint{5.425049in}{1.521005in}}%
\pgfpathlineto{\pgfqpoint{5.430152in}{1.521005in}}%
\pgfpathlineto{\pgfqpoint{5.435255in}{1.523715in}}%
\pgfpathlineto{\pgfqpoint{5.445462in}{1.523715in}}%
\pgfpathlineto{\pgfqpoint{5.450565in}{1.526584in}}%
\pgfpathlineto{\pgfqpoint{5.455668in}{1.526584in}}%
\pgfpathlineto{\pgfqpoint{5.460771in}{1.529453in}}%
\pgfpathlineto{\pgfqpoint{5.470977in}{1.529453in}}%
\pgfpathlineto{\pgfqpoint{5.476080in}{1.532163in}}%
\pgfpathlineto{\pgfqpoint{5.481183in}{1.532163in}}%
\pgfpathlineto{\pgfqpoint{5.486286in}{1.535032in}}%
\pgfpathlineto{\pgfqpoint{5.491389in}{1.535032in}}%
\pgfpathlineto{\pgfqpoint{5.496492in}{1.537742in}}%
\pgfpathlineto{\pgfqpoint{5.506698in}{1.537742in}}%
\pgfpathlineto{\pgfqpoint{5.511801in}{1.540611in}}%
\pgfpathlineto{\pgfqpoint{5.516904in}{1.540611in}}%
\pgfpathlineto{\pgfqpoint{5.522007in}{1.543480in}}%
\pgfpathlineto{\pgfqpoint{5.527110in}{1.543480in}}%
\pgfpathlineto{\pgfqpoint{5.532213in}{1.546190in}}%
\pgfpathlineto{\pgfqpoint{5.537316in}{1.546190in}}%
\pgfpathlineto{\pgfqpoint{5.542419in}{1.549059in}}%
\pgfpathlineto{\pgfqpoint{5.547522in}{1.549059in}}%
\pgfpathlineto{\pgfqpoint{5.552625in}{1.551769in}}%
\pgfpathlineto{\pgfqpoint{5.562831in}{1.551769in}}%
\pgfpathlineto{\pgfqpoint{5.567934in}{1.554638in}}%
\pgfpathlineto{\pgfqpoint{5.578140in}{1.554638in}}%
\pgfpathlineto{\pgfqpoint{5.583243in}{1.557507in}}%
\pgfpathlineto{\pgfqpoint{5.588346in}{1.557507in}}%
\pgfpathlineto{\pgfqpoint{5.588346in}{1.557507in}}%
\pgfusepath{stroke}%
\end{pgfscope}%
\begin{pgfscope}%
\pgfsetrectcap%
\pgfsetmiterjoin%
\pgfsetlinewidth{0.803000pt}%
\definecolor{currentstroke}{rgb}{0.000000,0.000000,0.000000}%
\pgfsetstrokecolor{currentstroke}%
\pgfsetdash{}{0pt}%
\pgfpathmoveto{\pgfqpoint{0.834028in}{0.582778in}}%
\pgfpathlineto{\pgfqpoint{0.834028in}{3.020000in}}%
\pgfusepath{stroke}%
\end{pgfscope}%
\begin{pgfscope}%
\pgfsetrectcap%
\pgfsetmiterjoin%
\pgfsetlinewidth{0.803000pt}%
\definecolor{currentstroke}{rgb}{0.000000,0.000000,0.000000}%
\pgfsetstrokecolor{currentstroke}%
\pgfsetdash{}{0pt}%
\pgfpathmoveto{\pgfqpoint{5.593472in}{0.582778in}}%
\pgfpathlineto{\pgfqpoint{5.593472in}{3.020000in}}%
\pgfusepath{stroke}%
\end{pgfscope}%
\begin{pgfscope}%
\pgfsetrectcap%
\pgfsetmiterjoin%
\pgfsetlinewidth{0.803000pt}%
\definecolor{currentstroke}{rgb}{0.000000,0.000000,0.000000}%
\pgfsetstrokecolor{currentstroke}%
\pgfsetdash{}{0pt}%
\pgfpathmoveto{\pgfqpoint{0.834028in}{0.582778in}}%
\pgfpathlineto{\pgfqpoint{5.593472in}{0.582778in}}%
\pgfusepath{stroke}%
\end{pgfscope}%
\begin{pgfscope}%
\pgfsetrectcap%
\pgfsetmiterjoin%
\pgfsetlinewidth{0.803000pt}%
\definecolor{currentstroke}{rgb}{0.000000,0.000000,0.000000}%
\pgfsetstrokecolor{currentstroke}%
\pgfsetdash{}{0pt}%
\pgfpathmoveto{\pgfqpoint{0.834028in}{3.020000in}}%
\pgfpathlineto{\pgfqpoint{5.593472in}{3.020000in}}%
\pgfusepath{stroke}%
\end{pgfscope}%
\begin{pgfscope}%
\definecolor{textcolor}{rgb}{0.000000,0.000000,0.000000}%
\pgfsetstrokecolor{textcolor}%
\pgfsetfillcolor{textcolor}%
\pgftext[x=3.150000in,y=3.430000in,,top]{\color{textcolor}\sffamily\fontsize{12.000000}{14.400000}\selectfont Measured Impedance and Phase of Inductor}%
\end{pgfscope}%
\end{pgfpicture}%
\makeatother%
\endgroup%

		\caption{Measured impedance magnitude and phase of an inductor as a function of frequency.}
		\label{fig:352}
	\end{figure}
	
	\subsubsection{}
%	Figure \ref{fig:353} shows the measured magnitude and phase of the return loss of the quartz
%	resonator versus frequency. A pronounced minimum of the return loss occurs at the
%	resonance frequency $f_0 \approx 9.8283\,\mathrm{MHz}$, accompanied by a rapid phase
%	transition of approximately $180^\circ$, which is characteristic of quartz resonance.
	
	The $3\,\mathrm{dB}$ bandwidth is determined from the frequencies where the return loss
	is $3\,\mathrm{dB}$ above the minimum. From the figure, a bandwidth of approximately
	$\Delta f_{3\mathrm{dB}} \approx 540\,\mathrm{Hz}$ is obtained.
	The quality factor of the quartz resonator is therefore
	\[
	Q = \frac{f_0}{\Delta f_{3\mathrm{dB}}} \approx \frac{9.8283\,\mathrm{MHz}}{540\,\mathrm{Hz}}
	\approx 18200.45
	\]
	indicating a very high Q and highly frequency-selective resonator.
	
	\begin{figure}[H]
		\centering
		%% Creator: Matplotlib, PGF backend
%%
%% To include the figure in your LaTeX document, write
%%   \input{<filename>.pgf}
%%
%% Make sure the required packages are loaded in your preamble
%%   \usepackage{pgf}
%%
%% Also ensure that all the required font packages are loaded; for instance,
%% the lmodern package is sometimes necessary when using math font.
%%   \usepackage{lmodern}
%%
%% Figures using additional raster images can only be included by \input if
%% they are in the same directory as the main LaTeX file. For loading figures
%% from other directories you can use the `import` package
%%   \usepackage{import}
%%
%% and then include the figures with
%%   \import{<path to file>}{<filename>.pgf}
%%
%% Matplotlib used the following preamble
%%
\begingroup%
\makeatletter%
\begin{pgfpicture}%
\pgfpathrectangle{\pgfpointorigin}{\pgfqpoint{6.300000in}{3.500000in}}%
\pgfusepath{use as bounding box, clip}%
\begin{pgfscope}%
\pgfsetbuttcap%
\pgfsetmiterjoin%
\definecolor{currentfill}{rgb}{1.000000,1.000000,1.000000}%
\pgfsetfillcolor{currentfill}%
\pgfsetlinewidth{0.000000pt}%
\definecolor{currentstroke}{rgb}{1.000000,1.000000,1.000000}%
\pgfsetstrokecolor{currentstroke}%
\pgfsetdash{}{0pt}%
\pgfpathmoveto{\pgfqpoint{0.000000in}{0.000000in}}%
\pgfpathlineto{\pgfqpoint{6.300000in}{0.000000in}}%
\pgfpathlineto{\pgfqpoint{6.300000in}{3.500000in}}%
\pgfpathlineto{\pgfqpoint{0.000000in}{3.500000in}}%
\pgfpathlineto{\pgfqpoint{0.000000in}{0.000000in}}%
\pgfpathclose%
\pgfusepath{fill}%
\end{pgfscope}%
\begin{pgfscope}%
\pgfsetbuttcap%
\pgfsetmiterjoin%
\definecolor{currentfill}{rgb}{1.000000,1.000000,1.000000}%
\pgfsetfillcolor{currentfill}%
\pgfsetlinewidth{0.000000pt}%
\definecolor{currentstroke}{rgb}{0.000000,0.000000,0.000000}%
\pgfsetstrokecolor{currentstroke}%
\pgfsetstrokeopacity{0.000000}%
\pgfsetdash{}{0pt}%
\pgfpathmoveto{\pgfqpoint{0.727161in}{0.565123in}}%
\pgfpathlineto{\pgfqpoint{5.649999in}{0.565123in}}%
\pgfpathlineto{\pgfqpoint{5.649999in}{3.350000in}}%
\pgfpathlineto{\pgfqpoint{0.727161in}{3.350000in}}%
\pgfpathlineto{\pgfqpoint{0.727161in}{0.565123in}}%
\pgfpathclose%
\pgfusepath{fill}%
\end{pgfscope}%
\begin{pgfscope}%
\pgfpathrectangle{\pgfqpoint{0.727161in}{0.565123in}}{\pgfqpoint{4.922838in}{2.784877in}}%
\pgfusepath{clip}%
\pgfsetbuttcap%
\pgfsetroundjoin%
\pgfsetlinewidth{0.803000pt}%
\definecolor{currentstroke}{rgb}{0.690196,0.690196,0.690196}%
\pgfsetstrokecolor{currentstroke}%
\pgfsetdash{{0.800000pt}{1.320000pt}}{0.000000pt}%
\pgfpathmoveto{\pgfqpoint{0.727161in}{0.565123in}}%
\pgfpathlineto{\pgfqpoint{0.727161in}{3.350000in}}%
\pgfusepath{stroke}%
\end{pgfscope}%
\begin{pgfscope}%
\pgfsetbuttcap%
\pgfsetroundjoin%
\definecolor{currentfill}{rgb}{0.000000,0.000000,0.000000}%
\pgfsetfillcolor{currentfill}%
\pgfsetlinewidth{0.803000pt}%
\definecolor{currentstroke}{rgb}{0.000000,0.000000,0.000000}%
\pgfsetstrokecolor{currentstroke}%
\pgfsetdash{}{0pt}%
\pgfsys@defobject{currentmarker}{\pgfqpoint{0.000000in}{-0.048611in}}{\pgfqpoint{0.000000in}{0.000000in}}{%
\pgfpathmoveto{\pgfqpoint{0.000000in}{0.000000in}}%
\pgfpathlineto{\pgfqpoint{0.000000in}{-0.048611in}}%
\pgfusepath{stroke,fill}%
}%
\begin{pgfscope}%
\pgfsys@transformshift{0.727161in}{0.565123in}%
\pgfsys@useobject{currentmarker}{}%
\end{pgfscope}%
\end{pgfscope}%
\begin{pgfscope}%
\definecolor{textcolor}{rgb}{0.000000,0.000000,0.000000}%
\pgfsetstrokecolor{textcolor}%
\pgfsetfillcolor{textcolor}%
\pgftext[x=0.727161in,y=0.467901in,,top]{\color{textcolor}\rmfamily\fontsize{10.000000}{12.000000}\selectfont \(\displaystyle {9.8270}\)}%
\end{pgfscope}%
\begin{pgfscope}%
\pgfpathrectangle{\pgfqpoint{0.727161in}{0.565123in}}{\pgfqpoint{4.922838in}{2.784877in}}%
\pgfusepath{clip}%
\pgfsetbuttcap%
\pgfsetroundjoin%
\pgfsetlinewidth{0.803000pt}%
\definecolor{currentstroke}{rgb}{0.690196,0.690196,0.690196}%
\pgfsetstrokecolor{currentstroke}%
\pgfsetdash{{0.800000pt}{1.320000pt}}{0.000000pt}%
\pgfpathmoveto{\pgfqpoint{1.612246in}{0.565123in}}%
\pgfpathlineto{\pgfqpoint{1.612246in}{3.350000in}}%
\pgfusepath{stroke}%
\end{pgfscope}%
\begin{pgfscope}%
\pgfsetbuttcap%
\pgfsetroundjoin%
\definecolor{currentfill}{rgb}{0.000000,0.000000,0.000000}%
\pgfsetfillcolor{currentfill}%
\pgfsetlinewidth{0.803000pt}%
\definecolor{currentstroke}{rgb}{0.000000,0.000000,0.000000}%
\pgfsetstrokecolor{currentstroke}%
\pgfsetdash{}{0pt}%
\pgfsys@defobject{currentmarker}{\pgfqpoint{0.000000in}{-0.048611in}}{\pgfqpoint{0.000000in}{0.000000in}}{%
\pgfpathmoveto{\pgfqpoint{0.000000in}{0.000000in}}%
\pgfpathlineto{\pgfqpoint{0.000000in}{-0.048611in}}%
\pgfusepath{stroke,fill}%
}%
\begin{pgfscope}%
\pgfsys@transformshift{1.612246in}{0.565123in}%
\pgfsys@useobject{currentmarker}{}%
\end{pgfscope}%
\end{pgfscope}%
\begin{pgfscope}%
\definecolor{textcolor}{rgb}{0.000000,0.000000,0.000000}%
\pgfsetstrokecolor{textcolor}%
\pgfsetfillcolor{textcolor}%
\pgftext[x=1.612246in,y=0.467901in,,top]{\color{textcolor}\rmfamily\fontsize{10.000000}{12.000000}\selectfont \(\displaystyle {9.8275}\)}%
\end{pgfscope}%
\begin{pgfscope}%
\pgfpathrectangle{\pgfqpoint{0.727161in}{0.565123in}}{\pgfqpoint{4.922838in}{2.784877in}}%
\pgfusepath{clip}%
\pgfsetbuttcap%
\pgfsetroundjoin%
\pgfsetlinewidth{0.803000pt}%
\definecolor{currentstroke}{rgb}{0.690196,0.690196,0.690196}%
\pgfsetstrokecolor{currentstroke}%
\pgfsetdash{{0.800000pt}{1.320000pt}}{0.000000pt}%
\pgfpathmoveto{\pgfqpoint{2.497330in}{0.565123in}}%
\pgfpathlineto{\pgfqpoint{2.497330in}{3.350000in}}%
\pgfusepath{stroke}%
\end{pgfscope}%
\begin{pgfscope}%
\pgfsetbuttcap%
\pgfsetroundjoin%
\definecolor{currentfill}{rgb}{0.000000,0.000000,0.000000}%
\pgfsetfillcolor{currentfill}%
\pgfsetlinewidth{0.803000pt}%
\definecolor{currentstroke}{rgb}{0.000000,0.000000,0.000000}%
\pgfsetstrokecolor{currentstroke}%
\pgfsetdash{}{0pt}%
\pgfsys@defobject{currentmarker}{\pgfqpoint{0.000000in}{-0.048611in}}{\pgfqpoint{0.000000in}{0.000000in}}{%
\pgfpathmoveto{\pgfqpoint{0.000000in}{0.000000in}}%
\pgfpathlineto{\pgfqpoint{0.000000in}{-0.048611in}}%
\pgfusepath{stroke,fill}%
}%
\begin{pgfscope}%
\pgfsys@transformshift{2.497330in}{0.565123in}%
\pgfsys@useobject{currentmarker}{}%
\end{pgfscope}%
\end{pgfscope}%
\begin{pgfscope}%
\definecolor{textcolor}{rgb}{0.000000,0.000000,0.000000}%
\pgfsetstrokecolor{textcolor}%
\pgfsetfillcolor{textcolor}%
\pgftext[x=2.497330in,y=0.467901in,,top]{\color{textcolor}\rmfamily\fontsize{10.000000}{12.000000}\selectfont \(\displaystyle {9.8280}\)}%
\end{pgfscope}%
\begin{pgfscope}%
\pgfpathrectangle{\pgfqpoint{0.727161in}{0.565123in}}{\pgfqpoint{4.922838in}{2.784877in}}%
\pgfusepath{clip}%
\pgfsetbuttcap%
\pgfsetroundjoin%
\pgfsetlinewidth{0.803000pt}%
\definecolor{currentstroke}{rgb}{0.690196,0.690196,0.690196}%
\pgfsetstrokecolor{currentstroke}%
\pgfsetdash{{0.800000pt}{1.320000pt}}{0.000000pt}%
\pgfpathmoveto{\pgfqpoint{3.382414in}{0.565123in}}%
\pgfpathlineto{\pgfqpoint{3.382414in}{3.350000in}}%
\pgfusepath{stroke}%
\end{pgfscope}%
\begin{pgfscope}%
\pgfsetbuttcap%
\pgfsetroundjoin%
\definecolor{currentfill}{rgb}{0.000000,0.000000,0.000000}%
\pgfsetfillcolor{currentfill}%
\pgfsetlinewidth{0.803000pt}%
\definecolor{currentstroke}{rgb}{0.000000,0.000000,0.000000}%
\pgfsetstrokecolor{currentstroke}%
\pgfsetdash{}{0pt}%
\pgfsys@defobject{currentmarker}{\pgfqpoint{0.000000in}{-0.048611in}}{\pgfqpoint{0.000000in}{0.000000in}}{%
\pgfpathmoveto{\pgfqpoint{0.000000in}{0.000000in}}%
\pgfpathlineto{\pgfqpoint{0.000000in}{-0.048611in}}%
\pgfusepath{stroke,fill}%
}%
\begin{pgfscope}%
\pgfsys@transformshift{3.382414in}{0.565123in}%
\pgfsys@useobject{currentmarker}{}%
\end{pgfscope}%
\end{pgfscope}%
\begin{pgfscope}%
\definecolor{textcolor}{rgb}{0.000000,0.000000,0.000000}%
\pgfsetstrokecolor{textcolor}%
\pgfsetfillcolor{textcolor}%
\pgftext[x=3.382414in,y=0.467901in,,top]{\color{textcolor}\rmfamily\fontsize{10.000000}{12.000000}\selectfont \(\displaystyle {9.8285}\)}%
\end{pgfscope}%
\begin{pgfscope}%
\pgfpathrectangle{\pgfqpoint{0.727161in}{0.565123in}}{\pgfqpoint{4.922838in}{2.784877in}}%
\pgfusepath{clip}%
\pgfsetbuttcap%
\pgfsetroundjoin%
\pgfsetlinewidth{0.803000pt}%
\definecolor{currentstroke}{rgb}{0.690196,0.690196,0.690196}%
\pgfsetstrokecolor{currentstroke}%
\pgfsetdash{{0.800000pt}{1.320000pt}}{0.000000pt}%
\pgfpathmoveto{\pgfqpoint{4.267498in}{0.565123in}}%
\pgfpathlineto{\pgfqpoint{4.267498in}{3.350000in}}%
\pgfusepath{stroke}%
\end{pgfscope}%
\begin{pgfscope}%
\pgfsetbuttcap%
\pgfsetroundjoin%
\definecolor{currentfill}{rgb}{0.000000,0.000000,0.000000}%
\pgfsetfillcolor{currentfill}%
\pgfsetlinewidth{0.803000pt}%
\definecolor{currentstroke}{rgb}{0.000000,0.000000,0.000000}%
\pgfsetstrokecolor{currentstroke}%
\pgfsetdash{}{0pt}%
\pgfsys@defobject{currentmarker}{\pgfqpoint{0.000000in}{-0.048611in}}{\pgfqpoint{0.000000in}{0.000000in}}{%
\pgfpathmoveto{\pgfqpoint{0.000000in}{0.000000in}}%
\pgfpathlineto{\pgfqpoint{0.000000in}{-0.048611in}}%
\pgfusepath{stroke,fill}%
}%
\begin{pgfscope}%
\pgfsys@transformshift{4.267498in}{0.565123in}%
\pgfsys@useobject{currentmarker}{}%
\end{pgfscope}%
\end{pgfscope}%
\begin{pgfscope}%
\definecolor{textcolor}{rgb}{0.000000,0.000000,0.000000}%
\pgfsetstrokecolor{textcolor}%
\pgfsetfillcolor{textcolor}%
\pgftext[x=4.267498in,y=0.467901in,,top]{\color{textcolor}\rmfamily\fontsize{10.000000}{12.000000}\selectfont \(\displaystyle {9.8290}\)}%
\end{pgfscope}%
\begin{pgfscope}%
\pgfpathrectangle{\pgfqpoint{0.727161in}{0.565123in}}{\pgfqpoint{4.922838in}{2.784877in}}%
\pgfusepath{clip}%
\pgfsetbuttcap%
\pgfsetroundjoin%
\pgfsetlinewidth{0.803000pt}%
\definecolor{currentstroke}{rgb}{0.690196,0.690196,0.690196}%
\pgfsetstrokecolor{currentstroke}%
\pgfsetdash{{0.800000pt}{1.320000pt}}{0.000000pt}%
\pgfpathmoveto{\pgfqpoint{5.152582in}{0.565123in}}%
\pgfpathlineto{\pgfqpoint{5.152582in}{3.350000in}}%
\pgfusepath{stroke}%
\end{pgfscope}%
\begin{pgfscope}%
\pgfsetbuttcap%
\pgfsetroundjoin%
\definecolor{currentfill}{rgb}{0.000000,0.000000,0.000000}%
\pgfsetfillcolor{currentfill}%
\pgfsetlinewidth{0.803000pt}%
\definecolor{currentstroke}{rgb}{0.000000,0.000000,0.000000}%
\pgfsetstrokecolor{currentstroke}%
\pgfsetdash{}{0pt}%
\pgfsys@defobject{currentmarker}{\pgfqpoint{0.000000in}{-0.048611in}}{\pgfqpoint{0.000000in}{0.000000in}}{%
\pgfpathmoveto{\pgfqpoint{0.000000in}{0.000000in}}%
\pgfpathlineto{\pgfqpoint{0.000000in}{-0.048611in}}%
\pgfusepath{stroke,fill}%
}%
\begin{pgfscope}%
\pgfsys@transformshift{5.152582in}{0.565123in}%
\pgfsys@useobject{currentmarker}{}%
\end{pgfscope}%
\end{pgfscope}%
\begin{pgfscope}%
\definecolor{textcolor}{rgb}{0.000000,0.000000,0.000000}%
\pgfsetstrokecolor{textcolor}%
\pgfsetfillcolor{textcolor}%
\pgftext[x=5.152582in,y=0.467901in,,top]{\color{textcolor}\rmfamily\fontsize{10.000000}{12.000000}\selectfont \(\displaystyle {9.8295}\)}%
\end{pgfscope}%
\begin{pgfscope}%
\definecolor{textcolor}{rgb}{0.000000,0.000000,0.000000}%
\pgfsetstrokecolor{textcolor}%
\pgfsetfillcolor{textcolor}%
\pgftext[x=3.188580in,y=0.288889in,,top]{\color{textcolor}\rmfamily\fontsize{10.000000}{12.000000}\selectfont Frequency (Hz)}%
\end{pgfscope}%
\begin{pgfscope}%
\definecolor{textcolor}{rgb}{0.000000,0.000000,0.000000}%
\pgfsetstrokecolor{textcolor}%
\pgfsetfillcolor{textcolor}%
\pgftext[x=5.649999in,y=0.302778in,right,top]{\color{textcolor}\rmfamily\fontsize{10.000000}{12.000000}\selectfont \(\displaystyle \times{10^{6}}{}\)}%
\end{pgfscope}%
\begin{pgfscope}%
\pgfpathrectangle{\pgfqpoint{0.727161in}{0.565123in}}{\pgfqpoint{4.922838in}{2.784877in}}%
\pgfusepath{clip}%
\pgfsetbuttcap%
\pgfsetroundjoin%
\pgfsetlinewidth{0.803000pt}%
\definecolor{currentstroke}{rgb}{0.690196,0.690196,0.690196}%
\pgfsetstrokecolor{currentstroke}%
\pgfsetdash{{0.800000pt}{1.320000pt}}{0.000000pt}%
\pgfpathmoveto{\pgfqpoint{0.727161in}{0.567413in}}%
\pgfpathlineto{\pgfqpoint{5.649999in}{0.567413in}}%
\pgfusepath{stroke}%
\end{pgfscope}%
\begin{pgfscope}%
\pgfsetbuttcap%
\pgfsetroundjoin%
\definecolor{currentfill}{rgb}{0.000000,0.000000,0.000000}%
\pgfsetfillcolor{currentfill}%
\pgfsetlinewidth{0.803000pt}%
\definecolor{currentstroke}{rgb}{0.000000,0.000000,0.000000}%
\pgfsetstrokecolor{currentstroke}%
\pgfsetdash{}{0pt}%
\pgfsys@defobject{currentmarker}{\pgfqpoint{-0.048611in}{0.000000in}}{\pgfqpoint{-0.000000in}{0.000000in}}{%
\pgfpathmoveto{\pgfqpoint{-0.000000in}{0.000000in}}%
\pgfpathlineto{\pgfqpoint{-0.048611in}{0.000000in}}%
\pgfusepath{stroke,fill}%
}%
\begin{pgfscope}%
\pgfsys@transformshift{0.727161in}{0.567413in}%
\pgfsys@useobject{currentmarker}{}%
\end{pgfscope}%
\end{pgfscope}%
\begin{pgfscope}%
\definecolor{textcolor}{rgb}{0.000000,0.000000,0.000000}%
\pgfsetstrokecolor{textcolor}%
\pgfsetfillcolor{textcolor}%
\pgftext[x=0.344444in, y=0.519188in, left, base]{\color{textcolor}\rmfamily\fontsize{10.000000}{12.000000}\selectfont \(\displaystyle {\ensuremath{-}4.0}\)}%
\end{pgfscope}%
\begin{pgfscope}%
\pgfpathrectangle{\pgfqpoint{0.727161in}{0.565123in}}{\pgfqpoint{4.922838in}{2.784877in}}%
\pgfusepath{clip}%
\pgfsetbuttcap%
\pgfsetroundjoin%
\pgfsetlinewidth{0.803000pt}%
\definecolor{currentstroke}{rgb}{0.690196,0.690196,0.690196}%
\pgfsetstrokecolor{currentstroke}%
\pgfsetdash{{0.800000pt}{1.320000pt}}{0.000000pt}%
\pgfpathmoveto{\pgfqpoint{0.727161in}{0.894507in}}%
\pgfpathlineto{\pgfqpoint{5.649999in}{0.894507in}}%
\pgfusepath{stroke}%
\end{pgfscope}%
\begin{pgfscope}%
\pgfsetbuttcap%
\pgfsetroundjoin%
\definecolor{currentfill}{rgb}{0.000000,0.000000,0.000000}%
\pgfsetfillcolor{currentfill}%
\pgfsetlinewidth{0.803000pt}%
\definecolor{currentstroke}{rgb}{0.000000,0.000000,0.000000}%
\pgfsetstrokecolor{currentstroke}%
\pgfsetdash{}{0pt}%
\pgfsys@defobject{currentmarker}{\pgfqpoint{-0.048611in}{0.000000in}}{\pgfqpoint{-0.000000in}{0.000000in}}{%
\pgfpathmoveto{\pgfqpoint{-0.000000in}{0.000000in}}%
\pgfpathlineto{\pgfqpoint{-0.048611in}{0.000000in}}%
\pgfusepath{stroke,fill}%
}%
\begin{pgfscope}%
\pgfsys@transformshift{0.727161in}{0.894507in}%
\pgfsys@useobject{currentmarker}{}%
\end{pgfscope}%
\end{pgfscope}%
\begin{pgfscope}%
\definecolor{textcolor}{rgb}{0.000000,0.000000,0.000000}%
\pgfsetstrokecolor{textcolor}%
\pgfsetfillcolor{textcolor}%
\pgftext[x=0.344444in, y=0.846282in, left, base]{\color{textcolor}\rmfamily\fontsize{10.000000}{12.000000}\selectfont \(\displaystyle {\ensuremath{-}3.5}\)}%
\end{pgfscope}%
\begin{pgfscope}%
\pgfpathrectangle{\pgfqpoint{0.727161in}{0.565123in}}{\pgfqpoint{4.922838in}{2.784877in}}%
\pgfusepath{clip}%
\pgfsetbuttcap%
\pgfsetroundjoin%
\pgfsetlinewidth{0.803000pt}%
\definecolor{currentstroke}{rgb}{0.690196,0.690196,0.690196}%
\pgfsetstrokecolor{currentstroke}%
\pgfsetdash{{0.800000pt}{1.320000pt}}{0.000000pt}%
\pgfpathmoveto{\pgfqpoint{0.727161in}{1.221601in}}%
\pgfpathlineto{\pgfqpoint{5.649999in}{1.221601in}}%
\pgfusepath{stroke}%
\end{pgfscope}%
\begin{pgfscope}%
\pgfsetbuttcap%
\pgfsetroundjoin%
\definecolor{currentfill}{rgb}{0.000000,0.000000,0.000000}%
\pgfsetfillcolor{currentfill}%
\pgfsetlinewidth{0.803000pt}%
\definecolor{currentstroke}{rgb}{0.000000,0.000000,0.000000}%
\pgfsetstrokecolor{currentstroke}%
\pgfsetdash{}{0pt}%
\pgfsys@defobject{currentmarker}{\pgfqpoint{-0.048611in}{0.000000in}}{\pgfqpoint{-0.000000in}{0.000000in}}{%
\pgfpathmoveto{\pgfqpoint{-0.000000in}{0.000000in}}%
\pgfpathlineto{\pgfqpoint{-0.048611in}{0.000000in}}%
\pgfusepath{stroke,fill}%
}%
\begin{pgfscope}%
\pgfsys@transformshift{0.727161in}{1.221601in}%
\pgfsys@useobject{currentmarker}{}%
\end{pgfscope}%
\end{pgfscope}%
\begin{pgfscope}%
\definecolor{textcolor}{rgb}{0.000000,0.000000,0.000000}%
\pgfsetstrokecolor{textcolor}%
\pgfsetfillcolor{textcolor}%
\pgftext[x=0.344444in, y=1.173375in, left, base]{\color{textcolor}\rmfamily\fontsize{10.000000}{12.000000}\selectfont \(\displaystyle {\ensuremath{-}3.0}\)}%
\end{pgfscope}%
\begin{pgfscope}%
\pgfpathrectangle{\pgfqpoint{0.727161in}{0.565123in}}{\pgfqpoint{4.922838in}{2.784877in}}%
\pgfusepath{clip}%
\pgfsetbuttcap%
\pgfsetroundjoin%
\pgfsetlinewidth{0.803000pt}%
\definecolor{currentstroke}{rgb}{0.690196,0.690196,0.690196}%
\pgfsetstrokecolor{currentstroke}%
\pgfsetdash{{0.800000pt}{1.320000pt}}{0.000000pt}%
\pgfpathmoveto{\pgfqpoint{0.727161in}{1.548694in}}%
\pgfpathlineto{\pgfqpoint{5.649999in}{1.548694in}}%
\pgfusepath{stroke}%
\end{pgfscope}%
\begin{pgfscope}%
\pgfsetbuttcap%
\pgfsetroundjoin%
\definecolor{currentfill}{rgb}{0.000000,0.000000,0.000000}%
\pgfsetfillcolor{currentfill}%
\pgfsetlinewidth{0.803000pt}%
\definecolor{currentstroke}{rgb}{0.000000,0.000000,0.000000}%
\pgfsetstrokecolor{currentstroke}%
\pgfsetdash{}{0pt}%
\pgfsys@defobject{currentmarker}{\pgfqpoint{-0.048611in}{0.000000in}}{\pgfqpoint{-0.000000in}{0.000000in}}{%
\pgfpathmoveto{\pgfqpoint{-0.000000in}{0.000000in}}%
\pgfpathlineto{\pgfqpoint{-0.048611in}{0.000000in}}%
\pgfusepath{stroke,fill}%
}%
\begin{pgfscope}%
\pgfsys@transformshift{0.727161in}{1.548694in}%
\pgfsys@useobject{currentmarker}{}%
\end{pgfscope}%
\end{pgfscope}%
\begin{pgfscope}%
\definecolor{textcolor}{rgb}{0.000000,0.000000,0.000000}%
\pgfsetstrokecolor{textcolor}%
\pgfsetfillcolor{textcolor}%
\pgftext[x=0.344444in, y=1.500469in, left, base]{\color{textcolor}\rmfamily\fontsize{10.000000}{12.000000}\selectfont \(\displaystyle {\ensuremath{-}2.5}\)}%
\end{pgfscope}%
\begin{pgfscope}%
\pgfpathrectangle{\pgfqpoint{0.727161in}{0.565123in}}{\pgfqpoint{4.922838in}{2.784877in}}%
\pgfusepath{clip}%
\pgfsetbuttcap%
\pgfsetroundjoin%
\pgfsetlinewidth{0.803000pt}%
\definecolor{currentstroke}{rgb}{0.690196,0.690196,0.690196}%
\pgfsetstrokecolor{currentstroke}%
\pgfsetdash{{0.800000pt}{1.320000pt}}{0.000000pt}%
\pgfpathmoveto{\pgfqpoint{0.727161in}{1.875788in}}%
\pgfpathlineto{\pgfqpoint{5.649999in}{1.875788in}}%
\pgfusepath{stroke}%
\end{pgfscope}%
\begin{pgfscope}%
\pgfsetbuttcap%
\pgfsetroundjoin%
\definecolor{currentfill}{rgb}{0.000000,0.000000,0.000000}%
\pgfsetfillcolor{currentfill}%
\pgfsetlinewidth{0.803000pt}%
\definecolor{currentstroke}{rgb}{0.000000,0.000000,0.000000}%
\pgfsetstrokecolor{currentstroke}%
\pgfsetdash{}{0pt}%
\pgfsys@defobject{currentmarker}{\pgfqpoint{-0.048611in}{0.000000in}}{\pgfqpoint{-0.000000in}{0.000000in}}{%
\pgfpathmoveto{\pgfqpoint{-0.000000in}{0.000000in}}%
\pgfpathlineto{\pgfqpoint{-0.048611in}{0.000000in}}%
\pgfusepath{stroke,fill}%
}%
\begin{pgfscope}%
\pgfsys@transformshift{0.727161in}{1.875788in}%
\pgfsys@useobject{currentmarker}{}%
\end{pgfscope}%
\end{pgfscope}%
\begin{pgfscope}%
\definecolor{textcolor}{rgb}{0.000000,0.000000,0.000000}%
\pgfsetstrokecolor{textcolor}%
\pgfsetfillcolor{textcolor}%
\pgftext[x=0.344444in, y=1.827563in, left, base]{\color{textcolor}\rmfamily\fontsize{10.000000}{12.000000}\selectfont \(\displaystyle {\ensuremath{-}2.0}\)}%
\end{pgfscope}%
\begin{pgfscope}%
\pgfpathrectangle{\pgfqpoint{0.727161in}{0.565123in}}{\pgfqpoint{4.922838in}{2.784877in}}%
\pgfusepath{clip}%
\pgfsetbuttcap%
\pgfsetroundjoin%
\pgfsetlinewidth{0.803000pt}%
\definecolor{currentstroke}{rgb}{0.690196,0.690196,0.690196}%
\pgfsetstrokecolor{currentstroke}%
\pgfsetdash{{0.800000pt}{1.320000pt}}{0.000000pt}%
\pgfpathmoveto{\pgfqpoint{0.727161in}{2.202882in}}%
\pgfpathlineto{\pgfqpoint{5.649999in}{2.202882in}}%
\pgfusepath{stroke}%
\end{pgfscope}%
\begin{pgfscope}%
\pgfsetbuttcap%
\pgfsetroundjoin%
\definecolor{currentfill}{rgb}{0.000000,0.000000,0.000000}%
\pgfsetfillcolor{currentfill}%
\pgfsetlinewidth{0.803000pt}%
\definecolor{currentstroke}{rgb}{0.000000,0.000000,0.000000}%
\pgfsetstrokecolor{currentstroke}%
\pgfsetdash{}{0pt}%
\pgfsys@defobject{currentmarker}{\pgfqpoint{-0.048611in}{0.000000in}}{\pgfqpoint{-0.000000in}{0.000000in}}{%
\pgfpathmoveto{\pgfqpoint{-0.000000in}{0.000000in}}%
\pgfpathlineto{\pgfqpoint{-0.048611in}{0.000000in}}%
\pgfusepath{stroke,fill}%
}%
\begin{pgfscope}%
\pgfsys@transformshift{0.727161in}{2.202882in}%
\pgfsys@useobject{currentmarker}{}%
\end{pgfscope}%
\end{pgfscope}%
\begin{pgfscope}%
\definecolor{textcolor}{rgb}{0.000000,0.000000,0.000000}%
\pgfsetstrokecolor{textcolor}%
\pgfsetfillcolor{textcolor}%
\pgftext[x=0.344444in, y=2.154657in, left, base]{\color{textcolor}\rmfamily\fontsize{10.000000}{12.000000}\selectfont \(\displaystyle {\ensuremath{-}1.5}\)}%
\end{pgfscope}%
\begin{pgfscope}%
\pgfpathrectangle{\pgfqpoint{0.727161in}{0.565123in}}{\pgfqpoint{4.922838in}{2.784877in}}%
\pgfusepath{clip}%
\pgfsetbuttcap%
\pgfsetroundjoin%
\pgfsetlinewidth{0.803000pt}%
\definecolor{currentstroke}{rgb}{0.690196,0.690196,0.690196}%
\pgfsetstrokecolor{currentstroke}%
\pgfsetdash{{0.800000pt}{1.320000pt}}{0.000000pt}%
\pgfpathmoveto{\pgfqpoint{0.727161in}{2.529976in}}%
\pgfpathlineto{\pgfqpoint{5.649999in}{2.529976in}}%
\pgfusepath{stroke}%
\end{pgfscope}%
\begin{pgfscope}%
\pgfsetbuttcap%
\pgfsetroundjoin%
\definecolor{currentfill}{rgb}{0.000000,0.000000,0.000000}%
\pgfsetfillcolor{currentfill}%
\pgfsetlinewidth{0.803000pt}%
\definecolor{currentstroke}{rgb}{0.000000,0.000000,0.000000}%
\pgfsetstrokecolor{currentstroke}%
\pgfsetdash{}{0pt}%
\pgfsys@defobject{currentmarker}{\pgfqpoint{-0.048611in}{0.000000in}}{\pgfqpoint{-0.000000in}{0.000000in}}{%
\pgfpathmoveto{\pgfqpoint{-0.000000in}{0.000000in}}%
\pgfpathlineto{\pgfqpoint{-0.048611in}{0.000000in}}%
\pgfusepath{stroke,fill}%
}%
\begin{pgfscope}%
\pgfsys@transformshift{0.727161in}{2.529976in}%
\pgfsys@useobject{currentmarker}{}%
\end{pgfscope}%
\end{pgfscope}%
\begin{pgfscope}%
\definecolor{textcolor}{rgb}{0.000000,0.000000,0.000000}%
\pgfsetstrokecolor{textcolor}%
\pgfsetfillcolor{textcolor}%
\pgftext[x=0.344444in, y=2.481751in, left, base]{\color{textcolor}\rmfamily\fontsize{10.000000}{12.000000}\selectfont \(\displaystyle {\ensuremath{-}1.0}\)}%
\end{pgfscope}%
\begin{pgfscope}%
\pgfpathrectangle{\pgfqpoint{0.727161in}{0.565123in}}{\pgfqpoint{4.922838in}{2.784877in}}%
\pgfusepath{clip}%
\pgfsetbuttcap%
\pgfsetroundjoin%
\pgfsetlinewidth{0.803000pt}%
\definecolor{currentstroke}{rgb}{0.690196,0.690196,0.690196}%
\pgfsetstrokecolor{currentstroke}%
\pgfsetdash{{0.800000pt}{1.320000pt}}{0.000000pt}%
\pgfpathmoveto{\pgfqpoint{0.727161in}{2.857070in}}%
\pgfpathlineto{\pgfqpoint{5.649999in}{2.857070in}}%
\pgfusepath{stroke}%
\end{pgfscope}%
\begin{pgfscope}%
\pgfsetbuttcap%
\pgfsetroundjoin%
\definecolor{currentfill}{rgb}{0.000000,0.000000,0.000000}%
\pgfsetfillcolor{currentfill}%
\pgfsetlinewidth{0.803000pt}%
\definecolor{currentstroke}{rgb}{0.000000,0.000000,0.000000}%
\pgfsetstrokecolor{currentstroke}%
\pgfsetdash{}{0pt}%
\pgfsys@defobject{currentmarker}{\pgfqpoint{-0.048611in}{0.000000in}}{\pgfqpoint{-0.000000in}{0.000000in}}{%
\pgfpathmoveto{\pgfqpoint{-0.000000in}{0.000000in}}%
\pgfpathlineto{\pgfqpoint{-0.048611in}{0.000000in}}%
\pgfusepath{stroke,fill}%
}%
\begin{pgfscope}%
\pgfsys@transformshift{0.727161in}{2.857070in}%
\pgfsys@useobject{currentmarker}{}%
\end{pgfscope}%
\end{pgfscope}%
\begin{pgfscope}%
\definecolor{textcolor}{rgb}{0.000000,0.000000,0.000000}%
\pgfsetstrokecolor{textcolor}%
\pgfsetfillcolor{textcolor}%
\pgftext[x=0.344444in, y=2.808844in, left, base]{\color{textcolor}\rmfamily\fontsize{10.000000}{12.000000}\selectfont \(\displaystyle {\ensuremath{-}0.5}\)}%
\end{pgfscope}%
\begin{pgfscope}%
\pgfpathrectangle{\pgfqpoint{0.727161in}{0.565123in}}{\pgfqpoint{4.922838in}{2.784877in}}%
\pgfusepath{clip}%
\pgfsetbuttcap%
\pgfsetroundjoin%
\pgfsetlinewidth{0.803000pt}%
\definecolor{currentstroke}{rgb}{0.690196,0.690196,0.690196}%
\pgfsetstrokecolor{currentstroke}%
\pgfsetdash{{0.800000pt}{1.320000pt}}{0.000000pt}%
\pgfpathmoveto{\pgfqpoint{0.727161in}{3.184163in}}%
\pgfpathlineto{\pgfqpoint{5.649999in}{3.184163in}}%
\pgfusepath{stroke}%
\end{pgfscope}%
\begin{pgfscope}%
\pgfsetbuttcap%
\pgfsetroundjoin%
\definecolor{currentfill}{rgb}{0.000000,0.000000,0.000000}%
\pgfsetfillcolor{currentfill}%
\pgfsetlinewidth{0.803000pt}%
\definecolor{currentstroke}{rgb}{0.000000,0.000000,0.000000}%
\pgfsetstrokecolor{currentstroke}%
\pgfsetdash{}{0pt}%
\pgfsys@defobject{currentmarker}{\pgfqpoint{-0.048611in}{0.000000in}}{\pgfqpoint{-0.000000in}{0.000000in}}{%
\pgfpathmoveto{\pgfqpoint{-0.000000in}{0.000000in}}%
\pgfpathlineto{\pgfqpoint{-0.048611in}{0.000000in}}%
\pgfusepath{stroke,fill}%
}%
\begin{pgfscope}%
\pgfsys@transformshift{0.727161in}{3.184163in}%
\pgfsys@useobject{currentmarker}{}%
\end{pgfscope}%
\end{pgfscope}%
\begin{pgfscope}%
\definecolor{textcolor}{rgb}{0.000000,0.000000,0.000000}%
\pgfsetstrokecolor{textcolor}%
\pgfsetfillcolor{textcolor}%
\pgftext[x=0.452469in, y=3.135938in, left, base]{\color{textcolor}\rmfamily\fontsize{10.000000}{12.000000}\selectfont \(\displaystyle {0.0}\)}%
\end{pgfscope}%
\begin{pgfscope}%
\definecolor{textcolor}{rgb}{0.000000,0.000000,0.000000}%
\pgfsetstrokecolor{textcolor}%
\pgfsetfillcolor{textcolor}%
\pgftext[x=0.288889in,y=1.957562in,,bottom,rotate=90.000000]{\color{textcolor}\rmfamily\fontsize{10.000000}{12.000000}\selectfont Return Loss (dB)}%
\end{pgfscope}%
\begin{pgfscope}%
\pgfpathrectangle{\pgfqpoint{0.727161in}{0.565123in}}{\pgfqpoint{4.922838in}{2.784877in}}%
\pgfusepath{clip}%
\pgfsetrectcap%
\pgfsetroundjoin%
\pgfsetlinewidth{1.003750pt}%
\definecolor{currentstroke}{rgb}{0.000000,0.000000,0.000000}%
\pgfsetstrokecolor{currentstroke}%
\pgfsetdash{}{0pt}%
\pgfpathmoveto{\pgfqpoint{0.727161in}{1.843079in}}%
\pgfpathlineto{\pgfqpoint{0.732472in}{3.223415in}}%
\pgfpathlineto{\pgfqpoint{0.743093in}{3.144912in}}%
\pgfpathlineto{\pgfqpoint{0.748403in}{3.184163in}}%
\pgfpathlineto{\pgfqpoint{0.764335in}{3.184163in}}%
\pgfpathlineto{\pgfqpoint{0.769645in}{3.144912in}}%
\pgfpathlineto{\pgfqpoint{0.774956in}{3.184163in}}%
\pgfpathlineto{\pgfqpoint{0.780266in}{3.144912in}}%
\pgfpathlineto{\pgfqpoint{0.785577in}{3.144912in}}%
\pgfpathlineto{\pgfqpoint{0.790887in}{3.184163in}}%
\pgfpathlineto{\pgfqpoint{0.796198in}{3.144912in}}%
\pgfpathlineto{\pgfqpoint{0.812129in}{3.144912in}}%
\pgfpathlineto{\pgfqpoint{0.817440in}{3.184163in}}%
\pgfpathlineto{\pgfqpoint{0.822750in}{3.144912in}}%
\pgfpathlineto{\pgfqpoint{0.828061in}{3.144912in}}%
\pgfpathlineto{\pgfqpoint{0.833371in}{3.184163in}}%
\pgfpathlineto{\pgfqpoint{0.838682in}{3.144912in}}%
\pgfpathlineto{\pgfqpoint{0.849303in}{3.144912in}}%
\pgfpathlineto{\pgfqpoint{0.854614in}{3.184163in}}%
\pgfpathlineto{\pgfqpoint{0.859924in}{3.144912in}}%
\pgfpathlineto{\pgfqpoint{0.870545in}{3.144912in}}%
\pgfpathlineto{\pgfqpoint{0.875856in}{3.184163in}}%
\pgfpathlineto{\pgfqpoint{0.881166in}{3.144912in}}%
\pgfpathlineto{\pgfqpoint{1.449390in}{3.144912in}}%
\pgfpathlineto{\pgfqpoint{1.454701in}{3.105661in}}%
\pgfpathlineto{\pgfqpoint{1.460011in}{3.105661in}}%
\pgfpathlineto{\pgfqpoint{1.465322in}{3.144912in}}%
\pgfpathlineto{\pgfqpoint{1.470632in}{3.105661in}}%
\pgfpathlineto{\pgfqpoint{1.757399in}{3.105661in}}%
\pgfpathlineto{\pgfqpoint{1.762710in}{3.066410in}}%
\pgfpathlineto{\pgfqpoint{1.911404in}{3.066410in}}%
\pgfpathlineto{\pgfqpoint{1.916714in}{3.033700in}}%
\pgfpathlineto{\pgfqpoint{2.038856in}{3.033700in}}%
\pgfpathlineto{\pgfqpoint{2.044167in}{2.994449in}}%
\pgfpathlineto{\pgfqpoint{2.113203in}{2.994449in}}%
\pgfpathlineto{\pgfqpoint{2.118514in}{2.955198in}}%
\pgfpathlineto{\pgfqpoint{2.187550in}{2.955198in}}%
\pgfpathlineto{\pgfqpoint{2.192861in}{2.915947in}}%
\pgfpathlineto{\pgfqpoint{2.240655in}{2.915947in}}%
\pgfpathlineto{\pgfqpoint{2.245966in}{2.876695in}}%
\pgfpathlineto{\pgfqpoint{2.283139in}{2.876695in}}%
\pgfpathlineto{\pgfqpoint{2.288450in}{2.837444in}}%
\pgfpathlineto{\pgfqpoint{2.315002in}{2.837444in}}%
\pgfpathlineto{\pgfqpoint{2.320313in}{2.798193in}}%
\pgfpathlineto{\pgfqpoint{2.325623in}{2.837444in}}%
\pgfpathlineto{\pgfqpoint{2.330934in}{2.798193in}}%
\pgfpathlineto{\pgfqpoint{2.357486in}{2.798193in}}%
\pgfpathlineto{\pgfqpoint{2.362797in}{2.758941in}}%
\pgfpathlineto{\pgfqpoint{2.384039in}{2.758941in}}%
\pgfpathlineto{\pgfqpoint{2.389349in}{2.726232in}}%
\pgfpathlineto{\pgfqpoint{2.410591in}{2.726232in}}%
\pgfpathlineto{\pgfqpoint{2.415902in}{2.686981in}}%
\pgfpathlineto{\pgfqpoint{2.431833in}{2.686981in}}%
\pgfpathlineto{\pgfqpoint{2.437144in}{2.647730in}}%
\pgfpathlineto{\pgfqpoint{2.458386in}{2.647730in}}%
\pgfpathlineto{\pgfqpoint{2.463696in}{2.608478in}}%
\pgfpathlineto{\pgfqpoint{2.479628in}{2.608478in}}%
\pgfpathlineto{\pgfqpoint{2.484938in}{2.569227in}}%
\pgfpathlineto{\pgfqpoint{2.495559in}{2.569227in}}%
\pgfpathlineto{\pgfqpoint{2.500870in}{2.529976in}}%
\pgfpathlineto{\pgfqpoint{2.511491in}{2.529976in}}%
\pgfpathlineto{\pgfqpoint{2.516801in}{2.490725in}}%
\pgfpathlineto{\pgfqpoint{2.527423in}{2.490725in}}%
\pgfpathlineto{\pgfqpoint{2.532733in}{2.458015in}}%
\pgfpathlineto{\pgfqpoint{2.543354in}{2.458015in}}%
\pgfpathlineto{\pgfqpoint{2.548665in}{2.418764in}}%
\pgfpathlineto{\pgfqpoint{2.553975in}{2.418764in}}%
\pgfpathlineto{\pgfqpoint{2.559286in}{2.379513in}}%
\pgfpathlineto{\pgfqpoint{2.564596in}{2.379513in}}%
\pgfpathlineto{\pgfqpoint{2.569907in}{2.340261in}}%
\pgfpathlineto{\pgfqpoint{2.580528in}{2.340261in}}%
\pgfpathlineto{\pgfqpoint{2.585838in}{2.301010in}}%
\pgfpathlineto{\pgfqpoint{2.591149in}{2.301010in}}%
\pgfpathlineto{\pgfqpoint{2.596459in}{2.261759in}}%
\pgfpathlineto{\pgfqpoint{2.601770in}{2.261759in}}%
\pgfpathlineto{\pgfqpoint{2.607080in}{2.222508in}}%
\pgfpathlineto{\pgfqpoint{2.612391in}{2.222508in}}%
\pgfpathlineto{\pgfqpoint{2.623012in}{2.150547in}}%
\pgfpathlineto{\pgfqpoint{2.628322in}{2.150547in}}%
\pgfpathlineto{\pgfqpoint{2.633633in}{2.111296in}}%
\pgfpathlineto{\pgfqpoint{2.638943in}{2.111296in}}%
\pgfpathlineto{\pgfqpoint{2.649564in}{2.032793in}}%
\pgfpathlineto{\pgfqpoint{2.654875in}{2.032793in}}%
\pgfpathlineto{\pgfqpoint{2.660185in}{1.993542in}}%
\pgfpathlineto{\pgfqpoint{2.665496in}{1.993542in}}%
\pgfpathlineto{\pgfqpoint{2.681427in}{1.882330in}}%
\pgfpathlineto{\pgfqpoint{2.686738in}{1.882330in}}%
\pgfpathlineto{\pgfqpoint{2.697359in}{1.803828in}}%
\pgfpathlineto{\pgfqpoint{2.702669in}{1.803828in}}%
\pgfpathlineto{\pgfqpoint{2.729222in}{1.614113in}}%
\pgfpathlineto{\pgfqpoint{2.734532in}{1.614113in}}%
\pgfpathlineto{\pgfqpoint{2.766395in}{1.378606in}}%
\pgfpathlineto{\pgfqpoint{2.771706in}{1.378606in}}%
\pgfpathlineto{\pgfqpoint{2.814190in}{1.071137in}}%
\pgfpathlineto{\pgfqpoint{2.819500in}{1.071137in}}%
\pgfpathlineto{\pgfqpoint{2.830121in}{0.999177in}}%
\pgfpathlineto{\pgfqpoint{2.835432in}{0.959926in}}%
\pgfpathlineto{\pgfqpoint{2.840742in}{0.959926in}}%
\pgfpathlineto{\pgfqpoint{2.851363in}{0.881423in}}%
\pgfpathlineto{\pgfqpoint{2.856674in}{0.881423in}}%
\pgfpathlineto{\pgfqpoint{2.867295in}{0.802921in}}%
\pgfpathlineto{\pgfqpoint{2.877916in}{0.802921in}}%
\pgfpathlineto{\pgfqpoint{2.883226in}{0.763669in}}%
\pgfpathlineto{\pgfqpoint{2.888537in}{0.763669in}}%
\pgfpathlineto{\pgfqpoint{2.893847in}{0.730960in}}%
\pgfpathlineto{\pgfqpoint{2.920400in}{0.730960in}}%
\pgfpathlineto{\pgfqpoint{2.925710in}{0.691709in}}%
\pgfpathlineto{\pgfqpoint{2.931021in}{0.730960in}}%
\pgfpathlineto{\pgfqpoint{2.946952in}{0.730960in}}%
\pgfpathlineto{\pgfqpoint{2.952263in}{0.763669in}}%
\pgfpathlineto{\pgfqpoint{2.957573in}{0.763669in}}%
\pgfpathlineto{\pgfqpoint{2.962884in}{0.802921in}}%
\pgfpathlineto{\pgfqpoint{2.968194in}{0.802921in}}%
\pgfpathlineto{\pgfqpoint{2.973505in}{0.842172in}}%
\pgfpathlineto{\pgfqpoint{2.978815in}{0.842172in}}%
\pgfpathlineto{\pgfqpoint{2.984126in}{0.881423in}}%
\pgfpathlineto{\pgfqpoint{2.989436in}{0.881423in}}%
\pgfpathlineto{\pgfqpoint{3.010678in}{1.038428in}}%
\pgfpathlineto{\pgfqpoint{3.015989in}{1.038428in}}%
\pgfpathlineto{\pgfqpoint{3.026610in}{1.110389in}}%
\pgfpathlineto{\pgfqpoint{3.058473in}{1.339354in}}%
\pgfpathlineto{\pgfqpoint{3.063783in}{1.339354in}}%
\pgfpathlineto{\pgfqpoint{3.106268in}{1.646823in}}%
\pgfpathlineto{\pgfqpoint{3.111578in}{1.646823in}}%
\pgfpathlineto{\pgfqpoint{3.132820in}{1.803828in}}%
\pgfpathlineto{\pgfqpoint{3.138131in}{1.803828in}}%
\pgfpathlineto{\pgfqpoint{3.159373in}{1.954291in}}%
\pgfpathlineto{\pgfqpoint{3.164683in}{1.954291in}}%
\pgfpathlineto{\pgfqpoint{3.175304in}{2.032793in}}%
\pgfpathlineto{\pgfqpoint{3.180615in}{2.032793in}}%
\pgfpathlineto{\pgfqpoint{3.191236in}{2.111296in}}%
\pgfpathlineto{\pgfqpoint{3.196546in}{2.111296in}}%
\pgfpathlineto{\pgfqpoint{3.201857in}{2.150547in}}%
\pgfpathlineto{\pgfqpoint{3.207167in}{2.150547in}}%
\pgfpathlineto{\pgfqpoint{3.212478in}{2.189798in}}%
\pgfpathlineto{\pgfqpoint{3.217788in}{2.189798in}}%
\pgfpathlineto{\pgfqpoint{3.228409in}{2.261759in}}%
\pgfpathlineto{\pgfqpoint{3.233720in}{2.261759in}}%
\pgfpathlineto{\pgfqpoint{3.239030in}{2.301010in}}%
\pgfpathlineto{\pgfqpoint{3.244341in}{2.301010in}}%
\pgfpathlineto{\pgfqpoint{3.249651in}{2.340261in}}%
\pgfpathlineto{\pgfqpoint{3.254962in}{2.340261in}}%
\pgfpathlineto{\pgfqpoint{3.260272in}{2.379513in}}%
\pgfpathlineto{\pgfqpoint{3.270893in}{2.379513in}}%
\pgfpathlineto{\pgfqpoint{3.276204in}{2.418764in}}%
\pgfpathlineto{\pgfqpoint{3.281514in}{2.418764in}}%
\pgfpathlineto{\pgfqpoint{3.286825in}{2.458015in}}%
\pgfpathlineto{\pgfqpoint{3.297446in}{2.458015in}}%
\pgfpathlineto{\pgfqpoint{3.302756in}{2.490725in}}%
\pgfpathlineto{\pgfqpoint{3.308067in}{2.490725in}}%
\pgfpathlineto{\pgfqpoint{3.313377in}{2.529976in}}%
\pgfpathlineto{\pgfqpoint{3.329309in}{2.529976in}}%
\pgfpathlineto{\pgfqpoint{3.334619in}{2.569227in}}%
\pgfpathlineto{\pgfqpoint{3.345240in}{2.569227in}}%
\pgfpathlineto{\pgfqpoint{3.350551in}{2.608478in}}%
\pgfpathlineto{\pgfqpoint{3.366482in}{2.608478in}}%
\pgfpathlineto{\pgfqpoint{3.371793in}{2.647730in}}%
\pgfpathlineto{\pgfqpoint{3.393035in}{2.647730in}}%
\pgfpathlineto{\pgfqpoint{3.398345in}{2.686981in}}%
\pgfpathlineto{\pgfqpoint{3.414277in}{2.686981in}}%
\pgfpathlineto{\pgfqpoint{3.419587in}{2.726232in}}%
\pgfpathlineto{\pgfqpoint{3.446140in}{2.726232in}}%
\pgfpathlineto{\pgfqpoint{3.451450in}{2.758941in}}%
\pgfpathlineto{\pgfqpoint{3.478003in}{2.758941in}}%
\pgfpathlineto{\pgfqpoint{3.483313in}{2.798193in}}%
\pgfpathlineto{\pgfqpoint{3.515176in}{2.798193in}}%
\pgfpathlineto{\pgfqpoint{3.520487in}{2.837444in}}%
\pgfpathlineto{\pgfqpoint{3.568281in}{2.837444in}}%
\pgfpathlineto{\pgfqpoint{3.573592in}{2.876695in}}%
\pgfpathlineto{\pgfqpoint{3.610765in}{2.876695in}}%
\pgfpathlineto{\pgfqpoint{3.616076in}{2.915947in}}%
\pgfpathlineto{\pgfqpoint{3.674492in}{2.915947in}}%
\pgfpathlineto{\pgfqpoint{3.679802in}{2.955198in}}%
\pgfpathlineto{\pgfqpoint{3.759460in}{2.955198in}}%
\pgfpathlineto{\pgfqpoint{3.764770in}{2.994449in}}%
\pgfpathlineto{\pgfqpoint{3.855049in}{2.994449in}}%
\pgfpathlineto{\pgfqpoint{3.860359in}{3.033700in}}%
\pgfpathlineto{\pgfqpoint{3.982501in}{3.033700in}}%
\pgfpathlineto{\pgfqpoint{3.987811in}{3.066410in}}%
\pgfpathlineto{\pgfqpoint{4.163058in}{3.066410in}}%
\pgfpathlineto{\pgfqpoint{4.168368in}{3.105661in}}%
\pgfpathlineto{\pgfqpoint{4.471067in}{3.105661in}}%
\pgfpathlineto{\pgfqpoint{4.476378in}{3.144912in}}%
\pgfpathlineto{\pgfqpoint{4.481688in}{3.105661in}}%
\pgfpathlineto{\pgfqpoint{4.486999in}{3.144912in}}%
\pgfpathlineto{\pgfqpoint{4.492309in}{3.105661in}}%
\pgfpathlineto{\pgfqpoint{4.497620in}{3.144912in}}%
\pgfpathlineto{\pgfqpoint{4.502930in}{3.105661in}}%
\pgfpathlineto{\pgfqpoint{4.508241in}{3.144912in}}%
\pgfpathlineto{\pgfqpoint{5.103017in}{3.144912in}}%
\pgfpathlineto{\pgfqpoint{5.108328in}{3.184163in}}%
\pgfpathlineto{\pgfqpoint{5.113638in}{3.144912in}}%
\pgfpathlineto{\pgfqpoint{5.129570in}{3.144912in}}%
\pgfpathlineto{\pgfqpoint{5.134880in}{3.184163in}}%
\pgfpathlineto{\pgfqpoint{5.140191in}{3.144912in}}%
\pgfpathlineto{\pgfqpoint{5.156122in}{3.144912in}}%
\pgfpathlineto{\pgfqpoint{5.161433in}{3.184163in}}%
\pgfpathlineto{\pgfqpoint{5.166743in}{3.184163in}}%
\pgfpathlineto{\pgfqpoint{5.172054in}{3.144912in}}%
\pgfpathlineto{\pgfqpoint{5.177364in}{3.144912in}}%
\pgfpathlineto{\pgfqpoint{5.182675in}{3.184163in}}%
\pgfpathlineto{\pgfqpoint{5.203917in}{3.184163in}}%
\pgfpathlineto{\pgfqpoint{5.209227in}{3.144912in}}%
\pgfpathlineto{\pgfqpoint{5.214538in}{3.184163in}}%
\pgfpathlineto{\pgfqpoint{5.246401in}{3.184163in}}%
\pgfpathlineto{\pgfqpoint{5.251711in}{3.144912in}}%
\pgfpathlineto{\pgfqpoint{5.257022in}{3.184163in}}%
\pgfpathlineto{\pgfqpoint{5.262332in}{3.184163in}}%
\pgfpathlineto{\pgfqpoint{5.267643in}{3.144912in}}%
\pgfpathlineto{\pgfqpoint{5.272953in}{3.144912in}}%
\pgfpathlineto{\pgfqpoint{5.278264in}{3.184163in}}%
\pgfpathlineto{\pgfqpoint{5.649999in}{3.184163in}}%
\pgfpathlineto{\pgfqpoint{5.649999in}{3.184163in}}%
\pgfusepath{stroke}%
\end{pgfscope}%
\begin{pgfscope}%
\pgfpathrectangle{\pgfqpoint{0.727161in}{0.565123in}}{\pgfqpoint{4.922838in}{2.784877in}}%
\pgfusepath{clip}%
\pgfsetbuttcap%
\pgfsetroundjoin%
\pgfsetlinewidth{1.003750pt}%
\definecolor{currentstroke}{rgb}{1.000000,0.000000,0.000000}%
\pgfsetstrokecolor{currentstroke}%
\pgfsetdash{{1.000000pt}{1.650000pt}}{0.000000pt}%
\pgfpathmoveto{\pgfqpoint{0.727161in}{2.654271in}}%
\pgfpathlineto{\pgfqpoint{5.649999in}{2.654271in}}%
\pgfusepath{stroke}%
\end{pgfscope}%
\begin{pgfscope}%
\pgfsetrectcap%
\pgfsetmiterjoin%
\pgfsetlinewidth{0.803000pt}%
\definecolor{currentstroke}{rgb}{0.000000,0.000000,0.000000}%
\pgfsetstrokecolor{currentstroke}%
\pgfsetdash{}{0pt}%
\pgfpathmoveto{\pgfqpoint{0.727161in}{0.565123in}}%
\pgfpathlineto{\pgfqpoint{0.727161in}{3.350000in}}%
\pgfusepath{stroke}%
\end{pgfscope}%
\begin{pgfscope}%
\pgfsetrectcap%
\pgfsetmiterjoin%
\pgfsetlinewidth{0.803000pt}%
\definecolor{currentstroke}{rgb}{0.000000,0.000000,0.000000}%
\pgfsetstrokecolor{currentstroke}%
\pgfsetdash{}{0pt}%
\pgfpathmoveto{\pgfqpoint{5.649999in}{0.565123in}}%
\pgfpathlineto{\pgfqpoint{5.649999in}{3.350000in}}%
\pgfusepath{stroke}%
\end{pgfscope}%
\begin{pgfscope}%
\pgfsetrectcap%
\pgfsetmiterjoin%
\pgfsetlinewidth{0.803000pt}%
\definecolor{currentstroke}{rgb}{0.000000,0.000000,0.000000}%
\pgfsetstrokecolor{currentstroke}%
\pgfsetdash{}{0pt}%
\pgfpathmoveto{\pgfqpoint{0.727161in}{0.565123in}}%
\pgfpathlineto{\pgfqpoint{5.649999in}{0.565123in}}%
\pgfusepath{stroke}%
\end{pgfscope}%
\begin{pgfscope}%
\pgfsetrectcap%
\pgfsetmiterjoin%
\pgfsetlinewidth{0.803000pt}%
\definecolor{currentstroke}{rgb}{0.000000,0.000000,0.000000}%
\pgfsetstrokecolor{currentstroke}%
\pgfsetdash{}{0pt}%
\pgfpathmoveto{\pgfqpoint{0.727161in}{3.350000in}}%
\pgfpathlineto{\pgfqpoint{5.649999in}{3.350000in}}%
\pgfusepath{stroke}%
\end{pgfscope}%
\begin{pgfscope}%
\pgfsetbuttcap%
\pgfsetmiterjoin%
\definecolor{currentfill}{rgb}{1.000000,1.000000,1.000000}%
\pgfsetfillcolor{currentfill}%
\pgfsetfillopacity{0.800000}%
\pgfsetlinewidth{1.003750pt}%
\definecolor{currentstroke}{rgb}{0.800000,0.800000,0.800000}%
\pgfsetstrokecolor{currentstroke}%
\pgfsetstrokeopacity{0.800000}%
\pgfsetdash{}{0pt}%
\pgfpathmoveto{\pgfqpoint{4.039658in}{1.631559in}}%
\pgfpathlineto{\pgfqpoint{5.552777in}{1.631559in}}%
\pgfpathquadraticcurveto{\pgfqpoint{5.580555in}{1.631559in}}{\pgfqpoint{5.580555in}{1.659336in}}%
\pgfpathlineto{\pgfqpoint{5.580555in}{2.255787in}}%
\pgfpathquadraticcurveto{\pgfqpoint{5.580555in}{2.283565in}}{\pgfqpoint{5.552777in}{2.283565in}}%
\pgfpathlineto{\pgfqpoint{4.039658in}{2.283565in}}%
\pgfpathquadraticcurveto{\pgfqpoint{4.011880in}{2.283565in}}{\pgfqpoint{4.011880in}{2.255787in}}%
\pgfpathlineto{\pgfqpoint{4.011880in}{1.659336in}}%
\pgfpathquadraticcurveto{\pgfqpoint{4.011880in}{1.631559in}}{\pgfqpoint{4.039658in}{1.631559in}}%
\pgfpathlineto{\pgfqpoint{4.039658in}{1.631559in}}%
\pgfpathclose%
\pgfusepath{stroke,fill}%
\end{pgfscope}%
\begin{pgfscope}%
\pgfsetrectcap%
\pgfsetroundjoin%
\pgfsetlinewidth{1.003750pt}%
\definecolor{currentstroke}{rgb}{0.000000,0.000000,0.000000}%
\pgfsetstrokecolor{currentstroke}%
\pgfsetdash{}{0pt}%
\pgfpathmoveto{\pgfqpoint{4.067436in}{2.172454in}}%
\pgfpathlineto{\pgfqpoint{4.206325in}{2.172454in}}%
\pgfpathlineto{\pgfqpoint{4.345213in}{2.172454in}}%
\pgfusepath{stroke}%
\end{pgfscope}%
\begin{pgfscope}%
\definecolor{textcolor}{rgb}{0.000000,0.000000,0.000000}%
\pgfsetstrokecolor{textcolor}%
\pgfsetfillcolor{textcolor}%
\pgftext[x=4.456325in,y=2.123843in,left,base]{\color{textcolor}\rmfamily\fontsize{10.000000}{12.000000}\selectfont Return Loss (dB)}%
\end{pgfscope}%
\begin{pgfscope}%
\pgfsetbuttcap%
\pgfsetroundjoin%
\pgfsetlinewidth{1.003750pt}%
\definecolor{currentstroke}{rgb}{1.000000,0.000000,0.000000}%
\pgfsetstrokecolor{currentstroke}%
\pgfsetdash{{1.000000pt}{1.650000pt}}{0.000000pt}%
\pgfpathmoveto{\pgfqpoint{4.067436in}{1.971065in}}%
\pgfpathlineto{\pgfqpoint{4.206325in}{1.971065in}}%
\pgfpathlineto{\pgfqpoint{4.345213in}{1.971065in}}%
\pgfusepath{stroke}%
\end{pgfscope}%
\begin{pgfscope}%
\definecolor{textcolor}{rgb}{0.000000,0.000000,0.000000}%
\pgfsetstrokecolor{textcolor}%
\pgfsetfillcolor{textcolor}%
\pgftext[x=4.456325in,y=1.922454in,left,base]{\color{textcolor}\rmfamily\fontsize{10.000000}{12.000000}\selectfont 3 dB level}%
\end{pgfscope}%
\begin{pgfscope}%
\pgfsetbuttcap%
\pgfsetroundjoin%
\pgfsetlinewidth{1.003750pt}%
\definecolor{currentstroke}{rgb}{0.000000,0.000000,0.000000}%
\pgfsetstrokecolor{currentstroke}%
\pgfsetdash{{3.700000pt}{1.600000pt}}{0.000000pt}%
\pgfpathmoveto{\pgfqpoint{4.067436in}{1.770447in}}%
\pgfpathlineto{\pgfqpoint{4.206325in}{1.770447in}}%
\pgfpathlineto{\pgfqpoint{4.345213in}{1.770447in}}%
\pgfusepath{stroke}%
\end{pgfscope}%
\begin{pgfscope}%
\definecolor{textcolor}{rgb}{0.000000,0.000000,0.000000}%
\pgfsetstrokecolor{textcolor}%
\pgfsetfillcolor{textcolor}%
\pgftext[x=4.456325in,y=1.721836in,left,base]{\color{textcolor}\rmfamily\fontsize{10.000000}{12.000000}\selectfont Phase (deg)}%
\end{pgfscope}%
\begin{pgfscope}%
\pgfsetbuttcap%
\pgfsetroundjoin%
\definecolor{currentfill}{rgb}{0.000000,0.000000,0.000000}%
\pgfsetfillcolor{currentfill}%
\pgfsetlinewidth{0.803000pt}%
\definecolor{currentstroke}{rgb}{0.000000,0.000000,0.000000}%
\pgfsetstrokecolor{currentstroke}%
\pgfsetdash{}{0pt}%
\pgfsys@defobject{currentmarker}{\pgfqpoint{0.000000in}{0.000000in}}{\pgfqpoint{0.048611in}{0.000000in}}{%
\pgfpathmoveto{\pgfqpoint{0.000000in}{0.000000in}}%
\pgfpathlineto{\pgfqpoint{0.048611in}{0.000000in}}%
\pgfusepath{stroke,fill}%
}%
\begin{pgfscope}%
\pgfsys@transformshift{5.649999in}{0.654024in}%
\pgfsys@useobject{currentmarker}{}%
\end{pgfscope}%
\end{pgfscope}%
\begin{pgfscope}%
\definecolor{textcolor}{rgb}{0.000000,0.000000,0.000000}%
\pgfsetstrokecolor{textcolor}%
\pgfsetfillcolor{textcolor}%
\pgftext[x=5.747222in, y=0.605799in, left, base]{\color{textcolor}\rmfamily\fontsize{10.000000}{12.000000}\selectfont \(\displaystyle {0}\)}%
\end{pgfscope}%
\begin{pgfscope}%
\pgfsetbuttcap%
\pgfsetroundjoin%
\definecolor{currentfill}{rgb}{0.000000,0.000000,0.000000}%
\pgfsetfillcolor{currentfill}%
\pgfsetlinewidth{0.803000pt}%
\definecolor{currentstroke}{rgb}{0.000000,0.000000,0.000000}%
\pgfsetstrokecolor{currentstroke}%
\pgfsetdash{}{0pt}%
\pgfsys@defobject{currentmarker}{\pgfqpoint{0.000000in}{0.000000in}}{\pgfqpoint{0.048611in}{0.000000in}}{%
\pgfpathmoveto{\pgfqpoint{0.000000in}{0.000000in}}%
\pgfpathlineto{\pgfqpoint{0.048611in}{0.000000in}}%
\pgfusepath{stroke,fill}%
}%
\begin{pgfscope}%
\pgfsys@transformshift{5.649999in}{1.010884in}%
\pgfsys@useobject{currentmarker}{}%
\end{pgfscope}%
\end{pgfscope}%
\begin{pgfscope}%
\definecolor{textcolor}{rgb}{0.000000,0.000000,0.000000}%
\pgfsetstrokecolor{textcolor}%
\pgfsetfillcolor{textcolor}%
\pgftext[x=5.747222in, y=0.962659in, left, base]{\color{textcolor}\rmfamily\fontsize{10.000000}{12.000000}\selectfont \(\displaystyle {25}\)}%
\end{pgfscope}%
\begin{pgfscope}%
\pgfsetbuttcap%
\pgfsetroundjoin%
\definecolor{currentfill}{rgb}{0.000000,0.000000,0.000000}%
\pgfsetfillcolor{currentfill}%
\pgfsetlinewidth{0.803000pt}%
\definecolor{currentstroke}{rgb}{0.000000,0.000000,0.000000}%
\pgfsetstrokecolor{currentstroke}%
\pgfsetdash{}{0pt}%
\pgfsys@defobject{currentmarker}{\pgfqpoint{0.000000in}{0.000000in}}{\pgfqpoint{0.048611in}{0.000000in}}{%
\pgfpathmoveto{\pgfqpoint{0.000000in}{0.000000in}}%
\pgfpathlineto{\pgfqpoint{0.048611in}{0.000000in}}%
\pgfusepath{stroke,fill}%
}%
\begin{pgfscope}%
\pgfsys@transformshift{5.649999in}{1.367744in}%
\pgfsys@useobject{currentmarker}{}%
\end{pgfscope}%
\end{pgfscope}%
\begin{pgfscope}%
\definecolor{textcolor}{rgb}{0.000000,0.000000,0.000000}%
\pgfsetstrokecolor{textcolor}%
\pgfsetfillcolor{textcolor}%
\pgftext[x=5.747222in, y=1.319519in, left, base]{\color{textcolor}\rmfamily\fontsize{10.000000}{12.000000}\selectfont \(\displaystyle {50}\)}%
\end{pgfscope}%
\begin{pgfscope}%
\pgfsetbuttcap%
\pgfsetroundjoin%
\definecolor{currentfill}{rgb}{0.000000,0.000000,0.000000}%
\pgfsetfillcolor{currentfill}%
\pgfsetlinewidth{0.803000pt}%
\definecolor{currentstroke}{rgb}{0.000000,0.000000,0.000000}%
\pgfsetstrokecolor{currentstroke}%
\pgfsetdash{}{0pt}%
\pgfsys@defobject{currentmarker}{\pgfqpoint{0.000000in}{0.000000in}}{\pgfqpoint{0.048611in}{0.000000in}}{%
\pgfpathmoveto{\pgfqpoint{0.000000in}{0.000000in}}%
\pgfpathlineto{\pgfqpoint{0.048611in}{0.000000in}}%
\pgfusepath{stroke,fill}%
}%
\begin{pgfscope}%
\pgfsys@transformshift{5.649999in}{1.724604in}%
\pgfsys@useobject{currentmarker}{}%
\end{pgfscope}%
\end{pgfscope}%
\begin{pgfscope}%
\definecolor{textcolor}{rgb}{0.000000,0.000000,0.000000}%
\pgfsetstrokecolor{textcolor}%
\pgfsetfillcolor{textcolor}%
\pgftext[x=5.747222in, y=1.676378in, left, base]{\color{textcolor}\rmfamily\fontsize{10.000000}{12.000000}\selectfont \(\displaystyle {75}\)}%
\end{pgfscope}%
\begin{pgfscope}%
\pgfsetbuttcap%
\pgfsetroundjoin%
\definecolor{currentfill}{rgb}{0.000000,0.000000,0.000000}%
\pgfsetfillcolor{currentfill}%
\pgfsetlinewidth{0.803000pt}%
\definecolor{currentstroke}{rgb}{0.000000,0.000000,0.000000}%
\pgfsetstrokecolor{currentstroke}%
\pgfsetdash{}{0pt}%
\pgfsys@defobject{currentmarker}{\pgfqpoint{0.000000in}{0.000000in}}{\pgfqpoint{0.048611in}{0.000000in}}{%
\pgfpathmoveto{\pgfqpoint{0.000000in}{0.000000in}}%
\pgfpathlineto{\pgfqpoint{0.048611in}{0.000000in}}%
\pgfusepath{stroke,fill}%
}%
\begin{pgfscope}%
\pgfsys@transformshift{5.649999in}{2.081463in}%
\pgfsys@useobject{currentmarker}{}%
\end{pgfscope}%
\end{pgfscope}%
\begin{pgfscope}%
\definecolor{textcolor}{rgb}{0.000000,0.000000,0.000000}%
\pgfsetstrokecolor{textcolor}%
\pgfsetfillcolor{textcolor}%
\pgftext[x=5.747222in, y=2.033238in, left, base]{\color{textcolor}\rmfamily\fontsize{10.000000}{12.000000}\selectfont \(\displaystyle {100}\)}%
\end{pgfscope}%
\begin{pgfscope}%
\pgfsetbuttcap%
\pgfsetroundjoin%
\definecolor{currentfill}{rgb}{0.000000,0.000000,0.000000}%
\pgfsetfillcolor{currentfill}%
\pgfsetlinewidth{0.803000pt}%
\definecolor{currentstroke}{rgb}{0.000000,0.000000,0.000000}%
\pgfsetstrokecolor{currentstroke}%
\pgfsetdash{}{0pt}%
\pgfsys@defobject{currentmarker}{\pgfqpoint{0.000000in}{0.000000in}}{\pgfqpoint{0.048611in}{0.000000in}}{%
\pgfpathmoveto{\pgfqpoint{0.000000in}{0.000000in}}%
\pgfpathlineto{\pgfqpoint{0.048611in}{0.000000in}}%
\pgfusepath{stroke,fill}%
}%
\begin{pgfscope}%
\pgfsys@transformshift{5.649999in}{2.438323in}%
\pgfsys@useobject{currentmarker}{}%
\end{pgfscope}%
\end{pgfscope}%
\begin{pgfscope}%
\definecolor{textcolor}{rgb}{0.000000,0.000000,0.000000}%
\pgfsetstrokecolor{textcolor}%
\pgfsetfillcolor{textcolor}%
\pgftext[x=5.747222in, y=2.390098in, left, base]{\color{textcolor}\rmfamily\fontsize{10.000000}{12.000000}\selectfont \(\displaystyle {125}\)}%
\end{pgfscope}%
\begin{pgfscope}%
\pgfsetbuttcap%
\pgfsetroundjoin%
\definecolor{currentfill}{rgb}{0.000000,0.000000,0.000000}%
\pgfsetfillcolor{currentfill}%
\pgfsetlinewidth{0.803000pt}%
\definecolor{currentstroke}{rgb}{0.000000,0.000000,0.000000}%
\pgfsetstrokecolor{currentstroke}%
\pgfsetdash{}{0pt}%
\pgfsys@defobject{currentmarker}{\pgfqpoint{0.000000in}{0.000000in}}{\pgfqpoint{0.048611in}{0.000000in}}{%
\pgfpathmoveto{\pgfqpoint{0.000000in}{0.000000in}}%
\pgfpathlineto{\pgfqpoint{0.048611in}{0.000000in}}%
\pgfusepath{stroke,fill}%
}%
\begin{pgfscope}%
\pgfsys@transformshift{5.649999in}{2.795183in}%
\pgfsys@useobject{currentmarker}{}%
\end{pgfscope}%
\end{pgfscope}%
\begin{pgfscope}%
\definecolor{textcolor}{rgb}{0.000000,0.000000,0.000000}%
\pgfsetstrokecolor{textcolor}%
\pgfsetfillcolor{textcolor}%
\pgftext[x=5.747222in, y=2.746958in, left, base]{\color{textcolor}\rmfamily\fontsize{10.000000}{12.000000}\selectfont \(\displaystyle {150}\)}%
\end{pgfscope}%
\begin{pgfscope}%
\pgfsetbuttcap%
\pgfsetroundjoin%
\definecolor{currentfill}{rgb}{0.000000,0.000000,0.000000}%
\pgfsetfillcolor{currentfill}%
\pgfsetlinewidth{0.803000pt}%
\definecolor{currentstroke}{rgb}{0.000000,0.000000,0.000000}%
\pgfsetstrokecolor{currentstroke}%
\pgfsetdash{}{0pt}%
\pgfsys@defobject{currentmarker}{\pgfqpoint{0.000000in}{0.000000in}}{\pgfqpoint{0.048611in}{0.000000in}}{%
\pgfpathmoveto{\pgfqpoint{0.000000in}{0.000000in}}%
\pgfpathlineto{\pgfqpoint{0.048611in}{0.000000in}}%
\pgfusepath{stroke,fill}%
}%
\begin{pgfscope}%
\pgfsys@transformshift{5.649999in}{3.152043in}%
\pgfsys@useobject{currentmarker}{}%
\end{pgfscope}%
\end{pgfscope}%
\begin{pgfscope}%
\definecolor{textcolor}{rgb}{0.000000,0.000000,0.000000}%
\pgfsetstrokecolor{textcolor}%
\pgfsetfillcolor{textcolor}%
\pgftext[x=5.747222in, y=3.103817in, left, base]{\color{textcolor}\rmfamily\fontsize{10.000000}{12.000000}\selectfont \(\displaystyle {175}\)}%
\end{pgfscope}%
\begin{pgfscope}%
\definecolor{textcolor}{rgb}{0.000000,0.000000,0.000000}%
\pgfsetstrokecolor{textcolor}%
\pgfsetfillcolor{textcolor}%
\pgftext[x=6.011111in,y=1.957562in,,top,rotate=90.000000]{\color{textcolor}\rmfamily\fontsize{10.000000}{12.000000}\selectfont Phase (deg)}%
\end{pgfscope}%
\begin{pgfscope}%
\pgfpathrectangle{\pgfqpoint{0.727161in}{0.565123in}}{\pgfqpoint{4.922838in}{2.784877in}}%
\pgfusepath{clip}%
\pgfsetbuttcap%
\pgfsetroundjoin%
\pgfsetlinewidth{1.003750pt}%
\definecolor{currentstroke}{rgb}{0.000000,0.000000,0.000000}%
\pgfsetstrokecolor{currentstroke}%
\pgfsetdash{{3.700000pt}{1.600000pt}}{0.000000pt}%
\pgfpathmoveto{\pgfqpoint{0.727161in}{1.533041in}}%
\pgfpathlineto{\pgfqpoint{0.732472in}{0.837307in}}%
\pgfpathlineto{\pgfqpoint{0.737782in}{0.842446in}}%
\pgfpathlineto{\pgfqpoint{0.743093in}{0.837307in}}%
\pgfpathlineto{\pgfqpoint{0.759024in}{0.837307in}}%
\pgfpathlineto{\pgfqpoint{0.764335in}{0.839877in}}%
\pgfpathlineto{\pgfqpoint{0.790887in}{0.839877in}}%
\pgfpathlineto{\pgfqpoint{0.796198in}{0.842446in}}%
\pgfpathlineto{\pgfqpoint{0.817440in}{0.842446in}}%
\pgfpathlineto{\pgfqpoint{0.822750in}{0.844873in}}%
\pgfpathlineto{\pgfqpoint{0.838682in}{0.844873in}}%
\pgfpathlineto{\pgfqpoint{0.843992in}{0.847442in}}%
\pgfpathlineto{\pgfqpoint{0.870545in}{0.847442in}}%
\pgfpathlineto{\pgfqpoint{0.875856in}{0.849869in}}%
\pgfpathlineto{\pgfqpoint{0.891787in}{0.849869in}}%
\pgfpathlineto{\pgfqpoint{0.897098in}{0.852438in}}%
\pgfpathlineto{\pgfqpoint{0.913029in}{0.852438in}}%
\pgfpathlineto{\pgfqpoint{0.918340in}{0.855008in}}%
\pgfpathlineto{\pgfqpoint{0.944892in}{0.855008in}}%
\pgfpathlineto{\pgfqpoint{0.950203in}{0.857434in}}%
\pgfpathlineto{\pgfqpoint{0.966134in}{0.857434in}}%
\pgfpathlineto{\pgfqpoint{0.971445in}{0.860004in}}%
\pgfpathlineto{\pgfqpoint{0.987376in}{0.860004in}}%
\pgfpathlineto{\pgfqpoint{0.992687in}{0.862430in}}%
\pgfpathlineto{\pgfqpoint{1.008618in}{0.862430in}}%
\pgfpathlineto{\pgfqpoint{1.013929in}{0.865000in}}%
\pgfpathlineto{\pgfqpoint{1.035171in}{0.865000in}}%
\pgfpathlineto{\pgfqpoint{1.040481in}{0.867569in}}%
\pgfpathlineto{\pgfqpoint{1.056413in}{0.867569in}}%
\pgfpathlineto{\pgfqpoint{1.061723in}{0.869996in}}%
\pgfpathlineto{\pgfqpoint{1.072344in}{0.869996in}}%
\pgfpathlineto{\pgfqpoint{1.077655in}{0.872565in}}%
\pgfpathlineto{\pgfqpoint{1.093586in}{0.872565in}}%
\pgfpathlineto{\pgfqpoint{1.098897in}{0.874992in}}%
\pgfpathlineto{\pgfqpoint{1.109518in}{0.874992in}}%
\pgfpathlineto{\pgfqpoint{1.114828in}{0.877561in}}%
\pgfpathlineto{\pgfqpoint{1.130760in}{0.877561in}}%
\pgfpathlineto{\pgfqpoint{1.136070in}{0.880131in}}%
\pgfpathlineto{\pgfqpoint{1.146691in}{0.880131in}}%
\pgfpathlineto{\pgfqpoint{1.152002in}{0.882557in}}%
\pgfpathlineto{\pgfqpoint{1.167933in}{0.882557in}}%
\pgfpathlineto{\pgfqpoint{1.173244in}{0.885127in}}%
\pgfpathlineto{\pgfqpoint{1.183865in}{0.885127in}}%
\pgfpathlineto{\pgfqpoint{1.189175in}{0.887553in}}%
\pgfpathlineto{\pgfqpoint{1.205107in}{0.887553in}}%
\pgfpathlineto{\pgfqpoint{1.210417in}{0.890123in}}%
\pgfpathlineto{\pgfqpoint{1.221038in}{0.890123in}}%
\pgfpathlineto{\pgfqpoint{1.226349in}{0.892692in}}%
\pgfpathlineto{\pgfqpoint{1.236970in}{0.892692in}}%
\pgfpathlineto{\pgfqpoint{1.242280in}{0.895119in}}%
\pgfpathlineto{\pgfqpoint{1.252901in}{0.895119in}}%
\pgfpathlineto{\pgfqpoint{1.258212in}{0.897688in}}%
\pgfpathlineto{\pgfqpoint{1.268833in}{0.897688in}}%
\pgfpathlineto{\pgfqpoint{1.274143in}{0.900115in}}%
\pgfpathlineto{\pgfqpoint{1.284764in}{0.900115in}}%
\pgfpathlineto{\pgfqpoint{1.290075in}{0.902684in}}%
\pgfpathlineto{\pgfqpoint{1.300696in}{0.902684in}}%
\pgfpathlineto{\pgfqpoint{1.306006in}{0.905254in}}%
\pgfpathlineto{\pgfqpoint{1.316627in}{0.905254in}}%
\pgfpathlineto{\pgfqpoint{1.321938in}{0.907680in}}%
\pgfpathlineto{\pgfqpoint{1.332559in}{0.907680in}}%
\pgfpathlineto{\pgfqpoint{1.337869in}{0.910250in}}%
\pgfpathlineto{\pgfqpoint{1.343180in}{0.910250in}}%
\pgfpathlineto{\pgfqpoint{1.348490in}{0.912676in}}%
\pgfpathlineto{\pgfqpoint{1.359111in}{0.912676in}}%
\pgfpathlineto{\pgfqpoint{1.364422in}{0.915246in}}%
\pgfpathlineto{\pgfqpoint{1.375043in}{0.915246in}}%
\pgfpathlineto{\pgfqpoint{1.380353in}{0.917815in}}%
\pgfpathlineto{\pgfqpoint{1.390974in}{0.917815in}}%
\pgfpathlineto{\pgfqpoint{1.396285in}{0.920242in}}%
\pgfpathlineto{\pgfqpoint{1.401595in}{0.920242in}}%
\pgfpathlineto{\pgfqpoint{1.406906in}{0.922811in}}%
\pgfpathlineto{\pgfqpoint{1.417527in}{0.922811in}}%
\pgfpathlineto{\pgfqpoint{1.422838in}{0.925238in}}%
\pgfpathlineto{\pgfqpoint{1.428148in}{0.925238in}}%
\pgfpathlineto{\pgfqpoint{1.433459in}{0.927807in}}%
\pgfpathlineto{\pgfqpoint{1.444080in}{0.927807in}}%
\pgfpathlineto{\pgfqpoint{1.449390in}{0.930234in}}%
\pgfpathlineto{\pgfqpoint{1.454701in}{0.930234in}}%
\pgfpathlineto{\pgfqpoint{1.460011in}{0.932803in}}%
\pgfpathlineto{\pgfqpoint{1.465322in}{0.932803in}}%
\pgfpathlineto{\pgfqpoint{1.470632in}{0.935372in}}%
\pgfpathlineto{\pgfqpoint{1.481253in}{0.935372in}}%
\pgfpathlineto{\pgfqpoint{1.486564in}{0.937799in}}%
\pgfpathlineto{\pgfqpoint{1.491874in}{0.937799in}}%
\pgfpathlineto{\pgfqpoint{1.497185in}{0.940369in}}%
\pgfpathlineto{\pgfqpoint{1.507806in}{0.940369in}}%
\pgfpathlineto{\pgfqpoint{1.513116in}{0.942795in}}%
\pgfpathlineto{\pgfqpoint{1.518427in}{0.942795in}}%
\pgfpathlineto{\pgfqpoint{1.523737in}{0.945365in}}%
\pgfpathlineto{\pgfqpoint{1.529048in}{0.945365in}}%
\pgfpathlineto{\pgfqpoint{1.534358in}{0.947934in}}%
\pgfpathlineto{\pgfqpoint{1.539669in}{0.947934in}}%
\pgfpathlineto{\pgfqpoint{1.544979in}{0.950361in}}%
\pgfpathlineto{\pgfqpoint{1.550290in}{0.950361in}}%
\pgfpathlineto{\pgfqpoint{1.555600in}{0.952930in}}%
\pgfpathlineto{\pgfqpoint{1.560911in}{0.952930in}}%
\pgfpathlineto{\pgfqpoint{1.566221in}{0.955357in}}%
\pgfpathlineto{\pgfqpoint{1.576842in}{0.955357in}}%
\pgfpathlineto{\pgfqpoint{1.582153in}{0.957926in}}%
\pgfpathlineto{\pgfqpoint{1.587463in}{0.957926in}}%
\pgfpathlineto{\pgfqpoint{1.592774in}{0.960495in}}%
\pgfpathlineto{\pgfqpoint{1.598084in}{0.960495in}}%
\pgfpathlineto{\pgfqpoint{1.603395in}{0.962922in}}%
\pgfpathlineto{\pgfqpoint{1.608705in}{0.962922in}}%
\pgfpathlineto{\pgfqpoint{1.614016in}{0.965491in}}%
\pgfpathlineto{\pgfqpoint{1.619326in}{0.965491in}}%
\pgfpathlineto{\pgfqpoint{1.629947in}{0.970488in}}%
\pgfpathlineto{\pgfqpoint{1.635258in}{0.970488in}}%
\pgfpathlineto{\pgfqpoint{1.640568in}{0.973057in}}%
\pgfpathlineto{\pgfqpoint{1.645879in}{0.973057in}}%
\pgfpathlineto{\pgfqpoint{1.651189in}{0.975484in}}%
\pgfpathlineto{\pgfqpoint{1.656500in}{0.975484in}}%
\pgfpathlineto{\pgfqpoint{1.661810in}{0.978053in}}%
\pgfpathlineto{\pgfqpoint{1.667121in}{0.978053in}}%
\pgfpathlineto{\pgfqpoint{1.672431in}{0.980480in}}%
\pgfpathlineto{\pgfqpoint{1.677742in}{0.980480in}}%
\pgfpathlineto{\pgfqpoint{1.683052in}{0.983049in}}%
\pgfpathlineto{\pgfqpoint{1.688363in}{0.983049in}}%
\pgfpathlineto{\pgfqpoint{1.698984in}{0.988045in}}%
\pgfpathlineto{\pgfqpoint{1.704294in}{0.988045in}}%
\pgfpathlineto{\pgfqpoint{1.709605in}{0.990614in}}%
\pgfpathlineto{\pgfqpoint{1.714915in}{0.990614in}}%
\pgfpathlineto{\pgfqpoint{1.725536in}{0.995610in}}%
\pgfpathlineto{\pgfqpoint{1.730847in}{0.995610in}}%
\pgfpathlineto{\pgfqpoint{1.741468in}{1.000606in}}%
\pgfpathlineto{\pgfqpoint{1.746778in}{1.000606in}}%
\pgfpathlineto{\pgfqpoint{1.752089in}{1.003176in}}%
\pgfpathlineto{\pgfqpoint{1.757399in}{1.003176in}}%
\pgfpathlineto{\pgfqpoint{1.768020in}{1.008172in}}%
\pgfpathlineto{\pgfqpoint{1.773331in}{1.008172in}}%
\pgfpathlineto{\pgfqpoint{1.783952in}{1.013168in}}%
\pgfpathlineto{\pgfqpoint{1.789262in}{1.013168in}}%
\pgfpathlineto{\pgfqpoint{1.799883in}{1.018164in}}%
\pgfpathlineto{\pgfqpoint{1.805194in}{1.018164in}}%
\pgfpathlineto{\pgfqpoint{1.821125in}{1.025729in}}%
\pgfpathlineto{\pgfqpoint{1.826436in}{1.025729in}}%
\pgfpathlineto{\pgfqpoint{1.842367in}{1.033295in}}%
\pgfpathlineto{\pgfqpoint{1.847678in}{1.033295in}}%
\pgfpathlineto{\pgfqpoint{1.863609in}{1.040860in}}%
\pgfpathlineto{\pgfqpoint{1.868920in}{1.040860in}}%
\pgfpathlineto{\pgfqpoint{1.890162in}{1.050852in}}%
\pgfpathlineto{\pgfqpoint{1.895472in}{1.050852in}}%
\pgfpathlineto{\pgfqpoint{1.922025in}{1.063414in}}%
\pgfpathlineto{\pgfqpoint{1.927335in}{1.063414in}}%
\pgfpathlineto{\pgfqpoint{1.969820in}{1.083541in}}%
\pgfpathlineto{\pgfqpoint{1.975130in}{1.083541in}}%
\pgfpathlineto{\pgfqpoint{2.022925in}{1.106094in}}%
\pgfpathlineto{\pgfqpoint{2.028235in}{1.111090in}}%
\pgfpathlineto{\pgfqpoint{2.065409in}{1.128791in}}%
\pgfpathlineto{\pgfqpoint{2.070719in}{1.133787in}}%
\pgfpathlineto{\pgfqpoint{2.097272in}{1.146348in}}%
\pgfpathlineto{\pgfqpoint{2.102582in}{1.151344in}}%
\pgfpathlineto{\pgfqpoint{2.113203in}{1.156340in}}%
\pgfpathlineto{\pgfqpoint{2.118514in}{1.161336in}}%
\pgfpathlineto{\pgfqpoint{2.129135in}{1.166332in}}%
\pgfpathlineto{\pgfqpoint{2.134445in}{1.171471in}}%
\pgfpathlineto{\pgfqpoint{2.145066in}{1.176467in}}%
\pgfpathlineto{\pgfqpoint{2.150377in}{1.181463in}}%
\pgfpathlineto{\pgfqpoint{2.160998in}{1.186459in}}%
\pgfpathlineto{\pgfqpoint{2.166308in}{1.191455in}}%
\pgfpathlineto{\pgfqpoint{2.171619in}{1.194024in}}%
\pgfpathlineto{\pgfqpoint{2.176929in}{1.199021in}}%
\pgfpathlineto{\pgfqpoint{2.182240in}{1.201590in}}%
\pgfpathlineto{\pgfqpoint{2.187550in}{1.206586in}}%
\pgfpathlineto{\pgfqpoint{2.192861in}{1.209155in}}%
\pgfpathlineto{\pgfqpoint{2.198171in}{1.214151in}}%
\pgfpathlineto{\pgfqpoint{2.203482in}{1.216578in}}%
\pgfpathlineto{\pgfqpoint{2.214103in}{1.226713in}}%
\pgfpathlineto{\pgfqpoint{2.219413in}{1.229139in}}%
\pgfpathlineto{\pgfqpoint{2.224724in}{1.234278in}}%
\pgfpathlineto{\pgfqpoint{2.230034in}{1.236705in}}%
\pgfpathlineto{\pgfqpoint{2.245966in}{1.251836in}}%
\pgfpathlineto{\pgfqpoint{2.251276in}{1.254262in}}%
\pgfpathlineto{\pgfqpoint{2.325623in}{1.324635in}}%
\pgfpathlineto{\pgfqpoint{2.330934in}{1.332201in}}%
\pgfpathlineto{\pgfqpoint{2.346865in}{1.347189in}}%
\pgfpathlineto{\pgfqpoint{2.352176in}{1.354754in}}%
\pgfpathlineto{\pgfqpoint{2.357486in}{1.359750in}}%
\pgfpathlineto{\pgfqpoint{2.362797in}{1.367316in}}%
\pgfpathlineto{\pgfqpoint{2.368107in}{1.372312in}}%
\pgfpathlineto{\pgfqpoint{2.373418in}{1.379877in}}%
\pgfpathlineto{\pgfqpoint{2.378728in}{1.384873in}}%
\pgfpathlineto{\pgfqpoint{2.384039in}{1.392439in}}%
\pgfpathlineto{\pgfqpoint{2.389349in}{1.397435in}}%
\pgfpathlineto{\pgfqpoint{2.405281in}{1.420131in}}%
\pgfpathlineto{\pgfqpoint{2.410591in}{1.425127in}}%
\pgfpathlineto{\pgfqpoint{2.447765in}{1.477799in}}%
\pgfpathlineto{\pgfqpoint{2.453075in}{1.487934in}}%
\pgfpathlineto{\pgfqpoint{2.463696in}{1.502922in}}%
\pgfpathlineto{\pgfqpoint{2.469007in}{1.513057in}}%
\pgfpathlineto{\pgfqpoint{2.474317in}{1.520480in}}%
\pgfpathlineto{\pgfqpoint{2.479628in}{1.530615in}}%
\pgfpathlineto{\pgfqpoint{2.484938in}{1.538180in}}%
\pgfpathlineto{\pgfqpoint{2.495559in}{1.558164in}}%
\pgfpathlineto{\pgfqpoint{2.500870in}{1.565730in}}%
\pgfpathlineto{\pgfqpoint{2.532733in}{1.625968in}}%
\pgfpathlineto{\pgfqpoint{2.548665in}{1.661225in}}%
\pgfpathlineto{\pgfqpoint{2.553975in}{1.671217in}}%
\pgfpathlineto{\pgfqpoint{2.607080in}{1.799259in}}%
\pgfpathlineto{\pgfqpoint{2.644254in}{1.907316in}}%
\pgfpathlineto{\pgfqpoint{2.686738in}{2.050488in}}%
\pgfpathlineto{\pgfqpoint{2.707980in}{2.133422in}}%
\pgfpathlineto{\pgfqpoint{2.745153in}{2.294152in}}%
\pgfpathlineto{\pgfqpoint{2.771706in}{2.422193in}}%
\pgfpathlineto{\pgfqpoint{2.830121in}{2.756214in}}%
\pgfpathlineto{\pgfqpoint{2.888537in}{3.150615in}}%
\pgfpathlineto{\pgfqpoint{2.899158in}{3.223415in}}%
\pgfpathlineto{\pgfqpoint{2.915089in}{3.223415in}}%
\pgfpathlineto{\pgfqpoint{2.920400in}{3.195722in}}%
\pgfpathlineto{\pgfqpoint{2.984126in}{2.751218in}}%
\pgfpathlineto{\pgfqpoint{3.010678in}{2.582923in}}%
\pgfpathlineto{\pgfqpoint{3.031920in}{2.457308in}}%
\pgfpathlineto{\pgfqpoint{3.058473in}{2.311709in}}%
\pgfpathlineto{\pgfqpoint{3.079715in}{2.206222in}}%
\pgfpathlineto{\pgfqpoint{3.106268in}{2.085603in}}%
\pgfpathlineto{\pgfqpoint{3.138131in}{1.954992in}}%
\pgfpathlineto{\pgfqpoint{3.159373in}{1.877197in}}%
\pgfpathlineto{\pgfqpoint{3.196546in}{1.756578in}}%
\pgfpathlineto{\pgfqpoint{3.239030in}{1.638529in}}%
\pgfpathlineto{\pgfqpoint{3.292135in}{1.515484in}}%
\pgfpathlineto{\pgfqpoint{3.302756in}{1.495357in}}%
\pgfpathlineto{\pgfqpoint{3.313377in}{1.472803in}}%
\pgfpathlineto{\pgfqpoint{3.350551in}{1.405000in}}%
\pgfpathlineto{\pgfqpoint{3.355861in}{1.397435in}}%
\pgfpathlineto{\pgfqpoint{3.361172in}{1.387442in}}%
\pgfpathlineto{\pgfqpoint{3.371793in}{1.372312in}}%
\pgfpathlineto{\pgfqpoint{3.377103in}{1.362320in}}%
\pgfpathlineto{\pgfqpoint{3.387724in}{1.347189in}}%
\pgfpathlineto{\pgfqpoint{3.393035in}{1.337197in}}%
\pgfpathlineto{\pgfqpoint{3.424898in}{1.291947in}}%
\pgfpathlineto{\pgfqpoint{3.430208in}{1.286951in}}%
\pgfpathlineto{\pgfqpoint{3.446140in}{1.264397in}}%
\pgfpathlineto{\pgfqpoint{3.451450in}{1.259258in}}%
\pgfpathlineto{\pgfqpoint{3.456761in}{1.251836in}}%
\pgfpathlineto{\pgfqpoint{3.462071in}{1.246697in}}%
\pgfpathlineto{\pgfqpoint{3.467382in}{1.239274in}}%
\pgfpathlineto{\pgfqpoint{3.472692in}{1.234278in}}%
\pgfpathlineto{\pgfqpoint{3.478003in}{1.226713in}}%
\pgfpathlineto{\pgfqpoint{3.488624in}{1.216578in}}%
\pgfpathlineto{\pgfqpoint{3.493934in}{1.209155in}}%
\pgfpathlineto{\pgfqpoint{3.515176in}{1.189028in}}%
\pgfpathlineto{\pgfqpoint{3.520487in}{1.181463in}}%
\pgfpathlineto{\pgfqpoint{3.557660in}{1.146348in}}%
\pgfpathlineto{\pgfqpoint{3.562971in}{1.143779in}}%
\pgfpathlineto{\pgfqpoint{3.594834in}{1.113660in}}%
\pgfpathlineto{\pgfqpoint{3.600144in}{1.111090in}}%
\pgfpathlineto{\pgfqpoint{3.610765in}{1.101098in}}%
\pgfpathlineto{\pgfqpoint{3.616076in}{1.098529in}}%
\pgfpathlineto{\pgfqpoint{3.621386in}{1.093533in}}%
\pgfpathlineto{\pgfqpoint{3.626697in}{1.091106in}}%
\pgfpathlineto{\pgfqpoint{3.632007in}{1.085967in}}%
\pgfpathlineto{\pgfqpoint{3.637318in}{1.083541in}}%
\pgfpathlineto{\pgfqpoint{3.647939in}{1.073406in}}%
\pgfpathlineto{\pgfqpoint{3.658560in}{1.068410in}}%
\pgfpathlineto{\pgfqpoint{3.663871in}{1.063414in}}%
\pgfpathlineto{\pgfqpoint{3.669181in}{1.060844in}}%
\pgfpathlineto{\pgfqpoint{3.674492in}{1.055848in}}%
\pgfpathlineto{\pgfqpoint{3.685113in}{1.050852in}}%
\pgfpathlineto{\pgfqpoint{3.690423in}{1.045856in}}%
\pgfpathlineto{\pgfqpoint{3.695734in}{1.043287in}}%
\pgfpathlineto{\pgfqpoint{3.701044in}{1.038291in}}%
\pgfpathlineto{\pgfqpoint{3.716976in}{1.030725in}}%
\pgfpathlineto{\pgfqpoint{3.722286in}{1.025729in}}%
\pgfpathlineto{\pgfqpoint{3.738218in}{1.018164in}}%
\pgfpathlineto{\pgfqpoint{3.743528in}{1.013168in}}%
\pgfpathlineto{\pgfqpoint{3.775391in}{0.998180in}}%
\pgfpathlineto{\pgfqpoint{3.780702in}{0.993041in}}%
\pgfpathlineto{\pgfqpoint{3.855049in}{0.957926in}}%
\pgfpathlineto{\pgfqpoint{3.870980in}{0.950361in}}%
\pgfpathlineto{\pgfqpoint{3.876291in}{0.950361in}}%
\pgfpathlineto{\pgfqpoint{3.913464in}{0.932803in}}%
\pgfpathlineto{\pgfqpoint{3.918775in}{0.932803in}}%
\pgfpathlineto{\pgfqpoint{3.940017in}{0.922811in}}%
\pgfpathlineto{\pgfqpoint{3.945327in}{0.922811in}}%
\pgfpathlineto{\pgfqpoint{3.966569in}{0.912676in}}%
\pgfpathlineto{\pgfqpoint{3.971880in}{0.912676in}}%
\pgfpathlineto{\pgfqpoint{3.987811in}{0.905254in}}%
\pgfpathlineto{\pgfqpoint{3.993122in}{0.905254in}}%
\pgfpathlineto{\pgfqpoint{4.009053in}{0.897688in}}%
\pgfpathlineto{\pgfqpoint{4.014364in}{0.897688in}}%
\pgfpathlineto{\pgfqpoint{4.024985in}{0.892692in}}%
\pgfpathlineto{\pgfqpoint{4.030295in}{0.892692in}}%
\pgfpathlineto{\pgfqpoint{4.040916in}{0.887553in}}%
\pgfpathlineto{\pgfqpoint{4.046227in}{0.887553in}}%
\pgfpathlineto{\pgfqpoint{4.062158in}{0.880131in}}%
\pgfpathlineto{\pgfqpoint{4.067469in}{0.880131in}}%
\pgfpathlineto{\pgfqpoint{4.072779in}{0.877561in}}%
\pgfpathlineto{\pgfqpoint{4.078090in}{0.877561in}}%
\pgfpathlineto{\pgfqpoint{4.088711in}{0.872565in}}%
\pgfpathlineto{\pgfqpoint{4.094021in}{0.872565in}}%
\pgfpathlineto{\pgfqpoint{4.104642in}{0.867569in}}%
\pgfpathlineto{\pgfqpoint{4.109953in}{0.867569in}}%
\pgfpathlineto{\pgfqpoint{4.115263in}{0.865000in}}%
\pgfpathlineto{\pgfqpoint{4.120574in}{0.865000in}}%
\pgfpathlineto{\pgfqpoint{4.131195in}{0.860004in}}%
\pgfpathlineto{\pgfqpoint{4.136505in}{0.860004in}}%
\pgfpathlineto{\pgfqpoint{4.141816in}{0.857434in}}%
\pgfpathlineto{\pgfqpoint{4.147126in}{0.857434in}}%
\pgfpathlineto{\pgfqpoint{4.152437in}{0.855008in}}%
\pgfpathlineto{\pgfqpoint{4.157747in}{0.855008in}}%
\pgfpathlineto{\pgfqpoint{4.163058in}{0.852438in}}%
\pgfpathlineto{\pgfqpoint{4.168368in}{0.852438in}}%
\pgfpathlineto{\pgfqpoint{4.173679in}{0.849869in}}%
\pgfpathlineto{\pgfqpoint{4.178989in}{0.849869in}}%
\pgfpathlineto{\pgfqpoint{4.189610in}{0.844873in}}%
\pgfpathlineto{\pgfqpoint{4.194921in}{0.844873in}}%
\pgfpathlineto{\pgfqpoint{4.200232in}{0.842446in}}%
\pgfpathlineto{\pgfqpoint{4.205542in}{0.842446in}}%
\pgfpathlineto{\pgfqpoint{4.210853in}{0.839877in}}%
\pgfpathlineto{\pgfqpoint{4.216163in}{0.839877in}}%
\pgfpathlineto{\pgfqpoint{4.221474in}{0.837307in}}%
\pgfpathlineto{\pgfqpoint{4.226784in}{0.837307in}}%
\pgfpathlineto{\pgfqpoint{4.232095in}{0.834881in}}%
\pgfpathlineto{\pgfqpoint{4.237405in}{0.834881in}}%
\pgfpathlineto{\pgfqpoint{4.242716in}{0.832311in}}%
\pgfpathlineto{\pgfqpoint{4.248026in}{0.832311in}}%
\pgfpathlineto{\pgfqpoint{4.253337in}{0.829885in}}%
\pgfpathlineto{\pgfqpoint{4.258647in}{0.829885in}}%
\pgfpathlineto{\pgfqpoint{4.263958in}{0.827315in}}%
\pgfpathlineto{\pgfqpoint{4.269268in}{0.827315in}}%
\pgfpathlineto{\pgfqpoint{4.274579in}{0.824746in}}%
\pgfpathlineto{\pgfqpoint{4.285200in}{0.824746in}}%
\pgfpathlineto{\pgfqpoint{4.290510in}{0.822319in}}%
\pgfpathlineto{\pgfqpoint{4.295821in}{0.822319in}}%
\pgfpathlineto{\pgfqpoint{4.301131in}{0.819750in}}%
\pgfpathlineto{\pgfqpoint{4.306442in}{0.819750in}}%
\pgfpathlineto{\pgfqpoint{4.311752in}{0.817323in}}%
\pgfpathlineto{\pgfqpoint{4.317063in}{0.817323in}}%
\pgfpathlineto{\pgfqpoint{4.322373in}{0.814754in}}%
\pgfpathlineto{\pgfqpoint{4.332994in}{0.814754in}}%
\pgfpathlineto{\pgfqpoint{4.338305in}{0.812327in}}%
\pgfpathlineto{\pgfqpoint{4.343615in}{0.812327in}}%
\pgfpathlineto{\pgfqpoint{4.348926in}{0.809758in}}%
\pgfpathlineto{\pgfqpoint{4.359547in}{0.809758in}}%
\pgfpathlineto{\pgfqpoint{4.364857in}{0.807188in}}%
\pgfpathlineto{\pgfqpoint{4.370168in}{0.807188in}}%
\pgfpathlineto{\pgfqpoint{4.375478in}{0.804762in}}%
\pgfpathlineto{\pgfqpoint{4.386099in}{0.804762in}}%
\pgfpathlineto{\pgfqpoint{4.391410in}{0.802192in}}%
\pgfpathlineto{\pgfqpoint{4.396720in}{0.802192in}}%
\pgfpathlineto{\pgfqpoint{4.402031in}{0.799766in}}%
\pgfpathlineto{\pgfqpoint{4.412652in}{0.799766in}}%
\pgfpathlineto{\pgfqpoint{4.417962in}{0.797196in}}%
\pgfpathlineto{\pgfqpoint{4.423273in}{0.797196in}}%
\pgfpathlineto{\pgfqpoint{4.428583in}{0.794627in}}%
\pgfpathlineto{\pgfqpoint{4.439204in}{0.794627in}}%
\pgfpathlineto{\pgfqpoint{4.444515in}{0.792200in}}%
\pgfpathlineto{\pgfqpoint{4.455136in}{0.792200in}}%
\pgfpathlineto{\pgfqpoint{4.460446in}{0.789631in}}%
\pgfpathlineto{\pgfqpoint{4.471067in}{0.789631in}}%
\pgfpathlineto{\pgfqpoint{4.476378in}{0.787204in}}%
\pgfpathlineto{\pgfqpoint{4.486999in}{0.787204in}}%
\pgfpathlineto{\pgfqpoint{4.492309in}{0.784635in}}%
\pgfpathlineto{\pgfqpoint{4.497620in}{0.784635in}}%
\pgfpathlineto{\pgfqpoint{4.502930in}{0.782066in}}%
\pgfpathlineto{\pgfqpoint{4.513551in}{0.782066in}}%
\pgfpathlineto{\pgfqpoint{4.518862in}{0.779639in}}%
\pgfpathlineto{\pgfqpoint{4.534793in}{0.779639in}}%
\pgfpathlineto{\pgfqpoint{4.540104in}{0.777069in}}%
\pgfpathlineto{\pgfqpoint{4.550725in}{0.777069in}}%
\pgfpathlineto{\pgfqpoint{4.556035in}{0.774643in}}%
\pgfpathlineto{\pgfqpoint{4.566656in}{0.774643in}}%
\pgfpathlineto{\pgfqpoint{4.571967in}{0.772073in}}%
\pgfpathlineto{\pgfqpoint{4.582588in}{0.772073in}}%
\pgfpathlineto{\pgfqpoint{4.587898in}{0.769504in}}%
\pgfpathlineto{\pgfqpoint{4.603830in}{0.769504in}}%
\pgfpathlineto{\pgfqpoint{4.609140in}{0.767077in}}%
\pgfpathlineto{\pgfqpoint{4.625072in}{0.767077in}}%
\pgfpathlineto{\pgfqpoint{4.630382in}{0.764508in}}%
\pgfpathlineto{\pgfqpoint{4.641003in}{0.764508in}}%
\pgfpathlineto{\pgfqpoint{4.646314in}{0.762081in}}%
\pgfpathlineto{\pgfqpoint{4.662245in}{0.762081in}}%
\pgfpathlineto{\pgfqpoint{4.667556in}{0.759512in}}%
\pgfpathlineto{\pgfqpoint{4.683487in}{0.759512in}}%
\pgfpathlineto{\pgfqpoint{4.688798in}{0.756943in}}%
\pgfpathlineto{\pgfqpoint{4.699419in}{0.756943in}}%
\pgfpathlineto{\pgfqpoint{4.704729in}{0.754516in}}%
\pgfpathlineto{\pgfqpoint{4.720661in}{0.754516in}}%
\pgfpathlineto{\pgfqpoint{4.725971in}{0.751947in}}%
\pgfpathlineto{\pgfqpoint{4.747213in}{0.751947in}}%
\pgfpathlineto{\pgfqpoint{4.752524in}{0.749520in}}%
\pgfpathlineto{\pgfqpoint{4.768456in}{0.749520in}}%
\pgfpathlineto{\pgfqpoint{4.773766in}{0.746951in}}%
\pgfpathlineto{\pgfqpoint{4.789698in}{0.746951in}}%
\pgfpathlineto{\pgfqpoint{4.795008in}{0.744381in}}%
\pgfpathlineto{\pgfqpoint{4.810940in}{0.744381in}}%
\pgfpathlineto{\pgfqpoint{4.816250in}{0.741954in}}%
\pgfpathlineto{\pgfqpoint{4.837492in}{0.741954in}}%
\pgfpathlineto{\pgfqpoint{4.842803in}{0.739385in}}%
\pgfpathlineto{\pgfqpoint{4.864045in}{0.739385in}}%
\pgfpathlineto{\pgfqpoint{4.869355in}{0.736958in}}%
\pgfpathlineto{\pgfqpoint{4.890597in}{0.736958in}}%
\pgfpathlineto{\pgfqpoint{4.895908in}{0.734389in}}%
\pgfpathlineto{\pgfqpoint{4.917150in}{0.734389in}}%
\pgfpathlineto{\pgfqpoint{4.922460in}{0.731820in}}%
\pgfpathlineto{\pgfqpoint{4.943702in}{0.731820in}}%
\pgfpathlineto{\pgfqpoint{4.949013in}{0.729393in}}%
\pgfpathlineto{\pgfqpoint{4.970255in}{0.729393in}}%
\pgfpathlineto{\pgfqpoint{4.975565in}{0.726824in}}%
\pgfpathlineto{\pgfqpoint{5.002118in}{0.726824in}}%
\pgfpathlineto{\pgfqpoint{5.007428in}{0.724397in}}%
\pgfpathlineto{\pgfqpoint{5.033981in}{0.724397in}}%
\pgfpathlineto{\pgfqpoint{5.039291in}{0.721828in}}%
\pgfpathlineto{\pgfqpoint{5.065844in}{0.721828in}}%
\pgfpathlineto{\pgfqpoint{5.071154in}{0.719258in}}%
\pgfpathlineto{\pgfqpoint{5.097707in}{0.719258in}}%
\pgfpathlineto{\pgfqpoint{5.103017in}{0.716832in}}%
\pgfpathlineto{\pgfqpoint{5.129570in}{0.716832in}}%
\pgfpathlineto{\pgfqpoint{5.134880in}{0.714262in}}%
\pgfpathlineto{\pgfqpoint{5.172054in}{0.714262in}}%
\pgfpathlineto{\pgfqpoint{5.177364in}{0.711836in}}%
\pgfpathlineto{\pgfqpoint{5.209227in}{0.711836in}}%
\pgfpathlineto{\pgfqpoint{5.214538in}{0.709266in}}%
\pgfpathlineto{\pgfqpoint{5.257022in}{0.709266in}}%
\pgfpathlineto{\pgfqpoint{5.262332in}{0.706839in}}%
\pgfpathlineto{\pgfqpoint{5.299506in}{0.706839in}}%
\pgfpathlineto{\pgfqpoint{5.304816in}{0.704270in}}%
\pgfpathlineto{\pgfqpoint{5.347301in}{0.704270in}}%
\pgfpathlineto{\pgfqpoint{5.352611in}{0.701701in}}%
\pgfpathlineto{\pgfqpoint{5.395095in}{0.701701in}}%
\pgfpathlineto{\pgfqpoint{5.400406in}{0.699274in}}%
\pgfpathlineto{\pgfqpoint{5.453511in}{0.699274in}}%
\pgfpathlineto{\pgfqpoint{5.458821in}{0.696705in}}%
\pgfpathlineto{\pgfqpoint{5.511926in}{0.696705in}}%
\pgfpathlineto{\pgfqpoint{5.517237in}{0.694278in}}%
\pgfpathlineto{\pgfqpoint{5.580963in}{0.694278in}}%
\pgfpathlineto{\pgfqpoint{5.586273in}{0.691709in}}%
\pgfpathlineto{\pgfqpoint{5.649999in}{0.691709in}}%
\pgfpathlineto{\pgfqpoint{5.649999in}{0.691709in}}%
\pgfusepath{stroke}%
\end{pgfscope}%
\begin{pgfscope}%
\pgfsetrectcap%
\pgfsetmiterjoin%
\pgfsetlinewidth{0.803000pt}%
\definecolor{currentstroke}{rgb}{0.000000,0.000000,0.000000}%
\pgfsetstrokecolor{currentstroke}%
\pgfsetdash{}{0pt}%
\pgfpathmoveto{\pgfqpoint{0.727161in}{0.565123in}}%
\pgfpathlineto{\pgfqpoint{0.727161in}{3.350000in}}%
\pgfusepath{stroke}%
\end{pgfscope}%
\begin{pgfscope}%
\pgfsetrectcap%
\pgfsetmiterjoin%
\pgfsetlinewidth{0.803000pt}%
\definecolor{currentstroke}{rgb}{0.000000,0.000000,0.000000}%
\pgfsetstrokecolor{currentstroke}%
\pgfsetdash{}{0pt}%
\pgfpathmoveto{\pgfqpoint{5.649999in}{0.565123in}}%
\pgfpathlineto{\pgfqpoint{5.649999in}{3.350000in}}%
\pgfusepath{stroke}%
\end{pgfscope}%
\begin{pgfscope}%
\pgfsetrectcap%
\pgfsetmiterjoin%
\pgfsetlinewidth{0.803000pt}%
\definecolor{currentstroke}{rgb}{0.000000,0.000000,0.000000}%
\pgfsetstrokecolor{currentstroke}%
\pgfsetdash{}{0pt}%
\pgfpathmoveto{\pgfqpoint{0.727161in}{0.565123in}}%
\pgfpathlineto{\pgfqpoint{5.649999in}{0.565123in}}%
\pgfusepath{stroke}%
\end{pgfscope}%
\begin{pgfscope}%
\pgfsetrectcap%
\pgfsetmiterjoin%
\pgfsetlinewidth{0.803000pt}%
\definecolor{currentstroke}{rgb}{0.000000,0.000000,0.000000}%
\pgfsetstrokecolor{currentstroke}%
\pgfsetdash{}{0pt}%
\pgfpathmoveto{\pgfqpoint{0.727161in}{3.350000in}}%
\pgfpathlineto{\pgfqpoint{5.649999in}{3.350000in}}%
\pgfusepath{stroke}%
\end{pgfscope}%
\end{pgfpicture}%
\makeatother%
\endgroup%

		\caption{Measured return loss and phase of a quartz resonatoras a function of frequency.}
		\label{fig:353}
	\end{figure}
	
	\subsubsection{}
	\begin{figure}[H]
		\centering
		%% Creator: Matplotlib, PGF backend
%%
%% To include the figure in your LaTeX document, write
%%   \input{<filename>.pgf}
%%
%% Make sure the required packages are loaded in your preamble
%%   \usepackage{pgf}
%%
%% Also ensure that all the required font packages are loaded; for instance,
%% the lmodern package is sometimes necessary when using math font.
%%   \usepackage{lmodern}
%%
%% Figures using additional raster images can only be included by \input if
%% they are in the same directory as the main LaTeX file. For loading figures
%% from other directories you can use the `import` package
%%   \usepackage{import}
%%
%% and then include the figures with
%%   \import{<path to file>}{<filename>.pgf}
%%
%% Matplotlib used the following preamble
%%
\begingroup%
\makeatletter%
\begin{pgfpicture}%
\pgfpathrectangle{\pgfpointorigin}{\pgfqpoint{6.300000in}{3.500000in}}%
\pgfusepath{use as bounding box, clip}%
\begin{pgfscope}%
\pgfsetbuttcap%
\pgfsetmiterjoin%
\definecolor{currentfill}{rgb}{1.000000,1.000000,1.000000}%
\pgfsetfillcolor{currentfill}%
\pgfsetlinewidth{0.000000pt}%
\definecolor{currentstroke}{rgb}{1.000000,1.000000,1.000000}%
\pgfsetstrokecolor{currentstroke}%
\pgfsetdash{}{0pt}%
\pgfpathmoveto{\pgfqpoint{0.000000in}{0.000000in}}%
\pgfpathlineto{\pgfqpoint{6.300000in}{0.000000in}}%
\pgfpathlineto{\pgfqpoint{6.300000in}{3.500000in}}%
\pgfpathlineto{\pgfqpoint{0.000000in}{3.500000in}}%
\pgfpathlineto{\pgfqpoint{0.000000in}{0.000000in}}%
\pgfpathclose%
\pgfusepath{fill}%
\end{pgfscope}%
\begin{pgfscope}%
\pgfsetbuttcap%
\pgfsetmiterjoin%
\definecolor{currentfill}{rgb}{1.000000,1.000000,1.000000}%
\pgfsetfillcolor{currentfill}%
\pgfsetlinewidth{0.000000pt}%
\definecolor{currentstroke}{rgb}{0.000000,0.000000,0.000000}%
\pgfsetstrokecolor{currentstroke}%
\pgfsetstrokeopacity{0.000000}%
\pgfsetdash{}{0pt}%
\pgfpathmoveto{\pgfqpoint{0.719445in}{0.565123in}}%
\pgfpathlineto{\pgfqpoint{6.150000in}{0.565123in}}%
\pgfpathlineto{\pgfqpoint{6.150000in}{3.150926in}}%
\pgfpathlineto{\pgfqpoint{0.719445in}{3.150926in}}%
\pgfpathlineto{\pgfqpoint{0.719445in}{0.565123in}}%
\pgfpathclose%
\pgfusepath{fill}%
\end{pgfscope}%
\begin{pgfscope}%
\pgfpathrectangle{\pgfqpoint{0.719445in}{0.565123in}}{\pgfqpoint{5.430555in}{2.585803in}}%
\pgfusepath{clip}%
\pgfsetbuttcap%
\pgfsetroundjoin%
\pgfsetlinewidth{0.803000pt}%
\definecolor{currentstroke}{rgb}{0.690196,0.690196,0.690196}%
\pgfsetstrokecolor{currentstroke}%
\pgfsetdash{{0.800000pt}{1.320000pt}}{0.000000pt}%
\pgfpathmoveto{\pgfqpoint{0.966289in}{0.565123in}}%
\pgfpathlineto{\pgfqpoint{0.966289in}{3.150926in}}%
\pgfusepath{stroke}%
\end{pgfscope}%
\begin{pgfscope}%
\pgfsetbuttcap%
\pgfsetroundjoin%
\definecolor{currentfill}{rgb}{0.000000,0.000000,0.000000}%
\pgfsetfillcolor{currentfill}%
\pgfsetlinewidth{0.803000pt}%
\definecolor{currentstroke}{rgb}{0.000000,0.000000,0.000000}%
\pgfsetstrokecolor{currentstroke}%
\pgfsetdash{}{0pt}%
\pgfsys@defobject{currentmarker}{\pgfqpoint{0.000000in}{-0.048611in}}{\pgfqpoint{0.000000in}{0.000000in}}{%
\pgfpathmoveto{\pgfqpoint{0.000000in}{0.000000in}}%
\pgfpathlineto{\pgfqpoint{0.000000in}{-0.048611in}}%
\pgfusepath{stroke,fill}%
}%
\begin{pgfscope}%
\pgfsys@transformshift{0.966289in}{0.565123in}%
\pgfsys@useobject{currentmarker}{}%
\end{pgfscope}%
\end{pgfscope}%
\begin{pgfscope}%
\definecolor{textcolor}{rgb}{0.000000,0.000000,0.000000}%
\pgfsetstrokecolor{textcolor}%
\pgfsetfillcolor{textcolor}%
\pgftext[x=0.966289in,y=0.467901in,,top]{\color{textcolor}\rmfamily\fontsize{10.000000}{12.000000}\selectfont \(\displaystyle {9.8270}\)}%
\end{pgfscope}%
\begin{pgfscope}%
\pgfpathrectangle{\pgfqpoint{0.719445in}{0.565123in}}{\pgfqpoint{5.430555in}{2.585803in}}%
\pgfusepath{clip}%
\pgfsetbuttcap%
\pgfsetroundjoin%
\pgfsetlinewidth{0.803000pt}%
\definecolor{currentstroke}{rgb}{0.690196,0.690196,0.690196}%
\pgfsetstrokecolor{currentstroke}%
\pgfsetdash{{0.800000pt}{1.320000pt}}{0.000000pt}%
\pgfpathmoveto{\pgfqpoint{1.853895in}{0.565123in}}%
\pgfpathlineto{\pgfqpoint{1.853895in}{3.150926in}}%
\pgfusepath{stroke}%
\end{pgfscope}%
\begin{pgfscope}%
\pgfsetbuttcap%
\pgfsetroundjoin%
\definecolor{currentfill}{rgb}{0.000000,0.000000,0.000000}%
\pgfsetfillcolor{currentfill}%
\pgfsetlinewidth{0.803000pt}%
\definecolor{currentstroke}{rgb}{0.000000,0.000000,0.000000}%
\pgfsetstrokecolor{currentstroke}%
\pgfsetdash{}{0pt}%
\pgfsys@defobject{currentmarker}{\pgfqpoint{0.000000in}{-0.048611in}}{\pgfqpoint{0.000000in}{0.000000in}}{%
\pgfpathmoveto{\pgfqpoint{0.000000in}{0.000000in}}%
\pgfpathlineto{\pgfqpoint{0.000000in}{-0.048611in}}%
\pgfusepath{stroke,fill}%
}%
\begin{pgfscope}%
\pgfsys@transformshift{1.853895in}{0.565123in}%
\pgfsys@useobject{currentmarker}{}%
\end{pgfscope}%
\end{pgfscope}%
\begin{pgfscope}%
\definecolor{textcolor}{rgb}{0.000000,0.000000,0.000000}%
\pgfsetstrokecolor{textcolor}%
\pgfsetfillcolor{textcolor}%
\pgftext[x=1.853895in,y=0.467901in,,top]{\color{textcolor}\rmfamily\fontsize{10.000000}{12.000000}\selectfont \(\displaystyle {9.8275}\)}%
\end{pgfscope}%
\begin{pgfscope}%
\pgfpathrectangle{\pgfqpoint{0.719445in}{0.565123in}}{\pgfqpoint{5.430555in}{2.585803in}}%
\pgfusepath{clip}%
\pgfsetbuttcap%
\pgfsetroundjoin%
\pgfsetlinewidth{0.803000pt}%
\definecolor{currentstroke}{rgb}{0.690196,0.690196,0.690196}%
\pgfsetstrokecolor{currentstroke}%
\pgfsetdash{{0.800000pt}{1.320000pt}}{0.000000pt}%
\pgfpathmoveto{\pgfqpoint{2.741502in}{0.565123in}}%
\pgfpathlineto{\pgfqpoint{2.741502in}{3.150926in}}%
\pgfusepath{stroke}%
\end{pgfscope}%
\begin{pgfscope}%
\pgfsetbuttcap%
\pgfsetroundjoin%
\definecolor{currentfill}{rgb}{0.000000,0.000000,0.000000}%
\pgfsetfillcolor{currentfill}%
\pgfsetlinewidth{0.803000pt}%
\definecolor{currentstroke}{rgb}{0.000000,0.000000,0.000000}%
\pgfsetstrokecolor{currentstroke}%
\pgfsetdash{}{0pt}%
\pgfsys@defobject{currentmarker}{\pgfqpoint{0.000000in}{-0.048611in}}{\pgfqpoint{0.000000in}{0.000000in}}{%
\pgfpathmoveto{\pgfqpoint{0.000000in}{0.000000in}}%
\pgfpathlineto{\pgfqpoint{0.000000in}{-0.048611in}}%
\pgfusepath{stroke,fill}%
}%
\begin{pgfscope}%
\pgfsys@transformshift{2.741502in}{0.565123in}%
\pgfsys@useobject{currentmarker}{}%
\end{pgfscope}%
\end{pgfscope}%
\begin{pgfscope}%
\definecolor{textcolor}{rgb}{0.000000,0.000000,0.000000}%
\pgfsetstrokecolor{textcolor}%
\pgfsetfillcolor{textcolor}%
\pgftext[x=2.741502in,y=0.467901in,,top]{\color{textcolor}\rmfamily\fontsize{10.000000}{12.000000}\selectfont \(\displaystyle {9.8280}\)}%
\end{pgfscope}%
\begin{pgfscope}%
\pgfpathrectangle{\pgfqpoint{0.719445in}{0.565123in}}{\pgfqpoint{5.430555in}{2.585803in}}%
\pgfusepath{clip}%
\pgfsetbuttcap%
\pgfsetroundjoin%
\pgfsetlinewidth{0.803000pt}%
\definecolor{currentstroke}{rgb}{0.690196,0.690196,0.690196}%
\pgfsetstrokecolor{currentstroke}%
\pgfsetdash{{0.800000pt}{1.320000pt}}{0.000000pt}%
\pgfpathmoveto{\pgfqpoint{3.629108in}{0.565123in}}%
\pgfpathlineto{\pgfqpoint{3.629108in}{3.150926in}}%
\pgfusepath{stroke}%
\end{pgfscope}%
\begin{pgfscope}%
\pgfsetbuttcap%
\pgfsetroundjoin%
\definecolor{currentfill}{rgb}{0.000000,0.000000,0.000000}%
\pgfsetfillcolor{currentfill}%
\pgfsetlinewidth{0.803000pt}%
\definecolor{currentstroke}{rgb}{0.000000,0.000000,0.000000}%
\pgfsetstrokecolor{currentstroke}%
\pgfsetdash{}{0pt}%
\pgfsys@defobject{currentmarker}{\pgfqpoint{0.000000in}{-0.048611in}}{\pgfqpoint{0.000000in}{0.000000in}}{%
\pgfpathmoveto{\pgfqpoint{0.000000in}{0.000000in}}%
\pgfpathlineto{\pgfqpoint{0.000000in}{-0.048611in}}%
\pgfusepath{stroke,fill}%
}%
\begin{pgfscope}%
\pgfsys@transformshift{3.629108in}{0.565123in}%
\pgfsys@useobject{currentmarker}{}%
\end{pgfscope}%
\end{pgfscope}%
\begin{pgfscope}%
\definecolor{textcolor}{rgb}{0.000000,0.000000,0.000000}%
\pgfsetstrokecolor{textcolor}%
\pgfsetfillcolor{textcolor}%
\pgftext[x=3.629108in,y=0.467901in,,top]{\color{textcolor}\rmfamily\fontsize{10.000000}{12.000000}\selectfont \(\displaystyle {9.8285}\)}%
\end{pgfscope}%
\begin{pgfscope}%
\pgfpathrectangle{\pgfqpoint{0.719445in}{0.565123in}}{\pgfqpoint{5.430555in}{2.585803in}}%
\pgfusepath{clip}%
\pgfsetbuttcap%
\pgfsetroundjoin%
\pgfsetlinewidth{0.803000pt}%
\definecolor{currentstroke}{rgb}{0.690196,0.690196,0.690196}%
\pgfsetstrokecolor{currentstroke}%
\pgfsetdash{{0.800000pt}{1.320000pt}}{0.000000pt}%
\pgfpathmoveto{\pgfqpoint{4.516715in}{0.565123in}}%
\pgfpathlineto{\pgfqpoint{4.516715in}{3.150926in}}%
\pgfusepath{stroke}%
\end{pgfscope}%
\begin{pgfscope}%
\pgfsetbuttcap%
\pgfsetroundjoin%
\definecolor{currentfill}{rgb}{0.000000,0.000000,0.000000}%
\pgfsetfillcolor{currentfill}%
\pgfsetlinewidth{0.803000pt}%
\definecolor{currentstroke}{rgb}{0.000000,0.000000,0.000000}%
\pgfsetstrokecolor{currentstroke}%
\pgfsetdash{}{0pt}%
\pgfsys@defobject{currentmarker}{\pgfqpoint{0.000000in}{-0.048611in}}{\pgfqpoint{0.000000in}{0.000000in}}{%
\pgfpathmoveto{\pgfqpoint{0.000000in}{0.000000in}}%
\pgfpathlineto{\pgfqpoint{0.000000in}{-0.048611in}}%
\pgfusepath{stroke,fill}%
}%
\begin{pgfscope}%
\pgfsys@transformshift{4.516715in}{0.565123in}%
\pgfsys@useobject{currentmarker}{}%
\end{pgfscope}%
\end{pgfscope}%
\begin{pgfscope}%
\definecolor{textcolor}{rgb}{0.000000,0.000000,0.000000}%
\pgfsetstrokecolor{textcolor}%
\pgfsetfillcolor{textcolor}%
\pgftext[x=4.516715in,y=0.467901in,,top]{\color{textcolor}\rmfamily\fontsize{10.000000}{12.000000}\selectfont \(\displaystyle {9.8290}\)}%
\end{pgfscope}%
\begin{pgfscope}%
\pgfpathrectangle{\pgfqpoint{0.719445in}{0.565123in}}{\pgfqpoint{5.430555in}{2.585803in}}%
\pgfusepath{clip}%
\pgfsetbuttcap%
\pgfsetroundjoin%
\pgfsetlinewidth{0.803000pt}%
\definecolor{currentstroke}{rgb}{0.690196,0.690196,0.690196}%
\pgfsetstrokecolor{currentstroke}%
\pgfsetdash{{0.800000pt}{1.320000pt}}{0.000000pt}%
\pgfpathmoveto{\pgfqpoint{5.404322in}{0.565123in}}%
\pgfpathlineto{\pgfqpoint{5.404322in}{3.150926in}}%
\pgfusepath{stroke}%
\end{pgfscope}%
\begin{pgfscope}%
\pgfsetbuttcap%
\pgfsetroundjoin%
\definecolor{currentfill}{rgb}{0.000000,0.000000,0.000000}%
\pgfsetfillcolor{currentfill}%
\pgfsetlinewidth{0.803000pt}%
\definecolor{currentstroke}{rgb}{0.000000,0.000000,0.000000}%
\pgfsetstrokecolor{currentstroke}%
\pgfsetdash{}{0pt}%
\pgfsys@defobject{currentmarker}{\pgfqpoint{0.000000in}{-0.048611in}}{\pgfqpoint{0.000000in}{0.000000in}}{%
\pgfpathmoveto{\pgfqpoint{0.000000in}{0.000000in}}%
\pgfpathlineto{\pgfqpoint{0.000000in}{-0.048611in}}%
\pgfusepath{stroke,fill}%
}%
\begin{pgfscope}%
\pgfsys@transformshift{5.404322in}{0.565123in}%
\pgfsys@useobject{currentmarker}{}%
\end{pgfscope}%
\end{pgfscope}%
\begin{pgfscope}%
\definecolor{textcolor}{rgb}{0.000000,0.000000,0.000000}%
\pgfsetstrokecolor{textcolor}%
\pgfsetfillcolor{textcolor}%
\pgftext[x=5.404322in,y=0.467901in,,top]{\color{textcolor}\rmfamily\fontsize{10.000000}{12.000000}\selectfont \(\displaystyle {9.8295}\)}%
\end{pgfscope}%
\begin{pgfscope}%
\definecolor{textcolor}{rgb}{0.000000,0.000000,0.000000}%
\pgfsetstrokecolor{textcolor}%
\pgfsetfillcolor{textcolor}%
\pgftext[x=3.434723in,y=0.288889in,,top]{\color{textcolor}\rmfamily\fontsize{10.000000}{12.000000}\selectfont Frequency (Hz)}%
\end{pgfscope}%
\begin{pgfscope}%
\definecolor{textcolor}{rgb}{0.000000,0.000000,0.000000}%
\pgfsetstrokecolor{textcolor}%
\pgfsetfillcolor{textcolor}%
\pgftext[x=6.150000in,y=0.302778in,right,top]{\color{textcolor}\rmfamily\fontsize{10.000000}{12.000000}\selectfont \(\displaystyle \times{10^{6}}{}\)}%
\end{pgfscope}%
\begin{pgfscope}%
\pgfpathrectangle{\pgfqpoint{0.719445in}{0.565123in}}{\pgfqpoint{5.430555in}{2.585803in}}%
\pgfusepath{clip}%
\pgfsetbuttcap%
\pgfsetroundjoin%
\pgfsetlinewidth{0.803000pt}%
\definecolor{currentstroke}{rgb}{0.690196,0.690196,0.690196}%
\pgfsetstrokecolor{currentstroke}%
\pgfsetdash{{0.800000pt}{1.320000pt}}{0.000000pt}%
\pgfpathmoveto{\pgfqpoint{0.719445in}{0.671123in}}%
\pgfpathlineto{\pgfqpoint{6.150000in}{0.671123in}}%
\pgfusepath{stroke}%
\end{pgfscope}%
\begin{pgfscope}%
\pgfsetbuttcap%
\pgfsetroundjoin%
\definecolor{currentfill}{rgb}{0.000000,0.000000,0.000000}%
\pgfsetfillcolor{currentfill}%
\pgfsetlinewidth{0.803000pt}%
\definecolor{currentstroke}{rgb}{0.000000,0.000000,0.000000}%
\pgfsetstrokecolor{currentstroke}%
\pgfsetdash{}{0pt}%
\pgfsys@defobject{currentmarker}{\pgfqpoint{-0.048611in}{0.000000in}}{\pgfqpoint{-0.000000in}{0.000000in}}{%
\pgfpathmoveto{\pgfqpoint{-0.000000in}{0.000000in}}%
\pgfpathlineto{\pgfqpoint{-0.048611in}{0.000000in}}%
\pgfusepath{stroke,fill}%
}%
\begin{pgfscope}%
\pgfsys@transformshift{0.719445in}{0.671123in}%
\pgfsys@useobject{currentmarker}{}%
\end{pgfscope}%
\end{pgfscope}%
\begin{pgfscope}%
\definecolor{textcolor}{rgb}{0.000000,0.000000,0.000000}%
\pgfsetstrokecolor{textcolor}%
\pgfsetfillcolor{textcolor}%
\pgftext[x=0.552778in, y=0.622898in, left, base]{\color{textcolor}\rmfamily\fontsize{10.000000}{12.000000}\selectfont \(\displaystyle {0}\)}%
\end{pgfscope}%
\begin{pgfscope}%
\pgfpathrectangle{\pgfqpoint{0.719445in}{0.565123in}}{\pgfqpoint{5.430555in}{2.585803in}}%
\pgfusepath{clip}%
\pgfsetbuttcap%
\pgfsetroundjoin%
\pgfsetlinewidth{0.803000pt}%
\definecolor{currentstroke}{rgb}{0.690196,0.690196,0.690196}%
\pgfsetstrokecolor{currentstroke}%
\pgfsetdash{{0.800000pt}{1.320000pt}}{0.000000pt}%
\pgfpathmoveto{\pgfqpoint{0.719445in}{1.215299in}}%
\pgfpathlineto{\pgfqpoint{6.150000in}{1.215299in}}%
\pgfusepath{stroke}%
\end{pgfscope}%
\begin{pgfscope}%
\pgfsetbuttcap%
\pgfsetroundjoin%
\definecolor{currentfill}{rgb}{0.000000,0.000000,0.000000}%
\pgfsetfillcolor{currentfill}%
\pgfsetlinewidth{0.803000pt}%
\definecolor{currentstroke}{rgb}{0.000000,0.000000,0.000000}%
\pgfsetstrokecolor{currentstroke}%
\pgfsetdash{}{0pt}%
\pgfsys@defobject{currentmarker}{\pgfqpoint{-0.048611in}{0.000000in}}{\pgfqpoint{-0.000000in}{0.000000in}}{%
\pgfpathmoveto{\pgfqpoint{-0.000000in}{0.000000in}}%
\pgfpathlineto{\pgfqpoint{-0.048611in}{0.000000in}}%
\pgfusepath{stroke,fill}%
}%
\begin{pgfscope}%
\pgfsys@transformshift{0.719445in}{1.215299in}%
\pgfsys@useobject{currentmarker}{}%
\end{pgfscope}%
\end{pgfscope}%
\begin{pgfscope}%
\definecolor{textcolor}{rgb}{0.000000,0.000000,0.000000}%
\pgfsetstrokecolor{textcolor}%
\pgfsetfillcolor{textcolor}%
\pgftext[x=0.413889in, y=1.167074in, left, base]{\color{textcolor}\rmfamily\fontsize{10.000000}{12.000000}\selectfont \(\displaystyle {500}\)}%
\end{pgfscope}%
\begin{pgfscope}%
\pgfpathrectangle{\pgfqpoint{0.719445in}{0.565123in}}{\pgfqpoint{5.430555in}{2.585803in}}%
\pgfusepath{clip}%
\pgfsetbuttcap%
\pgfsetroundjoin%
\pgfsetlinewidth{0.803000pt}%
\definecolor{currentstroke}{rgb}{0.690196,0.690196,0.690196}%
\pgfsetstrokecolor{currentstroke}%
\pgfsetdash{{0.800000pt}{1.320000pt}}{0.000000pt}%
\pgfpathmoveto{\pgfqpoint{0.719445in}{1.759475in}}%
\pgfpathlineto{\pgfqpoint{6.150000in}{1.759475in}}%
\pgfusepath{stroke}%
\end{pgfscope}%
\begin{pgfscope}%
\pgfsetbuttcap%
\pgfsetroundjoin%
\definecolor{currentfill}{rgb}{0.000000,0.000000,0.000000}%
\pgfsetfillcolor{currentfill}%
\pgfsetlinewidth{0.803000pt}%
\definecolor{currentstroke}{rgb}{0.000000,0.000000,0.000000}%
\pgfsetstrokecolor{currentstroke}%
\pgfsetdash{}{0pt}%
\pgfsys@defobject{currentmarker}{\pgfqpoint{-0.048611in}{0.000000in}}{\pgfqpoint{-0.000000in}{0.000000in}}{%
\pgfpathmoveto{\pgfqpoint{-0.000000in}{0.000000in}}%
\pgfpathlineto{\pgfqpoint{-0.048611in}{0.000000in}}%
\pgfusepath{stroke,fill}%
}%
\begin{pgfscope}%
\pgfsys@transformshift{0.719445in}{1.759475in}%
\pgfsys@useobject{currentmarker}{}%
\end{pgfscope}%
\end{pgfscope}%
\begin{pgfscope}%
\definecolor{textcolor}{rgb}{0.000000,0.000000,0.000000}%
\pgfsetstrokecolor{textcolor}%
\pgfsetfillcolor{textcolor}%
\pgftext[x=0.344444in, y=1.711249in, left, base]{\color{textcolor}\rmfamily\fontsize{10.000000}{12.000000}\selectfont \(\displaystyle {1000}\)}%
\end{pgfscope}%
\begin{pgfscope}%
\pgfpathrectangle{\pgfqpoint{0.719445in}{0.565123in}}{\pgfqpoint{5.430555in}{2.585803in}}%
\pgfusepath{clip}%
\pgfsetbuttcap%
\pgfsetroundjoin%
\pgfsetlinewidth{0.803000pt}%
\definecolor{currentstroke}{rgb}{0.690196,0.690196,0.690196}%
\pgfsetstrokecolor{currentstroke}%
\pgfsetdash{{0.800000pt}{1.320000pt}}{0.000000pt}%
\pgfpathmoveto{\pgfqpoint{0.719445in}{2.303650in}}%
\pgfpathlineto{\pgfqpoint{6.150000in}{2.303650in}}%
\pgfusepath{stroke}%
\end{pgfscope}%
\begin{pgfscope}%
\pgfsetbuttcap%
\pgfsetroundjoin%
\definecolor{currentfill}{rgb}{0.000000,0.000000,0.000000}%
\pgfsetfillcolor{currentfill}%
\pgfsetlinewidth{0.803000pt}%
\definecolor{currentstroke}{rgb}{0.000000,0.000000,0.000000}%
\pgfsetstrokecolor{currentstroke}%
\pgfsetdash{}{0pt}%
\pgfsys@defobject{currentmarker}{\pgfqpoint{-0.048611in}{0.000000in}}{\pgfqpoint{-0.000000in}{0.000000in}}{%
\pgfpathmoveto{\pgfqpoint{-0.000000in}{0.000000in}}%
\pgfpathlineto{\pgfqpoint{-0.048611in}{0.000000in}}%
\pgfusepath{stroke,fill}%
}%
\begin{pgfscope}%
\pgfsys@transformshift{0.719445in}{2.303650in}%
\pgfsys@useobject{currentmarker}{}%
\end{pgfscope}%
\end{pgfscope}%
\begin{pgfscope}%
\definecolor{textcolor}{rgb}{0.000000,0.000000,0.000000}%
\pgfsetstrokecolor{textcolor}%
\pgfsetfillcolor{textcolor}%
\pgftext[x=0.344444in, y=2.255425in, left, base]{\color{textcolor}\rmfamily\fontsize{10.000000}{12.000000}\selectfont \(\displaystyle {1500}\)}%
\end{pgfscope}%
\begin{pgfscope}%
\pgfpathrectangle{\pgfqpoint{0.719445in}{0.565123in}}{\pgfqpoint{5.430555in}{2.585803in}}%
\pgfusepath{clip}%
\pgfsetbuttcap%
\pgfsetroundjoin%
\pgfsetlinewidth{0.803000pt}%
\definecolor{currentstroke}{rgb}{0.690196,0.690196,0.690196}%
\pgfsetstrokecolor{currentstroke}%
\pgfsetdash{{0.800000pt}{1.320000pt}}{0.000000pt}%
\pgfpathmoveto{\pgfqpoint{0.719445in}{2.847826in}}%
\pgfpathlineto{\pgfqpoint{6.150000in}{2.847826in}}%
\pgfusepath{stroke}%
\end{pgfscope}%
\begin{pgfscope}%
\pgfsetbuttcap%
\pgfsetroundjoin%
\definecolor{currentfill}{rgb}{0.000000,0.000000,0.000000}%
\pgfsetfillcolor{currentfill}%
\pgfsetlinewidth{0.803000pt}%
\definecolor{currentstroke}{rgb}{0.000000,0.000000,0.000000}%
\pgfsetstrokecolor{currentstroke}%
\pgfsetdash{}{0pt}%
\pgfsys@defobject{currentmarker}{\pgfqpoint{-0.048611in}{0.000000in}}{\pgfqpoint{-0.000000in}{0.000000in}}{%
\pgfpathmoveto{\pgfqpoint{-0.000000in}{0.000000in}}%
\pgfpathlineto{\pgfqpoint{-0.048611in}{0.000000in}}%
\pgfusepath{stroke,fill}%
}%
\begin{pgfscope}%
\pgfsys@transformshift{0.719445in}{2.847826in}%
\pgfsys@useobject{currentmarker}{}%
\end{pgfscope}%
\end{pgfscope}%
\begin{pgfscope}%
\definecolor{textcolor}{rgb}{0.000000,0.000000,0.000000}%
\pgfsetstrokecolor{textcolor}%
\pgfsetfillcolor{textcolor}%
\pgftext[x=0.344444in, y=2.799601in, left, base]{\color{textcolor}\rmfamily\fontsize{10.000000}{12.000000}\selectfont \(\displaystyle {2000}\)}%
\end{pgfscope}%
\begin{pgfscope}%
\definecolor{textcolor}{rgb}{0.000000,0.000000,0.000000}%
\pgfsetstrokecolor{textcolor}%
\pgfsetfillcolor{textcolor}%
\pgftext[x=0.288889in,y=1.858025in,,bottom,rotate=90.000000]{\color{textcolor}\rmfamily\fontsize{10.000000}{12.000000}\selectfont Impedance |Z| (\(\displaystyle \Omega)\)}%
\end{pgfscope}%
\begin{pgfscope}%
\pgfpathrectangle{\pgfqpoint{0.719445in}{0.565123in}}{\pgfqpoint{5.430555in}{2.585803in}}%
\pgfusepath{clip}%
\pgfsetrectcap%
\pgfsetroundjoin%
\pgfsetlinewidth{2.007500pt}%
\definecolor{currentstroke}{rgb}{0.000000,0.000000,0.000000}%
\pgfsetstrokecolor{currentstroke}%
\pgfsetdash{}{0pt}%
\pgfpathmoveto{\pgfqpoint{0.966289in}{0.760912in}}%
\pgfpathlineto{\pgfqpoint{0.971614in}{1.154351in}}%
\pgfpathlineto{\pgfqpoint{0.976940in}{1.141618in}}%
\pgfpathlineto{\pgfqpoint{0.982266in}{1.154351in}}%
\pgfpathlineto{\pgfqpoint{0.998243in}{1.154569in}}%
\pgfpathlineto{\pgfqpoint{1.003568in}{1.147930in}}%
\pgfpathlineto{\pgfqpoint{1.030196in}{1.147930in}}%
\pgfpathlineto{\pgfqpoint{1.035522in}{1.141400in}}%
\pgfpathlineto{\pgfqpoint{1.056825in}{1.141618in}}%
\pgfpathlineto{\pgfqpoint{1.062150in}{1.135087in}}%
\pgfpathlineto{\pgfqpoint{1.078127in}{1.135087in}}%
\pgfpathlineto{\pgfqpoint{1.083453in}{1.129102in}}%
\pgfpathlineto{\pgfqpoint{1.110081in}{1.129102in}}%
\pgfpathlineto{\pgfqpoint{1.115407in}{1.123333in}}%
\pgfpathlineto{\pgfqpoint{1.131384in}{1.123116in}}%
\pgfpathlineto{\pgfqpoint{1.136709in}{1.117347in}}%
\pgfpathlineto{\pgfqpoint{1.152686in}{1.117347in}}%
\pgfpathlineto{\pgfqpoint{1.158012in}{1.111688in}}%
\pgfpathlineto{\pgfqpoint{1.184640in}{1.111688in}}%
\pgfpathlineto{\pgfqpoint{1.189966in}{1.106246in}}%
\pgfpathlineto{\pgfqpoint{1.205942in}{1.106246in}}%
\pgfpathlineto{\pgfqpoint{1.211268in}{1.100913in}}%
\pgfpathlineto{\pgfqpoint{1.227245in}{1.100913in}}%
\pgfpathlineto{\pgfqpoint{1.232571in}{1.095689in}}%
\pgfpathlineto{\pgfqpoint{1.248548in}{1.095689in}}%
\pgfpathlineto{\pgfqpoint{1.253873in}{1.090574in}}%
\pgfpathlineto{\pgfqpoint{1.275176in}{1.090574in}}%
\pgfpathlineto{\pgfqpoint{1.280501in}{1.085567in}}%
\pgfpathlineto{\pgfqpoint{1.296478in}{1.085567in}}%
\pgfpathlineto{\pgfqpoint{1.301804in}{1.080670in}}%
\pgfpathlineto{\pgfqpoint{1.312455in}{1.080670in}}%
\pgfpathlineto{\pgfqpoint{1.317781in}{1.075881in}}%
\pgfpathlineto{\pgfqpoint{1.333758in}{1.075881in}}%
\pgfpathlineto{\pgfqpoint{1.339083in}{1.071310in}}%
\pgfpathlineto{\pgfqpoint{1.349735in}{1.071310in}}%
\pgfpathlineto{\pgfqpoint{1.355060in}{1.066739in}}%
\pgfpathlineto{\pgfqpoint{1.371037in}{1.066739in}}%
\pgfpathlineto{\pgfqpoint{1.376363in}{1.062277in}}%
\pgfpathlineto{\pgfqpoint{1.387014in}{1.062277in}}%
\pgfpathlineto{\pgfqpoint{1.392340in}{1.057923in}}%
\pgfpathlineto{\pgfqpoint{1.408317in}{1.057923in}}%
\pgfpathlineto{\pgfqpoint{1.413642in}{1.053679in}}%
\pgfpathlineto{\pgfqpoint{1.424294in}{1.053679in}}%
\pgfpathlineto{\pgfqpoint{1.429619in}{1.049543in}}%
\pgfpathlineto{\pgfqpoint{1.445596in}{1.049543in}}%
\pgfpathlineto{\pgfqpoint{1.450922in}{1.045407in}}%
\pgfpathlineto{\pgfqpoint{1.461573in}{1.045407in}}%
\pgfpathlineto{\pgfqpoint{1.466899in}{1.041380in}}%
\pgfpathlineto{\pgfqpoint{1.477550in}{1.041380in}}%
\pgfpathlineto{\pgfqpoint{1.482876in}{1.037571in}}%
\pgfpathlineto{\pgfqpoint{1.493527in}{1.037571in}}%
\pgfpathlineto{\pgfqpoint{1.498853in}{1.033653in}}%
\pgfpathlineto{\pgfqpoint{1.509504in}{1.033653in}}%
\pgfpathlineto{\pgfqpoint{1.514830in}{1.029953in}}%
\pgfpathlineto{\pgfqpoint{1.525481in}{1.029953in}}%
\pgfpathlineto{\pgfqpoint{1.530806in}{1.026252in}}%
\pgfpathlineto{\pgfqpoint{1.541458in}{1.026252in}}%
\pgfpathlineto{\pgfqpoint{1.546783in}{1.022661in}}%
\pgfpathlineto{\pgfqpoint{1.557435in}{1.022661in}}%
\pgfpathlineto{\pgfqpoint{1.562760in}{1.019069in}}%
\pgfpathlineto{\pgfqpoint{1.573412in}{1.019069in}}%
\pgfpathlineto{\pgfqpoint{1.578737in}{1.015695in}}%
\pgfpathlineto{\pgfqpoint{1.584063in}{1.015695in}}%
\pgfpathlineto{\pgfqpoint{1.589389in}{1.012213in}}%
\pgfpathlineto{\pgfqpoint{1.600040in}{1.012213in}}%
\pgfpathlineto{\pgfqpoint{1.605365in}{1.008948in}}%
\pgfpathlineto{\pgfqpoint{1.616017in}{1.008948in}}%
\pgfpathlineto{\pgfqpoint{1.621342in}{1.005682in}}%
\pgfpathlineto{\pgfqpoint{1.631994in}{1.005682in}}%
\pgfpathlineto{\pgfqpoint{1.637319in}{1.002417in}}%
\pgfpathlineto{\pgfqpoint{1.642645in}{1.002417in}}%
\pgfpathlineto{\pgfqpoint{1.647971in}{0.999261in}}%
\pgfpathlineto{\pgfqpoint{1.658622in}{0.999261in}}%
\pgfpathlineto{\pgfqpoint{1.663947in}{0.996214in}}%
\pgfpathlineto{\pgfqpoint{1.669273in}{0.996214in}}%
\pgfpathlineto{\pgfqpoint{1.674599in}{0.993166in}}%
\pgfpathlineto{\pgfqpoint{1.685250in}{0.993166in}}%
\pgfpathlineto{\pgfqpoint{1.690576in}{0.990228in}}%
\pgfpathlineto{\pgfqpoint{1.695901in}{0.990010in}}%
\pgfpathlineto{\pgfqpoint{1.701227in}{0.987072in}}%
\pgfpathlineto{\pgfqpoint{1.706553in}{0.987289in}}%
\pgfpathlineto{\pgfqpoint{1.711878in}{0.984242in}}%
\pgfpathlineto{\pgfqpoint{1.722530in}{0.984242in}}%
\pgfpathlineto{\pgfqpoint{1.727855in}{0.981412in}}%
\pgfpathlineto{\pgfqpoint{1.733181in}{0.981412in}}%
\pgfpathlineto{\pgfqpoint{1.738506in}{0.978583in}}%
\pgfpathlineto{\pgfqpoint{1.749158in}{0.978583in}}%
\pgfpathlineto{\pgfqpoint{1.754483in}{0.975862in}}%
\pgfpathlineto{\pgfqpoint{1.759809in}{0.975862in}}%
\pgfpathlineto{\pgfqpoint{1.765135in}{0.973141in}}%
\pgfpathlineto{\pgfqpoint{1.770460in}{0.973141in}}%
\pgfpathlineto{\pgfqpoint{1.775786in}{0.970529in}}%
\pgfpathlineto{\pgfqpoint{1.781112in}{0.970529in}}%
\pgfpathlineto{\pgfqpoint{1.786437in}{0.967917in}}%
\pgfpathlineto{\pgfqpoint{1.791763in}{0.967917in}}%
\pgfpathlineto{\pgfqpoint{1.797088in}{0.965413in}}%
\pgfpathlineto{\pgfqpoint{1.802414in}{0.965413in}}%
\pgfpathlineto{\pgfqpoint{1.807740in}{0.962910in}}%
\pgfpathlineto{\pgfqpoint{1.818391in}{0.962910in}}%
\pgfpathlineto{\pgfqpoint{1.823717in}{0.960407in}}%
\pgfpathlineto{\pgfqpoint{1.829042in}{0.960407in}}%
\pgfpathlineto{\pgfqpoint{1.834368in}{0.958013in}}%
\pgfpathlineto{\pgfqpoint{1.839694in}{0.958013in}}%
\pgfpathlineto{\pgfqpoint{1.845019in}{0.955618in}}%
\pgfpathlineto{\pgfqpoint{1.850345in}{0.955618in}}%
\pgfpathlineto{\pgfqpoint{1.855670in}{0.953333in}}%
\pgfpathlineto{\pgfqpoint{1.860996in}{0.953333in}}%
\pgfpathlineto{\pgfqpoint{1.871647in}{0.948762in}}%
\pgfpathlineto{\pgfqpoint{1.876973in}{0.948762in}}%
\pgfpathlineto{\pgfqpoint{1.882299in}{0.946476in}}%
\pgfpathlineto{\pgfqpoint{1.887624in}{0.946476in}}%
\pgfpathlineto{\pgfqpoint{1.892950in}{0.944299in}}%
\pgfpathlineto{\pgfqpoint{1.898276in}{0.944299in}}%
\pgfpathlineto{\pgfqpoint{1.903601in}{0.942123in}}%
\pgfpathlineto{\pgfqpoint{1.908927in}{0.942123in}}%
\pgfpathlineto{\pgfqpoint{1.914253in}{0.939946in}}%
\pgfpathlineto{\pgfqpoint{1.919578in}{0.939946in}}%
\pgfpathlineto{\pgfqpoint{1.924904in}{0.937878in}}%
\pgfpathlineto{\pgfqpoint{1.930229in}{0.937878in}}%
\pgfpathlineto{\pgfqpoint{1.940881in}{0.933742in}}%
\pgfpathlineto{\pgfqpoint{1.946206in}{0.933742in}}%
\pgfpathlineto{\pgfqpoint{1.951532in}{0.931675in}}%
\pgfpathlineto{\pgfqpoint{1.956858in}{0.931675in}}%
\pgfpathlineto{\pgfqpoint{1.967509in}{0.927757in}}%
\pgfpathlineto{\pgfqpoint{1.972835in}{0.927757in}}%
\pgfpathlineto{\pgfqpoint{1.983486in}{0.923947in}}%
\pgfpathlineto{\pgfqpoint{1.988811in}{0.923947in}}%
\pgfpathlineto{\pgfqpoint{1.994137in}{0.922097in}}%
\pgfpathlineto{\pgfqpoint{1.999463in}{0.922097in}}%
\pgfpathlineto{\pgfqpoint{2.010114in}{0.918288in}}%
\pgfpathlineto{\pgfqpoint{2.015440in}{0.918288in}}%
\pgfpathlineto{\pgfqpoint{2.026091in}{0.914696in}}%
\pgfpathlineto{\pgfqpoint{2.031417in}{0.914696in}}%
\pgfpathlineto{\pgfqpoint{2.042068in}{0.911214in}}%
\pgfpathlineto{\pgfqpoint{2.047394in}{0.911214in}}%
\pgfpathlineto{\pgfqpoint{2.063370in}{0.906207in}}%
\pgfpathlineto{\pgfqpoint{2.068696in}{0.906207in}}%
\pgfpathlineto{\pgfqpoint{2.084673in}{0.901418in}}%
\pgfpathlineto{\pgfqpoint{2.089999in}{0.901418in}}%
\pgfpathlineto{\pgfqpoint{2.105976in}{0.896739in}}%
\pgfpathlineto{\pgfqpoint{2.111301in}{0.896739in}}%
\pgfpathlineto{\pgfqpoint{2.132604in}{0.890861in}}%
\pgfpathlineto{\pgfqpoint{2.137929in}{0.890861in}}%
\pgfpathlineto{\pgfqpoint{2.164558in}{0.883678in}}%
\pgfpathlineto{\pgfqpoint{2.169883in}{0.883678in}}%
\pgfpathlineto{\pgfqpoint{2.212488in}{0.873339in}}%
\pgfpathlineto{\pgfqpoint{2.217814in}{0.873339in}}%
\pgfpathlineto{\pgfqpoint{2.468119in}{0.817724in}}%
\pgfpathlineto{\pgfqpoint{2.478770in}{0.815765in}}%
\pgfpathlineto{\pgfqpoint{2.505399in}{0.809779in}}%
\pgfpathlineto{\pgfqpoint{2.649191in}{0.778326in}}%
\pgfpathlineto{\pgfqpoint{2.665168in}{0.775170in}}%
\pgfpathlineto{\pgfqpoint{2.766355in}{0.753511in}}%
\pgfpathlineto{\pgfqpoint{2.942101in}{0.716508in}}%
\pgfpathlineto{\pgfqpoint{3.000683in}{0.704427in}}%
\pgfpathlineto{\pgfqpoint{3.117847in}{0.683857in}}%
\pgfpathlineto{\pgfqpoint{3.144475in}{0.682660in}}%
\pgfpathlineto{\pgfqpoint{3.181755in}{0.683639in}}%
\pgfpathlineto{\pgfqpoint{3.213709in}{0.687775in}}%
\pgfpathlineto{\pgfqpoint{3.256314in}{0.695720in}}%
\pgfpathlineto{\pgfqpoint{3.432060in}{0.736316in}}%
\pgfpathlineto{\pgfqpoint{3.634434in}{0.791168in}}%
\pgfpathlineto{\pgfqpoint{3.650411in}{0.796284in}}%
\pgfpathlineto{\pgfqpoint{3.698342in}{0.810759in}}%
\pgfpathlineto{\pgfqpoint{3.714319in}{0.815983in}}%
\pgfpathlineto{\pgfqpoint{3.730296in}{0.820989in}}%
\pgfpathlineto{\pgfqpoint{3.756924in}{0.829152in}}%
\pgfpathlineto{\pgfqpoint{3.783552in}{0.838403in}}%
\pgfpathlineto{\pgfqpoint{3.836808in}{0.856687in}}%
\pgfpathlineto{\pgfqpoint{3.858111in}{0.864306in}}%
\pgfpathlineto{\pgfqpoint{3.874088in}{0.869203in}}%
\pgfpathlineto{\pgfqpoint{3.895390in}{0.877801in}}%
\pgfpathlineto{\pgfqpoint{3.906042in}{0.880413in}}%
\pgfpathlineto{\pgfqpoint{3.911367in}{0.883134in}}%
\pgfpathlineto{\pgfqpoint{3.916693in}{0.884440in}}%
\pgfpathlineto{\pgfqpoint{3.927344in}{0.888902in}}%
\pgfpathlineto{\pgfqpoint{3.932670in}{0.890317in}}%
\pgfpathlineto{\pgfqpoint{3.937996in}{0.893147in}}%
\pgfpathlineto{\pgfqpoint{3.943321in}{0.894671in}}%
\pgfpathlineto{\pgfqpoint{3.948647in}{0.897718in}}%
\pgfpathlineto{\pgfqpoint{3.964624in}{0.902289in}}%
\pgfpathlineto{\pgfqpoint{3.969949in}{0.905554in}}%
\pgfpathlineto{\pgfqpoint{3.985926in}{0.910561in}}%
\pgfpathlineto{\pgfqpoint{3.991252in}{0.913934in}}%
\pgfpathlineto{\pgfqpoint{4.023206in}{0.925144in}}%
\pgfpathlineto{\pgfqpoint{4.028531in}{0.929063in}}%
\pgfpathlineto{\pgfqpoint{4.071137in}{0.945605in}}%
\pgfpathlineto{\pgfqpoint{4.151021in}{0.980650in}}%
\pgfpathlineto{\pgfqpoint{4.161672in}{0.986310in}}%
\pgfpathlineto{\pgfqpoint{4.166998in}{0.986310in}}%
\pgfpathlineto{\pgfqpoint{4.188301in}{0.998282in}}%
\pgfpathlineto{\pgfqpoint{4.193626in}{0.998282in}}%
\pgfpathlineto{\pgfqpoint{4.214929in}{1.011124in}}%
\pgfpathlineto{\pgfqpoint{4.220254in}{1.011124in}}%
\pgfpathlineto{\pgfqpoint{4.236231in}{1.022008in}}%
\pgfpathlineto{\pgfqpoint{4.241557in}{1.022008in}}%
\pgfpathlineto{\pgfqpoint{4.257534in}{1.033000in}}%
\pgfpathlineto{\pgfqpoint{4.262860in}{1.033000in}}%
\pgfpathlineto{\pgfqpoint{4.273511in}{1.040619in}}%
\pgfpathlineto{\pgfqpoint{4.278837in}{1.040619in}}%
\pgfpathlineto{\pgfqpoint{4.289488in}{1.048672in}}%
\pgfpathlineto{\pgfqpoint{4.294813in}{1.048672in}}%
\pgfpathlineto{\pgfqpoint{4.310790in}{1.061406in}}%
\pgfpathlineto{\pgfqpoint{4.316116in}{1.061406in}}%
\pgfpathlineto{\pgfqpoint{4.321442in}{1.065759in}}%
\pgfpathlineto{\pgfqpoint{4.326767in}{1.065759in}}%
\pgfpathlineto{\pgfqpoint{4.337419in}{1.074902in}}%
\pgfpathlineto{\pgfqpoint{4.342744in}{1.074902in}}%
\pgfpathlineto{\pgfqpoint{4.353395in}{1.084479in}}%
\pgfpathlineto{\pgfqpoint{4.358721in}{1.084479in}}%
\pgfpathlineto{\pgfqpoint{4.364047in}{1.089377in}}%
\pgfpathlineto{\pgfqpoint{4.369372in}{1.089377in}}%
\pgfpathlineto{\pgfqpoint{4.380024in}{1.099716in}}%
\pgfpathlineto{\pgfqpoint{4.385349in}{1.099716in}}%
\pgfpathlineto{\pgfqpoint{4.390675in}{1.104940in}}%
\pgfpathlineto{\pgfqpoint{4.396001in}{1.104940in}}%
\pgfpathlineto{\pgfqpoint{4.401326in}{1.110382in}}%
\pgfpathlineto{\pgfqpoint{4.406652in}{1.110382in}}%
\pgfpathlineto{\pgfqpoint{4.411978in}{1.116041in}}%
\pgfpathlineto{\pgfqpoint{4.417303in}{1.116803in}}%
\pgfpathlineto{\pgfqpoint{4.422629in}{1.122571in}}%
\pgfpathlineto{\pgfqpoint{4.427954in}{1.122571in}}%
\pgfpathlineto{\pgfqpoint{4.438606in}{1.134543in}}%
\pgfpathlineto{\pgfqpoint{4.443931in}{1.134543in}}%
\pgfpathlineto{\pgfqpoint{4.449257in}{1.140747in}}%
\pgfpathlineto{\pgfqpoint{4.454583in}{1.140747in}}%
\pgfpathlineto{\pgfqpoint{4.459908in}{1.147168in}}%
\pgfpathlineto{\pgfqpoint{4.465234in}{1.147168in}}%
\pgfpathlineto{\pgfqpoint{4.470560in}{1.153698in}}%
\pgfpathlineto{\pgfqpoint{4.475885in}{1.153698in}}%
\pgfpathlineto{\pgfqpoint{4.481211in}{1.160446in}}%
\pgfpathlineto{\pgfqpoint{4.486536in}{1.160446in}}%
\pgfpathlineto{\pgfqpoint{4.491862in}{1.167411in}}%
\pgfpathlineto{\pgfqpoint{4.497188in}{1.167411in}}%
\pgfpathlineto{\pgfqpoint{4.502513in}{1.174486in}}%
\pgfpathlineto{\pgfqpoint{4.507839in}{1.174486in}}%
\pgfpathlineto{\pgfqpoint{4.513165in}{1.181778in}}%
\pgfpathlineto{\pgfqpoint{4.518490in}{1.181778in}}%
\pgfpathlineto{\pgfqpoint{4.523816in}{1.189287in}}%
\pgfpathlineto{\pgfqpoint{4.534467in}{1.189287in}}%
\pgfpathlineto{\pgfqpoint{4.539793in}{1.197123in}}%
\pgfpathlineto{\pgfqpoint{4.545119in}{1.197123in}}%
\pgfpathlineto{\pgfqpoint{4.550444in}{1.205068in}}%
\pgfpathlineto{\pgfqpoint{4.555770in}{1.205068in}}%
\pgfpathlineto{\pgfqpoint{4.561095in}{1.213340in}}%
\pgfpathlineto{\pgfqpoint{4.566421in}{1.213340in}}%
\pgfpathlineto{\pgfqpoint{4.571747in}{1.221829in}}%
\pgfpathlineto{\pgfqpoint{4.582398in}{1.221829in}}%
\pgfpathlineto{\pgfqpoint{4.587724in}{1.230536in}}%
\pgfpathlineto{\pgfqpoint{4.593049in}{1.230536in}}%
\pgfpathlineto{\pgfqpoint{4.598375in}{1.239569in}}%
\pgfpathlineto{\pgfqpoint{4.609026in}{1.239569in}}%
\pgfpathlineto{\pgfqpoint{4.614352in}{1.248929in}}%
\pgfpathlineto{\pgfqpoint{4.619677in}{1.248929in}}%
\pgfpathlineto{\pgfqpoint{4.625003in}{1.258506in}}%
\pgfpathlineto{\pgfqpoint{4.635654in}{1.258506in}}%
\pgfpathlineto{\pgfqpoint{4.640980in}{1.268519in}}%
\pgfpathlineto{\pgfqpoint{4.646306in}{1.268519in}}%
\pgfpathlineto{\pgfqpoint{4.651631in}{1.278859in}}%
\pgfpathlineto{\pgfqpoint{4.662283in}{1.278859in}}%
\pgfpathlineto{\pgfqpoint{4.667608in}{1.289416in}}%
\pgfpathlineto{\pgfqpoint{4.672934in}{1.289416in}}%
\pgfpathlineto{\pgfqpoint{4.678259in}{1.300517in}}%
\pgfpathlineto{\pgfqpoint{4.688911in}{1.300517in}}%
\pgfpathlineto{\pgfqpoint{4.694236in}{1.311945in}}%
\pgfpathlineto{\pgfqpoint{4.704888in}{1.311945in}}%
\pgfpathlineto{\pgfqpoint{4.710213in}{1.323808in}}%
\pgfpathlineto{\pgfqpoint{4.720865in}{1.323808in}}%
\pgfpathlineto{\pgfqpoint{4.726190in}{1.337738in}}%
\pgfpathlineto{\pgfqpoint{4.731516in}{1.336106in}}%
\pgfpathlineto{\pgfqpoint{4.736842in}{1.337738in}}%
\pgfpathlineto{\pgfqpoint{4.742167in}{1.348840in}}%
\pgfpathlineto{\pgfqpoint{4.747493in}{1.350581in}}%
\pgfpathlineto{\pgfqpoint{4.752818in}{1.362009in}}%
\pgfpathlineto{\pgfqpoint{4.758144in}{1.363968in}}%
\pgfpathlineto{\pgfqpoint{4.763470in}{1.363968in}}%
\pgfpathlineto{\pgfqpoint{4.768795in}{1.377899in}}%
\pgfpathlineto{\pgfqpoint{4.784772in}{1.377899in}}%
\pgfpathlineto{\pgfqpoint{4.790098in}{1.392265in}}%
\pgfpathlineto{\pgfqpoint{4.800749in}{1.392265in}}%
\pgfpathlineto{\pgfqpoint{4.806075in}{1.407393in}}%
\pgfpathlineto{\pgfqpoint{4.816726in}{1.407393in}}%
\pgfpathlineto{\pgfqpoint{4.822052in}{1.423065in}}%
\pgfpathlineto{\pgfqpoint{4.832703in}{1.423065in}}%
\pgfpathlineto{\pgfqpoint{4.838029in}{1.439390in}}%
\pgfpathlineto{\pgfqpoint{4.854006in}{1.439390in}}%
\pgfpathlineto{\pgfqpoint{4.859331in}{1.456478in}}%
\pgfpathlineto{\pgfqpoint{4.875308in}{1.456478in}}%
\pgfpathlineto{\pgfqpoint{4.880634in}{1.474327in}}%
\pgfpathlineto{\pgfqpoint{4.891285in}{1.474327in}}%
\pgfpathlineto{\pgfqpoint{4.896611in}{1.493046in}}%
\pgfpathlineto{\pgfqpoint{4.912588in}{1.493046in}}%
\pgfpathlineto{\pgfqpoint{4.917913in}{1.512637in}}%
\pgfpathlineto{\pgfqpoint{4.933890in}{1.512637in}}%
\pgfpathlineto{\pgfqpoint{4.939216in}{1.533098in}}%
\pgfpathlineto{\pgfqpoint{4.949867in}{1.533098in}}%
\pgfpathlineto{\pgfqpoint{4.955193in}{1.554647in}}%
\pgfpathlineto{\pgfqpoint{4.971170in}{1.554647in}}%
\pgfpathlineto{\pgfqpoint{4.976495in}{1.577285in}}%
\pgfpathlineto{\pgfqpoint{4.997798in}{1.577285in}}%
\pgfpathlineto{\pgfqpoint{5.003124in}{1.601119in}}%
\pgfpathlineto{\pgfqpoint{5.019100in}{1.601119in}}%
\pgfpathlineto{\pgfqpoint{5.024426in}{1.626260in}}%
\pgfpathlineto{\pgfqpoint{5.040403in}{1.626260in}}%
\pgfpathlineto{\pgfqpoint{5.045729in}{1.652707in}}%
\pgfpathlineto{\pgfqpoint{5.061706in}{1.652707in}}%
\pgfpathlineto{\pgfqpoint{5.067031in}{1.680787in}}%
\pgfpathlineto{\pgfqpoint{5.088334in}{1.680787in}}%
\pgfpathlineto{\pgfqpoint{5.093659in}{1.710390in}}%
\pgfpathlineto{\pgfqpoint{5.114962in}{1.710390in}}%
\pgfpathlineto{\pgfqpoint{5.120288in}{1.741734in}}%
\pgfpathlineto{\pgfqpoint{5.141590in}{1.741734in}}%
\pgfpathlineto{\pgfqpoint{5.146916in}{1.775147in}}%
\pgfpathlineto{\pgfqpoint{5.168218in}{1.775147in}}%
\pgfpathlineto{\pgfqpoint{5.173544in}{1.810627in}}%
\pgfpathlineto{\pgfqpoint{5.194847in}{1.810627in}}%
\pgfpathlineto{\pgfqpoint{5.200172in}{1.848502in}}%
\pgfpathlineto{\pgfqpoint{5.221475in}{1.848502in}}%
\pgfpathlineto{\pgfqpoint{5.226800in}{1.888880in}}%
\pgfpathlineto{\pgfqpoint{5.253429in}{1.888880in}}%
\pgfpathlineto{\pgfqpoint{5.258754in}{1.932196in}}%
\pgfpathlineto{\pgfqpoint{5.285382in}{1.932196in}}%
\pgfpathlineto{\pgfqpoint{5.290708in}{1.978669in}}%
\pgfpathlineto{\pgfqpoint{5.317336in}{1.978669in}}%
\pgfpathlineto{\pgfqpoint{5.322662in}{2.028624in}}%
\pgfpathlineto{\pgfqpoint{5.349290in}{2.028624in}}%
\pgfpathlineto{\pgfqpoint{5.354616in}{2.082606in}}%
\pgfpathlineto{\pgfqpoint{5.359941in}{2.088048in}}%
\pgfpathlineto{\pgfqpoint{5.365267in}{2.082606in}}%
\pgfpathlineto{\pgfqpoint{5.381244in}{2.082606in}}%
\pgfpathlineto{\pgfqpoint{5.386570in}{2.147145in}}%
\pgfpathlineto{\pgfqpoint{5.391895in}{2.140942in}}%
\pgfpathlineto{\pgfqpoint{5.407872in}{2.140942in}}%
\pgfpathlineto{\pgfqpoint{5.413198in}{2.147145in}}%
\pgfpathlineto{\pgfqpoint{5.418523in}{2.147145in}}%
\pgfpathlineto{\pgfqpoint{5.423849in}{2.140942in}}%
\pgfpathlineto{\pgfqpoint{5.429175in}{2.204393in}}%
\pgfpathlineto{\pgfqpoint{5.434500in}{2.211358in}}%
\pgfpathlineto{\pgfqpoint{5.455803in}{2.211358in}}%
\pgfpathlineto{\pgfqpoint{5.461129in}{2.204393in}}%
\pgfpathlineto{\pgfqpoint{5.466454in}{2.281448in}}%
\pgfpathlineto{\pgfqpoint{5.498408in}{2.281448in}}%
\pgfpathlineto{\pgfqpoint{5.503734in}{2.273394in}}%
\pgfpathlineto{\pgfqpoint{5.509059in}{2.281448in}}%
\pgfpathlineto{\pgfqpoint{5.514385in}{2.358177in}}%
\pgfpathlineto{\pgfqpoint{5.519711in}{2.349034in}}%
\pgfpathlineto{\pgfqpoint{5.525036in}{2.349034in}}%
\pgfpathlineto{\pgfqpoint{5.530362in}{2.358177in}}%
\pgfpathlineto{\pgfqpoint{5.551664in}{2.358177in}}%
\pgfpathlineto{\pgfqpoint{5.556990in}{2.442524in}}%
\pgfpathlineto{\pgfqpoint{5.599595in}{2.442524in}}%
\pgfpathlineto{\pgfqpoint{5.604921in}{2.535904in}}%
\pgfpathlineto{\pgfqpoint{5.647526in}{2.535904in}}%
\pgfpathlineto{\pgfqpoint{5.652852in}{2.639515in}}%
\pgfpathlineto{\pgfqpoint{5.706108in}{2.639515in}}%
\pgfpathlineto{\pgfqpoint{5.711434in}{2.755316in}}%
\pgfpathlineto{\pgfqpoint{5.764690in}{2.755316in}}%
\pgfpathlineto{\pgfqpoint{5.770016in}{2.885700in}}%
\pgfpathlineto{\pgfqpoint{5.833923in}{2.885700in}}%
\pgfpathlineto{\pgfqpoint{5.839249in}{3.033390in}}%
\pgfpathlineto{\pgfqpoint{5.903157in}{3.033390in}}%
\pgfpathlineto{\pgfqpoint{5.903157in}{3.033390in}}%
\pgfusepath{stroke}%
\end{pgfscope}%
\begin{pgfscope}%
\pgfsetrectcap%
\pgfsetmiterjoin%
\pgfsetlinewidth{0.803000pt}%
\definecolor{currentstroke}{rgb}{0.000000,0.000000,0.000000}%
\pgfsetstrokecolor{currentstroke}%
\pgfsetdash{}{0pt}%
\pgfpathmoveto{\pgfqpoint{0.719445in}{0.565123in}}%
\pgfpathlineto{\pgfqpoint{0.719445in}{3.150926in}}%
\pgfusepath{stroke}%
\end{pgfscope}%
\begin{pgfscope}%
\pgfsetrectcap%
\pgfsetmiterjoin%
\pgfsetlinewidth{0.803000pt}%
\definecolor{currentstroke}{rgb}{0.000000,0.000000,0.000000}%
\pgfsetstrokecolor{currentstroke}%
\pgfsetdash{}{0pt}%
\pgfpathmoveto{\pgfqpoint{6.150000in}{0.565123in}}%
\pgfpathlineto{\pgfqpoint{6.150000in}{3.150926in}}%
\pgfusepath{stroke}%
\end{pgfscope}%
\begin{pgfscope}%
\pgfsetrectcap%
\pgfsetmiterjoin%
\pgfsetlinewidth{0.803000pt}%
\definecolor{currentstroke}{rgb}{0.000000,0.000000,0.000000}%
\pgfsetstrokecolor{currentstroke}%
\pgfsetdash{}{0pt}%
\pgfpathmoveto{\pgfqpoint{0.719445in}{0.565123in}}%
\pgfpathlineto{\pgfqpoint{6.150000in}{0.565123in}}%
\pgfusepath{stroke}%
\end{pgfscope}%
\begin{pgfscope}%
\pgfsetrectcap%
\pgfsetmiterjoin%
\pgfsetlinewidth{0.803000pt}%
\definecolor{currentstroke}{rgb}{0.000000,0.000000,0.000000}%
\pgfsetstrokecolor{currentstroke}%
\pgfsetdash{}{0pt}%
\pgfpathmoveto{\pgfqpoint{0.719445in}{3.150926in}}%
\pgfpathlineto{\pgfqpoint{6.150000in}{3.150926in}}%
\pgfusepath{stroke}%
\end{pgfscope}%
\begin{pgfscope}%
\definecolor{textcolor}{rgb}{0.000000,0.000000,0.000000}%
\pgfsetstrokecolor{textcolor}%
\pgfsetfillcolor{textcolor}%
\pgftext[x=3.434723in,y=3.234260in,,base]{\color{textcolor}\rmfamily\fontsize{12.000000}{14.400000}\selectfont Measured Quartz Impedance vs Frequency}%
\end{pgfscope}%
\begin{pgfscope}%
\pgfpathrectangle{\pgfqpoint{0.719445in}{0.565123in}}{\pgfqpoint{5.430555in}{2.585803in}}%
\pgfusepath{clip}%
\pgfsetbuttcap%
\pgfsetroundjoin%
\definecolor{currentfill}{rgb}{0.000000,0.000000,0.000000}%
\pgfsetfillcolor{currentfill}%
\pgfsetlinewidth{1.003750pt}%
\definecolor{currentstroke}{rgb}{0.000000,0.000000,0.000000}%
\pgfsetstrokecolor{currentstroke}%
\pgfsetdash{}{0pt}%
\pgfsys@defobject{currentmarker}{\pgfqpoint{-0.041667in}{-0.041667in}}{\pgfqpoint{0.041667in}{0.041667in}}{%
\pgfpathmoveto{\pgfqpoint{0.000000in}{-0.041667in}}%
\pgfpathcurveto{\pgfqpoint{0.011050in}{-0.041667in}}{\pgfqpoint{0.021649in}{-0.037276in}}{\pgfqpoint{0.029463in}{-0.029463in}}%
\pgfpathcurveto{\pgfqpoint{0.037276in}{-0.021649in}}{\pgfqpoint{0.041667in}{-0.011050in}}{\pgfqpoint{0.041667in}{0.000000in}}%
\pgfpathcurveto{\pgfqpoint{0.041667in}{0.011050in}}{\pgfqpoint{0.037276in}{0.021649in}}{\pgfqpoint{0.029463in}{0.029463in}}%
\pgfpathcurveto{\pgfqpoint{0.021649in}{0.037276in}}{\pgfqpoint{0.011050in}{0.041667in}}{\pgfqpoint{0.000000in}{0.041667in}}%
\pgfpathcurveto{\pgfqpoint{-0.011050in}{0.041667in}}{\pgfqpoint{-0.021649in}{0.037276in}}{\pgfqpoint{-0.029463in}{0.029463in}}%
\pgfpathcurveto{\pgfqpoint{-0.037276in}{0.021649in}}{\pgfqpoint{-0.041667in}{0.011050in}}{\pgfqpoint{-0.041667in}{0.000000in}}%
\pgfpathcurveto{\pgfqpoint{-0.041667in}{-0.011050in}}{\pgfqpoint{-0.037276in}{-0.021649in}}{\pgfqpoint{-0.029463in}{-0.029463in}}%
\pgfpathcurveto{\pgfqpoint{-0.021649in}{-0.037276in}}{\pgfqpoint{-0.011050in}{-0.041667in}}{\pgfqpoint{0.000000in}{-0.041667in}}%
\pgfpathlineto{\pgfqpoint{0.000000in}{-0.041667in}}%
\pgfpathclose%
\pgfusepath{stroke,fill}%
}%
\begin{pgfscope}%
\pgfsys@transformshift{3.144475in}{0.682660in}%
\pgfsys@useobject{currentmarker}{}%
\end{pgfscope}%
\end{pgfscope}%
\begin{pgfscope}%
\pgfsetbuttcap%
\pgfsetmiterjoin%
\definecolor{currentfill}{rgb}{1.000000,1.000000,1.000000}%
\pgfsetfillcolor{currentfill}%
\pgfsetfillopacity{0.800000}%
\pgfsetlinewidth{1.003750pt}%
\definecolor{currentstroke}{rgb}{0.800000,0.800000,0.800000}%
\pgfsetstrokecolor{currentstroke}%
\pgfsetstrokeopacity{0.800000}%
\pgfsetdash{}{0pt}%
\pgfpathmoveto{\pgfqpoint{0.816668in}{2.637809in}}%
\pgfpathlineto{\pgfqpoint{2.202857in}{2.637809in}}%
\pgfpathquadraticcurveto{\pgfqpoint{2.230635in}{2.637809in}}{\pgfqpoint{2.230635in}{2.665587in}}%
\pgfpathlineto{\pgfqpoint{2.230635in}{3.053704in}}%
\pgfpathquadraticcurveto{\pgfqpoint{2.230635in}{3.081482in}}{\pgfqpoint{2.202857in}{3.081482in}}%
\pgfpathlineto{\pgfqpoint{0.816668in}{3.081482in}}%
\pgfpathquadraticcurveto{\pgfqpoint{0.788890in}{3.081482in}}{\pgfqpoint{0.788890in}{3.053704in}}%
\pgfpathlineto{\pgfqpoint{0.788890in}{2.665587in}}%
\pgfpathquadraticcurveto{\pgfqpoint{0.788890in}{2.637809in}}{\pgfqpoint{0.816668in}{2.637809in}}%
\pgfpathlineto{\pgfqpoint{0.816668in}{2.637809in}}%
\pgfpathclose%
\pgfusepath{stroke,fill}%
\end{pgfscope}%
\begin{pgfscope}%
\pgfsetrectcap%
\pgfsetroundjoin%
\pgfsetlinewidth{2.007500pt}%
\definecolor{currentstroke}{rgb}{0.000000,0.000000,0.000000}%
\pgfsetstrokecolor{currentstroke}%
\pgfsetdash{}{0pt}%
\pgfpathmoveto{\pgfqpoint{0.844445in}{2.970371in}}%
\pgfpathlineto{\pgfqpoint{0.983334in}{2.970371in}}%
\pgfpathlineto{\pgfqpoint{1.122223in}{2.970371in}}%
\pgfusepath{stroke}%
\end{pgfscope}%
\begin{pgfscope}%
\definecolor{textcolor}{rgb}{0.000000,0.000000,0.000000}%
\pgfsetstrokecolor{textcolor}%
\pgfsetfillcolor{textcolor}%
\pgftext[x=1.233334in,y=2.921760in,left,base]{\color{textcolor}\rmfamily\fontsize{10.000000}{12.000000}\selectfont |Z| (Quartz)}%
\end{pgfscope}%
\begin{pgfscope}%
\pgfsetbuttcap%
\pgfsetroundjoin%
\definecolor{currentfill}{rgb}{0.000000,0.000000,0.000000}%
\pgfsetfillcolor{currentfill}%
\pgfsetlinewidth{1.003750pt}%
\definecolor{currentstroke}{rgb}{0.000000,0.000000,0.000000}%
\pgfsetstrokecolor{currentstroke}%
\pgfsetdash{}{0pt}%
\pgfsys@defobject{currentmarker}{\pgfqpoint{-0.041667in}{-0.041667in}}{\pgfqpoint{0.041667in}{0.041667in}}{%
\pgfpathmoveto{\pgfqpoint{0.000000in}{-0.041667in}}%
\pgfpathcurveto{\pgfqpoint{0.011050in}{-0.041667in}}{\pgfqpoint{0.021649in}{-0.037276in}}{\pgfqpoint{0.029463in}{-0.029463in}}%
\pgfpathcurveto{\pgfqpoint{0.037276in}{-0.021649in}}{\pgfqpoint{0.041667in}{-0.011050in}}{\pgfqpoint{0.041667in}{0.000000in}}%
\pgfpathcurveto{\pgfqpoint{0.041667in}{0.011050in}}{\pgfqpoint{0.037276in}{0.021649in}}{\pgfqpoint{0.029463in}{0.029463in}}%
\pgfpathcurveto{\pgfqpoint{0.021649in}{0.037276in}}{\pgfqpoint{0.011050in}{0.041667in}}{\pgfqpoint{0.000000in}{0.041667in}}%
\pgfpathcurveto{\pgfqpoint{-0.011050in}{0.041667in}}{\pgfqpoint{-0.021649in}{0.037276in}}{\pgfqpoint{-0.029463in}{0.029463in}}%
\pgfpathcurveto{\pgfqpoint{-0.037276in}{0.021649in}}{\pgfqpoint{-0.041667in}{0.011050in}}{\pgfqpoint{-0.041667in}{0.000000in}}%
\pgfpathcurveto{\pgfqpoint{-0.041667in}{-0.011050in}}{\pgfqpoint{-0.037276in}{-0.021649in}}{\pgfqpoint{-0.029463in}{-0.029463in}}%
\pgfpathcurveto{\pgfqpoint{-0.021649in}{-0.037276in}}{\pgfqpoint{-0.011050in}{-0.041667in}}{\pgfqpoint{0.000000in}{-0.041667in}}%
\pgfpathlineto{\pgfqpoint{0.000000in}{-0.041667in}}%
\pgfpathclose%
\pgfusepath{stroke,fill}%
}%
\begin{pgfscope}%
\pgfsys@transformshift{0.983334in}{2.756829in}%
\pgfsys@useobject{currentmarker}{}%
\end{pgfscope}%
\end{pgfscope}%
\begin{pgfscope}%
\definecolor{textcolor}{rgb}{0.000000,0.000000,0.000000}%
\pgfsetstrokecolor{textcolor}%
\pgfsetfillcolor{textcolor}%
\pgftext[x=1.233334in,y=2.720371in,left,base]{\color{textcolor}\rmfamily\fontsize{10.000000}{12.000000}\selectfont Resonance}%
\end{pgfscope}%
\end{pgfpicture}%
\makeatother%
\endgroup%

		\caption{Measured return loss and phase of a quartz resonatoras a function of frequency.}
		\label{fig:354}
	\end{figure}
	Figure~\ref{fig:354} shows the measured impedance magnitude $|Z|$ of the quartz resonator versus
	frequency. A distinct minimum of the impedance occurs at the resonance frequency
	$f_0 \approx 9.828227\,\mathrm{MHz}$.
	
	At resonance, the inductive and capacitive reactances of the quartz cancel each other,
	leaving only the small series resistance. As a result, the impedance reaches
	its lowest value. Away from resonance, reactive effects dominate, causing the
	impedance to increase rapidly.
	
	
	\subsubsection{}\label{ch:355}
	The unknown component is a inductance. This is evident from the impedance behavior as seen in fig. \ref{fig:355}. Decreasing at low frequencies and increasing linearly at high frequencies. At the intersection point where $|Z| = Z_0 = 50\,\Omega$, the frequency is $f \approx 1.8 \times 10^7\,\mathrm{Hz}$.
	\begin{figure}[H]
		\centering
		%% Creator: Matplotlib, PGF backend
%%
%% To include the figure in your LaTeX document, write
%%   \input{<filename>.pgf}
%%
%% Make sure the required packages are loaded in your preamble
%%   \usepackage{pgf}
%%
%% Also ensure that all the required font packages are loaded; for instance,
%% the lmodern package is sometimes necessary when using math font.
%%   \usepackage{lmodern}
%%
%% Figures using additional raster images can only be included by \input if
%% they are in the same directory as the main LaTeX file. For loading figures
%% from other directories you can use the `import` package
%%   \usepackage{import}
%%
%% and then include the figures with
%%   \import{<path to file>}{<filename>.pgf}
%%
%% Matplotlib used the following preamble
%%
\begingroup%
\makeatletter%
\begin{pgfpicture}%
\pgfpathrectangle{\pgfpointorigin}{\pgfqpoint{6.300000in}{3.500000in}}%
\pgfusepath{use as bounding box, clip}%
\begin{pgfscope}%
\pgfsetbuttcap%
\pgfsetmiterjoin%
\definecolor{currentfill}{rgb}{1.000000,1.000000,1.000000}%
\pgfsetfillcolor{currentfill}%
\pgfsetlinewidth{0.000000pt}%
\definecolor{currentstroke}{rgb}{1.000000,1.000000,1.000000}%
\pgfsetstrokecolor{currentstroke}%
\pgfsetdash{}{0pt}%
\pgfpathmoveto{\pgfqpoint{0.000000in}{0.000000in}}%
\pgfpathlineto{\pgfqpoint{6.300000in}{0.000000in}}%
\pgfpathlineto{\pgfqpoint{6.300000in}{3.500000in}}%
\pgfpathlineto{\pgfqpoint{0.000000in}{3.500000in}}%
\pgfpathlineto{\pgfqpoint{0.000000in}{0.000000in}}%
\pgfpathclose%
\pgfusepath{fill}%
\end{pgfscope}%
\begin{pgfscope}%
\pgfsetbuttcap%
\pgfsetmiterjoin%
\definecolor{currentfill}{rgb}{1.000000,1.000000,1.000000}%
\pgfsetfillcolor{currentfill}%
\pgfsetlinewidth{0.000000pt}%
\definecolor{currentstroke}{rgb}{0.000000,0.000000,0.000000}%
\pgfsetstrokecolor{currentstroke}%
\pgfsetstrokeopacity{0.000000}%
\pgfsetdash{}{0pt}%
\pgfpathmoveto{\pgfqpoint{0.650001in}{0.565123in}}%
\pgfpathlineto{\pgfqpoint{6.150000in}{0.565123in}}%
\pgfpathlineto{\pgfqpoint{6.150000in}{3.350000in}}%
\pgfpathlineto{\pgfqpoint{0.650001in}{3.350000in}}%
\pgfpathlineto{\pgfqpoint{0.650001in}{0.565123in}}%
\pgfpathclose%
\pgfusepath{fill}%
\end{pgfscope}%
\begin{pgfscope}%
\pgfpathrectangle{\pgfqpoint{0.650001in}{0.565123in}}{\pgfqpoint{5.499999in}{2.784877in}}%
\pgfusepath{clip}%
\pgfsetbuttcap%
\pgfsetroundjoin%
\pgfsetlinewidth{0.803000pt}%
\definecolor{currentstroke}{rgb}{0.690196,0.690196,0.690196}%
\pgfsetstrokecolor{currentstroke}%
\pgfsetdash{{0.800000pt}{1.320000pt}}{0.000000pt}%
\pgfpathmoveto{\pgfqpoint{1.552334in}{0.565123in}}%
\pgfpathlineto{\pgfqpoint{1.552334in}{3.350000in}}%
\pgfusepath{stroke}%
\end{pgfscope}%
\begin{pgfscope}%
\pgfsetbuttcap%
\pgfsetroundjoin%
\definecolor{currentfill}{rgb}{0.000000,0.000000,0.000000}%
\pgfsetfillcolor{currentfill}%
\pgfsetlinewidth{0.803000pt}%
\definecolor{currentstroke}{rgb}{0.000000,0.000000,0.000000}%
\pgfsetstrokecolor{currentstroke}%
\pgfsetdash{}{0pt}%
\pgfsys@defobject{currentmarker}{\pgfqpoint{0.000000in}{-0.048611in}}{\pgfqpoint{0.000000in}{0.000000in}}{%
\pgfpathmoveto{\pgfqpoint{0.000000in}{0.000000in}}%
\pgfpathlineto{\pgfqpoint{0.000000in}{-0.048611in}}%
\pgfusepath{stroke,fill}%
}%
\begin{pgfscope}%
\pgfsys@transformshift{1.552334in}{0.565123in}%
\pgfsys@useobject{currentmarker}{}%
\end{pgfscope}%
\end{pgfscope}%
\begin{pgfscope}%
\definecolor{textcolor}{rgb}{0.000000,0.000000,0.000000}%
\pgfsetstrokecolor{textcolor}%
\pgfsetfillcolor{textcolor}%
\pgftext[x=1.552334in,y=0.467901in,,top]{\color{textcolor}\rmfamily\fontsize{10.000000}{12.000000}\selectfont \(\displaystyle {0.5}\)}%
\end{pgfscope}%
\begin{pgfscope}%
\pgfpathrectangle{\pgfqpoint{0.650001in}{0.565123in}}{\pgfqpoint{5.499999in}{2.784877in}}%
\pgfusepath{clip}%
\pgfsetbuttcap%
\pgfsetroundjoin%
\pgfsetlinewidth{0.803000pt}%
\definecolor{currentstroke}{rgb}{0.690196,0.690196,0.690196}%
\pgfsetstrokecolor{currentstroke}%
\pgfsetdash{{0.800000pt}{1.320000pt}}{0.000000pt}%
\pgfpathmoveto{\pgfqpoint{2.473082in}{0.565123in}}%
\pgfpathlineto{\pgfqpoint{2.473082in}{3.350000in}}%
\pgfusepath{stroke}%
\end{pgfscope}%
\begin{pgfscope}%
\pgfsetbuttcap%
\pgfsetroundjoin%
\definecolor{currentfill}{rgb}{0.000000,0.000000,0.000000}%
\pgfsetfillcolor{currentfill}%
\pgfsetlinewidth{0.803000pt}%
\definecolor{currentstroke}{rgb}{0.000000,0.000000,0.000000}%
\pgfsetstrokecolor{currentstroke}%
\pgfsetdash{}{0pt}%
\pgfsys@defobject{currentmarker}{\pgfqpoint{0.000000in}{-0.048611in}}{\pgfqpoint{0.000000in}{0.000000in}}{%
\pgfpathmoveto{\pgfqpoint{0.000000in}{0.000000in}}%
\pgfpathlineto{\pgfqpoint{0.000000in}{-0.048611in}}%
\pgfusepath{stroke,fill}%
}%
\begin{pgfscope}%
\pgfsys@transformshift{2.473082in}{0.565123in}%
\pgfsys@useobject{currentmarker}{}%
\end{pgfscope}%
\end{pgfscope}%
\begin{pgfscope}%
\definecolor{textcolor}{rgb}{0.000000,0.000000,0.000000}%
\pgfsetstrokecolor{textcolor}%
\pgfsetfillcolor{textcolor}%
\pgftext[x=2.473082in,y=0.467901in,,top]{\color{textcolor}\rmfamily\fontsize{10.000000}{12.000000}\selectfont \(\displaystyle {1.0}\)}%
\end{pgfscope}%
\begin{pgfscope}%
\pgfpathrectangle{\pgfqpoint{0.650001in}{0.565123in}}{\pgfqpoint{5.499999in}{2.784877in}}%
\pgfusepath{clip}%
\pgfsetbuttcap%
\pgfsetroundjoin%
\pgfsetlinewidth{0.803000pt}%
\definecolor{currentstroke}{rgb}{0.690196,0.690196,0.690196}%
\pgfsetstrokecolor{currentstroke}%
\pgfsetdash{{0.800000pt}{1.320000pt}}{0.000000pt}%
\pgfpathmoveto{\pgfqpoint{3.393830in}{0.565123in}}%
\pgfpathlineto{\pgfqpoint{3.393830in}{3.350000in}}%
\pgfusepath{stroke}%
\end{pgfscope}%
\begin{pgfscope}%
\pgfsetbuttcap%
\pgfsetroundjoin%
\definecolor{currentfill}{rgb}{0.000000,0.000000,0.000000}%
\pgfsetfillcolor{currentfill}%
\pgfsetlinewidth{0.803000pt}%
\definecolor{currentstroke}{rgb}{0.000000,0.000000,0.000000}%
\pgfsetstrokecolor{currentstroke}%
\pgfsetdash{}{0pt}%
\pgfsys@defobject{currentmarker}{\pgfqpoint{0.000000in}{-0.048611in}}{\pgfqpoint{0.000000in}{0.000000in}}{%
\pgfpathmoveto{\pgfqpoint{0.000000in}{0.000000in}}%
\pgfpathlineto{\pgfqpoint{0.000000in}{-0.048611in}}%
\pgfusepath{stroke,fill}%
}%
\begin{pgfscope}%
\pgfsys@transformshift{3.393830in}{0.565123in}%
\pgfsys@useobject{currentmarker}{}%
\end{pgfscope}%
\end{pgfscope}%
\begin{pgfscope}%
\definecolor{textcolor}{rgb}{0.000000,0.000000,0.000000}%
\pgfsetstrokecolor{textcolor}%
\pgfsetfillcolor{textcolor}%
\pgftext[x=3.393830in,y=0.467901in,,top]{\color{textcolor}\rmfamily\fontsize{10.000000}{12.000000}\selectfont \(\displaystyle {1.5}\)}%
\end{pgfscope}%
\begin{pgfscope}%
\pgfpathrectangle{\pgfqpoint{0.650001in}{0.565123in}}{\pgfqpoint{5.499999in}{2.784877in}}%
\pgfusepath{clip}%
\pgfsetbuttcap%
\pgfsetroundjoin%
\pgfsetlinewidth{0.803000pt}%
\definecolor{currentstroke}{rgb}{0.690196,0.690196,0.690196}%
\pgfsetstrokecolor{currentstroke}%
\pgfsetdash{{0.800000pt}{1.320000pt}}{0.000000pt}%
\pgfpathmoveto{\pgfqpoint{4.314578in}{0.565123in}}%
\pgfpathlineto{\pgfqpoint{4.314578in}{3.350000in}}%
\pgfusepath{stroke}%
\end{pgfscope}%
\begin{pgfscope}%
\pgfsetbuttcap%
\pgfsetroundjoin%
\definecolor{currentfill}{rgb}{0.000000,0.000000,0.000000}%
\pgfsetfillcolor{currentfill}%
\pgfsetlinewidth{0.803000pt}%
\definecolor{currentstroke}{rgb}{0.000000,0.000000,0.000000}%
\pgfsetstrokecolor{currentstroke}%
\pgfsetdash{}{0pt}%
\pgfsys@defobject{currentmarker}{\pgfqpoint{0.000000in}{-0.048611in}}{\pgfqpoint{0.000000in}{0.000000in}}{%
\pgfpathmoveto{\pgfqpoint{0.000000in}{0.000000in}}%
\pgfpathlineto{\pgfqpoint{0.000000in}{-0.048611in}}%
\pgfusepath{stroke,fill}%
}%
\begin{pgfscope}%
\pgfsys@transformshift{4.314578in}{0.565123in}%
\pgfsys@useobject{currentmarker}{}%
\end{pgfscope}%
\end{pgfscope}%
\begin{pgfscope}%
\definecolor{textcolor}{rgb}{0.000000,0.000000,0.000000}%
\pgfsetstrokecolor{textcolor}%
\pgfsetfillcolor{textcolor}%
\pgftext[x=4.314578in,y=0.467901in,,top]{\color{textcolor}\rmfamily\fontsize{10.000000}{12.000000}\selectfont \(\displaystyle {2.0}\)}%
\end{pgfscope}%
\begin{pgfscope}%
\pgfpathrectangle{\pgfqpoint{0.650001in}{0.565123in}}{\pgfqpoint{5.499999in}{2.784877in}}%
\pgfusepath{clip}%
\pgfsetbuttcap%
\pgfsetroundjoin%
\pgfsetlinewidth{0.803000pt}%
\definecolor{currentstroke}{rgb}{0.690196,0.690196,0.690196}%
\pgfsetstrokecolor{currentstroke}%
\pgfsetdash{{0.800000pt}{1.320000pt}}{0.000000pt}%
\pgfpathmoveto{\pgfqpoint{5.235326in}{0.565123in}}%
\pgfpathlineto{\pgfqpoint{5.235326in}{3.350000in}}%
\pgfusepath{stroke}%
\end{pgfscope}%
\begin{pgfscope}%
\pgfsetbuttcap%
\pgfsetroundjoin%
\definecolor{currentfill}{rgb}{0.000000,0.000000,0.000000}%
\pgfsetfillcolor{currentfill}%
\pgfsetlinewidth{0.803000pt}%
\definecolor{currentstroke}{rgb}{0.000000,0.000000,0.000000}%
\pgfsetstrokecolor{currentstroke}%
\pgfsetdash{}{0pt}%
\pgfsys@defobject{currentmarker}{\pgfqpoint{0.000000in}{-0.048611in}}{\pgfqpoint{0.000000in}{0.000000in}}{%
\pgfpathmoveto{\pgfqpoint{0.000000in}{0.000000in}}%
\pgfpathlineto{\pgfqpoint{0.000000in}{-0.048611in}}%
\pgfusepath{stroke,fill}%
}%
\begin{pgfscope}%
\pgfsys@transformshift{5.235326in}{0.565123in}%
\pgfsys@useobject{currentmarker}{}%
\end{pgfscope}%
\end{pgfscope}%
\begin{pgfscope}%
\definecolor{textcolor}{rgb}{0.000000,0.000000,0.000000}%
\pgfsetstrokecolor{textcolor}%
\pgfsetfillcolor{textcolor}%
\pgftext[x=5.235326in,y=0.467901in,,top]{\color{textcolor}\rmfamily\fontsize{10.000000}{12.000000}\selectfont \(\displaystyle {2.5}\)}%
\end{pgfscope}%
\begin{pgfscope}%
\definecolor{textcolor}{rgb}{0.000000,0.000000,0.000000}%
\pgfsetstrokecolor{textcolor}%
\pgfsetfillcolor{textcolor}%
\pgftext[x=3.400000in,y=0.288889in,,top]{\color{textcolor}\rmfamily\fontsize{10.000000}{12.000000}\selectfont Frequency (Hz)}%
\end{pgfscope}%
\begin{pgfscope}%
\definecolor{textcolor}{rgb}{0.000000,0.000000,0.000000}%
\pgfsetstrokecolor{textcolor}%
\pgfsetfillcolor{textcolor}%
\pgftext[x=6.150000in,y=0.302778in,right,top]{\color{textcolor}\rmfamily\fontsize{10.000000}{12.000000}\selectfont \(\displaystyle \times{10^{7}}{}\)}%
\end{pgfscope}%
\begin{pgfscope}%
\pgfpathrectangle{\pgfqpoint{0.650001in}{0.565123in}}{\pgfqpoint{5.499999in}{2.784877in}}%
\pgfusepath{clip}%
\pgfsetbuttcap%
\pgfsetroundjoin%
\pgfsetlinewidth{0.803000pt}%
\definecolor{currentstroke}{rgb}{0.690196,0.690196,0.690196}%
\pgfsetstrokecolor{currentstroke}%
\pgfsetdash{{0.800000pt}{1.320000pt}}{0.000000pt}%
\pgfpathmoveto{\pgfqpoint{0.650001in}{0.602968in}}%
\pgfpathlineto{\pgfqpoint{6.150000in}{0.602968in}}%
\pgfusepath{stroke}%
\end{pgfscope}%
\begin{pgfscope}%
\pgfsetbuttcap%
\pgfsetroundjoin%
\definecolor{currentfill}{rgb}{0.000000,0.000000,0.000000}%
\pgfsetfillcolor{currentfill}%
\pgfsetlinewidth{0.803000pt}%
\definecolor{currentstroke}{rgb}{0.000000,0.000000,0.000000}%
\pgfsetstrokecolor{currentstroke}%
\pgfsetdash{}{0pt}%
\pgfsys@defobject{currentmarker}{\pgfqpoint{-0.048611in}{0.000000in}}{\pgfqpoint{-0.000000in}{0.000000in}}{%
\pgfpathmoveto{\pgfqpoint{-0.000000in}{0.000000in}}%
\pgfpathlineto{\pgfqpoint{-0.048611in}{0.000000in}}%
\pgfusepath{stroke,fill}%
}%
\begin{pgfscope}%
\pgfsys@transformshift{0.650001in}{0.602968in}%
\pgfsys@useobject{currentmarker}{}%
\end{pgfscope}%
\end{pgfscope}%
\begin{pgfscope}%
\definecolor{textcolor}{rgb}{0.000000,0.000000,0.000000}%
\pgfsetstrokecolor{textcolor}%
\pgfsetfillcolor{textcolor}%
\pgftext[x=0.483334in, y=0.554743in, left, base]{\color{textcolor}\rmfamily\fontsize{10.000000}{12.000000}\selectfont \(\displaystyle {0}\)}%
\end{pgfscope}%
\begin{pgfscope}%
\pgfpathrectangle{\pgfqpoint{0.650001in}{0.565123in}}{\pgfqpoint{5.499999in}{2.784877in}}%
\pgfusepath{clip}%
\pgfsetbuttcap%
\pgfsetroundjoin%
\pgfsetlinewidth{0.803000pt}%
\definecolor{currentstroke}{rgb}{0.690196,0.690196,0.690196}%
\pgfsetstrokecolor{currentstroke}%
\pgfsetdash{{0.800000pt}{1.320000pt}}{0.000000pt}%
\pgfpathmoveto{\pgfqpoint{0.650001in}{1.124970in}}%
\pgfpathlineto{\pgfqpoint{6.150000in}{1.124970in}}%
\pgfusepath{stroke}%
\end{pgfscope}%
\begin{pgfscope}%
\pgfsetbuttcap%
\pgfsetroundjoin%
\definecolor{currentfill}{rgb}{0.000000,0.000000,0.000000}%
\pgfsetfillcolor{currentfill}%
\pgfsetlinewidth{0.803000pt}%
\definecolor{currentstroke}{rgb}{0.000000,0.000000,0.000000}%
\pgfsetstrokecolor{currentstroke}%
\pgfsetdash{}{0pt}%
\pgfsys@defobject{currentmarker}{\pgfqpoint{-0.048611in}{0.000000in}}{\pgfqpoint{-0.000000in}{0.000000in}}{%
\pgfpathmoveto{\pgfqpoint{-0.000000in}{0.000000in}}%
\pgfpathlineto{\pgfqpoint{-0.048611in}{0.000000in}}%
\pgfusepath{stroke,fill}%
}%
\begin{pgfscope}%
\pgfsys@transformshift{0.650001in}{1.124970in}%
\pgfsys@useobject{currentmarker}{}%
\end{pgfscope}%
\end{pgfscope}%
\begin{pgfscope}%
\definecolor{textcolor}{rgb}{0.000000,0.000000,0.000000}%
\pgfsetstrokecolor{textcolor}%
\pgfsetfillcolor{textcolor}%
\pgftext[x=0.413889in, y=1.076744in, left, base]{\color{textcolor}\rmfamily\fontsize{10.000000}{12.000000}\selectfont \(\displaystyle {20}\)}%
\end{pgfscope}%
\begin{pgfscope}%
\pgfpathrectangle{\pgfqpoint{0.650001in}{0.565123in}}{\pgfqpoint{5.499999in}{2.784877in}}%
\pgfusepath{clip}%
\pgfsetbuttcap%
\pgfsetroundjoin%
\pgfsetlinewidth{0.803000pt}%
\definecolor{currentstroke}{rgb}{0.690196,0.690196,0.690196}%
\pgfsetstrokecolor{currentstroke}%
\pgfsetdash{{0.800000pt}{1.320000pt}}{0.000000pt}%
\pgfpathmoveto{\pgfqpoint{0.650001in}{1.646971in}}%
\pgfpathlineto{\pgfqpoint{6.150000in}{1.646971in}}%
\pgfusepath{stroke}%
\end{pgfscope}%
\begin{pgfscope}%
\pgfsetbuttcap%
\pgfsetroundjoin%
\definecolor{currentfill}{rgb}{0.000000,0.000000,0.000000}%
\pgfsetfillcolor{currentfill}%
\pgfsetlinewidth{0.803000pt}%
\definecolor{currentstroke}{rgb}{0.000000,0.000000,0.000000}%
\pgfsetstrokecolor{currentstroke}%
\pgfsetdash{}{0pt}%
\pgfsys@defobject{currentmarker}{\pgfqpoint{-0.048611in}{0.000000in}}{\pgfqpoint{-0.000000in}{0.000000in}}{%
\pgfpathmoveto{\pgfqpoint{-0.000000in}{0.000000in}}%
\pgfpathlineto{\pgfqpoint{-0.048611in}{0.000000in}}%
\pgfusepath{stroke,fill}%
}%
\begin{pgfscope}%
\pgfsys@transformshift{0.650001in}{1.646971in}%
\pgfsys@useobject{currentmarker}{}%
\end{pgfscope}%
\end{pgfscope}%
\begin{pgfscope}%
\definecolor{textcolor}{rgb}{0.000000,0.000000,0.000000}%
\pgfsetstrokecolor{textcolor}%
\pgfsetfillcolor{textcolor}%
\pgftext[x=0.413889in, y=1.598746in, left, base]{\color{textcolor}\rmfamily\fontsize{10.000000}{12.000000}\selectfont \(\displaystyle {40}\)}%
\end{pgfscope}%
\begin{pgfscope}%
\pgfpathrectangle{\pgfqpoint{0.650001in}{0.565123in}}{\pgfqpoint{5.499999in}{2.784877in}}%
\pgfusepath{clip}%
\pgfsetbuttcap%
\pgfsetroundjoin%
\pgfsetlinewidth{0.803000pt}%
\definecolor{currentstroke}{rgb}{0.690196,0.690196,0.690196}%
\pgfsetstrokecolor{currentstroke}%
\pgfsetdash{{0.800000pt}{1.320000pt}}{0.000000pt}%
\pgfpathmoveto{\pgfqpoint{0.650001in}{2.168972in}}%
\pgfpathlineto{\pgfqpoint{6.150000in}{2.168972in}}%
\pgfusepath{stroke}%
\end{pgfscope}%
\begin{pgfscope}%
\pgfsetbuttcap%
\pgfsetroundjoin%
\definecolor{currentfill}{rgb}{0.000000,0.000000,0.000000}%
\pgfsetfillcolor{currentfill}%
\pgfsetlinewidth{0.803000pt}%
\definecolor{currentstroke}{rgb}{0.000000,0.000000,0.000000}%
\pgfsetstrokecolor{currentstroke}%
\pgfsetdash{}{0pt}%
\pgfsys@defobject{currentmarker}{\pgfqpoint{-0.048611in}{0.000000in}}{\pgfqpoint{-0.000000in}{0.000000in}}{%
\pgfpathmoveto{\pgfqpoint{-0.000000in}{0.000000in}}%
\pgfpathlineto{\pgfqpoint{-0.048611in}{0.000000in}}%
\pgfusepath{stroke,fill}%
}%
\begin{pgfscope}%
\pgfsys@transformshift{0.650001in}{2.168972in}%
\pgfsys@useobject{currentmarker}{}%
\end{pgfscope}%
\end{pgfscope}%
\begin{pgfscope}%
\definecolor{textcolor}{rgb}{0.000000,0.000000,0.000000}%
\pgfsetstrokecolor{textcolor}%
\pgfsetfillcolor{textcolor}%
\pgftext[x=0.413889in, y=2.120747in, left, base]{\color{textcolor}\rmfamily\fontsize{10.000000}{12.000000}\selectfont \(\displaystyle {60}\)}%
\end{pgfscope}%
\begin{pgfscope}%
\pgfpathrectangle{\pgfqpoint{0.650001in}{0.565123in}}{\pgfqpoint{5.499999in}{2.784877in}}%
\pgfusepath{clip}%
\pgfsetbuttcap%
\pgfsetroundjoin%
\pgfsetlinewidth{0.803000pt}%
\definecolor{currentstroke}{rgb}{0.690196,0.690196,0.690196}%
\pgfsetstrokecolor{currentstroke}%
\pgfsetdash{{0.800000pt}{1.320000pt}}{0.000000pt}%
\pgfpathmoveto{\pgfqpoint{0.650001in}{2.690973in}}%
\pgfpathlineto{\pgfqpoint{6.150000in}{2.690973in}}%
\pgfusepath{stroke}%
\end{pgfscope}%
\begin{pgfscope}%
\pgfsetbuttcap%
\pgfsetroundjoin%
\definecolor{currentfill}{rgb}{0.000000,0.000000,0.000000}%
\pgfsetfillcolor{currentfill}%
\pgfsetlinewidth{0.803000pt}%
\definecolor{currentstroke}{rgb}{0.000000,0.000000,0.000000}%
\pgfsetstrokecolor{currentstroke}%
\pgfsetdash{}{0pt}%
\pgfsys@defobject{currentmarker}{\pgfqpoint{-0.048611in}{0.000000in}}{\pgfqpoint{-0.000000in}{0.000000in}}{%
\pgfpathmoveto{\pgfqpoint{-0.000000in}{0.000000in}}%
\pgfpathlineto{\pgfqpoint{-0.048611in}{0.000000in}}%
\pgfusepath{stroke,fill}%
}%
\begin{pgfscope}%
\pgfsys@transformshift{0.650001in}{2.690973in}%
\pgfsys@useobject{currentmarker}{}%
\end{pgfscope}%
\end{pgfscope}%
\begin{pgfscope}%
\definecolor{textcolor}{rgb}{0.000000,0.000000,0.000000}%
\pgfsetstrokecolor{textcolor}%
\pgfsetfillcolor{textcolor}%
\pgftext[x=0.413889in, y=2.642748in, left, base]{\color{textcolor}\rmfamily\fontsize{10.000000}{12.000000}\selectfont \(\displaystyle {80}\)}%
\end{pgfscope}%
\begin{pgfscope}%
\pgfpathrectangle{\pgfqpoint{0.650001in}{0.565123in}}{\pgfqpoint{5.499999in}{2.784877in}}%
\pgfusepath{clip}%
\pgfsetbuttcap%
\pgfsetroundjoin%
\pgfsetlinewidth{0.803000pt}%
\definecolor{currentstroke}{rgb}{0.690196,0.690196,0.690196}%
\pgfsetstrokecolor{currentstroke}%
\pgfsetdash{{0.800000pt}{1.320000pt}}{0.000000pt}%
\pgfpathmoveto{\pgfqpoint{0.650001in}{3.212975in}}%
\pgfpathlineto{\pgfqpoint{6.150000in}{3.212975in}}%
\pgfusepath{stroke}%
\end{pgfscope}%
\begin{pgfscope}%
\pgfsetbuttcap%
\pgfsetroundjoin%
\definecolor{currentfill}{rgb}{0.000000,0.000000,0.000000}%
\pgfsetfillcolor{currentfill}%
\pgfsetlinewidth{0.803000pt}%
\definecolor{currentstroke}{rgb}{0.000000,0.000000,0.000000}%
\pgfsetstrokecolor{currentstroke}%
\pgfsetdash{}{0pt}%
\pgfsys@defobject{currentmarker}{\pgfqpoint{-0.048611in}{0.000000in}}{\pgfqpoint{-0.000000in}{0.000000in}}{%
\pgfpathmoveto{\pgfqpoint{-0.000000in}{0.000000in}}%
\pgfpathlineto{\pgfqpoint{-0.048611in}{0.000000in}}%
\pgfusepath{stroke,fill}%
}%
\begin{pgfscope}%
\pgfsys@transformshift{0.650001in}{3.212975in}%
\pgfsys@useobject{currentmarker}{}%
\end{pgfscope}%
\end{pgfscope}%
\begin{pgfscope}%
\definecolor{textcolor}{rgb}{0.000000,0.000000,0.000000}%
\pgfsetstrokecolor{textcolor}%
\pgfsetfillcolor{textcolor}%
\pgftext[x=0.344444in, y=3.164749in, left, base]{\color{textcolor}\rmfamily\fontsize{10.000000}{12.000000}\selectfont \(\displaystyle {100}\)}%
\end{pgfscope}%
\begin{pgfscope}%
\definecolor{textcolor}{rgb}{0.000000,0.000000,0.000000}%
\pgfsetstrokecolor{textcolor}%
\pgfsetfillcolor{textcolor}%
\pgftext[x=0.288889in,y=1.957562in,,bottom,rotate=90.000000]{\color{textcolor}\rmfamily\fontsize{10.000000}{12.000000}\selectfont Impedance |Z| (\(\displaystyle \Omega\))}%
\end{pgfscope}%
\begin{pgfscope}%
\pgfpathrectangle{\pgfqpoint{0.650001in}{0.565123in}}{\pgfqpoint{5.499999in}{2.784877in}}%
\pgfusepath{clip}%
\pgfsetrectcap%
\pgfsetroundjoin%
\pgfsetlinewidth{1.003750pt}%
\definecolor{currentstroke}{rgb}{0.000000,0.000000,0.000000}%
\pgfsetstrokecolor{currentstroke}%
\pgfsetdash{}{0pt}%
\pgfpathmoveto{\pgfqpoint{0.650001in}{1.357260in}}%
\pgfpathlineto{\pgfqpoint{0.655934in}{0.738689in}}%
\pgfpathlineto{\pgfqpoint{0.661867in}{0.730859in}}%
\pgfpathlineto{\pgfqpoint{0.673733in}{0.720419in}}%
\pgfpathlineto{\pgfqpoint{0.679666in}{0.720419in}}%
\pgfpathlineto{\pgfqpoint{0.685599in}{0.717809in}}%
\pgfpathlineto{\pgfqpoint{0.691533in}{0.707369in}}%
\pgfpathlineto{\pgfqpoint{0.697466in}{0.707369in}}%
\pgfpathlineto{\pgfqpoint{0.703399in}{0.704759in}}%
\pgfpathlineto{\pgfqpoint{0.721198in}{0.704759in}}%
\pgfpathlineto{\pgfqpoint{0.727131in}{0.699539in}}%
\pgfpathlineto{\pgfqpoint{0.744931in}{0.699539in}}%
\pgfpathlineto{\pgfqpoint{0.750864in}{0.702149in}}%
\pgfpathlineto{\pgfqpoint{0.756797in}{0.696929in}}%
\pgfpathlineto{\pgfqpoint{0.762730in}{0.696929in}}%
\pgfpathlineto{\pgfqpoint{0.768663in}{0.694319in}}%
\pgfpathlineto{\pgfqpoint{0.774596in}{0.699539in}}%
\pgfpathlineto{\pgfqpoint{0.780529in}{0.694319in}}%
\pgfpathlineto{\pgfqpoint{0.786462in}{0.694319in}}%
\pgfpathlineto{\pgfqpoint{0.792395in}{0.691709in}}%
\pgfpathlineto{\pgfqpoint{0.798329in}{0.694319in}}%
\pgfpathlineto{\pgfqpoint{0.810195in}{0.694319in}}%
\pgfpathlineto{\pgfqpoint{0.816128in}{0.696929in}}%
\pgfpathlineto{\pgfqpoint{0.827994in}{0.696929in}}%
\pgfpathlineto{\pgfqpoint{0.833927in}{0.699539in}}%
\pgfpathlineto{\pgfqpoint{0.839860in}{0.699539in}}%
\pgfpathlineto{\pgfqpoint{0.845794in}{0.696929in}}%
\pgfpathlineto{\pgfqpoint{0.851727in}{0.699539in}}%
\pgfpathlineto{\pgfqpoint{0.857660in}{0.699539in}}%
\pgfpathlineto{\pgfqpoint{0.869526in}{0.704759in}}%
\pgfpathlineto{\pgfqpoint{0.875459in}{0.704759in}}%
\pgfpathlineto{\pgfqpoint{0.881392in}{0.707369in}}%
\pgfpathlineto{\pgfqpoint{0.887325in}{0.707369in}}%
\pgfpathlineto{\pgfqpoint{0.893258in}{0.704759in}}%
\pgfpathlineto{\pgfqpoint{0.905125in}{0.709979in}}%
\pgfpathlineto{\pgfqpoint{0.911058in}{0.709979in}}%
\pgfpathlineto{\pgfqpoint{0.928857in}{0.717809in}}%
\pgfpathlineto{\pgfqpoint{0.934790in}{0.715199in}}%
\pgfpathlineto{\pgfqpoint{0.940723in}{0.717809in}}%
\pgfpathlineto{\pgfqpoint{0.946657in}{0.717809in}}%
\pgfpathlineto{\pgfqpoint{0.970389in}{0.728249in}}%
\pgfpathlineto{\pgfqpoint{0.976322in}{0.728249in}}%
\pgfpathlineto{\pgfqpoint{1.000055in}{0.738689in}}%
\pgfpathlineto{\pgfqpoint{1.011921in}{0.738689in}}%
\pgfpathlineto{\pgfqpoint{1.029720in}{0.746519in}}%
\pgfpathlineto{\pgfqpoint{1.035653in}{0.746519in}}%
\pgfpathlineto{\pgfqpoint{1.047520in}{0.751739in}}%
\pgfpathlineto{\pgfqpoint{1.053453in}{0.751739in}}%
\pgfpathlineto{\pgfqpoint{1.059386in}{0.756959in}}%
\pgfpathlineto{\pgfqpoint{1.065319in}{0.759569in}}%
\pgfpathlineto{\pgfqpoint{1.071252in}{0.759569in}}%
\pgfpathlineto{\pgfqpoint{1.100918in}{0.772619in}}%
\pgfpathlineto{\pgfqpoint{1.106851in}{0.772619in}}%
\pgfpathlineto{\pgfqpoint{1.124650in}{0.780449in}}%
\pgfpathlineto{\pgfqpoint{1.130583in}{0.780449in}}%
\pgfpathlineto{\pgfqpoint{1.142449in}{0.785669in}}%
\pgfpathlineto{\pgfqpoint{1.148383in}{0.785669in}}%
\pgfpathlineto{\pgfqpoint{1.154316in}{0.788279in}}%
\pgfpathlineto{\pgfqpoint{1.160249in}{0.788279in}}%
\pgfpathlineto{\pgfqpoint{1.166182in}{0.793499in}}%
\pgfpathlineto{\pgfqpoint{1.172115in}{0.796109in}}%
\pgfpathlineto{\pgfqpoint{1.178048in}{0.796109in}}%
\pgfpathlineto{\pgfqpoint{1.195847in}{0.803939in}}%
\pgfpathlineto{\pgfqpoint{1.201781in}{0.803939in}}%
\pgfpathlineto{\pgfqpoint{1.207714in}{0.809159in}}%
\pgfpathlineto{\pgfqpoint{1.213647in}{0.811769in}}%
\pgfpathlineto{\pgfqpoint{1.219580in}{0.811769in}}%
\pgfpathlineto{\pgfqpoint{1.261112in}{0.830039in}}%
\pgfpathlineto{\pgfqpoint{1.267045in}{0.830039in}}%
\pgfpathlineto{\pgfqpoint{1.290777in}{0.840479in}}%
\pgfpathlineto{\pgfqpoint{1.296710in}{0.840479in}}%
\pgfpathlineto{\pgfqpoint{1.338242in}{0.858749in}}%
\pgfpathlineto{\pgfqpoint{1.344175in}{0.858749in}}%
\pgfpathlineto{\pgfqpoint{1.391640in}{0.879629in}}%
\pgfpathlineto{\pgfqpoint{1.397573in}{0.879629in}}%
\pgfpathlineto{\pgfqpoint{1.403507in}{0.884849in}}%
\pgfpathlineto{\pgfqpoint{1.409440in}{0.887459in}}%
\pgfpathlineto{\pgfqpoint{1.415373in}{0.887459in}}%
\pgfpathlineto{\pgfqpoint{1.439105in}{0.897899in}}%
\pgfpathlineto{\pgfqpoint{1.445038in}{0.897899in}}%
\pgfpathlineto{\pgfqpoint{1.450971in}{0.900509in}}%
\pgfpathlineto{\pgfqpoint{1.456905in}{0.905729in}}%
\pgfpathlineto{\pgfqpoint{1.462838in}{0.905729in}}%
\pgfpathlineto{\pgfqpoint{1.480637in}{0.913559in}}%
\pgfpathlineto{\pgfqpoint{1.486570in}{0.913559in}}%
\pgfpathlineto{\pgfqpoint{1.492503in}{0.916169in}}%
\pgfpathlineto{\pgfqpoint{1.498436in}{0.921389in}}%
\pgfpathlineto{\pgfqpoint{1.510303in}{0.926609in}}%
\pgfpathlineto{\pgfqpoint{1.516236in}{0.926609in}}%
\pgfpathlineto{\pgfqpoint{1.539968in}{0.937049in}}%
\pgfpathlineto{\pgfqpoint{1.545901in}{0.937049in}}%
\pgfpathlineto{\pgfqpoint{1.551834in}{0.939659in}}%
\pgfpathlineto{\pgfqpoint{1.557768in}{0.944879in}}%
\pgfpathlineto{\pgfqpoint{1.563701in}{0.947489in}}%
\pgfpathlineto{\pgfqpoint{1.569634in}{0.947489in}}%
\pgfpathlineto{\pgfqpoint{1.593366in}{0.957929in}}%
\pgfpathlineto{\pgfqpoint{1.599299in}{0.957929in}}%
\pgfpathlineto{\pgfqpoint{1.605233in}{0.960539in}}%
\pgfpathlineto{\pgfqpoint{1.611166in}{0.965759in}}%
\pgfpathlineto{\pgfqpoint{1.623032in}{0.970979in}}%
\pgfpathlineto{\pgfqpoint{1.628965in}{0.970979in}}%
\pgfpathlineto{\pgfqpoint{1.646764in}{0.978809in}}%
\pgfpathlineto{\pgfqpoint{1.652697in}{0.978809in}}%
\pgfpathlineto{\pgfqpoint{1.664564in}{0.984029in}}%
\pgfpathlineto{\pgfqpoint{1.670497in}{0.989249in}}%
\pgfpathlineto{\pgfqpoint{1.676430in}{0.991859in}}%
\pgfpathlineto{\pgfqpoint{1.682363in}{0.991859in}}%
\pgfpathlineto{\pgfqpoint{1.712029in}{1.004909in}}%
\pgfpathlineto{\pgfqpoint{1.717962in}{1.004909in}}%
\pgfpathlineto{\pgfqpoint{1.723895in}{1.010129in}}%
\pgfpathlineto{\pgfqpoint{1.729828in}{1.010129in}}%
\pgfpathlineto{\pgfqpoint{1.735761in}{1.015349in}}%
\pgfpathlineto{\pgfqpoint{1.741694in}{1.017959in}}%
\pgfpathlineto{\pgfqpoint{1.747627in}{1.017959in}}%
\pgfpathlineto{\pgfqpoint{1.783226in}{1.033619in}}%
\pgfpathlineto{\pgfqpoint{1.789159in}{1.033619in}}%
\pgfpathlineto{\pgfqpoint{1.848490in}{1.059720in}}%
\pgfpathlineto{\pgfqpoint{1.854423in}{1.059720in}}%
\pgfpathlineto{\pgfqpoint{1.890022in}{1.075380in}}%
\pgfpathlineto{\pgfqpoint{1.895955in}{1.075380in}}%
\pgfpathlineto{\pgfqpoint{1.925621in}{1.088430in}}%
\pgfpathlineto{\pgfqpoint{1.931554in}{1.093650in}}%
\pgfpathlineto{\pgfqpoint{1.937487in}{1.093650in}}%
\pgfpathlineto{\pgfqpoint{1.979019in}{1.111920in}}%
\pgfpathlineto{\pgfqpoint{1.984952in}{1.111920in}}%
\pgfpathlineto{\pgfqpoint{2.026484in}{1.130190in}}%
\pgfpathlineto{\pgfqpoint{2.032417in}{1.130190in}}%
\pgfpathlineto{\pgfqpoint{2.079882in}{1.151070in}}%
\pgfpathlineto{\pgfqpoint{2.085815in}{1.151070in}}%
\pgfpathlineto{\pgfqpoint{2.145146in}{1.177170in}}%
\pgfpathlineto{\pgfqpoint{2.151079in}{1.177170in}}%
\pgfpathlineto{\pgfqpoint{2.198544in}{1.198050in}}%
\pgfpathlineto{\pgfqpoint{2.204477in}{1.203270in}}%
\pgfpathlineto{\pgfqpoint{2.216344in}{1.203270in}}%
\pgfpathlineto{\pgfqpoint{2.228210in}{1.208490in}}%
\pgfpathlineto{\pgfqpoint{2.234143in}{1.213710in}}%
\pgfpathlineto{\pgfqpoint{2.240076in}{1.216320in}}%
\pgfpathlineto{\pgfqpoint{2.246009in}{1.216320in}}%
\pgfpathlineto{\pgfqpoint{2.311273in}{1.245030in}}%
\pgfpathlineto{\pgfqpoint{2.317207in}{1.245030in}}%
\pgfpathlineto{\pgfqpoint{2.418070in}{1.289400in}}%
\pgfpathlineto{\pgfqpoint{2.424003in}{1.289400in}}%
\pgfpathlineto{\pgfqpoint{2.429936in}{1.294620in}}%
\pgfpathlineto{\pgfqpoint{2.441802in}{1.299840in}}%
\pgfpathlineto{\pgfqpoint{2.447735in}{1.299840in}}%
\pgfpathlineto{\pgfqpoint{2.536732in}{1.338990in}}%
\pgfpathlineto{\pgfqpoint{2.542665in}{1.338990in}}%
\pgfpathlineto{\pgfqpoint{2.607929in}{1.367700in}}%
\pgfpathlineto{\pgfqpoint{2.613862in}{1.367700in}}%
\pgfpathlineto{\pgfqpoint{2.673194in}{1.393800in}}%
\pgfpathlineto{\pgfqpoint{2.679127in}{1.393800in}}%
\pgfpathlineto{\pgfqpoint{2.696926in}{1.401630in}}%
\pgfpathlineto{\pgfqpoint{2.702859in}{1.406850in}}%
\pgfpathlineto{\pgfqpoint{2.738458in}{1.422510in}}%
\pgfpathlineto{\pgfqpoint{2.744391in}{1.422510in}}%
\pgfpathlineto{\pgfqpoint{2.785923in}{1.440780in}}%
\pgfpathlineto{\pgfqpoint{2.791856in}{1.440780in}}%
\pgfpathlineto{\pgfqpoint{2.797789in}{1.446000in}}%
\pgfpathlineto{\pgfqpoint{2.833388in}{1.461660in}}%
\pgfpathlineto{\pgfqpoint{2.839321in}{1.461660in}}%
\pgfpathlineto{\pgfqpoint{2.857120in}{1.469490in}}%
\pgfpathlineto{\pgfqpoint{2.863053in}{1.474711in}}%
\pgfpathlineto{\pgfqpoint{2.874920in}{1.479931in}}%
\pgfpathlineto{\pgfqpoint{2.880853in}{1.479931in}}%
\pgfpathlineto{\pgfqpoint{2.916451in}{1.495591in}}%
\pgfpathlineto{\pgfqpoint{2.922384in}{1.500811in}}%
\pgfpathlineto{\pgfqpoint{2.928318in}{1.500811in}}%
\pgfpathlineto{\pgfqpoint{2.957983in}{1.513861in}}%
\pgfpathlineto{\pgfqpoint{2.963916in}{1.513861in}}%
\pgfpathlineto{\pgfqpoint{2.969849in}{1.516471in}}%
\pgfpathlineto{\pgfqpoint{2.975782in}{1.521691in}}%
\pgfpathlineto{\pgfqpoint{2.999515in}{1.532131in}}%
\pgfpathlineto{\pgfqpoint{3.005448in}{1.532131in}}%
\pgfpathlineto{\pgfqpoint{3.017314in}{1.537351in}}%
\pgfpathlineto{\pgfqpoint{3.023247in}{1.542571in}}%
\pgfpathlineto{\pgfqpoint{3.035114in}{1.547791in}}%
\pgfpathlineto{\pgfqpoint{3.041047in}{1.547791in}}%
\pgfpathlineto{\pgfqpoint{3.058846in}{1.555621in}}%
\pgfpathlineto{\pgfqpoint{3.064779in}{1.560841in}}%
\pgfpathlineto{\pgfqpoint{3.070712in}{1.563451in}}%
\pgfpathlineto{\pgfqpoint{3.076645in}{1.563451in}}%
\pgfpathlineto{\pgfqpoint{3.094445in}{1.571281in}}%
\pgfpathlineto{\pgfqpoint{3.100378in}{1.576501in}}%
\pgfpathlineto{\pgfqpoint{3.106311in}{1.576501in}}%
\pgfpathlineto{\pgfqpoint{3.130044in}{1.586941in}}%
\pgfpathlineto{\pgfqpoint{3.135977in}{1.592161in}}%
\pgfpathlineto{\pgfqpoint{3.141910in}{1.592161in}}%
\pgfpathlineto{\pgfqpoint{3.165642in}{1.602601in}}%
\pgfpathlineto{\pgfqpoint{3.171575in}{1.602601in}}%
\pgfpathlineto{\pgfqpoint{3.177508in}{1.607821in}}%
\pgfpathlineto{\pgfqpoint{3.195308in}{1.615651in}}%
\pgfpathlineto{\pgfqpoint{3.201241in}{1.615651in}}%
\pgfpathlineto{\pgfqpoint{3.207174in}{1.620871in}}%
\pgfpathlineto{\pgfqpoint{3.224973in}{1.628701in}}%
\pgfpathlineto{\pgfqpoint{3.230907in}{1.628701in}}%
\pgfpathlineto{\pgfqpoint{3.236840in}{1.633921in}}%
\pgfpathlineto{\pgfqpoint{3.254639in}{1.641751in}}%
\pgfpathlineto{\pgfqpoint{3.260572in}{1.641751in}}%
\pgfpathlineto{\pgfqpoint{3.266505in}{1.644361in}}%
\pgfpathlineto{\pgfqpoint{3.272438in}{1.649581in}}%
\pgfpathlineto{\pgfqpoint{3.284305in}{1.654801in}}%
\pgfpathlineto{\pgfqpoint{3.290238in}{1.654801in}}%
\pgfpathlineto{\pgfqpoint{3.296171in}{1.660021in}}%
\pgfpathlineto{\pgfqpoint{3.313970in}{1.667851in}}%
\pgfpathlineto{\pgfqpoint{3.319903in}{1.667851in}}%
\pgfpathlineto{\pgfqpoint{3.325836in}{1.673071in}}%
\pgfpathlineto{\pgfqpoint{3.337703in}{1.678291in}}%
\pgfpathlineto{\pgfqpoint{3.343636in}{1.678291in}}%
\pgfpathlineto{\pgfqpoint{3.349569in}{1.683511in}}%
\pgfpathlineto{\pgfqpoint{3.361435in}{1.688731in}}%
\pgfpathlineto{\pgfqpoint{3.367368in}{1.688731in}}%
\pgfpathlineto{\pgfqpoint{3.373301in}{1.691341in}}%
\pgfpathlineto{\pgfqpoint{3.379234in}{1.696561in}}%
\pgfpathlineto{\pgfqpoint{3.391101in}{1.701781in}}%
\pgfpathlineto{\pgfqpoint{3.397034in}{1.701781in}}%
\pgfpathlineto{\pgfqpoint{3.402967in}{1.707001in}}%
\pgfpathlineto{\pgfqpoint{3.414833in}{1.712221in}}%
\pgfpathlineto{\pgfqpoint{3.420766in}{1.712221in}}%
\pgfpathlineto{\pgfqpoint{3.426699in}{1.717441in}}%
\pgfpathlineto{\pgfqpoint{3.438566in}{1.722661in}}%
\pgfpathlineto{\pgfqpoint{3.444499in}{1.722661in}}%
\pgfpathlineto{\pgfqpoint{3.450432in}{1.727881in}}%
\pgfpathlineto{\pgfqpoint{3.462298in}{1.733101in}}%
\pgfpathlineto{\pgfqpoint{3.468231in}{1.733101in}}%
\pgfpathlineto{\pgfqpoint{3.474164in}{1.738321in}}%
\pgfpathlineto{\pgfqpoint{3.486031in}{1.743541in}}%
\pgfpathlineto{\pgfqpoint{3.491964in}{1.748761in}}%
\pgfpathlineto{\pgfqpoint{3.497897in}{1.748761in}}%
\pgfpathlineto{\pgfqpoint{3.509763in}{1.753981in}}%
\pgfpathlineto{\pgfqpoint{3.515696in}{1.753981in}}%
\pgfpathlineto{\pgfqpoint{3.521629in}{1.759201in}}%
\pgfpathlineto{\pgfqpoint{3.527562in}{1.761811in}}%
\pgfpathlineto{\pgfqpoint{3.533495in}{1.767031in}}%
\pgfpathlineto{\pgfqpoint{3.539429in}{1.769641in}}%
\pgfpathlineto{\pgfqpoint{3.545362in}{1.769641in}}%
\pgfpathlineto{\pgfqpoint{3.551295in}{1.772251in}}%
\pgfpathlineto{\pgfqpoint{3.557228in}{1.777471in}}%
\pgfpathlineto{\pgfqpoint{3.563161in}{1.780081in}}%
\pgfpathlineto{\pgfqpoint{3.569094in}{1.780081in}}%
\pgfpathlineto{\pgfqpoint{3.575027in}{1.782691in}}%
\pgfpathlineto{\pgfqpoint{3.580960in}{1.787911in}}%
\pgfpathlineto{\pgfqpoint{3.586894in}{1.787911in}}%
\pgfpathlineto{\pgfqpoint{3.592827in}{1.790521in}}%
\pgfpathlineto{\pgfqpoint{3.598760in}{1.795741in}}%
\pgfpathlineto{\pgfqpoint{3.604693in}{1.798351in}}%
\pgfpathlineto{\pgfqpoint{3.610626in}{1.798351in}}%
\pgfpathlineto{\pgfqpoint{3.616559in}{1.800961in}}%
\pgfpathlineto{\pgfqpoint{3.622492in}{1.806181in}}%
\pgfpathlineto{\pgfqpoint{3.628425in}{1.808791in}}%
\pgfpathlineto{\pgfqpoint{3.634358in}{1.808791in}}%
\pgfpathlineto{\pgfqpoint{3.640292in}{1.814011in}}%
\pgfpathlineto{\pgfqpoint{3.646225in}{1.816621in}}%
\pgfpathlineto{\pgfqpoint{3.652158in}{1.821841in}}%
\pgfpathlineto{\pgfqpoint{3.658091in}{1.821841in}}%
\pgfpathlineto{\pgfqpoint{3.669957in}{1.827061in}}%
\pgfpathlineto{\pgfqpoint{3.675890in}{1.827061in}}%
\pgfpathlineto{\pgfqpoint{3.681823in}{1.832281in}}%
\pgfpathlineto{\pgfqpoint{3.687757in}{1.834891in}}%
\pgfpathlineto{\pgfqpoint{3.693690in}{1.840111in}}%
\pgfpathlineto{\pgfqpoint{3.699623in}{1.840111in}}%
\pgfpathlineto{\pgfqpoint{3.705556in}{1.842721in}}%
\pgfpathlineto{\pgfqpoint{3.711489in}{1.847941in}}%
\pgfpathlineto{\pgfqpoint{3.717422in}{1.847941in}}%
\pgfpathlineto{\pgfqpoint{3.723355in}{1.850551in}}%
\pgfpathlineto{\pgfqpoint{3.729288in}{1.855771in}}%
\pgfpathlineto{\pgfqpoint{3.735221in}{1.858381in}}%
\pgfpathlineto{\pgfqpoint{3.741155in}{1.858381in}}%
\pgfpathlineto{\pgfqpoint{3.747088in}{1.863601in}}%
\pgfpathlineto{\pgfqpoint{3.753021in}{1.866211in}}%
\pgfpathlineto{\pgfqpoint{3.758954in}{1.866211in}}%
\pgfpathlineto{\pgfqpoint{3.764887in}{1.871431in}}%
\pgfpathlineto{\pgfqpoint{3.770820in}{1.874041in}}%
\pgfpathlineto{\pgfqpoint{3.776753in}{1.879261in}}%
\pgfpathlineto{\pgfqpoint{3.782686in}{1.879261in}}%
\pgfpathlineto{\pgfqpoint{3.788619in}{1.881871in}}%
\pgfpathlineto{\pgfqpoint{3.794553in}{1.887091in}}%
\pgfpathlineto{\pgfqpoint{3.800486in}{1.887091in}}%
\pgfpathlineto{\pgfqpoint{3.806419in}{1.889702in}}%
\pgfpathlineto{\pgfqpoint{3.812352in}{1.894922in}}%
\pgfpathlineto{\pgfqpoint{3.818285in}{1.894922in}}%
\pgfpathlineto{\pgfqpoint{3.824218in}{1.897532in}}%
\pgfpathlineto{\pgfqpoint{3.830151in}{1.902752in}}%
\pgfpathlineto{\pgfqpoint{3.836084in}{1.902752in}}%
\pgfpathlineto{\pgfqpoint{3.842018in}{1.905362in}}%
\pgfpathlineto{\pgfqpoint{3.847951in}{1.910582in}}%
\pgfpathlineto{\pgfqpoint{3.853884in}{1.910582in}}%
\pgfpathlineto{\pgfqpoint{3.859817in}{1.913192in}}%
\pgfpathlineto{\pgfqpoint{3.865750in}{1.918412in}}%
\pgfpathlineto{\pgfqpoint{3.871683in}{1.918412in}}%
\pgfpathlineto{\pgfqpoint{3.877616in}{1.921022in}}%
\pgfpathlineto{\pgfqpoint{3.883549in}{1.926242in}}%
\pgfpathlineto{\pgfqpoint{3.889482in}{1.926242in}}%
\pgfpathlineto{\pgfqpoint{3.895416in}{1.931462in}}%
\pgfpathlineto{\pgfqpoint{3.901349in}{1.934072in}}%
\pgfpathlineto{\pgfqpoint{3.907282in}{1.934072in}}%
\pgfpathlineto{\pgfqpoint{3.913215in}{1.939292in}}%
\pgfpathlineto{\pgfqpoint{3.919148in}{1.941902in}}%
\pgfpathlineto{\pgfqpoint{3.925081in}{1.941902in}}%
\pgfpathlineto{\pgfqpoint{3.931014in}{1.947122in}}%
\pgfpathlineto{\pgfqpoint{3.936947in}{1.949732in}}%
\pgfpathlineto{\pgfqpoint{3.942881in}{1.949732in}}%
\pgfpathlineto{\pgfqpoint{3.954747in}{1.960172in}}%
\pgfpathlineto{\pgfqpoint{3.960680in}{1.960172in}}%
\pgfpathlineto{\pgfqpoint{3.966613in}{1.962782in}}%
\pgfpathlineto{\pgfqpoint{3.972546in}{1.968002in}}%
\pgfpathlineto{\pgfqpoint{3.978479in}{1.968002in}}%
\pgfpathlineto{\pgfqpoint{3.984412in}{1.970612in}}%
\pgfpathlineto{\pgfqpoint{3.990345in}{1.975832in}}%
\pgfpathlineto{\pgfqpoint{3.996279in}{1.975832in}}%
\pgfpathlineto{\pgfqpoint{4.002212in}{1.981052in}}%
\pgfpathlineto{\pgfqpoint{4.008145in}{1.983662in}}%
\pgfpathlineto{\pgfqpoint{4.014078in}{1.983662in}}%
\pgfpathlineto{\pgfqpoint{4.025944in}{1.994102in}}%
\pgfpathlineto{\pgfqpoint{4.031877in}{1.994102in}}%
\pgfpathlineto{\pgfqpoint{4.037810in}{1.996712in}}%
\pgfpathlineto{\pgfqpoint{4.043744in}{1.996712in}}%
\pgfpathlineto{\pgfqpoint{4.049677in}{2.001932in}}%
\pgfpathlineto{\pgfqpoint{4.055610in}{2.004542in}}%
\pgfpathlineto{\pgfqpoint{4.061543in}{2.009762in}}%
\pgfpathlineto{\pgfqpoint{4.067476in}{2.009762in}}%
\pgfpathlineto{\pgfqpoint{4.073409in}{2.014982in}}%
\pgfpathlineto{\pgfqpoint{4.079342in}{2.014982in}}%
\pgfpathlineto{\pgfqpoint{4.085275in}{2.017592in}}%
\pgfpathlineto{\pgfqpoint{4.097142in}{2.028032in}}%
\pgfpathlineto{\pgfqpoint{4.103075in}{2.028032in}}%
\pgfpathlineto{\pgfqpoint{4.109008in}{2.030642in}}%
\pgfpathlineto{\pgfqpoint{4.114941in}{2.030642in}}%
\pgfpathlineto{\pgfqpoint{4.126807in}{2.041082in}}%
\pgfpathlineto{\pgfqpoint{4.132740in}{2.041082in}}%
\pgfpathlineto{\pgfqpoint{4.138673in}{2.046302in}}%
\pgfpathlineto{\pgfqpoint{4.144607in}{2.048912in}}%
\pgfpathlineto{\pgfqpoint{4.150540in}{2.048912in}}%
\pgfpathlineto{\pgfqpoint{4.156473in}{2.054132in}}%
\pgfpathlineto{\pgfqpoint{4.162406in}{2.054132in}}%
\pgfpathlineto{\pgfqpoint{4.168339in}{2.059352in}}%
\pgfpathlineto{\pgfqpoint{4.174272in}{2.061962in}}%
\pgfpathlineto{\pgfqpoint{4.180205in}{2.061962in}}%
\pgfpathlineto{\pgfqpoint{4.192071in}{2.072402in}}%
\pgfpathlineto{\pgfqpoint{4.198005in}{2.072402in}}%
\pgfpathlineto{\pgfqpoint{4.203938in}{2.077622in}}%
\pgfpathlineto{\pgfqpoint{4.209871in}{2.080232in}}%
\pgfpathlineto{\pgfqpoint{4.215804in}{2.080232in}}%
\pgfpathlineto{\pgfqpoint{4.221737in}{2.085452in}}%
\pgfpathlineto{\pgfqpoint{4.227670in}{2.085452in}}%
\pgfpathlineto{\pgfqpoint{4.239536in}{2.095892in}}%
\pgfpathlineto{\pgfqpoint{4.245469in}{2.095892in}}%
\pgfpathlineto{\pgfqpoint{4.251403in}{2.098502in}}%
\pgfpathlineto{\pgfqpoint{4.257336in}{2.103722in}}%
\pgfpathlineto{\pgfqpoint{4.263269in}{2.103722in}}%
\pgfpathlineto{\pgfqpoint{4.275135in}{2.114162in}}%
\pgfpathlineto{\pgfqpoint{4.281068in}{2.114162in}}%
\pgfpathlineto{\pgfqpoint{4.287001in}{2.119382in}}%
\pgfpathlineto{\pgfqpoint{4.292934in}{2.119382in}}%
\pgfpathlineto{\pgfqpoint{4.298868in}{2.121992in}}%
\pgfpathlineto{\pgfqpoint{4.304801in}{2.127212in}}%
\pgfpathlineto{\pgfqpoint{4.310734in}{2.127212in}}%
\pgfpathlineto{\pgfqpoint{4.322600in}{2.137652in}}%
\pgfpathlineto{\pgfqpoint{4.328533in}{2.137652in}}%
\pgfpathlineto{\pgfqpoint{4.340399in}{2.148092in}}%
\pgfpathlineto{\pgfqpoint{4.346332in}{2.148092in}}%
\pgfpathlineto{\pgfqpoint{4.352266in}{2.150702in}}%
\pgfpathlineto{\pgfqpoint{4.358199in}{2.150702in}}%
\pgfpathlineto{\pgfqpoint{4.370065in}{2.161142in}}%
\pgfpathlineto{\pgfqpoint{4.375998in}{2.161142in}}%
\pgfpathlineto{\pgfqpoint{4.387864in}{2.171582in}}%
\pgfpathlineto{\pgfqpoint{4.393797in}{2.171582in}}%
\pgfpathlineto{\pgfqpoint{4.399731in}{2.176802in}}%
\pgfpathlineto{\pgfqpoint{4.405664in}{2.176802in}}%
\pgfpathlineto{\pgfqpoint{4.417530in}{2.187242in}}%
\pgfpathlineto{\pgfqpoint{4.423463in}{2.187242in}}%
\pgfpathlineto{\pgfqpoint{4.429396in}{2.189852in}}%
\pgfpathlineto{\pgfqpoint{4.435329in}{2.189852in}}%
\pgfpathlineto{\pgfqpoint{4.447195in}{2.200292in}}%
\pgfpathlineto{\pgfqpoint{4.453129in}{2.200292in}}%
\pgfpathlineto{\pgfqpoint{4.464995in}{2.210732in}}%
\pgfpathlineto{\pgfqpoint{4.470928in}{2.210732in}}%
\pgfpathlineto{\pgfqpoint{4.476861in}{2.215952in}}%
\pgfpathlineto{\pgfqpoint{4.482794in}{2.215952in}}%
\pgfpathlineto{\pgfqpoint{4.494660in}{2.226392in}}%
\pgfpathlineto{\pgfqpoint{4.500594in}{2.226392in}}%
\pgfpathlineto{\pgfqpoint{4.506527in}{2.231612in}}%
\pgfpathlineto{\pgfqpoint{4.512460in}{2.231612in}}%
\pgfpathlineto{\pgfqpoint{4.524326in}{2.242052in}}%
\pgfpathlineto{\pgfqpoint{4.530259in}{2.242052in}}%
\pgfpathlineto{\pgfqpoint{4.536192in}{2.247272in}}%
\pgfpathlineto{\pgfqpoint{4.542125in}{2.247272in}}%
\pgfpathlineto{\pgfqpoint{4.553992in}{2.257712in}}%
\pgfpathlineto{\pgfqpoint{4.559925in}{2.257712in}}%
\pgfpathlineto{\pgfqpoint{4.565858in}{2.262932in}}%
\pgfpathlineto{\pgfqpoint{4.571791in}{2.262932in}}%
\pgfpathlineto{\pgfqpoint{4.583657in}{2.273372in}}%
\pgfpathlineto{\pgfqpoint{4.589590in}{2.273372in}}%
\pgfpathlineto{\pgfqpoint{4.595523in}{2.278592in}}%
\pgfpathlineto{\pgfqpoint{4.601457in}{2.278592in}}%
\pgfpathlineto{\pgfqpoint{4.607390in}{2.283812in}}%
\pgfpathlineto{\pgfqpoint{4.613323in}{2.283812in}}%
\pgfpathlineto{\pgfqpoint{4.625189in}{2.294252in}}%
\pgfpathlineto{\pgfqpoint{4.631122in}{2.294252in}}%
\pgfpathlineto{\pgfqpoint{4.637055in}{2.299472in}}%
\pgfpathlineto{\pgfqpoint{4.642988in}{2.299472in}}%
\pgfpathlineto{\pgfqpoint{4.654855in}{2.309913in}}%
\pgfpathlineto{\pgfqpoint{4.660788in}{2.309913in}}%
\pgfpathlineto{\pgfqpoint{4.666721in}{2.315133in}}%
\pgfpathlineto{\pgfqpoint{4.672654in}{2.315133in}}%
\pgfpathlineto{\pgfqpoint{4.678587in}{2.320353in}}%
\pgfpathlineto{\pgfqpoint{4.684520in}{2.320353in}}%
\pgfpathlineto{\pgfqpoint{4.690453in}{2.325573in}}%
\pgfpathlineto{\pgfqpoint{4.696386in}{2.333403in}}%
\pgfpathlineto{\pgfqpoint{4.702319in}{2.333403in}}%
\pgfpathlineto{\pgfqpoint{4.708253in}{2.338623in}}%
\pgfpathlineto{\pgfqpoint{4.714186in}{2.338623in}}%
\pgfpathlineto{\pgfqpoint{4.720119in}{2.343843in}}%
\pgfpathlineto{\pgfqpoint{4.726052in}{2.343843in}}%
\pgfpathlineto{\pgfqpoint{4.737918in}{2.354283in}}%
\pgfpathlineto{\pgfqpoint{4.743851in}{2.354283in}}%
\pgfpathlineto{\pgfqpoint{4.749784in}{2.359503in}}%
\pgfpathlineto{\pgfqpoint{4.755718in}{2.359503in}}%
\pgfpathlineto{\pgfqpoint{4.761651in}{2.364723in}}%
\pgfpathlineto{\pgfqpoint{4.767584in}{2.364723in}}%
\pgfpathlineto{\pgfqpoint{4.773517in}{2.372553in}}%
\pgfpathlineto{\pgfqpoint{4.779450in}{2.377773in}}%
\pgfpathlineto{\pgfqpoint{4.785383in}{2.377773in}}%
\pgfpathlineto{\pgfqpoint{4.791316in}{2.382993in}}%
\pgfpathlineto{\pgfqpoint{4.797249in}{2.382993in}}%
\pgfpathlineto{\pgfqpoint{4.803182in}{2.388213in}}%
\pgfpathlineto{\pgfqpoint{4.809116in}{2.388213in}}%
\pgfpathlineto{\pgfqpoint{4.815049in}{2.393433in}}%
\pgfpathlineto{\pgfqpoint{4.820982in}{2.401263in}}%
\pgfpathlineto{\pgfqpoint{4.826915in}{2.401263in}}%
\pgfpathlineto{\pgfqpoint{4.832848in}{2.406483in}}%
\pgfpathlineto{\pgfqpoint{4.838781in}{2.406483in}}%
\pgfpathlineto{\pgfqpoint{4.844714in}{2.411703in}}%
\pgfpathlineto{\pgfqpoint{4.850647in}{2.411703in}}%
\pgfpathlineto{\pgfqpoint{4.856581in}{2.416923in}}%
\pgfpathlineto{\pgfqpoint{4.862514in}{2.416923in}}%
\pgfpathlineto{\pgfqpoint{4.868447in}{2.424753in}}%
\pgfpathlineto{\pgfqpoint{4.874380in}{2.429973in}}%
\pgfpathlineto{\pgfqpoint{4.880313in}{2.429973in}}%
\pgfpathlineto{\pgfqpoint{4.886246in}{2.435193in}}%
\pgfpathlineto{\pgfqpoint{4.892179in}{2.435193in}}%
\pgfpathlineto{\pgfqpoint{4.898112in}{2.440413in}}%
\pgfpathlineto{\pgfqpoint{4.904045in}{2.440413in}}%
\pgfpathlineto{\pgfqpoint{4.909979in}{2.448243in}}%
\pgfpathlineto{\pgfqpoint{4.915912in}{2.448243in}}%
\pgfpathlineto{\pgfqpoint{4.921845in}{2.453463in}}%
\pgfpathlineto{\pgfqpoint{4.927778in}{2.453463in}}%
\pgfpathlineto{\pgfqpoint{4.933711in}{2.458683in}}%
\pgfpathlineto{\pgfqpoint{4.939644in}{2.466513in}}%
\pgfpathlineto{\pgfqpoint{4.945577in}{2.466513in}}%
\pgfpathlineto{\pgfqpoint{4.951510in}{2.471733in}}%
\pgfpathlineto{\pgfqpoint{4.957444in}{2.471733in}}%
\pgfpathlineto{\pgfqpoint{4.963377in}{2.476953in}}%
\pgfpathlineto{\pgfqpoint{4.969310in}{2.476953in}}%
\pgfpathlineto{\pgfqpoint{4.975243in}{2.484783in}}%
\pgfpathlineto{\pgfqpoint{4.981176in}{2.484783in}}%
\pgfpathlineto{\pgfqpoint{4.987109in}{2.490003in}}%
\pgfpathlineto{\pgfqpoint{4.993042in}{2.490003in}}%
\pgfpathlineto{\pgfqpoint{4.998975in}{2.495223in}}%
\pgfpathlineto{\pgfqpoint{5.004908in}{2.503053in}}%
\pgfpathlineto{\pgfqpoint{5.016775in}{2.503053in}}%
\pgfpathlineto{\pgfqpoint{5.022708in}{2.508273in}}%
\pgfpathlineto{\pgfqpoint{5.028641in}{2.516103in}}%
\pgfpathlineto{\pgfqpoint{5.034574in}{2.516103in}}%
\pgfpathlineto{\pgfqpoint{5.040507in}{2.521323in}}%
\pgfpathlineto{\pgfqpoint{5.046440in}{2.521323in}}%
\pgfpathlineto{\pgfqpoint{5.052373in}{2.529153in}}%
\pgfpathlineto{\pgfqpoint{5.058307in}{2.529153in}}%
\pgfpathlineto{\pgfqpoint{5.064240in}{2.534373in}}%
\pgfpathlineto{\pgfqpoint{5.070173in}{2.534373in}}%
\pgfpathlineto{\pgfqpoint{5.076106in}{2.539593in}}%
\pgfpathlineto{\pgfqpoint{5.082039in}{2.547423in}}%
\pgfpathlineto{\pgfqpoint{5.093905in}{2.547423in}}%
\pgfpathlineto{\pgfqpoint{5.099838in}{2.552643in}}%
\pgfpathlineto{\pgfqpoint{5.105771in}{2.560473in}}%
\pgfpathlineto{\pgfqpoint{5.111705in}{2.560473in}}%
\pgfpathlineto{\pgfqpoint{5.117638in}{2.565693in}}%
\pgfpathlineto{\pgfqpoint{5.123571in}{2.565693in}}%
\pgfpathlineto{\pgfqpoint{5.129504in}{2.573523in}}%
\pgfpathlineto{\pgfqpoint{5.135437in}{2.573523in}}%
\pgfpathlineto{\pgfqpoint{5.141370in}{2.578743in}}%
\pgfpathlineto{\pgfqpoint{5.147303in}{2.578743in}}%
\pgfpathlineto{\pgfqpoint{5.153236in}{2.586573in}}%
\pgfpathlineto{\pgfqpoint{5.159169in}{2.586573in}}%
\pgfpathlineto{\pgfqpoint{5.165103in}{2.591793in}}%
\pgfpathlineto{\pgfqpoint{5.171036in}{2.591793in}}%
\pgfpathlineto{\pgfqpoint{5.176969in}{2.599623in}}%
\pgfpathlineto{\pgfqpoint{5.182902in}{2.599623in}}%
\pgfpathlineto{\pgfqpoint{5.188835in}{2.607453in}}%
\pgfpathlineto{\pgfqpoint{5.194768in}{2.607453in}}%
\pgfpathlineto{\pgfqpoint{5.200701in}{2.612673in}}%
\pgfpathlineto{\pgfqpoint{5.206634in}{2.612673in}}%
\pgfpathlineto{\pgfqpoint{5.212568in}{2.620503in}}%
\pgfpathlineto{\pgfqpoint{5.218501in}{2.620503in}}%
\pgfpathlineto{\pgfqpoint{5.224434in}{2.625723in}}%
\pgfpathlineto{\pgfqpoint{5.230367in}{2.625723in}}%
\pgfpathlineto{\pgfqpoint{5.236300in}{2.633553in}}%
\pgfpathlineto{\pgfqpoint{5.242233in}{2.633553in}}%
\pgfpathlineto{\pgfqpoint{5.248166in}{2.641383in}}%
\pgfpathlineto{\pgfqpoint{5.254099in}{2.641383in}}%
\pgfpathlineto{\pgfqpoint{5.260032in}{2.646603in}}%
\pgfpathlineto{\pgfqpoint{5.265966in}{2.646603in}}%
\pgfpathlineto{\pgfqpoint{5.271899in}{2.654433in}}%
\pgfpathlineto{\pgfqpoint{5.277832in}{2.654433in}}%
\pgfpathlineto{\pgfqpoint{5.283765in}{2.662263in}}%
\pgfpathlineto{\pgfqpoint{5.289698in}{2.667483in}}%
\pgfpathlineto{\pgfqpoint{5.301564in}{2.667483in}}%
\pgfpathlineto{\pgfqpoint{5.307497in}{2.675313in}}%
\pgfpathlineto{\pgfqpoint{5.313431in}{2.675313in}}%
\pgfpathlineto{\pgfqpoint{5.319364in}{2.683143in}}%
\pgfpathlineto{\pgfqpoint{5.325297in}{2.683143in}}%
\pgfpathlineto{\pgfqpoint{5.331230in}{2.688363in}}%
\pgfpathlineto{\pgfqpoint{5.337163in}{2.688363in}}%
\pgfpathlineto{\pgfqpoint{5.343096in}{2.696193in}}%
\pgfpathlineto{\pgfqpoint{5.349029in}{2.696193in}}%
\pgfpathlineto{\pgfqpoint{5.354962in}{2.704023in}}%
\pgfpathlineto{\pgfqpoint{5.360895in}{2.704023in}}%
\pgfpathlineto{\pgfqpoint{5.366829in}{2.711853in}}%
\pgfpathlineto{\pgfqpoint{5.372762in}{2.711853in}}%
\pgfpathlineto{\pgfqpoint{5.378695in}{2.717073in}}%
\pgfpathlineto{\pgfqpoint{5.384628in}{2.717073in}}%
\pgfpathlineto{\pgfqpoint{5.390561in}{2.724904in}}%
\pgfpathlineto{\pgfqpoint{5.396494in}{2.724904in}}%
\pgfpathlineto{\pgfqpoint{5.402427in}{2.732734in}}%
\pgfpathlineto{\pgfqpoint{5.408360in}{2.732734in}}%
\pgfpathlineto{\pgfqpoint{5.414294in}{2.740564in}}%
\pgfpathlineto{\pgfqpoint{5.420227in}{2.740564in}}%
\pgfpathlineto{\pgfqpoint{5.426160in}{2.748394in}}%
\pgfpathlineto{\pgfqpoint{5.432093in}{2.748394in}}%
\pgfpathlineto{\pgfqpoint{5.438026in}{2.756224in}}%
\pgfpathlineto{\pgfqpoint{5.443959in}{2.756224in}}%
\pgfpathlineto{\pgfqpoint{5.449892in}{2.761444in}}%
\pgfpathlineto{\pgfqpoint{5.455825in}{2.761444in}}%
\pgfpathlineto{\pgfqpoint{5.461758in}{2.769274in}}%
\pgfpathlineto{\pgfqpoint{5.467692in}{2.769274in}}%
\pgfpathlineto{\pgfqpoint{5.473625in}{2.777104in}}%
\pgfpathlineto{\pgfqpoint{5.479558in}{2.777104in}}%
\pgfpathlineto{\pgfqpoint{5.485491in}{2.784934in}}%
\pgfpathlineto{\pgfqpoint{5.491424in}{2.784934in}}%
\pgfpathlineto{\pgfqpoint{5.497357in}{2.792764in}}%
\pgfpathlineto{\pgfqpoint{5.503290in}{2.792764in}}%
\pgfpathlineto{\pgfqpoint{5.509223in}{2.800594in}}%
\pgfpathlineto{\pgfqpoint{5.521090in}{2.800594in}}%
\pgfpathlineto{\pgfqpoint{5.527023in}{2.808424in}}%
\pgfpathlineto{\pgfqpoint{5.532956in}{2.808424in}}%
\pgfpathlineto{\pgfqpoint{5.538889in}{2.816254in}}%
\pgfpathlineto{\pgfqpoint{5.544822in}{2.816254in}}%
\pgfpathlineto{\pgfqpoint{5.550755in}{2.824084in}}%
\pgfpathlineto{\pgfqpoint{5.556688in}{2.824084in}}%
\pgfpathlineto{\pgfqpoint{5.562621in}{2.831914in}}%
\pgfpathlineto{\pgfqpoint{5.568555in}{2.831914in}}%
\pgfpathlineto{\pgfqpoint{5.574488in}{2.839744in}}%
\pgfpathlineto{\pgfqpoint{5.580421in}{2.839744in}}%
\pgfpathlineto{\pgfqpoint{5.586354in}{2.847574in}}%
\pgfpathlineto{\pgfqpoint{5.592287in}{2.847574in}}%
\pgfpathlineto{\pgfqpoint{5.598220in}{2.855404in}}%
\pgfpathlineto{\pgfqpoint{5.604153in}{2.855404in}}%
\pgfpathlineto{\pgfqpoint{5.610086in}{2.863234in}}%
\pgfpathlineto{\pgfqpoint{5.616019in}{2.863234in}}%
\pgfpathlineto{\pgfqpoint{5.621953in}{2.871064in}}%
\pgfpathlineto{\pgfqpoint{5.633819in}{2.871064in}}%
\pgfpathlineto{\pgfqpoint{5.639752in}{2.878894in}}%
\pgfpathlineto{\pgfqpoint{5.645685in}{2.878894in}}%
\pgfpathlineto{\pgfqpoint{5.651618in}{2.886724in}}%
\pgfpathlineto{\pgfqpoint{5.657551in}{2.886724in}}%
\pgfpathlineto{\pgfqpoint{5.663484in}{2.894554in}}%
\pgfpathlineto{\pgfqpoint{5.669418in}{2.894554in}}%
\pgfpathlineto{\pgfqpoint{5.675351in}{2.904994in}}%
\pgfpathlineto{\pgfqpoint{5.681284in}{2.904994in}}%
\pgfpathlineto{\pgfqpoint{5.687217in}{2.912824in}}%
\pgfpathlineto{\pgfqpoint{5.699083in}{2.912824in}}%
\pgfpathlineto{\pgfqpoint{5.705016in}{2.920654in}}%
\pgfpathlineto{\pgfqpoint{5.710949in}{2.920654in}}%
\pgfpathlineto{\pgfqpoint{5.716882in}{2.928484in}}%
\pgfpathlineto{\pgfqpoint{5.722816in}{2.928484in}}%
\pgfpathlineto{\pgfqpoint{5.728749in}{2.936314in}}%
\pgfpathlineto{\pgfqpoint{5.734682in}{2.936314in}}%
\pgfpathlineto{\pgfqpoint{5.740615in}{2.946754in}}%
\pgfpathlineto{\pgfqpoint{5.746548in}{2.946754in}}%
\pgfpathlineto{\pgfqpoint{5.752481in}{2.954584in}}%
\pgfpathlineto{\pgfqpoint{5.758414in}{2.954584in}}%
\pgfpathlineto{\pgfqpoint{5.764347in}{2.962414in}}%
\pgfpathlineto{\pgfqpoint{5.776214in}{2.962414in}}%
\pgfpathlineto{\pgfqpoint{5.782147in}{2.970244in}}%
\pgfpathlineto{\pgfqpoint{5.788080in}{2.970244in}}%
\pgfpathlineto{\pgfqpoint{5.794013in}{2.980684in}}%
\pgfpathlineto{\pgfqpoint{5.799946in}{2.980684in}}%
\pgfpathlineto{\pgfqpoint{5.805879in}{2.988514in}}%
\pgfpathlineto{\pgfqpoint{5.817745in}{2.988514in}}%
\pgfpathlineto{\pgfqpoint{5.823679in}{2.996344in}}%
\pgfpathlineto{\pgfqpoint{5.829612in}{2.996344in}}%
\pgfpathlineto{\pgfqpoint{5.835545in}{3.006784in}}%
\pgfpathlineto{\pgfqpoint{5.841478in}{3.006784in}}%
\pgfpathlineto{\pgfqpoint{5.847411in}{3.014614in}}%
\pgfpathlineto{\pgfqpoint{5.853344in}{3.014614in}}%
\pgfpathlineto{\pgfqpoint{5.859277in}{3.022444in}}%
\pgfpathlineto{\pgfqpoint{5.871144in}{3.022444in}}%
\pgfpathlineto{\pgfqpoint{5.877077in}{3.032884in}}%
\pgfpathlineto{\pgfqpoint{5.883010in}{3.032884in}}%
\pgfpathlineto{\pgfqpoint{5.888943in}{3.040714in}}%
\pgfpathlineto{\pgfqpoint{5.900809in}{3.040714in}}%
\pgfpathlineto{\pgfqpoint{5.906742in}{3.051154in}}%
\pgfpathlineto{\pgfqpoint{5.912675in}{3.051154in}}%
\pgfpathlineto{\pgfqpoint{5.918608in}{3.058984in}}%
\pgfpathlineto{\pgfqpoint{5.924542in}{3.058984in}}%
\pgfpathlineto{\pgfqpoint{5.930475in}{3.069424in}}%
\pgfpathlineto{\pgfqpoint{5.936408in}{3.069424in}}%
\pgfpathlineto{\pgfqpoint{5.942341in}{3.077254in}}%
\pgfpathlineto{\pgfqpoint{5.948274in}{3.077254in}}%
\pgfpathlineto{\pgfqpoint{5.954207in}{3.087694in}}%
\pgfpathlineto{\pgfqpoint{5.966073in}{3.087694in}}%
\pgfpathlineto{\pgfqpoint{5.972006in}{3.095524in}}%
\pgfpathlineto{\pgfqpoint{5.977940in}{3.095524in}}%
\pgfpathlineto{\pgfqpoint{5.983873in}{3.105964in}}%
\pgfpathlineto{\pgfqpoint{5.995739in}{3.105964in}}%
\pgfpathlineto{\pgfqpoint{6.001672in}{3.116404in}}%
\pgfpathlineto{\pgfqpoint{6.007605in}{3.116404in}}%
\pgfpathlineto{\pgfqpoint{6.013538in}{3.124234in}}%
\pgfpathlineto{\pgfqpoint{6.019471in}{3.124234in}}%
\pgfpathlineto{\pgfqpoint{6.025405in}{3.134674in}}%
\pgfpathlineto{\pgfqpoint{6.037271in}{3.134674in}}%
\pgfpathlineto{\pgfqpoint{6.043204in}{3.145115in}}%
\pgfpathlineto{\pgfqpoint{6.049137in}{3.145115in}}%
\pgfpathlineto{\pgfqpoint{6.055070in}{3.152945in}}%
\pgfpathlineto{\pgfqpoint{6.061003in}{3.152945in}}%
\pgfpathlineto{\pgfqpoint{6.066936in}{3.163385in}}%
\pgfpathlineto{\pgfqpoint{6.078803in}{3.163385in}}%
\pgfpathlineto{\pgfqpoint{6.084736in}{3.173825in}}%
\pgfpathlineto{\pgfqpoint{6.090669in}{3.173825in}}%
\pgfpathlineto{\pgfqpoint{6.096602in}{3.181655in}}%
\pgfpathlineto{\pgfqpoint{6.102535in}{3.181655in}}%
\pgfpathlineto{\pgfqpoint{6.108468in}{3.192095in}}%
\pgfpathlineto{\pgfqpoint{6.120334in}{3.192095in}}%
\pgfpathlineto{\pgfqpoint{6.126268in}{3.202535in}}%
\pgfpathlineto{\pgfqpoint{6.132201in}{3.202535in}}%
\pgfpathlineto{\pgfqpoint{6.138134in}{3.212975in}}%
\pgfpathlineto{\pgfqpoint{6.144067in}{3.212975in}}%
\pgfpathlineto{\pgfqpoint{6.150000in}{3.223415in}}%
\pgfpathlineto{\pgfqpoint{6.150000in}{3.223415in}}%
\pgfusepath{stroke}%
\end{pgfscope}%
\begin{pgfscope}%
\pgfpathrectangle{\pgfqpoint{0.650001in}{0.565123in}}{\pgfqpoint{5.499999in}{2.784877in}}%
\pgfusepath{clip}%
\pgfsetbuttcap%
\pgfsetroundjoin%
\pgfsetlinewidth{1.505625pt}%
\definecolor{currentstroke}{rgb}{0.000000,0.000000,0.000000}%
\pgfsetstrokecolor{currentstroke}%
\pgfsetdash{{1.500000pt}{2.475000pt}}{0.000000pt}%
\pgfpathmoveto{\pgfqpoint{0.650001in}{1.907972in}}%
\pgfpathlineto{\pgfqpoint{6.150000in}{1.907972in}}%
\pgfusepath{stroke}%
\end{pgfscope}%
\begin{pgfscope}%
\pgfsetrectcap%
\pgfsetmiterjoin%
\pgfsetlinewidth{0.803000pt}%
\definecolor{currentstroke}{rgb}{0.000000,0.000000,0.000000}%
\pgfsetstrokecolor{currentstroke}%
\pgfsetdash{}{0pt}%
\pgfpathmoveto{\pgfqpoint{0.650001in}{0.565123in}}%
\pgfpathlineto{\pgfqpoint{0.650001in}{3.350000in}}%
\pgfusepath{stroke}%
\end{pgfscope}%
\begin{pgfscope}%
\pgfsetrectcap%
\pgfsetmiterjoin%
\pgfsetlinewidth{0.803000pt}%
\definecolor{currentstroke}{rgb}{0.000000,0.000000,0.000000}%
\pgfsetstrokecolor{currentstroke}%
\pgfsetdash{}{0pt}%
\pgfpathmoveto{\pgfqpoint{6.150000in}{0.565123in}}%
\pgfpathlineto{\pgfqpoint{6.150000in}{3.350000in}}%
\pgfusepath{stroke}%
\end{pgfscope}%
\begin{pgfscope}%
\pgfsetrectcap%
\pgfsetmiterjoin%
\pgfsetlinewidth{0.803000pt}%
\definecolor{currentstroke}{rgb}{0.000000,0.000000,0.000000}%
\pgfsetstrokecolor{currentstroke}%
\pgfsetdash{}{0pt}%
\pgfpathmoveto{\pgfqpoint{0.650001in}{0.565123in}}%
\pgfpathlineto{\pgfqpoint{6.150000in}{0.565123in}}%
\pgfusepath{stroke}%
\end{pgfscope}%
\begin{pgfscope}%
\pgfsetrectcap%
\pgfsetmiterjoin%
\pgfsetlinewidth{0.803000pt}%
\definecolor{currentstroke}{rgb}{0.000000,0.000000,0.000000}%
\pgfsetstrokecolor{currentstroke}%
\pgfsetdash{}{0pt}%
\pgfpathmoveto{\pgfqpoint{0.650001in}{3.350000in}}%
\pgfpathlineto{\pgfqpoint{6.150000in}{3.350000in}}%
\pgfusepath{stroke}%
\end{pgfscope}%
\begin{pgfscope}%
\pgfpathrectangle{\pgfqpoint{0.650001in}{0.565123in}}{\pgfqpoint{5.499999in}{2.784877in}}%
\pgfusepath{clip}%
\pgfsetbuttcap%
\pgfsetroundjoin%
\definecolor{currentfill}{rgb}{0.000000,0.000000,0.000000}%
\pgfsetfillcolor{currentfill}%
\pgfsetlinewidth{1.003750pt}%
\definecolor{currentstroke}{rgb}{0.000000,0.000000,0.000000}%
\pgfsetstrokecolor{currentstroke}%
\pgfsetdash{}{0pt}%
\pgfsys@defobject{currentmarker}{\pgfqpoint{-0.041667in}{-0.041667in}}{\pgfqpoint{0.041667in}{0.041667in}}{%
\pgfpathmoveto{\pgfqpoint{0.000000in}{-0.041667in}}%
\pgfpathcurveto{\pgfqpoint{0.011050in}{-0.041667in}}{\pgfqpoint{0.021649in}{-0.037276in}}{\pgfqpoint{0.029463in}{-0.029463in}}%
\pgfpathcurveto{\pgfqpoint{0.037276in}{-0.021649in}}{\pgfqpoint{0.041667in}{-0.011050in}}{\pgfqpoint{0.041667in}{0.000000in}}%
\pgfpathcurveto{\pgfqpoint{0.041667in}{0.011050in}}{\pgfqpoint{0.037276in}{0.021649in}}{\pgfqpoint{0.029463in}{0.029463in}}%
\pgfpathcurveto{\pgfqpoint{0.021649in}{0.037276in}}{\pgfqpoint{0.011050in}{0.041667in}}{\pgfqpoint{0.000000in}{0.041667in}}%
\pgfpathcurveto{\pgfqpoint{-0.011050in}{0.041667in}}{\pgfqpoint{-0.021649in}{0.037276in}}{\pgfqpoint{-0.029463in}{0.029463in}}%
\pgfpathcurveto{\pgfqpoint{-0.037276in}{0.021649in}}{\pgfqpoint{-0.041667in}{0.011050in}}{\pgfqpoint{-0.041667in}{0.000000in}}%
\pgfpathcurveto{\pgfqpoint{-0.041667in}{-0.011050in}}{\pgfqpoint{-0.037276in}{-0.021649in}}{\pgfqpoint{-0.029463in}{-0.029463in}}%
\pgfpathcurveto{\pgfqpoint{-0.021649in}{-0.037276in}}{\pgfqpoint{-0.011050in}{-0.041667in}}{\pgfqpoint{0.000000in}{-0.041667in}}%
\pgfpathlineto{\pgfqpoint{0.000000in}{-0.041667in}}%
\pgfpathclose%
\pgfusepath{stroke,fill}%
}%
\begin{pgfscope}%
\pgfsys@transformshift{3.842018in}{1.905362in}%
\pgfsys@useobject{currentmarker}{}%
\end{pgfscope}%
\end{pgfscope}%
\begin{pgfscope}%
\pgfsetbuttcap%
\pgfsetmiterjoin%
\definecolor{currentfill}{rgb}{1.000000,1.000000,1.000000}%
\pgfsetfillcolor{currentfill}%
\pgfsetfillopacity{0.800000}%
\pgfsetlinewidth{1.003750pt}%
\definecolor{currentstroke}{rgb}{0.800000,0.800000,0.800000}%
\pgfsetstrokecolor{currentstroke}%
\pgfsetstrokeopacity{0.800000}%
\pgfsetdash{}{0pt}%
\pgfpathmoveto{\pgfqpoint{0.747223in}{2.657871in}}%
\pgfpathlineto{\pgfqpoint{1.959261in}{2.657871in}}%
\pgfpathquadraticcurveto{\pgfqpoint{1.987039in}{2.657871in}}{\pgfqpoint{1.987039in}{2.685648in}}%
\pgfpathlineto{\pgfqpoint{1.987039in}{3.252778in}}%
\pgfpathquadraticcurveto{\pgfqpoint{1.987039in}{3.280556in}}{\pgfqpoint{1.959261in}{3.280556in}}%
\pgfpathlineto{\pgfqpoint{0.747223in}{3.280556in}}%
\pgfpathquadraticcurveto{\pgfqpoint{0.719445in}{3.280556in}}{\pgfqpoint{0.719445in}{3.252778in}}%
\pgfpathlineto{\pgfqpoint{0.719445in}{2.685648in}}%
\pgfpathquadraticcurveto{\pgfqpoint{0.719445in}{2.657871in}}{\pgfqpoint{0.747223in}{2.657871in}}%
\pgfpathlineto{\pgfqpoint{0.747223in}{2.657871in}}%
\pgfpathclose%
\pgfusepath{stroke,fill}%
\end{pgfscope}%
\begin{pgfscope}%
\pgfsetrectcap%
\pgfsetroundjoin%
\pgfsetlinewidth{1.003750pt}%
\definecolor{currentstroke}{rgb}{0.000000,0.000000,0.000000}%
\pgfsetstrokecolor{currentstroke}%
\pgfsetdash{}{0pt}%
\pgfpathmoveto{\pgfqpoint{0.775001in}{3.176389in}}%
\pgfpathlineto{\pgfqpoint{0.913890in}{3.176389in}}%
\pgfpathlineto{\pgfqpoint{1.052778in}{3.176389in}}%
\pgfusepath{stroke}%
\end{pgfscope}%
\begin{pgfscope}%
\definecolor{textcolor}{rgb}{0.000000,0.000000,0.000000}%
\pgfsetstrokecolor{textcolor}%
\pgfsetfillcolor{textcolor}%
\pgftext[x=1.163890in,y=3.127778in,left,base]{\color{textcolor}\rmfamily\fontsize{10.000000}{12.000000}\selectfont |Z|}%
\end{pgfscope}%
\begin{pgfscope}%
\pgfsetbuttcap%
\pgfsetroundjoin%
\pgfsetlinewidth{1.505625pt}%
\definecolor{currentstroke}{rgb}{0.000000,0.000000,0.000000}%
\pgfsetstrokecolor{currentstroke}%
\pgfsetdash{{1.500000pt}{2.475000pt}}{0.000000pt}%
\pgfpathmoveto{\pgfqpoint{0.775001in}{2.982716in}}%
\pgfpathlineto{\pgfqpoint{0.913890in}{2.982716in}}%
\pgfpathlineto{\pgfqpoint{1.052778in}{2.982716in}}%
\pgfusepath{stroke}%
\end{pgfscope}%
\begin{pgfscope}%
\definecolor{textcolor}{rgb}{0.000000,0.000000,0.000000}%
\pgfsetstrokecolor{textcolor}%
\pgfsetfillcolor{textcolor}%
\pgftext[x=1.163890in,y=2.934105in,left,base]{\color{textcolor}\rmfamily\fontsize{10.000000}{12.000000}\selectfont Z\_0 = 50 \(\displaystyle \Omega\)}%
\end{pgfscope}%
\begin{pgfscope}%
\pgfsetbuttcap%
\pgfsetroundjoin%
\definecolor{currentfill}{rgb}{0.000000,0.000000,0.000000}%
\pgfsetfillcolor{currentfill}%
\pgfsetlinewidth{1.003750pt}%
\definecolor{currentstroke}{rgb}{0.000000,0.000000,0.000000}%
\pgfsetstrokecolor{currentstroke}%
\pgfsetdash{}{0pt}%
\pgfsys@defobject{currentmarker}{\pgfqpoint{-0.041667in}{-0.041667in}}{\pgfqpoint{0.041667in}{0.041667in}}{%
\pgfpathmoveto{\pgfqpoint{0.000000in}{-0.041667in}}%
\pgfpathcurveto{\pgfqpoint{0.011050in}{-0.041667in}}{\pgfqpoint{0.021649in}{-0.037276in}}{\pgfqpoint{0.029463in}{-0.029463in}}%
\pgfpathcurveto{\pgfqpoint{0.037276in}{-0.021649in}}{\pgfqpoint{0.041667in}{-0.011050in}}{\pgfqpoint{0.041667in}{0.000000in}}%
\pgfpathcurveto{\pgfqpoint{0.041667in}{0.011050in}}{\pgfqpoint{0.037276in}{0.021649in}}{\pgfqpoint{0.029463in}{0.029463in}}%
\pgfpathcurveto{\pgfqpoint{0.021649in}{0.037276in}}{\pgfqpoint{0.011050in}{0.041667in}}{\pgfqpoint{0.000000in}{0.041667in}}%
\pgfpathcurveto{\pgfqpoint{-0.011050in}{0.041667in}}{\pgfqpoint{-0.021649in}{0.037276in}}{\pgfqpoint{-0.029463in}{0.029463in}}%
\pgfpathcurveto{\pgfqpoint{-0.037276in}{0.021649in}}{\pgfqpoint{-0.041667in}{0.011050in}}{\pgfqpoint{-0.041667in}{0.000000in}}%
\pgfpathcurveto{\pgfqpoint{-0.041667in}{-0.011050in}}{\pgfqpoint{-0.037276in}{-0.021649in}}{\pgfqpoint{-0.029463in}{-0.029463in}}%
\pgfpathcurveto{\pgfqpoint{-0.021649in}{-0.037276in}}{\pgfqpoint{-0.011050in}{-0.041667in}}{\pgfqpoint{0.000000in}{-0.041667in}}%
\pgfpathlineto{\pgfqpoint{0.000000in}{-0.041667in}}%
\pgfpathclose%
\pgfusepath{stroke,fill}%
}%
\begin{pgfscope}%
\pgfsys@transformshift{0.913890in}{2.776891in}%
\pgfsys@useobject{currentmarker}{}%
\end{pgfscope}%
\end{pgfscope}%
\begin{pgfscope}%
\definecolor{textcolor}{rgb}{0.000000,0.000000,0.000000}%
\pgfsetstrokecolor{textcolor}%
\pgfsetfillcolor{textcolor}%
\pgftext[x=1.163890in,y=2.740432in,left,base]{\color{textcolor}\rmfamily\fontsize{10.000000}{12.000000}\selectfont |Z| = Z\_0}%
\end{pgfscope}%
\end{pgfpicture}%
\makeatother%
\endgroup%

		\caption{Measured return loss and phase of a quartz resonatoras a function of frequency.}
		\label{fig:355}
	\end{figure}
	The measured series resistance is 1 $\Omega$. To compare this value to the theoretical value the equation for our reflection at -0.18 dB and 90 degree needs to be evaluated:
	\[r=\frac{Z_L - Z_0}{Z_L + Z_0} \Leftrightarrow Z_L = Z_0 \frac{1 + r}{1 - r} \approx 1 + \mathrm{j} \, 50\Omega\]
	\[L = \frac{X_L}{2\pi f} = \frac{50}{2\pi \cdot 17.434 \, \mathrm{MHz}} \approx 456,44 \, \mathrm{nH}\]
	This result matches the measurement very well.
	
	\subsubsection{}
		
	\begin{figure}[H]
		\centering
		%% Creator: Matplotlib, PGF backend
%%
%% To include the figure in your LaTeX document, write
%%   \input{<filename>.pgf}
%%
%% Make sure the required packages are loaded in your preamble
%%   \usepackage{pgf}
%%
%% Also ensure that all the required font packages are loaded; for instance,
%% the lmodern package is sometimes necessary when using math font.
%%   \usepackage{lmodern}
%%
%% Figures using additional raster images can only be included by \input if
%% they are in the same directory as the main LaTeX file. For loading figures
%% from other directories you can use the `import` package
%%   \usepackage{import}
%%
%% and then include the figures with
%%   \import{<path to file>}{<filename>.pgf}
%%
%% Matplotlib used the following preamble
%%
\begingroup%
\makeatletter%
\begin{pgfpicture}%
\pgfpathrectangle{\pgfpointorigin}{\pgfqpoint{6.300000in}{3.500000in}}%
\pgfusepath{use as bounding box, clip}%
\begin{pgfscope}%
\pgfsetbuttcap%
\pgfsetmiterjoin%
\definecolor{currentfill}{rgb}{1.000000,1.000000,1.000000}%
\pgfsetfillcolor{currentfill}%
\pgfsetlinewidth{0.000000pt}%
\definecolor{currentstroke}{rgb}{1.000000,1.000000,1.000000}%
\pgfsetstrokecolor{currentstroke}%
\pgfsetdash{}{0pt}%
\pgfpathmoveto{\pgfqpoint{0.000000in}{0.000000in}}%
\pgfpathlineto{\pgfqpoint{6.300000in}{0.000000in}}%
\pgfpathlineto{\pgfqpoint{6.300000in}{3.500000in}}%
\pgfpathlineto{\pgfqpoint{0.000000in}{3.500000in}}%
\pgfpathlineto{\pgfqpoint{0.000000in}{0.000000in}}%
\pgfpathclose%
\pgfusepath{fill}%
\end{pgfscope}%
\begin{pgfscope}%
\pgfsetbuttcap%
\pgfsetmiterjoin%
\definecolor{currentfill}{rgb}{1.000000,1.000000,1.000000}%
\pgfsetfillcolor{currentfill}%
\pgfsetlinewidth{0.000000pt}%
\definecolor{currentstroke}{rgb}{0.000000,0.000000,0.000000}%
\pgfsetstrokecolor{currentstroke}%
\pgfsetstrokeopacity{0.000000}%
\pgfsetdash{}{0pt}%
\pgfpathmoveto{\pgfqpoint{0.650001in}{0.565123in}}%
\pgfpathlineto{\pgfqpoint{6.150000in}{0.565123in}}%
\pgfpathlineto{\pgfqpoint{6.150000in}{3.350000in}}%
\pgfpathlineto{\pgfqpoint{0.650001in}{3.350000in}}%
\pgfpathlineto{\pgfqpoint{0.650001in}{0.565123in}}%
\pgfpathclose%
\pgfusepath{fill}%
\end{pgfscope}%
\begin{pgfscope}%
\pgfpathrectangle{\pgfqpoint{0.650001in}{0.565123in}}{\pgfqpoint{5.499999in}{2.784877in}}%
\pgfusepath{clip}%
\pgfsetbuttcap%
\pgfsetroundjoin%
\pgfsetlinewidth{0.803000pt}%
\definecolor{currentstroke}{rgb}{0.690196,0.690196,0.690196}%
\pgfsetstrokecolor{currentstroke}%
\pgfsetdash{{0.800000pt}{1.320000pt}}{0.000000pt}%
\pgfpathmoveto{\pgfqpoint{1.234432in}{0.565123in}}%
\pgfpathlineto{\pgfqpoint{1.234432in}{3.350000in}}%
\pgfusepath{stroke}%
\end{pgfscope}%
\begin{pgfscope}%
\pgfsetbuttcap%
\pgfsetroundjoin%
\definecolor{currentfill}{rgb}{0.000000,0.000000,0.000000}%
\pgfsetfillcolor{currentfill}%
\pgfsetlinewidth{0.803000pt}%
\definecolor{currentstroke}{rgb}{0.000000,0.000000,0.000000}%
\pgfsetstrokecolor{currentstroke}%
\pgfsetdash{}{0pt}%
\pgfsys@defobject{currentmarker}{\pgfqpoint{0.000000in}{-0.048611in}}{\pgfqpoint{0.000000in}{0.000000in}}{%
\pgfpathmoveto{\pgfqpoint{0.000000in}{0.000000in}}%
\pgfpathlineto{\pgfqpoint{0.000000in}{-0.048611in}}%
\pgfusepath{stroke,fill}%
}%
\begin{pgfscope}%
\pgfsys@transformshift{1.234432in}{0.565123in}%
\pgfsys@useobject{currentmarker}{}%
\end{pgfscope}%
\end{pgfscope}%
\begin{pgfscope}%
\definecolor{textcolor}{rgb}{0.000000,0.000000,0.000000}%
\pgfsetstrokecolor{textcolor}%
\pgfsetfillcolor{textcolor}%
\pgftext[x=1.234432in,y=0.467901in,,top]{\color{textcolor}\rmfamily\fontsize{10.000000}{12.000000}\selectfont \(\displaystyle {0.2}\)}%
\end{pgfscope}%
\begin{pgfscope}%
\pgfpathrectangle{\pgfqpoint{0.650001in}{0.565123in}}{\pgfqpoint{5.499999in}{2.784877in}}%
\pgfusepath{clip}%
\pgfsetbuttcap%
\pgfsetroundjoin%
\pgfsetlinewidth{0.803000pt}%
\definecolor{currentstroke}{rgb}{0.690196,0.690196,0.690196}%
\pgfsetstrokecolor{currentstroke}%
\pgfsetdash{{0.800000pt}{1.320000pt}}{0.000000pt}%
\pgfpathmoveto{\pgfqpoint{1.849623in}{0.565123in}}%
\pgfpathlineto{\pgfqpoint{1.849623in}{3.350000in}}%
\pgfusepath{stroke}%
\end{pgfscope}%
\begin{pgfscope}%
\pgfsetbuttcap%
\pgfsetroundjoin%
\definecolor{currentfill}{rgb}{0.000000,0.000000,0.000000}%
\pgfsetfillcolor{currentfill}%
\pgfsetlinewidth{0.803000pt}%
\definecolor{currentstroke}{rgb}{0.000000,0.000000,0.000000}%
\pgfsetstrokecolor{currentstroke}%
\pgfsetdash{}{0pt}%
\pgfsys@defobject{currentmarker}{\pgfqpoint{0.000000in}{-0.048611in}}{\pgfqpoint{0.000000in}{0.000000in}}{%
\pgfpathmoveto{\pgfqpoint{0.000000in}{0.000000in}}%
\pgfpathlineto{\pgfqpoint{0.000000in}{-0.048611in}}%
\pgfusepath{stroke,fill}%
}%
\begin{pgfscope}%
\pgfsys@transformshift{1.849623in}{0.565123in}%
\pgfsys@useobject{currentmarker}{}%
\end{pgfscope}%
\end{pgfscope}%
\begin{pgfscope}%
\definecolor{textcolor}{rgb}{0.000000,0.000000,0.000000}%
\pgfsetstrokecolor{textcolor}%
\pgfsetfillcolor{textcolor}%
\pgftext[x=1.849623in,y=0.467901in,,top]{\color{textcolor}\rmfamily\fontsize{10.000000}{12.000000}\selectfont \(\displaystyle {0.4}\)}%
\end{pgfscope}%
\begin{pgfscope}%
\pgfpathrectangle{\pgfqpoint{0.650001in}{0.565123in}}{\pgfqpoint{5.499999in}{2.784877in}}%
\pgfusepath{clip}%
\pgfsetbuttcap%
\pgfsetroundjoin%
\pgfsetlinewidth{0.803000pt}%
\definecolor{currentstroke}{rgb}{0.690196,0.690196,0.690196}%
\pgfsetstrokecolor{currentstroke}%
\pgfsetdash{{0.800000pt}{1.320000pt}}{0.000000pt}%
\pgfpathmoveto{\pgfqpoint{2.464814in}{0.565123in}}%
\pgfpathlineto{\pgfqpoint{2.464814in}{3.350000in}}%
\pgfusepath{stroke}%
\end{pgfscope}%
\begin{pgfscope}%
\pgfsetbuttcap%
\pgfsetroundjoin%
\definecolor{currentfill}{rgb}{0.000000,0.000000,0.000000}%
\pgfsetfillcolor{currentfill}%
\pgfsetlinewidth{0.803000pt}%
\definecolor{currentstroke}{rgb}{0.000000,0.000000,0.000000}%
\pgfsetstrokecolor{currentstroke}%
\pgfsetdash{}{0pt}%
\pgfsys@defobject{currentmarker}{\pgfqpoint{0.000000in}{-0.048611in}}{\pgfqpoint{0.000000in}{0.000000in}}{%
\pgfpathmoveto{\pgfqpoint{0.000000in}{0.000000in}}%
\pgfpathlineto{\pgfqpoint{0.000000in}{-0.048611in}}%
\pgfusepath{stroke,fill}%
}%
\begin{pgfscope}%
\pgfsys@transformshift{2.464814in}{0.565123in}%
\pgfsys@useobject{currentmarker}{}%
\end{pgfscope}%
\end{pgfscope}%
\begin{pgfscope}%
\definecolor{textcolor}{rgb}{0.000000,0.000000,0.000000}%
\pgfsetstrokecolor{textcolor}%
\pgfsetfillcolor{textcolor}%
\pgftext[x=2.464814in,y=0.467901in,,top]{\color{textcolor}\rmfamily\fontsize{10.000000}{12.000000}\selectfont \(\displaystyle {0.6}\)}%
\end{pgfscope}%
\begin{pgfscope}%
\pgfpathrectangle{\pgfqpoint{0.650001in}{0.565123in}}{\pgfqpoint{5.499999in}{2.784877in}}%
\pgfusepath{clip}%
\pgfsetbuttcap%
\pgfsetroundjoin%
\pgfsetlinewidth{0.803000pt}%
\definecolor{currentstroke}{rgb}{0.690196,0.690196,0.690196}%
\pgfsetstrokecolor{currentstroke}%
\pgfsetdash{{0.800000pt}{1.320000pt}}{0.000000pt}%
\pgfpathmoveto{\pgfqpoint{3.080005in}{0.565123in}}%
\pgfpathlineto{\pgfqpoint{3.080005in}{3.350000in}}%
\pgfusepath{stroke}%
\end{pgfscope}%
\begin{pgfscope}%
\pgfsetbuttcap%
\pgfsetroundjoin%
\definecolor{currentfill}{rgb}{0.000000,0.000000,0.000000}%
\pgfsetfillcolor{currentfill}%
\pgfsetlinewidth{0.803000pt}%
\definecolor{currentstroke}{rgb}{0.000000,0.000000,0.000000}%
\pgfsetstrokecolor{currentstroke}%
\pgfsetdash{}{0pt}%
\pgfsys@defobject{currentmarker}{\pgfqpoint{0.000000in}{-0.048611in}}{\pgfqpoint{0.000000in}{0.000000in}}{%
\pgfpathmoveto{\pgfqpoint{0.000000in}{0.000000in}}%
\pgfpathlineto{\pgfqpoint{0.000000in}{-0.048611in}}%
\pgfusepath{stroke,fill}%
}%
\begin{pgfscope}%
\pgfsys@transformshift{3.080005in}{0.565123in}%
\pgfsys@useobject{currentmarker}{}%
\end{pgfscope}%
\end{pgfscope}%
\begin{pgfscope}%
\definecolor{textcolor}{rgb}{0.000000,0.000000,0.000000}%
\pgfsetstrokecolor{textcolor}%
\pgfsetfillcolor{textcolor}%
\pgftext[x=3.080005in,y=0.467901in,,top]{\color{textcolor}\rmfamily\fontsize{10.000000}{12.000000}\selectfont \(\displaystyle {0.8}\)}%
\end{pgfscope}%
\begin{pgfscope}%
\pgfpathrectangle{\pgfqpoint{0.650001in}{0.565123in}}{\pgfqpoint{5.499999in}{2.784877in}}%
\pgfusepath{clip}%
\pgfsetbuttcap%
\pgfsetroundjoin%
\pgfsetlinewidth{0.803000pt}%
\definecolor{currentstroke}{rgb}{0.690196,0.690196,0.690196}%
\pgfsetstrokecolor{currentstroke}%
\pgfsetdash{{0.800000pt}{1.320000pt}}{0.000000pt}%
\pgfpathmoveto{\pgfqpoint{3.695196in}{0.565123in}}%
\pgfpathlineto{\pgfqpoint{3.695196in}{3.350000in}}%
\pgfusepath{stroke}%
\end{pgfscope}%
\begin{pgfscope}%
\pgfsetbuttcap%
\pgfsetroundjoin%
\definecolor{currentfill}{rgb}{0.000000,0.000000,0.000000}%
\pgfsetfillcolor{currentfill}%
\pgfsetlinewidth{0.803000pt}%
\definecolor{currentstroke}{rgb}{0.000000,0.000000,0.000000}%
\pgfsetstrokecolor{currentstroke}%
\pgfsetdash{}{0pt}%
\pgfsys@defobject{currentmarker}{\pgfqpoint{0.000000in}{-0.048611in}}{\pgfqpoint{0.000000in}{0.000000in}}{%
\pgfpathmoveto{\pgfqpoint{0.000000in}{0.000000in}}%
\pgfpathlineto{\pgfqpoint{0.000000in}{-0.048611in}}%
\pgfusepath{stroke,fill}%
}%
\begin{pgfscope}%
\pgfsys@transformshift{3.695196in}{0.565123in}%
\pgfsys@useobject{currentmarker}{}%
\end{pgfscope}%
\end{pgfscope}%
\begin{pgfscope}%
\definecolor{textcolor}{rgb}{0.000000,0.000000,0.000000}%
\pgfsetstrokecolor{textcolor}%
\pgfsetfillcolor{textcolor}%
\pgftext[x=3.695196in,y=0.467901in,,top]{\color{textcolor}\rmfamily\fontsize{10.000000}{12.000000}\selectfont \(\displaystyle {1.0}\)}%
\end{pgfscope}%
\begin{pgfscope}%
\pgfpathrectangle{\pgfqpoint{0.650001in}{0.565123in}}{\pgfqpoint{5.499999in}{2.784877in}}%
\pgfusepath{clip}%
\pgfsetbuttcap%
\pgfsetroundjoin%
\pgfsetlinewidth{0.803000pt}%
\definecolor{currentstroke}{rgb}{0.690196,0.690196,0.690196}%
\pgfsetstrokecolor{currentstroke}%
\pgfsetdash{{0.800000pt}{1.320000pt}}{0.000000pt}%
\pgfpathmoveto{\pgfqpoint{4.310387in}{0.565123in}}%
\pgfpathlineto{\pgfqpoint{4.310387in}{3.350000in}}%
\pgfusepath{stroke}%
\end{pgfscope}%
\begin{pgfscope}%
\pgfsetbuttcap%
\pgfsetroundjoin%
\definecolor{currentfill}{rgb}{0.000000,0.000000,0.000000}%
\pgfsetfillcolor{currentfill}%
\pgfsetlinewidth{0.803000pt}%
\definecolor{currentstroke}{rgb}{0.000000,0.000000,0.000000}%
\pgfsetstrokecolor{currentstroke}%
\pgfsetdash{}{0pt}%
\pgfsys@defobject{currentmarker}{\pgfqpoint{0.000000in}{-0.048611in}}{\pgfqpoint{0.000000in}{0.000000in}}{%
\pgfpathmoveto{\pgfqpoint{0.000000in}{0.000000in}}%
\pgfpathlineto{\pgfqpoint{0.000000in}{-0.048611in}}%
\pgfusepath{stroke,fill}%
}%
\begin{pgfscope}%
\pgfsys@transformshift{4.310387in}{0.565123in}%
\pgfsys@useobject{currentmarker}{}%
\end{pgfscope}%
\end{pgfscope}%
\begin{pgfscope}%
\definecolor{textcolor}{rgb}{0.000000,0.000000,0.000000}%
\pgfsetstrokecolor{textcolor}%
\pgfsetfillcolor{textcolor}%
\pgftext[x=4.310387in,y=0.467901in,,top]{\color{textcolor}\rmfamily\fontsize{10.000000}{12.000000}\selectfont \(\displaystyle {1.2}\)}%
\end{pgfscope}%
\begin{pgfscope}%
\pgfpathrectangle{\pgfqpoint{0.650001in}{0.565123in}}{\pgfqpoint{5.499999in}{2.784877in}}%
\pgfusepath{clip}%
\pgfsetbuttcap%
\pgfsetroundjoin%
\pgfsetlinewidth{0.803000pt}%
\definecolor{currentstroke}{rgb}{0.690196,0.690196,0.690196}%
\pgfsetstrokecolor{currentstroke}%
\pgfsetdash{{0.800000pt}{1.320000pt}}{0.000000pt}%
\pgfpathmoveto{\pgfqpoint{4.925578in}{0.565123in}}%
\pgfpathlineto{\pgfqpoint{4.925578in}{3.350000in}}%
\pgfusepath{stroke}%
\end{pgfscope}%
\begin{pgfscope}%
\pgfsetbuttcap%
\pgfsetroundjoin%
\definecolor{currentfill}{rgb}{0.000000,0.000000,0.000000}%
\pgfsetfillcolor{currentfill}%
\pgfsetlinewidth{0.803000pt}%
\definecolor{currentstroke}{rgb}{0.000000,0.000000,0.000000}%
\pgfsetstrokecolor{currentstroke}%
\pgfsetdash{}{0pt}%
\pgfsys@defobject{currentmarker}{\pgfqpoint{0.000000in}{-0.048611in}}{\pgfqpoint{0.000000in}{0.000000in}}{%
\pgfpathmoveto{\pgfqpoint{0.000000in}{0.000000in}}%
\pgfpathlineto{\pgfqpoint{0.000000in}{-0.048611in}}%
\pgfusepath{stroke,fill}%
}%
\begin{pgfscope}%
\pgfsys@transformshift{4.925578in}{0.565123in}%
\pgfsys@useobject{currentmarker}{}%
\end{pgfscope}%
\end{pgfscope}%
\begin{pgfscope}%
\definecolor{textcolor}{rgb}{0.000000,0.000000,0.000000}%
\pgfsetstrokecolor{textcolor}%
\pgfsetfillcolor{textcolor}%
\pgftext[x=4.925578in,y=0.467901in,,top]{\color{textcolor}\rmfamily\fontsize{10.000000}{12.000000}\selectfont \(\displaystyle {1.4}\)}%
\end{pgfscope}%
\begin{pgfscope}%
\pgfpathrectangle{\pgfqpoint{0.650001in}{0.565123in}}{\pgfqpoint{5.499999in}{2.784877in}}%
\pgfusepath{clip}%
\pgfsetbuttcap%
\pgfsetroundjoin%
\pgfsetlinewidth{0.803000pt}%
\definecolor{currentstroke}{rgb}{0.690196,0.690196,0.690196}%
\pgfsetstrokecolor{currentstroke}%
\pgfsetdash{{0.800000pt}{1.320000pt}}{0.000000pt}%
\pgfpathmoveto{\pgfqpoint{5.540769in}{0.565123in}}%
\pgfpathlineto{\pgfqpoint{5.540769in}{3.350000in}}%
\pgfusepath{stroke}%
\end{pgfscope}%
\begin{pgfscope}%
\pgfsetbuttcap%
\pgfsetroundjoin%
\definecolor{currentfill}{rgb}{0.000000,0.000000,0.000000}%
\pgfsetfillcolor{currentfill}%
\pgfsetlinewidth{0.803000pt}%
\definecolor{currentstroke}{rgb}{0.000000,0.000000,0.000000}%
\pgfsetstrokecolor{currentstroke}%
\pgfsetdash{}{0pt}%
\pgfsys@defobject{currentmarker}{\pgfqpoint{0.000000in}{-0.048611in}}{\pgfqpoint{0.000000in}{0.000000in}}{%
\pgfpathmoveto{\pgfqpoint{0.000000in}{0.000000in}}%
\pgfpathlineto{\pgfqpoint{0.000000in}{-0.048611in}}%
\pgfusepath{stroke,fill}%
}%
\begin{pgfscope}%
\pgfsys@transformshift{5.540769in}{0.565123in}%
\pgfsys@useobject{currentmarker}{}%
\end{pgfscope}%
\end{pgfscope}%
\begin{pgfscope}%
\definecolor{textcolor}{rgb}{0.000000,0.000000,0.000000}%
\pgfsetstrokecolor{textcolor}%
\pgfsetfillcolor{textcolor}%
\pgftext[x=5.540769in,y=0.467901in,,top]{\color{textcolor}\rmfamily\fontsize{10.000000}{12.000000}\selectfont \(\displaystyle {1.6}\)}%
\end{pgfscope}%
\begin{pgfscope}%
\definecolor{textcolor}{rgb}{0.000000,0.000000,0.000000}%
\pgfsetstrokecolor{textcolor}%
\pgfsetfillcolor{textcolor}%
\pgftext[x=3.400000in,y=0.288889in,,top]{\color{textcolor}\rmfamily\fontsize{10.000000}{12.000000}\selectfont Frequency (Hz)}%
\end{pgfscope}%
\begin{pgfscope}%
\definecolor{textcolor}{rgb}{0.000000,0.000000,0.000000}%
\pgfsetstrokecolor{textcolor}%
\pgfsetfillcolor{textcolor}%
\pgftext[x=6.150000in,y=0.302778in,right,top]{\color{textcolor}\rmfamily\fontsize{10.000000}{12.000000}\selectfont \(\displaystyle \times{10^{8}}{}\)}%
\end{pgfscope}%
\begin{pgfscope}%
\pgfpathrectangle{\pgfqpoint{0.650001in}{0.565123in}}{\pgfqpoint{5.499999in}{2.784877in}}%
\pgfusepath{clip}%
\pgfsetbuttcap%
\pgfsetroundjoin%
\pgfsetlinewidth{0.803000pt}%
\definecolor{currentstroke}{rgb}{0.690196,0.690196,0.690196}%
\pgfsetstrokecolor{currentstroke}%
\pgfsetdash{{0.800000pt}{1.320000pt}}{0.000000pt}%
\pgfpathmoveto{\pgfqpoint{0.650001in}{0.674290in}}%
\pgfpathlineto{\pgfqpoint{6.150000in}{0.674290in}}%
\pgfusepath{stroke}%
\end{pgfscope}%
\begin{pgfscope}%
\pgfsetbuttcap%
\pgfsetroundjoin%
\definecolor{currentfill}{rgb}{0.000000,0.000000,0.000000}%
\pgfsetfillcolor{currentfill}%
\pgfsetlinewidth{0.803000pt}%
\definecolor{currentstroke}{rgb}{0.000000,0.000000,0.000000}%
\pgfsetstrokecolor{currentstroke}%
\pgfsetdash{}{0pt}%
\pgfsys@defobject{currentmarker}{\pgfqpoint{-0.048611in}{0.000000in}}{\pgfqpoint{-0.000000in}{0.000000in}}{%
\pgfpathmoveto{\pgfqpoint{-0.000000in}{0.000000in}}%
\pgfpathlineto{\pgfqpoint{-0.048611in}{0.000000in}}%
\pgfusepath{stroke,fill}%
}%
\begin{pgfscope}%
\pgfsys@transformshift{0.650001in}{0.674290in}%
\pgfsys@useobject{currentmarker}{}%
\end{pgfscope}%
\end{pgfscope}%
\begin{pgfscope}%
\definecolor{textcolor}{rgb}{0.000000,0.000000,0.000000}%
\pgfsetstrokecolor{textcolor}%
\pgfsetfillcolor{textcolor}%
\pgftext[x=0.483334in, y=0.626064in, left, base]{\color{textcolor}\rmfamily\fontsize{10.000000}{12.000000}\selectfont \(\displaystyle {0}\)}%
\end{pgfscope}%
\begin{pgfscope}%
\pgfpathrectangle{\pgfqpoint{0.650001in}{0.565123in}}{\pgfqpoint{5.499999in}{2.784877in}}%
\pgfusepath{clip}%
\pgfsetbuttcap%
\pgfsetroundjoin%
\pgfsetlinewidth{0.803000pt}%
\definecolor{currentstroke}{rgb}{0.690196,0.690196,0.690196}%
\pgfsetstrokecolor{currentstroke}%
\pgfsetdash{{0.800000pt}{1.320000pt}}{0.000000pt}%
\pgfpathmoveto{\pgfqpoint{0.650001in}{1.028335in}}%
\pgfpathlineto{\pgfqpoint{6.150000in}{1.028335in}}%
\pgfusepath{stroke}%
\end{pgfscope}%
\begin{pgfscope}%
\pgfsetbuttcap%
\pgfsetroundjoin%
\definecolor{currentfill}{rgb}{0.000000,0.000000,0.000000}%
\pgfsetfillcolor{currentfill}%
\pgfsetlinewidth{0.803000pt}%
\definecolor{currentstroke}{rgb}{0.000000,0.000000,0.000000}%
\pgfsetstrokecolor{currentstroke}%
\pgfsetdash{}{0pt}%
\pgfsys@defobject{currentmarker}{\pgfqpoint{-0.048611in}{0.000000in}}{\pgfqpoint{-0.000000in}{0.000000in}}{%
\pgfpathmoveto{\pgfqpoint{-0.000000in}{0.000000in}}%
\pgfpathlineto{\pgfqpoint{-0.048611in}{0.000000in}}%
\pgfusepath{stroke,fill}%
}%
\begin{pgfscope}%
\pgfsys@transformshift{0.650001in}{1.028335in}%
\pgfsys@useobject{currentmarker}{}%
\end{pgfscope}%
\end{pgfscope}%
\begin{pgfscope}%
\definecolor{textcolor}{rgb}{0.000000,0.000000,0.000000}%
\pgfsetstrokecolor{textcolor}%
\pgfsetfillcolor{textcolor}%
\pgftext[x=0.413889in, y=0.980109in, left, base]{\color{textcolor}\rmfamily\fontsize{10.000000}{12.000000}\selectfont \(\displaystyle {25}\)}%
\end{pgfscope}%
\begin{pgfscope}%
\pgfpathrectangle{\pgfqpoint{0.650001in}{0.565123in}}{\pgfqpoint{5.499999in}{2.784877in}}%
\pgfusepath{clip}%
\pgfsetbuttcap%
\pgfsetroundjoin%
\pgfsetlinewidth{0.803000pt}%
\definecolor{currentstroke}{rgb}{0.690196,0.690196,0.690196}%
\pgfsetstrokecolor{currentstroke}%
\pgfsetdash{{0.800000pt}{1.320000pt}}{0.000000pt}%
\pgfpathmoveto{\pgfqpoint{0.650001in}{1.382380in}}%
\pgfpathlineto{\pgfqpoint{6.150000in}{1.382380in}}%
\pgfusepath{stroke}%
\end{pgfscope}%
\begin{pgfscope}%
\pgfsetbuttcap%
\pgfsetroundjoin%
\definecolor{currentfill}{rgb}{0.000000,0.000000,0.000000}%
\pgfsetfillcolor{currentfill}%
\pgfsetlinewidth{0.803000pt}%
\definecolor{currentstroke}{rgb}{0.000000,0.000000,0.000000}%
\pgfsetstrokecolor{currentstroke}%
\pgfsetdash{}{0pt}%
\pgfsys@defobject{currentmarker}{\pgfqpoint{-0.048611in}{0.000000in}}{\pgfqpoint{-0.000000in}{0.000000in}}{%
\pgfpathmoveto{\pgfqpoint{-0.000000in}{0.000000in}}%
\pgfpathlineto{\pgfqpoint{-0.048611in}{0.000000in}}%
\pgfusepath{stroke,fill}%
}%
\begin{pgfscope}%
\pgfsys@transformshift{0.650001in}{1.382380in}%
\pgfsys@useobject{currentmarker}{}%
\end{pgfscope}%
\end{pgfscope}%
\begin{pgfscope}%
\definecolor{textcolor}{rgb}{0.000000,0.000000,0.000000}%
\pgfsetstrokecolor{textcolor}%
\pgfsetfillcolor{textcolor}%
\pgftext[x=0.413889in, y=1.334155in, left, base]{\color{textcolor}\rmfamily\fontsize{10.000000}{12.000000}\selectfont \(\displaystyle {50}\)}%
\end{pgfscope}%
\begin{pgfscope}%
\pgfpathrectangle{\pgfqpoint{0.650001in}{0.565123in}}{\pgfqpoint{5.499999in}{2.784877in}}%
\pgfusepath{clip}%
\pgfsetbuttcap%
\pgfsetroundjoin%
\pgfsetlinewidth{0.803000pt}%
\definecolor{currentstroke}{rgb}{0.690196,0.690196,0.690196}%
\pgfsetstrokecolor{currentstroke}%
\pgfsetdash{{0.800000pt}{1.320000pt}}{0.000000pt}%
\pgfpathmoveto{\pgfqpoint{0.650001in}{1.736425in}}%
\pgfpathlineto{\pgfqpoint{6.150000in}{1.736425in}}%
\pgfusepath{stroke}%
\end{pgfscope}%
\begin{pgfscope}%
\pgfsetbuttcap%
\pgfsetroundjoin%
\definecolor{currentfill}{rgb}{0.000000,0.000000,0.000000}%
\pgfsetfillcolor{currentfill}%
\pgfsetlinewidth{0.803000pt}%
\definecolor{currentstroke}{rgb}{0.000000,0.000000,0.000000}%
\pgfsetstrokecolor{currentstroke}%
\pgfsetdash{}{0pt}%
\pgfsys@defobject{currentmarker}{\pgfqpoint{-0.048611in}{0.000000in}}{\pgfqpoint{-0.000000in}{0.000000in}}{%
\pgfpathmoveto{\pgfqpoint{-0.000000in}{0.000000in}}%
\pgfpathlineto{\pgfqpoint{-0.048611in}{0.000000in}}%
\pgfusepath{stroke,fill}%
}%
\begin{pgfscope}%
\pgfsys@transformshift{0.650001in}{1.736425in}%
\pgfsys@useobject{currentmarker}{}%
\end{pgfscope}%
\end{pgfscope}%
\begin{pgfscope}%
\definecolor{textcolor}{rgb}{0.000000,0.000000,0.000000}%
\pgfsetstrokecolor{textcolor}%
\pgfsetfillcolor{textcolor}%
\pgftext[x=0.413889in, y=1.688200in, left, base]{\color{textcolor}\rmfamily\fontsize{10.000000}{12.000000}\selectfont \(\displaystyle {75}\)}%
\end{pgfscope}%
\begin{pgfscope}%
\pgfpathrectangle{\pgfqpoint{0.650001in}{0.565123in}}{\pgfqpoint{5.499999in}{2.784877in}}%
\pgfusepath{clip}%
\pgfsetbuttcap%
\pgfsetroundjoin%
\pgfsetlinewidth{0.803000pt}%
\definecolor{currentstroke}{rgb}{0.690196,0.690196,0.690196}%
\pgfsetstrokecolor{currentstroke}%
\pgfsetdash{{0.800000pt}{1.320000pt}}{0.000000pt}%
\pgfpathmoveto{\pgfqpoint{0.650001in}{2.090470in}}%
\pgfpathlineto{\pgfqpoint{6.150000in}{2.090470in}}%
\pgfusepath{stroke}%
\end{pgfscope}%
\begin{pgfscope}%
\pgfsetbuttcap%
\pgfsetroundjoin%
\definecolor{currentfill}{rgb}{0.000000,0.000000,0.000000}%
\pgfsetfillcolor{currentfill}%
\pgfsetlinewidth{0.803000pt}%
\definecolor{currentstroke}{rgb}{0.000000,0.000000,0.000000}%
\pgfsetstrokecolor{currentstroke}%
\pgfsetdash{}{0pt}%
\pgfsys@defobject{currentmarker}{\pgfqpoint{-0.048611in}{0.000000in}}{\pgfqpoint{-0.000000in}{0.000000in}}{%
\pgfpathmoveto{\pgfqpoint{-0.000000in}{0.000000in}}%
\pgfpathlineto{\pgfqpoint{-0.048611in}{0.000000in}}%
\pgfusepath{stroke,fill}%
}%
\begin{pgfscope}%
\pgfsys@transformshift{0.650001in}{2.090470in}%
\pgfsys@useobject{currentmarker}{}%
\end{pgfscope}%
\end{pgfscope}%
\begin{pgfscope}%
\definecolor{textcolor}{rgb}{0.000000,0.000000,0.000000}%
\pgfsetstrokecolor{textcolor}%
\pgfsetfillcolor{textcolor}%
\pgftext[x=0.344444in, y=2.042245in, left, base]{\color{textcolor}\rmfamily\fontsize{10.000000}{12.000000}\selectfont \(\displaystyle {100}\)}%
\end{pgfscope}%
\begin{pgfscope}%
\pgfpathrectangle{\pgfqpoint{0.650001in}{0.565123in}}{\pgfqpoint{5.499999in}{2.784877in}}%
\pgfusepath{clip}%
\pgfsetbuttcap%
\pgfsetroundjoin%
\pgfsetlinewidth{0.803000pt}%
\definecolor{currentstroke}{rgb}{0.690196,0.690196,0.690196}%
\pgfsetstrokecolor{currentstroke}%
\pgfsetdash{{0.800000pt}{1.320000pt}}{0.000000pt}%
\pgfpathmoveto{\pgfqpoint{0.650001in}{2.444515in}}%
\pgfpathlineto{\pgfqpoint{6.150000in}{2.444515in}}%
\pgfusepath{stroke}%
\end{pgfscope}%
\begin{pgfscope}%
\pgfsetbuttcap%
\pgfsetroundjoin%
\definecolor{currentfill}{rgb}{0.000000,0.000000,0.000000}%
\pgfsetfillcolor{currentfill}%
\pgfsetlinewidth{0.803000pt}%
\definecolor{currentstroke}{rgb}{0.000000,0.000000,0.000000}%
\pgfsetstrokecolor{currentstroke}%
\pgfsetdash{}{0pt}%
\pgfsys@defobject{currentmarker}{\pgfqpoint{-0.048611in}{0.000000in}}{\pgfqpoint{-0.000000in}{0.000000in}}{%
\pgfpathmoveto{\pgfqpoint{-0.000000in}{0.000000in}}%
\pgfpathlineto{\pgfqpoint{-0.048611in}{0.000000in}}%
\pgfusepath{stroke,fill}%
}%
\begin{pgfscope}%
\pgfsys@transformshift{0.650001in}{2.444515in}%
\pgfsys@useobject{currentmarker}{}%
\end{pgfscope}%
\end{pgfscope}%
\begin{pgfscope}%
\definecolor{textcolor}{rgb}{0.000000,0.000000,0.000000}%
\pgfsetstrokecolor{textcolor}%
\pgfsetfillcolor{textcolor}%
\pgftext[x=0.344444in, y=2.396290in, left, base]{\color{textcolor}\rmfamily\fontsize{10.000000}{12.000000}\selectfont \(\displaystyle {125}\)}%
\end{pgfscope}%
\begin{pgfscope}%
\pgfpathrectangle{\pgfqpoint{0.650001in}{0.565123in}}{\pgfqpoint{5.499999in}{2.784877in}}%
\pgfusepath{clip}%
\pgfsetbuttcap%
\pgfsetroundjoin%
\pgfsetlinewidth{0.803000pt}%
\definecolor{currentstroke}{rgb}{0.690196,0.690196,0.690196}%
\pgfsetstrokecolor{currentstroke}%
\pgfsetdash{{0.800000pt}{1.320000pt}}{0.000000pt}%
\pgfpathmoveto{\pgfqpoint{0.650001in}{2.798561in}}%
\pgfpathlineto{\pgfqpoint{6.150000in}{2.798561in}}%
\pgfusepath{stroke}%
\end{pgfscope}%
\begin{pgfscope}%
\pgfsetbuttcap%
\pgfsetroundjoin%
\definecolor{currentfill}{rgb}{0.000000,0.000000,0.000000}%
\pgfsetfillcolor{currentfill}%
\pgfsetlinewidth{0.803000pt}%
\definecolor{currentstroke}{rgb}{0.000000,0.000000,0.000000}%
\pgfsetstrokecolor{currentstroke}%
\pgfsetdash{}{0pt}%
\pgfsys@defobject{currentmarker}{\pgfqpoint{-0.048611in}{0.000000in}}{\pgfqpoint{-0.000000in}{0.000000in}}{%
\pgfpathmoveto{\pgfqpoint{-0.000000in}{0.000000in}}%
\pgfpathlineto{\pgfqpoint{-0.048611in}{0.000000in}}%
\pgfusepath{stroke,fill}%
}%
\begin{pgfscope}%
\pgfsys@transformshift{0.650001in}{2.798561in}%
\pgfsys@useobject{currentmarker}{}%
\end{pgfscope}%
\end{pgfscope}%
\begin{pgfscope}%
\definecolor{textcolor}{rgb}{0.000000,0.000000,0.000000}%
\pgfsetstrokecolor{textcolor}%
\pgfsetfillcolor{textcolor}%
\pgftext[x=0.344444in, y=2.750335in, left, base]{\color{textcolor}\rmfamily\fontsize{10.000000}{12.000000}\selectfont \(\displaystyle {150}\)}%
\end{pgfscope}%
\begin{pgfscope}%
\pgfpathrectangle{\pgfqpoint{0.650001in}{0.565123in}}{\pgfqpoint{5.499999in}{2.784877in}}%
\pgfusepath{clip}%
\pgfsetbuttcap%
\pgfsetroundjoin%
\pgfsetlinewidth{0.803000pt}%
\definecolor{currentstroke}{rgb}{0.690196,0.690196,0.690196}%
\pgfsetstrokecolor{currentstroke}%
\pgfsetdash{{0.800000pt}{1.320000pt}}{0.000000pt}%
\pgfpathmoveto{\pgfqpoint{0.650001in}{3.152606in}}%
\pgfpathlineto{\pgfqpoint{6.150000in}{3.152606in}}%
\pgfusepath{stroke}%
\end{pgfscope}%
\begin{pgfscope}%
\pgfsetbuttcap%
\pgfsetroundjoin%
\definecolor{currentfill}{rgb}{0.000000,0.000000,0.000000}%
\pgfsetfillcolor{currentfill}%
\pgfsetlinewidth{0.803000pt}%
\definecolor{currentstroke}{rgb}{0.000000,0.000000,0.000000}%
\pgfsetstrokecolor{currentstroke}%
\pgfsetdash{}{0pt}%
\pgfsys@defobject{currentmarker}{\pgfqpoint{-0.048611in}{0.000000in}}{\pgfqpoint{-0.000000in}{0.000000in}}{%
\pgfpathmoveto{\pgfqpoint{-0.000000in}{0.000000in}}%
\pgfpathlineto{\pgfqpoint{-0.048611in}{0.000000in}}%
\pgfusepath{stroke,fill}%
}%
\begin{pgfscope}%
\pgfsys@transformshift{0.650001in}{3.152606in}%
\pgfsys@useobject{currentmarker}{}%
\end{pgfscope}%
\end{pgfscope}%
\begin{pgfscope}%
\definecolor{textcolor}{rgb}{0.000000,0.000000,0.000000}%
\pgfsetstrokecolor{textcolor}%
\pgfsetfillcolor{textcolor}%
\pgftext[x=0.344444in, y=3.104380in, left, base]{\color{textcolor}\rmfamily\fontsize{10.000000}{12.000000}\selectfont \(\displaystyle {175}\)}%
\end{pgfscope}%
\begin{pgfscope}%
\definecolor{textcolor}{rgb}{0.000000,0.000000,0.000000}%
\pgfsetstrokecolor{textcolor}%
\pgfsetfillcolor{textcolor}%
\pgftext[x=0.288889in,y=1.957562in,,bottom,rotate=90.000000]{\color{textcolor}\rmfamily\fontsize{10.000000}{12.000000}\selectfont Phase of Return Loss (deg)}%
\end{pgfscope}%
\begin{pgfscope}%
\pgfpathrectangle{\pgfqpoint{0.650001in}{0.565123in}}{\pgfqpoint{5.499999in}{2.784877in}}%
\pgfusepath{clip}%
\pgfsetrectcap%
\pgfsetroundjoin%
\pgfsetlinewidth{1.003750pt}%
\definecolor{currentstroke}{rgb}{0.000000,0.000000,0.000000}%
\pgfsetstrokecolor{currentstroke}%
\pgfsetdash{}{0pt}%
\pgfpathmoveto{\pgfqpoint{0.650001in}{1.576397in}}%
\pgfpathlineto{\pgfqpoint{0.655934in}{0.709128in}}%
\pgfpathlineto{\pgfqpoint{0.661867in}{0.716633in}}%
\pgfpathlineto{\pgfqpoint{0.679666in}{0.771440in}}%
\pgfpathlineto{\pgfqpoint{0.804262in}{1.187655in}}%
\pgfpathlineto{\pgfqpoint{0.810195in}{1.205074in}}%
\pgfpathlineto{\pgfqpoint{0.816128in}{1.227450in}}%
\pgfpathlineto{\pgfqpoint{0.827994in}{1.264837in}}%
\pgfpathlineto{\pgfqpoint{0.916991in}{1.566342in}}%
\pgfpathlineto{\pgfqpoint{0.928857in}{1.608686in}}%
\pgfpathlineto{\pgfqpoint{0.958523in}{1.710934in}}%
\pgfpathlineto{\pgfqpoint{0.970389in}{1.753278in}}%
\pgfpathlineto{\pgfqpoint{0.994121in}{1.835416in}}%
\pgfpathlineto{\pgfqpoint{1.017854in}{1.922653in}}%
\pgfpathlineto{\pgfqpoint{1.029720in}{1.964997in}}%
\pgfpathlineto{\pgfqpoint{1.053453in}{2.052233in}}%
\pgfpathlineto{\pgfqpoint{1.065319in}{2.094577in}}%
\pgfpathlineto{\pgfqpoint{1.219580in}{2.682717in}}%
\pgfpathlineto{\pgfqpoint{1.231446in}{2.730017in}}%
\pgfpathlineto{\pgfqpoint{1.243312in}{2.777318in}}%
\pgfpathlineto{\pgfqpoint{1.267045in}{2.874609in}}%
\pgfpathlineto{\pgfqpoint{1.278911in}{2.921910in}}%
\pgfpathlineto{\pgfqpoint{1.296710in}{2.994135in}}%
\pgfpathlineto{\pgfqpoint{1.320443in}{3.086328in}}%
\pgfpathlineto{\pgfqpoint{1.338242in}{3.163652in}}%
\pgfpathlineto{\pgfqpoint{1.350108in}{3.210952in}}%
\pgfpathlineto{\pgfqpoint{1.356042in}{3.223415in}}%
\pgfpathlineto{\pgfqpoint{1.379774in}{3.223415in}}%
\pgfpathlineto{\pgfqpoint{1.385707in}{3.220866in}}%
\pgfpathlineto{\pgfqpoint{1.415373in}{3.098791in}}%
\pgfpathlineto{\pgfqpoint{1.427239in}{3.051490in}}%
\pgfpathlineto{\pgfqpoint{1.450971in}{2.954340in}}%
\pgfpathlineto{\pgfqpoint{1.468771in}{2.884523in}}%
\pgfpathlineto{\pgfqpoint{1.492503in}{2.787373in}}%
\pgfpathlineto{\pgfqpoint{1.516236in}{2.695179in}}%
\pgfpathlineto{\pgfqpoint{1.528102in}{2.647737in}}%
\pgfpathlineto{\pgfqpoint{1.545901in}{2.578061in}}%
\pgfpathlineto{\pgfqpoint{1.557768in}{2.530761in}}%
\pgfpathlineto{\pgfqpoint{1.676430in}{2.084664in}}%
\pgfpathlineto{\pgfqpoint{1.688296in}{2.042320in}}%
\pgfpathlineto{\pgfqpoint{1.712029in}{1.955083in}}%
\pgfpathlineto{\pgfqpoint{1.729828in}{1.892771in}}%
\pgfpathlineto{\pgfqpoint{1.747627in}{1.828052in}}%
\pgfpathlineto{\pgfqpoint{1.801025in}{1.646073in}}%
\pgfpathlineto{\pgfqpoint{1.812892in}{1.603729in}}%
\pgfpathlineto{\pgfqpoint{1.872223in}{1.406880in}}%
\pgfpathlineto{\pgfqpoint{1.884089in}{1.369493in}}%
\pgfpathlineto{\pgfqpoint{1.913755in}{1.272343in}}%
\pgfpathlineto{\pgfqpoint{1.925621in}{1.234956in}}%
\pgfpathlineto{\pgfqpoint{1.943420in}{1.177600in}}%
\pgfpathlineto{\pgfqpoint{1.955286in}{1.140213in}}%
\pgfpathlineto{\pgfqpoint{1.973086in}{1.082999in}}%
\pgfpathlineto{\pgfqpoint{1.990885in}{1.028193in}}%
\pgfpathlineto{\pgfqpoint{2.008684in}{0.970838in}}%
\pgfpathlineto{\pgfqpoint{2.044283in}{0.863633in}}%
\pgfpathlineto{\pgfqpoint{2.056149in}{0.826246in}}%
\pgfpathlineto{\pgfqpoint{2.085815in}{0.741558in}}%
\pgfpathlineto{\pgfqpoint{2.097681in}{0.714226in}}%
\pgfpathlineto{\pgfqpoint{2.103614in}{0.704171in}}%
\pgfpathlineto{\pgfqpoint{2.115481in}{0.694258in}}%
\pgfpathlineto{\pgfqpoint{2.121414in}{0.691709in}}%
\pgfpathlineto{\pgfqpoint{2.133280in}{0.691709in}}%
\pgfpathlineto{\pgfqpoint{2.145146in}{0.696665in}}%
\pgfpathlineto{\pgfqpoint{2.157012in}{0.706720in}}%
\pgfpathlineto{\pgfqpoint{2.162945in}{0.716633in}}%
\pgfpathlineto{\pgfqpoint{2.180745in}{0.763934in}}%
\pgfpathlineto{\pgfqpoint{2.234143in}{0.923537in}}%
\pgfpathlineto{\pgfqpoint{2.251942in}{0.980751in}}%
\pgfpathlineto{\pgfqpoint{2.269742in}{1.035557in}}%
\pgfpathlineto{\pgfqpoint{2.287541in}{1.092913in}}%
\pgfpathlineto{\pgfqpoint{2.305340in}{1.147719in}}%
\pgfpathlineto{\pgfqpoint{2.329073in}{1.225042in}}%
\pgfpathlineto{\pgfqpoint{2.340939in}{1.262288in}}%
\pgfpathlineto{\pgfqpoint{2.388404in}{1.419342in}}%
\pgfpathlineto{\pgfqpoint{2.400270in}{1.456729in}}%
\pgfpathlineto{\pgfqpoint{2.489267in}{1.758234in}}%
\pgfpathlineto{\pgfqpoint{2.507066in}{1.822954in}}%
\pgfpathlineto{\pgfqpoint{2.524866in}{1.885266in}}%
\pgfpathlineto{\pgfqpoint{2.542665in}{1.950127in}}%
\pgfpathlineto{\pgfqpoint{2.554531in}{1.992471in}}%
\pgfpathlineto{\pgfqpoint{2.596063in}{2.146976in}}%
\pgfpathlineto{\pgfqpoint{2.607929in}{2.189320in}}%
\pgfpathlineto{\pgfqpoint{2.667260in}{2.416050in}}%
\pgfpathlineto{\pgfqpoint{2.679127in}{2.463351in}}%
\pgfpathlineto{\pgfqpoint{2.702859in}{2.555544in}}%
\pgfpathlineto{\pgfqpoint{2.714725in}{2.602986in}}%
\pgfpathlineto{\pgfqpoint{2.732525in}{2.672662in}}%
\pgfpathlineto{\pgfqpoint{2.803722in}{2.969210in}}%
\pgfpathlineto{\pgfqpoint{2.815588in}{3.016652in}}%
\pgfpathlineto{\pgfqpoint{2.863053in}{3.213501in}}%
\pgfpathlineto{\pgfqpoint{2.868986in}{3.223415in}}%
\pgfpathlineto{\pgfqpoint{2.892719in}{3.223415in}}%
\pgfpathlineto{\pgfqpoint{2.898652in}{3.210952in}}%
\pgfpathlineto{\pgfqpoint{2.910518in}{3.163652in}}%
\pgfpathlineto{\pgfqpoint{2.952050in}{2.991728in}}%
\pgfpathlineto{\pgfqpoint{2.963916in}{2.944286in}}%
\pgfpathlineto{\pgfqpoint{2.993582in}{2.822211in}}%
\pgfpathlineto{\pgfqpoint{3.005448in}{2.774910in}}%
\pgfpathlineto{\pgfqpoint{3.017314in}{2.727610in}}%
\pgfpathlineto{\pgfqpoint{3.035114in}{2.657792in}}%
\pgfpathlineto{\pgfqpoint{3.046980in}{2.610492in}}%
\pgfpathlineto{\pgfqpoint{3.076645in}{2.495781in}}%
\pgfpathlineto{\pgfqpoint{3.088512in}{2.448481in}}%
\pgfpathlineto{\pgfqpoint{3.147843in}{2.226707in}}%
\pgfpathlineto{\pgfqpoint{3.165642in}{2.164395in}}%
\pgfpathlineto{\pgfqpoint{3.189375in}{2.077158in}}%
\pgfpathlineto{\pgfqpoint{3.213107in}{1.995020in}}%
\pgfpathlineto{\pgfqpoint{3.224973in}{1.952534in}}%
\pgfpathlineto{\pgfqpoint{3.302104in}{1.690966in}}%
\pgfpathlineto{\pgfqpoint{3.355502in}{1.514085in}}%
\pgfpathlineto{\pgfqpoint{3.367368in}{1.476698in}}%
\pgfpathlineto{\pgfqpoint{3.385168in}{1.419342in}}%
\pgfpathlineto{\pgfqpoint{3.402967in}{1.364536in}}%
\pgfpathlineto{\pgfqpoint{3.414833in}{1.327149in}}%
\pgfpathlineto{\pgfqpoint{3.450432in}{1.219944in}}%
\pgfpathlineto{\pgfqpoint{3.462298in}{1.182557in}}%
\pgfpathlineto{\pgfqpoint{3.545362in}{0.936000in}}%
\pgfpathlineto{\pgfqpoint{3.557228in}{0.898613in}}%
\pgfpathlineto{\pgfqpoint{3.598760in}{0.778945in}}%
\pgfpathlineto{\pgfqpoint{3.622492in}{0.719183in}}%
\pgfpathlineto{\pgfqpoint{3.628425in}{0.709128in}}%
\pgfpathlineto{\pgfqpoint{3.634358in}{0.701764in}}%
\pgfpathlineto{\pgfqpoint{3.640292in}{0.696665in}}%
\pgfpathlineto{\pgfqpoint{3.652158in}{0.691709in}}%
\pgfpathlineto{\pgfqpoint{3.664024in}{0.691709in}}%
\pgfpathlineto{\pgfqpoint{3.675890in}{0.696665in}}%
\pgfpathlineto{\pgfqpoint{3.687757in}{0.706720in}}%
\pgfpathlineto{\pgfqpoint{3.693690in}{0.716633in}}%
\pgfpathlineto{\pgfqpoint{3.705556in}{0.744107in}}%
\pgfpathlineto{\pgfqpoint{3.723355in}{0.791408in}}%
\pgfpathlineto{\pgfqpoint{3.824218in}{1.090363in}}%
\pgfpathlineto{\pgfqpoint{3.836084in}{1.127751in}}%
\pgfpathlineto{\pgfqpoint{3.865750in}{1.217536in}}%
\pgfpathlineto{\pgfqpoint{3.877616in}{1.254924in}}%
\pgfpathlineto{\pgfqpoint{3.889482in}{1.292311in}}%
\pgfpathlineto{\pgfqpoint{3.907282in}{1.349525in}}%
\pgfpathlineto{\pgfqpoint{3.925081in}{1.404331in}}%
\pgfpathlineto{\pgfqpoint{3.954747in}{1.501622in}}%
\pgfpathlineto{\pgfqpoint{3.966613in}{1.539009in}}%
\pgfpathlineto{\pgfqpoint{4.043744in}{1.800578in}}%
\pgfpathlineto{\pgfqpoint{4.055610in}{1.842922in}}%
\pgfpathlineto{\pgfqpoint{4.079342in}{1.925202in}}%
\pgfpathlineto{\pgfqpoint{4.091208in}{1.967546in}}%
\pgfpathlineto{\pgfqpoint{4.103075in}{2.009890in}}%
\pgfpathlineto{\pgfqpoint{4.126807in}{2.097126in}}%
\pgfpathlineto{\pgfqpoint{4.138673in}{2.139470in}}%
\pgfpathlineto{\pgfqpoint{4.209871in}{2.411094in}}%
\pgfpathlineto{\pgfqpoint{4.233603in}{2.508243in}}%
\pgfpathlineto{\pgfqpoint{4.245469in}{2.555544in}}%
\pgfpathlineto{\pgfqpoint{4.292934in}{2.752393in}}%
\pgfpathlineto{\pgfqpoint{4.304801in}{2.799835in}}%
\pgfpathlineto{\pgfqpoint{4.405664in}{3.220866in}}%
\pgfpathlineto{\pgfqpoint{4.411597in}{3.223415in}}%
\pgfpathlineto{\pgfqpoint{4.435329in}{3.223415in}}%
\pgfpathlineto{\pgfqpoint{4.571791in}{2.652836in}}%
\pgfpathlineto{\pgfqpoint{4.583657in}{2.605393in}}%
\pgfpathlineto{\pgfqpoint{4.601457in}{2.533168in}}%
\pgfpathlineto{\pgfqpoint{4.619256in}{2.463351in}}%
\pgfpathlineto{\pgfqpoint{4.631122in}{2.416050in}}%
\pgfpathlineto{\pgfqpoint{4.773517in}{1.882858in}}%
\pgfpathlineto{\pgfqpoint{4.797249in}{1.800578in}}%
\pgfpathlineto{\pgfqpoint{4.809116in}{1.758234in}}%
\pgfpathlineto{\pgfqpoint{4.880313in}{1.521449in}}%
\pgfpathlineto{\pgfqpoint{4.898112in}{1.466643in}}%
\pgfpathlineto{\pgfqpoint{4.915912in}{1.409429in}}%
\pgfpathlineto{\pgfqpoint{4.927778in}{1.372042in}}%
\pgfpathlineto{\pgfqpoint{4.945577in}{1.314686in}}%
\pgfpathlineto{\pgfqpoint{4.963377in}{1.259880in}}%
\pgfpathlineto{\pgfqpoint{4.975243in}{1.222493in}}%
\pgfpathlineto{\pgfqpoint{4.993042in}{1.167687in}}%
\pgfpathlineto{\pgfqpoint{5.010842in}{1.110332in}}%
\pgfpathlineto{\pgfqpoint{5.034574in}{1.038106in}}%
\pgfpathlineto{\pgfqpoint{5.046440in}{1.000719in}}%
\pgfpathlineto{\pgfqpoint{5.093905in}{0.858676in}}%
\pgfpathlineto{\pgfqpoint{5.099838in}{0.838708in}}%
\pgfpathlineto{\pgfqpoint{5.111705in}{0.806419in}}%
\pgfpathlineto{\pgfqpoint{5.129504in}{0.756570in}}%
\pgfpathlineto{\pgfqpoint{5.147303in}{0.714226in}}%
\pgfpathlineto{\pgfqpoint{5.153236in}{0.704171in}}%
\pgfpathlineto{\pgfqpoint{5.165103in}{0.694258in}}%
\pgfpathlineto{\pgfqpoint{5.171036in}{0.691709in}}%
\pgfpathlineto{\pgfqpoint{5.182902in}{0.691709in}}%
\pgfpathlineto{\pgfqpoint{5.188835in}{0.694258in}}%
\pgfpathlineto{\pgfqpoint{5.194768in}{0.694258in}}%
\pgfpathlineto{\pgfqpoint{5.206634in}{0.704171in}}%
\pgfpathlineto{\pgfqpoint{5.218501in}{0.724139in}}%
\pgfpathlineto{\pgfqpoint{5.236300in}{0.771440in}}%
\pgfpathlineto{\pgfqpoint{5.319364in}{1.018138in}}%
\pgfpathlineto{\pgfqpoint{5.331230in}{1.055525in}}%
\pgfpathlineto{\pgfqpoint{5.354962in}{1.127751in}}%
\pgfpathlineto{\pgfqpoint{5.372762in}{1.185106in}}%
\pgfpathlineto{\pgfqpoint{5.396494in}{1.257331in}}%
\pgfpathlineto{\pgfqpoint{5.414294in}{1.314686in}}%
\pgfpathlineto{\pgfqpoint{5.426160in}{1.352074in}}%
\pgfpathlineto{\pgfqpoint{5.449892in}{1.429255in}}%
\pgfpathlineto{\pgfqpoint{5.461758in}{1.466643in}}%
\pgfpathlineto{\pgfqpoint{5.491424in}{1.563934in}}%
\pgfpathlineto{\pgfqpoint{5.503290in}{1.601180in}}%
\pgfpathlineto{\pgfqpoint{5.568555in}{1.822954in}}%
\pgfpathlineto{\pgfqpoint{5.580421in}{1.865439in}}%
\pgfpathlineto{\pgfqpoint{5.598220in}{1.927609in}}%
\pgfpathlineto{\pgfqpoint{5.716882in}{2.381212in}}%
\pgfpathlineto{\pgfqpoint{5.734682in}{2.453437in}}%
\pgfpathlineto{\pgfqpoint{5.752481in}{2.523255in}}%
\pgfpathlineto{\pgfqpoint{5.770281in}{2.595480in}}%
\pgfpathlineto{\pgfqpoint{5.788080in}{2.665298in}}%
\pgfpathlineto{\pgfqpoint{5.912675in}{3.186028in}}%
\pgfpathlineto{\pgfqpoint{5.918608in}{3.208403in}}%
\pgfpathlineto{\pgfqpoint{5.924542in}{3.223415in}}%
\pgfpathlineto{\pgfqpoint{5.948274in}{3.223415in}}%
\pgfpathlineto{\pgfqpoint{5.966073in}{3.161103in}}%
\pgfpathlineto{\pgfqpoint{6.061003in}{2.764855in}}%
\pgfpathlineto{\pgfqpoint{6.072869in}{2.717555in}}%
\pgfpathlineto{\pgfqpoint{6.096602in}{2.620405in}}%
\pgfpathlineto{\pgfqpoint{6.114401in}{2.550587in}}%
\pgfpathlineto{\pgfqpoint{6.132201in}{2.478362in}}%
\pgfpathlineto{\pgfqpoint{6.150000in}{2.408544in}}%
\pgfpathlineto{\pgfqpoint{6.150000in}{2.408544in}}%
\pgfusepath{stroke}%
\end{pgfscope}%
\begin{pgfscope}%
\pgfpathrectangle{\pgfqpoint{0.650001in}{0.565123in}}{\pgfqpoint{5.499999in}{2.784877in}}%
\pgfusepath{clip}%
\pgfsetbuttcap%
\pgfsetroundjoin%
\pgfsetlinewidth{1.003750pt}%
\definecolor{currentstroke}{rgb}{0.000000,0.000000,0.000000}%
\pgfsetstrokecolor{currentstroke}%
\pgfsetdash{{3.700000pt}{1.600000pt}}{0.000000pt}%
\pgfpathmoveto{\pgfqpoint{0.650001in}{1.496524in}}%
\pgfpathlineto{\pgfqpoint{0.655934in}{0.741558in}}%
\pgfpathlineto{\pgfqpoint{0.667800in}{0.793957in}}%
\pgfpathlineto{\pgfqpoint{0.697466in}{0.933451in}}%
\pgfpathlineto{\pgfqpoint{0.715265in}{1.018138in}}%
\pgfpathlineto{\pgfqpoint{0.810195in}{1.461686in}}%
\pgfpathlineto{\pgfqpoint{0.827994in}{1.546374in}}%
\pgfpathlineto{\pgfqpoint{0.875459in}{1.770697in}}%
\pgfpathlineto{\pgfqpoint{0.905125in}{1.915147in}}%
\pgfpathlineto{\pgfqpoint{0.928857in}{2.029858in}}%
\pgfpathlineto{\pgfqpoint{0.964456in}{2.204331in}}%
\pgfpathlineto{\pgfqpoint{0.982255in}{2.291426in}}%
\pgfpathlineto{\pgfqpoint{1.000055in}{2.378663in}}%
\pgfpathlineto{\pgfqpoint{1.154316in}{3.161103in}}%
\pgfpathlineto{\pgfqpoint{1.166182in}{3.220866in}}%
\pgfpathlineto{\pgfqpoint{1.172115in}{3.223415in}}%
\pgfpathlineto{\pgfqpoint{1.189914in}{3.223415in}}%
\pgfpathlineto{\pgfqpoint{1.195847in}{3.201039in}}%
\pgfpathlineto{\pgfqpoint{1.207714in}{3.141135in}}%
\pgfpathlineto{\pgfqpoint{1.373841in}{2.308987in}}%
\pgfpathlineto{\pgfqpoint{1.409440in}{2.139470in}}%
\pgfpathlineto{\pgfqpoint{1.433172in}{2.024901in}}%
\pgfpathlineto{\pgfqpoint{1.593366in}{1.287213in}}%
\pgfpathlineto{\pgfqpoint{1.611166in}{1.210031in}}%
\pgfpathlineto{\pgfqpoint{1.640831in}{1.077901in}}%
\pgfpathlineto{\pgfqpoint{1.664564in}{0.975794in}}%
\pgfpathlineto{\pgfqpoint{1.694229in}{0.843665in}}%
\pgfpathlineto{\pgfqpoint{1.717962in}{0.746515in}}%
\pgfpathlineto{\pgfqpoint{1.729828in}{0.709128in}}%
\pgfpathlineto{\pgfqpoint{1.735761in}{0.699214in}}%
\pgfpathlineto{\pgfqpoint{1.741694in}{0.694258in}}%
\pgfpathlineto{\pgfqpoint{1.747627in}{0.691709in}}%
\pgfpathlineto{\pgfqpoint{1.753560in}{0.691709in}}%
\pgfpathlineto{\pgfqpoint{1.759494in}{0.694258in}}%
\pgfpathlineto{\pgfqpoint{1.765427in}{0.699214in}}%
\pgfpathlineto{\pgfqpoint{1.771360in}{0.711677in}}%
\pgfpathlineto{\pgfqpoint{1.783226in}{0.749064in}}%
\pgfpathlineto{\pgfqpoint{1.795092in}{0.796364in}}%
\pgfpathlineto{\pgfqpoint{1.812892in}{0.873688in}}%
\pgfpathlineto{\pgfqpoint{1.842557in}{1.005676in}}%
\pgfpathlineto{\pgfqpoint{1.860357in}{1.082999in}}%
\pgfpathlineto{\pgfqpoint{1.895955in}{1.242461in}}%
\pgfpathlineto{\pgfqpoint{1.907821in}{1.294718in}}%
\pgfpathlineto{\pgfqpoint{1.973086in}{1.593816in}}%
\pgfpathlineto{\pgfqpoint{1.984952in}{1.646073in}}%
\pgfpathlineto{\pgfqpoint{2.002751in}{1.730760in}}%
\pgfpathlineto{\pgfqpoint{2.056149in}{1.982557in}}%
\pgfpathlineto{\pgfqpoint{2.085815in}{2.127008in}}%
\pgfpathlineto{\pgfqpoint{2.109547in}{2.241577in}}%
\pgfpathlineto{\pgfqpoint{2.287541in}{3.146233in}}%
\pgfpathlineto{\pgfqpoint{2.299407in}{3.208403in}}%
\pgfpathlineto{\pgfqpoint{2.305340in}{3.223415in}}%
\pgfpathlineto{\pgfqpoint{2.323140in}{3.223415in}}%
\pgfpathlineto{\pgfqpoint{2.329073in}{3.210952in}}%
\pgfpathlineto{\pgfqpoint{2.346872in}{3.118759in}}%
\pgfpathlineto{\pgfqpoint{2.406203in}{2.814705in}}%
\pgfpathlineto{\pgfqpoint{2.424003in}{2.722512in}}%
\pgfpathlineto{\pgfqpoint{2.489267in}{2.398631in}}%
\pgfpathlineto{\pgfqpoint{2.518932in}{2.256588in}}%
\pgfpathlineto{\pgfqpoint{2.542665in}{2.142019in}}%
\pgfpathlineto{\pgfqpoint{2.649461in}{1.651029in}}%
\pgfpathlineto{\pgfqpoint{2.661327in}{1.598772in}}%
\pgfpathlineto{\pgfqpoint{2.690993in}{1.466643in}}%
\pgfpathlineto{\pgfqpoint{2.708792in}{1.389461in}}%
\pgfpathlineto{\pgfqpoint{2.738458in}{1.257331in}}%
\pgfpathlineto{\pgfqpoint{2.756257in}{1.180149in}}%
\pgfpathlineto{\pgfqpoint{2.821521in}{0.903569in}}%
\pgfpathlineto{\pgfqpoint{2.839321in}{0.826246in}}%
\pgfpathlineto{\pgfqpoint{2.857120in}{0.754021in}}%
\pgfpathlineto{\pgfqpoint{2.863053in}{0.731645in}}%
\pgfpathlineto{\pgfqpoint{2.868986in}{0.714226in}}%
\pgfpathlineto{\pgfqpoint{2.874920in}{0.701764in}}%
\pgfpathlineto{\pgfqpoint{2.880853in}{0.694258in}}%
\pgfpathlineto{\pgfqpoint{2.886786in}{0.691709in}}%
\pgfpathlineto{\pgfqpoint{2.892719in}{0.691709in}}%
\pgfpathlineto{\pgfqpoint{2.904585in}{0.696665in}}%
\pgfpathlineto{\pgfqpoint{2.910518in}{0.704171in}}%
\pgfpathlineto{\pgfqpoint{2.922384in}{0.736602in}}%
\pgfpathlineto{\pgfqpoint{2.934251in}{0.783902in}}%
\pgfpathlineto{\pgfqpoint{2.981716in}{0.985708in}}%
\pgfpathlineto{\pgfqpoint{3.005448in}{1.090363in}}%
\pgfpathlineto{\pgfqpoint{3.023247in}{1.167687in}}%
\pgfpathlineto{\pgfqpoint{3.058846in}{1.327149in}}%
\pgfpathlineto{\pgfqpoint{3.070712in}{1.379406in}}%
\pgfpathlineto{\pgfqpoint{3.195308in}{1.960040in}}%
\pgfpathlineto{\pgfqpoint{3.236840in}{2.164395in}}%
\pgfpathlineto{\pgfqpoint{3.254639in}{2.251632in}}%
\pgfpathlineto{\pgfqpoint{3.343636in}{2.705093in}}%
\pgfpathlineto{\pgfqpoint{3.379234in}{2.894436in}}%
\pgfpathlineto{\pgfqpoint{3.414833in}{3.078823in}}%
\pgfpathlineto{\pgfqpoint{3.438566in}{3.203447in}}%
\pgfpathlineto{\pgfqpoint{3.444499in}{3.223415in}}%
\pgfpathlineto{\pgfqpoint{3.462298in}{3.223415in}}%
\pgfpathlineto{\pgfqpoint{3.468231in}{3.210952in}}%
\pgfpathlineto{\pgfqpoint{3.497897in}{3.056447in}}%
\pgfpathlineto{\pgfqpoint{3.521629in}{2.931823in}}%
\pgfpathlineto{\pgfqpoint{3.669957in}{2.189320in}}%
\pgfpathlineto{\pgfqpoint{3.747088in}{1.828052in}}%
\pgfpathlineto{\pgfqpoint{3.758954in}{1.770697in}}%
\pgfpathlineto{\pgfqpoint{3.764887in}{1.745772in}}%
\pgfpathlineto{\pgfqpoint{3.776753in}{1.690966in}}%
\pgfpathlineto{\pgfqpoint{3.812352in}{1.528955in}}%
\pgfpathlineto{\pgfqpoint{3.830151in}{1.451773in}}%
\pgfpathlineto{\pgfqpoint{3.859817in}{1.319643in}}%
\pgfpathlineto{\pgfqpoint{3.895416in}{1.167687in}}%
\pgfpathlineto{\pgfqpoint{3.919148in}{1.063031in}}%
\pgfpathlineto{\pgfqpoint{3.948814in}{0.936000in}}%
\pgfpathlineto{\pgfqpoint{3.966613in}{0.858676in}}%
\pgfpathlineto{\pgfqpoint{3.990345in}{0.758977in}}%
\pgfpathlineto{\pgfqpoint{4.002212in}{0.719183in}}%
\pgfpathlineto{\pgfqpoint{4.008145in}{0.704171in}}%
\pgfpathlineto{\pgfqpoint{4.014078in}{0.696665in}}%
\pgfpathlineto{\pgfqpoint{4.020011in}{0.691709in}}%
\pgfpathlineto{\pgfqpoint{4.031877in}{0.691709in}}%
\pgfpathlineto{\pgfqpoint{4.043744in}{0.701764in}}%
\pgfpathlineto{\pgfqpoint{4.049677in}{0.714226in}}%
\pgfpathlineto{\pgfqpoint{4.061543in}{0.756570in}}%
\pgfpathlineto{\pgfqpoint{4.073409in}{0.803870in}}%
\pgfpathlineto{\pgfqpoint{4.097142in}{0.905977in}}%
\pgfpathlineto{\pgfqpoint{4.126807in}{1.038106in}}%
\pgfpathlineto{\pgfqpoint{4.144607in}{1.115288in}}%
\pgfpathlineto{\pgfqpoint{4.174272in}{1.247418in}}%
\pgfpathlineto{\pgfqpoint{4.310734in}{1.880309in}}%
\pgfpathlineto{\pgfqpoint{4.346332in}{2.054782in}}%
\pgfpathlineto{\pgfqpoint{4.358199in}{2.111996in}}%
\pgfpathlineto{\pgfqpoint{4.464995in}{2.657792in}}%
\pgfpathlineto{\pgfqpoint{4.488727in}{2.782416in}}%
\pgfpathlineto{\pgfqpoint{4.530259in}{2.996684in}}%
\pgfpathlineto{\pgfqpoint{4.559925in}{3.153597in}}%
\pgfpathlineto{\pgfqpoint{4.571791in}{3.213501in}}%
\pgfpathlineto{\pgfqpoint{4.577724in}{3.223415in}}%
\pgfpathlineto{\pgfqpoint{4.595523in}{3.223415in}}%
\pgfpathlineto{\pgfqpoint{4.601457in}{3.201039in}}%
\pgfpathlineto{\pgfqpoint{4.755718in}{2.426105in}}%
\pgfpathlineto{\pgfqpoint{4.904045in}{1.743223in}}%
\pgfpathlineto{\pgfqpoint{4.939644in}{1.591267in}}%
\pgfpathlineto{\pgfqpoint{4.957444in}{1.514085in}}%
\pgfpathlineto{\pgfqpoint{5.052373in}{1.112881in}}%
\pgfpathlineto{\pgfqpoint{5.070173in}{1.035557in}}%
\pgfpathlineto{\pgfqpoint{5.141370in}{0.739009in}}%
\pgfpathlineto{\pgfqpoint{5.147303in}{0.719183in}}%
\pgfpathlineto{\pgfqpoint{5.159169in}{0.696665in}}%
\pgfpathlineto{\pgfqpoint{5.171036in}{0.691709in}}%
\pgfpathlineto{\pgfqpoint{5.176969in}{0.691709in}}%
\pgfpathlineto{\pgfqpoint{5.188835in}{0.701764in}}%
\pgfpathlineto{\pgfqpoint{5.194768in}{0.714226in}}%
\pgfpathlineto{\pgfqpoint{5.200701in}{0.731645in}}%
\pgfpathlineto{\pgfqpoint{5.218501in}{0.801321in}}%
\pgfpathlineto{\pgfqpoint{5.242233in}{0.903569in}}%
\pgfpathlineto{\pgfqpoint{5.283765in}{1.090363in}}%
\pgfpathlineto{\pgfqpoint{5.295631in}{1.142762in}}%
\pgfpathlineto{\pgfqpoint{5.408360in}{1.666041in}}%
\pgfpathlineto{\pgfqpoint{5.449892in}{1.870396in}}%
\pgfpathlineto{\pgfqpoint{5.461758in}{1.927609in}}%
\pgfpathlineto{\pgfqpoint{5.503290in}{2.142019in}}%
\pgfpathlineto{\pgfqpoint{5.538889in}{2.331363in}}%
\pgfpathlineto{\pgfqpoint{5.556688in}{2.426105in}}%
\pgfpathlineto{\pgfqpoint{5.598220in}{2.647737in}}%
\pgfpathlineto{\pgfqpoint{5.616019in}{2.742480in}}%
\pgfpathlineto{\pgfqpoint{5.699083in}{3.190984in}}%
\pgfpathlineto{\pgfqpoint{5.705016in}{3.220866in}}%
\pgfpathlineto{\pgfqpoint{5.710949in}{3.223415in}}%
\pgfpathlineto{\pgfqpoint{5.728749in}{3.223415in}}%
\pgfpathlineto{\pgfqpoint{5.764347in}{3.034071in}}%
\pgfpathlineto{\pgfqpoint{5.900809in}{2.351331in}}%
\pgfpathlineto{\pgfqpoint{5.977940in}{1.997427in}}%
\pgfpathlineto{\pgfqpoint{6.007605in}{1.870396in}}%
\pgfpathlineto{\pgfqpoint{6.025405in}{1.793072in}}%
\pgfpathlineto{\pgfqpoint{6.090669in}{1.521449in}}%
\pgfpathlineto{\pgfqpoint{6.102535in}{1.474148in}}%
\pgfpathlineto{\pgfqpoint{6.150000in}{1.274750in}}%
\pgfpathlineto{\pgfqpoint{6.150000in}{1.274750in}}%
\pgfusepath{stroke}%
\end{pgfscope}%
\begin{pgfscope}%
\pgfsetrectcap%
\pgfsetmiterjoin%
\pgfsetlinewidth{0.803000pt}%
\definecolor{currentstroke}{rgb}{0.000000,0.000000,0.000000}%
\pgfsetstrokecolor{currentstroke}%
\pgfsetdash{}{0pt}%
\pgfpathmoveto{\pgfqpoint{0.650001in}{0.565123in}}%
\pgfpathlineto{\pgfqpoint{0.650001in}{3.350000in}}%
\pgfusepath{stroke}%
\end{pgfscope}%
\begin{pgfscope}%
\pgfsetrectcap%
\pgfsetmiterjoin%
\pgfsetlinewidth{0.803000pt}%
\definecolor{currentstroke}{rgb}{0.000000,0.000000,0.000000}%
\pgfsetstrokecolor{currentstroke}%
\pgfsetdash{}{0pt}%
\pgfpathmoveto{\pgfqpoint{6.150000in}{0.565123in}}%
\pgfpathlineto{\pgfqpoint{6.150000in}{3.350000in}}%
\pgfusepath{stroke}%
\end{pgfscope}%
\begin{pgfscope}%
\pgfsetrectcap%
\pgfsetmiterjoin%
\pgfsetlinewidth{0.803000pt}%
\definecolor{currentstroke}{rgb}{0.000000,0.000000,0.000000}%
\pgfsetstrokecolor{currentstroke}%
\pgfsetdash{}{0pt}%
\pgfpathmoveto{\pgfqpoint{0.650001in}{0.565123in}}%
\pgfpathlineto{\pgfqpoint{6.150000in}{0.565123in}}%
\pgfusepath{stroke}%
\end{pgfscope}%
\begin{pgfscope}%
\pgfsetrectcap%
\pgfsetmiterjoin%
\pgfsetlinewidth{0.803000pt}%
\definecolor{currentstroke}{rgb}{0.000000,0.000000,0.000000}%
\pgfsetstrokecolor{currentstroke}%
\pgfsetdash{}{0pt}%
\pgfpathmoveto{\pgfqpoint{0.650001in}{3.350000in}}%
\pgfpathlineto{\pgfqpoint{6.150000in}{3.350000in}}%
\pgfusepath{stroke}%
\end{pgfscope}%
\begin{pgfscope}%
\pgfsetbuttcap%
\pgfsetmiterjoin%
\definecolor{currentfill}{rgb}{1.000000,1.000000,1.000000}%
\pgfsetfillcolor{currentfill}%
\pgfsetfillopacity{0.800000}%
\pgfsetlinewidth{1.003750pt}%
\definecolor{currentstroke}{rgb}{0.800000,0.800000,0.800000}%
\pgfsetstrokecolor{currentstroke}%
\pgfsetstrokeopacity{0.800000}%
\pgfsetdash{}{0pt}%
\pgfpathmoveto{\pgfqpoint{4.013038in}{1.743056in}}%
\pgfpathlineto{\pgfqpoint{6.052778in}{1.743056in}}%
\pgfpathquadraticcurveto{\pgfqpoint{6.080556in}{1.743056in}}{\pgfqpoint{6.080556in}{1.770833in}}%
\pgfpathlineto{\pgfqpoint{6.080556in}{2.144290in}}%
\pgfpathquadraticcurveto{\pgfqpoint{6.080556in}{2.172068in}}{\pgfqpoint{6.052778in}{2.172068in}}%
\pgfpathlineto{\pgfqpoint{4.013038in}{2.172068in}}%
\pgfpathquadraticcurveto{\pgfqpoint{3.985260in}{2.172068in}}{\pgfqpoint{3.985260in}{2.144290in}}%
\pgfpathlineto{\pgfqpoint{3.985260in}{1.770833in}}%
\pgfpathquadraticcurveto{\pgfqpoint{3.985260in}{1.743056in}}{\pgfqpoint{4.013038in}{1.743056in}}%
\pgfpathlineto{\pgfqpoint{4.013038in}{1.743056in}}%
\pgfpathclose%
\pgfusepath{stroke,fill}%
\end{pgfscope}%
\begin{pgfscope}%
\pgfsetrectcap%
\pgfsetroundjoin%
\pgfsetlinewidth{1.003750pt}%
\definecolor{currentstroke}{rgb}{0.000000,0.000000,0.000000}%
\pgfsetstrokecolor{currentstroke}%
\pgfsetdash{}{0pt}%
\pgfpathmoveto{\pgfqpoint{4.040816in}{2.067901in}}%
\pgfpathlineto{\pgfqpoint{4.179704in}{2.067901in}}%
\pgfpathlineto{\pgfqpoint{4.318593in}{2.067901in}}%
\pgfusepath{stroke}%
\end{pgfscope}%
\begin{pgfscope}%
\definecolor{textcolor}{rgb}{0.000000,0.000000,0.000000}%
\pgfsetstrokecolor{textcolor}%
\pgfsetfillcolor{textcolor}%
\pgftext[x=4.429704in,y=2.019290in,left,base]{\color{textcolor}\rmfamily\fontsize{10.000000}{12.000000}\selectfont Cable 1 – Measured phase}%
\end{pgfscope}%
\begin{pgfscope}%
\pgfsetbuttcap%
\pgfsetroundjoin%
\pgfsetlinewidth{1.003750pt}%
\definecolor{currentstroke}{rgb}{0.000000,0.000000,0.000000}%
\pgfsetstrokecolor{currentstroke}%
\pgfsetdash{{3.700000pt}{1.600000pt}}{0.000000pt}%
\pgfpathmoveto{\pgfqpoint{4.040816in}{1.874228in}}%
\pgfpathlineto{\pgfqpoint{4.179704in}{1.874228in}}%
\pgfpathlineto{\pgfqpoint{4.318593in}{1.874228in}}%
\pgfusepath{stroke}%
\end{pgfscope}%
\begin{pgfscope}%
\definecolor{textcolor}{rgb}{0.000000,0.000000,0.000000}%
\pgfsetstrokecolor{textcolor}%
\pgfsetfillcolor{textcolor}%
\pgftext[x=4.429704in,y=1.825617in,left,base]{\color{textcolor}\rmfamily\fontsize{10.000000}{12.000000}\selectfont Cable 2 – Measured phase}%
\end{pgfscope}%
\end{pgfpicture}%
\makeatother%
\endgroup%

		\caption{Measured return loss and phase of a quartz resonatoras a function of frequency.}
		\label{fig:356}
	\end{figure}
	From the formula of the refection coefficient $r = r_L \cdot e^{-j2\beta l}$ and the fact that we can distinctly have maxima at a phase of 180 degree as seen in fig. \ref{fig:356}, $-j2\beta l$ is set to $\pi$. in addition it is known that $\beta = \frac{2\pi}{\lambda}$ which results in a length of $l = \frac{\lambda}{4}$. After considering a velocity factor of 0.66 and the relation $\lambda_0 = \frac{c_0}{f}$ to connect the frequency to the wave length the resulting equation to calculate the lengths of the cables is $l = \frac{c_0 \cdot VF}{4 \cdot f}$. 
	Cable 1 has  the first peak at 24 339 327 Hz leading to a length of 2.7 m. Cable 2 has a peak at 18 359 836 Hz thus a length of 2.03 m.

	
	\subsubsection{}
	Figure \ref{fig:357} shows the return loss of a matching network. A resonance can be observed at 14.238 MHz with a loss of -25.51 dB. The SWR is defined as $\frac{1+\rho}{1-\rho}$. With the measured return loss the SWR is 1.11.
	\begin{figure}[H]
		\centering
		%% Creator: Matplotlib, PGF backend
%%
%% To include the figure in your LaTeX document, write
%%   \input{<filename>.pgf}
%%
%% Make sure the required packages are loaded in your preamble
%%   \usepackage{pgf}
%%
%% Also ensure that all the required font packages are loaded; for instance,
%% the lmodern package is sometimes necessary when using math font.
%%   \usepackage{lmodern}
%%
%% Figures using additional raster images can only be included by \input if
%% they are in the same directory as the main LaTeX file. For loading figures
%% from other directories you can use the `import` package
%%   \usepackage{import}
%%
%% and then include the figures with
%%   \import{<path to file>}{<filename>.pgf}
%%
%% Matplotlib used the following preamble
%%
\begingroup%
\makeatletter%
\begin{pgfpicture}%
\pgfpathrectangle{\pgfpointorigin}{\pgfqpoint{6.300000in}{3.500000in}}%
\pgfusepath{use as bounding box, clip}%
\begin{pgfscope}%
\pgfsetbuttcap%
\pgfsetmiterjoin%
\definecolor{currentfill}{rgb}{1.000000,1.000000,1.000000}%
\pgfsetfillcolor{currentfill}%
\pgfsetlinewidth{0.000000pt}%
\definecolor{currentstroke}{rgb}{1.000000,1.000000,1.000000}%
\pgfsetstrokecolor{currentstroke}%
\pgfsetdash{}{0pt}%
\pgfpathmoveto{\pgfqpoint{0.000000in}{0.000000in}}%
\pgfpathlineto{\pgfqpoint{6.300000in}{0.000000in}}%
\pgfpathlineto{\pgfqpoint{6.300000in}{3.500000in}}%
\pgfpathlineto{\pgfqpoint{0.000000in}{3.500000in}}%
\pgfpathlineto{\pgfqpoint{0.000000in}{0.000000in}}%
\pgfpathclose%
\pgfusepath{fill}%
\end{pgfscope}%
\begin{pgfscope}%
\pgfsetbuttcap%
\pgfsetmiterjoin%
\definecolor{currentfill}{rgb}{1.000000,1.000000,1.000000}%
\pgfsetfillcolor{currentfill}%
\pgfsetlinewidth{0.000000pt}%
\definecolor{currentstroke}{rgb}{0.000000,0.000000,0.000000}%
\pgfsetstrokecolor{currentstroke}%
\pgfsetstrokeopacity{0.000000}%
\pgfsetdash{}{0pt}%
\pgfpathmoveto{\pgfqpoint{0.494137in}{0.565123in}}%
\pgfpathlineto{\pgfqpoint{6.150000in}{0.565123in}}%
\pgfpathlineto{\pgfqpoint{6.150000in}{3.150926in}}%
\pgfpathlineto{\pgfqpoint{0.494137in}{3.150926in}}%
\pgfpathlineto{\pgfqpoint{0.494137in}{0.565123in}}%
\pgfpathclose%
\pgfusepath{fill}%
\end{pgfscope}%
\begin{pgfscope}%
\pgfpathrectangle{\pgfqpoint{0.494137in}{0.565123in}}{\pgfqpoint{5.655863in}{2.585803in}}%
\pgfusepath{clip}%
\pgfsetbuttcap%
\pgfsetroundjoin%
\pgfsetlinewidth{0.803000pt}%
\definecolor{currentstroke}{rgb}{0.690196,0.690196,0.690196}%
\pgfsetstrokecolor{currentstroke}%
\pgfsetdash{{0.800000pt}{1.320000pt}}{0.000000pt}%
\pgfpathmoveto{\pgfqpoint{0.619238in}{0.565123in}}%
\pgfpathlineto{\pgfqpoint{0.619238in}{3.150926in}}%
\pgfusepath{stroke}%
\end{pgfscope}%
\begin{pgfscope}%
\pgfsetbuttcap%
\pgfsetroundjoin%
\definecolor{currentfill}{rgb}{0.000000,0.000000,0.000000}%
\pgfsetfillcolor{currentfill}%
\pgfsetlinewidth{0.803000pt}%
\definecolor{currentstroke}{rgb}{0.000000,0.000000,0.000000}%
\pgfsetstrokecolor{currentstroke}%
\pgfsetdash{}{0pt}%
\pgfsys@defobject{currentmarker}{\pgfqpoint{0.000000in}{-0.048611in}}{\pgfqpoint{0.000000in}{0.000000in}}{%
\pgfpathmoveto{\pgfqpoint{0.000000in}{0.000000in}}%
\pgfpathlineto{\pgfqpoint{0.000000in}{-0.048611in}}%
\pgfusepath{stroke,fill}%
}%
\begin{pgfscope}%
\pgfsys@transformshift{0.619238in}{0.565123in}%
\pgfsys@useobject{currentmarker}{}%
\end{pgfscope}%
\end{pgfscope}%
\begin{pgfscope}%
\definecolor{textcolor}{rgb}{0.000000,0.000000,0.000000}%
\pgfsetstrokecolor{textcolor}%
\pgfsetfillcolor{textcolor}%
\pgftext[x=0.619238in,y=0.467901in,,top]{\color{textcolor}\rmfamily\fontsize{10.000000}{12.000000}\selectfont \(\displaystyle {0.0}\)}%
\end{pgfscope}%
\begin{pgfscope}%
\pgfpathrectangle{\pgfqpoint{0.494137in}{0.565123in}}{\pgfqpoint{5.655863in}{2.585803in}}%
\pgfusepath{clip}%
\pgfsetbuttcap%
\pgfsetroundjoin%
\pgfsetlinewidth{0.803000pt}%
\definecolor{currentstroke}{rgb}{0.690196,0.690196,0.690196}%
\pgfsetstrokecolor{currentstroke}%
\pgfsetdash{{0.800000pt}{1.320000pt}}{0.000000pt}%
\pgfpathmoveto{\pgfqpoint{1.279154in}{0.565123in}}%
\pgfpathlineto{\pgfqpoint{1.279154in}{3.150926in}}%
\pgfusepath{stroke}%
\end{pgfscope}%
\begin{pgfscope}%
\pgfsetbuttcap%
\pgfsetroundjoin%
\definecolor{currentfill}{rgb}{0.000000,0.000000,0.000000}%
\pgfsetfillcolor{currentfill}%
\pgfsetlinewidth{0.803000pt}%
\definecolor{currentstroke}{rgb}{0.000000,0.000000,0.000000}%
\pgfsetstrokecolor{currentstroke}%
\pgfsetdash{}{0pt}%
\pgfsys@defobject{currentmarker}{\pgfqpoint{0.000000in}{-0.048611in}}{\pgfqpoint{0.000000in}{0.000000in}}{%
\pgfpathmoveto{\pgfqpoint{0.000000in}{0.000000in}}%
\pgfpathlineto{\pgfqpoint{0.000000in}{-0.048611in}}%
\pgfusepath{stroke,fill}%
}%
\begin{pgfscope}%
\pgfsys@transformshift{1.279154in}{0.565123in}%
\pgfsys@useobject{currentmarker}{}%
\end{pgfscope}%
\end{pgfscope}%
\begin{pgfscope}%
\definecolor{textcolor}{rgb}{0.000000,0.000000,0.000000}%
\pgfsetstrokecolor{textcolor}%
\pgfsetfillcolor{textcolor}%
\pgftext[x=1.279154in,y=0.467901in,,top]{\color{textcolor}\rmfamily\fontsize{10.000000}{12.000000}\selectfont \(\displaystyle {0.5}\)}%
\end{pgfscope}%
\begin{pgfscope}%
\pgfpathrectangle{\pgfqpoint{0.494137in}{0.565123in}}{\pgfqpoint{5.655863in}{2.585803in}}%
\pgfusepath{clip}%
\pgfsetbuttcap%
\pgfsetroundjoin%
\pgfsetlinewidth{0.803000pt}%
\definecolor{currentstroke}{rgb}{0.690196,0.690196,0.690196}%
\pgfsetstrokecolor{currentstroke}%
\pgfsetdash{{0.800000pt}{1.320000pt}}{0.000000pt}%
\pgfpathmoveto{\pgfqpoint{1.939070in}{0.565123in}}%
\pgfpathlineto{\pgfqpoint{1.939070in}{3.150926in}}%
\pgfusepath{stroke}%
\end{pgfscope}%
\begin{pgfscope}%
\pgfsetbuttcap%
\pgfsetroundjoin%
\definecolor{currentfill}{rgb}{0.000000,0.000000,0.000000}%
\pgfsetfillcolor{currentfill}%
\pgfsetlinewidth{0.803000pt}%
\definecolor{currentstroke}{rgb}{0.000000,0.000000,0.000000}%
\pgfsetstrokecolor{currentstroke}%
\pgfsetdash{}{0pt}%
\pgfsys@defobject{currentmarker}{\pgfqpoint{0.000000in}{-0.048611in}}{\pgfqpoint{0.000000in}{0.000000in}}{%
\pgfpathmoveto{\pgfqpoint{0.000000in}{0.000000in}}%
\pgfpathlineto{\pgfqpoint{0.000000in}{-0.048611in}}%
\pgfusepath{stroke,fill}%
}%
\begin{pgfscope}%
\pgfsys@transformshift{1.939070in}{0.565123in}%
\pgfsys@useobject{currentmarker}{}%
\end{pgfscope}%
\end{pgfscope}%
\begin{pgfscope}%
\definecolor{textcolor}{rgb}{0.000000,0.000000,0.000000}%
\pgfsetstrokecolor{textcolor}%
\pgfsetfillcolor{textcolor}%
\pgftext[x=1.939070in,y=0.467901in,,top]{\color{textcolor}\rmfamily\fontsize{10.000000}{12.000000}\selectfont \(\displaystyle {1.0}\)}%
\end{pgfscope}%
\begin{pgfscope}%
\pgfpathrectangle{\pgfqpoint{0.494137in}{0.565123in}}{\pgfqpoint{5.655863in}{2.585803in}}%
\pgfusepath{clip}%
\pgfsetbuttcap%
\pgfsetroundjoin%
\pgfsetlinewidth{0.803000pt}%
\definecolor{currentstroke}{rgb}{0.690196,0.690196,0.690196}%
\pgfsetstrokecolor{currentstroke}%
\pgfsetdash{{0.800000pt}{1.320000pt}}{0.000000pt}%
\pgfpathmoveto{\pgfqpoint{2.598987in}{0.565123in}}%
\pgfpathlineto{\pgfqpoint{2.598987in}{3.150926in}}%
\pgfusepath{stroke}%
\end{pgfscope}%
\begin{pgfscope}%
\pgfsetbuttcap%
\pgfsetroundjoin%
\definecolor{currentfill}{rgb}{0.000000,0.000000,0.000000}%
\pgfsetfillcolor{currentfill}%
\pgfsetlinewidth{0.803000pt}%
\definecolor{currentstroke}{rgb}{0.000000,0.000000,0.000000}%
\pgfsetstrokecolor{currentstroke}%
\pgfsetdash{}{0pt}%
\pgfsys@defobject{currentmarker}{\pgfqpoint{0.000000in}{-0.048611in}}{\pgfqpoint{0.000000in}{0.000000in}}{%
\pgfpathmoveto{\pgfqpoint{0.000000in}{0.000000in}}%
\pgfpathlineto{\pgfqpoint{0.000000in}{-0.048611in}}%
\pgfusepath{stroke,fill}%
}%
\begin{pgfscope}%
\pgfsys@transformshift{2.598987in}{0.565123in}%
\pgfsys@useobject{currentmarker}{}%
\end{pgfscope}%
\end{pgfscope}%
\begin{pgfscope}%
\definecolor{textcolor}{rgb}{0.000000,0.000000,0.000000}%
\pgfsetstrokecolor{textcolor}%
\pgfsetfillcolor{textcolor}%
\pgftext[x=2.598987in,y=0.467901in,,top]{\color{textcolor}\rmfamily\fontsize{10.000000}{12.000000}\selectfont \(\displaystyle {1.5}\)}%
\end{pgfscope}%
\begin{pgfscope}%
\pgfpathrectangle{\pgfqpoint{0.494137in}{0.565123in}}{\pgfqpoint{5.655863in}{2.585803in}}%
\pgfusepath{clip}%
\pgfsetbuttcap%
\pgfsetroundjoin%
\pgfsetlinewidth{0.803000pt}%
\definecolor{currentstroke}{rgb}{0.690196,0.690196,0.690196}%
\pgfsetstrokecolor{currentstroke}%
\pgfsetdash{{0.800000pt}{1.320000pt}}{0.000000pt}%
\pgfpathmoveto{\pgfqpoint{3.258903in}{0.565123in}}%
\pgfpathlineto{\pgfqpoint{3.258903in}{3.150926in}}%
\pgfusepath{stroke}%
\end{pgfscope}%
\begin{pgfscope}%
\pgfsetbuttcap%
\pgfsetroundjoin%
\definecolor{currentfill}{rgb}{0.000000,0.000000,0.000000}%
\pgfsetfillcolor{currentfill}%
\pgfsetlinewidth{0.803000pt}%
\definecolor{currentstroke}{rgb}{0.000000,0.000000,0.000000}%
\pgfsetstrokecolor{currentstroke}%
\pgfsetdash{}{0pt}%
\pgfsys@defobject{currentmarker}{\pgfqpoint{0.000000in}{-0.048611in}}{\pgfqpoint{0.000000in}{0.000000in}}{%
\pgfpathmoveto{\pgfqpoint{0.000000in}{0.000000in}}%
\pgfpathlineto{\pgfqpoint{0.000000in}{-0.048611in}}%
\pgfusepath{stroke,fill}%
}%
\begin{pgfscope}%
\pgfsys@transformshift{3.258903in}{0.565123in}%
\pgfsys@useobject{currentmarker}{}%
\end{pgfscope}%
\end{pgfscope}%
\begin{pgfscope}%
\definecolor{textcolor}{rgb}{0.000000,0.000000,0.000000}%
\pgfsetstrokecolor{textcolor}%
\pgfsetfillcolor{textcolor}%
\pgftext[x=3.258903in,y=0.467901in,,top]{\color{textcolor}\rmfamily\fontsize{10.000000}{12.000000}\selectfont \(\displaystyle {2.0}\)}%
\end{pgfscope}%
\begin{pgfscope}%
\pgfpathrectangle{\pgfqpoint{0.494137in}{0.565123in}}{\pgfqpoint{5.655863in}{2.585803in}}%
\pgfusepath{clip}%
\pgfsetbuttcap%
\pgfsetroundjoin%
\pgfsetlinewidth{0.803000pt}%
\definecolor{currentstroke}{rgb}{0.690196,0.690196,0.690196}%
\pgfsetstrokecolor{currentstroke}%
\pgfsetdash{{0.800000pt}{1.320000pt}}{0.000000pt}%
\pgfpathmoveto{\pgfqpoint{3.918819in}{0.565123in}}%
\pgfpathlineto{\pgfqpoint{3.918819in}{3.150926in}}%
\pgfusepath{stroke}%
\end{pgfscope}%
\begin{pgfscope}%
\pgfsetbuttcap%
\pgfsetroundjoin%
\definecolor{currentfill}{rgb}{0.000000,0.000000,0.000000}%
\pgfsetfillcolor{currentfill}%
\pgfsetlinewidth{0.803000pt}%
\definecolor{currentstroke}{rgb}{0.000000,0.000000,0.000000}%
\pgfsetstrokecolor{currentstroke}%
\pgfsetdash{}{0pt}%
\pgfsys@defobject{currentmarker}{\pgfqpoint{0.000000in}{-0.048611in}}{\pgfqpoint{0.000000in}{0.000000in}}{%
\pgfpathmoveto{\pgfqpoint{0.000000in}{0.000000in}}%
\pgfpathlineto{\pgfqpoint{0.000000in}{-0.048611in}}%
\pgfusepath{stroke,fill}%
}%
\begin{pgfscope}%
\pgfsys@transformshift{3.918819in}{0.565123in}%
\pgfsys@useobject{currentmarker}{}%
\end{pgfscope}%
\end{pgfscope}%
\begin{pgfscope}%
\definecolor{textcolor}{rgb}{0.000000,0.000000,0.000000}%
\pgfsetstrokecolor{textcolor}%
\pgfsetfillcolor{textcolor}%
\pgftext[x=3.918819in,y=0.467901in,,top]{\color{textcolor}\rmfamily\fontsize{10.000000}{12.000000}\selectfont \(\displaystyle {2.5}\)}%
\end{pgfscope}%
\begin{pgfscope}%
\pgfpathrectangle{\pgfqpoint{0.494137in}{0.565123in}}{\pgfqpoint{5.655863in}{2.585803in}}%
\pgfusepath{clip}%
\pgfsetbuttcap%
\pgfsetroundjoin%
\pgfsetlinewidth{0.803000pt}%
\definecolor{currentstroke}{rgb}{0.690196,0.690196,0.690196}%
\pgfsetstrokecolor{currentstroke}%
\pgfsetdash{{0.800000pt}{1.320000pt}}{0.000000pt}%
\pgfpathmoveto{\pgfqpoint{4.578735in}{0.565123in}}%
\pgfpathlineto{\pgfqpoint{4.578735in}{3.150926in}}%
\pgfusepath{stroke}%
\end{pgfscope}%
\begin{pgfscope}%
\pgfsetbuttcap%
\pgfsetroundjoin%
\definecolor{currentfill}{rgb}{0.000000,0.000000,0.000000}%
\pgfsetfillcolor{currentfill}%
\pgfsetlinewidth{0.803000pt}%
\definecolor{currentstroke}{rgb}{0.000000,0.000000,0.000000}%
\pgfsetstrokecolor{currentstroke}%
\pgfsetdash{}{0pt}%
\pgfsys@defobject{currentmarker}{\pgfqpoint{0.000000in}{-0.048611in}}{\pgfqpoint{0.000000in}{0.000000in}}{%
\pgfpathmoveto{\pgfqpoint{0.000000in}{0.000000in}}%
\pgfpathlineto{\pgfqpoint{0.000000in}{-0.048611in}}%
\pgfusepath{stroke,fill}%
}%
\begin{pgfscope}%
\pgfsys@transformshift{4.578735in}{0.565123in}%
\pgfsys@useobject{currentmarker}{}%
\end{pgfscope}%
\end{pgfscope}%
\begin{pgfscope}%
\definecolor{textcolor}{rgb}{0.000000,0.000000,0.000000}%
\pgfsetstrokecolor{textcolor}%
\pgfsetfillcolor{textcolor}%
\pgftext[x=4.578735in,y=0.467901in,,top]{\color{textcolor}\rmfamily\fontsize{10.000000}{12.000000}\selectfont \(\displaystyle {3.0}\)}%
\end{pgfscope}%
\begin{pgfscope}%
\pgfpathrectangle{\pgfqpoint{0.494137in}{0.565123in}}{\pgfqpoint{5.655863in}{2.585803in}}%
\pgfusepath{clip}%
\pgfsetbuttcap%
\pgfsetroundjoin%
\pgfsetlinewidth{0.803000pt}%
\definecolor{currentstroke}{rgb}{0.690196,0.690196,0.690196}%
\pgfsetstrokecolor{currentstroke}%
\pgfsetdash{{0.800000pt}{1.320000pt}}{0.000000pt}%
\pgfpathmoveto{\pgfqpoint{5.238651in}{0.565123in}}%
\pgfpathlineto{\pgfqpoint{5.238651in}{3.150926in}}%
\pgfusepath{stroke}%
\end{pgfscope}%
\begin{pgfscope}%
\pgfsetbuttcap%
\pgfsetroundjoin%
\definecolor{currentfill}{rgb}{0.000000,0.000000,0.000000}%
\pgfsetfillcolor{currentfill}%
\pgfsetlinewidth{0.803000pt}%
\definecolor{currentstroke}{rgb}{0.000000,0.000000,0.000000}%
\pgfsetstrokecolor{currentstroke}%
\pgfsetdash{}{0pt}%
\pgfsys@defobject{currentmarker}{\pgfqpoint{0.000000in}{-0.048611in}}{\pgfqpoint{0.000000in}{0.000000in}}{%
\pgfpathmoveto{\pgfqpoint{0.000000in}{0.000000in}}%
\pgfpathlineto{\pgfqpoint{0.000000in}{-0.048611in}}%
\pgfusepath{stroke,fill}%
}%
\begin{pgfscope}%
\pgfsys@transformshift{5.238651in}{0.565123in}%
\pgfsys@useobject{currentmarker}{}%
\end{pgfscope}%
\end{pgfscope}%
\begin{pgfscope}%
\definecolor{textcolor}{rgb}{0.000000,0.000000,0.000000}%
\pgfsetstrokecolor{textcolor}%
\pgfsetfillcolor{textcolor}%
\pgftext[x=5.238651in,y=0.467901in,,top]{\color{textcolor}\rmfamily\fontsize{10.000000}{12.000000}\selectfont \(\displaystyle {3.5}\)}%
\end{pgfscope}%
\begin{pgfscope}%
\pgfpathrectangle{\pgfqpoint{0.494137in}{0.565123in}}{\pgfqpoint{5.655863in}{2.585803in}}%
\pgfusepath{clip}%
\pgfsetbuttcap%
\pgfsetroundjoin%
\pgfsetlinewidth{0.803000pt}%
\definecolor{currentstroke}{rgb}{0.690196,0.690196,0.690196}%
\pgfsetstrokecolor{currentstroke}%
\pgfsetdash{{0.800000pt}{1.320000pt}}{0.000000pt}%
\pgfpathmoveto{\pgfqpoint{5.898567in}{0.565123in}}%
\pgfpathlineto{\pgfqpoint{5.898567in}{3.150926in}}%
\pgfusepath{stroke}%
\end{pgfscope}%
\begin{pgfscope}%
\pgfsetbuttcap%
\pgfsetroundjoin%
\definecolor{currentfill}{rgb}{0.000000,0.000000,0.000000}%
\pgfsetfillcolor{currentfill}%
\pgfsetlinewidth{0.803000pt}%
\definecolor{currentstroke}{rgb}{0.000000,0.000000,0.000000}%
\pgfsetstrokecolor{currentstroke}%
\pgfsetdash{}{0pt}%
\pgfsys@defobject{currentmarker}{\pgfqpoint{0.000000in}{-0.048611in}}{\pgfqpoint{0.000000in}{0.000000in}}{%
\pgfpathmoveto{\pgfqpoint{0.000000in}{0.000000in}}%
\pgfpathlineto{\pgfqpoint{0.000000in}{-0.048611in}}%
\pgfusepath{stroke,fill}%
}%
\begin{pgfscope}%
\pgfsys@transformshift{5.898567in}{0.565123in}%
\pgfsys@useobject{currentmarker}{}%
\end{pgfscope}%
\end{pgfscope}%
\begin{pgfscope}%
\definecolor{textcolor}{rgb}{0.000000,0.000000,0.000000}%
\pgfsetstrokecolor{textcolor}%
\pgfsetfillcolor{textcolor}%
\pgftext[x=5.898567in,y=0.467901in,,top]{\color{textcolor}\rmfamily\fontsize{10.000000}{12.000000}\selectfont \(\displaystyle {4.0}\)}%
\end{pgfscope}%
\begin{pgfscope}%
\definecolor{textcolor}{rgb}{0.000000,0.000000,0.000000}%
\pgfsetstrokecolor{textcolor}%
\pgfsetfillcolor{textcolor}%
\pgftext[x=3.322068in,y=0.288889in,,top]{\color{textcolor}\rmfamily\fontsize{10.000000}{12.000000}\selectfont Return Loss (dB)}%
\end{pgfscope}%
\begin{pgfscope}%
\definecolor{textcolor}{rgb}{0.000000,0.000000,0.000000}%
\pgfsetstrokecolor{textcolor}%
\pgfsetfillcolor{textcolor}%
\pgftext[x=6.150000in,y=0.302778in,right,top]{\color{textcolor}\rmfamily\fontsize{10.000000}{12.000000}\selectfont \(\displaystyle \times{10^{7}}{}\)}%
\end{pgfscope}%
\begin{pgfscope}%
\pgfpathrectangle{\pgfqpoint{0.494137in}{0.565123in}}{\pgfqpoint{5.655863in}{2.585803in}}%
\pgfusepath{clip}%
\pgfsetbuttcap%
\pgfsetroundjoin%
\pgfsetlinewidth{0.803000pt}%
\definecolor{currentstroke}{rgb}{0.690196,0.690196,0.690196}%
\pgfsetstrokecolor{currentstroke}%
\pgfsetdash{{0.800000pt}{1.320000pt}}{0.000000pt}%
\pgfpathmoveto{\pgfqpoint{0.494137in}{0.727883in}}%
\pgfpathlineto{\pgfqpoint{6.150000in}{0.727883in}}%
\pgfusepath{stroke}%
\end{pgfscope}%
\begin{pgfscope}%
\pgfsetbuttcap%
\pgfsetroundjoin%
\definecolor{currentfill}{rgb}{0.000000,0.000000,0.000000}%
\pgfsetfillcolor{currentfill}%
\pgfsetlinewidth{0.803000pt}%
\definecolor{currentstroke}{rgb}{0.000000,0.000000,0.000000}%
\pgfsetstrokecolor{currentstroke}%
\pgfsetdash{}{0pt}%
\pgfsys@defobject{currentmarker}{\pgfqpoint{-0.048611in}{0.000000in}}{\pgfqpoint{-0.000000in}{0.000000in}}{%
\pgfpathmoveto{\pgfqpoint{-0.000000in}{0.000000in}}%
\pgfpathlineto{\pgfqpoint{-0.048611in}{0.000000in}}%
\pgfusepath{stroke,fill}%
}%
\begin{pgfscope}%
\pgfsys@transformshift{0.494137in}{0.727883in}%
\pgfsys@useobject{currentmarker}{}%
\end{pgfscope}%
\end{pgfscope}%
\begin{pgfscope}%
\definecolor{textcolor}{rgb}{0.000000,0.000000,0.000000}%
\pgfsetstrokecolor{textcolor}%
\pgfsetfillcolor{textcolor}%
\pgftext[x=0.150000in, y=0.679658in, left, base]{\color{textcolor}\rmfamily\fontsize{10.000000}{12.000000}\selectfont \(\displaystyle {\ensuremath{-}25}\)}%
\end{pgfscope}%
\begin{pgfscope}%
\pgfpathrectangle{\pgfqpoint{0.494137in}{0.565123in}}{\pgfqpoint{5.655863in}{2.585803in}}%
\pgfusepath{clip}%
\pgfsetbuttcap%
\pgfsetroundjoin%
\pgfsetlinewidth{0.803000pt}%
\definecolor{currentstroke}{rgb}{0.690196,0.690196,0.690196}%
\pgfsetstrokecolor{currentstroke}%
\pgfsetdash{{0.800000pt}{1.320000pt}}{0.000000pt}%
\pgfpathmoveto{\pgfqpoint{0.494137in}{1.171250in}}%
\pgfpathlineto{\pgfqpoint{6.150000in}{1.171250in}}%
\pgfusepath{stroke}%
\end{pgfscope}%
\begin{pgfscope}%
\pgfsetbuttcap%
\pgfsetroundjoin%
\definecolor{currentfill}{rgb}{0.000000,0.000000,0.000000}%
\pgfsetfillcolor{currentfill}%
\pgfsetlinewidth{0.803000pt}%
\definecolor{currentstroke}{rgb}{0.000000,0.000000,0.000000}%
\pgfsetstrokecolor{currentstroke}%
\pgfsetdash{}{0pt}%
\pgfsys@defobject{currentmarker}{\pgfqpoint{-0.048611in}{0.000000in}}{\pgfqpoint{-0.000000in}{0.000000in}}{%
\pgfpathmoveto{\pgfqpoint{-0.000000in}{0.000000in}}%
\pgfpathlineto{\pgfqpoint{-0.048611in}{0.000000in}}%
\pgfusepath{stroke,fill}%
}%
\begin{pgfscope}%
\pgfsys@transformshift{0.494137in}{1.171250in}%
\pgfsys@useobject{currentmarker}{}%
\end{pgfscope}%
\end{pgfscope}%
\begin{pgfscope}%
\definecolor{textcolor}{rgb}{0.000000,0.000000,0.000000}%
\pgfsetstrokecolor{textcolor}%
\pgfsetfillcolor{textcolor}%
\pgftext[x=0.150000in, y=1.123025in, left, base]{\color{textcolor}\rmfamily\fontsize{10.000000}{12.000000}\selectfont \(\displaystyle {\ensuremath{-}20}\)}%
\end{pgfscope}%
\begin{pgfscope}%
\pgfpathrectangle{\pgfqpoint{0.494137in}{0.565123in}}{\pgfqpoint{5.655863in}{2.585803in}}%
\pgfusepath{clip}%
\pgfsetbuttcap%
\pgfsetroundjoin%
\pgfsetlinewidth{0.803000pt}%
\definecolor{currentstroke}{rgb}{0.690196,0.690196,0.690196}%
\pgfsetstrokecolor{currentstroke}%
\pgfsetdash{{0.800000pt}{1.320000pt}}{0.000000pt}%
\pgfpathmoveto{\pgfqpoint{0.494137in}{1.614617in}}%
\pgfpathlineto{\pgfqpoint{6.150000in}{1.614617in}}%
\pgfusepath{stroke}%
\end{pgfscope}%
\begin{pgfscope}%
\pgfsetbuttcap%
\pgfsetroundjoin%
\definecolor{currentfill}{rgb}{0.000000,0.000000,0.000000}%
\pgfsetfillcolor{currentfill}%
\pgfsetlinewidth{0.803000pt}%
\definecolor{currentstroke}{rgb}{0.000000,0.000000,0.000000}%
\pgfsetstrokecolor{currentstroke}%
\pgfsetdash{}{0pt}%
\pgfsys@defobject{currentmarker}{\pgfqpoint{-0.048611in}{0.000000in}}{\pgfqpoint{-0.000000in}{0.000000in}}{%
\pgfpathmoveto{\pgfqpoint{-0.000000in}{0.000000in}}%
\pgfpathlineto{\pgfqpoint{-0.048611in}{0.000000in}}%
\pgfusepath{stroke,fill}%
}%
\begin{pgfscope}%
\pgfsys@transformshift{0.494137in}{1.614617in}%
\pgfsys@useobject{currentmarker}{}%
\end{pgfscope}%
\end{pgfscope}%
\begin{pgfscope}%
\definecolor{textcolor}{rgb}{0.000000,0.000000,0.000000}%
\pgfsetstrokecolor{textcolor}%
\pgfsetfillcolor{textcolor}%
\pgftext[x=0.150000in, y=1.566391in, left, base]{\color{textcolor}\rmfamily\fontsize{10.000000}{12.000000}\selectfont \(\displaystyle {\ensuremath{-}15}\)}%
\end{pgfscope}%
\begin{pgfscope}%
\pgfpathrectangle{\pgfqpoint{0.494137in}{0.565123in}}{\pgfqpoint{5.655863in}{2.585803in}}%
\pgfusepath{clip}%
\pgfsetbuttcap%
\pgfsetroundjoin%
\pgfsetlinewidth{0.803000pt}%
\definecolor{currentstroke}{rgb}{0.690196,0.690196,0.690196}%
\pgfsetstrokecolor{currentstroke}%
\pgfsetdash{{0.800000pt}{1.320000pt}}{0.000000pt}%
\pgfpathmoveto{\pgfqpoint{0.494137in}{2.057983in}}%
\pgfpathlineto{\pgfqpoint{6.150000in}{2.057983in}}%
\pgfusepath{stroke}%
\end{pgfscope}%
\begin{pgfscope}%
\pgfsetbuttcap%
\pgfsetroundjoin%
\definecolor{currentfill}{rgb}{0.000000,0.000000,0.000000}%
\pgfsetfillcolor{currentfill}%
\pgfsetlinewidth{0.803000pt}%
\definecolor{currentstroke}{rgb}{0.000000,0.000000,0.000000}%
\pgfsetstrokecolor{currentstroke}%
\pgfsetdash{}{0pt}%
\pgfsys@defobject{currentmarker}{\pgfqpoint{-0.048611in}{0.000000in}}{\pgfqpoint{-0.000000in}{0.000000in}}{%
\pgfpathmoveto{\pgfqpoint{-0.000000in}{0.000000in}}%
\pgfpathlineto{\pgfqpoint{-0.048611in}{0.000000in}}%
\pgfusepath{stroke,fill}%
}%
\begin{pgfscope}%
\pgfsys@transformshift{0.494137in}{2.057983in}%
\pgfsys@useobject{currentmarker}{}%
\end{pgfscope}%
\end{pgfscope}%
\begin{pgfscope}%
\definecolor{textcolor}{rgb}{0.000000,0.000000,0.000000}%
\pgfsetstrokecolor{textcolor}%
\pgfsetfillcolor{textcolor}%
\pgftext[x=0.150000in, y=2.009758in, left, base]{\color{textcolor}\rmfamily\fontsize{10.000000}{12.000000}\selectfont \(\displaystyle {\ensuremath{-}10}\)}%
\end{pgfscope}%
\begin{pgfscope}%
\pgfpathrectangle{\pgfqpoint{0.494137in}{0.565123in}}{\pgfqpoint{5.655863in}{2.585803in}}%
\pgfusepath{clip}%
\pgfsetbuttcap%
\pgfsetroundjoin%
\pgfsetlinewidth{0.803000pt}%
\definecolor{currentstroke}{rgb}{0.690196,0.690196,0.690196}%
\pgfsetstrokecolor{currentstroke}%
\pgfsetdash{{0.800000pt}{1.320000pt}}{0.000000pt}%
\pgfpathmoveto{\pgfqpoint{0.494137in}{2.501350in}}%
\pgfpathlineto{\pgfqpoint{6.150000in}{2.501350in}}%
\pgfusepath{stroke}%
\end{pgfscope}%
\begin{pgfscope}%
\pgfsetbuttcap%
\pgfsetroundjoin%
\definecolor{currentfill}{rgb}{0.000000,0.000000,0.000000}%
\pgfsetfillcolor{currentfill}%
\pgfsetlinewidth{0.803000pt}%
\definecolor{currentstroke}{rgb}{0.000000,0.000000,0.000000}%
\pgfsetstrokecolor{currentstroke}%
\pgfsetdash{}{0pt}%
\pgfsys@defobject{currentmarker}{\pgfqpoint{-0.048611in}{0.000000in}}{\pgfqpoint{-0.000000in}{0.000000in}}{%
\pgfpathmoveto{\pgfqpoint{-0.000000in}{0.000000in}}%
\pgfpathlineto{\pgfqpoint{-0.048611in}{0.000000in}}%
\pgfusepath{stroke,fill}%
}%
\begin{pgfscope}%
\pgfsys@transformshift{0.494137in}{2.501350in}%
\pgfsys@useobject{currentmarker}{}%
\end{pgfscope}%
\end{pgfscope}%
\begin{pgfscope}%
\definecolor{textcolor}{rgb}{0.000000,0.000000,0.000000}%
\pgfsetstrokecolor{textcolor}%
\pgfsetfillcolor{textcolor}%
\pgftext[x=0.219445in, y=2.453125in, left, base]{\color{textcolor}\rmfamily\fontsize{10.000000}{12.000000}\selectfont \(\displaystyle {\ensuremath{-}5}\)}%
\end{pgfscope}%
\begin{pgfscope}%
\pgfpathrectangle{\pgfqpoint{0.494137in}{0.565123in}}{\pgfqpoint{5.655863in}{2.585803in}}%
\pgfusepath{clip}%
\pgfsetbuttcap%
\pgfsetroundjoin%
\pgfsetlinewidth{0.803000pt}%
\definecolor{currentstroke}{rgb}{0.690196,0.690196,0.690196}%
\pgfsetstrokecolor{currentstroke}%
\pgfsetdash{{0.800000pt}{1.320000pt}}{0.000000pt}%
\pgfpathmoveto{\pgfqpoint{0.494137in}{2.944716in}}%
\pgfpathlineto{\pgfqpoint{6.150000in}{2.944716in}}%
\pgfusepath{stroke}%
\end{pgfscope}%
\begin{pgfscope}%
\pgfsetbuttcap%
\pgfsetroundjoin%
\definecolor{currentfill}{rgb}{0.000000,0.000000,0.000000}%
\pgfsetfillcolor{currentfill}%
\pgfsetlinewidth{0.803000pt}%
\definecolor{currentstroke}{rgb}{0.000000,0.000000,0.000000}%
\pgfsetstrokecolor{currentstroke}%
\pgfsetdash{}{0pt}%
\pgfsys@defobject{currentmarker}{\pgfqpoint{-0.048611in}{0.000000in}}{\pgfqpoint{-0.000000in}{0.000000in}}{%
\pgfpathmoveto{\pgfqpoint{-0.000000in}{0.000000in}}%
\pgfpathlineto{\pgfqpoint{-0.048611in}{0.000000in}}%
\pgfusepath{stroke,fill}%
}%
\begin{pgfscope}%
\pgfsys@transformshift{0.494137in}{2.944716in}%
\pgfsys@useobject{currentmarker}{}%
\end{pgfscope}%
\end{pgfscope}%
\begin{pgfscope}%
\definecolor{textcolor}{rgb}{0.000000,0.000000,0.000000}%
\pgfsetstrokecolor{textcolor}%
\pgfsetfillcolor{textcolor}%
\pgftext[x=0.327470in, y=2.896491in, left, base]{\color{textcolor}\rmfamily\fontsize{10.000000}{12.000000}\selectfont \(\displaystyle {0}\)}%
\end{pgfscope}%
\begin{pgfscope}%
\pgfpathrectangle{\pgfqpoint{0.494137in}{0.565123in}}{\pgfqpoint{5.655863in}{2.585803in}}%
\pgfusepath{clip}%
\pgfsetrectcap%
\pgfsetroundjoin%
\pgfsetlinewidth{2.007500pt}%
\definecolor{currentstroke}{rgb}{0.121569,0.466667,0.705882}%
\pgfsetstrokecolor{currentstroke}%
\pgfsetdash{}{0pt}%
\pgfpathmoveto{\pgfqpoint{0.751221in}{2.924322in}}%
\pgfpathlineto{\pgfqpoint{0.756768in}{3.033390in}}%
\pgfpathlineto{\pgfqpoint{0.762314in}{3.033390in}}%
\pgfpathlineto{\pgfqpoint{0.773408in}{3.022749in}}%
\pgfpathlineto{\pgfqpoint{0.784501in}{3.022749in}}%
\pgfpathlineto{\pgfqpoint{0.790047in}{3.017429in}}%
\pgfpathlineto{\pgfqpoint{0.806687in}{3.017429in}}%
\pgfpathlineto{\pgfqpoint{0.812234in}{3.012108in}}%
\pgfpathlineto{\pgfqpoint{0.828874in}{3.012108in}}%
\pgfpathlineto{\pgfqpoint{0.839967in}{3.002354in}}%
\pgfpathlineto{\pgfqpoint{0.851060in}{3.002354in}}%
\pgfpathlineto{\pgfqpoint{0.856607in}{2.997034in}}%
\pgfpathlineto{\pgfqpoint{0.867700in}{2.997034in}}%
\pgfpathlineto{\pgfqpoint{0.873246in}{2.991713in}}%
\pgfpathlineto{\pgfqpoint{0.884340in}{2.991713in}}%
\pgfpathlineto{\pgfqpoint{0.889886in}{2.986393in}}%
\pgfpathlineto{\pgfqpoint{0.900979in}{2.986393in}}%
\pgfpathlineto{\pgfqpoint{0.906526in}{2.981072in}}%
\pgfpathlineto{\pgfqpoint{0.917619in}{2.981072in}}%
\pgfpathlineto{\pgfqpoint{0.923166in}{2.975752in}}%
\pgfpathlineto{\pgfqpoint{0.934259in}{2.975752in}}%
\pgfpathlineto{\pgfqpoint{0.939805in}{2.970432in}}%
\pgfpathlineto{\pgfqpoint{0.950899in}{2.970432in}}%
\pgfpathlineto{\pgfqpoint{0.956445in}{2.965111in}}%
\pgfpathlineto{\pgfqpoint{0.961992in}{2.965111in}}%
\pgfpathlineto{\pgfqpoint{0.973085in}{2.955357in}}%
\pgfpathlineto{\pgfqpoint{0.978632in}{2.955357in}}%
\pgfpathlineto{\pgfqpoint{0.984178in}{2.950037in}}%
\pgfpathlineto{\pgfqpoint{0.989725in}{2.950037in}}%
\pgfpathlineto{\pgfqpoint{0.995271in}{2.944716in}}%
\pgfpathlineto{\pgfqpoint{1.006365in}{2.944716in}}%
\pgfpathlineto{\pgfqpoint{1.011911in}{2.939396in}}%
\pgfpathlineto{\pgfqpoint{1.017458in}{2.939396in}}%
\pgfpathlineto{\pgfqpoint{1.023004in}{2.934076in}}%
\pgfpathlineto{\pgfqpoint{1.028551in}{2.934076in}}%
\pgfpathlineto{\pgfqpoint{1.034098in}{2.928755in}}%
\pgfpathlineto{\pgfqpoint{1.039644in}{2.928755in}}%
\pgfpathlineto{\pgfqpoint{1.045191in}{2.924322in}}%
\pgfpathlineto{\pgfqpoint{1.050737in}{2.924322in}}%
\pgfpathlineto{\pgfqpoint{1.056284in}{2.919001in}}%
\pgfpathlineto{\pgfqpoint{1.061831in}{2.919001in}}%
\pgfpathlineto{\pgfqpoint{1.067377in}{2.913681in}}%
\pgfpathlineto{\pgfqpoint{1.072924in}{2.913681in}}%
\pgfpathlineto{\pgfqpoint{1.078470in}{2.908360in}}%
\pgfpathlineto{\pgfqpoint{1.084017in}{2.908360in}}%
\pgfpathlineto{\pgfqpoint{1.089564in}{2.903040in}}%
\pgfpathlineto{\pgfqpoint{1.100657in}{2.903040in}}%
\pgfpathlineto{\pgfqpoint{1.111750in}{2.892399in}}%
\pgfpathlineto{\pgfqpoint{1.117297in}{2.892399in}}%
\pgfpathlineto{\pgfqpoint{1.122843in}{2.887079in}}%
\pgfpathlineto{\pgfqpoint{1.128390in}{2.887079in}}%
\pgfpathlineto{\pgfqpoint{1.133936in}{2.882645in}}%
\pgfpathlineto{\pgfqpoint{1.139483in}{2.882645in}}%
\pgfpathlineto{\pgfqpoint{1.145030in}{2.877325in}}%
\pgfpathlineto{\pgfqpoint{1.150576in}{2.877325in}}%
\pgfpathlineto{\pgfqpoint{1.161669in}{2.866684in}}%
\pgfpathlineto{\pgfqpoint{1.167216in}{2.866684in}}%
\pgfpathlineto{\pgfqpoint{1.172762in}{2.861363in}}%
\pgfpathlineto{\pgfqpoint{1.178309in}{2.861363in}}%
\pgfpathlineto{\pgfqpoint{1.183856in}{2.856043in}}%
\pgfpathlineto{\pgfqpoint{1.189402in}{2.856043in}}%
\pgfpathlineto{\pgfqpoint{1.200495in}{2.846289in}}%
\pgfpathlineto{\pgfqpoint{1.206042in}{2.846289in}}%
\pgfpathlineto{\pgfqpoint{1.217135in}{2.835648in}}%
\pgfpathlineto{\pgfqpoint{1.222682in}{2.835648in}}%
\pgfpathlineto{\pgfqpoint{1.228228in}{2.830328in}}%
\pgfpathlineto{\pgfqpoint{1.233775in}{2.830328in}}%
\pgfpathlineto{\pgfqpoint{1.244868in}{2.819687in}}%
\pgfpathlineto{\pgfqpoint{1.250415in}{2.819687in}}%
\pgfpathlineto{\pgfqpoint{1.261508in}{2.809933in}}%
\pgfpathlineto{\pgfqpoint{1.267055in}{2.809933in}}%
\pgfpathlineto{\pgfqpoint{1.278148in}{2.799292in}}%
\pgfpathlineto{\pgfqpoint{1.283694in}{2.799292in}}%
\pgfpathlineto{\pgfqpoint{1.300334in}{2.783331in}}%
\pgfpathlineto{\pgfqpoint{1.305881in}{2.783331in}}%
\pgfpathlineto{\pgfqpoint{1.316974in}{2.772690in}}%
\pgfpathlineto{\pgfqpoint{1.322521in}{2.772690in}}%
\pgfpathlineto{\pgfqpoint{1.333614in}{2.762936in}}%
\pgfpathlineto{\pgfqpoint{1.339160in}{2.757616in}}%
\pgfpathlineto{\pgfqpoint{1.344707in}{2.757616in}}%
\pgfpathlineto{\pgfqpoint{1.366893in}{2.736334in}}%
\pgfpathlineto{\pgfqpoint{1.372440in}{2.736334in}}%
\pgfpathlineto{\pgfqpoint{1.383533in}{2.726580in}}%
\pgfpathlineto{\pgfqpoint{1.394626in}{2.715939in}}%
\pgfpathlineto{\pgfqpoint{1.400173in}{2.715939in}}%
\pgfpathlineto{\pgfqpoint{1.427906in}{2.690224in}}%
\pgfpathlineto{\pgfqpoint{1.433452in}{2.690224in}}%
\pgfpathlineto{\pgfqpoint{1.483372in}{2.643227in}}%
\pgfpathlineto{\pgfqpoint{1.488918in}{2.643227in}}%
\pgfpathlineto{\pgfqpoint{1.555478in}{2.580269in}}%
\pgfpathlineto{\pgfqpoint{1.566571in}{2.570515in}}%
\pgfpathlineto{\pgfqpoint{1.599850in}{2.539479in}}%
\pgfpathlineto{\pgfqpoint{1.605397in}{2.528839in}}%
\pgfpathlineto{\pgfqpoint{1.644223in}{2.492482in}}%
\pgfpathlineto{\pgfqpoint{1.649770in}{2.481842in}}%
\pgfpathlineto{\pgfqpoint{1.677503in}{2.456126in}}%
\pgfpathlineto{\pgfqpoint{1.683049in}{2.445486in}}%
\pgfpathlineto{\pgfqpoint{1.705236in}{2.424204in}}%
\pgfpathlineto{\pgfqpoint{1.710782in}{2.414450in}}%
\pgfpathlineto{\pgfqpoint{1.727422in}{2.398489in}}%
\pgfpathlineto{\pgfqpoint{1.732969in}{2.387848in}}%
\pgfpathlineto{\pgfqpoint{1.744062in}{2.378094in}}%
\pgfpathlineto{\pgfqpoint{1.749608in}{2.367453in}}%
\pgfpathlineto{\pgfqpoint{1.766248in}{2.351492in}}%
\pgfpathlineto{\pgfqpoint{1.771795in}{2.341738in}}%
\pgfpathlineto{\pgfqpoint{1.777341in}{2.336417in}}%
\pgfpathlineto{\pgfqpoint{1.782888in}{2.325777in}}%
\pgfpathlineto{\pgfqpoint{1.793981in}{2.315136in}}%
\pgfpathlineto{\pgfqpoint{1.799528in}{2.305382in}}%
\pgfpathlineto{\pgfqpoint{1.805074in}{2.300061in}}%
\pgfpathlineto{\pgfqpoint{1.810621in}{2.289421in}}%
\pgfpathlineto{\pgfqpoint{1.821714in}{2.278780in}}%
\pgfpathlineto{\pgfqpoint{1.827261in}{2.269026in}}%
\pgfpathlineto{\pgfqpoint{1.832807in}{2.263705in}}%
\pgfpathlineto{\pgfqpoint{1.838354in}{2.253064in}}%
\pgfpathlineto{\pgfqpoint{1.843901in}{2.247744in}}%
\pgfpathlineto{\pgfqpoint{1.849447in}{2.237103in}}%
\pgfpathlineto{\pgfqpoint{1.854994in}{2.231783in}}%
\pgfpathlineto{\pgfqpoint{1.860540in}{2.222029in}}%
\pgfpathlineto{\pgfqpoint{1.866087in}{2.216708in}}%
\pgfpathlineto{\pgfqpoint{1.871634in}{2.206068in}}%
\pgfpathlineto{\pgfqpoint{1.877180in}{2.200747in}}%
\pgfpathlineto{\pgfqpoint{1.888273in}{2.180352in}}%
\pgfpathlineto{\pgfqpoint{1.893820in}{2.175032in}}%
\pgfpathlineto{\pgfqpoint{1.899367in}{2.164391in}}%
\pgfpathlineto{\pgfqpoint{1.904913in}{2.164391in}}%
\pgfpathlineto{\pgfqpoint{1.916006in}{2.143996in}}%
\pgfpathlineto{\pgfqpoint{1.921553in}{2.138676in}}%
\pgfpathlineto{\pgfqpoint{1.938193in}{2.107640in}}%
\pgfpathlineto{\pgfqpoint{1.943739in}{2.102320in}}%
\pgfpathlineto{\pgfqpoint{1.954832in}{2.081038in}}%
\pgfpathlineto{\pgfqpoint{1.960379in}{2.076605in}}%
\pgfpathlineto{\pgfqpoint{1.988112in}{2.024287in}}%
\pgfpathlineto{\pgfqpoint{1.993659in}{2.018967in}}%
\pgfpathlineto{\pgfqpoint{2.043578in}{1.924973in}}%
\pgfpathlineto{\pgfqpoint{2.054671in}{1.904578in}}%
\pgfpathlineto{\pgfqpoint{2.060218in}{1.888617in}}%
\pgfpathlineto{\pgfqpoint{2.076858in}{1.857581in}}%
\pgfpathlineto{\pgfqpoint{2.093497in}{1.826546in}}%
\pgfpathlineto{\pgfqpoint{2.099044in}{1.810585in}}%
\pgfpathlineto{\pgfqpoint{2.115684in}{1.779549in}}%
\pgfpathlineto{\pgfqpoint{2.121230in}{1.764474in}}%
\pgfpathlineto{\pgfqpoint{2.132324in}{1.743193in}}%
\pgfpathlineto{\pgfqpoint{2.137870in}{1.728118in}}%
\pgfpathlineto{\pgfqpoint{2.143417in}{1.717478in}}%
\pgfpathlineto{\pgfqpoint{2.148963in}{1.701516in}}%
\pgfpathlineto{\pgfqpoint{2.154510in}{1.691762in}}%
\pgfpathlineto{\pgfqpoint{2.160056in}{1.675801in}}%
\pgfpathlineto{\pgfqpoint{2.165603in}{1.665160in}}%
\pgfpathlineto{\pgfqpoint{2.171150in}{1.650086in}}%
\pgfpathlineto{\pgfqpoint{2.176696in}{1.639445in}}%
\pgfpathlineto{\pgfqpoint{2.182243in}{1.623484in}}%
\pgfpathlineto{\pgfqpoint{2.187789in}{1.613730in}}%
\pgfpathlineto{\pgfqpoint{2.204429in}{1.566733in}}%
\pgfpathlineto{\pgfqpoint{2.209976in}{1.556092in}}%
\pgfpathlineto{\pgfqpoint{2.254349in}{1.431063in}}%
\pgfpathlineto{\pgfqpoint{2.265442in}{1.400027in}}%
\pgfpathlineto{\pgfqpoint{2.270988in}{1.379632in}}%
\pgfpathlineto{\pgfqpoint{2.287628in}{1.332635in}}%
\pgfpathlineto{\pgfqpoint{2.293175in}{1.311354in}}%
\pgfpathlineto{\pgfqpoint{2.298721in}{1.296279in}}%
\pgfpathlineto{\pgfqpoint{2.304268in}{1.274998in}}%
\pgfpathlineto{\pgfqpoint{2.309815in}{1.259923in}}%
\pgfpathlineto{\pgfqpoint{2.315361in}{1.238642in}}%
\pgfpathlineto{\pgfqpoint{2.320908in}{1.223567in}}%
\pgfpathlineto{\pgfqpoint{2.332001in}{1.181891in}}%
\pgfpathlineto{\pgfqpoint{2.337548in}{1.165929in}}%
\pgfpathlineto{\pgfqpoint{2.431840in}{0.822764in}}%
\pgfpathlineto{\pgfqpoint{2.465119in}{0.728770in}}%
\pgfpathlineto{\pgfqpoint{2.476212in}{0.708375in}}%
\pgfpathlineto{\pgfqpoint{2.481759in}{0.703055in}}%
\pgfpathlineto{\pgfqpoint{2.487306in}{0.692414in}}%
\pgfpathlineto{\pgfqpoint{2.498399in}{0.682660in}}%
\pgfpathlineto{\pgfqpoint{2.520585in}{0.682660in}}%
\pgfpathlineto{\pgfqpoint{2.537225in}{0.697734in}}%
\pgfpathlineto{\pgfqpoint{2.542772in}{0.708375in}}%
\pgfpathlineto{\pgfqpoint{2.548318in}{0.713695in}}%
\pgfpathlineto{\pgfqpoint{2.553865in}{0.728770in}}%
\pgfpathlineto{\pgfqpoint{2.570505in}{0.760692in}}%
\pgfpathlineto{\pgfqpoint{2.576051in}{0.775767in}}%
\pgfpathlineto{\pgfqpoint{2.581598in}{0.786408in}}%
\pgfpathlineto{\pgfqpoint{2.587144in}{0.802369in}}%
\pgfpathlineto{\pgfqpoint{2.592691in}{0.822764in}}%
\pgfpathlineto{\pgfqpoint{2.598238in}{0.833404in}}%
\pgfpathlineto{\pgfqpoint{2.603784in}{0.853799in}}%
\pgfpathlineto{\pgfqpoint{2.620424in}{0.900796in}}%
\pgfpathlineto{\pgfqpoint{2.625971in}{0.911437in}}%
\pgfpathlineto{\pgfqpoint{2.631517in}{0.931832in}}%
\pgfpathlineto{\pgfqpoint{2.648157in}{0.978829in}}%
\pgfpathlineto{\pgfqpoint{2.653704in}{0.999224in}}%
\pgfpathlineto{\pgfqpoint{2.670343in}{1.046220in}}%
\pgfpathlineto{\pgfqpoint{2.675890in}{1.056861in}}%
\pgfpathlineto{\pgfqpoint{2.698076in}{1.118933in}}%
\pgfpathlineto{\pgfqpoint{2.703623in}{1.129573in}}%
\pgfpathlineto{\pgfqpoint{2.714716in}{1.160609in}}%
\pgfpathlineto{\pgfqpoint{2.720263in}{1.171250in}}%
\pgfpathlineto{\pgfqpoint{2.725809in}{1.187211in}}%
\pgfpathlineto{\pgfqpoint{2.731356in}{1.196965in}}%
\pgfpathlineto{\pgfqpoint{2.736902in}{1.212926in}}%
\pgfpathlineto{\pgfqpoint{2.742449in}{1.223567in}}%
\pgfpathlineto{\pgfqpoint{2.747996in}{1.238642in}}%
\pgfpathlineto{\pgfqpoint{2.759089in}{1.259923in}}%
\pgfpathlineto{\pgfqpoint{2.764635in}{1.274998in}}%
\pgfpathlineto{\pgfqpoint{2.820101in}{1.379632in}}%
\pgfpathlineto{\pgfqpoint{2.831195in}{1.400027in}}%
\pgfpathlineto{\pgfqpoint{2.836741in}{1.405347in}}%
\pgfpathlineto{\pgfqpoint{2.853381in}{1.436383in}}%
\pgfpathlineto{\pgfqpoint{2.858928in}{1.441704in}}%
\pgfpathlineto{\pgfqpoint{2.875567in}{1.472739in}}%
\pgfpathlineto{\pgfqpoint{2.886661in}{1.483380in}}%
\pgfpathlineto{\pgfqpoint{2.897754in}{1.503775in}}%
\pgfpathlineto{\pgfqpoint{2.903300in}{1.509095in}}%
\pgfpathlineto{\pgfqpoint{2.908847in}{1.519736in}}%
\pgfpathlineto{\pgfqpoint{2.919940in}{1.530377in}}%
\pgfpathlineto{\pgfqpoint{2.925487in}{1.540131in}}%
\pgfpathlineto{\pgfqpoint{2.931033in}{1.545451in}}%
\pgfpathlineto{\pgfqpoint{2.936580in}{1.556092in}}%
\pgfpathlineto{\pgfqpoint{2.953220in}{1.572053in}}%
\pgfpathlineto{\pgfqpoint{2.958766in}{1.581807in}}%
\pgfpathlineto{\pgfqpoint{2.969859in}{1.592448in}}%
\pgfpathlineto{\pgfqpoint{2.975406in}{1.603089in}}%
\pgfpathlineto{\pgfqpoint{3.008686in}{1.634125in}}%
\pgfpathlineto{\pgfqpoint{3.014232in}{1.644765in}}%
\pgfpathlineto{\pgfqpoint{3.064152in}{1.691762in}}%
\pgfpathlineto{\pgfqpoint{3.075245in}{1.701516in}}%
\pgfpathlineto{\pgfqpoint{3.086338in}{1.712157in}}%
\pgfpathlineto{\pgfqpoint{3.091885in}{1.712157in}}%
\pgfpathlineto{\pgfqpoint{3.125164in}{1.743193in}}%
\pgfpathlineto{\pgfqpoint{3.130711in}{1.743193in}}%
\pgfpathlineto{\pgfqpoint{3.152897in}{1.764474in}}%
\pgfpathlineto{\pgfqpoint{3.158444in}{1.764474in}}%
\pgfpathlineto{\pgfqpoint{3.169537in}{1.774228in}}%
\pgfpathlineto{\pgfqpoint{3.175083in}{1.779549in}}%
\pgfpathlineto{\pgfqpoint{3.180630in}{1.779549in}}%
\pgfpathlineto{\pgfqpoint{3.197270in}{1.795510in}}%
\pgfpathlineto{\pgfqpoint{3.202816in}{1.795510in}}%
\pgfpathlineto{\pgfqpoint{3.213910in}{1.806151in}}%
\pgfpathlineto{\pgfqpoint{3.219456in}{1.806151in}}%
\pgfpathlineto{\pgfqpoint{3.230549in}{1.815905in}}%
\pgfpathlineto{\pgfqpoint{3.236096in}{1.815905in}}%
\pgfpathlineto{\pgfqpoint{3.241643in}{1.821225in}}%
\pgfpathlineto{\pgfqpoint{3.247189in}{1.821225in}}%
\pgfpathlineto{\pgfqpoint{3.258282in}{1.831866in}}%
\pgfpathlineto{\pgfqpoint{3.263829in}{1.831866in}}%
\pgfpathlineto{\pgfqpoint{3.269376in}{1.837187in}}%
\pgfpathlineto{\pgfqpoint{3.274922in}{1.837187in}}%
\pgfpathlineto{\pgfqpoint{3.286015in}{1.846941in}}%
\pgfpathlineto{\pgfqpoint{3.291562in}{1.846941in}}%
\pgfpathlineto{\pgfqpoint{3.297109in}{1.852261in}}%
\pgfpathlineto{\pgfqpoint{3.302655in}{1.852261in}}%
\pgfpathlineto{\pgfqpoint{3.308202in}{1.857581in}}%
\pgfpathlineto{\pgfqpoint{3.313748in}{1.857581in}}%
\pgfpathlineto{\pgfqpoint{3.319295in}{1.862902in}}%
\pgfpathlineto{\pgfqpoint{3.324842in}{1.862902in}}%
\pgfpathlineto{\pgfqpoint{3.330388in}{1.868222in}}%
\pgfpathlineto{\pgfqpoint{3.335935in}{1.878863in}}%
\pgfpathlineto{\pgfqpoint{3.341481in}{1.873543in}}%
\pgfpathlineto{\pgfqpoint{3.347028in}{1.878863in}}%
\pgfpathlineto{\pgfqpoint{3.352575in}{1.878863in}}%
\pgfpathlineto{\pgfqpoint{3.363668in}{1.888617in}}%
\pgfpathlineto{\pgfqpoint{3.369214in}{1.888617in}}%
\pgfpathlineto{\pgfqpoint{3.374761in}{1.893937in}}%
\pgfpathlineto{\pgfqpoint{3.385854in}{1.893937in}}%
\pgfpathlineto{\pgfqpoint{3.396947in}{1.904578in}}%
\pgfpathlineto{\pgfqpoint{3.402494in}{1.904578in}}%
\pgfpathlineto{\pgfqpoint{3.408040in}{1.909899in}}%
\pgfpathlineto{\pgfqpoint{3.413587in}{1.909899in}}%
\pgfpathlineto{\pgfqpoint{3.419134in}{1.915219in}}%
\pgfpathlineto{\pgfqpoint{3.430227in}{1.915219in}}%
\pgfpathlineto{\pgfqpoint{3.435773in}{1.920539in}}%
\pgfpathlineto{\pgfqpoint{3.441320in}{1.920539in}}%
\pgfpathlineto{\pgfqpoint{3.446867in}{1.924973in}}%
\pgfpathlineto{\pgfqpoint{3.457960in}{1.924973in}}%
\pgfpathlineto{\pgfqpoint{3.463506in}{1.930294in}}%
\pgfpathlineto{\pgfqpoint{3.469053in}{1.930294in}}%
\pgfpathlineto{\pgfqpoint{3.474600in}{1.935614in}}%
\pgfpathlineto{\pgfqpoint{3.485693in}{1.935614in}}%
\pgfpathlineto{\pgfqpoint{3.491239in}{1.940934in}}%
\pgfpathlineto{\pgfqpoint{3.496786in}{1.940934in}}%
\pgfpathlineto{\pgfqpoint{3.502333in}{1.946255in}}%
\pgfpathlineto{\pgfqpoint{3.513426in}{1.946255in}}%
\pgfpathlineto{\pgfqpoint{3.518972in}{1.951575in}}%
\pgfpathlineto{\pgfqpoint{3.530066in}{1.951575in}}%
\pgfpathlineto{\pgfqpoint{3.535612in}{1.956896in}}%
\pgfpathlineto{\pgfqpoint{3.546705in}{1.956896in}}%
\pgfpathlineto{\pgfqpoint{3.552252in}{1.961329in}}%
\pgfpathlineto{\pgfqpoint{3.563345in}{1.961329in}}%
\pgfpathlineto{\pgfqpoint{3.568892in}{1.966650in}}%
\pgfpathlineto{\pgfqpoint{3.579985in}{1.966650in}}%
\pgfpathlineto{\pgfqpoint{3.585532in}{1.971970in}}%
\pgfpathlineto{\pgfqpoint{3.596625in}{1.971970in}}%
\pgfpathlineto{\pgfqpoint{3.602171in}{1.977290in}}%
\pgfpathlineto{\pgfqpoint{3.618811in}{1.977290in}}%
\pgfpathlineto{\pgfqpoint{3.624358in}{1.982611in}}%
\pgfpathlineto{\pgfqpoint{3.635451in}{1.982611in}}%
\pgfpathlineto{\pgfqpoint{3.640998in}{1.987931in}}%
\pgfpathlineto{\pgfqpoint{3.652091in}{1.987931in}}%
\pgfpathlineto{\pgfqpoint{3.657637in}{1.993252in}}%
\pgfpathlineto{\pgfqpoint{3.674277in}{1.993252in}}%
\pgfpathlineto{\pgfqpoint{3.679824in}{1.998572in}}%
\pgfpathlineto{\pgfqpoint{3.690917in}{1.998572in}}%
\pgfpathlineto{\pgfqpoint{3.696463in}{2.003006in}}%
\pgfpathlineto{\pgfqpoint{3.713103in}{2.003006in}}%
\pgfpathlineto{\pgfqpoint{3.718650in}{2.008326in}}%
\pgfpathlineto{\pgfqpoint{3.735290in}{2.008326in}}%
\pgfpathlineto{\pgfqpoint{3.740836in}{2.013646in}}%
\pgfpathlineto{\pgfqpoint{3.757476in}{2.013646in}}%
\pgfpathlineto{\pgfqpoint{3.763023in}{2.018967in}}%
\pgfpathlineto{\pgfqpoint{3.779662in}{2.018967in}}%
\pgfpathlineto{\pgfqpoint{3.785209in}{2.024287in}}%
\pgfpathlineto{\pgfqpoint{3.801849in}{2.024287in}}%
\pgfpathlineto{\pgfqpoint{3.807395in}{2.029608in}}%
\pgfpathlineto{\pgfqpoint{3.824035in}{2.029608in}}%
\pgfpathlineto{\pgfqpoint{3.829582in}{2.034928in}}%
\pgfpathlineto{\pgfqpoint{3.846222in}{2.034928in}}%
\pgfpathlineto{\pgfqpoint{3.851768in}{2.039362in}}%
\pgfpathlineto{\pgfqpoint{3.868408in}{2.039362in}}%
\pgfpathlineto{\pgfqpoint{3.873955in}{2.044682in}}%
\pgfpathlineto{\pgfqpoint{3.890594in}{2.044682in}}%
\pgfpathlineto{\pgfqpoint{3.896141in}{2.050003in}}%
\pgfpathlineto{\pgfqpoint{3.907234in}{2.050003in}}%
\pgfpathlineto{\pgfqpoint{3.912781in}{2.055323in}}%
\pgfpathlineto{\pgfqpoint{3.929420in}{2.055323in}}%
\pgfpathlineto{\pgfqpoint{3.934967in}{2.060643in}}%
\pgfpathlineto{\pgfqpoint{3.946060in}{2.060643in}}%
\pgfpathlineto{\pgfqpoint{3.951607in}{2.065964in}}%
\pgfpathlineto{\pgfqpoint{3.968247in}{2.065964in}}%
\pgfpathlineto{\pgfqpoint{3.973793in}{2.071284in}}%
\pgfpathlineto{\pgfqpoint{3.984886in}{2.071284in}}%
\pgfpathlineto{\pgfqpoint{3.990433in}{2.076605in}}%
\pgfpathlineto{\pgfqpoint{4.001526in}{2.076605in}}%
\pgfpathlineto{\pgfqpoint{4.007073in}{2.081038in}}%
\pgfpathlineto{\pgfqpoint{4.018166in}{2.081038in}}%
\pgfpathlineto{\pgfqpoint{4.023713in}{2.086359in}}%
\pgfpathlineto{\pgfqpoint{4.040352in}{2.086359in}}%
\pgfpathlineto{\pgfqpoint{4.045899in}{2.091679in}}%
\pgfpathlineto{\pgfqpoint{4.051446in}{2.091679in}}%
\pgfpathlineto{\pgfqpoint{4.056992in}{2.096999in}}%
\pgfpathlineto{\pgfqpoint{4.068085in}{2.096999in}}%
\pgfpathlineto{\pgfqpoint{4.073632in}{2.102320in}}%
\pgfpathlineto{\pgfqpoint{4.084725in}{2.102320in}}%
\pgfpathlineto{\pgfqpoint{4.090272in}{2.107640in}}%
\pgfpathlineto{\pgfqpoint{4.095818in}{2.107640in}}%
\pgfpathlineto{\pgfqpoint{4.101365in}{2.112961in}}%
\pgfpathlineto{\pgfqpoint{4.112458in}{2.112961in}}%
\pgfpathlineto{\pgfqpoint{4.118005in}{2.117394in}}%
\pgfpathlineto{\pgfqpoint{4.123551in}{2.117394in}}%
\pgfpathlineto{\pgfqpoint{4.129098in}{2.122715in}}%
\pgfpathlineto{\pgfqpoint{4.134645in}{2.122715in}}%
\pgfpathlineto{\pgfqpoint{4.140191in}{2.128035in}}%
\pgfpathlineto{\pgfqpoint{4.151284in}{2.128035in}}%
\pgfpathlineto{\pgfqpoint{4.156831in}{2.133355in}}%
\pgfpathlineto{\pgfqpoint{4.162377in}{2.133355in}}%
\pgfpathlineto{\pgfqpoint{4.167924in}{2.138676in}}%
\pgfpathlineto{\pgfqpoint{4.173471in}{2.138676in}}%
\pgfpathlineto{\pgfqpoint{4.179017in}{2.143996in}}%
\pgfpathlineto{\pgfqpoint{4.190110in}{2.143996in}}%
\pgfpathlineto{\pgfqpoint{4.201204in}{2.154637in}}%
\pgfpathlineto{\pgfqpoint{4.206750in}{2.154637in}}%
\pgfpathlineto{\pgfqpoint{4.212297in}{2.159071in}}%
\pgfpathlineto{\pgfqpoint{4.217843in}{2.159071in}}%
\pgfpathlineto{\pgfqpoint{4.223390in}{2.164391in}}%
\pgfpathlineto{\pgfqpoint{4.228937in}{2.164391in}}%
\pgfpathlineto{\pgfqpoint{4.234483in}{2.169712in}}%
\pgfpathlineto{\pgfqpoint{4.240030in}{2.169712in}}%
\pgfpathlineto{\pgfqpoint{4.251123in}{2.180352in}}%
\pgfpathlineto{\pgfqpoint{4.256670in}{2.180352in}}%
\pgfpathlineto{\pgfqpoint{4.262216in}{2.185673in}}%
\pgfpathlineto{\pgfqpoint{4.267763in}{2.185673in}}%
\pgfpathlineto{\pgfqpoint{4.278856in}{2.195427in}}%
\pgfpathlineto{\pgfqpoint{4.284403in}{2.195427in}}%
\pgfpathlineto{\pgfqpoint{4.289949in}{2.200747in}}%
\pgfpathlineto{\pgfqpoint{4.295496in}{2.200747in}}%
\pgfpathlineto{\pgfqpoint{4.306589in}{2.211388in}}%
\pgfpathlineto{\pgfqpoint{4.312136in}{2.211388in}}%
\pgfpathlineto{\pgfqpoint{4.323229in}{2.222029in}}%
\pgfpathlineto{\pgfqpoint{4.328775in}{2.222029in}}%
\pgfpathlineto{\pgfqpoint{4.339869in}{2.231783in}}%
\pgfpathlineto{\pgfqpoint{4.345415in}{2.242424in}}%
\pgfpathlineto{\pgfqpoint{4.350962in}{2.237103in}}%
\pgfpathlineto{\pgfqpoint{4.362055in}{2.247744in}}%
\pgfpathlineto{\pgfqpoint{4.367602in}{2.247744in}}%
\pgfpathlineto{\pgfqpoint{4.389788in}{2.269026in}}%
\pgfpathlineto{\pgfqpoint{4.395335in}{2.269026in}}%
\pgfpathlineto{\pgfqpoint{4.406428in}{2.278780in}}%
\pgfpathlineto{\pgfqpoint{4.411974in}{2.289421in}}%
\pgfpathlineto{\pgfqpoint{4.417521in}{2.289421in}}%
\pgfpathlineto{\pgfqpoint{4.428614in}{2.300061in}}%
\pgfpathlineto{\pgfqpoint{4.434161in}{2.300061in}}%
\pgfpathlineto{\pgfqpoint{4.439707in}{2.305382in}}%
\pgfpathlineto{\pgfqpoint{4.445254in}{2.315136in}}%
\pgfpathlineto{\pgfqpoint{4.450800in}{2.315136in}}%
\pgfpathlineto{\pgfqpoint{4.456347in}{2.325777in}}%
\pgfpathlineto{\pgfqpoint{4.461894in}{2.325777in}}%
\pgfpathlineto{\pgfqpoint{4.489627in}{2.351492in}}%
\pgfpathlineto{\pgfqpoint{4.495173in}{2.351492in}}%
\pgfpathlineto{\pgfqpoint{4.500720in}{2.362133in}}%
\pgfpathlineto{\pgfqpoint{4.528453in}{2.387848in}}%
\pgfpathlineto{\pgfqpoint{4.533999in}{2.387848in}}%
\pgfpathlineto{\pgfqpoint{4.611652in}{2.461447in}}%
\pgfpathlineto{\pgfqpoint{4.622745in}{2.471201in}}%
\pgfpathlineto{\pgfqpoint{4.633838in}{2.481842in}}%
\pgfpathlineto{\pgfqpoint{4.639385in}{2.492482in}}%
\pgfpathlineto{\pgfqpoint{4.689304in}{2.539479in}}%
\pgfpathlineto{\pgfqpoint{4.700397in}{2.549233in}}%
\pgfpathlineto{\pgfqpoint{4.733677in}{2.580269in}}%
\pgfpathlineto{\pgfqpoint{4.739223in}{2.590910in}}%
\pgfpathlineto{\pgfqpoint{4.805783in}{2.653868in}}%
\pgfpathlineto{\pgfqpoint{4.816876in}{2.663622in}}%
\pgfpathlineto{\pgfqpoint{4.888982in}{2.731900in}}%
\pgfpathlineto{\pgfqpoint{4.900075in}{2.741655in}}%
\pgfpathlineto{\pgfqpoint{4.938901in}{2.778011in}}%
\pgfpathlineto{\pgfqpoint{4.944447in}{2.778011in}}%
\pgfpathlineto{\pgfqpoint{4.977727in}{2.809933in}}%
\pgfpathlineto{\pgfqpoint{4.983274in}{2.809933in}}%
\pgfpathlineto{\pgfqpoint{4.994367in}{2.819687in}}%
\pgfpathlineto{\pgfqpoint{5.016553in}{2.840969in}}%
\pgfpathlineto{\pgfqpoint{5.022100in}{2.840969in}}%
\pgfpathlineto{\pgfqpoint{5.044286in}{2.861363in}}%
\pgfpathlineto{\pgfqpoint{5.049833in}{2.861363in}}%
\pgfpathlineto{\pgfqpoint{5.072019in}{2.882645in}}%
\pgfpathlineto{\pgfqpoint{5.077566in}{2.882645in}}%
\pgfpathlineto{\pgfqpoint{5.088659in}{2.892399in}}%
\pgfpathlineto{\pgfqpoint{5.094206in}{2.897720in}}%
\pgfpathlineto{\pgfqpoint{5.099752in}{2.897720in}}%
\pgfpathlineto{\pgfqpoint{5.110845in}{2.908360in}}%
\pgfpathlineto{\pgfqpoint{5.116392in}{2.908360in}}%
\pgfpathlineto{\pgfqpoint{5.127485in}{2.919001in}}%
\pgfpathlineto{\pgfqpoint{5.133032in}{2.919001in}}%
\pgfpathlineto{\pgfqpoint{5.144125in}{2.928755in}}%
\pgfpathlineto{\pgfqpoint{5.149671in}{2.928755in}}%
\pgfpathlineto{\pgfqpoint{5.155218in}{2.934076in}}%
\pgfpathlineto{\pgfqpoint{5.160765in}{2.934076in}}%
\pgfpathlineto{\pgfqpoint{5.166311in}{2.939396in}}%
\pgfpathlineto{\pgfqpoint{5.171858in}{2.939396in}}%
\pgfpathlineto{\pgfqpoint{5.177404in}{2.944716in}}%
\pgfpathlineto{\pgfqpoint{5.182951in}{2.944716in}}%
\pgfpathlineto{\pgfqpoint{5.188498in}{2.950037in}}%
\pgfpathlineto{\pgfqpoint{5.199591in}{2.950037in}}%
\pgfpathlineto{\pgfqpoint{5.205137in}{2.955357in}}%
\pgfpathlineto{\pgfqpoint{5.210684in}{2.955357in}}%
\pgfpathlineto{\pgfqpoint{5.216231in}{2.960678in}}%
\pgfpathlineto{\pgfqpoint{5.221777in}{2.960678in}}%
\pgfpathlineto{\pgfqpoint{5.227324in}{2.965111in}}%
\pgfpathlineto{\pgfqpoint{5.238417in}{2.965111in}}%
\pgfpathlineto{\pgfqpoint{5.243964in}{2.970432in}}%
\pgfpathlineto{\pgfqpoint{5.249510in}{2.970432in}}%
\pgfpathlineto{\pgfqpoint{5.255057in}{2.975752in}}%
\pgfpathlineto{\pgfqpoint{5.266150in}{2.975752in}}%
\pgfpathlineto{\pgfqpoint{5.271697in}{2.981072in}}%
\pgfpathlineto{\pgfqpoint{5.282790in}{2.981072in}}%
\pgfpathlineto{\pgfqpoint{5.288336in}{2.986393in}}%
\pgfpathlineto{\pgfqpoint{5.299430in}{2.986393in}}%
\pgfpathlineto{\pgfqpoint{5.304976in}{2.991713in}}%
\pgfpathlineto{\pgfqpoint{5.321616in}{2.991713in}}%
\pgfpathlineto{\pgfqpoint{5.327163in}{2.997034in}}%
\pgfpathlineto{\pgfqpoint{5.343802in}{2.997034in}}%
\pgfpathlineto{\pgfqpoint{5.349349in}{3.002354in}}%
\pgfpathlineto{\pgfqpoint{5.377082in}{3.002354in}}%
\pgfpathlineto{\pgfqpoint{5.382629in}{3.006788in}}%
\pgfpathlineto{\pgfqpoint{5.415908in}{3.006788in}}%
\pgfpathlineto{\pgfqpoint{5.421455in}{3.012108in}}%
\pgfpathlineto{\pgfqpoint{5.554573in}{3.012108in}}%
\pgfpathlineto{\pgfqpoint{5.560120in}{3.006788in}}%
\pgfpathlineto{\pgfqpoint{5.604492in}{3.006788in}}%
\pgfpathlineto{\pgfqpoint{5.610039in}{3.002354in}}%
\pgfpathlineto{\pgfqpoint{5.643318in}{3.002354in}}%
\pgfpathlineto{\pgfqpoint{5.648865in}{2.997034in}}%
\pgfpathlineto{\pgfqpoint{5.671051in}{2.997034in}}%
\pgfpathlineto{\pgfqpoint{5.676598in}{2.991713in}}%
\pgfpathlineto{\pgfqpoint{5.698784in}{2.991713in}}%
\pgfpathlineto{\pgfqpoint{5.704331in}{2.986393in}}%
\pgfpathlineto{\pgfqpoint{5.726517in}{2.986393in}}%
\pgfpathlineto{\pgfqpoint{5.732064in}{2.981072in}}%
\pgfpathlineto{\pgfqpoint{5.748704in}{2.981072in}}%
\pgfpathlineto{\pgfqpoint{5.754250in}{2.975752in}}%
\pgfpathlineto{\pgfqpoint{5.765344in}{2.975752in}}%
\pgfpathlineto{\pgfqpoint{5.770890in}{2.970432in}}%
\pgfpathlineto{\pgfqpoint{5.793077in}{2.970432in}}%
\pgfpathlineto{\pgfqpoint{5.798623in}{2.965111in}}%
\pgfpathlineto{\pgfqpoint{5.809716in}{2.965111in}}%
\pgfpathlineto{\pgfqpoint{5.815263in}{2.960678in}}%
\pgfpathlineto{\pgfqpoint{5.831903in}{2.960678in}}%
\pgfpathlineto{\pgfqpoint{5.837449in}{2.955357in}}%
\pgfpathlineto{\pgfqpoint{5.848543in}{2.955357in}}%
\pgfpathlineto{\pgfqpoint{5.854089in}{2.950037in}}%
\pgfpathlineto{\pgfqpoint{5.865182in}{2.950037in}}%
\pgfpathlineto{\pgfqpoint{5.870729in}{2.944716in}}%
\pgfpathlineto{\pgfqpoint{5.881822in}{2.944716in}}%
\pgfpathlineto{\pgfqpoint{5.887369in}{2.939396in}}%
\pgfpathlineto{\pgfqpoint{5.892915in}{2.944716in}}%
\pgfpathlineto{\pgfqpoint{5.892915in}{2.944716in}}%
\pgfusepath{stroke}%
\end{pgfscope}%
\begin{pgfscope}%
\pgfsetrectcap%
\pgfsetmiterjoin%
\pgfsetlinewidth{0.803000pt}%
\definecolor{currentstroke}{rgb}{0.000000,0.000000,0.000000}%
\pgfsetstrokecolor{currentstroke}%
\pgfsetdash{}{0pt}%
\pgfpathmoveto{\pgfqpoint{0.494137in}{0.565123in}}%
\pgfpathlineto{\pgfqpoint{0.494137in}{3.150926in}}%
\pgfusepath{stroke}%
\end{pgfscope}%
\begin{pgfscope}%
\pgfsetrectcap%
\pgfsetmiterjoin%
\pgfsetlinewidth{0.803000pt}%
\definecolor{currentstroke}{rgb}{0.000000,0.000000,0.000000}%
\pgfsetstrokecolor{currentstroke}%
\pgfsetdash{}{0pt}%
\pgfpathmoveto{\pgfqpoint{6.150000in}{0.565123in}}%
\pgfpathlineto{\pgfqpoint{6.150000in}{3.150926in}}%
\pgfusepath{stroke}%
\end{pgfscope}%
\begin{pgfscope}%
\pgfsetrectcap%
\pgfsetmiterjoin%
\pgfsetlinewidth{0.803000pt}%
\definecolor{currentstroke}{rgb}{0.000000,0.000000,0.000000}%
\pgfsetstrokecolor{currentstroke}%
\pgfsetdash{}{0pt}%
\pgfpathmoveto{\pgfqpoint{0.494137in}{0.565123in}}%
\pgfpathlineto{\pgfqpoint{6.150000in}{0.565123in}}%
\pgfusepath{stroke}%
\end{pgfscope}%
\begin{pgfscope}%
\pgfsetrectcap%
\pgfsetmiterjoin%
\pgfsetlinewidth{0.803000pt}%
\definecolor{currentstroke}{rgb}{0.000000,0.000000,0.000000}%
\pgfsetstrokecolor{currentstroke}%
\pgfsetdash{}{0pt}%
\pgfpathmoveto{\pgfqpoint{0.494137in}{3.150926in}}%
\pgfpathlineto{\pgfqpoint{6.150000in}{3.150926in}}%
\pgfusepath{stroke}%
\end{pgfscope}%
\begin{pgfscope}%
\definecolor{textcolor}{rgb}{0.000000,0.000000,0.000000}%
\pgfsetstrokecolor{textcolor}%
\pgfsetfillcolor{textcolor}%
\pgftext[x=3.322068in,y=3.234260in,,base]{\color{textcolor}\rmfamily\fontsize{12.000000}{14.400000}\selectfont Measured Return Loss of Matching Network}%
\end{pgfscope}%
\begin{pgfscope}%
\pgfpathrectangle{\pgfqpoint{0.494137in}{0.565123in}}{\pgfqpoint{5.655863in}{2.585803in}}%
\pgfusepath{clip}%
\pgfsetbuttcap%
\pgfsetroundjoin%
\definecolor{currentfill}{rgb}{0.121569,0.466667,0.705882}%
\pgfsetfillcolor{currentfill}%
\pgfsetlinewidth{1.003750pt}%
\definecolor{currentstroke}{rgb}{0.121569,0.466667,0.705882}%
\pgfsetstrokecolor{currentstroke}%
\pgfsetdash{}{0pt}%
\pgfsys@defobject{currentmarker}{\pgfqpoint{-0.041667in}{-0.041667in}}{\pgfqpoint{0.041667in}{0.041667in}}{%
\pgfpathmoveto{\pgfqpoint{0.000000in}{-0.041667in}}%
\pgfpathcurveto{\pgfqpoint{0.011050in}{-0.041667in}}{\pgfqpoint{0.021649in}{-0.037276in}}{\pgfqpoint{0.029463in}{-0.029463in}}%
\pgfpathcurveto{\pgfqpoint{0.037276in}{-0.021649in}}{\pgfqpoint{0.041667in}{-0.011050in}}{\pgfqpoint{0.041667in}{0.000000in}}%
\pgfpathcurveto{\pgfqpoint{0.041667in}{0.011050in}}{\pgfqpoint{0.037276in}{0.021649in}}{\pgfqpoint{0.029463in}{0.029463in}}%
\pgfpathcurveto{\pgfqpoint{0.021649in}{0.037276in}}{\pgfqpoint{0.011050in}{0.041667in}}{\pgfqpoint{0.000000in}{0.041667in}}%
\pgfpathcurveto{\pgfqpoint{-0.011050in}{0.041667in}}{\pgfqpoint{-0.021649in}{0.037276in}}{\pgfqpoint{-0.029463in}{0.029463in}}%
\pgfpathcurveto{\pgfqpoint{-0.037276in}{0.021649in}}{\pgfqpoint{-0.041667in}{0.011050in}}{\pgfqpoint{-0.041667in}{0.000000in}}%
\pgfpathcurveto{\pgfqpoint{-0.041667in}{-0.011050in}}{\pgfqpoint{-0.037276in}{-0.021649in}}{\pgfqpoint{-0.029463in}{-0.029463in}}%
\pgfpathcurveto{\pgfqpoint{-0.021649in}{-0.037276in}}{\pgfqpoint{-0.011050in}{-0.041667in}}{\pgfqpoint{0.000000in}{-0.041667in}}%
\pgfpathlineto{\pgfqpoint{0.000000in}{-0.041667in}}%
\pgfpathclose%
\pgfusepath{stroke,fill}%
}%
\begin{pgfscope}%
\pgfsys@transformshift{2.498399in}{0.682660in}%
\pgfsys@useobject{currentmarker}{}%
\end{pgfscope}%
\end{pgfscope}%
\begin{pgfscope}%
\pgfsetbuttcap%
\pgfsetmiterjoin%
\definecolor{currentfill}{rgb}{1.000000,1.000000,1.000000}%
\pgfsetfillcolor{currentfill}%
\pgfsetfillopacity{0.800000}%
\pgfsetlinewidth{1.003750pt}%
\definecolor{currentstroke}{rgb}{0.800000,0.800000,0.800000}%
\pgfsetstrokecolor{currentstroke}%
\pgfsetstrokeopacity{0.800000}%
\pgfsetdash{}{0pt}%
\pgfpathmoveto{\pgfqpoint{0.591359in}{0.634568in}}%
\pgfpathlineto{\pgfqpoint{2.104478in}{0.634568in}}%
\pgfpathquadraticcurveto{\pgfqpoint{2.132256in}{0.634568in}}{\pgfqpoint{2.132256in}{0.662346in}}%
\pgfpathlineto{\pgfqpoint{2.132256in}{1.050463in}}%
\pgfpathquadraticcurveto{\pgfqpoint{2.132256in}{1.078241in}}{\pgfqpoint{2.104478in}{1.078241in}}%
\pgfpathlineto{\pgfqpoint{0.591359in}{1.078241in}}%
\pgfpathquadraticcurveto{\pgfqpoint{0.563581in}{1.078241in}}{\pgfqpoint{0.563581in}{1.050463in}}%
\pgfpathlineto{\pgfqpoint{0.563581in}{0.662346in}}%
\pgfpathquadraticcurveto{\pgfqpoint{0.563581in}{0.634568in}}{\pgfqpoint{0.591359in}{0.634568in}}%
\pgfpathlineto{\pgfqpoint{0.591359in}{0.634568in}}%
\pgfpathclose%
\pgfusepath{stroke,fill}%
\end{pgfscope}%
\begin{pgfscope}%
\pgfsetrectcap%
\pgfsetroundjoin%
\pgfsetlinewidth{2.007500pt}%
\definecolor{currentstroke}{rgb}{0.121569,0.466667,0.705882}%
\pgfsetstrokecolor{currentstroke}%
\pgfsetdash{}{0pt}%
\pgfpathmoveto{\pgfqpoint{0.619137in}{0.967129in}}%
\pgfpathlineto{\pgfqpoint{0.758025in}{0.967129in}}%
\pgfpathlineto{\pgfqpoint{0.896914in}{0.967129in}}%
\pgfusepath{stroke}%
\end{pgfscope}%
\begin{pgfscope}%
\definecolor{textcolor}{rgb}{0.000000,0.000000,0.000000}%
\pgfsetstrokecolor{textcolor}%
\pgfsetfillcolor{textcolor}%
\pgftext[x=1.008025in,y=0.918518in,left,base]{\color{textcolor}\rmfamily\fontsize{10.000000}{12.000000}\selectfont Return Loss (dB)}%
\end{pgfscope}%
\begin{pgfscope}%
\pgfsetbuttcap%
\pgfsetroundjoin%
\definecolor{currentfill}{rgb}{0.121569,0.466667,0.705882}%
\pgfsetfillcolor{currentfill}%
\pgfsetlinewidth{1.003750pt}%
\definecolor{currentstroke}{rgb}{0.121569,0.466667,0.705882}%
\pgfsetstrokecolor{currentstroke}%
\pgfsetdash{}{0pt}%
\pgfsys@defobject{currentmarker}{\pgfqpoint{-0.041667in}{-0.041667in}}{\pgfqpoint{0.041667in}{0.041667in}}{%
\pgfpathmoveto{\pgfqpoint{0.000000in}{-0.041667in}}%
\pgfpathcurveto{\pgfqpoint{0.011050in}{-0.041667in}}{\pgfqpoint{0.021649in}{-0.037276in}}{\pgfqpoint{0.029463in}{-0.029463in}}%
\pgfpathcurveto{\pgfqpoint{0.037276in}{-0.021649in}}{\pgfqpoint{0.041667in}{-0.011050in}}{\pgfqpoint{0.041667in}{0.000000in}}%
\pgfpathcurveto{\pgfqpoint{0.041667in}{0.011050in}}{\pgfqpoint{0.037276in}{0.021649in}}{\pgfqpoint{0.029463in}{0.029463in}}%
\pgfpathcurveto{\pgfqpoint{0.021649in}{0.037276in}}{\pgfqpoint{0.011050in}{0.041667in}}{\pgfqpoint{0.000000in}{0.041667in}}%
\pgfpathcurveto{\pgfqpoint{-0.011050in}{0.041667in}}{\pgfqpoint{-0.021649in}{0.037276in}}{\pgfqpoint{-0.029463in}{0.029463in}}%
\pgfpathcurveto{\pgfqpoint{-0.037276in}{0.021649in}}{\pgfqpoint{-0.041667in}{0.011050in}}{\pgfqpoint{-0.041667in}{0.000000in}}%
\pgfpathcurveto{\pgfqpoint{-0.041667in}{-0.011050in}}{\pgfqpoint{-0.037276in}{-0.021649in}}{\pgfqpoint{-0.029463in}{-0.029463in}}%
\pgfpathcurveto{\pgfqpoint{-0.021649in}{-0.037276in}}{\pgfqpoint{-0.011050in}{-0.041667in}}{\pgfqpoint{0.000000in}{-0.041667in}}%
\pgfpathlineto{\pgfqpoint{0.000000in}{-0.041667in}}%
\pgfpathclose%
\pgfusepath{stroke,fill}%
}%
\begin{pgfscope}%
\pgfsys@transformshift{0.758025in}{0.753588in}%
\pgfsys@useobject{currentmarker}{}%
\end{pgfscope}%
\end{pgfscope}%
\begin{pgfscope}%
\definecolor{textcolor}{rgb}{0.000000,0.000000,0.000000}%
\pgfsetstrokecolor{textcolor}%
\pgfsetfillcolor{textcolor}%
\pgftext[x=1.008025in,y=0.717129in,left,base]{\color{textcolor}\rmfamily\fontsize{10.000000}{12.000000}\selectfont Resonance}%
\end{pgfscope}%
\end{pgfpicture}%
\makeatother%
\endgroup%

		\caption{Measured return loss and phase of a quartz resonatoras a function of frequency.}
		\label{fig:357}
	\end{figure}
	
	\subsubsection{}
	\begin{figure}[H]
		\centering
		%% Creator: Matplotlib, PGF backend
%%
%% To include the figure in your LaTeX document, write
%%   \input{<filename>.pgf}
%%
%% Make sure the required packages are loaded in your preamble
%%   \usepackage{pgf}
%%
%% Also ensure that all the required font packages are loaded; for instance,
%% the lmodern package is sometimes necessary when using math font.
%%   \usepackage{lmodern}
%%
%% Figures using additional raster images can only be included by \input if
%% they are in the same directory as the main LaTeX file. For loading figures
%% from other directories you can use the `import` package
%%   \usepackage{import}
%%
%% and then include the figures with
%%   \import{<path to file>}{<filename>.pgf}
%%
%% Matplotlib used the following preamble
%%
\begingroup%
\makeatletter%
\begin{pgfpicture}%
\pgfpathrectangle{\pgfpointorigin}{\pgfqpoint{6.300000in}{3.500000in}}%
\pgfusepath{use as bounding box, clip}%
\begin{pgfscope}%
\pgfsetbuttcap%
\pgfsetmiterjoin%
\definecolor{currentfill}{rgb}{1.000000,1.000000,1.000000}%
\pgfsetfillcolor{currentfill}%
\pgfsetlinewidth{0.000000pt}%
\definecolor{currentstroke}{rgb}{1.000000,1.000000,1.000000}%
\pgfsetstrokecolor{currentstroke}%
\pgfsetdash{}{0pt}%
\pgfpathmoveto{\pgfqpoint{0.000000in}{0.000000in}}%
\pgfpathlineto{\pgfqpoint{6.300000in}{0.000000in}}%
\pgfpathlineto{\pgfqpoint{6.300000in}{3.500000in}}%
\pgfpathlineto{\pgfqpoint{0.000000in}{3.500000in}}%
\pgfpathlineto{\pgfqpoint{0.000000in}{0.000000in}}%
\pgfpathclose%
\pgfusepath{fill}%
\end{pgfscope}%
\begin{pgfscope}%
\pgfsetbuttcap%
\pgfsetmiterjoin%
\definecolor{currentfill}{rgb}{1.000000,1.000000,1.000000}%
\pgfsetfillcolor{currentfill}%
\pgfsetlinewidth{0.000000pt}%
\definecolor{currentstroke}{rgb}{0.000000,0.000000,0.000000}%
\pgfsetstrokecolor{currentstroke}%
\pgfsetstrokeopacity{0.000000}%
\pgfsetdash{}{0pt}%
\pgfpathmoveto{\pgfqpoint{0.494137in}{0.565123in}}%
\pgfpathlineto{\pgfqpoint{6.150000in}{0.565123in}}%
\pgfpathlineto{\pgfqpoint{6.150000in}{3.336317in}}%
\pgfpathlineto{\pgfqpoint{0.494137in}{3.336317in}}%
\pgfpathlineto{\pgfqpoint{0.494137in}{0.565123in}}%
\pgfpathclose%
\pgfusepath{fill}%
\end{pgfscope}%
\begin{pgfscope}%
\pgfpathrectangle{\pgfqpoint{0.494137in}{0.565123in}}{\pgfqpoint{5.655863in}{2.771193in}}%
\pgfusepath{clip}%
\pgfsetbuttcap%
\pgfsetroundjoin%
\pgfsetlinewidth{0.803000pt}%
\definecolor{currentstroke}{rgb}{0.690196,0.690196,0.690196}%
\pgfsetstrokecolor{currentstroke}%
\pgfsetdash{{0.800000pt}{1.320000pt}}{0.000000pt}%
\pgfpathmoveto{\pgfqpoint{1.120449in}{0.565123in}}%
\pgfpathlineto{\pgfqpoint{1.120449in}{3.336317in}}%
\pgfusepath{stroke}%
\end{pgfscope}%
\begin{pgfscope}%
\pgfsetbuttcap%
\pgfsetroundjoin%
\definecolor{currentfill}{rgb}{0.000000,0.000000,0.000000}%
\pgfsetfillcolor{currentfill}%
\pgfsetlinewidth{0.803000pt}%
\definecolor{currentstroke}{rgb}{0.000000,0.000000,0.000000}%
\pgfsetstrokecolor{currentstroke}%
\pgfsetdash{}{0pt}%
\pgfsys@defobject{currentmarker}{\pgfqpoint{0.000000in}{-0.048611in}}{\pgfqpoint{0.000000in}{0.000000in}}{%
\pgfpathmoveto{\pgfqpoint{0.000000in}{0.000000in}}%
\pgfpathlineto{\pgfqpoint{0.000000in}{-0.048611in}}%
\pgfusepath{stroke,fill}%
}%
\begin{pgfscope}%
\pgfsys@transformshift{1.120449in}{0.565123in}%
\pgfsys@useobject{currentmarker}{}%
\end{pgfscope}%
\end{pgfscope}%
\begin{pgfscope}%
\definecolor{textcolor}{rgb}{0.000000,0.000000,0.000000}%
\pgfsetstrokecolor{textcolor}%
\pgfsetfillcolor{textcolor}%
\pgftext[x=1.120449in,y=0.467901in,,top]{\color{textcolor}\rmfamily\fontsize{10.000000}{12.000000}\selectfont \(\displaystyle {0.2}\)}%
\end{pgfscope}%
\begin{pgfscope}%
\pgfpathrectangle{\pgfqpoint{0.494137in}{0.565123in}}{\pgfqpoint{5.655863in}{2.771193in}}%
\pgfusepath{clip}%
\pgfsetbuttcap%
\pgfsetroundjoin%
\pgfsetlinewidth{0.803000pt}%
\definecolor{currentstroke}{rgb}{0.690196,0.690196,0.690196}%
\pgfsetstrokecolor{currentstroke}%
\pgfsetdash{{0.800000pt}{1.320000pt}}{0.000000pt}%
\pgfpathmoveto{\pgfqpoint{1.749908in}{0.565123in}}%
\pgfpathlineto{\pgfqpoint{1.749908in}{3.336317in}}%
\pgfusepath{stroke}%
\end{pgfscope}%
\begin{pgfscope}%
\pgfsetbuttcap%
\pgfsetroundjoin%
\definecolor{currentfill}{rgb}{0.000000,0.000000,0.000000}%
\pgfsetfillcolor{currentfill}%
\pgfsetlinewidth{0.803000pt}%
\definecolor{currentstroke}{rgb}{0.000000,0.000000,0.000000}%
\pgfsetstrokecolor{currentstroke}%
\pgfsetdash{}{0pt}%
\pgfsys@defobject{currentmarker}{\pgfqpoint{0.000000in}{-0.048611in}}{\pgfqpoint{0.000000in}{0.000000in}}{%
\pgfpathmoveto{\pgfqpoint{0.000000in}{0.000000in}}%
\pgfpathlineto{\pgfqpoint{0.000000in}{-0.048611in}}%
\pgfusepath{stroke,fill}%
}%
\begin{pgfscope}%
\pgfsys@transformshift{1.749908in}{0.565123in}%
\pgfsys@useobject{currentmarker}{}%
\end{pgfscope}%
\end{pgfscope}%
\begin{pgfscope}%
\definecolor{textcolor}{rgb}{0.000000,0.000000,0.000000}%
\pgfsetstrokecolor{textcolor}%
\pgfsetfillcolor{textcolor}%
\pgftext[x=1.749908in,y=0.467901in,,top]{\color{textcolor}\rmfamily\fontsize{10.000000}{12.000000}\selectfont \(\displaystyle {0.4}\)}%
\end{pgfscope}%
\begin{pgfscope}%
\pgfpathrectangle{\pgfqpoint{0.494137in}{0.565123in}}{\pgfqpoint{5.655863in}{2.771193in}}%
\pgfusepath{clip}%
\pgfsetbuttcap%
\pgfsetroundjoin%
\pgfsetlinewidth{0.803000pt}%
\definecolor{currentstroke}{rgb}{0.690196,0.690196,0.690196}%
\pgfsetstrokecolor{currentstroke}%
\pgfsetdash{{0.800000pt}{1.320000pt}}{0.000000pt}%
\pgfpathmoveto{\pgfqpoint{2.379367in}{0.565123in}}%
\pgfpathlineto{\pgfqpoint{2.379367in}{3.336317in}}%
\pgfusepath{stroke}%
\end{pgfscope}%
\begin{pgfscope}%
\pgfsetbuttcap%
\pgfsetroundjoin%
\definecolor{currentfill}{rgb}{0.000000,0.000000,0.000000}%
\pgfsetfillcolor{currentfill}%
\pgfsetlinewidth{0.803000pt}%
\definecolor{currentstroke}{rgb}{0.000000,0.000000,0.000000}%
\pgfsetstrokecolor{currentstroke}%
\pgfsetdash{}{0pt}%
\pgfsys@defobject{currentmarker}{\pgfqpoint{0.000000in}{-0.048611in}}{\pgfqpoint{0.000000in}{0.000000in}}{%
\pgfpathmoveto{\pgfqpoint{0.000000in}{0.000000in}}%
\pgfpathlineto{\pgfqpoint{0.000000in}{-0.048611in}}%
\pgfusepath{stroke,fill}%
}%
\begin{pgfscope}%
\pgfsys@transformshift{2.379367in}{0.565123in}%
\pgfsys@useobject{currentmarker}{}%
\end{pgfscope}%
\end{pgfscope}%
\begin{pgfscope}%
\definecolor{textcolor}{rgb}{0.000000,0.000000,0.000000}%
\pgfsetstrokecolor{textcolor}%
\pgfsetfillcolor{textcolor}%
\pgftext[x=2.379367in,y=0.467901in,,top]{\color{textcolor}\rmfamily\fontsize{10.000000}{12.000000}\selectfont \(\displaystyle {0.6}\)}%
\end{pgfscope}%
\begin{pgfscope}%
\pgfpathrectangle{\pgfqpoint{0.494137in}{0.565123in}}{\pgfqpoint{5.655863in}{2.771193in}}%
\pgfusepath{clip}%
\pgfsetbuttcap%
\pgfsetroundjoin%
\pgfsetlinewidth{0.803000pt}%
\definecolor{currentstroke}{rgb}{0.690196,0.690196,0.690196}%
\pgfsetstrokecolor{currentstroke}%
\pgfsetdash{{0.800000pt}{1.320000pt}}{0.000000pt}%
\pgfpathmoveto{\pgfqpoint{3.008827in}{0.565123in}}%
\pgfpathlineto{\pgfqpoint{3.008827in}{3.336317in}}%
\pgfusepath{stroke}%
\end{pgfscope}%
\begin{pgfscope}%
\pgfsetbuttcap%
\pgfsetroundjoin%
\definecolor{currentfill}{rgb}{0.000000,0.000000,0.000000}%
\pgfsetfillcolor{currentfill}%
\pgfsetlinewidth{0.803000pt}%
\definecolor{currentstroke}{rgb}{0.000000,0.000000,0.000000}%
\pgfsetstrokecolor{currentstroke}%
\pgfsetdash{}{0pt}%
\pgfsys@defobject{currentmarker}{\pgfqpoint{0.000000in}{-0.048611in}}{\pgfqpoint{0.000000in}{0.000000in}}{%
\pgfpathmoveto{\pgfqpoint{0.000000in}{0.000000in}}%
\pgfpathlineto{\pgfqpoint{0.000000in}{-0.048611in}}%
\pgfusepath{stroke,fill}%
}%
\begin{pgfscope}%
\pgfsys@transformshift{3.008827in}{0.565123in}%
\pgfsys@useobject{currentmarker}{}%
\end{pgfscope}%
\end{pgfscope}%
\begin{pgfscope}%
\definecolor{textcolor}{rgb}{0.000000,0.000000,0.000000}%
\pgfsetstrokecolor{textcolor}%
\pgfsetfillcolor{textcolor}%
\pgftext[x=3.008827in,y=0.467901in,,top]{\color{textcolor}\rmfamily\fontsize{10.000000}{12.000000}\selectfont \(\displaystyle {0.8}\)}%
\end{pgfscope}%
\begin{pgfscope}%
\pgfpathrectangle{\pgfqpoint{0.494137in}{0.565123in}}{\pgfqpoint{5.655863in}{2.771193in}}%
\pgfusepath{clip}%
\pgfsetbuttcap%
\pgfsetroundjoin%
\pgfsetlinewidth{0.803000pt}%
\definecolor{currentstroke}{rgb}{0.690196,0.690196,0.690196}%
\pgfsetstrokecolor{currentstroke}%
\pgfsetdash{{0.800000pt}{1.320000pt}}{0.000000pt}%
\pgfpathmoveto{\pgfqpoint{3.638286in}{0.565123in}}%
\pgfpathlineto{\pgfqpoint{3.638286in}{3.336317in}}%
\pgfusepath{stroke}%
\end{pgfscope}%
\begin{pgfscope}%
\pgfsetbuttcap%
\pgfsetroundjoin%
\definecolor{currentfill}{rgb}{0.000000,0.000000,0.000000}%
\pgfsetfillcolor{currentfill}%
\pgfsetlinewidth{0.803000pt}%
\definecolor{currentstroke}{rgb}{0.000000,0.000000,0.000000}%
\pgfsetstrokecolor{currentstroke}%
\pgfsetdash{}{0pt}%
\pgfsys@defobject{currentmarker}{\pgfqpoint{0.000000in}{-0.048611in}}{\pgfqpoint{0.000000in}{0.000000in}}{%
\pgfpathmoveto{\pgfqpoint{0.000000in}{0.000000in}}%
\pgfpathlineto{\pgfqpoint{0.000000in}{-0.048611in}}%
\pgfusepath{stroke,fill}%
}%
\begin{pgfscope}%
\pgfsys@transformshift{3.638286in}{0.565123in}%
\pgfsys@useobject{currentmarker}{}%
\end{pgfscope}%
\end{pgfscope}%
\begin{pgfscope}%
\definecolor{textcolor}{rgb}{0.000000,0.000000,0.000000}%
\pgfsetstrokecolor{textcolor}%
\pgfsetfillcolor{textcolor}%
\pgftext[x=3.638286in,y=0.467901in,,top]{\color{textcolor}\rmfamily\fontsize{10.000000}{12.000000}\selectfont \(\displaystyle {1.0}\)}%
\end{pgfscope}%
\begin{pgfscope}%
\pgfpathrectangle{\pgfqpoint{0.494137in}{0.565123in}}{\pgfqpoint{5.655863in}{2.771193in}}%
\pgfusepath{clip}%
\pgfsetbuttcap%
\pgfsetroundjoin%
\pgfsetlinewidth{0.803000pt}%
\definecolor{currentstroke}{rgb}{0.690196,0.690196,0.690196}%
\pgfsetstrokecolor{currentstroke}%
\pgfsetdash{{0.800000pt}{1.320000pt}}{0.000000pt}%
\pgfpathmoveto{\pgfqpoint{4.267745in}{0.565123in}}%
\pgfpathlineto{\pgfqpoint{4.267745in}{3.336317in}}%
\pgfusepath{stroke}%
\end{pgfscope}%
\begin{pgfscope}%
\pgfsetbuttcap%
\pgfsetroundjoin%
\definecolor{currentfill}{rgb}{0.000000,0.000000,0.000000}%
\pgfsetfillcolor{currentfill}%
\pgfsetlinewidth{0.803000pt}%
\definecolor{currentstroke}{rgb}{0.000000,0.000000,0.000000}%
\pgfsetstrokecolor{currentstroke}%
\pgfsetdash{}{0pt}%
\pgfsys@defobject{currentmarker}{\pgfqpoint{0.000000in}{-0.048611in}}{\pgfqpoint{0.000000in}{0.000000in}}{%
\pgfpathmoveto{\pgfqpoint{0.000000in}{0.000000in}}%
\pgfpathlineto{\pgfqpoint{0.000000in}{-0.048611in}}%
\pgfusepath{stroke,fill}%
}%
\begin{pgfscope}%
\pgfsys@transformshift{4.267745in}{0.565123in}%
\pgfsys@useobject{currentmarker}{}%
\end{pgfscope}%
\end{pgfscope}%
\begin{pgfscope}%
\definecolor{textcolor}{rgb}{0.000000,0.000000,0.000000}%
\pgfsetstrokecolor{textcolor}%
\pgfsetfillcolor{textcolor}%
\pgftext[x=4.267745in,y=0.467901in,,top]{\color{textcolor}\rmfamily\fontsize{10.000000}{12.000000}\selectfont \(\displaystyle {1.2}\)}%
\end{pgfscope}%
\begin{pgfscope}%
\pgfpathrectangle{\pgfqpoint{0.494137in}{0.565123in}}{\pgfqpoint{5.655863in}{2.771193in}}%
\pgfusepath{clip}%
\pgfsetbuttcap%
\pgfsetroundjoin%
\pgfsetlinewidth{0.803000pt}%
\definecolor{currentstroke}{rgb}{0.690196,0.690196,0.690196}%
\pgfsetstrokecolor{currentstroke}%
\pgfsetdash{{0.800000pt}{1.320000pt}}{0.000000pt}%
\pgfpathmoveto{\pgfqpoint{4.897205in}{0.565123in}}%
\pgfpathlineto{\pgfqpoint{4.897205in}{3.336317in}}%
\pgfusepath{stroke}%
\end{pgfscope}%
\begin{pgfscope}%
\pgfsetbuttcap%
\pgfsetroundjoin%
\definecolor{currentfill}{rgb}{0.000000,0.000000,0.000000}%
\pgfsetfillcolor{currentfill}%
\pgfsetlinewidth{0.803000pt}%
\definecolor{currentstroke}{rgb}{0.000000,0.000000,0.000000}%
\pgfsetstrokecolor{currentstroke}%
\pgfsetdash{}{0pt}%
\pgfsys@defobject{currentmarker}{\pgfqpoint{0.000000in}{-0.048611in}}{\pgfqpoint{0.000000in}{0.000000in}}{%
\pgfpathmoveto{\pgfqpoint{0.000000in}{0.000000in}}%
\pgfpathlineto{\pgfqpoint{0.000000in}{-0.048611in}}%
\pgfusepath{stroke,fill}%
}%
\begin{pgfscope}%
\pgfsys@transformshift{4.897205in}{0.565123in}%
\pgfsys@useobject{currentmarker}{}%
\end{pgfscope}%
\end{pgfscope}%
\begin{pgfscope}%
\definecolor{textcolor}{rgb}{0.000000,0.000000,0.000000}%
\pgfsetstrokecolor{textcolor}%
\pgfsetfillcolor{textcolor}%
\pgftext[x=4.897205in,y=0.467901in,,top]{\color{textcolor}\rmfamily\fontsize{10.000000}{12.000000}\selectfont \(\displaystyle {1.4}\)}%
\end{pgfscope}%
\begin{pgfscope}%
\pgfpathrectangle{\pgfqpoint{0.494137in}{0.565123in}}{\pgfqpoint{5.655863in}{2.771193in}}%
\pgfusepath{clip}%
\pgfsetbuttcap%
\pgfsetroundjoin%
\pgfsetlinewidth{0.803000pt}%
\definecolor{currentstroke}{rgb}{0.690196,0.690196,0.690196}%
\pgfsetstrokecolor{currentstroke}%
\pgfsetdash{{0.800000pt}{1.320000pt}}{0.000000pt}%
\pgfpathmoveto{\pgfqpoint{5.526664in}{0.565123in}}%
\pgfpathlineto{\pgfqpoint{5.526664in}{3.336317in}}%
\pgfusepath{stroke}%
\end{pgfscope}%
\begin{pgfscope}%
\pgfsetbuttcap%
\pgfsetroundjoin%
\definecolor{currentfill}{rgb}{0.000000,0.000000,0.000000}%
\pgfsetfillcolor{currentfill}%
\pgfsetlinewidth{0.803000pt}%
\definecolor{currentstroke}{rgb}{0.000000,0.000000,0.000000}%
\pgfsetstrokecolor{currentstroke}%
\pgfsetdash{}{0pt}%
\pgfsys@defobject{currentmarker}{\pgfqpoint{0.000000in}{-0.048611in}}{\pgfqpoint{0.000000in}{0.000000in}}{%
\pgfpathmoveto{\pgfqpoint{0.000000in}{0.000000in}}%
\pgfpathlineto{\pgfqpoint{0.000000in}{-0.048611in}}%
\pgfusepath{stroke,fill}%
}%
\begin{pgfscope}%
\pgfsys@transformshift{5.526664in}{0.565123in}%
\pgfsys@useobject{currentmarker}{}%
\end{pgfscope}%
\end{pgfscope}%
\begin{pgfscope}%
\definecolor{textcolor}{rgb}{0.000000,0.000000,0.000000}%
\pgfsetstrokecolor{textcolor}%
\pgfsetfillcolor{textcolor}%
\pgftext[x=5.526664in,y=0.467901in,,top]{\color{textcolor}\rmfamily\fontsize{10.000000}{12.000000}\selectfont \(\displaystyle {1.6}\)}%
\end{pgfscope}%
\begin{pgfscope}%
\definecolor{textcolor}{rgb}{0.000000,0.000000,0.000000}%
\pgfsetstrokecolor{textcolor}%
\pgfsetfillcolor{textcolor}%
\pgftext[x=3.322068in,y=0.288889in,,top]{\color{textcolor}\rmfamily\fontsize{10.000000}{12.000000}\selectfont Return Loss (dB)}%
\end{pgfscope}%
\begin{pgfscope}%
\definecolor{textcolor}{rgb}{0.000000,0.000000,0.000000}%
\pgfsetstrokecolor{textcolor}%
\pgfsetfillcolor{textcolor}%
\pgftext[x=6.150000in,y=0.302778in,right,top]{\color{textcolor}\rmfamily\fontsize{10.000000}{12.000000}\selectfont \(\displaystyle \times{10^{8}}{}\)}%
\end{pgfscope}%
\begin{pgfscope}%
\pgfpathrectangle{\pgfqpoint{0.494137in}{0.565123in}}{\pgfqpoint{5.655863in}{2.771193in}}%
\pgfusepath{clip}%
\pgfsetbuttcap%
\pgfsetroundjoin%
\pgfsetlinewidth{0.803000pt}%
\definecolor{currentstroke}{rgb}{0.690196,0.690196,0.690196}%
\pgfsetstrokecolor{currentstroke}%
\pgfsetdash{{0.800000pt}{1.320000pt}}{0.000000pt}%
\pgfpathmoveto{\pgfqpoint{0.494137in}{0.662823in}}%
\pgfpathlineto{\pgfqpoint{6.150000in}{0.662823in}}%
\pgfusepath{stroke}%
\end{pgfscope}%
\begin{pgfscope}%
\pgfsetbuttcap%
\pgfsetroundjoin%
\definecolor{currentfill}{rgb}{0.000000,0.000000,0.000000}%
\pgfsetfillcolor{currentfill}%
\pgfsetlinewidth{0.803000pt}%
\definecolor{currentstroke}{rgb}{0.000000,0.000000,0.000000}%
\pgfsetstrokecolor{currentstroke}%
\pgfsetdash{}{0pt}%
\pgfsys@defobject{currentmarker}{\pgfqpoint{-0.048611in}{0.000000in}}{\pgfqpoint{-0.000000in}{0.000000in}}{%
\pgfpathmoveto{\pgfqpoint{-0.000000in}{0.000000in}}%
\pgfpathlineto{\pgfqpoint{-0.048611in}{0.000000in}}%
\pgfusepath{stroke,fill}%
}%
\begin{pgfscope}%
\pgfsys@transformshift{0.494137in}{0.662823in}%
\pgfsys@useobject{currentmarker}{}%
\end{pgfscope}%
\end{pgfscope}%
\begin{pgfscope}%
\definecolor{textcolor}{rgb}{0.000000,0.000000,0.000000}%
\pgfsetstrokecolor{textcolor}%
\pgfsetfillcolor{textcolor}%
\pgftext[x=0.150000in, y=0.614597in, left, base]{\color{textcolor}\rmfamily\fontsize{10.000000}{12.000000}\selectfont \(\displaystyle {\ensuremath{-}12}\)}%
\end{pgfscope}%
\begin{pgfscope}%
\pgfpathrectangle{\pgfqpoint{0.494137in}{0.565123in}}{\pgfqpoint{5.655863in}{2.771193in}}%
\pgfusepath{clip}%
\pgfsetbuttcap%
\pgfsetroundjoin%
\pgfsetlinewidth{0.803000pt}%
\definecolor{currentstroke}{rgb}{0.690196,0.690196,0.690196}%
\pgfsetstrokecolor{currentstroke}%
\pgfsetdash{{0.800000pt}{1.320000pt}}{0.000000pt}%
\pgfpathmoveto{\pgfqpoint{0.494137in}{1.039676in}}%
\pgfpathlineto{\pgfqpoint{6.150000in}{1.039676in}}%
\pgfusepath{stroke}%
\end{pgfscope}%
\begin{pgfscope}%
\pgfsetbuttcap%
\pgfsetroundjoin%
\definecolor{currentfill}{rgb}{0.000000,0.000000,0.000000}%
\pgfsetfillcolor{currentfill}%
\pgfsetlinewidth{0.803000pt}%
\definecolor{currentstroke}{rgb}{0.000000,0.000000,0.000000}%
\pgfsetstrokecolor{currentstroke}%
\pgfsetdash{}{0pt}%
\pgfsys@defobject{currentmarker}{\pgfqpoint{-0.048611in}{0.000000in}}{\pgfqpoint{-0.000000in}{0.000000in}}{%
\pgfpathmoveto{\pgfqpoint{-0.000000in}{0.000000in}}%
\pgfpathlineto{\pgfqpoint{-0.048611in}{0.000000in}}%
\pgfusepath{stroke,fill}%
}%
\begin{pgfscope}%
\pgfsys@transformshift{0.494137in}{1.039676in}%
\pgfsys@useobject{currentmarker}{}%
\end{pgfscope}%
\end{pgfscope}%
\begin{pgfscope}%
\definecolor{textcolor}{rgb}{0.000000,0.000000,0.000000}%
\pgfsetstrokecolor{textcolor}%
\pgfsetfillcolor{textcolor}%
\pgftext[x=0.150000in, y=0.991451in, left, base]{\color{textcolor}\rmfamily\fontsize{10.000000}{12.000000}\selectfont \(\displaystyle {\ensuremath{-}10}\)}%
\end{pgfscope}%
\begin{pgfscope}%
\pgfpathrectangle{\pgfqpoint{0.494137in}{0.565123in}}{\pgfqpoint{5.655863in}{2.771193in}}%
\pgfusepath{clip}%
\pgfsetbuttcap%
\pgfsetroundjoin%
\pgfsetlinewidth{0.803000pt}%
\definecolor{currentstroke}{rgb}{0.690196,0.690196,0.690196}%
\pgfsetstrokecolor{currentstroke}%
\pgfsetdash{{0.800000pt}{1.320000pt}}{0.000000pt}%
\pgfpathmoveto{\pgfqpoint{0.494137in}{1.416530in}}%
\pgfpathlineto{\pgfqpoint{6.150000in}{1.416530in}}%
\pgfusepath{stroke}%
\end{pgfscope}%
\begin{pgfscope}%
\pgfsetbuttcap%
\pgfsetroundjoin%
\definecolor{currentfill}{rgb}{0.000000,0.000000,0.000000}%
\pgfsetfillcolor{currentfill}%
\pgfsetlinewidth{0.803000pt}%
\definecolor{currentstroke}{rgb}{0.000000,0.000000,0.000000}%
\pgfsetstrokecolor{currentstroke}%
\pgfsetdash{}{0pt}%
\pgfsys@defobject{currentmarker}{\pgfqpoint{-0.048611in}{0.000000in}}{\pgfqpoint{-0.000000in}{0.000000in}}{%
\pgfpathmoveto{\pgfqpoint{-0.000000in}{0.000000in}}%
\pgfpathlineto{\pgfqpoint{-0.048611in}{0.000000in}}%
\pgfusepath{stroke,fill}%
}%
\begin{pgfscope}%
\pgfsys@transformshift{0.494137in}{1.416530in}%
\pgfsys@useobject{currentmarker}{}%
\end{pgfscope}%
\end{pgfscope}%
\begin{pgfscope}%
\definecolor{textcolor}{rgb}{0.000000,0.000000,0.000000}%
\pgfsetstrokecolor{textcolor}%
\pgfsetfillcolor{textcolor}%
\pgftext[x=0.219445in, y=1.368305in, left, base]{\color{textcolor}\rmfamily\fontsize{10.000000}{12.000000}\selectfont \(\displaystyle {\ensuremath{-}8}\)}%
\end{pgfscope}%
\begin{pgfscope}%
\pgfpathrectangle{\pgfqpoint{0.494137in}{0.565123in}}{\pgfqpoint{5.655863in}{2.771193in}}%
\pgfusepath{clip}%
\pgfsetbuttcap%
\pgfsetroundjoin%
\pgfsetlinewidth{0.803000pt}%
\definecolor{currentstroke}{rgb}{0.690196,0.690196,0.690196}%
\pgfsetstrokecolor{currentstroke}%
\pgfsetdash{{0.800000pt}{1.320000pt}}{0.000000pt}%
\pgfpathmoveto{\pgfqpoint{0.494137in}{1.793384in}}%
\pgfpathlineto{\pgfqpoint{6.150000in}{1.793384in}}%
\pgfusepath{stroke}%
\end{pgfscope}%
\begin{pgfscope}%
\pgfsetbuttcap%
\pgfsetroundjoin%
\definecolor{currentfill}{rgb}{0.000000,0.000000,0.000000}%
\pgfsetfillcolor{currentfill}%
\pgfsetlinewidth{0.803000pt}%
\definecolor{currentstroke}{rgb}{0.000000,0.000000,0.000000}%
\pgfsetstrokecolor{currentstroke}%
\pgfsetdash{}{0pt}%
\pgfsys@defobject{currentmarker}{\pgfqpoint{-0.048611in}{0.000000in}}{\pgfqpoint{-0.000000in}{0.000000in}}{%
\pgfpathmoveto{\pgfqpoint{-0.000000in}{0.000000in}}%
\pgfpathlineto{\pgfqpoint{-0.048611in}{0.000000in}}%
\pgfusepath{stroke,fill}%
}%
\begin{pgfscope}%
\pgfsys@transformshift{0.494137in}{1.793384in}%
\pgfsys@useobject{currentmarker}{}%
\end{pgfscope}%
\end{pgfscope}%
\begin{pgfscope}%
\definecolor{textcolor}{rgb}{0.000000,0.000000,0.000000}%
\pgfsetstrokecolor{textcolor}%
\pgfsetfillcolor{textcolor}%
\pgftext[x=0.219445in, y=1.745158in, left, base]{\color{textcolor}\rmfamily\fontsize{10.000000}{12.000000}\selectfont \(\displaystyle {\ensuremath{-}6}\)}%
\end{pgfscope}%
\begin{pgfscope}%
\pgfpathrectangle{\pgfqpoint{0.494137in}{0.565123in}}{\pgfqpoint{5.655863in}{2.771193in}}%
\pgfusepath{clip}%
\pgfsetbuttcap%
\pgfsetroundjoin%
\pgfsetlinewidth{0.803000pt}%
\definecolor{currentstroke}{rgb}{0.690196,0.690196,0.690196}%
\pgfsetstrokecolor{currentstroke}%
\pgfsetdash{{0.800000pt}{1.320000pt}}{0.000000pt}%
\pgfpathmoveto{\pgfqpoint{0.494137in}{2.170237in}}%
\pgfpathlineto{\pgfqpoint{6.150000in}{2.170237in}}%
\pgfusepath{stroke}%
\end{pgfscope}%
\begin{pgfscope}%
\pgfsetbuttcap%
\pgfsetroundjoin%
\definecolor{currentfill}{rgb}{0.000000,0.000000,0.000000}%
\pgfsetfillcolor{currentfill}%
\pgfsetlinewidth{0.803000pt}%
\definecolor{currentstroke}{rgb}{0.000000,0.000000,0.000000}%
\pgfsetstrokecolor{currentstroke}%
\pgfsetdash{}{0pt}%
\pgfsys@defobject{currentmarker}{\pgfqpoint{-0.048611in}{0.000000in}}{\pgfqpoint{-0.000000in}{0.000000in}}{%
\pgfpathmoveto{\pgfqpoint{-0.000000in}{0.000000in}}%
\pgfpathlineto{\pgfqpoint{-0.048611in}{0.000000in}}%
\pgfusepath{stroke,fill}%
}%
\begin{pgfscope}%
\pgfsys@transformshift{0.494137in}{2.170237in}%
\pgfsys@useobject{currentmarker}{}%
\end{pgfscope}%
\end{pgfscope}%
\begin{pgfscope}%
\definecolor{textcolor}{rgb}{0.000000,0.000000,0.000000}%
\pgfsetstrokecolor{textcolor}%
\pgfsetfillcolor{textcolor}%
\pgftext[x=0.219445in, y=2.122012in, left, base]{\color{textcolor}\rmfamily\fontsize{10.000000}{12.000000}\selectfont \(\displaystyle {\ensuremath{-}4}\)}%
\end{pgfscope}%
\begin{pgfscope}%
\pgfpathrectangle{\pgfqpoint{0.494137in}{0.565123in}}{\pgfqpoint{5.655863in}{2.771193in}}%
\pgfusepath{clip}%
\pgfsetbuttcap%
\pgfsetroundjoin%
\pgfsetlinewidth{0.803000pt}%
\definecolor{currentstroke}{rgb}{0.690196,0.690196,0.690196}%
\pgfsetstrokecolor{currentstroke}%
\pgfsetdash{{0.800000pt}{1.320000pt}}{0.000000pt}%
\pgfpathmoveto{\pgfqpoint{0.494137in}{2.547091in}}%
\pgfpathlineto{\pgfqpoint{6.150000in}{2.547091in}}%
\pgfusepath{stroke}%
\end{pgfscope}%
\begin{pgfscope}%
\pgfsetbuttcap%
\pgfsetroundjoin%
\definecolor{currentfill}{rgb}{0.000000,0.000000,0.000000}%
\pgfsetfillcolor{currentfill}%
\pgfsetlinewidth{0.803000pt}%
\definecolor{currentstroke}{rgb}{0.000000,0.000000,0.000000}%
\pgfsetstrokecolor{currentstroke}%
\pgfsetdash{}{0pt}%
\pgfsys@defobject{currentmarker}{\pgfqpoint{-0.048611in}{0.000000in}}{\pgfqpoint{-0.000000in}{0.000000in}}{%
\pgfpathmoveto{\pgfqpoint{-0.000000in}{0.000000in}}%
\pgfpathlineto{\pgfqpoint{-0.048611in}{0.000000in}}%
\pgfusepath{stroke,fill}%
}%
\begin{pgfscope}%
\pgfsys@transformshift{0.494137in}{2.547091in}%
\pgfsys@useobject{currentmarker}{}%
\end{pgfscope}%
\end{pgfscope}%
\begin{pgfscope}%
\definecolor{textcolor}{rgb}{0.000000,0.000000,0.000000}%
\pgfsetstrokecolor{textcolor}%
\pgfsetfillcolor{textcolor}%
\pgftext[x=0.219445in, y=2.498866in, left, base]{\color{textcolor}\rmfamily\fontsize{10.000000}{12.000000}\selectfont \(\displaystyle {\ensuremath{-}2}\)}%
\end{pgfscope}%
\begin{pgfscope}%
\pgfpathrectangle{\pgfqpoint{0.494137in}{0.565123in}}{\pgfqpoint{5.655863in}{2.771193in}}%
\pgfusepath{clip}%
\pgfsetbuttcap%
\pgfsetroundjoin%
\pgfsetlinewidth{0.803000pt}%
\definecolor{currentstroke}{rgb}{0.690196,0.690196,0.690196}%
\pgfsetstrokecolor{currentstroke}%
\pgfsetdash{{0.800000pt}{1.320000pt}}{0.000000pt}%
\pgfpathmoveto{\pgfqpoint{0.494137in}{2.923945in}}%
\pgfpathlineto{\pgfqpoint{6.150000in}{2.923945in}}%
\pgfusepath{stroke}%
\end{pgfscope}%
\begin{pgfscope}%
\pgfsetbuttcap%
\pgfsetroundjoin%
\definecolor{currentfill}{rgb}{0.000000,0.000000,0.000000}%
\pgfsetfillcolor{currentfill}%
\pgfsetlinewidth{0.803000pt}%
\definecolor{currentstroke}{rgb}{0.000000,0.000000,0.000000}%
\pgfsetstrokecolor{currentstroke}%
\pgfsetdash{}{0pt}%
\pgfsys@defobject{currentmarker}{\pgfqpoint{-0.048611in}{0.000000in}}{\pgfqpoint{-0.000000in}{0.000000in}}{%
\pgfpathmoveto{\pgfqpoint{-0.000000in}{0.000000in}}%
\pgfpathlineto{\pgfqpoint{-0.048611in}{0.000000in}}%
\pgfusepath{stroke,fill}%
}%
\begin{pgfscope}%
\pgfsys@transformshift{0.494137in}{2.923945in}%
\pgfsys@useobject{currentmarker}{}%
\end{pgfscope}%
\end{pgfscope}%
\begin{pgfscope}%
\definecolor{textcolor}{rgb}{0.000000,0.000000,0.000000}%
\pgfsetstrokecolor{textcolor}%
\pgfsetfillcolor{textcolor}%
\pgftext[x=0.327470in, y=2.875719in, left, base]{\color{textcolor}\rmfamily\fontsize{10.000000}{12.000000}\selectfont \(\displaystyle {0}\)}%
\end{pgfscope}%
\begin{pgfscope}%
\pgfpathrectangle{\pgfqpoint{0.494137in}{0.565123in}}{\pgfqpoint{5.655863in}{2.771193in}}%
\pgfusepath{clip}%
\pgfsetbuttcap%
\pgfsetroundjoin%
\pgfsetlinewidth{0.803000pt}%
\definecolor{currentstroke}{rgb}{0.690196,0.690196,0.690196}%
\pgfsetstrokecolor{currentstroke}%
\pgfsetdash{{0.800000pt}{1.320000pt}}{0.000000pt}%
\pgfpathmoveto{\pgfqpoint{0.494137in}{3.300798in}}%
\pgfpathlineto{\pgfqpoint{6.150000in}{3.300798in}}%
\pgfusepath{stroke}%
\end{pgfscope}%
\begin{pgfscope}%
\pgfsetbuttcap%
\pgfsetroundjoin%
\definecolor{currentfill}{rgb}{0.000000,0.000000,0.000000}%
\pgfsetfillcolor{currentfill}%
\pgfsetlinewidth{0.803000pt}%
\definecolor{currentstroke}{rgb}{0.000000,0.000000,0.000000}%
\pgfsetstrokecolor{currentstroke}%
\pgfsetdash{}{0pt}%
\pgfsys@defobject{currentmarker}{\pgfqpoint{-0.048611in}{0.000000in}}{\pgfqpoint{-0.000000in}{0.000000in}}{%
\pgfpathmoveto{\pgfqpoint{-0.000000in}{0.000000in}}%
\pgfpathlineto{\pgfqpoint{-0.048611in}{0.000000in}}%
\pgfusepath{stroke,fill}%
}%
\begin{pgfscope}%
\pgfsys@transformshift{0.494137in}{3.300798in}%
\pgfsys@useobject{currentmarker}{}%
\end{pgfscope}%
\end{pgfscope}%
\begin{pgfscope}%
\definecolor{textcolor}{rgb}{0.000000,0.000000,0.000000}%
\pgfsetstrokecolor{textcolor}%
\pgfsetfillcolor{textcolor}%
\pgftext[x=0.327470in, y=3.252573in, left, base]{\color{textcolor}\rmfamily\fontsize{10.000000}{12.000000}\selectfont \(\displaystyle {2}\)}%
\end{pgfscope}%
\begin{pgfscope}%
\pgfpathrectangle{\pgfqpoint{0.494137in}{0.565123in}}{\pgfqpoint{5.655863in}{2.771193in}}%
\pgfusepath{clip}%
\pgfsetrectcap%
\pgfsetroundjoin%
\pgfsetlinewidth{1.003750pt}%
\definecolor{currentstroke}{rgb}{0.000000,0.000000,0.000000}%
\pgfsetstrokecolor{currentstroke}%
\pgfsetdash{}{0pt}%
\pgfpathmoveto{\pgfqpoint{0.494137in}{3.210353in}}%
\pgfpathlineto{\pgfqpoint{0.500238in}{2.946556in}}%
\pgfpathlineto{\pgfqpoint{0.512440in}{2.946556in}}%
\pgfpathlineto{\pgfqpoint{0.518542in}{2.923945in}}%
\pgfpathlineto{\pgfqpoint{0.524643in}{2.935250in}}%
\pgfpathlineto{\pgfqpoint{0.530744in}{2.923945in}}%
\pgfpathlineto{\pgfqpoint{0.536845in}{2.935250in}}%
\pgfpathlineto{\pgfqpoint{0.542947in}{2.923945in}}%
\pgfpathlineto{\pgfqpoint{0.555149in}{2.923945in}}%
\pgfpathlineto{\pgfqpoint{0.561250in}{2.935250in}}%
\pgfpathlineto{\pgfqpoint{0.567352in}{2.935250in}}%
\pgfpathlineto{\pgfqpoint{0.573453in}{2.923945in}}%
\pgfpathlineto{\pgfqpoint{0.823604in}{2.923945in}}%
\pgfpathlineto{\pgfqpoint{0.829706in}{2.935250in}}%
\pgfpathlineto{\pgfqpoint{0.835807in}{2.923945in}}%
\pgfpathlineto{\pgfqpoint{0.841908in}{2.935250in}}%
\pgfpathlineto{\pgfqpoint{0.854111in}{2.935250in}}%
\pgfpathlineto{\pgfqpoint{0.860212in}{2.923945in}}%
\pgfpathlineto{\pgfqpoint{0.866313in}{2.935250in}}%
\pgfpathlineto{\pgfqpoint{0.872414in}{2.935250in}}%
\pgfpathlineto{\pgfqpoint{0.878516in}{2.923945in}}%
\pgfpathlineto{\pgfqpoint{0.988338in}{2.923945in}}%
\pgfpathlineto{\pgfqpoint{0.994439in}{2.935250in}}%
\pgfpathlineto{\pgfqpoint{1.000541in}{2.923945in}}%
\pgfpathlineto{\pgfqpoint{1.006642in}{2.935250in}}%
\pgfpathlineto{\pgfqpoint{1.012743in}{2.923945in}}%
\pgfpathlineto{\pgfqpoint{1.018844in}{2.923945in}}%
\pgfpathlineto{\pgfqpoint{1.024946in}{2.935250in}}%
\pgfpathlineto{\pgfqpoint{1.079857in}{2.935250in}}%
\pgfpathlineto{\pgfqpoint{1.085958in}{2.923945in}}%
\pgfpathlineto{\pgfqpoint{1.092060in}{2.935250in}}%
\pgfpathlineto{\pgfqpoint{1.098161in}{2.923945in}}%
\pgfpathlineto{\pgfqpoint{1.104262in}{2.935250in}}%
\pgfpathlineto{\pgfqpoint{1.140870in}{2.935250in}}%
\pgfpathlineto{\pgfqpoint{1.146971in}{2.923945in}}%
\pgfpathlineto{\pgfqpoint{1.153072in}{2.935250in}}%
\pgfpathlineto{\pgfqpoint{1.159173in}{2.923945in}}%
\pgfpathlineto{\pgfqpoint{1.220186in}{2.923945in}}%
\pgfpathlineto{\pgfqpoint{1.226287in}{2.935250in}}%
\pgfpathlineto{\pgfqpoint{1.336110in}{2.935250in}}%
\pgfpathlineto{\pgfqpoint{1.342211in}{2.923945in}}%
\pgfpathlineto{\pgfqpoint{1.348312in}{2.923945in}}%
\pgfpathlineto{\pgfqpoint{1.354413in}{2.935250in}}%
\pgfpathlineto{\pgfqpoint{1.482540in}{2.935250in}}%
\pgfpathlineto{\pgfqpoint{1.488641in}{2.923945in}}%
\pgfpathlineto{\pgfqpoint{1.506945in}{2.923945in}}%
\pgfpathlineto{\pgfqpoint{1.513046in}{2.935250in}}%
\pgfpathlineto{\pgfqpoint{1.561856in}{2.935250in}}%
\pgfpathlineto{\pgfqpoint{1.567957in}{2.946556in}}%
\pgfpathlineto{\pgfqpoint{1.586261in}{2.946556in}}%
\pgfpathlineto{\pgfqpoint{1.592362in}{2.935250in}}%
\pgfpathlineto{\pgfqpoint{1.647274in}{2.935250in}}%
\pgfpathlineto{\pgfqpoint{1.653375in}{2.946556in}}%
\pgfpathlineto{\pgfqpoint{1.696084in}{2.946556in}}%
\pgfpathlineto{\pgfqpoint{1.702185in}{2.935250in}}%
\pgfpathlineto{\pgfqpoint{1.726590in}{2.935250in}}%
\pgfpathlineto{\pgfqpoint{1.732691in}{2.946556in}}%
\pgfpathlineto{\pgfqpoint{1.738793in}{2.935250in}}%
\pgfpathlineto{\pgfqpoint{1.744894in}{2.946556in}}%
\pgfpathlineto{\pgfqpoint{1.787603in}{2.946556in}}%
\pgfpathlineto{\pgfqpoint{1.793704in}{2.935250in}}%
\pgfpathlineto{\pgfqpoint{1.836413in}{2.935250in}}%
\pgfpathlineto{\pgfqpoint{1.842514in}{2.946556in}}%
\pgfpathlineto{\pgfqpoint{1.879121in}{2.946556in}}%
\pgfpathlineto{\pgfqpoint{1.885223in}{2.935250in}}%
\pgfpathlineto{\pgfqpoint{1.909628in}{2.935250in}}%
\pgfpathlineto{\pgfqpoint{1.915729in}{2.946556in}}%
\pgfpathlineto{\pgfqpoint{1.964539in}{2.946556in}}%
\pgfpathlineto{\pgfqpoint{1.970640in}{2.935250in}}%
\pgfpathlineto{\pgfqpoint{2.037754in}{2.935250in}}%
\pgfpathlineto{\pgfqpoint{2.043855in}{2.923945in}}%
\pgfpathlineto{\pgfqpoint{2.080463in}{2.923945in}}%
\pgfpathlineto{\pgfqpoint{2.086564in}{2.912639in}}%
\pgfpathlineto{\pgfqpoint{2.104868in}{2.912639in}}%
\pgfpathlineto{\pgfqpoint{2.110969in}{2.901333in}}%
\pgfpathlineto{\pgfqpoint{2.123172in}{2.901333in}}%
\pgfpathlineto{\pgfqpoint{2.129273in}{2.880606in}}%
\pgfpathlineto{\pgfqpoint{2.141475in}{2.880606in}}%
\pgfpathlineto{\pgfqpoint{2.147577in}{2.869301in}}%
\pgfpathlineto{\pgfqpoint{2.159779in}{2.869301in}}%
\pgfpathlineto{\pgfqpoint{2.165880in}{2.857995in}}%
\pgfpathlineto{\pgfqpoint{2.178083in}{2.857995in}}%
\pgfpathlineto{\pgfqpoint{2.184184in}{2.846690in}}%
\pgfpathlineto{\pgfqpoint{2.190285in}{2.846690in}}%
\pgfpathlineto{\pgfqpoint{2.202488in}{2.824078in}}%
\pgfpathlineto{\pgfqpoint{2.208589in}{2.824078in}}%
\pgfpathlineto{\pgfqpoint{2.214690in}{2.812773in}}%
\pgfpathlineto{\pgfqpoint{2.220792in}{2.812773in}}%
\pgfpathlineto{\pgfqpoint{2.226893in}{2.801467in}}%
\pgfpathlineto{\pgfqpoint{2.232994in}{2.801467in}}%
\pgfpathlineto{\pgfqpoint{2.239095in}{2.792046in}}%
\pgfpathlineto{\pgfqpoint{2.251298in}{2.792046in}}%
\pgfpathlineto{\pgfqpoint{2.257399in}{2.780740in}}%
\pgfpathlineto{\pgfqpoint{2.263501in}{2.780740in}}%
\pgfpathlineto{\pgfqpoint{2.269602in}{2.769435in}}%
\pgfpathlineto{\pgfqpoint{2.275703in}{2.746823in}}%
\pgfpathlineto{\pgfqpoint{2.281804in}{2.746823in}}%
\pgfpathlineto{\pgfqpoint{2.287906in}{2.735518in}}%
\pgfpathlineto{\pgfqpoint{2.294007in}{2.735518in}}%
\pgfpathlineto{\pgfqpoint{2.300108in}{2.724212in}}%
\pgfpathlineto{\pgfqpoint{2.312311in}{2.724212in}}%
\pgfpathlineto{\pgfqpoint{2.318412in}{2.714791in}}%
\pgfpathlineto{\pgfqpoint{2.336716in}{2.714791in}}%
\pgfpathlineto{\pgfqpoint{2.342817in}{2.703485in}}%
\pgfpathlineto{\pgfqpoint{2.385526in}{2.703485in}}%
\pgfpathlineto{\pgfqpoint{2.391627in}{2.714791in}}%
\pgfpathlineto{\pgfqpoint{2.403829in}{2.714791in}}%
\pgfpathlineto{\pgfqpoint{2.409931in}{2.703485in}}%
\pgfpathlineto{\pgfqpoint{2.428234in}{2.703485in}}%
\pgfpathlineto{\pgfqpoint{2.434336in}{2.714791in}}%
\pgfpathlineto{\pgfqpoint{2.464842in}{2.714791in}}%
\pgfpathlineto{\pgfqpoint{2.470943in}{2.724212in}}%
\pgfpathlineto{\pgfqpoint{2.483146in}{2.724212in}}%
\pgfpathlineto{\pgfqpoint{2.489247in}{2.714791in}}%
\pgfpathlineto{\pgfqpoint{2.519753in}{2.714791in}}%
\pgfpathlineto{\pgfqpoint{2.525854in}{2.703485in}}%
\pgfpathlineto{\pgfqpoint{2.544158in}{2.703485in}}%
\pgfpathlineto{\pgfqpoint{2.550259in}{2.692180in}}%
\pgfpathlineto{\pgfqpoint{2.562462in}{2.692180in}}%
\pgfpathlineto{\pgfqpoint{2.568563in}{2.680874in}}%
\pgfpathlineto{\pgfqpoint{2.574665in}{2.680874in}}%
\pgfpathlineto{\pgfqpoint{2.580766in}{2.669568in}}%
\pgfpathlineto{\pgfqpoint{2.592968in}{2.669568in}}%
\pgfpathlineto{\pgfqpoint{2.599070in}{2.658263in}}%
\pgfpathlineto{\pgfqpoint{2.611272in}{2.658263in}}%
\pgfpathlineto{\pgfqpoint{2.617373in}{2.646957in}}%
\pgfpathlineto{\pgfqpoint{2.623475in}{2.646957in}}%
\pgfpathlineto{\pgfqpoint{2.629576in}{2.637536in}}%
\pgfpathlineto{\pgfqpoint{2.635677in}{2.637536in}}%
\pgfpathlineto{\pgfqpoint{2.653981in}{2.603619in}}%
\pgfpathlineto{\pgfqpoint{2.660082in}{2.581008in}}%
\pgfpathlineto{\pgfqpoint{2.678386in}{2.548975in}}%
\pgfpathlineto{\pgfqpoint{2.684487in}{2.526364in}}%
\pgfpathlineto{\pgfqpoint{2.696690in}{2.503753in}}%
\pgfpathlineto{\pgfqpoint{2.702791in}{2.481141in}}%
\pgfpathlineto{\pgfqpoint{2.708892in}{2.471720in}}%
\pgfpathlineto{\pgfqpoint{2.714993in}{2.449109in}}%
\pgfpathlineto{\pgfqpoint{2.739398in}{2.403886in}}%
\pgfpathlineto{\pgfqpoint{2.745500in}{2.383160in}}%
\pgfpathlineto{\pgfqpoint{2.769905in}{2.337937in}}%
\pgfpathlineto{\pgfqpoint{2.776006in}{2.315326in}}%
\pgfpathlineto{\pgfqpoint{2.782107in}{2.305905in}}%
\pgfpathlineto{\pgfqpoint{2.788208in}{2.283293in}}%
\pgfpathlineto{\pgfqpoint{2.794310in}{2.283293in}}%
\pgfpathlineto{\pgfqpoint{2.824816in}{2.226765in}}%
\pgfpathlineto{\pgfqpoint{2.849221in}{2.140089in}}%
\pgfpathlineto{\pgfqpoint{2.855322in}{2.117478in}}%
\pgfpathlineto{\pgfqpoint{2.861424in}{2.083561in}}%
\pgfpathlineto{\pgfqpoint{2.867525in}{2.062834in}}%
\pgfpathlineto{\pgfqpoint{2.873626in}{2.051528in}}%
\pgfpathlineto{\pgfqpoint{2.879727in}{2.028917in}}%
\pgfpathlineto{\pgfqpoint{2.891930in}{2.006306in}}%
\pgfpathlineto{\pgfqpoint{2.916335in}{2.006306in}}%
\pgfpathlineto{\pgfqpoint{2.922436in}{2.017611in}}%
\pgfpathlineto{\pgfqpoint{2.971246in}{2.017611in}}%
\pgfpathlineto{\pgfqpoint{2.977347in}{2.006306in}}%
\pgfpathlineto{\pgfqpoint{2.995651in}{2.006306in}}%
\pgfpathlineto{\pgfqpoint{3.013955in}{1.974273in}}%
\pgfpathlineto{\pgfqpoint{3.044461in}{1.863101in}}%
\pgfpathlineto{\pgfqpoint{3.050562in}{1.829185in}}%
\pgfpathlineto{\pgfqpoint{3.068866in}{1.763235in}}%
\pgfpathlineto{\pgfqpoint{3.074967in}{1.731203in}}%
\pgfpathlineto{\pgfqpoint{3.087170in}{1.685980in}}%
\pgfpathlineto{\pgfqpoint{3.093271in}{1.653948in}}%
\pgfpathlineto{\pgfqpoint{3.099372in}{1.642642in}}%
\pgfpathlineto{\pgfqpoint{3.105474in}{1.608725in}}%
\pgfpathlineto{\pgfqpoint{3.123777in}{1.542776in}}%
\pgfpathlineto{\pgfqpoint{3.135980in}{1.476827in}}%
\pgfpathlineto{\pgfqpoint{3.142081in}{1.431604in}}%
\pgfpathlineto{\pgfqpoint{3.148182in}{1.399572in}}%
\pgfpathlineto{\pgfqpoint{3.154284in}{1.343043in}}%
\pgfpathlineto{\pgfqpoint{3.160385in}{1.299705in}}%
\pgfpathlineto{\pgfqpoint{3.178689in}{1.145195in}}%
\pgfpathlineto{\pgfqpoint{3.184790in}{1.111278in}}%
\pgfpathlineto{\pgfqpoint{3.190891in}{1.067940in}}%
\pgfpathlineto{\pgfqpoint{3.196993in}{1.045329in}}%
\pgfpathlineto{\pgfqpoint{3.215296in}{1.011412in}}%
\pgfpathlineto{\pgfqpoint{3.221398in}{1.011412in}}%
\pgfpathlineto{\pgfqpoint{3.227499in}{1.000107in}}%
\pgfpathlineto{\pgfqpoint{3.233600in}{1.000107in}}%
\pgfpathlineto{\pgfqpoint{3.245803in}{0.979380in}}%
\pgfpathlineto{\pgfqpoint{3.264106in}{0.945463in}}%
\pgfpathlineto{\pgfqpoint{3.270208in}{0.945463in}}%
\pgfpathlineto{\pgfqpoint{3.276309in}{0.934157in}}%
\pgfpathlineto{\pgfqpoint{3.294613in}{0.934157in}}%
\pgfpathlineto{\pgfqpoint{3.319018in}{0.979380in}}%
\pgfpathlineto{\pgfqpoint{3.325119in}{0.979380in}}%
\pgfpathlineto{\pgfqpoint{3.331220in}{0.968074in}}%
\pgfpathlineto{\pgfqpoint{3.337321in}{0.934157in}}%
\pgfpathlineto{\pgfqpoint{3.343423in}{0.890819in}}%
\pgfpathlineto{\pgfqpoint{3.355625in}{0.779647in}}%
\pgfpathlineto{\pgfqpoint{3.361726in}{0.736309in}}%
\pgfpathlineto{\pgfqpoint{3.367828in}{0.702392in}}%
\pgfpathlineto{\pgfqpoint{3.373929in}{0.691087in}}%
\pgfpathlineto{\pgfqpoint{3.380030in}{0.691087in}}%
\pgfpathlineto{\pgfqpoint{3.386131in}{0.702392in}}%
\pgfpathlineto{\pgfqpoint{3.398334in}{0.768342in}}%
\pgfpathlineto{\pgfqpoint{3.404435in}{0.824870in}}%
\pgfpathlineto{\pgfqpoint{3.410536in}{0.856902in}}%
\pgfpathlineto{\pgfqpoint{3.416638in}{0.913430in}}%
\pgfpathlineto{\pgfqpoint{3.422739in}{0.956768in}}%
\pgfpathlineto{\pgfqpoint{3.428840in}{0.990685in}}%
\pgfpathlineto{\pgfqpoint{3.434941in}{1.034023in}}%
\pgfpathlineto{\pgfqpoint{3.441043in}{1.067940in}}%
\pgfpathlineto{\pgfqpoint{3.447144in}{1.088667in}}%
\pgfpathlineto{\pgfqpoint{3.465448in}{1.188533in}}%
\pgfpathlineto{\pgfqpoint{3.471549in}{1.211145in}}%
\pgfpathlineto{\pgfqpoint{3.489853in}{1.311011in}}%
\pgfpathlineto{\pgfqpoint{3.502055in}{1.354349in}}%
\pgfpathlineto{\pgfqpoint{3.508157in}{1.388266in}}%
\pgfpathlineto{\pgfqpoint{3.526460in}{1.454215in}}%
\pgfpathlineto{\pgfqpoint{3.532562in}{1.465521in}}%
\pgfpathlineto{\pgfqpoint{3.538663in}{1.488132in}}%
\pgfpathlineto{\pgfqpoint{3.550865in}{1.508859in}}%
\pgfpathlineto{\pgfqpoint{3.563068in}{1.531470in}}%
\pgfpathlineto{\pgfqpoint{3.569169in}{1.531470in}}%
\pgfpathlineto{\pgfqpoint{3.575270in}{1.542776in}}%
\pgfpathlineto{\pgfqpoint{3.581372in}{1.542776in}}%
\pgfpathlineto{\pgfqpoint{3.593574in}{1.565387in}}%
\pgfpathlineto{\pgfqpoint{3.599675in}{1.565387in}}%
\pgfpathlineto{\pgfqpoint{3.605777in}{1.574808in}}%
\pgfpathlineto{\pgfqpoint{3.660688in}{1.574808in}}%
\pgfpathlineto{\pgfqpoint{3.666789in}{1.586114in}}%
\pgfpathlineto{\pgfqpoint{3.691194in}{1.586114in}}%
\pgfpathlineto{\pgfqpoint{3.697295in}{1.597420in}}%
\pgfpathlineto{\pgfqpoint{3.764409in}{1.597420in}}%
\pgfpathlineto{\pgfqpoint{3.788814in}{1.642642in}}%
\pgfpathlineto{\pgfqpoint{3.794916in}{1.663369in}}%
\pgfpathlineto{\pgfqpoint{3.801017in}{1.697286in}}%
\pgfpathlineto{\pgfqpoint{3.807118in}{1.719897in}}%
\pgfpathlineto{\pgfqpoint{3.813219in}{1.751930in}}%
\pgfpathlineto{\pgfqpoint{3.819321in}{1.797152in}}%
\pgfpathlineto{\pgfqpoint{3.825422in}{1.817879in}}%
\pgfpathlineto{\pgfqpoint{3.843726in}{1.917745in}}%
\pgfpathlineto{\pgfqpoint{3.855928in}{1.962968in}}%
\pgfpathlineto{\pgfqpoint{3.874232in}{1.995000in}}%
\pgfpathlineto{\pgfqpoint{3.880333in}{1.995000in}}%
\pgfpathlineto{\pgfqpoint{3.886434in}{2.006306in}}%
\pgfpathlineto{\pgfqpoint{3.892536in}{2.006306in}}%
\pgfpathlineto{\pgfqpoint{3.898637in}{2.017611in}}%
\pgfpathlineto{\pgfqpoint{3.910839in}{2.017611in}}%
\pgfpathlineto{\pgfqpoint{3.916941in}{2.028917in}}%
\pgfpathlineto{\pgfqpoint{3.965751in}{2.028917in}}%
\pgfpathlineto{\pgfqpoint{4.002358in}{2.094866in}}%
\pgfpathlineto{\pgfqpoint{4.008459in}{2.094866in}}%
\pgfpathlineto{\pgfqpoint{4.014561in}{2.106172in}}%
\pgfpathlineto{\pgfqpoint{4.020662in}{2.106172in}}%
\pgfpathlineto{\pgfqpoint{4.026763in}{2.117478in}}%
\pgfpathlineto{\pgfqpoint{4.032864in}{2.117478in}}%
\pgfpathlineto{\pgfqpoint{4.038966in}{2.128783in}}%
\pgfpathlineto{\pgfqpoint{4.045067in}{2.128783in}}%
\pgfpathlineto{\pgfqpoint{4.057270in}{2.149510in}}%
\pgfpathlineto{\pgfqpoint{4.069472in}{2.149510in}}%
\pgfpathlineto{\pgfqpoint{4.075573in}{2.160816in}}%
\pgfpathlineto{\pgfqpoint{4.081675in}{2.160816in}}%
\pgfpathlineto{\pgfqpoint{4.087776in}{2.172121in}}%
\pgfpathlineto{\pgfqpoint{4.093877in}{2.172121in}}%
\pgfpathlineto{\pgfqpoint{4.167092in}{2.305905in}}%
\pgfpathlineto{\pgfqpoint{4.173193in}{2.305905in}}%
\pgfpathlineto{\pgfqpoint{4.179295in}{2.315326in}}%
\pgfpathlineto{\pgfqpoint{4.185396in}{2.337937in}}%
\pgfpathlineto{\pgfqpoint{4.197598in}{2.337937in}}%
\pgfpathlineto{\pgfqpoint{4.203700in}{2.349243in}}%
\pgfpathlineto{\pgfqpoint{4.209801in}{2.349243in}}%
\pgfpathlineto{\pgfqpoint{4.215902in}{2.360548in}}%
\pgfpathlineto{\pgfqpoint{4.240307in}{2.360548in}}%
\pgfpathlineto{\pgfqpoint{4.246408in}{2.371854in}}%
\pgfpathlineto{\pgfqpoint{4.252510in}{2.371854in}}%
\pgfpathlineto{\pgfqpoint{4.258611in}{2.360548in}}%
\pgfpathlineto{\pgfqpoint{4.264712in}{2.360548in}}%
\pgfpathlineto{\pgfqpoint{4.270813in}{2.349243in}}%
\pgfpathlineto{\pgfqpoint{4.283016in}{2.349243in}}%
\pgfpathlineto{\pgfqpoint{4.295218in}{2.326631in}}%
\pgfpathlineto{\pgfqpoint{4.307421in}{2.326631in}}%
\pgfpathlineto{\pgfqpoint{4.313522in}{2.315326in}}%
\pgfpathlineto{\pgfqpoint{4.319623in}{2.315326in}}%
\pgfpathlineto{\pgfqpoint{4.325725in}{2.305905in}}%
\pgfpathlineto{\pgfqpoint{4.331826in}{2.305905in}}%
\pgfpathlineto{\pgfqpoint{4.344028in}{2.283293in}}%
\pgfpathlineto{\pgfqpoint{4.350130in}{2.283293in}}%
\pgfpathlineto{\pgfqpoint{4.356231in}{2.271988in}}%
\pgfpathlineto{\pgfqpoint{4.362332in}{2.271988in}}%
\pgfpathlineto{\pgfqpoint{4.368434in}{2.283293in}}%
\pgfpathlineto{\pgfqpoint{4.374535in}{2.271988in}}%
\pgfpathlineto{\pgfqpoint{4.411142in}{2.271988in}}%
\pgfpathlineto{\pgfqpoint{4.417244in}{2.283293in}}%
\pgfpathlineto{\pgfqpoint{4.423345in}{2.283293in}}%
\pgfpathlineto{\pgfqpoint{4.429446in}{2.294599in}}%
\pgfpathlineto{\pgfqpoint{4.435547in}{2.294599in}}%
\pgfpathlineto{\pgfqpoint{4.441649in}{2.305905in}}%
\pgfpathlineto{\pgfqpoint{4.453851in}{2.305905in}}%
\pgfpathlineto{\pgfqpoint{4.459952in}{2.315326in}}%
\pgfpathlineto{\pgfqpoint{4.490459in}{2.315326in}}%
\pgfpathlineto{\pgfqpoint{4.496560in}{2.305905in}}%
\pgfpathlineto{\pgfqpoint{4.514864in}{2.305905in}}%
\pgfpathlineto{\pgfqpoint{4.520965in}{2.294599in}}%
\pgfpathlineto{\pgfqpoint{4.527066in}{2.294599in}}%
\pgfpathlineto{\pgfqpoint{4.533167in}{2.283293in}}%
\pgfpathlineto{\pgfqpoint{4.563674in}{2.283293in}}%
\pgfpathlineto{\pgfqpoint{4.569775in}{2.294599in}}%
\pgfpathlineto{\pgfqpoint{4.594180in}{2.294599in}}%
\pgfpathlineto{\pgfqpoint{4.600281in}{2.315326in}}%
\pgfpathlineto{\pgfqpoint{4.606382in}{2.315326in}}%
\pgfpathlineto{\pgfqpoint{4.612484in}{2.326631in}}%
\pgfpathlineto{\pgfqpoint{4.630787in}{2.326631in}}%
\pgfpathlineto{\pgfqpoint{4.642990in}{2.349243in}}%
\pgfpathlineto{\pgfqpoint{4.661294in}{2.349243in}}%
\pgfpathlineto{\pgfqpoint{4.667395in}{2.360548in}}%
\pgfpathlineto{\pgfqpoint{4.685699in}{2.360548in}}%
\pgfpathlineto{\pgfqpoint{4.697901in}{2.383160in}}%
\pgfpathlineto{\pgfqpoint{4.716205in}{2.383160in}}%
\pgfpathlineto{\pgfqpoint{4.722306in}{2.392581in}}%
\pgfpathlineto{\pgfqpoint{4.728408in}{2.392581in}}%
\pgfpathlineto{\pgfqpoint{4.734509in}{2.403886in}}%
\pgfpathlineto{\pgfqpoint{4.740610in}{2.403886in}}%
\pgfpathlineto{\pgfqpoint{4.746711in}{2.415192in}}%
\pgfpathlineto{\pgfqpoint{4.758914in}{2.415192in}}%
\pgfpathlineto{\pgfqpoint{4.765015in}{2.426498in}}%
\pgfpathlineto{\pgfqpoint{4.771116in}{2.426498in}}%
\pgfpathlineto{\pgfqpoint{4.783319in}{2.449109in}}%
\pgfpathlineto{\pgfqpoint{4.789420in}{2.449109in}}%
\pgfpathlineto{\pgfqpoint{4.795521in}{2.460415in}}%
\pgfpathlineto{\pgfqpoint{4.819926in}{2.460415in}}%
\pgfpathlineto{\pgfqpoint{4.826028in}{2.471720in}}%
\pgfpathlineto{\pgfqpoint{4.868736in}{2.471720in}}%
\pgfpathlineto{\pgfqpoint{4.874838in}{2.481141in}}%
\pgfpathlineto{\pgfqpoint{4.887040in}{2.481141in}}%
\pgfpathlineto{\pgfqpoint{4.893141in}{2.471720in}}%
\pgfpathlineto{\pgfqpoint{4.899243in}{2.471720in}}%
\pgfpathlineto{\pgfqpoint{4.905344in}{2.481141in}}%
\pgfpathlineto{\pgfqpoint{4.923648in}{2.481141in}}%
\pgfpathlineto{\pgfqpoint{4.929749in}{2.492447in}}%
\pgfpathlineto{\pgfqpoint{4.954154in}{2.492447in}}%
\pgfpathlineto{\pgfqpoint{4.960255in}{2.503753in}}%
\pgfpathlineto{\pgfqpoint{4.984660in}{2.503753in}}%
\pgfpathlineto{\pgfqpoint{4.996863in}{2.526364in}}%
\pgfpathlineto{\pgfqpoint{5.045673in}{2.526364in}}%
\pgfpathlineto{\pgfqpoint{5.051774in}{2.537670in}}%
\pgfpathlineto{\pgfqpoint{5.076179in}{2.537670in}}%
\pgfpathlineto{\pgfqpoint{5.082280in}{2.526364in}}%
\pgfpathlineto{\pgfqpoint{5.124989in}{2.526364in}}%
\pgfpathlineto{\pgfqpoint{5.131090in}{2.537670in}}%
\pgfpathlineto{\pgfqpoint{5.143293in}{2.537670in}}%
\pgfpathlineto{\pgfqpoint{5.149394in}{2.548975in}}%
\pgfpathlineto{\pgfqpoint{5.173799in}{2.548975in}}%
\pgfpathlineto{\pgfqpoint{5.179900in}{2.558396in}}%
\pgfpathlineto{\pgfqpoint{5.192103in}{2.558396in}}%
\pgfpathlineto{\pgfqpoint{5.198204in}{2.569702in}}%
\pgfpathlineto{\pgfqpoint{5.210407in}{2.569702in}}%
\pgfpathlineto{\pgfqpoint{5.216508in}{2.581008in}}%
\pgfpathlineto{\pgfqpoint{5.259217in}{2.581008in}}%
\pgfpathlineto{\pgfqpoint{5.265318in}{2.592313in}}%
\pgfpathlineto{\pgfqpoint{5.271419in}{2.592313in}}%
\pgfpathlineto{\pgfqpoint{5.277521in}{2.603619in}}%
\pgfpathlineto{\pgfqpoint{5.289723in}{2.603619in}}%
\pgfpathlineto{\pgfqpoint{5.295824in}{2.614925in}}%
\pgfpathlineto{\pgfqpoint{5.308027in}{2.614925in}}%
\pgfpathlineto{\pgfqpoint{5.314128in}{2.626230in}}%
\pgfpathlineto{\pgfqpoint{5.320229in}{2.626230in}}%
\pgfpathlineto{\pgfqpoint{5.326331in}{2.637536in}}%
\pgfpathlineto{\pgfqpoint{5.332432in}{2.637536in}}%
\pgfpathlineto{\pgfqpoint{5.338533in}{2.646957in}}%
\pgfpathlineto{\pgfqpoint{5.350736in}{2.646957in}}%
\pgfpathlineto{\pgfqpoint{5.356837in}{2.669568in}}%
\pgfpathlineto{\pgfqpoint{5.375141in}{2.669568in}}%
\pgfpathlineto{\pgfqpoint{5.381242in}{2.680874in}}%
\pgfpathlineto{\pgfqpoint{5.430052in}{2.680874in}}%
\pgfpathlineto{\pgfqpoint{5.436153in}{2.692180in}}%
\pgfpathlineto{\pgfqpoint{5.448356in}{2.692180in}}%
\pgfpathlineto{\pgfqpoint{5.454457in}{2.703485in}}%
\pgfpathlineto{\pgfqpoint{5.472761in}{2.703485in}}%
\pgfpathlineto{\pgfqpoint{5.478862in}{2.714791in}}%
\pgfpathlineto{\pgfqpoint{5.503267in}{2.714791in}}%
\pgfpathlineto{\pgfqpoint{5.509368in}{2.724212in}}%
\pgfpathlineto{\pgfqpoint{5.539874in}{2.724212in}}%
\pgfpathlineto{\pgfqpoint{5.545976in}{2.746823in}}%
\pgfpathlineto{\pgfqpoint{5.637495in}{2.746823in}}%
\pgfpathlineto{\pgfqpoint{5.643596in}{2.758129in}}%
\pgfpathlineto{\pgfqpoint{5.716811in}{2.758129in}}%
\pgfpathlineto{\pgfqpoint{5.722912in}{2.769435in}}%
\pgfpathlineto{\pgfqpoint{5.759520in}{2.769435in}}%
\pgfpathlineto{\pgfqpoint{5.765621in}{2.780740in}}%
\pgfpathlineto{\pgfqpoint{5.777823in}{2.780740in}}%
\pgfpathlineto{\pgfqpoint{5.783925in}{2.792046in}}%
\pgfpathlineto{\pgfqpoint{5.960861in}{2.792046in}}%
\pgfpathlineto{\pgfqpoint{5.966962in}{2.801467in}}%
\pgfpathlineto{\pgfqpoint{6.027975in}{2.801467in}}%
\pgfpathlineto{\pgfqpoint{6.034076in}{2.792046in}}%
\pgfpathlineto{\pgfqpoint{6.150000in}{2.792046in}}%
\pgfpathlineto{\pgfqpoint{6.150000in}{2.792046in}}%
\pgfusepath{stroke}%
\end{pgfscope}%
\begin{pgfscope}%
\pgfsetrectcap%
\pgfsetmiterjoin%
\pgfsetlinewidth{0.803000pt}%
\definecolor{currentstroke}{rgb}{0.000000,0.000000,0.000000}%
\pgfsetstrokecolor{currentstroke}%
\pgfsetdash{}{0pt}%
\pgfpathmoveto{\pgfqpoint{0.494137in}{0.565123in}}%
\pgfpathlineto{\pgfqpoint{0.494137in}{3.336317in}}%
\pgfusepath{stroke}%
\end{pgfscope}%
\begin{pgfscope}%
\pgfsetrectcap%
\pgfsetmiterjoin%
\pgfsetlinewidth{0.803000pt}%
\definecolor{currentstroke}{rgb}{0.000000,0.000000,0.000000}%
\pgfsetstrokecolor{currentstroke}%
\pgfsetdash{}{0pt}%
\pgfpathmoveto{\pgfqpoint{6.150000in}{0.565123in}}%
\pgfpathlineto{\pgfqpoint{6.150000in}{3.336317in}}%
\pgfusepath{stroke}%
\end{pgfscope}%
\begin{pgfscope}%
\pgfsetrectcap%
\pgfsetmiterjoin%
\pgfsetlinewidth{0.803000pt}%
\definecolor{currentstroke}{rgb}{0.000000,0.000000,0.000000}%
\pgfsetstrokecolor{currentstroke}%
\pgfsetdash{}{0pt}%
\pgfpathmoveto{\pgfqpoint{0.494137in}{0.565123in}}%
\pgfpathlineto{\pgfqpoint{6.150000in}{0.565123in}}%
\pgfusepath{stroke}%
\end{pgfscope}%
\begin{pgfscope}%
\pgfsetrectcap%
\pgfsetmiterjoin%
\pgfsetlinewidth{0.803000pt}%
\definecolor{currentstroke}{rgb}{0.000000,0.000000,0.000000}%
\pgfsetstrokecolor{currentstroke}%
\pgfsetdash{}{0pt}%
\pgfpathmoveto{\pgfqpoint{0.494137in}{3.336317in}}%
\pgfpathlineto{\pgfqpoint{6.150000in}{3.336317in}}%
\pgfusepath{stroke}%
\end{pgfscope}%
\begin{pgfscope}%
\pgfsetbuttcap%
\pgfsetmiterjoin%
\definecolor{currentfill}{rgb}{1.000000,1.000000,1.000000}%
\pgfsetfillcolor{currentfill}%
\pgfsetfillopacity{0.800000}%
\pgfsetlinewidth{1.003750pt}%
\definecolor{currentstroke}{rgb}{0.800000,0.800000,0.800000}%
\pgfsetstrokecolor{currentstroke}%
\pgfsetstrokeopacity{0.800000}%
\pgfsetdash{}{0pt}%
\pgfpathmoveto{\pgfqpoint{4.539659in}{3.016872in}}%
\pgfpathlineto{\pgfqpoint{6.052778in}{3.016872in}}%
\pgfpathquadraticcurveto{\pgfqpoint{6.080556in}{3.016872in}}{\pgfqpoint{6.080556in}{3.044650in}}%
\pgfpathlineto{\pgfqpoint{6.080556in}{3.239094in}}%
\pgfpathquadraticcurveto{\pgfqpoint{6.080556in}{3.266872in}}{\pgfqpoint{6.052778in}{3.266872in}}%
\pgfpathlineto{\pgfqpoint{4.539659in}{3.266872in}}%
\pgfpathquadraticcurveto{\pgfqpoint{4.511881in}{3.266872in}}{\pgfqpoint{4.511881in}{3.239094in}}%
\pgfpathlineto{\pgfqpoint{4.511881in}{3.044650in}}%
\pgfpathquadraticcurveto{\pgfqpoint{4.511881in}{3.016872in}}{\pgfqpoint{4.539659in}{3.016872in}}%
\pgfpathlineto{\pgfqpoint{4.539659in}{3.016872in}}%
\pgfpathclose%
\pgfusepath{stroke,fill}%
\end{pgfscope}%
\begin{pgfscope}%
\pgfsetrectcap%
\pgfsetroundjoin%
\pgfsetlinewidth{1.003750pt}%
\definecolor{currentstroke}{rgb}{0.000000,0.000000,0.000000}%
\pgfsetstrokecolor{currentstroke}%
\pgfsetdash{}{0pt}%
\pgfpathmoveto{\pgfqpoint{4.567436in}{3.155761in}}%
\pgfpathlineto{\pgfqpoint{4.706325in}{3.155761in}}%
\pgfpathlineto{\pgfqpoint{4.845214in}{3.155761in}}%
\pgfusepath{stroke}%
\end{pgfscope}%
\begin{pgfscope}%
\definecolor{textcolor}{rgb}{0.000000,0.000000,0.000000}%
\pgfsetstrokecolor{textcolor}%
\pgfsetfillcolor{textcolor}%
\pgftext[x=4.956325in,y=3.107150in,left,base]{\color{textcolor}\rmfamily\fontsize{10.000000}{12.000000}\selectfont Return Loss (dB)}%
\end{pgfscope}%
\end{pgfpicture}%
\makeatother%
\endgroup%

		\caption{Measured return loss and phase of a quartz resonatoras a function of frequency.}
		\label{fig:358}
	\end{figure}
	Figure \ref{fig:358} shows a minimum in the reflection coefficient at approximately \(91.9\,\mathrm{MHz}\), where the magnitude is measured as \(-11.85\,\mathrm{dB}\). Substituting this value into the equation below gives the standing-wave ratio (SWR):
	
	\[
	\mathrm{SWR} = \frac{1 + 10^{\left(\frac{|S_{11}|_{\mathrm{dB}}}{20}\right)}}{1 - 10^{\left(\frac{|S_{11}|_{\mathrm{dB}}}{20}\right)}} = 1.69
	\]
	
	The impedance can be determined just like shown in chapter \ref{ch:355}, the reflection coefficient phase at resonance was measured at \(152.5^\circ\). Thus the calculated value results in $Z_L=30.78 + j7.75 \Omega$.
	The resonance frequency lies within the FM broadcast band (\(87\)--\(108\,\mathrm{MHz}\)), the antenna is suitable for use in this frequency range.

	
	
	
	\subsubsection{}
	For a monopole antenna, the approximate length is given by the quarter-wavelength rule $l = \frac{\lambda}{4}$.
	Substituting the free-space wavelength relation \(\lambda = c_0 / f\) yields:
	\[
	l = \frac{c_0}{4 f} = \frac{c_0}{4 \cdot 91.601\,\mathrm{MHz}} \approx 0.816\,\mathrm{m} = 0.818\,\mathrm{m}
	\]
	
	\subsubsection{}
	
	\begin{figure}[H]
		\centering
		\input{3510.pgf}
		\caption{\url{https://tenor.com/de/view/will-smith-gif-3197096261065537476}}
		\label{fig:3510}
	\end{figure}
	
	\subsubsection{}
	
	All three attenuators can be used for the entire frequency range of 8 GHz as can be seen in fig. \ref{fig:3510}. The given cutoff frequency was given on the elements much higher at 18 GHz.
	
	\subsubsection{}
	Since the ground creates a nearly perfect reflection due to having a impedance of zero, the signal passes through the attenuator two times resulting in a doubled damping thus a doubled return loss: 12 dB, 20 dB, 40 dB.
	\subsubsection{}
	In figure \ref{fig:3513} the center frequency of the LC band-pass can be determined at 8.502 MHz with a value of -0.18 dB. The return loss is measured at $\sim 37$ dB. The 3 dB threshold in transmission is measured at 7495021 Hz.
	\begin{figure}[H]
		\centering
		%% Creator: Matplotlib, PGF backend
%%
%% To include the figure in your LaTeX document, write
%%   \input{<filename>.pgf}
%%
%% Make sure the required packages are loaded in your preamble
%%   \usepackage{pgf}
%%
%% Also ensure that all the required font packages are loaded; for instance,
%% the lmodern package is sometimes necessary when using math font.
%%   \usepackage{lmodern}
%%
%% Figures using additional raster images can only be included by \input if
%% they are in the same directory as the main LaTeX file. For loading figures
%% from other directories you can use the `import` package
%%   \usepackage{import}
%%
%% and then include the figures with
%%   \import{<path to file>}{<filename>.pgf}
%%
%% Matplotlib used the following preamble
%%
\begingroup%
\makeatletter%
\begin{pgfpicture}%
\pgfpathrectangle{\pgfpointorigin}{\pgfqpoint{6.300000in}{3.500000in}}%
\pgfusepath{use as bounding box, clip}%
\begin{pgfscope}%
\pgfsetbuttcap%
\pgfsetmiterjoin%
\definecolor{currentfill}{rgb}{1.000000,1.000000,1.000000}%
\pgfsetfillcolor{currentfill}%
\pgfsetlinewidth{0.000000pt}%
\definecolor{currentstroke}{rgb}{1.000000,1.000000,1.000000}%
\pgfsetstrokecolor{currentstroke}%
\pgfsetdash{}{0pt}%
\pgfpathmoveto{\pgfqpoint{0.000000in}{0.000000in}}%
\pgfpathlineto{\pgfqpoint{6.300000in}{0.000000in}}%
\pgfpathlineto{\pgfqpoint{6.300000in}{3.500000in}}%
\pgfpathlineto{\pgfqpoint{0.000000in}{3.500000in}}%
\pgfpathlineto{\pgfqpoint{0.000000in}{0.000000in}}%
\pgfpathclose%
\pgfusepath{fill}%
\end{pgfscope}%
\begin{pgfscope}%
\pgfsetbuttcap%
\pgfsetmiterjoin%
\definecolor{currentfill}{rgb}{1.000000,1.000000,1.000000}%
\pgfsetfillcolor{currentfill}%
\pgfsetlinewidth{0.000000pt}%
\definecolor{currentstroke}{rgb}{0.000000,0.000000,0.000000}%
\pgfsetstrokecolor{currentstroke}%
\pgfsetstrokeopacity{0.000000}%
\pgfsetdash{}{0pt}%
\pgfpathmoveto{\pgfqpoint{0.688581in}{0.565123in}}%
\pgfpathlineto{\pgfqpoint{6.150000in}{0.565123in}}%
\pgfpathlineto{\pgfqpoint{6.150000in}{3.150926in}}%
\pgfpathlineto{\pgfqpoint{0.688581in}{3.150926in}}%
\pgfpathlineto{\pgfqpoint{0.688581in}{0.565123in}}%
\pgfpathclose%
\pgfusepath{fill}%
\end{pgfscope}%
\begin{pgfscope}%
\pgfpathrectangle{\pgfqpoint{0.688581in}{0.565123in}}{\pgfqpoint{5.461419in}{2.585803in}}%
\pgfusepath{clip}%
\pgfsetbuttcap%
\pgfsetroundjoin%
\pgfsetlinewidth{0.803000pt}%
\definecolor{currentstroke}{rgb}{0.690196,0.690196,0.690196}%
\pgfsetstrokecolor{currentstroke}%
\pgfsetdash{{0.800000pt}{1.320000pt}}{0.000000pt}%
\pgfpathmoveto{\pgfqpoint{0.852584in}{0.565123in}}%
\pgfpathlineto{\pgfqpoint{0.852584in}{3.150926in}}%
\pgfusepath{stroke}%
\end{pgfscope}%
\begin{pgfscope}%
\pgfsetbuttcap%
\pgfsetroundjoin%
\definecolor{currentfill}{rgb}{0.000000,0.000000,0.000000}%
\pgfsetfillcolor{currentfill}%
\pgfsetlinewidth{0.803000pt}%
\definecolor{currentstroke}{rgb}{0.000000,0.000000,0.000000}%
\pgfsetstrokecolor{currentstroke}%
\pgfsetdash{}{0pt}%
\pgfsys@defobject{currentmarker}{\pgfqpoint{0.000000in}{-0.048611in}}{\pgfqpoint{0.000000in}{0.000000in}}{%
\pgfpathmoveto{\pgfqpoint{0.000000in}{0.000000in}}%
\pgfpathlineto{\pgfqpoint{0.000000in}{-0.048611in}}%
\pgfusepath{stroke,fill}%
}%
\begin{pgfscope}%
\pgfsys@transformshift{0.852584in}{0.565123in}%
\pgfsys@useobject{currentmarker}{}%
\end{pgfscope}%
\end{pgfscope}%
\begin{pgfscope}%
\definecolor{textcolor}{rgb}{0.000000,0.000000,0.000000}%
\pgfsetstrokecolor{textcolor}%
\pgfsetfillcolor{textcolor}%
\pgftext[x=0.852584in,y=0.467901in,,top]{\color{textcolor}\rmfamily\fontsize{10.000000}{12.000000}\selectfont \(\displaystyle {0}\)}%
\end{pgfscope}%
\begin{pgfscope}%
\pgfpathrectangle{\pgfqpoint{0.688581in}{0.565123in}}{\pgfqpoint{5.461419in}{2.585803in}}%
\pgfusepath{clip}%
\pgfsetbuttcap%
\pgfsetroundjoin%
\pgfsetlinewidth{0.803000pt}%
\definecolor{currentstroke}{rgb}{0.690196,0.690196,0.690196}%
\pgfsetstrokecolor{currentstroke}%
\pgfsetdash{{0.800000pt}{1.320000pt}}{0.000000pt}%
\pgfpathmoveto{\pgfqpoint{1.695013in}{0.565123in}}%
\pgfpathlineto{\pgfqpoint{1.695013in}{3.150926in}}%
\pgfusepath{stroke}%
\end{pgfscope}%
\begin{pgfscope}%
\pgfsetbuttcap%
\pgfsetroundjoin%
\definecolor{currentfill}{rgb}{0.000000,0.000000,0.000000}%
\pgfsetfillcolor{currentfill}%
\pgfsetlinewidth{0.803000pt}%
\definecolor{currentstroke}{rgb}{0.000000,0.000000,0.000000}%
\pgfsetstrokecolor{currentstroke}%
\pgfsetdash{}{0pt}%
\pgfsys@defobject{currentmarker}{\pgfqpoint{0.000000in}{-0.048611in}}{\pgfqpoint{0.000000in}{0.000000in}}{%
\pgfpathmoveto{\pgfqpoint{0.000000in}{0.000000in}}%
\pgfpathlineto{\pgfqpoint{0.000000in}{-0.048611in}}%
\pgfusepath{stroke,fill}%
}%
\begin{pgfscope}%
\pgfsys@transformshift{1.695013in}{0.565123in}%
\pgfsys@useobject{currentmarker}{}%
\end{pgfscope}%
\end{pgfscope}%
\begin{pgfscope}%
\definecolor{textcolor}{rgb}{0.000000,0.000000,0.000000}%
\pgfsetstrokecolor{textcolor}%
\pgfsetfillcolor{textcolor}%
\pgftext[x=1.695013in,y=0.467901in,,top]{\color{textcolor}\rmfamily\fontsize{10.000000}{12.000000}\selectfont \(\displaystyle {10}\)}%
\end{pgfscope}%
\begin{pgfscope}%
\pgfpathrectangle{\pgfqpoint{0.688581in}{0.565123in}}{\pgfqpoint{5.461419in}{2.585803in}}%
\pgfusepath{clip}%
\pgfsetbuttcap%
\pgfsetroundjoin%
\pgfsetlinewidth{0.803000pt}%
\definecolor{currentstroke}{rgb}{0.690196,0.690196,0.690196}%
\pgfsetstrokecolor{currentstroke}%
\pgfsetdash{{0.800000pt}{1.320000pt}}{0.000000pt}%
\pgfpathmoveto{\pgfqpoint{2.537441in}{0.565123in}}%
\pgfpathlineto{\pgfqpoint{2.537441in}{3.150926in}}%
\pgfusepath{stroke}%
\end{pgfscope}%
\begin{pgfscope}%
\pgfsetbuttcap%
\pgfsetroundjoin%
\definecolor{currentfill}{rgb}{0.000000,0.000000,0.000000}%
\pgfsetfillcolor{currentfill}%
\pgfsetlinewidth{0.803000pt}%
\definecolor{currentstroke}{rgb}{0.000000,0.000000,0.000000}%
\pgfsetstrokecolor{currentstroke}%
\pgfsetdash{}{0pt}%
\pgfsys@defobject{currentmarker}{\pgfqpoint{0.000000in}{-0.048611in}}{\pgfqpoint{0.000000in}{0.000000in}}{%
\pgfpathmoveto{\pgfqpoint{0.000000in}{0.000000in}}%
\pgfpathlineto{\pgfqpoint{0.000000in}{-0.048611in}}%
\pgfusepath{stroke,fill}%
}%
\begin{pgfscope}%
\pgfsys@transformshift{2.537441in}{0.565123in}%
\pgfsys@useobject{currentmarker}{}%
\end{pgfscope}%
\end{pgfscope}%
\begin{pgfscope}%
\definecolor{textcolor}{rgb}{0.000000,0.000000,0.000000}%
\pgfsetstrokecolor{textcolor}%
\pgfsetfillcolor{textcolor}%
\pgftext[x=2.537441in,y=0.467901in,,top]{\color{textcolor}\rmfamily\fontsize{10.000000}{12.000000}\selectfont \(\displaystyle {20}\)}%
\end{pgfscope}%
\begin{pgfscope}%
\pgfpathrectangle{\pgfqpoint{0.688581in}{0.565123in}}{\pgfqpoint{5.461419in}{2.585803in}}%
\pgfusepath{clip}%
\pgfsetbuttcap%
\pgfsetroundjoin%
\pgfsetlinewidth{0.803000pt}%
\definecolor{currentstroke}{rgb}{0.690196,0.690196,0.690196}%
\pgfsetstrokecolor{currentstroke}%
\pgfsetdash{{0.800000pt}{1.320000pt}}{0.000000pt}%
\pgfpathmoveto{\pgfqpoint{3.379870in}{0.565123in}}%
\pgfpathlineto{\pgfqpoint{3.379870in}{3.150926in}}%
\pgfusepath{stroke}%
\end{pgfscope}%
\begin{pgfscope}%
\pgfsetbuttcap%
\pgfsetroundjoin%
\definecolor{currentfill}{rgb}{0.000000,0.000000,0.000000}%
\pgfsetfillcolor{currentfill}%
\pgfsetlinewidth{0.803000pt}%
\definecolor{currentstroke}{rgb}{0.000000,0.000000,0.000000}%
\pgfsetstrokecolor{currentstroke}%
\pgfsetdash{}{0pt}%
\pgfsys@defobject{currentmarker}{\pgfqpoint{0.000000in}{-0.048611in}}{\pgfqpoint{0.000000in}{0.000000in}}{%
\pgfpathmoveto{\pgfqpoint{0.000000in}{0.000000in}}%
\pgfpathlineto{\pgfqpoint{0.000000in}{-0.048611in}}%
\pgfusepath{stroke,fill}%
}%
\begin{pgfscope}%
\pgfsys@transformshift{3.379870in}{0.565123in}%
\pgfsys@useobject{currentmarker}{}%
\end{pgfscope}%
\end{pgfscope}%
\begin{pgfscope}%
\definecolor{textcolor}{rgb}{0.000000,0.000000,0.000000}%
\pgfsetstrokecolor{textcolor}%
\pgfsetfillcolor{textcolor}%
\pgftext[x=3.379870in,y=0.467901in,,top]{\color{textcolor}\rmfamily\fontsize{10.000000}{12.000000}\selectfont \(\displaystyle {30}\)}%
\end{pgfscope}%
\begin{pgfscope}%
\pgfpathrectangle{\pgfqpoint{0.688581in}{0.565123in}}{\pgfqpoint{5.461419in}{2.585803in}}%
\pgfusepath{clip}%
\pgfsetbuttcap%
\pgfsetroundjoin%
\pgfsetlinewidth{0.803000pt}%
\definecolor{currentstroke}{rgb}{0.690196,0.690196,0.690196}%
\pgfsetstrokecolor{currentstroke}%
\pgfsetdash{{0.800000pt}{1.320000pt}}{0.000000pt}%
\pgfpathmoveto{\pgfqpoint{4.222298in}{0.565123in}}%
\pgfpathlineto{\pgfqpoint{4.222298in}{3.150926in}}%
\pgfusepath{stroke}%
\end{pgfscope}%
\begin{pgfscope}%
\pgfsetbuttcap%
\pgfsetroundjoin%
\definecolor{currentfill}{rgb}{0.000000,0.000000,0.000000}%
\pgfsetfillcolor{currentfill}%
\pgfsetlinewidth{0.803000pt}%
\definecolor{currentstroke}{rgb}{0.000000,0.000000,0.000000}%
\pgfsetstrokecolor{currentstroke}%
\pgfsetdash{}{0pt}%
\pgfsys@defobject{currentmarker}{\pgfqpoint{0.000000in}{-0.048611in}}{\pgfqpoint{0.000000in}{0.000000in}}{%
\pgfpathmoveto{\pgfqpoint{0.000000in}{0.000000in}}%
\pgfpathlineto{\pgfqpoint{0.000000in}{-0.048611in}}%
\pgfusepath{stroke,fill}%
}%
\begin{pgfscope}%
\pgfsys@transformshift{4.222298in}{0.565123in}%
\pgfsys@useobject{currentmarker}{}%
\end{pgfscope}%
\end{pgfscope}%
\begin{pgfscope}%
\definecolor{textcolor}{rgb}{0.000000,0.000000,0.000000}%
\pgfsetstrokecolor{textcolor}%
\pgfsetfillcolor{textcolor}%
\pgftext[x=4.222298in,y=0.467901in,,top]{\color{textcolor}\rmfamily\fontsize{10.000000}{12.000000}\selectfont \(\displaystyle {40}\)}%
\end{pgfscope}%
\begin{pgfscope}%
\pgfpathrectangle{\pgfqpoint{0.688581in}{0.565123in}}{\pgfqpoint{5.461419in}{2.585803in}}%
\pgfusepath{clip}%
\pgfsetbuttcap%
\pgfsetroundjoin%
\pgfsetlinewidth{0.803000pt}%
\definecolor{currentstroke}{rgb}{0.690196,0.690196,0.690196}%
\pgfsetstrokecolor{currentstroke}%
\pgfsetdash{{0.800000pt}{1.320000pt}}{0.000000pt}%
\pgfpathmoveto{\pgfqpoint{5.064727in}{0.565123in}}%
\pgfpathlineto{\pgfqpoint{5.064727in}{3.150926in}}%
\pgfusepath{stroke}%
\end{pgfscope}%
\begin{pgfscope}%
\pgfsetbuttcap%
\pgfsetroundjoin%
\definecolor{currentfill}{rgb}{0.000000,0.000000,0.000000}%
\pgfsetfillcolor{currentfill}%
\pgfsetlinewidth{0.803000pt}%
\definecolor{currentstroke}{rgb}{0.000000,0.000000,0.000000}%
\pgfsetstrokecolor{currentstroke}%
\pgfsetdash{}{0pt}%
\pgfsys@defobject{currentmarker}{\pgfqpoint{0.000000in}{-0.048611in}}{\pgfqpoint{0.000000in}{0.000000in}}{%
\pgfpathmoveto{\pgfqpoint{0.000000in}{0.000000in}}%
\pgfpathlineto{\pgfqpoint{0.000000in}{-0.048611in}}%
\pgfusepath{stroke,fill}%
}%
\begin{pgfscope}%
\pgfsys@transformshift{5.064727in}{0.565123in}%
\pgfsys@useobject{currentmarker}{}%
\end{pgfscope}%
\end{pgfscope}%
\begin{pgfscope}%
\definecolor{textcolor}{rgb}{0.000000,0.000000,0.000000}%
\pgfsetstrokecolor{textcolor}%
\pgfsetfillcolor{textcolor}%
\pgftext[x=5.064727in,y=0.467901in,,top]{\color{textcolor}\rmfamily\fontsize{10.000000}{12.000000}\selectfont \(\displaystyle {50}\)}%
\end{pgfscope}%
\begin{pgfscope}%
\pgfpathrectangle{\pgfqpoint{0.688581in}{0.565123in}}{\pgfqpoint{5.461419in}{2.585803in}}%
\pgfusepath{clip}%
\pgfsetbuttcap%
\pgfsetroundjoin%
\pgfsetlinewidth{0.803000pt}%
\definecolor{currentstroke}{rgb}{0.690196,0.690196,0.690196}%
\pgfsetstrokecolor{currentstroke}%
\pgfsetdash{{0.800000pt}{1.320000pt}}{0.000000pt}%
\pgfpathmoveto{\pgfqpoint{5.907155in}{0.565123in}}%
\pgfpathlineto{\pgfqpoint{5.907155in}{3.150926in}}%
\pgfusepath{stroke}%
\end{pgfscope}%
\begin{pgfscope}%
\pgfsetbuttcap%
\pgfsetroundjoin%
\definecolor{currentfill}{rgb}{0.000000,0.000000,0.000000}%
\pgfsetfillcolor{currentfill}%
\pgfsetlinewidth{0.803000pt}%
\definecolor{currentstroke}{rgb}{0.000000,0.000000,0.000000}%
\pgfsetstrokecolor{currentstroke}%
\pgfsetdash{}{0pt}%
\pgfsys@defobject{currentmarker}{\pgfqpoint{0.000000in}{-0.048611in}}{\pgfqpoint{0.000000in}{0.000000in}}{%
\pgfpathmoveto{\pgfqpoint{0.000000in}{0.000000in}}%
\pgfpathlineto{\pgfqpoint{0.000000in}{-0.048611in}}%
\pgfusepath{stroke,fill}%
}%
\begin{pgfscope}%
\pgfsys@transformshift{5.907155in}{0.565123in}%
\pgfsys@useobject{currentmarker}{}%
\end{pgfscope}%
\end{pgfscope}%
\begin{pgfscope}%
\definecolor{textcolor}{rgb}{0.000000,0.000000,0.000000}%
\pgfsetstrokecolor{textcolor}%
\pgfsetfillcolor{textcolor}%
\pgftext[x=5.907155in,y=0.467901in,,top]{\color{textcolor}\rmfamily\fontsize{10.000000}{12.000000}\selectfont \(\displaystyle {60}\)}%
\end{pgfscope}%
\begin{pgfscope}%
\definecolor{textcolor}{rgb}{0.000000,0.000000,0.000000}%
\pgfsetstrokecolor{textcolor}%
\pgfsetfillcolor{textcolor}%
\pgftext[x=3.419290in,y=0.288889in,,top]{\color{textcolor}\rmfamily\fontsize{10.000000}{12.000000}\selectfont Frequency (MHz)}%
\end{pgfscope}%
\begin{pgfscope}%
\pgfpathrectangle{\pgfqpoint{0.688581in}{0.565123in}}{\pgfqpoint{5.461419in}{2.585803in}}%
\pgfusepath{clip}%
\pgfsetbuttcap%
\pgfsetroundjoin%
\pgfsetlinewidth{0.803000pt}%
\definecolor{currentstroke}{rgb}{0.690196,0.690196,0.690196}%
\pgfsetstrokecolor{currentstroke}%
\pgfsetdash{{0.800000pt}{1.320000pt}}{0.000000pt}%
\pgfpathmoveto{\pgfqpoint{0.688581in}{0.793431in}}%
\pgfpathlineto{\pgfqpoint{6.150000in}{0.793431in}}%
\pgfusepath{stroke}%
\end{pgfscope}%
\begin{pgfscope}%
\pgfsetbuttcap%
\pgfsetroundjoin%
\definecolor{currentfill}{rgb}{0.000000,0.000000,0.000000}%
\pgfsetfillcolor{currentfill}%
\pgfsetlinewidth{0.803000pt}%
\definecolor{currentstroke}{rgb}{0.000000,0.000000,0.000000}%
\pgfsetstrokecolor{currentstroke}%
\pgfsetdash{}{0pt}%
\pgfsys@defobject{currentmarker}{\pgfqpoint{-0.048611in}{0.000000in}}{\pgfqpoint{-0.000000in}{0.000000in}}{%
\pgfpathmoveto{\pgfqpoint{-0.000000in}{0.000000in}}%
\pgfpathlineto{\pgfqpoint{-0.048611in}{0.000000in}}%
\pgfusepath{stroke,fill}%
}%
\begin{pgfscope}%
\pgfsys@transformshift{0.688581in}{0.793431in}%
\pgfsys@useobject{currentmarker}{}%
\end{pgfscope}%
\end{pgfscope}%
\begin{pgfscope}%
\definecolor{textcolor}{rgb}{0.000000,0.000000,0.000000}%
\pgfsetstrokecolor{textcolor}%
\pgfsetfillcolor{textcolor}%
\pgftext[x=0.344444in, y=0.745205in, left, base]{\color{textcolor}\rmfamily\fontsize{10.000000}{12.000000}\selectfont \(\displaystyle {\ensuremath{-}35}\)}%
\end{pgfscope}%
\begin{pgfscope}%
\pgfpathrectangle{\pgfqpoint{0.688581in}{0.565123in}}{\pgfqpoint{5.461419in}{2.585803in}}%
\pgfusepath{clip}%
\pgfsetbuttcap%
\pgfsetroundjoin%
\pgfsetlinewidth{0.803000pt}%
\definecolor{currentstroke}{rgb}{0.690196,0.690196,0.690196}%
\pgfsetstrokecolor{currentstroke}%
\pgfsetdash{{0.800000pt}{1.320000pt}}{0.000000pt}%
\pgfpathmoveto{\pgfqpoint{0.688581in}{1.108120in}}%
\pgfpathlineto{\pgfqpoint{6.150000in}{1.108120in}}%
\pgfusepath{stroke}%
\end{pgfscope}%
\begin{pgfscope}%
\pgfsetbuttcap%
\pgfsetroundjoin%
\definecolor{currentfill}{rgb}{0.000000,0.000000,0.000000}%
\pgfsetfillcolor{currentfill}%
\pgfsetlinewidth{0.803000pt}%
\definecolor{currentstroke}{rgb}{0.000000,0.000000,0.000000}%
\pgfsetstrokecolor{currentstroke}%
\pgfsetdash{}{0pt}%
\pgfsys@defobject{currentmarker}{\pgfqpoint{-0.048611in}{0.000000in}}{\pgfqpoint{-0.000000in}{0.000000in}}{%
\pgfpathmoveto{\pgfqpoint{-0.000000in}{0.000000in}}%
\pgfpathlineto{\pgfqpoint{-0.048611in}{0.000000in}}%
\pgfusepath{stroke,fill}%
}%
\begin{pgfscope}%
\pgfsys@transformshift{0.688581in}{1.108120in}%
\pgfsys@useobject{currentmarker}{}%
\end{pgfscope}%
\end{pgfscope}%
\begin{pgfscope}%
\definecolor{textcolor}{rgb}{0.000000,0.000000,0.000000}%
\pgfsetstrokecolor{textcolor}%
\pgfsetfillcolor{textcolor}%
\pgftext[x=0.344444in, y=1.059895in, left, base]{\color{textcolor}\rmfamily\fontsize{10.000000}{12.000000}\selectfont \(\displaystyle {\ensuremath{-}30}\)}%
\end{pgfscope}%
\begin{pgfscope}%
\pgfpathrectangle{\pgfqpoint{0.688581in}{0.565123in}}{\pgfqpoint{5.461419in}{2.585803in}}%
\pgfusepath{clip}%
\pgfsetbuttcap%
\pgfsetroundjoin%
\pgfsetlinewidth{0.803000pt}%
\definecolor{currentstroke}{rgb}{0.690196,0.690196,0.690196}%
\pgfsetstrokecolor{currentstroke}%
\pgfsetdash{{0.800000pt}{1.320000pt}}{0.000000pt}%
\pgfpathmoveto{\pgfqpoint{0.688581in}{1.422809in}}%
\pgfpathlineto{\pgfqpoint{6.150000in}{1.422809in}}%
\pgfusepath{stroke}%
\end{pgfscope}%
\begin{pgfscope}%
\pgfsetbuttcap%
\pgfsetroundjoin%
\definecolor{currentfill}{rgb}{0.000000,0.000000,0.000000}%
\pgfsetfillcolor{currentfill}%
\pgfsetlinewidth{0.803000pt}%
\definecolor{currentstroke}{rgb}{0.000000,0.000000,0.000000}%
\pgfsetstrokecolor{currentstroke}%
\pgfsetdash{}{0pt}%
\pgfsys@defobject{currentmarker}{\pgfqpoint{-0.048611in}{0.000000in}}{\pgfqpoint{-0.000000in}{0.000000in}}{%
\pgfpathmoveto{\pgfqpoint{-0.000000in}{0.000000in}}%
\pgfpathlineto{\pgfqpoint{-0.048611in}{0.000000in}}%
\pgfusepath{stroke,fill}%
}%
\begin{pgfscope}%
\pgfsys@transformshift{0.688581in}{1.422809in}%
\pgfsys@useobject{currentmarker}{}%
\end{pgfscope}%
\end{pgfscope}%
\begin{pgfscope}%
\definecolor{textcolor}{rgb}{0.000000,0.000000,0.000000}%
\pgfsetstrokecolor{textcolor}%
\pgfsetfillcolor{textcolor}%
\pgftext[x=0.344444in, y=1.374584in, left, base]{\color{textcolor}\rmfamily\fontsize{10.000000}{12.000000}\selectfont \(\displaystyle {\ensuremath{-}25}\)}%
\end{pgfscope}%
\begin{pgfscope}%
\pgfpathrectangle{\pgfqpoint{0.688581in}{0.565123in}}{\pgfqpoint{5.461419in}{2.585803in}}%
\pgfusepath{clip}%
\pgfsetbuttcap%
\pgfsetroundjoin%
\pgfsetlinewidth{0.803000pt}%
\definecolor{currentstroke}{rgb}{0.690196,0.690196,0.690196}%
\pgfsetstrokecolor{currentstroke}%
\pgfsetdash{{0.800000pt}{1.320000pt}}{0.000000pt}%
\pgfpathmoveto{\pgfqpoint{0.688581in}{1.737499in}}%
\pgfpathlineto{\pgfqpoint{6.150000in}{1.737499in}}%
\pgfusepath{stroke}%
\end{pgfscope}%
\begin{pgfscope}%
\pgfsetbuttcap%
\pgfsetroundjoin%
\definecolor{currentfill}{rgb}{0.000000,0.000000,0.000000}%
\pgfsetfillcolor{currentfill}%
\pgfsetlinewidth{0.803000pt}%
\definecolor{currentstroke}{rgb}{0.000000,0.000000,0.000000}%
\pgfsetstrokecolor{currentstroke}%
\pgfsetdash{}{0pt}%
\pgfsys@defobject{currentmarker}{\pgfqpoint{-0.048611in}{0.000000in}}{\pgfqpoint{-0.000000in}{0.000000in}}{%
\pgfpathmoveto{\pgfqpoint{-0.000000in}{0.000000in}}%
\pgfpathlineto{\pgfqpoint{-0.048611in}{0.000000in}}%
\pgfusepath{stroke,fill}%
}%
\begin{pgfscope}%
\pgfsys@transformshift{0.688581in}{1.737499in}%
\pgfsys@useobject{currentmarker}{}%
\end{pgfscope}%
\end{pgfscope}%
\begin{pgfscope}%
\definecolor{textcolor}{rgb}{0.000000,0.000000,0.000000}%
\pgfsetstrokecolor{textcolor}%
\pgfsetfillcolor{textcolor}%
\pgftext[x=0.344444in, y=1.689273in, left, base]{\color{textcolor}\rmfamily\fontsize{10.000000}{12.000000}\selectfont \(\displaystyle {\ensuremath{-}20}\)}%
\end{pgfscope}%
\begin{pgfscope}%
\pgfpathrectangle{\pgfqpoint{0.688581in}{0.565123in}}{\pgfqpoint{5.461419in}{2.585803in}}%
\pgfusepath{clip}%
\pgfsetbuttcap%
\pgfsetroundjoin%
\pgfsetlinewidth{0.803000pt}%
\definecolor{currentstroke}{rgb}{0.690196,0.690196,0.690196}%
\pgfsetstrokecolor{currentstroke}%
\pgfsetdash{{0.800000pt}{1.320000pt}}{0.000000pt}%
\pgfpathmoveto{\pgfqpoint{0.688581in}{2.052188in}}%
\pgfpathlineto{\pgfqpoint{6.150000in}{2.052188in}}%
\pgfusepath{stroke}%
\end{pgfscope}%
\begin{pgfscope}%
\pgfsetbuttcap%
\pgfsetroundjoin%
\definecolor{currentfill}{rgb}{0.000000,0.000000,0.000000}%
\pgfsetfillcolor{currentfill}%
\pgfsetlinewidth{0.803000pt}%
\definecolor{currentstroke}{rgb}{0.000000,0.000000,0.000000}%
\pgfsetstrokecolor{currentstroke}%
\pgfsetdash{}{0pt}%
\pgfsys@defobject{currentmarker}{\pgfqpoint{-0.048611in}{0.000000in}}{\pgfqpoint{-0.000000in}{0.000000in}}{%
\pgfpathmoveto{\pgfqpoint{-0.000000in}{0.000000in}}%
\pgfpathlineto{\pgfqpoint{-0.048611in}{0.000000in}}%
\pgfusepath{stroke,fill}%
}%
\begin{pgfscope}%
\pgfsys@transformshift{0.688581in}{2.052188in}%
\pgfsys@useobject{currentmarker}{}%
\end{pgfscope}%
\end{pgfscope}%
\begin{pgfscope}%
\definecolor{textcolor}{rgb}{0.000000,0.000000,0.000000}%
\pgfsetstrokecolor{textcolor}%
\pgfsetfillcolor{textcolor}%
\pgftext[x=0.344444in, y=2.003963in, left, base]{\color{textcolor}\rmfamily\fontsize{10.000000}{12.000000}\selectfont \(\displaystyle {\ensuremath{-}15}\)}%
\end{pgfscope}%
\begin{pgfscope}%
\pgfpathrectangle{\pgfqpoint{0.688581in}{0.565123in}}{\pgfqpoint{5.461419in}{2.585803in}}%
\pgfusepath{clip}%
\pgfsetbuttcap%
\pgfsetroundjoin%
\pgfsetlinewidth{0.803000pt}%
\definecolor{currentstroke}{rgb}{0.690196,0.690196,0.690196}%
\pgfsetstrokecolor{currentstroke}%
\pgfsetdash{{0.800000pt}{1.320000pt}}{0.000000pt}%
\pgfpathmoveto{\pgfqpoint{0.688581in}{2.366878in}}%
\pgfpathlineto{\pgfqpoint{6.150000in}{2.366878in}}%
\pgfusepath{stroke}%
\end{pgfscope}%
\begin{pgfscope}%
\pgfsetbuttcap%
\pgfsetroundjoin%
\definecolor{currentfill}{rgb}{0.000000,0.000000,0.000000}%
\pgfsetfillcolor{currentfill}%
\pgfsetlinewidth{0.803000pt}%
\definecolor{currentstroke}{rgb}{0.000000,0.000000,0.000000}%
\pgfsetstrokecolor{currentstroke}%
\pgfsetdash{}{0pt}%
\pgfsys@defobject{currentmarker}{\pgfqpoint{-0.048611in}{0.000000in}}{\pgfqpoint{-0.000000in}{0.000000in}}{%
\pgfpathmoveto{\pgfqpoint{-0.000000in}{0.000000in}}%
\pgfpathlineto{\pgfqpoint{-0.048611in}{0.000000in}}%
\pgfusepath{stroke,fill}%
}%
\begin{pgfscope}%
\pgfsys@transformshift{0.688581in}{2.366878in}%
\pgfsys@useobject{currentmarker}{}%
\end{pgfscope}%
\end{pgfscope}%
\begin{pgfscope}%
\definecolor{textcolor}{rgb}{0.000000,0.000000,0.000000}%
\pgfsetstrokecolor{textcolor}%
\pgfsetfillcolor{textcolor}%
\pgftext[x=0.344444in, y=2.318652in, left, base]{\color{textcolor}\rmfamily\fontsize{10.000000}{12.000000}\selectfont \(\displaystyle {\ensuremath{-}10}\)}%
\end{pgfscope}%
\begin{pgfscope}%
\pgfpathrectangle{\pgfqpoint{0.688581in}{0.565123in}}{\pgfqpoint{5.461419in}{2.585803in}}%
\pgfusepath{clip}%
\pgfsetbuttcap%
\pgfsetroundjoin%
\pgfsetlinewidth{0.803000pt}%
\definecolor{currentstroke}{rgb}{0.690196,0.690196,0.690196}%
\pgfsetstrokecolor{currentstroke}%
\pgfsetdash{{0.800000pt}{1.320000pt}}{0.000000pt}%
\pgfpathmoveto{\pgfqpoint{0.688581in}{2.681567in}}%
\pgfpathlineto{\pgfqpoint{6.150000in}{2.681567in}}%
\pgfusepath{stroke}%
\end{pgfscope}%
\begin{pgfscope}%
\pgfsetbuttcap%
\pgfsetroundjoin%
\definecolor{currentfill}{rgb}{0.000000,0.000000,0.000000}%
\pgfsetfillcolor{currentfill}%
\pgfsetlinewidth{0.803000pt}%
\definecolor{currentstroke}{rgb}{0.000000,0.000000,0.000000}%
\pgfsetstrokecolor{currentstroke}%
\pgfsetdash{}{0pt}%
\pgfsys@defobject{currentmarker}{\pgfqpoint{-0.048611in}{0.000000in}}{\pgfqpoint{-0.000000in}{0.000000in}}{%
\pgfpathmoveto{\pgfqpoint{-0.000000in}{0.000000in}}%
\pgfpathlineto{\pgfqpoint{-0.048611in}{0.000000in}}%
\pgfusepath{stroke,fill}%
}%
\begin{pgfscope}%
\pgfsys@transformshift{0.688581in}{2.681567in}%
\pgfsys@useobject{currentmarker}{}%
\end{pgfscope}%
\end{pgfscope}%
\begin{pgfscope}%
\definecolor{textcolor}{rgb}{0.000000,0.000000,0.000000}%
\pgfsetstrokecolor{textcolor}%
\pgfsetfillcolor{textcolor}%
\pgftext[x=0.413889in, y=2.633342in, left, base]{\color{textcolor}\rmfamily\fontsize{10.000000}{12.000000}\selectfont \(\displaystyle {\ensuremath{-}5}\)}%
\end{pgfscope}%
\begin{pgfscope}%
\pgfpathrectangle{\pgfqpoint{0.688581in}{0.565123in}}{\pgfqpoint{5.461419in}{2.585803in}}%
\pgfusepath{clip}%
\pgfsetbuttcap%
\pgfsetroundjoin%
\pgfsetlinewidth{0.803000pt}%
\definecolor{currentstroke}{rgb}{0.690196,0.690196,0.690196}%
\pgfsetstrokecolor{currentstroke}%
\pgfsetdash{{0.800000pt}{1.320000pt}}{0.000000pt}%
\pgfpathmoveto{\pgfqpoint{0.688581in}{2.996256in}}%
\pgfpathlineto{\pgfqpoint{6.150000in}{2.996256in}}%
\pgfusepath{stroke}%
\end{pgfscope}%
\begin{pgfscope}%
\pgfsetbuttcap%
\pgfsetroundjoin%
\definecolor{currentfill}{rgb}{0.000000,0.000000,0.000000}%
\pgfsetfillcolor{currentfill}%
\pgfsetlinewidth{0.803000pt}%
\definecolor{currentstroke}{rgb}{0.000000,0.000000,0.000000}%
\pgfsetstrokecolor{currentstroke}%
\pgfsetdash{}{0pt}%
\pgfsys@defobject{currentmarker}{\pgfqpoint{-0.048611in}{0.000000in}}{\pgfqpoint{-0.000000in}{0.000000in}}{%
\pgfpathmoveto{\pgfqpoint{-0.000000in}{0.000000in}}%
\pgfpathlineto{\pgfqpoint{-0.048611in}{0.000000in}}%
\pgfusepath{stroke,fill}%
}%
\begin{pgfscope}%
\pgfsys@transformshift{0.688581in}{2.996256in}%
\pgfsys@useobject{currentmarker}{}%
\end{pgfscope}%
\end{pgfscope}%
\begin{pgfscope}%
\definecolor{textcolor}{rgb}{0.000000,0.000000,0.000000}%
\pgfsetstrokecolor{textcolor}%
\pgfsetfillcolor{textcolor}%
\pgftext[x=0.521914in, y=2.948031in, left, base]{\color{textcolor}\rmfamily\fontsize{10.000000}{12.000000}\selectfont \(\displaystyle {0}\)}%
\end{pgfscope}%
\begin{pgfscope}%
\definecolor{textcolor}{rgb}{0.000000,0.000000,0.000000}%
\pgfsetstrokecolor{textcolor}%
\pgfsetfillcolor{textcolor}%
\pgftext[x=0.288889in,y=1.858025in,,bottom,rotate=90.000000]{\color{textcolor}\rmfamily\fontsize{10.000000}{12.000000}\selectfont Transmission Loss (dB)}%
\end{pgfscope}%
\begin{pgfscope}%
\pgfpathrectangle{\pgfqpoint{0.688581in}{0.565123in}}{\pgfqpoint{5.461419in}{2.585803in}}%
\pgfusepath{clip}%
\pgfsetrectcap%
\pgfsetroundjoin%
\pgfsetlinewidth{2.007500pt}%
\definecolor{currentstroke}{rgb}{0.121569,0.466667,0.705882}%
\pgfsetstrokecolor{currentstroke}%
\pgfsetdash{}{0pt}%
\pgfpathmoveto{\pgfqpoint{0.936827in}{1.279940in}}%
\pgfpathlineto{\pgfqpoint{0.942183in}{1.072875in}}%
\pgfpathlineto{\pgfqpoint{0.963607in}{1.298192in}}%
\pgfpathlineto{\pgfqpoint{0.979675in}{1.442320in}}%
\pgfpathlineto{\pgfqpoint{1.017166in}{1.719247in}}%
\pgfpathlineto{\pgfqpoint{1.043945in}{1.877850in}}%
\pgfpathlineto{\pgfqpoint{1.065369in}{1.992397in}}%
\pgfpathlineto{\pgfqpoint{1.097505in}{2.143448in}}%
\pgfpathlineto{\pgfqpoint{1.140352in}{2.317157in}}%
\pgfpathlineto{\pgfqpoint{1.151064in}{2.357437in}}%
\pgfpathlineto{\pgfqpoint{1.167131in}{2.412822in}}%
\pgfpathlineto{\pgfqpoint{1.177843in}{2.446179in}}%
\pgfpathlineto{\pgfqpoint{1.193911in}{2.497788in}}%
\pgfpathlineto{\pgfqpoint{1.247470in}{2.645692in}}%
\pgfpathlineto{\pgfqpoint{1.263538in}{2.685973in}}%
\pgfpathlineto{\pgfqpoint{1.311741in}{2.786044in}}%
\pgfpathlineto{\pgfqpoint{1.317097in}{2.792967in}}%
\pgfpathlineto{\pgfqpoint{1.322453in}{2.804296in}}%
\pgfpathlineto{\pgfqpoint{1.333164in}{2.818772in}}%
\pgfpathlineto{\pgfqpoint{1.338520in}{2.830100in}}%
\pgfpathlineto{\pgfqpoint{1.386724in}{2.896815in}}%
\pgfpathlineto{\pgfqpoint{1.392079in}{2.900591in}}%
\pgfpathlineto{\pgfqpoint{1.397435in}{2.907514in}}%
\pgfpathlineto{\pgfqpoint{1.402791in}{2.911290in}}%
\pgfpathlineto{\pgfqpoint{1.408147in}{2.918843in}}%
\pgfpathlineto{\pgfqpoint{1.418859in}{2.926395in}}%
\pgfpathlineto{\pgfqpoint{1.424215in}{2.933319in}}%
\pgfpathlineto{\pgfqpoint{1.467062in}{2.962899in}}%
\pgfpathlineto{\pgfqpoint{1.472418in}{2.962899in}}%
\pgfpathlineto{\pgfqpoint{1.483130in}{2.970452in}}%
\pgfpathlineto{\pgfqpoint{1.488486in}{2.970452in}}%
\pgfpathlineto{\pgfqpoint{1.493842in}{2.974228in}}%
\pgfpathlineto{\pgfqpoint{1.504554in}{2.974228in}}%
\pgfpathlineto{\pgfqpoint{1.509909in}{2.978004in}}%
\pgfpathlineto{\pgfqpoint{1.515265in}{2.978004in}}%
\pgfpathlineto{\pgfqpoint{1.520621in}{2.981781in}}%
\pgfpathlineto{\pgfqpoint{1.563468in}{2.981781in}}%
\pgfpathlineto{\pgfqpoint{1.568824in}{2.984928in}}%
\pgfpathlineto{\pgfqpoint{1.574180in}{2.981781in}}%
\pgfpathlineto{\pgfqpoint{1.595604in}{2.981781in}}%
\pgfpathlineto{\pgfqpoint{1.600960in}{2.978004in}}%
\pgfpathlineto{\pgfqpoint{1.611672in}{2.978004in}}%
\pgfpathlineto{\pgfqpoint{1.617028in}{2.974228in}}%
\pgfpathlineto{\pgfqpoint{1.627739in}{2.974228in}}%
\pgfpathlineto{\pgfqpoint{1.633095in}{2.970452in}}%
\pgfpathlineto{\pgfqpoint{1.638451in}{2.970452in}}%
\pgfpathlineto{\pgfqpoint{1.643807in}{2.966676in}}%
\pgfpathlineto{\pgfqpoint{1.649163in}{2.966676in}}%
\pgfpathlineto{\pgfqpoint{1.654519in}{2.962899in}}%
\pgfpathlineto{\pgfqpoint{1.659875in}{2.962899in}}%
\pgfpathlineto{\pgfqpoint{1.665231in}{2.959123in}}%
\pgfpathlineto{\pgfqpoint{1.670587in}{2.959123in}}%
\pgfpathlineto{\pgfqpoint{1.681298in}{2.952200in}}%
\pgfpathlineto{\pgfqpoint{1.686654in}{2.952200in}}%
\pgfpathlineto{\pgfqpoint{1.692010in}{2.948424in}}%
\pgfpathlineto{\pgfqpoint{1.697366in}{2.948424in}}%
\pgfpathlineto{\pgfqpoint{1.713434in}{2.937095in}}%
\pgfpathlineto{\pgfqpoint{1.718790in}{2.937095in}}%
\pgfpathlineto{\pgfqpoint{1.740213in}{2.922619in}}%
\pgfpathlineto{\pgfqpoint{1.745569in}{2.922619in}}%
\pgfpathlineto{\pgfqpoint{1.766993in}{2.907514in}}%
\pgfpathlineto{\pgfqpoint{1.772349in}{2.907514in}}%
\pgfpathlineto{\pgfqpoint{1.804484in}{2.885486in}}%
\pgfpathlineto{\pgfqpoint{1.809840in}{2.885486in}}%
\pgfpathlineto{\pgfqpoint{1.858043in}{2.852129in}}%
\pgfpathlineto{\pgfqpoint{1.863399in}{2.852129in}}%
\pgfpathlineto{\pgfqpoint{1.906247in}{2.822548in}}%
\pgfpathlineto{\pgfqpoint{1.911602in}{2.822548in}}%
\pgfpathlineto{\pgfqpoint{1.954450in}{2.792967in}}%
\pgfpathlineto{\pgfqpoint{1.959806in}{2.792967in}}%
\pgfpathlineto{\pgfqpoint{1.975873in}{2.782268in}}%
\pgfpathlineto{\pgfqpoint{2.008009in}{2.760239in}}%
\pgfpathlineto{\pgfqpoint{2.013365in}{2.760239in}}%
\pgfpathlineto{\pgfqpoint{2.050856in}{2.734435in}}%
\pgfpathlineto{\pgfqpoint{2.056212in}{2.734435in}}%
\pgfpathlineto{\pgfqpoint{2.109771in}{2.697301in}}%
\pgfpathlineto{\pgfqpoint{2.115127in}{2.697301in}}%
\pgfpathlineto{\pgfqpoint{2.216889in}{2.626811in}}%
\pgfpathlineto{\pgfqpoint{2.232957in}{2.616112in}}%
\pgfpathlineto{\pgfqpoint{2.334719in}{2.545621in}}%
\pgfpathlineto{\pgfqpoint{2.350787in}{2.534922in}}%
\pgfpathlineto{\pgfqpoint{2.398990in}{2.501565in}}%
\pgfpathlineto{\pgfqpoint{2.404346in}{2.501565in}}%
\pgfpathlineto{\pgfqpoint{2.495396in}{2.438627in}}%
\pgfpathlineto{\pgfqpoint{2.500752in}{2.438627in}}%
\pgfpathlineto{\pgfqpoint{2.516820in}{2.427927in}}%
\pgfpathlineto{\pgfqpoint{2.559667in}{2.398347in}}%
\pgfpathlineto{\pgfqpoint{2.565023in}{2.398347in}}%
\pgfpathlineto{\pgfqpoint{2.607870in}{2.368766in}}%
\pgfpathlineto{\pgfqpoint{2.613226in}{2.368766in}}%
\pgfpathlineto{\pgfqpoint{2.650718in}{2.342961in}}%
\pgfpathlineto{\pgfqpoint{2.656074in}{2.342961in}}%
\pgfpathlineto{\pgfqpoint{2.688209in}{2.320933in}}%
\pgfpathlineto{\pgfqpoint{2.693565in}{2.320933in}}%
\pgfpathlineto{\pgfqpoint{2.725700in}{2.298275in}}%
\pgfpathlineto{\pgfqpoint{2.731056in}{2.298275in}}%
\pgfpathlineto{\pgfqpoint{2.747124in}{2.287576in}}%
\pgfpathlineto{\pgfqpoint{2.752480in}{2.283800in}}%
\pgfpathlineto{\pgfqpoint{2.757836in}{2.283800in}}%
\pgfpathlineto{\pgfqpoint{2.784615in}{2.265548in}}%
\pgfpathlineto{\pgfqpoint{2.789971in}{2.265548in}}%
\pgfpathlineto{\pgfqpoint{2.822107in}{2.243519in}}%
\pgfpathlineto{\pgfqpoint{2.827463in}{2.243519in}}%
\pgfpathlineto{\pgfqpoint{2.848886in}{2.228414in}}%
\pgfpathlineto{\pgfqpoint{2.854242in}{2.228414in}}%
\pgfpathlineto{\pgfqpoint{2.875666in}{2.213939in}}%
\pgfpathlineto{\pgfqpoint{2.881022in}{2.213939in}}%
\pgfpathlineto{\pgfqpoint{2.902445in}{2.198833in}}%
\pgfpathlineto{\pgfqpoint{2.907801in}{2.198833in}}%
\pgfpathlineto{\pgfqpoint{2.929225in}{2.184358in}}%
\pgfpathlineto{\pgfqpoint{2.934581in}{2.184358in}}%
\pgfpathlineto{\pgfqpoint{2.956004in}{2.169253in}}%
\pgfpathlineto{\pgfqpoint{2.961360in}{2.169253in}}%
\pgfpathlineto{\pgfqpoint{2.966716in}{2.161700in}}%
\pgfpathlineto{\pgfqpoint{2.972072in}{2.161700in}}%
\pgfpathlineto{\pgfqpoint{2.982784in}{2.154777in}}%
\pgfpathlineto{\pgfqpoint{2.988140in}{2.154777in}}%
\pgfpathlineto{\pgfqpoint{2.998852in}{2.147224in}}%
\pgfpathlineto{\pgfqpoint{3.004208in}{2.147224in}}%
\pgfpathlineto{\pgfqpoint{3.030987in}{2.128972in}}%
\pgfpathlineto{\pgfqpoint{3.036343in}{2.128972in}}%
\pgfpathlineto{\pgfqpoint{3.041699in}{2.125196in}}%
\pgfpathlineto{\pgfqpoint{3.047055in}{2.125196in}}%
\pgfpathlineto{\pgfqpoint{3.063123in}{2.113867in}}%
\pgfpathlineto{\pgfqpoint{3.068479in}{2.113867in}}%
\pgfpathlineto{\pgfqpoint{3.084546in}{2.103168in}}%
\pgfpathlineto{\pgfqpoint{3.089902in}{2.103168in}}%
\pgfpathlineto{\pgfqpoint{3.095258in}{2.099392in}}%
\pgfpathlineto{\pgfqpoint{3.100614in}{2.099392in}}%
\pgfpathlineto{\pgfqpoint{3.132749in}{2.077363in}}%
\pgfpathlineto{\pgfqpoint{3.138105in}{2.077363in}}%
\pgfpathlineto{\pgfqpoint{3.143461in}{2.073587in}}%
\pgfpathlineto{\pgfqpoint{3.148817in}{2.073587in}}%
\pgfpathlineto{\pgfqpoint{3.164885in}{2.062258in}}%
\pgfpathlineto{\pgfqpoint{3.170241in}{2.062258in}}%
\pgfpathlineto{\pgfqpoint{3.180953in}{2.054706in}}%
\pgfpathlineto{\pgfqpoint{3.186309in}{2.054706in}}%
\pgfpathlineto{\pgfqpoint{3.197020in}{2.047783in}}%
\pgfpathlineto{\pgfqpoint{3.202376in}{2.047783in}}%
\pgfpathlineto{\pgfqpoint{3.213088in}{2.040230in}}%
\pgfpathlineto{\pgfqpoint{3.218444in}{2.040230in}}%
\pgfpathlineto{\pgfqpoint{3.234512in}{2.028901in}}%
\pgfpathlineto{\pgfqpoint{3.239868in}{2.028901in}}%
\pgfpathlineto{\pgfqpoint{3.255935in}{2.018202in}}%
\pgfpathlineto{\pgfqpoint{3.261291in}{2.018202in}}%
\pgfpathlineto{\pgfqpoint{3.272003in}{2.010649in}}%
\pgfpathlineto{\pgfqpoint{3.277359in}{2.010649in}}%
\pgfpathlineto{\pgfqpoint{3.288071in}{2.003097in}}%
\pgfpathlineto{\pgfqpoint{3.293427in}{2.003097in}}%
\pgfpathlineto{\pgfqpoint{3.298783in}{1.999320in}}%
\pgfpathlineto{\pgfqpoint{3.304138in}{1.999320in}}%
\pgfpathlineto{\pgfqpoint{3.314850in}{1.992397in}}%
\pgfpathlineto{\pgfqpoint{3.320206in}{1.992397in}}%
\pgfpathlineto{\pgfqpoint{3.330918in}{1.984845in}}%
\pgfpathlineto{\pgfqpoint{3.336274in}{1.984845in}}%
\pgfpathlineto{\pgfqpoint{3.352342in}{1.973516in}}%
\pgfpathlineto{\pgfqpoint{3.357698in}{1.973516in}}%
\pgfpathlineto{\pgfqpoint{3.363053in}{1.970369in}}%
\pgfpathlineto{\pgfqpoint{3.368409in}{1.970369in}}%
\pgfpathlineto{\pgfqpoint{3.379121in}{1.962816in}}%
\pgfpathlineto{\pgfqpoint{3.384477in}{1.962816in}}%
\pgfpathlineto{\pgfqpoint{3.395189in}{1.955264in}}%
\pgfpathlineto{\pgfqpoint{3.400545in}{1.955264in}}%
\pgfpathlineto{\pgfqpoint{3.416613in}{1.943935in}}%
\pgfpathlineto{\pgfqpoint{3.421968in}{1.943935in}}%
\pgfpathlineto{\pgfqpoint{3.427324in}{1.940788in}}%
\pgfpathlineto{\pgfqpoint{3.432680in}{1.940788in}}%
\pgfpathlineto{\pgfqpoint{3.448748in}{1.929459in}}%
\pgfpathlineto{\pgfqpoint{3.454104in}{1.929459in}}%
\pgfpathlineto{\pgfqpoint{3.459460in}{1.925683in}}%
\pgfpathlineto{\pgfqpoint{3.464816in}{1.925683in}}%
\pgfpathlineto{\pgfqpoint{3.491595in}{1.907431in}}%
\pgfpathlineto{\pgfqpoint{3.496951in}{1.911207in}}%
\pgfpathlineto{\pgfqpoint{3.513019in}{1.899878in}}%
\pgfpathlineto{\pgfqpoint{3.518375in}{1.899878in}}%
\pgfpathlineto{\pgfqpoint{3.523731in}{1.896102in}}%
\pgfpathlineto{\pgfqpoint{3.529087in}{1.896102in}}%
\pgfpathlineto{\pgfqpoint{3.539798in}{1.888550in}}%
\pgfpathlineto{\pgfqpoint{3.545154in}{1.888550in}}%
\pgfpathlineto{\pgfqpoint{3.555866in}{1.881626in}}%
\pgfpathlineto{\pgfqpoint{3.561222in}{1.881626in}}%
\pgfpathlineto{\pgfqpoint{3.566578in}{1.877850in}}%
\pgfpathlineto{\pgfqpoint{3.571934in}{1.877850in}}%
\pgfpathlineto{\pgfqpoint{3.582646in}{1.870298in}}%
\pgfpathlineto{\pgfqpoint{3.588002in}{1.870298in}}%
\pgfpathlineto{\pgfqpoint{3.598713in}{1.862745in}}%
\pgfpathlineto{\pgfqpoint{3.604069in}{1.862745in}}%
\pgfpathlineto{\pgfqpoint{3.609425in}{1.859598in}}%
\pgfpathlineto{\pgfqpoint{3.614781in}{1.859598in}}%
\pgfpathlineto{\pgfqpoint{3.625493in}{1.852046in}}%
\pgfpathlineto{\pgfqpoint{3.630849in}{1.852046in}}%
\pgfpathlineto{\pgfqpoint{3.636205in}{1.848269in}}%
\pgfpathlineto{\pgfqpoint{3.641561in}{1.848269in}}%
\pgfpathlineto{\pgfqpoint{3.652272in}{1.840717in}}%
\pgfpathlineto{\pgfqpoint{3.657628in}{1.840717in}}%
\pgfpathlineto{\pgfqpoint{3.668340in}{1.833164in}}%
\pgfpathlineto{\pgfqpoint{3.673696in}{1.833164in}}%
\pgfpathlineto{\pgfqpoint{3.679052in}{1.830017in}}%
\pgfpathlineto{\pgfqpoint{3.684408in}{1.830017in}}%
\pgfpathlineto{\pgfqpoint{3.689764in}{1.826241in}}%
\pgfpathlineto{\pgfqpoint{3.695120in}{1.826241in}}%
\pgfpathlineto{\pgfqpoint{3.705832in}{1.818689in}}%
\pgfpathlineto{\pgfqpoint{3.711187in}{1.818689in}}%
\pgfpathlineto{\pgfqpoint{3.727255in}{1.807360in}}%
\pgfpathlineto{\pgfqpoint{3.732611in}{1.807360in}}%
\pgfpathlineto{\pgfqpoint{3.743323in}{1.800437in}}%
\pgfpathlineto{\pgfqpoint{3.748679in}{1.800437in}}%
\pgfpathlineto{\pgfqpoint{3.754035in}{1.796660in}}%
\pgfpathlineto{\pgfqpoint{3.759391in}{1.796660in}}%
\pgfpathlineto{\pgfqpoint{3.764747in}{1.792884in}}%
\pgfpathlineto{\pgfqpoint{3.770102in}{1.792884in}}%
\pgfpathlineto{\pgfqpoint{3.780814in}{1.785332in}}%
\pgfpathlineto{\pgfqpoint{3.786170in}{1.785332in}}%
\pgfpathlineto{\pgfqpoint{3.791526in}{1.781555in}}%
\pgfpathlineto{\pgfqpoint{3.796882in}{1.781555in}}%
\pgfpathlineto{\pgfqpoint{3.807594in}{1.774632in}}%
\pgfpathlineto{\pgfqpoint{3.812950in}{1.774632in}}%
\pgfpathlineto{\pgfqpoint{3.823662in}{1.767080in}}%
\pgfpathlineto{\pgfqpoint{3.829017in}{1.767080in}}%
\pgfpathlineto{\pgfqpoint{3.834373in}{1.763303in}}%
\pgfpathlineto{\pgfqpoint{3.839729in}{1.763303in}}%
\pgfpathlineto{\pgfqpoint{3.850441in}{1.755751in}}%
\pgfpathlineto{\pgfqpoint{3.855797in}{1.755751in}}%
\pgfpathlineto{\pgfqpoint{3.866509in}{1.748828in}}%
\pgfpathlineto{\pgfqpoint{3.877221in}{1.748828in}}%
\pgfpathlineto{\pgfqpoint{3.893288in}{1.737499in}}%
\pgfpathlineto{\pgfqpoint{3.904000in}{1.737499in}}%
\pgfpathlineto{\pgfqpoint{3.914712in}{1.729946in}}%
\pgfpathlineto{\pgfqpoint{3.920068in}{1.729946in}}%
\pgfpathlineto{\pgfqpoint{3.925424in}{1.726170in}}%
\pgfpathlineto{\pgfqpoint{3.930780in}{1.726170in}}%
\pgfpathlineto{\pgfqpoint{3.936136in}{1.723023in}}%
\pgfpathlineto{\pgfqpoint{3.941491in}{1.723023in}}%
\pgfpathlineto{\pgfqpoint{3.946847in}{1.719247in}}%
\pgfpathlineto{\pgfqpoint{3.952203in}{1.719247in}}%
\pgfpathlineto{\pgfqpoint{3.968271in}{1.707918in}}%
\pgfpathlineto{\pgfqpoint{3.973627in}{1.707918in}}%
\pgfpathlineto{\pgfqpoint{3.978983in}{1.704142in}}%
\pgfpathlineto{\pgfqpoint{3.984339in}{1.704142in}}%
\pgfpathlineto{\pgfqpoint{3.989695in}{1.700365in}}%
\pgfpathlineto{\pgfqpoint{3.995051in}{1.700365in}}%
\pgfpathlineto{\pgfqpoint{4.000406in}{1.696589in}}%
\pgfpathlineto{\pgfqpoint{4.005762in}{1.696589in}}%
\pgfpathlineto{\pgfqpoint{4.016474in}{1.689666in}}%
\pgfpathlineto{\pgfqpoint{4.027186in}{1.689666in}}%
\pgfpathlineto{\pgfqpoint{4.043254in}{1.678337in}}%
\pgfpathlineto{\pgfqpoint{4.048610in}{1.678337in}}%
\pgfpathlineto{\pgfqpoint{4.053966in}{1.674561in}}%
\pgfpathlineto{\pgfqpoint{4.059321in}{1.674561in}}%
\pgfpathlineto{\pgfqpoint{4.070033in}{1.667638in}}%
\pgfpathlineto{\pgfqpoint{4.075389in}{1.667638in}}%
\pgfpathlineto{\pgfqpoint{4.080745in}{1.663861in}}%
\pgfpathlineto{\pgfqpoint{4.086101in}{1.663861in}}%
\pgfpathlineto{\pgfqpoint{4.096813in}{1.656309in}}%
\pgfpathlineto{\pgfqpoint{4.107525in}{1.656309in}}%
\pgfpathlineto{\pgfqpoint{4.118236in}{1.648756in}}%
\pgfpathlineto{\pgfqpoint{4.123592in}{1.648756in}}%
\pgfpathlineto{\pgfqpoint{4.128948in}{1.644980in}}%
\pgfpathlineto{\pgfqpoint{4.134304in}{1.644980in}}%
\pgfpathlineto{\pgfqpoint{4.139660in}{1.641833in}}%
\pgfpathlineto{\pgfqpoint{4.145016in}{1.641833in}}%
\pgfpathlineto{\pgfqpoint{4.155728in}{1.634281in}}%
\pgfpathlineto{\pgfqpoint{4.166440in}{1.634281in}}%
\pgfpathlineto{\pgfqpoint{4.171796in}{1.630504in}}%
\pgfpathlineto{\pgfqpoint{4.177151in}{1.630504in}}%
\pgfpathlineto{\pgfqpoint{4.187863in}{1.622952in}}%
\pgfpathlineto{\pgfqpoint{4.193219in}{1.622952in}}%
\pgfpathlineto{\pgfqpoint{4.209287in}{1.612252in}}%
\pgfpathlineto{\pgfqpoint{4.219999in}{1.612252in}}%
\pgfpathlineto{\pgfqpoint{4.225355in}{1.608476in}}%
\pgfpathlineto{\pgfqpoint{4.230710in}{1.608476in}}%
\pgfpathlineto{\pgfqpoint{4.236066in}{1.604700in}}%
\pgfpathlineto{\pgfqpoint{4.241422in}{1.604700in}}%
\pgfpathlineto{\pgfqpoint{4.246778in}{1.600924in}}%
\pgfpathlineto{\pgfqpoint{4.252134in}{1.600924in}}%
\pgfpathlineto{\pgfqpoint{4.268202in}{1.589595in}}%
\pgfpathlineto{\pgfqpoint{4.273558in}{1.589595in}}%
\pgfpathlineto{\pgfqpoint{4.278914in}{1.586448in}}%
\pgfpathlineto{\pgfqpoint{4.289625in}{1.586448in}}%
\pgfpathlineto{\pgfqpoint{4.294981in}{1.582672in}}%
\pgfpathlineto{\pgfqpoint{4.300337in}{1.582672in}}%
\pgfpathlineto{\pgfqpoint{4.305693in}{1.578895in}}%
\pgfpathlineto{\pgfqpoint{4.311049in}{1.578895in}}%
\pgfpathlineto{\pgfqpoint{4.316405in}{1.575119in}}%
\pgfpathlineto{\pgfqpoint{4.321761in}{1.575119in}}%
\pgfpathlineto{\pgfqpoint{4.348540in}{1.556867in}}%
\pgfpathlineto{\pgfqpoint{4.369964in}{1.556867in}}%
\pgfpathlineto{\pgfqpoint{4.386032in}{1.545538in}}%
\pgfpathlineto{\pgfqpoint{4.391388in}{1.545538in}}%
\pgfpathlineto{\pgfqpoint{4.396744in}{1.541762in}}%
\pgfpathlineto{\pgfqpoint{4.402100in}{1.541762in}}%
\pgfpathlineto{\pgfqpoint{4.412811in}{1.534209in}}%
\pgfpathlineto{\pgfqpoint{4.418167in}{1.537986in}}%
\pgfpathlineto{\pgfqpoint{4.423523in}{1.531062in}}%
\pgfpathlineto{\pgfqpoint{4.434235in}{1.531062in}}%
\pgfpathlineto{\pgfqpoint{4.439591in}{1.527286in}}%
\pgfpathlineto{\pgfqpoint{4.444947in}{1.527286in}}%
\pgfpathlineto{\pgfqpoint{4.450303in}{1.523510in}}%
\pgfpathlineto{\pgfqpoint{4.455659in}{1.527286in}}%
\pgfpathlineto{\pgfqpoint{4.461015in}{1.519734in}}%
\pgfpathlineto{\pgfqpoint{4.466370in}{1.515957in}}%
\pgfpathlineto{\pgfqpoint{4.471726in}{1.519734in}}%
\pgfpathlineto{\pgfqpoint{4.477082in}{1.512181in}}%
\pgfpathlineto{\pgfqpoint{4.482438in}{1.512181in}}%
\pgfpathlineto{\pgfqpoint{4.487794in}{1.508405in}}%
\pgfpathlineto{\pgfqpoint{4.493150in}{1.508405in}}%
\pgfpathlineto{\pgfqpoint{4.498506in}{1.515957in}}%
\pgfpathlineto{\pgfqpoint{4.503862in}{1.508405in}}%
\pgfpathlineto{\pgfqpoint{4.509218in}{1.508405in}}%
\pgfpathlineto{\pgfqpoint{4.514574in}{1.504629in}}%
\pgfpathlineto{\pgfqpoint{4.519930in}{1.504629in}}%
\pgfpathlineto{\pgfqpoint{4.530641in}{1.497705in}}%
\pgfpathlineto{\pgfqpoint{4.535997in}{1.497705in}}%
\pgfpathlineto{\pgfqpoint{4.541353in}{1.493929in}}%
\pgfpathlineto{\pgfqpoint{4.557421in}{1.493929in}}%
\pgfpathlineto{\pgfqpoint{4.562777in}{1.482600in}}%
\pgfpathlineto{\pgfqpoint{4.568133in}{1.490153in}}%
\pgfpathlineto{\pgfqpoint{4.573489in}{1.486377in}}%
\pgfpathlineto{\pgfqpoint{4.578844in}{1.490153in}}%
\pgfpathlineto{\pgfqpoint{4.584200in}{1.482600in}}%
\pgfpathlineto{\pgfqpoint{4.594912in}{1.475677in}}%
\pgfpathlineto{\pgfqpoint{4.600268in}{1.475677in}}%
\pgfpathlineto{\pgfqpoint{4.605624in}{1.486377in}}%
\pgfpathlineto{\pgfqpoint{4.610980in}{1.471901in}}%
\pgfpathlineto{\pgfqpoint{4.616336in}{1.468125in}}%
\pgfpathlineto{\pgfqpoint{4.621692in}{1.468125in}}%
\pgfpathlineto{\pgfqpoint{4.632404in}{1.475677in}}%
\pgfpathlineto{\pgfqpoint{4.637759in}{1.464348in}}%
\pgfpathlineto{\pgfqpoint{4.643115in}{1.460572in}}%
\pgfpathlineto{\pgfqpoint{4.648471in}{1.468125in}}%
\pgfpathlineto{\pgfqpoint{4.653827in}{1.460572in}}%
\pgfpathlineto{\pgfqpoint{4.659183in}{1.468125in}}%
\pgfpathlineto{\pgfqpoint{4.664539in}{1.456796in}}%
\pgfpathlineto{\pgfqpoint{4.669895in}{1.464348in}}%
\pgfpathlineto{\pgfqpoint{4.691319in}{1.449873in}}%
\pgfpathlineto{\pgfqpoint{4.696674in}{1.442320in}}%
\pgfpathlineto{\pgfqpoint{4.702030in}{1.446096in}}%
\pgfpathlineto{\pgfqpoint{4.707386in}{1.442320in}}%
\pgfpathlineto{\pgfqpoint{4.712742in}{1.442320in}}%
\pgfpathlineto{\pgfqpoint{4.718098in}{1.438544in}}%
\pgfpathlineto{\pgfqpoint{4.723454in}{1.438544in}}%
\pgfpathlineto{\pgfqpoint{4.728810in}{1.434768in}}%
\pgfpathlineto{\pgfqpoint{4.734166in}{1.438544in}}%
\pgfpathlineto{\pgfqpoint{4.739522in}{1.438544in}}%
\pgfpathlineto{\pgfqpoint{4.744878in}{1.430991in}}%
\pgfpathlineto{\pgfqpoint{4.750234in}{1.434768in}}%
\pgfpathlineto{\pgfqpoint{4.755589in}{1.427215in}}%
\pgfpathlineto{\pgfqpoint{4.766301in}{1.427215in}}%
\pgfpathlineto{\pgfqpoint{4.771657in}{1.420292in}}%
\pgfpathlineto{\pgfqpoint{4.777013in}{1.416516in}}%
\pgfpathlineto{\pgfqpoint{4.782369in}{1.420292in}}%
\pgfpathlineto{\pgfqpoint{4.787725in}{1.420292in}}%
\pgfpathlineto{\pgfqpoint{4.793081in}{1.416516in}}%
\pgfpathlineto{\pgfqpoint{4.798437in}{1.420292in}}%
\pgfpathlineto{\pgfqpoint{4.803793in}{1.408963in}}%
\pgfpathlineto{\pgfqpoint{4.814504in}{1.416516in}}%
\pgfpathlineto{\pgfqpoint{4.819860in}{1.416516in}}%
\pgfpathlineto{\pgfqpoint{4.830572in}{1.408963in}}%
\pgfpathlineto{\pgfqpoint{4.835928in}{1.401410in}}%
\pgfpathlineto{\pgfqpoint{4.841284in}{1.405187in}}%
\pgfpathlineto{\pgfqpoint{4.846640in}{1.416516in}}%
\pgfpathlineto{\pgfqpoint{4.857352in}{1.408963in}}%
\pgfpathlineto{\pgfqpoint{4.862708in}{1.412739in}}%
\pgfpathlineto{\pgfqpoint{4.868064in}{1.397634in}}%
\pgfpathlineto{\pgfqpoint{4.878775in}{1.405187in}}%
\pgfpathlineto{\pgfqpoint{4.884131in}{1.390711in}}%
\pgfpathlineto{\pgfqpoint{4.900199in}{1.390711in}}%
\pgfpathlineto{\pgfqpoint{4.905555in}{1.386935in}}%
\pgfpathlineto{\pgfqpoint{4.910911in}{1.375606in}}%
\pgfpathlineto{\pgfqpoint{4.916267in}{1.379382in}}%
\pgfpathlineto{\pgfqpoint{4.921623in}{1.375606in}}%
\pgfpathlineto{\pgfqpoint{4.926978in}{1.383158in}}%
\pgfpathlineto{\pgfqpoint{4.932334in}{1.371830in}}%
\pgfpathlineto{\pgfqpoint{4.937690in}{1.386935in}}%
\pgfpathlineto{\pgfqpoint{4.953758in}{1.386935in}}%
\pgfpathlineto{\pgfqpoint{4.959114in}{1.375606in}}%
\pgfpathlineto{\pgfqpoint{4.964470in}{1.375606in}}%
\pgfpathlineto{\pgfqpoint{4.969826in}{1.379382in}}%
\pgfpathlineto{\pgfqpoint{4.975182in}{1.379382in}}%
\pgfpathlineto{\pgfqpoint{4.980538in}{1.383158in}}%
\pgfpathlineto{\pgfqpoint{4.985893in}{1.368053in}}%
\pgfpathlineto{\pgfqpoint{4.991249in}{1.357354in}}%
\pgfpathlineto{\pgfqpoint{4.996605in}{1.357354in}}%
\pgfpathlineto{\pgfqpoint{5.001961in}{1.346025in}}%
\pgfpathlineto{\pgfqpoint{5.007317in}{1.349801in}}%
\pgfpathlineto{\pgfqpoint{5.012673in}{1.357354in}}%
\pgfpathlineto{\pgfqpoint{5.023385in}{1.357354in}}%
\pgfpathlineto{\pgfqpoint{5.028741in}{1.361130in}}%
\pgfpathlineto{\pgfqpoint{5.034097in}{1.361130in}}%
\pgfpathlineto{\pgfqpoint{5.039453in}{1.349801in}}%
\pgfpathlineto{\pgfqpoint{5.055520in}{1.349801in}}%
\pgfpathlineto{\pgfqpoint{5.060876in}{1.339102in}}%
\pgfpathlineto{\pgfqpoint{5.066232in}{1.346025in}}%
\pgfpathlineto{\pgfqpoint{5.071588in}{1.349801in}}%
\pgfpathlineto{\pgfqpoint{5.076944in}{1.371830in}}%
\pgfpathlineto{\pgfqpoint{5.082300in}{1.342249in}}%
\pgfpathlineto{\pgfqpoint{5.087656in}{1.342249in}}%
\pgfpathlineto{\pgfqpoint{5.093012in}{1.357354in}}%
\pgfpathlineto{\pgfqpoint{5.098368in}{1.339102in}}%
\pgfpathlineto{\pgfqpoint{5.103723in}{1.353578in}}%
\pgfpathlineto{\pgfqpoint{5.109079in}{1.349801in}}%
\pgfpathlineto{\pgfqpoint{5.114435in}{1.339102in}}%
\pgfpathlineto{\pgfqpoint{5.119791in}{1.339102in}}%
\pgfpathlineto{\pgfqpoint{5.125147in}{1.331549in}}%
\pgfpathlineto{\pgfqpoint{5.130503in}{1.353578in}}%
\pgfpathlineto{\pgfqpoint{5.135859in}{1.346025in}}%
\pgfpathlineto{\pgfqpoint{5.146571in}{1.316444in}}%
\pgfpathlineto{\pgfqpoint{5.151927in}{1.316444in}}%
\pgfpathlineto{\pgfqpoint{5.157283in}{1.331549in}}%
\pgfpathlineto{\pgfqpoint{5.162638in}{1.320221in}}%
\pgfpathlineto{\pgfqpoint{5.167994in}{1.316444in}}%
\pgfpathlineto{\pgfqpoint{5.173350in}{1.316444in}}%
\pgfpathlineto{\pgfqpoint{5.178706in}{1.331549in}}%
\pgfpathlineto{\pgfqpoint{5.184062in}{1.331549in}}%
\pgfpathlineto{\pgfqpoint{5.194774in}{1.316444in}}%
\pgfpathlineto{\pgfqpoint{5.200130in}{1.316444in}}%
\pgfpathlineto{\pgfqpoint{5.205486in}{1.323997in}}%
\pgfpathlineto{\pgfqpoint{5.210842in}{1.323997in}}%
\pgfpathlineto{\pgfqpoint{5.221553in}{1.301969in}}%
\pgfpathlineto{\pgfqpoint{5.226909in}{1.309521in}}%
\pgfpathlineto{\pgfqpoint{5.232265in}{1.298192in}}%
\pgfpathlineto{\pgfqpoint{5.237621in}{1.320221in}}%
\pgfpathlineto{\pgfqpoint{5.242977in}{1.316444in}}%
\pgfpathlineto{\pgfqpoint{5.248333in}{1.305745in}}%
\pgfpathlineto{\pgfqpoint{5.253689in}{1.313297in}}%
\pgfpathlineto{\pgfqpoint{5.259045in}{1.316444in}}%
\pgfpathlineto{\pgfqpoint{5.264401in}{1.313297in}}%
\pgfpathlineto{\pgfqpoint{5.269757in}{1.320221in}}%
\pgfpathlineto{\pgfqpoint{5.275112in}{1.294416in}}%
\pgfpathlineto{\pgfqpoint{5.280468in}{1.301969in}}%
\pgfpathlineto{\pgfqpoint{5.285824in}{1.286864in}}%
\pgfpathlineto{\pgfqpoint{5.291180in}{1.294416in}}%
\pgfpathlineto{\pgfqpoint{5.296536in}{1.320221in}}%
\pgfpathlineto{\pgfqpoint{5.301892in}{1.309521in}}%
\pgfpathlineto{\pgfqpoint{5.307248in}{1.301969in}}%
\pgfpathlineto{\pgfqpoint{5.312604in}{1.301969in}}%
\pgfpathlineto{\pgfqpoint{5.317960in}{1.316444in}}%
\pgfpathlineto{\pgfqpoint{5.323316in}{1.323997in}}%
\pgfpathlineto{\pgfqpoint{5.328672in}{1.301969in}}%
\pgfpathlineto{\pgfqpoint{5.334027in}{1.286864in}}%
\pgfpathlineto{\pgfqpoint{5.339383in}{1.283717in}}%
\pgfpathlineto{\pgfqpoint{5.350095in}{1.331549in}}%
\pgfpathlineto{\pgfqpoint{5.355451in}{1.309521in}}%
\pgfpathlineto{\pgfqpoint{5.366163in}{1.276164in}}%
\pgfpathlineto{\pgfqpoint{5.371519in}{1.313297in}}%
\pgfpathlineto{\pgfqpoint{5.376875in}{1.283717in}}%
\pgfpathlineto{\pgfqpoint{5.382231in}{1.261059in}}%
\pgfpathlineto{\pgfqpoint{5.387587in}{1.283717in}}%
\pgfpathlineto{\pgfqpoint{5.392942in}{1.279940in}}%
\pgfpathlineto{\pgfqpoint{5.398298in}{1.286864in}}%
\pgfpathlineto{\pgfqpoint{5.403654in}{1.283717in}}%
\pgfpathlineto{\pgfqpoint{5.409010in}{1.268612in}}%
\pgfpathlineto{\pgfqpoint{5.414366in}{1.264835in}}%
\pgfpathlineto{\pgfqpoint{5.419722in}{1.283717in}}%
\pgfpathlineto{\pgfqpoint{5.425078in}{1.283717in}}%
\pgfpathlineto{\pgfqpoint{5.430434in}{1.305745in}}%
\pgfpathlineto{\pgfqpoint{5.435790in}{1.250360in}}%
\pgfpathlineto{\pgfqpoint{5.441146in}{1.254136in}}%
\pgfpathlineto{\pgfqpoint{5.446502in}{1.276164in}}%
\pgfpathlineto{\pgfqpoint{5.451857in}{1.257912in}}%
\pgfpathlineto{\pgfqpoint{5.457213in}{1.264835in}}%
\pgfpathlineto{\pgfqpoint{5.462569in}{1.228331in}}%
\pgfpathlineto{\pgfqpoint{5.467925in}{1.264835in}}%
\pgfpathlineto{\pgfqpoint{5.473281in}{1.268612in}}%
\pgfpathlineto{\pgfqpoint{5.478637in}{1.283717in}}%
\pgfpathlineto{\pgfqpoint{5.483993in}{1.283717in}}%
\pgfpathlineto{\pgfqpoint{5.489349in}{1.261059in}}%
\pgfpathlineto{\pgfqpoint{5.494705in}{1.272388in}}%
\pgfpathlineto{\pgfqpoint{5.500061in}{1.268612in}}%
\pgfpathlineto{\pgfqpoint{5.505417in}{1.272388in}}%
\pgfpathlineto{\pgfqpoint{5.510772in}{1.283717in}}%
\pgfpathlineto{\pgfqpoint{5.516128in}{1.246583in}}%
\pgfpathlineto{\pgfqpoint{5.521484in}{1.224555in}}%
\pgfpathlineto{\pgfqpoint{5.526840in}{1.217002in}}%
\pgfpathlineto{\pgfqpoint{5.537552in}{1.276164in}}%
\pgfpathlineto{\pgfqpoint{5.542908in}{1.254136in}}%
\pgfpathlineto{\pgfqpoint{5.548264in}{1.257912in}}%
\pgfpathlineto{\pgfqpoint{5.553620in}{1.264835in}}%
\pgfpathlineto{\pgfqpoint{5.558976in}{1.264835in}}%
\pgfpathlineto{\pgfqpoint{5.564331in}{1.272388in}}%
\pgfpathlineto{\pgfqpoint{5.569687in}{1.264835in}}%
\pgfpathlineto{\pgfqpoint{5.575043in}{1.250360in}}%
\pgfpathlineto{\pgfqpoint{5.580399in}{1.246583in}}%
\pgfpathlineto{\pgfqpoint{5.585755in}{1.254136in}}%
\pgfpathlineto{\pgfqpoint{5.591111in}{1.254136in}}%
\pgfpathlineto{\pgfqpoint{5.596467in}{1.242807in}}%
\pgfpathlineto{\pgfqpoint{5.601823in}{1.217002in}}%
\pgfpathlineto{\pgfqpoint{5.607179in}{1.213226in}}%
\pgfpathlineto{\pgfqpoint{5.612535in}{1.217002in}}%
\pgfpathlineto{\pgfqpoint{5.617891in}{1.246583in}}%
\pgfpathlineto{\pgfqpoint{5.623246in}{1.209450in}}%
\pgfpathlineto{\pgfqpoint{5.628602in}{1.217002in}}%
\pgfpathlineto{\pgfqpoint{5.633958in}{1.202527in}}%
\pgfpathlineto{\pgfqpoint{5.644670in}{1.242807in}}%
\pgfpathlineto{\pgfqpoint{5.655382in}{1.191198in}}%
\pgfpathlineto{\pgfqpoint{5.660738in}{1.246583in}}%
\pgfpathlineto{\pgfqpoint{5.666094in}{1.191198in}}%
\pgfpathlineto{\pgfqpoint{5.671450in}{1.213226in}}%
\pgfpathlineto{\pgfqpoint{5.676806in}{1.209450in}}%
\pgfpathlineto{\pgfqpoint{5.682161in}{1.176722in}}%
\pgfpathlineto{\pgfqpoint{5.687517in}{1.198750in}}%
\pgfpathlineto{\pgfqpoint{5.692873in}{1.172946in}}%
\pgfpathlineto{\pgfqpoint{5.698229in}{1.220779in}}%
\pgfpathlineto{\pgfqpoint{5.703585in}{1.194974in}}%
\pgfpathlineto{\pgfqpoint{5.708941in}{1.224555in}}%
\pgfpathlineto{\pgfqpoint{5.714297in}{1.217002in}}%
\pgfpathlineto{\pgfqpoint{5.719653in}{1.176722in}}%
\pgfpathlineto{\pgfqpoint{5.730365in}{1.217002in}}%
\pgfpathlineto{\pgfqpoint{5.735721in}{1.198750in}}%
\pgfpathlineto{\pgfqpoint{5.741076in}{1.231478in}}%
\pgfpathlineto{\pgfqpoint{5.746432in}{1.205674in}}%
\pgfpathlineto{\pgfqpoint{5.751788in}{1.191198in}}%
\pgfpathlineto{\pgfqpoint{5.757144in}{1.235254in}}%
\pgfpathlineto{\pgfqpoint{5.762500in}{1.194974in}}%
\pgfpathlineto{\pgfqpoint{5.767856in}{1.220779in}}%
\pgfpathlineto{\pgfqpoint{5.773212in}{1.209450in}}%
\pgfpathlineto{\pgfqpoint{5.778568in}{1.205674in}}%
\pgfpathlineto{\pgfqpoint{5.783924in}{1.228331in}}%
\pgfpathlineto{\pgfqpoint{5.789280in}{1.209450in}}%
\pgfpathlineto{\pgfqpoint{5.794636in}{1.213226in}}%
\pgfpathlineto{\pgfqpoint{5.799991in}{1.198750in}}%
\pgfpathlineto{\pgfqpoint{5.805347in}{1.202527in}}%
\pgfpathlineto{\pgfqpoint{5.810703in}{1.194974in}}%
\pgfpathlineto{\pgfqpoint{5.816059in}{1.183645in}}%
\pgfpathlineto{\pgfqpoint{5.821415in}{1.235254in}}%
\pgfpathlineto{\pgfqpoint{5.826771in}{1.198750in}}%
\pgfpathlineto{\pgfqpoint{5.832127in}{1.194974in}}%
\pgfpathlineto{\pgfqpoint{5.837483in}{1.209450in}}%
\pgfpathlineto{\pgfqpoint{5.842839in}{1.198750in}}%
\pgfpathlineto{\pgfqpoint{5.848195in}{1.213226in}}%
\pgfpathlineto{\pgfqpoint{5.853551in}{1.191198in}}%
\pgfpathlineto{\pgfqpoint{5.858906in}{1.187422in}}%
\pgfpathlineto{\pgfqpoint{5.864262in}{1.172946in}}%
\pgfpathlineto{\pgfqpoint{5.869618in}{1.154065in}}%
\pgfpathlineto{\pgfqpoint{5.874974in}{1.172946in}}%
\pgfpathlineto{\pgfqpoint{5.880330in}{1.139589in}}%
\pgfpathlineto{\pgfqpoint{5.885686in}{1.157841in}}%
\pgfpathlineto{\pgfqpoint{5.891042in}{1.154065in}}%
\pgfpathlineto{\pgfqpoint{5.896398in}{1.183645in}}%
\pgfpathlineto{\pgfqpoint{5.901754in}{1.205674in}}%
\pgfpathlineto{\pgfqpoint{5.901754in}{1.205674in}}%
\pgfusepath{stroke}%
\end{pgfscope}%
\begin{pgfscope}%
\pgfpathrectangle{\pgfqpoint{0.688581in}{0.565123in}}{\pgfqpoint{5.461419in}{2.585803in}}%
\pgfusepath{clip}%
\pgfsetrectcap%
\pgfsetroundjoin%
\pgfsetlinewidth{2.007500pt}%
\definecolor{currentstroke}{rgb}{1.000000,0.498039,0.054902}%
\pgfsetstrokecolor{currentstroke}%
\pgfsetdash{}{0pt}%
\pgfpathmoveto{\pgfqpoint{0.936827in}{3.014508in}}%
\pgfpathlineto{\pgfqpoint{0.942183in}{3.033390in}}%
\pgfpathlineto{\pgfqpoint{0.947539in}{3.033390in}}%
\pgfpathlineto{\pgfqpoint{0.952895in}{3.029613in}}%
\pgfpathlineto{\pgfqpoint{0.963607in}{3.029613in}}%
\pgfpathlineto{\pgfqpoint{0.968963in}{3.024578in}}%
\pgfpathlineto{\pgfqpoint{0.974319in}{3.022061in}}%
\pgfpathlineto{\pgfqpoint{0.979675in}{3.018285in}}%
\pgfpathlineto{\pgfqpoint{0.985030in}{3.018285in}}%
\pgfpathlineto{\pgfqpoint{0.990386in}{3.014508in}}%
\pgfpathlineto{\pgfqpoint{0.995742in}{3.014508in}}%
\pgfpathlineto{\pgfqpoint{1.001098in}{3.010732in}}%
\pgfpathlineto{\pgfqpoint{1.006454in}{3.010732in}}%
\pgfpathlineto{\pgfqpoint{1.011810in}{3.007585in}}%
\pgfpathlineto{\pgfqpoint{1.017166in}{3.007585in}}%
\pgfpathlineto{\pgfqpoint{1.022522in}{3.000033in}}%
\pgfpathlineto{\pgfqpoint{1.027878in}{3.000033in}}%
\pgfpathlineto{\pgfqpoint{1.033234in}{2.996256in}}%
\pgfpathlineto{\pgfqpoint{1.038590in}{2.996256in}}%
\pgfpathlineto{\pgfqpoint{1.043945in}{2.992480in}}%
\pgfpathlineto{\pgfqpoint{1.049301in}{2.992480in}}%
\pgfpathlineto{\pgfqpoint{1.054657in}{2.988704in}}%
\pgfpathlineto{\pgfqpoint{1.060013in}{2.988704in}}%
\pgfpathlineto{\pgfqpoint{1.065369in}{2.984928in}}%
\pgfpathlineto{\pgfqpoint{1.070725in}{2.984928in}}%
\pgfpathlineto{\pgfqpoint{1.076081in}{2.981781in}}%
\pgfpathlineto{\pgfqpoint{1.081437in}{2.981781in}}%
\pgfpathlineto{\pgfqpoint{1.086793in}{2.978004in}}%
\pgfpathlineto{\pgfqpoint{1.092149in}{2.978004in}}%
\pgfpathlineto{\pgfqpoint{1.097505in}{2.974228in}}%
\pgfpathlineto{\pgfqpoint{1.102860in}{2.974228in}}%
\pgfpathlineto{\pgfqpoint{1.124284in}{2.959123in}}%
\pgfpathlineto{\pgfqpoint{1.129640in}{2.959123in}}%
\pgfpathlineto{\pgfqpoint{1.140352in}{2.952200in}}%
\pgfpathlineto{\pgfqpoint{1.145708in}{2.952200in}}%
\pgfpathlineto{\pgfqpoint{1.172487in}{2.933319in}}%
\pgfpathlineto{\pgfqpoint{1.177843in}{2.933319in}}%
\pgfpathlineto{\pgfqpoint{1.193911in}{2.922619in}}%
\pgfpathlineto{\pgfqpoint{1.226046in}{2.899332in}}%
\pgfpathlineto{\pgfqpoint{1.231402in}{2.889262in}}%
\pgfpathlineto{\pgfqpoint{1.252826in}{2.874786in}}%
\pgfpathlineto{\pgfqpoint{1.258182in}{2.867234in}}%
\pgfpathlineto{\pgfqpoint{1.263538in}{2.863457in}}%
\pgfpathlineto{\pgfqpoint{1.268894in}{2.856534in}}%
\pgfpathlineto{\pgfqpoint{1.274249in}{2.852758in}}%
\pgfpathlineto{\pgfqpoint{1.311741in}{2.801149in}}%
\pgfpathlineto{\pgfqpoint{1.317097in}{2.789820in}}%
\pgfpathlineto{\pgfqpoint{1.322453in}{2.782268in}}%
\pgfpathlineto{\pgfqpoint{1.333164in}{2.760239in}}%
\pgfpathlineto{\pgfqpoint{1.338520in}{2.753316in}}%
\pgfpathlineto{\pgfqpoint{1.354588in}{2.717442in}}%
\pgfpathlineto{\pgfqpoint{1.359944in}{2.701707in}}%
\pgfpathlineto{\pgfqpoint{1.365300in}{2.690378in}}%
\pgfpathlineto{\pgfqpoint{1.397435in}{2.602265in}}%
\pgfpathlineto{\pgfqpoint{1.424215in}{2.505970in}}%
\pgfpathlineto{\pgfqpoint{1.440283in}{2.436109in}}%
\pgfpathlineto{\pgfqpoint{1.467062in}{2.288835in}}%
\pgfpathlineto{\pgfqpoint{1.483130in}{2.181840in}}%
\pgfpathlineto{\pgfqpoint{1.499198in}{2.052818in}}%
\pgfpathlineto{\pgfqpoint{1.509909in}{1.946453in}}%
\pgfpathlineto{\pgfqpoint{1.520621in}{1.817430in}}%
\pgfpathlineto{\pgfqpoint{1.531333in}{1.644351in}}%
\pgfpathlineto{\pgfqpoint{1.536689in}{1.541133in}}%
\pgfpathlineto{\pgfqpoint{1.542045in}{1.415886in}}%
\pgfpathlineto{\pgfqpoint{1.547401in}{1.257283in}}%
\pgfpathlineto{\pgfqpoint{1.552757in}{1.050846in}}%
\pgfpathlineto{\pgfqpoint{1.558113in}{0.796577in}}%
\pgfpathlineto{\pgfqpoint{1.563468in}{0.682660in}}%
\pgfpathlineto{\pgfqpoint{1.568824in}{0.841263in}}%
\pgfpathlineto{\pgfqpoint{1.579536in}{1.257283in}}%
\pgfpathlineto{\pgfqpoint{1.590248in}{1.515328in}}%
\pgfpathlineto{\pgfqpoint{1.600960in}{1.688407in}}%
\pgfpathlineto{\pgfqpoint{1.611672in}{1.817430in}}%
\pgfpathlineto{\pgfqpoint{1.622383in}{1.916872in}}%
\pgfpathlineto{\pgfqpoint{1.633095in}{2.001838in}}%
\pgfpathlineto{\pgfqpoint{1.643807in}{2.071699in}}%
\pgfpathlineto{\pgfqpoint{1.665231in}{2.185616in}}%
\pgfpathlineto{\pgfqpoint{1.692010in}{2.288835in}}%
\pgfpathlineto{\pgfqpoint{1.713434in}{2.358696in}}%
\pgfpathlineto{\pgfqpoint{1.729502in}{2.403382in}}%
\pgfpathlineto{\pgfqpoint{1.734858in}{2.414081in}}%
\pgfpathlineto{\pgfqpoint{1.740213in}{2.429186in}}%
\pgfpathlineto{\pgfqpoint{1.783061in}{2.513523in}}%
\pgfpathlineto{\pgfqpoint{1.788417in}{2.521075in}}%
\pgfpathlineto{\pgfqpoint{1.793773in}{2.531775in}}%
\pgfpathlineto{\pgfqpoint{1.799128in}{2.539327in}}%
\pgfpathlineto{\pgfqpoint{1.804484in}{2.550656in}}%
\pgfpathlineto{\pgfqpoint{1.820552in}{2.572684in}}%
\pgfpathlineto{\pgfqpoint{1.847332in}{2.609188in}}%
\pgfpathlineto{\pgfqpoint{1.852687in}{2.612965in}}%
\pgfpathlineto{\pgfqpoint{1.868755in}{2.634993in}}%
\pgfpathlineto{\pgfqpoint{1.874111in}{2.638769in}}%
\pgfpathlineto{\pgfqpoint{1.884823in}{2.653874in}}%
\pgfpathlineto{\pgfqpoint{1.890179in}{2.657021in}}%
\pgfpathlineto{\pgfqpoint{1.895535in}{2.664574in}}%
\pgfpathlineto{\pgfqpoint{1.900891in}{2.665203in}}%
\pgfpathlineto{\pgfqpoint{1.906247in}{2.672126in}}%
\pgfpathlineto{\pgfqpoint{1.911602in}{2.675903in}}%
\pgfpathlineto{\pgfqpoint{1.916958in}{2.682826in}}%
\pgfpathlineto{\pgfqpoint{1.922314in}{2.686602in}}%
\pgfpathlineto{\pgfqpoint{1.927670in}{2.694155in}}%
\pgfpathlineto{\pgfqpoint{1.938382in}{2.701707in}}%
\pgfpathlineto{\pgfqpoint{1.943738in}{2.708630in}}%
\pgfpathlineto{\pgfqpoint{1.954450in}{2.716183in}}%
\pgfpathlineto{\pgfqpoint{1.959806in}{2.723735in}}%
\pgfpathlineto{\pgfqpoint{2.040144in}{2.779121in}}%
\pgfpathlineto{\pgfqpoint{2.056212in}{2.789820in}}%
\pgfpathlineto{\pgfqpoint{2.066924in}{2.797373in}}%
\pgfpathlineto{\pgfqpoint{2.072280in}{2.797373in}}%
\pgfpathlineto{\pgfqpoint{2.082992in}{2.804925in}}%
\pgfpathlineto{\pgfqpoint{2.088347in}{2.804925in}}%
\pgfpathlineto{\pgfqpoint{2.104415in}{2.815625in}}%
\pgfpathlineto{\pgfqpoint{2.109771in}{2.815625in}}%
\pgfpathlineto{\pgfqpoint{2.125839in}{2.826953in}}%
\pgfpathlineto{\pgfqpoint{2.131195in}{2.826953in}}%
\pgfpathlineto{\pgfqpoint{2.141907in}{2.833877in}}%
\pgfpathlineto{\pgfqpoint{2.147262in}{2.833877in}}%
\pgfpathlineto{\pgfqpoint{2.157974in}{2.841429in}}%
\pgfpathlineto{\pgfqpoint{2.163330in}{2.841429in}}%
\pgfpathlineto{\pgfqpoint{2.174042in}{2.848982in}}%
\pgfpathlineto{\pgfqpoint{2.179398in}{2.848982in}}%
\pgfpathlineto{\pgfqpoint{2.190110in}{2.856534in}}%
\pgfpathlineto{\pgfqpoint{2.195466in}{2.856534in}}%
\pgfpathlineto{\pgfqpoint{2.200821in}{2.859681in}}%
\pgfpathlineto{\pgfqpoint{2.206177in}{2.859681in}}%
\pgfpathlineto{\pgfqpoint{2.211533in}{2.863457in}}%
\pgfpathlineto{\pgfqpoint{2.216889in}{2.863457in}}%
\pgfpathlineto{\pgfqpoint{2.222245in}{2.867234in}}%
\pgfpathlineto{\pgfqpoint{2.227601in}{2.867234in}}%
\pgfpathlineto{\pgfqpoint{2.232957in}{2.871010in}}%
\pgfpathlineto{\pgfqpoint{2.238313in}{2.871010in}}%
\pgfpathlineto{\pgfqpoint{2.243669in}{2.874786in}}%
\pgfpathlineto{\pgfqpoint{2.249025in}{2.874786in}}%
\pgfpathlineto{\pgfqpoint{2.254381in}{2.878563in}}%
\pgfpathlineto{\pgfqpoint{2.265092in}{2.878563in}}%
\pgfpathlineto{\pgfqpoint{2.270448in}{2.882339in}}%
\pgfpathlineto{\pgfqpoint{2.275804in}{2.882339in}}%
\pgfpathlineto{\pgfqpoint{2.286516in}{2.889262in}}%
\pgfpathlineto{\pgfqpoint{2.302584in}{2.889262in}}%
\pgfpathlineto{\pgfqpoint{2.307940in}{2.891780in}}%
\pgfpathlineto{\pgfqpoint{2.313296in}{2.891150in}}%
\pgfpathlineto{\pgfqpoint{2.318651in}{2.893038in}}%
\pgfpathlineto{\pgfqpoint{2.324007in}{2.893038in}}%
\pgfpathlineto{\pgfqpoint{2.329363in}{2.896815in}}%
\pgfpathlineto{\pgfqpoint{2.340075in}{2.896815in}}%
\pgfpathlineto{\pgfqpoint{2.350787in}{2.903738in}}%
\pgfpathlineto{\pgfqpoint{2.356143in}{2.904367in}}%
\pgfpathlineto{\pgfqpoint{2.361499in}{2.908143in}}%
\pgfpathlineto{\pgfqpoint{2.372211in}{2.907514in}}%
\pgfpathlineto{\pgfqpoint{2.377566in}{2.904996in}}%
\pgfpathlineto{\pgfqpoint{2.382922in}{2.911290in}}%
\pgfpathlineto{\pgfqpoint{2.409702in}{2.911290in}}%
\pgfpathlineto{\pgfqpoint{2.415058in}{2.915067in}}%
\pgfpathlineto{\pgfqpoint{2.436481in}{2.915067in}}%
\pgfpathlineto{\pgfqpoint{2.441837in}{2.918843in}}%
\pgfpathlineto{\pgfqpoint{2.473973in}{2.918843in}}%
\pgfpathlineto{\pgfqpoint{2.479329in}{2.922619in}}%
\pgfpathlineto{\pgfqpoint{2.495396in}{2.922619in}}%
\pgfpathlineto{\pgfqpoint{2.506108in}{2.926395in}}%
\pgfpathlineto{\pgfqpoint{2.522176in}{2.926395in}}%
\pgfpathlineto{\pgfqpoint{2.527532in}{2.930172in}}%
\pgfpathlineto{\pgfqpoint{2.586447in}{2.930172in}}%
\pgfpathlineto{\pgfqpoint{2.597159in}{2.933319in}}%
\pgfpathlineto{\pgfqpoint{2.613226in}{2.933319in}}%
\pgfpathlineto{\pgfqpoint{2.618582in}{2.937095in}}%
\pgfpathlineto{\pgfqpoint{2.656074in}{2.937095in}}%
\pgfpathlineto{\pgfqpoint{2.666785in}{2.940871in}}%
\pgfpathlineto{\pgfqpoint{2.704277in}{2.940871in}}%
\pgfpathlineto{\pgfqpoint{2.709633in}{2.937095in}}%
\pgfpathlineto{\pgfqpoint{2.725700in}{2.937724in}}%
\pgfpathlineto{\pgfqpoint{2.731056in}{2.940871in}}%
\pgfpathlineto{\pgfqpoint{2.784615in}{2.940871in}}%
\pgfpathlineto{\pgfqpoint{2.789971in}{2.944018in}}%
\pgfpathlineto{\pgfqpoint{2.800683in}{2.944647in}}%
\pgfpathlineto{\pgfqpoint{2.811395in}{2.944647in}}%
\pgfpathlineto{\pgfqpoint{2.816751in}{2.940871in}}%
\pgfpathlineto{\pgfqpoint{2.832819in}{2.940871in}}%
\pgfpathlineto{\pgfqpoint{2.838175in}{2.942130in}}%
\pgfpathlineto{\pgfqpoint{2.843530in}{2.944647in}}%
\pgfpathlineto{\pgfqpoint{2.881022in}{2.944647in}}%
\pgfpathlineto{\pgfqpoint{2.886378in}{2.940871in}}%
\pgfpathlineto{\pgfqpoint{2.891734in}{2.940871in}}%
\pgfpathlineto{\pgfqpoint{2.897089in}{2.944647in}}%
\pgfpathlineto{\pgfqpoint{2.939937in}{2.944647in}}%
\pgfpathlineto{\pgfqpoint{2.945293in}{2.947165in}}%
\pgfpathlineto{\pgfqpoint{2.950649in}{2.948424in}}%
\pgfpathlineto{\pgfqpoint{2.956004in}{2.944647in}}%
\pgfpathlineto{\pgfqpoint{2.961360in}{2.948424in}}%
\pgfpathlineto{\pgfqpoint{2.966716in}{2.944647in}}%
\pgfpathlineto{\pgfqpoint{2.982784in}{2.944647in}}%
\pgfpathlineto{\pgfqpoint{2.988140in}{2.948424in}}%
\pgfpathlineto{\pgfqpoint{3.020275in}{2.948424in}}%
\pgfpathlineto{\pgfqpoint{3.030987in}{2.952200in}}%
\pgfpathlineto{\pgfqpoint{3.036343in}{2.948424in}}%
\pgfpathlineto{\pgfqpoint{3.063123in}{2.948424in}}%
\pgfpathlineto{\pgfqpoint{3.068479in}{2.952200in}}%
\pgfpathlineto{\pgfqpoint{3.105970in}{2.952200in}}%
\pgfpathlineto{\pgfqpoint{3.111326in}{2.950312in}}%
\pgfpathlineto{\pgfqpoint{3.116682in}{2.952200in}}%
\pgfpathlineto{\pgfqpoint{3.164885in}{2.952200in}}%
\pgfpathlineto{\pgfqpoint{3.170241in}{2.955976in}}%
\pgfpathlineto{\pgfqpoint{3.197020in}{2.955976in}}%
\pgfpathlineto{\pgfqpoint{3.202376in}{2.952200in}}%
\pgfpathlineto{\pgfqpoint{3.213088in}{2.955976in}}%
\pgfpathlineto{\pgfqpoint{3.245223in}{2.955976in}}%
\pgfpathlineto{\pgfqpoint{3.250579in}{2.959123in}}%
\pgfpathlineto{\pgfqpoint{3.282715in}{2.959123in}}%
\pgfpathlineto{\pgfqpoint{3.293427in}{2.962899in}}%
\pgfpathlineto{\pgfqpoint{3.325562in}{2.962899in}}%
\pgfpathlineto{\pgfqpoint{3.330918in}{2.966676in}}%
\pgfpathlineto{\pgfqpoint{3.363053in}{2.966676in}}%
\pgfpathlineto{\pgfqpoint{3.368409in}{2.967934in}}%
\pgfpathlineto{\pgfqpoint{3.373765in}{2.970452in}}%
\pgfpathlineto{\pgfqpoint{3.405901in}{2.970452in}}%
\pgfpathlineto{\pgfqpoint{3.411257in}{2.966676in}}%
\pgfpathlineto{\pgfqpoint{3.416613in}{2.969822in}}%
\pgfpathlineto{\pgfqpoint{3.427324in}{2.970452in}}%
\pgfpathlineto{\pgfqpoint{3.438036in}{2.970452in}}%
\pgfpathlineto{\pgfqpoint{3.443392in}{2.974228in}}%
\pgfpathlineto{\pgfqpoint{3.454104in}{2.974228in}}%
\pgfpathlineto{\pgfqpoint{3.459460in}{2.975487in}}%
\pgfpathlineto{\pgfqpoint{3.464816in}{2.978004in}}%
\pgfpathlineto{\pgfqpoint{3.507663in}{2.978004in}}%
\pgfpathlineto{\pgfqpoint{3.513019in}{2.981151in}}%
\pgfpathlineto{\pgfqpoint{3.523731in}{2.981781in}}%
\pgfpathlineto{\pgfqpoint{3.571934in}{2.981781in}}%
\pgfpathlineto{\pgfqpoint{3.577290in}{2.984298in}}%
\pgfpathlineto{\pgfqpoint{3.588002in}{2.984928in}}%
\pgfpathlineto{\pgfqpoint{3.636205in}{2.984928in}}%
\pgfpathlineto{\pgfqpoint{3.641561in}{2.988704in}}%
\pgfpathlineto{\pgfqpoint{3.657628in}{2.988704in}}%
\pgfpathlineto{\pgfqpoint{3.662984in}{2.991221in}}%
\pgfpathlineto{\pgfqpoint{3.668340in}{2.992480in}}%
\pgfpathlineto{\pgfqpoint{3.737967in}{2.992480in}}%
\pgfpathlineto{\pgfqpoint{3.743323in}{2.996256in}}%
\pgfpathlineto{\pgfqpoint{3.770102in}{2.996256in}}%
\pgfpathlineto{\pgfqpoint{3.775458in}{2.999403in}}%
\pgfpathlineto{\pgfqpoint{3.786170in}{3.000033in}}%
\pgfpathlineto{\pgfqpoint{3.834373in}{3.000033in}}%
\pgfpathlineto{\pgfqpoint{3.839729in}{3.003809in}}%
\pgfpathlineto{\pgfqpoint{3.914712in}{3.003809in}}%
\pgfpathlineto{\pgfqpoint{3.920068in}{3.007585in}}%
\pgfpathlineto{\pgfqpoint{3.968271in}{3.007585in}}%
\pgfpathlineto{\pgfqpoint{3.973627in}{3.010732in}}%
\pgfpathlineto{\pgfqpoint{4.027186in}{3.010732in}}%
\pgfpathlineto{\pgfqpoint{4.032542in}{3.014508in}}%
\pgfpathlineto{\pgfqpoint{4.043254in}{3.014508in}}%
\pgfpathlineto{\pgfqpoint{4.048610in}{3.010732in}}%
\pgfpathlineto{\pgfqpoint{4.053966in}{3.014508in}}%
\pgfpathlineto{\pgfqpoint{4.059321in}{3.010732in}}%
\pgfpathlineto{\pgfqpoint{4.091457in}{3.010732in}}%
\pgfpathlineto{\pgfqpoint{4.096813in}{3.014508in}}%
\pgfpathlineto{\pgfqpoint{4.145016in}{3.014508in}}%
\pgfpathlineto{\pgfqpoint{4.150372in}{3.018285in}}%
\pgfpathlineto{\pgfqpoint{4.198575in}{3.018285in}}%
\pgfpathlineto{\pgfqpoint{4.203931in}{3.019543in}}%
\pgfpathlineto{\pgfqpoint{4.209287in}{3.018285in}}%
\pgfpathlineto{\pgfqpoint{4.246778in}{3.018285in}}%
\pgfpathlineto{\pgfqpoint{4.252134in}{3.022061in}}%
\pgfpathlineto{\pgfqpoint{4.284270in}{3.022061in}}%
\pgfpathlineto{\pgfqpoint{4.289625in}{3.019543in}}%
\pgfpathlineto{\pgfqpoint{4.294981in}{3.018285in}}%
\pgfpathlineto{\pgfqpoint{4.300337in}{3.022061in}}%
\pgfpathlineto{\pgfqpoint{4.348540in}{3.022061in}}%
\pgfpathlineto{\pgfqpoint{4.353896in}{3.023949in}}%
\pgfpathlineto{\pgfqpoint{4.364608in}{3.023949in}}%
\pgfpathlineto{\pgfqpoint{4.369964in}{3.022061in}}%
\pgfpathlineto{\pgfqpoint{4.402100in}{3.022061in}}%
\pgfpathlineto{\pgfqpoint{4.407455in}{3.025837in}}%
\pgfpathlineto{\pgfqpoint{4.423523in}{3.025837in}}%
\pgfpathlineto{\pgfqpoint{4.428879in}{3.024578in}}%
\pgfpathlineto{\pgfqpoint{4.434235in}{3.022061in}}%
\pgfpathlineto{\pgfqpoint{4.450303in}{3.022061in}}%
\pgfpathlineto{\pgfqpoint{4.455659in}{3.025837in}}%
\pgfpathlineto{\pgfqpoint{4.482438in}{3.025837in}}%
\pgfpathlineto{\pgfqpoint{4.487794in}{3.022061in}}%
\pgfpathlineto{\pgfqpoint{4.509218in}{3.022061in}}%
\pgfpathlineto{\pgfqpoint{4.514574in}{3.025837in}}%
\pgfpathlineto{\pgfqpoint{4.530641in}{3.025837in}}%
\pgfpathlineto{\pgfqpoint{4.541353in}{3.022061in}}%
\pgfpathlineto{\pgfqpoint{4.557421in}{3.022061in}}%
\pgfpathlineto{\pgfqpoint{4.562777in}{3.025837in}}%
\pgfpathlineto{\pgfqpoint{4.589556in}{3.025208in}}%
\pgfpathlineto{\pgfqpoint{4.594912in}{3.022061in}}%
\pgfpathlineto{\pgfqpoint{4.621692in}{3.022061in}}%
\pgfpathlineto{\pgfqpoint{4.627048in}{3.025837in}}%
\pgfpathlineto{\pgfqpoint{4.637759in}{3.025837in}}%
\pgfpathlineto{\pgfqpoint{4.643115in}{3.023320in}}%
\pgfpathlineto{\pgfqpoint{4.648471in}{3.022061in}}%
\pgfpathlineto{\pgfqpoint{4.675251in}{3.022061in}}%
\pgfpathlineto{\pgfqpoint{4.680607in}{3.025837in}}%
\pgfpathlineto{\pgfqpoint{4.685963in}{3.025208in}}%
\pgfpathlineto{\pgfqpoint{4.691319in}{3.022061in}}%
\pgfpathlineto{\pgfqpoint{4.728810in}{3.022061in}}%
\pgfpathlineto{\pgfqpoint{4.734166in}{3.019543in}}%
\pgfpathlineto{\pgfqpoint{4.739522in}{3.018285in}}%
\pgfpathlineto{\pgfqpoint{4.771657in}{3.018285in}}%
\pgfpathlineto{\pgfqpoint{4.777013in}{3.022061in}}%
\pgfpathlineto{\pgfqpoint{4.782369in}{3.022061in}}%
\pgfpathlineto{\pgfqpoint{4.787725in}{3.020802in}}%
\pgfpathlineto{\pgfqpoint{4.793081in}{3.018285in}}%
\pgfpathlineto{\pgfqpoint{4.830572in}{3.018285in}}%
\pgfpathlineto{\pgfqpoint{4.835928in}{3.015138in}}%
\pgfpathlineto{\pgfqpoint{4.841284in}{3.018285in}}%
\pgfpathlineto{\pgfqpoint{4.884131in}{3.018285in}}%
\pgfpathlineto{\pgfqpoint{4.894843in}{3.014508in}}%
\pgfpathlineto{\pgfqpoint{4.905555in}{3.014508in}}%
\pgfpathlineto{\pgfqpoint{4.910911in}{3.018285in}}%
\pgfpathlineto{\pgfqpoint{4.916267in}{3.014508in}}%
\pgfpathlineto{\pgfqpoint{4.932334in}{3.014508in}}%
\pgfpathlineto{\pgfqpoint{4.937690in}{3.010732in}}%
\pgfpathlineto{\pgfqpoint{4.985893in}{3.010732in}}%
\pgfpathlineto{\pgfqpoint{4.996605in}{3.007585in}}%
\pgfpathlineto{\pgfqpoint{5.039453in}{3.006956in}}%
\pgfpathlineto{\pgfqpoint{5.044808in}{3.003809in}}%
\pgfpathlineto{\pgfqpoint{5.093012in}{3.003809in}}%
\pgfpathlineto{\pgfqpoint{5.098368in}{3.000662in}}%
\pgfpathlineto{\pgfqpoint{5.109079in}{3.000033in}}%
\pgfpathlineto{\pgfqpoint{5.178706in}{3.000033in}}%
\pgfpathlineto{\pgfqpoint{5.184062in}{2.996886in}}%
\pgfpathlineto{\pgfqpoint{5.194774in}{2.996256in}}%
\pgfpathlineto{\pgfqpoint{5.200130in}{2.996256in}}%
\pgfpathlineto{\pgfqpoint{5.205486in}{2.992480in}}%
\pgfpathlineto{\pgfqpoint{5.226909in}{2.992480in}}%
\pgfpathlineto{\pgfqpoint{5.232265in}{2.991221in}}%
\pgfpathlineto{\pgfqpoint{5.237621in}{2.988704in}}%
\pgfpathlineto{\pgfqpoint{5.285824in}{2.988704in}}%
\pgfpathlineto{\pgfqpoint{5.291180in}{2.984928in}}%
\pgfpathlineto{\pgfqpoint{5.317960in}{2.984298in}}%
\pgfpathlineto{\pgfqpoint{5.323316in}{2.982410in}}%
\pgfpathlineto{\pgfqpoint{5.328672in}{2.984928in}}%
\pgfpathlineto{\pgfqpoint{5.339383in}{2.984928in}}%
\pgfpathlineto{\pgfqpoint{5.344739in}{2.981781in}}%
\pgfpathlineto{\pgfqpoint{5.392942in}{2.981781in}}%
\pgfpathlineto{\pgfqpoint{5.403654in}{2.974228in}}%
\pgfpathlineto{\pgfqpoint{5.414366in}{2.976746in}}%
\pgfpathlineto{\pgfqpoint{5.419722in}{2.975487in}}%
\pgfpathlineto{\pgfqpoint{5.425078in}{2.978004in}}%
\pgfpathlineto{\pgfqpoint{5.430434in}{2.978004in}}%
\pgfpathlineto{\pgfqpoint{5.441146in}{2.971711in}}%
\pgfpathlineto{\pgfqpoint{5.451857in}{2.974228in}}%
\pgfpathlineto{\pgfqpoint{5.467925in}{2.974228in}}%
\pgfpathlineto{\pgfqpoint{5.473281in}{2.970452in}}%
\pgfpathlineto{\pgfqpoint{5.510772in}{2.970452in}}%
\pgfpathlineto{\pgfqpoint{5.516128in}{2.966676in}}%
\pgfpathlineto{\pgfqpoint{5.548264in}{2.966676in}}%
\pgfpathlineto{\pgfqpoint{5.553620in}{2.962899in}}%
\pgfpathlineto{\pgfqpoint{5.580399in}{2.962270in}}%
\pgfpathlineto{\pgfqpoint{5.585755in}{2.960382in}}%
\pgfpathlineto{\pgfqpoint{5.591111in}{2.962899in}}%
\pgfpathlineto{\pgfqpoint{5.607179in}{2.962899in}}%
\pgfpathlineto{\pgfqpoint{5.612535in}{2.959123in}}%
\pgfpathlineto{\pgfqpoint{5.639314in}{2.959123in}}%
\pgfpathlineto{\pgfqpoint{5.644670in}{2.955976in}}%
\pgfpathlineto{\pgfqpoint{5.650026in}{2.959123in}}%
\pgfpathlineto{\pgfqpoint{5.666094in}{2.959123in}}%
\pgfpathlineto{\pgfqpoint{5.671450in}{2.955976in}}%
\pgfpathlineto{\pgfqpoint{5.708941in}{2.955976in}}%
\pgfpathlineto{\pgfqpoint{5.714297in}{2.952200in}}%
\pgfpathlineto{\pgfqpoint{5.719653in}{2.954717in}}%
\pgfpathlineto{\pgfqpoint{5.725009in}{2.955976in}}%
\pgfpathlineto{\pgfqpoint{5.805347in}{2.955976in}}%
\pgfpathlineto{\pgfqpoint{5.810703in}{2.958494in}}%
\pgfpathlineto{\pgfqpoint{5.816059in}{2.955976in}}%
\pgfpathlineto{\pgfqpoint{5.853551in}{2.956606in}}%
\pgfpathlineto{\pgfqpoint{5.858906in}{2.959123in}}%
\pgfpathlineto{\pgfqpoint{5.874974in}{2.959123in}}%
\pgfpathlineto{\pgfqpoint{5.880330in}{2.955976in}}%
\pgfpathlineto{\pgfqpoint{5.891042in}{2.955976in}}%
\pgfpathlineto{\pgfqpoint{5.896398in}{2.958494in}}%
\pgfpathlineto{\pgfqpoint{5.901754in}{2.959123in}}%
\pgfpathlineto{\pgfqpoint{5.901754in}{2.959123in}}%
\pgfusepath{stroke}%
\end{pgfscope}%
\begin{pgfscope}%
\pgfpathrectangle{\pgfqpoint{0.688581in}{0.565123in}}{\pgfqpoint{5.461419in}{2.585803in}}%
\pgfusepath{clip}%
\pgfsetbuttcap%
\pgfsetroundjoin%
\pgfsetlinewidth{1.505625pt}%
\definecolor{currentstroke}{rgb}{1.000000,0.000000,0.000000}%
\pgfsetstrokecolor{currentstroke}%
\pgfsetdash{{5.550000pt}{2.400000pt}}{0.000000pt}%
\pgfpathmoveto{\pgfqpoint{1.568824in}{0.565123in}}%
\pgfpathlineto{\pgfqpoint{1.568824in}{3.150926in}}%
\pgfusepath{stroke}%
\end{pgfscope}%
\begin{pgfscope}%
\pgfsetrectcap%
\pgfsetmiterjoin%
\pgfsetlinewidth{0.803000pt}%
\definecolor{currentstroke}{rgb}{0.000000,0.000000,0.000000}%
\pgfsetstrokecolor{currentstroke}%
\pgfsetdash{}{0pt}%
\pgfpathmoveto{\pgfqpoint{0.688581in}{0.565123in}}%
\pgfpathlineto{\pgfqpoint{0.688581in}{3.150926in}}%
\pgfusepath{stroke}%
\end{pgfscope}%
\begin{pgfscope}%
\pgfsetrectcap%
\pgfsetmiterjoin%
\pgfsetlinewidth{0.803000pt}%
\definecolor{currentstroke}{rgb}{0.000000,0.000000,0.000000}%
\pgfsetstrokecolor{currentstroke}%
\pgfsetdash{}{0pt}%
\pgfpathmoveto{\pgfqpoint{6.150000in}{0.565123in}}%
\pgfpathlineto{\pgfqpoint{6.150000in}{3.150926in}}%
\pgfusepath{stroke}%
\end{pgfscope}%
\begin{pgfscope}%
\pgfsetrectcap%
\pgfsetmiterjoin%
\pgfsetlinewidth{0.803000pt}%
\definecolor{currentstroke}{rgb}{0.000000,0.000000,0.000000}%
\pgfsetstrokecolor{currentstroke}%
\pgfsetdash{}{0pt}%
\pgfpathmoveto{\pgfqpoint{0.688581in}{0.565123in}}%
\pgfpathlineto{\pgfqpoint{6.150000in}{0.565123in}}%
\pgfusepath{stroke}%
\end{pgfscope}%
\begin{pgfscope}%
\pgfsetrectcap%
\pgfsetmiterjoin%
\pgfsetlinewidth{0.803000pt}%
\definecolor{currentstroke}{rgb}{0.000000,0.000000,0.000000}%
\pgfsetstrokecolor{currentstroke}%
\pgfsetdash{}{0pt}%
\pgfpathmoveto{\pgfqpoint{0.688581in}{3.150926in}}%
\pgfpathlineto{\pgfqpoint{6.150000in}{3.150926in}}%
\pgfusepath{stroke}%
\end{pgfscope}%
\begin{pgfscope}%
\definecolor{textcolor}{rgb}{0.000000,0.000000,0.000000}%
\pgfsetstrokecolor{textcolor}%
\pgfsetfillcolor{textcolor}%
\pgftext[x=3.419290in,y=3.234260in,,base]{\color{textcolor}\rmfamily\fontsize{12.000000}{14.400000}\selectfont LC Band-Pass Filter Transmission Loss}%
\end{pgfscope}%
\begin{pgfscope}%
\pgfsetbuttcap%
\pgfsetmiterjoin%
\definecolor{currentfill}{rgb}{1.000000,1.000000,1.000000}%
\pgfsetfillcolor{currentfill}%
\pgfsetfillopacity{0.800000}%
\pgfsetlinewidth{1.003750pt}%
\definecolor{currentstroke}{rgb}{0.800000,0.800000,0.800000}%
\pgfsetstrokecolor{currentstroke}%
\pgfsetstrokeopacity{0.800000}%
\pgfsetdash{}{0pt}%
\pgfpathmoveto{\pgfqpoint{2.479281in}{0.634568in}}%
\pgfpathlineto{\pgfqpoint{4.359300in}{0.634568in}}%
\pgfpathquadraticcurveto{\pgfqpoint{4.387077in}{0.634568in}}{\pgfqpoint{4.387077in}{0.662346in}}%
\pgfpathlineto{\pgfqpoint{4.387077in}{1.258796in}}%
\pgfpathquadraticcurveto{\pgfqpoint{4.387077in}{1.286574in}}{\pgfqpoint{4.359300in}{1.286574in}}%
\pgfpathlineto{\pgfqpoint{2.479281in}{1.286574in}}%
\pgfpathquadraticcurveto{\pgfqpoint{2.451504in}{1.286574in}}{\pgfqpoint{2.451504in}{1.258796in}}%
\pgfpathlineto{\pgfqpoint{2.451504in}{0.662346in}}%
\pgfpathquadraticcurveto{\pgfqpoint{2.451504in}{0.634568in}}{\pgfqpoint{2.479281in}{0.634568in}}%
\pgfpathlineto{\pgfqpoint{2.479281in}{0.634568in}}%
\pgfpathclose%
\pgfusepath{stroke,fill}%
\end{pgfscope}%
\begin{pgfscope}%
\pgfsetrectcap%
\pgfsetroundjoin%
\pgfsetlinewidth{2.007500pt}%
\definecolor{currentstroke}{rgb}{0.121569,0.466667,0.705882}%
\pgfsetstrokecolor{currentstroke}%
\pgfsetdash{}{0pt}%
\pgfpathmoveto{\pgfqpoint{2.507059in}{1.175463in}}%
\pgfpathlineto{\pgfqpoint{2.645948in}{1.175463in}}%
\pgfpathlineto{\pgfqpoint{2.784837in}{1.175463in}}%
\pgfusepath{stroke}%
\end{pgfscope}%
\begin{pgfscope}%
\definecolor{textcolor}{rgb}{0.000000,0.000000,0.000000}%
\pgfsetstrokecolor{textcolor}%
\pgfsetfillcolor{textcolor}%
\pgftext[x=2.895948in,y=1.126852in,left,base]{\color{textcolor}\rmfamily\fontsize{10.000000}{12.000000}\selectfont Transmission Loss (dB)}%
\end{pgfscope}%
\begin{pgfscope}%
\pgfsetrectcap%
\pgfsetroundjoin%
\pgfsetlinewidth{2.007500pt}%
\definecolor{currentstroke}{rgb}{1.000000,0.498039,0.054902}%
\pgfsetstrokecolor{currentstroke}%
\pgfsetdash{}{0pt}%
\pgfpathmoveto{\pgfqpoint{2.507059in}{0.967129in}}%
\pgfpathlineto{\pgfqpoint{2.645948in}{0.967129in}}%
\pgfpathlineto{\pgfqpoint{2.784837in}{0.967129in}}%
\pgfusepath{stroke}%
\end{pgfscope}%
\begin{pgfscope}%
\definecolor{textcolor}{rgb}{0.000000,0.000000,0.000000}%
\pgfsetstrokecolor{textcolor}%
\pgfsetfillcolor{textcolor}%
\pgftext[x=2.895948in,y=0.918518in,left,base]{\color{textcolor}\rmfamily\fontsize{10.000000}{12.000000}\selectfont Transmission Loss (dB)}%
\end{pgfscope}%
\begin{pgfscope}%
\pgfsetbuttcap%
\pgfsetroundjoin%
\pgfsetlinewidth{1.505625pt}%
\definecolor{currentstroke}{rgb}{1.000000,0.000000,0.000000}%
\pgfsetstrokecolor{currentstroke}%
\pgfsetdash{{5.550000pt}{2.400000pt}}{0.000000pt}%
\pgfpathmoveto{\pgfqpoint{2.507059in}{0.765741in}}%
\pgfpathlineto{\pgfqpoint{2.645948in}{0.765741in}}%
\pgfpathlineto{\pgfqpoint{2.784837in}{0.765741in}}%
\pgfusepath{stroke}%
\end{pgfscope}%
\begin{pgfscope}%
\definecolor{textcolor}{rgb}{0.000000,0.000000,0.000000}%
\pgfsetstrokecolor{textcolor}%
\pgfsetfillcolor{textcolor}%
\pgftext[x=2.895948in,y=0.717129in,left,base]{\color{textcolor}\rmfamily\fontsize{10.000000}{12.000000}\selectfont Center Frequency}%
\end{pgfscope}%
\end{pgfpicture}%
\makeatother%
\endgroup%

		\caption{Measured return loss and phase of a quartz resonatoras a function of frequency.}
		\label{fig:3513}
	\end{figure}
	\subsubsection{}
			\begin{figure}[H]
		\centering
		%% Creator: Matplotlib, PGF backend
%%
%% To include the figure in your LaTeX document, write
%%   \input{<filename>.pgf}
%%
%% Make sure the required packages are loaded in your preamble
%%   \usepackage{pgf}
%%
%% Also ensure that all the required font packages are loaded; for instance,
%% the lmodern package is sometimes necessary when using math font.
%%   \usepackage{lmodern}
%%
%% Figures using additional raster images can only be included by \input if
%% they are in the same directory as the main LaTeX file. For loading figures
%% from other directories you can use the `import` package
%%   \usepackage{import}
%%
%% and then include the figures with
%%   \import{<path to file>}{<filename>.pgf}
%%
%% Matplotlib used the following preamble
%%
\begingroup%
\makeatletter%
\begin{pgfpicture}%
\pgfpathrectangle{\pgfpointorigin}{\pgfqpoint{6.300000in}{3.500000in}}%
\pgfusepath{use as bounding box, clip}%
\begin{pgfscope}%
\pgfsetbuttcap%
\pgfsetmiterjoin%
\definecolor{currentfill}{rgb}{1.000000,1.000000,1.000000}%
\pgfsetfillcolor{currentfill}%
\pgfsetlinewidth{0.000000pt}%
\definecolor{currentstroke}{rgb}{1.000000,1.000000,1.000000}%
\pgfsetstrokecolor{currentstroke}%
\pgfsetdash{}{0pt}%
\pgfpathmoveto{\pgfqpoint{0.000000in}{0.000000in}}%
\pgfpathlineto{\pgfqpoint{6.300000in}{0.000000in}}%
\pgfpathlineto{\pgfqpoint{6.300000in}{3.500000in}}%
\pgfpathlineto{\pgfqpoint{0.000000in}{3.500000in}}%
\pgfpathlineto{\pgfqpoint{0.000000in}{0.000000in}}%
\pgfpathclose%
\pgfusepath{fill}%
\end{pgfscope}%
\begin{pgfscope}%
\pgfsetbuttcap%
\pgfsetmiterjoin%
\definecolor{currentfill}{rgb}{1.000000,1.000000,1.000000}%
\pgfsetfillcolor{currentfill}%
\pgfsetlinewidth{0.000000pt}%
\definecolor{currentstroke}{rgb}{0.000000,0.000000,0.000000}%
\pgfsetstrokecolor{currentstroke}%
\pgfsetstrokeopacity{0.000000}%
\pgfsetdash{}{0pt}%
\pgfpathmoveto{\pgfqpoint{0.688581in}{0.565123in}}%
\pgfpathlineto{\pgfqpoint{5.991820in}{0.565123in}}%
\pgfpathlineto{\pgfqpoint{5.991820in}{3.150926in}}%
\pgfpathlineto{\pgfqpoint{0.688581in}{3.150926in}}%
\pgfpathlineto{\pgfqpoint{0.688581in}{0.565123in}}%
\pgfpathclose%
\pgfusepath{fill}%
\end{pgfscope}%
\begin{pgfscope}%
\pgfpathrectangle{\pgfqpoint{0.688581in}{0.565123in}}{\pgfqpoint{5.303240in}{2.585803in}}%
\pgfusepath{clip}%
\pgfsetbuttcap%
\pgfsetroundjoin%
\pgfsetlinewidth{0.803000pt}%
\definecolor{currentstroke}{rgb}{0.690196,0.690196,0.690196}%
\pgfsetstrokecolor{currentstroke}%
\pgfsetdash{{0.800000pt}{1.320000pt}}{0.000000pt}%
\pgfpathmoveto{\pgfqpoint{0.688581in}{0.565123in}}%
\pgfpathlineto{\pgfqpoint{0.688581in}{3.150926in}}%
\pgfusepath{stroke}%
\end{pgfscope}%
\begin{pgfscope}%
\pgfsetbuttcap%
\pgfsetroundjoin%
\definecolor{currentfill}{rgb}{0.000000,0.000000,0.000000}%
\pgfsetfillcolor{currentfill}%
\pgfsetlinewidth{0.803000pt}%
\definecolor{currentstroke}{rgb}{0.000000,0.000000,0.000000}%
\pgfsetstrokecolor{currentstroke}%
\pgfsetdash{}{0pt}%
\pgfsys@defobject{currentmarker}{\pgfqpoint{0.000000in}{-0.048611in}}{\pgfqpoint{0.000000in}{0.000000in}}{%
\pgfpathmoveto{\pgfqpoint{0.000000in}{0.000000in}}%
\pgfpathlineto{\pgfqpoint{0.000000in}{-0.048611in}}%
\pgfusepath{stroke,fill}%
}%
\begin{pgfscope}%
\pgfsys@transformshift{0.688581in}{0.565123in}%
\pgfsys@useobject{currentmarker}{}%
\end{pgfscope}%
\end{pgfscope}%
\begin{pgfscope}%
\definecolor{textcolor}{rgb}{0.000000,0.000000,0.000000}%
\pgfsetstrokecolor{textcolor}%
\pgfsetfillcolor{textcolor}%
\pgftext[x=0.688581in,y=0.467901in,,top]{\color{textcolor}\rmfamily\fontsize{10.000000}{12.000000}\selectfont \(\displaystyle {10.00}\)}%
\end{pgfscope}%
\begin{pgfscope}%
\pgfpathrectangle{\pgfqpoint{0.688581in}{0.565123in}}{\pgfqpoint{5.303240in}{2.585803in}}%
\pgfusepath{clip}%
\pgfsetbuttcap%
\pgfsetroundjoin%
\pgfsetlinewidth{0.803000pt}%
\definecolor{currentstroke}{rgb}{0.690196,0.690196,0.690196}%
\pgfsetstrokecolor{currentstroke}%
\pgfsetdash{{0.800000pt}{1.320000pt}}{0.000000pt}%
\pgfpathmoveto{\pgfqpoint{1.351486in}{0.565123in}}%
\pgfpathlineto{\pgfqpoint{1.351486in}{3.150926in}}%
\pgfusepath{stroke}%
\end{pgfscope}%
\begin{pgfscope}%
\pgfsetbuttcap%
\pgfsetroundjoin%
\definecolor{currentfill}{rgb}{0.000000,0.000000,0.000000}%
\pgfsetfillcolor{currentfill}%
\pgfsetlinewidth{0.803000pt}%
\definecolor{currentstroke}{rgb}{0.000000,0.000000,0.000000}%
\pgfsetstrokecolor{currentstroke}%
\pgfsetdash{}{0pt}%
\pgfsys@defobject{currentmarker}{\pgfqpoint{0.000000in}{-0.048611in}}{\pgfqpoint{0.000000in}{0.000000in}}{%
\pgfpathmoveto{\pgfqpoint{0.000000in}{0.000000in}}%
\pgfpathlineto{\pgfqpoint{0.000000in}{-0.048611in}}%
\pgfusepath{stroke,fill}%
}%
\begin{pgfscope}%
\pgfsys@transformshift{1.351486in}{0.565123in}%
\pgfsys@useobject{currentmarker}{}%
\end{pgfscope}%
\end{pgfscope}%
\begin{pgfscope}%
\definecolor{textcolor}{rgb}{0.000000,0.000000,0.000000}%
\pgfsetstrokecolor{textcolor}%
\pgfsetfillcolor{textcolor}%
\pgftext[x=1.351486in,y=0.467901in,,top]{\color{textcolor}\rmfamily\fontsize{10.000000}{12.000000}\selectfont \(\displaystyle {10.25}\)}%
\end{pgfscope}%
\begin{pgfscope}%
\pgfpathrectangle{\pgfqpoint{0.688581in}{0.565123in}}{\pgfqpoint{5.303240in}{2.585803in}}%
\pgfusepath{clip}%
\pgfsetbuttcap%
\pgfsetroundjoin%
\pgfsetlinewidth{0.803000pt}%
\definecolor{currentstroke}{rgb}{0.690196,0.690196,0.690196}%
\pgfsetstrokecolor{currentstroke}%
\pgfsetdash{{0.800000pt}{1.320000pt}}{0.000000pt}%
\pgfpathmoveto{\pgfqpoint{2.014391in}{0.565123in}}%
\pgfpathlineto{\pgfqpoint{2.014391in}{3.150926in}}%
\pgfusepath{stroke}%
\end{pgfscope}%
\begin{pgfscope}%
\pgfsetbuttcap%
\pgfsetroundjoin%
\definecolor{currentfill}{rgb}{0.000000,0.000000,0.000000}%
\pgfsetfillcolor{currentfill}%
\pgfsetlinewidth{0.803000pt}%
\definecolor{currentstroke}{rgb}{0.000000,0.000000,0.000000}%
\pgfsetstrokecolor{currentstroke}%
\pgfsetdash{}{0pt}%
\pgfsys@defobject{currentmarker}{\pgfqpoint{0.000000in}{-0.048611in}}{\pgfqpoint{0.000000in}{0.000000in}}{%
\pgfpathmoveto{\pgfqpoint{0.000000in}{0.000000in}}%
\pgfpathlineto{\pgfqpoint{0.000000in}{-0.048611in}}%
\pgfusepath{stroke,fill}%
}%
\begin{pgfscope}%
\pgfsys@transformshift{2.014391in}{0.565123in}%
\pgfsys@useobject{currentmarker}{}%
\end{pgfscope}%
\end{pgfscope}%
\begin{pgfscope}%
\definecolor{textcolor}{rgb}{0.000000,0.000000,0.000000}%
\pgfsetstrokecolor{textcolor}%
\pgfsetfillcolor{textcolor}%
\pgftext[x=2.014391in,y=0.467901in,,top]{\color{textcolor}\rmfamily\fontsize{10.000000}{12.000000}\selectfont \(\displaystyle {10.50}\)}%
\end{pgfscope}%
\begin{pgfscope}%
\pgfpathrectangle{\pgfqpoint{0.688581in}{0.565123in}}{\pgfqpoint{5.303240in}{2.585803in}}%
\pgfusepath{clip}%
\pgfsetbuttcap%
\pgfsetroundjoin%
\pgfsetlinewidth{0.803000pt}%
\definecolor{currentstroke}{rgb}{0.690196,0.690196,0.690196}%
\pgfsetstrokecolor{currentstroke}%
\pgfsetdash{{0.800000pt}{1.320000pt}}{0.000000pt}%
\pgfpathmoveto{\pgfqpoint{2.677296in}{0.565123in}}%
\pgfpathlineto{\pgfqpoint{2.677296in}{3.150926in}}%
\pgfusepath{stroke}%
\end{pgfscope}%
\begin{pgfscope}%
\pgfsetbuttcap%
\pgfsetroundjoin%
\definecolor{currentfill}{rgb}{0.000000,0.000000,0.000000}%
\pgfsetfillcolor{currentfill}%
\pgfsetlinewidth{0.803000pt}%
\definecolor{currentstroke}{rgb}{0.000000,0.000000,0.000000}%
\pgfsetstrokecolor{currentstroke}%
\pgfsetdash{}{0pt}%
\pgfsys@defobject{currentmarker}{\pgfqpoint{0.000000in}{-0.048611in}}{\pgfqpoint{0.000000in}{0.000000in}}{%
\pgfpathmoveto{\pgfqpoint{0.000000in}{0.000000in}}%
\pgfpathlineto{\pgfqpoint{0.000000in}{-0.048611in}}%
\pgfusepath{stroke,fill}%
}%
\begin{pgfscope}%
\pgfsys@transformshift{2.677296in}{0.565123in}%
\pgfsys@useobject{currentmarker}{}%
\end{pgfscope}%
\end{pgfscope}%
\begin{pgfscope}%
\definecolor{textcolor}{rgb}{0.000000,0.000000,0.000000}%
\pgfsetstrokecolor{textcolor}%
\pgfsetfillcolor{textcolor}%
\pgftext[x=2.677296in,y=0.467901in,,top]{\color{textcolor}\rmfamily\fontsize{10.000000}{12.000000}\selectfont \(\displaystyle {10.75}\)}%
\end{pgfscope}%
\begin{pgfscope}%
\pgfpathrectangle{\pgfqpoint{0.688581in}{0.565123in}}{\pgfqpoint{5.303240in}{2.585803in}}%
\pgfusepath{clip}%
\pgfsetbuttcap%
\pgfsetroundjoin%
\pgfsetlinewidth{0.803000pt}%
\definecolor{currentstroke}{rgb}{0.690196,0.690196,0.690196}%
\pgfsetstrokecolor{currentstroke}%
\pgfsetdash{{0.800000pt}{1.320000pt}}{0.000000pt}%
\pgfpathmoveto{\pgfqpoint{3.340201in}{0.565123in}}%
\pgfpathlineto{\pgfqpoint{3.340201in}{3.150926in}}%
\pgfusepath{stroke}%
\end{pgfscope}%
\begin{pgfscope}%
\pgfsetbuttcap%
\pgfsetroundjoin%
\definecolor{currentfill}{rgb}{0.000000,0.000000,0.000000}%
\pgfsetfillcolor{currentfill}%
\pgfsetlinewidth{0.803000pt}%
\definecolor{currentstroke}{rgb}{0.000000,0.000000,0.000000}%
\pgfsetstrokecolor{currentstroke}%
\pgfsetdash{}{0pt}%
\pgfsys@defobject{currentmarker}{\pgfqpoint{0.000000in}{-0.048611in}}{\pgfqpoint{0.000000in}{0.000000in}}{%
\pgfpathmoveto{\pgfqpoint{0.000000in}{0.000000in}}%
\pgfpathlineto{\pgfqpoint{0.000000in}{-0.048611in}}%
\pgfusepath{stroke,fill}%
}%
\begin{pgfscope}%
\pgfsys@transformshift{3.340201in}{0.565123in}%
\pgfsys@useobject{currentmarker}{}%
\end{pgfscope}%
\end{pgfscope}%
\begin{pgfscope}%
\definecolor{textcolor}{rgb}{0.000000,0.000000,0.000000}%
\pgfsetstrokecolor{textcolor}%
\pgfsetfillcolor{textcolor}%
\pgftext[x=3.340201in,y=0.467901in,,top]{\color{textcolor}\rmfamily\fontsize{10.000000}{12.000000}\selectfont \(\displaystyle {11.00}\)}%
\end{pgfscope}%
\begin{pgfscope}%
\pgfpathrectangle{\pgfqpoint{0.688581in}{0.565123in}}{\pgfqpoint{5.303240in}{2.585803in}}%
\pgfusepath{clip}%
\pgfsetbuttcap%
\pgfsetroundjoin%
\pgfsetlinewidth{0.803000pt}%
\definecolor{currentstroke}{rgb}{0.690196,0.690196,0.690196}%
\pgfsetstrokecolor{currentstroke}%
\pgfsetdash{{0.800000pt}{1.320000pt}}{0.000000pt}%
\pgfpathmoveto{\pgfqpoint{4.003106in}{0.565123in}}%
\pgfpathlineto{\pgfqpoint{4.003106in}{3.150926in}}%
\pgfusepath{stroke}%
\end{pgfscope}%
\begin{pgfscope}%
\pgfsetbuttcap%
\pgfsetroundjoin%
\definecolor{currentfill}{rgb}{0.000000,0.000000,0.000000}%
\pgfsetfillcolor{currentfill}%
\pgfsetlinewidth{0.803000pt}%
\definecolor{currentstroke}{rgb}{0.000000,0.000000,0.000000}%
\pgfsetstrokecolor{currentstroke}%
\pgfsetdash{}{0pt}%
\pgfsys@defobject{currentmarker}{\pgfqpoint{0.000000in}{-0.048611in}}{\pgfqpoint{0.000000in}{0.000000in}}{%
\pgfpathmoveto{\pgfqpoint{0.000000in}{0.000000in}}%
\pgfpathlineto{\pgfqpoint{0.000000in}{-0.048611in}}%
\pgfusepath{stroke,fill}%
}%
\begin{pgfscope}%
\pgfsys@transformshift{4.003106in}{0.565123in}%
\pgfsys@useobject{currentmarker}{}%
\end{pgfscope}%
\end{pgfscope}%
\begin{pgfscope}%
\definecolor{textcolor}{rgb}{0.000000,0.000000,0.000000}%
\pgfsetstrokecolor{textcolor}%
\pgfsetfillcolor{textcolor}%
\pgftext[x=4.003106in,y=0.467901in,,top]{\color{textcolor}\rmfamily\fontsize{10.000000}{12.000000}\selectfont \(\displaystyle {11.25}\)}%
\end{pgfscope}%
\begin{pgfscope}%
\pgfpathrectangle{\pgfqpoint{0.688581in}{0.565123in}}{\pgfqpoint{5.303240in}{2.585803in}}%
\pgfusepath{clip}%
\pgfsetbuttcap%
\pgfsetroundjoin%
\pgfsetlinewidth{0.803000pt}%
\definecolor{currentstroke}{rgb}{0.690196,0.690196,0.690196}%
\pgfsetstrokecolor{currentstroke}%
\pgfsetdash{{0.800000pt}{1.320000pt}}{0.000000pt}%
\pgfpathmoveto{\pgfqpoint{4.666011in}{0.565123in}}%
\pgfpathlineto{\pgfqpoint{4.666011in}{3.150926in}}%
\pgfusepath{stroke}%
\end{pgfscope}%
\begin{pgfscope}%
\pgfsetbuttcap%
\pgfsetroundjoin%
\definecolor{currentfill}{rgb}{0.000000,0.000000,0.000000}%
\pgfsetfillcolor{currentfill}%
\pgfsetlinewidth{0.803000pt}%
\definecolor{currentstroke}{rgb}{0.000000,0.000000,0.000000}%
\pgfsetstrokecolor{currentstroke}%
\pgfsetdash{}{0pt}%
\pgfsys@defobject{currentmarker}{\pgfqpoint{0.000000in}{-0.048611in}}{\pgfqpoint{0.000000in}{0.000000in}}{%
\pgfpathmoveto{\pgfqpoint{0.000000in}{0.000000in}}%
\pgfpathlineto{\pgfqpoint{0.000000in}{-0.048611in}}%
\pgfusepath{stroke,fill}%
}%
\begin{pgfscope}%
\pgfsys@transformshift{4.666011in}{0.565123in}%
\pgfsys@useobject{currentmarker}{}%
\end{pgfscope}%
\end{pgfscope}%
\begin{pgfscope}%
\definecolor{textcolor}{rgb}{0.000000,0.000000,0.000000}%
\pgfsetstrokecolor{textcolor}%
\pgfsetfillcolor{textcolor}%
\pgftext[x=4.666011in,y=0.467901in,,top]{\color{textcolor}\rmfamily\fontsize{10.000000}{12.000000}\selectfont \(\displaystyle {11.50}\)}%
\end{pgfscope}%
\begin{pgfscope}%
\pgfpathrectangle{\pgfqpoint{0.688581in}{0.565123in}}{\pgfqpoint{5.303240in}{2.585803in}}%
\pgfusepath{clip}%
\pgfsetbuttcap%
\pgfsetroundjoin%
\pgfsetlinewidth{0.803000pt}%
\definecolor{currentstroke}{rgb}{0.690196,0.690196,0.690196}%
\pgfsetstrokecolor{currentstroke}%
\pgfsetdash{{0.800000pt}{1.320000pt}}{0.000000pt}%
\pgfpathmoveto{\pgfqpoint{5.328916in}{0.565123in}}%
\pgfpathlineto{\pgfqpoint{5.328916in}{3.150926in}}%
\pgfusepath{stroke}%
\end{pgfscope}%
\begin{pgfscope}%
\pgfsetbuttcap%
\pgfsetroundjoin%
\definecolor{currentfill}{rgb}{0.000000,0.000000,0.000000}%
\pgfsetfillcolor{currentfill}%
\pgfsetlinewidth{0.803000pt}%
\definecolor{currentstroke}{rgb}{0.000000,0.000000,0.000000}%
\pgfsetstrokecolor{currentstroke}%
\pgfsetdash{}{0pt}%
\pgfsys@defobject{currentmarker}{\pgfqpoint{0.000000in}{-0.048611in}}{\pgfqpoint{0.000000in}{0.000000in}}{%
\pgfpathmoveto{\pgfqpoint{0.000000in}{0.000000in}}%
\pgfpathlineto{\pgfqpoint{0.000000in}{-0.048611in}}%
\pgfusepath{stroke,fill}%
}%
\begin{pgfscope}%
\pgfsys@transformshift{5.328916in}{0.565123in}%
\pgfsys@useobject{currentmarker}{}%
\end{pgfscope}%
\end{pgfscope}%
\begin{pgfscope}%
\definecolor{textcolor}{rgb}{0.000000,0.000000,0.000000}%
\pgfsetstrokecolor{textcolor}%
\pgfsetfillcolor{textcolor}%
\pgftext[x=5.328916in,y=0.467901in,,top]{\color{textcolor}\rmfamily\fontsize{10.000000}{12.000000}\selectfont \(\displaystyle {11.75}\)}%
\end{pgfscope}%
\begin{pgfscope}%
\pgfpathrectangle{\pgfqpoint{0.688581in}{0.565123in}}{\pgfqpoint{5.303240in}{2.585803in}}%
\pgfusepath{clip}%
\pgfsetbuttcap%
\pgfsetroundjoin%
\pgfsetlinewidth{0.803000pt}%
\definecolor{currentstroke}{rgb}{0.690196,0.690196,0.690196}%
\pgfsetstrokecolor{currentstroke}%
\pgfsetdash{{0.800000pt}{1.320000pt}}{0.000000pt}%
\pgfpathmoveto{\pgfqpoint{5.991820in}{0.565123in}}%
\pgfpathlineto{\pgfqpoint{5.991820in}{3.150926in}}%
\pgfusepath{stroke}%
\end{pgfscope}%
\begin{pgfscope}%
\pgfsetbuttcap%
\pgfsetroundjoin%
\definecolor{currentfill}{rgb}{0.000000,0.000000,0.000000}%
\pgfsetfillcolor{currentfill}%
\pgfsetlinewidth{0.803000pt}%
\definecolor{currentstroke}{rgb}{0.000000,0.000000,0.000000}%
\pgfsetstrokecolor{currentstroke}%
\pgfsetdash{}{0pt}%
\pgfsys@defobject{currentmarker}{\pgfqpoint{0.000000in}{-0.048611in}}{\pgfqpoint{0.000000in}{0.000000in}}{%
\pgfpathmoveto{\pgfqpoint{0.000000in}{0.000000in}}%
\pgfpathlineto{\pgfqpoint{0.000000in}{-0.048611in}}%
\pgfusepath{stroke,fill}%
}%
\begin{pgfscope}%
\pgfsys@transformshift{5.991820in}{0.565123in}%
\pgfsys@useobject{currentmarker}{}%
\end{pgfscope}%
\end{pgfscope}%
\begin{pgfscope}%
\definecolor{textcolor}{rgb}{0.000000,0.000000,0.000000}%
\pgfsetstrokecolor{textcolor}%
\pgfsetfillcolor{textcolor}%
\pgftext[x=5.991820in,y=0.467901in,,top]{\color{textcolor}\rmfamily\fontsize{10.000000}{12.000000}\selectfont \(\displaystyle {12.00}\)}%
\end{pgfscope}%
\begin{pgfscope}%
\definecolor{textcolor}{rgb}{0.000000,0.000000,0.000000}%
\pgfsetstrokecolor{textcolor}%
\pgfsetfillcolor{textcolor}%
\pgftext[x=3.340201in,y=0.288889in,,top]{\color{textcolor}\rmfamily\fontsize{10.000000}{12.000000}\selectfont Frequency (MHz)}%
\end{pgfscope}%
\begin{pgfscope}%
\pgfpathrectangle{\pgfqpoint{0.688581in}{0.565123in}}{\pgfqpoint{5.303240in}{2.585803in}}%
\pgfusepath{clip}%
\pgfsetbuttcap%
\pgfsetroundjoin%
\pgfsetlinewidth{0.803000pt}%
\definecolor{currentstroke}{rgb}{0.690196,0.690196,0.690196}%
\pgfsetstrokecolor{currentstroke}%
\pgfsetdash{{0.800000pt}{1.320000pt}}{0.000000pt}%
\pgfpathmoveto{\pgfqpoint{0.688581in}{0.701275in}}%
\pgfpathlineto{\pgfqpoint{5.991820in}{0.701275in}}%
\pgfusepath{stroke}%
\end{pgfscope}%
\begin{pgfscope}%
\pgfsetbuttcap%
\pgfsetroundjoin%
\definecolor{currentfill}{rgb}{0.000000,0.000000,0.000000}%
\pgfsetfillcolor{currentfill}%
\pgfsetlinewidth{0.803000pt}%
\definecolor{currentstroke}{rgb}{0.000000,0.000000,0.000000}%
\pgfsetstrokecolor{currentstroke}%
\pgfsetdash{}{0pt}%
\pgfsys@defobject{currentmarker}{\pgfqpoint{-0.048611in}{0.000000in}}{\pgfqpoint{-0.000000in}{0.000000in}}{%
\pgfpathmoveto{\pgfqpoint{-0.000000in}{0.000000in}}%
\pgfpathlineto{\pgfqpoint{-0.048611in}{0.000000in}}%
\pgfusepath{stroke,fill}%
}%
\begin{pgfscope}%
\pgfsys@transformshift{0.688581in}{0.701275in}%
\pgfsys@useobject{currentmarker}{}%
\end{pgfscope}%
\end{pgfscope}%
\begin{pgfscope}%
\definecolor{textcolor}{rgb}{0.000000,0.000000,0.000000}%
\pgfsetstrokecolor{textcolor}%
\pgfsetfillcolor{textcolor}%
\pgftext[x=0.344444in, y=0.653050in, left, base]{\color{textcolor}\rmfamily\fontsize{10.000000}{12.000000}\selectfont \(\displaystyle {\ensuremath{-}50}\)}%
\end{pgfscope}%
\begin{pgfscope}%
\pgfpathrectangle{\pgfqpoint{0.688581in}{0.565123in}}{\pgfqpoint{5.303240in}{2.585803in}}%
\pgfusepath{clip}%
\pgfsetbuttcap%
\pgfsetroundjoin%
\pgfsetlinewidth{0.803000pt}%
\definecolor{currentstroke}{rgb}{0.690196,0.690196,0.690196}%
\pgfsetstrokecolor{currentstroke}%
\pgfsetdash{{0.800000pt}{1.320000pt}}{0.000000pt}%
\pgfpathmoveto{\pgfqpoint{0.688581in}{0.992135in}}%
\pgfpathlineto{\pgfqpoint{5.991820in}{0.992135in}}%
\pgfusepath{stroke}%
\end{pgfscope}%
\begin{pgfscope}%
\pgfsetbuttcap%
\pgfsetroundjoin%
\definecolor{currentfill}{rgb}{0.000000,0.000000,0.000000}%
\pgfsetfillcolor{currentfill}%
\pgfsetlinewidth{0.803000pt}%
\definecolor{currentstroke}{rgb}{0.000000,0.000000,0.000000}%
\pgfsetstrokecolor{currentstroke}%
\pgfsetdash{}{0pt}%
\pgfsys@defobject{currentmarker}{\pgfqpoint{-0.048611in}{0.000000in}}{\pgfqpoint{-0.000000in}{0.000000in}}{%
\pgfpathmoveto{\pgfqpoint{-0.000000in}{0.000000in}}%
\pgfpathlineto{\pgfqpoint{-0.048611in}{0.000000in}}%
\pgfusepath{stroke,fill}%
}%
\begin{pgfscope}%
\pgfsys@transformshift{0.688581in}{0.992135in}%
\pgfsys@useobject{currentmarker}{}%
\end{pgfscope}%
\end{pgfscope}%
\begin{pgfscope}%
\definecolor{textcolor}{rgb}{0.000000,0.000000,0.000000}%
\pgfsetstrokecolor{textcolor}%
\pgfsetfillcolor{textcolor}%
\pgftext[x=0.344444in, y=0.943910in, left, base]{\color{textcolor}\rmfamily\fontsize{10.000000}{12.000000}\selectfont \(\displaystyle {\ensuremath{-}45}\)}%
\end{pgfscope}%
\begin{pgfscope}%
\pgfpathrectangle{\pgfqpoint{0.688581in}{0.565123in}}{\pgfqpoint{5.303240in}{2.585803in}}%
\pgfusepath{clip}%
\pgfsetbuttcap%
\pgfsetroundjoin%
\pgfsetlinewidth{0.803000pt}%
\definecolor{currentstroke}{rgb}{0.690196,0.690196,0.690196}%
\pgfsetstrokecolor{currentstroke}%
\pgfsetdash{{0.800000pt}{1.320000pt}}{0.000000pt}%
\pgfpathmoveto{\pgfqpoint{0.688581in}{1.282995in}}%
\pgfpathlineto{\pgfqpoint{5.991820in}{1.282995in}}%
\pgfusepath{stroke}%
\end{pgfscope}%
\begin{pgfscope}%
\pgfsetbuttcap%
\pgfsetroundjoin%
\definecolor{currentfill}{rgb}{0.000000,0.000000,0.000000}%
\pgfsetfillcolor{currentfill}%
\pgfsetlinewidth{0.803000pt}%
\definecolor{currentstroke}{rgb}{0.000000,0.000000,0.000000}%
\pgfsetstrokecolor{currentstroke}%
\pgfsetdash{}{0pt}%
\pgfsys@defobject{currentmarker}{\pgfqpoint{-0.048611in}{0.000000in}}{\pgfqpoint{-0.000000in}{0.000000in}}{%
\pgfpathmoveto{\pgfqpoint{-0.000000in}{0.000000in}}%
\pgfpathlineto{\pgfqpoint{-0.048611in}{0.000000in}}%
\pgfusepath{stroke,fill}%
}%
\begin{pgfscope}%
\pgfsys@transformshift{0.688581in}{1.282995in}%
\pgfsys@useobject{currentmarker}{}%
\end{pgfscope}%
\end{pgfscope}%
\begin{pgfscope}%
\definecolor{textcolor}{rgb}{0.000000,0.000000,0.000000}%
\pgfsetstrokecolor{textcolor}%
\pgfsetfillcolor{textcolor}%
\pgftext[x=0.344444in, y=1.234769in, left, base]{\color{textcolor}\rmfamily\fontsize{10.000000}{12.000000}\selectfont \(\displaystyle {\ensuremath{-}40}\)}%
\end{pgfscope}%
\begin{pgfscope}%
\pgfpathrectangle{\pgfqpoint{0.688581in}{0.565123in}}{\pgfqpoint{5.303240in}{2.585803in}}%
\pgfusepath{clip}%
\pgfsetbuttcap%
\pgfsetroundjoin%
\pgfsetlinewidth{0.803000pt}%
\definecolor{currentstroke}{rgb}{0.690196,0.690196,0.690196}%
\pgfsetstrokecolor{currentstroke}%
\pgfsetdash{{0.800000pt}{1.320000pt}}{0.000000pt}%
\pgfpathmoveto{\pgfqpoint{0.688581in}{1.573855in}}%
\pgfpathlineto{\pgfqpoint{5.991820in}{1.573855in}}%
\pgfusepath{stroke}%
\end{pgfscope}%
\begin{pgfscope}%
\pgfsetbuttcap%
\pgfsetroundjoin%
\definecolor{currentfill}{rgb}{0.000000,0.000000,0.000000}%
\pgfsetfillcolor{currentfill}%
\pgfsetlinewidth{0.803000pt}%
\definecolor{currentstroke}{rgb}{0.000000,0.000000,0.000000}%
\pgfsetstrokecolor{currentstroke}%
\pgfsetdash{}{0pt}%
\pgfsys@defobject{currentmarker}{\pgfqpoint{-0.048611in}{0.000000in}}{\pgfqpoint{-0.000000in}{0.000000in}}{%
\pgfpathmoveto{\pgfqpoint{-0.000000in}{0.000000in}}%
\pgfpathlineto{\pgfqpoint{-0.048611in}{0.000000in}}%
\pgfusepath{stroke,fill}%
}%
\begin{pgfscope}%
\pgfsys@transformshift{0.688581in}{1.573855in}%
\pgfsys@useobject{currentmarker}{}%
\end{pgfscope}%
\end{pgfscope}%
\begin{pgfscope}%
\definecolor{textcolor}{rgb}{0.000000,0.000000,0.000000}%
\pgfsetstrokecolor{textcolor}%
\pgfsetfillcolor{textcolor}%
\pgftext[x=0.344444in, y=1.525629in, left, base]{\color{textcolor}\rmfamily\fontsize{10.000000}{12.000000}\selectfont \(\displaystyle {\ensuremath{-}35}\)}%
\end{pgfscope}%
\begin{pgfscope}%
\pgfpathrectangle{\pgfqpoint{0.688581in}{0.565123in}}{\pgfqpoint{5.303240in}{2.585803in}}%
\pgfusepath{clip}%
\pgfsetbuttcap%
\pgfsetroundjoin%
\pgfsetlinewidth{0.803000pt}%
\definecolor{currentstroke}{rgb}{0.690196,0.690196,0.690196}%
\pgfsetstrokecolor{currentstroke}%
\pgfsetdash{{0.800000pt}{1.320000pt}}{0.000000pt}%
\pgfpathmoveto{\pgfqpoint{0.688581in}{1.864715in}}%
\pgfpathlineto{\pgfqpoint{5.991820in}{1.864715in}}%
\pgfusepath{stroke}%
\end{pgfscope}%
\begin{pgfscope}%
\pgfsetbuttcap%
\pgfsetroundjoin%
\definecolor{currentfill}{rgb}{0.000000,0.000000,0.000000}%
\pgfsetfillcolor{currentfill}%
\pgfsetlinewidth{0.803000pt}%
\definecolor{currentstroke}{rgb}{0.000000,0.000000,0.000000}%
\pgfsetstrokecolor{currentstroke}%
\pgfsetdash{}{0pt}%
\pgfsys@defobject{currentmarker}{\pgfqpoint{-0.048611in}{0.000000in}}{\pgfqpoint{-0.000000in}{0.000000in}}{%
\pgfpathmoveto{\pgfqpoint{-0.000000in}{0.000000in}}%
\pgfpathlineto{\pgfqpoint{-0.048611in}{0.000000in}}%
\pgfusepath{stroke,fill}%
}%
\begin{pgfscope}%
\pgfsys@transformshift{0.688581in}{1.864715in}%
\pgfsys@useobject{currentmarker}{}%
\end{pgfscope}%
\end{pgfscope}%
\begin{pgfscope}%
\definecolor{textcolor}{rgb}{0.000000,0.000000,0.000000}%
\pgfsetstrokecolor{textcolor}%
\pgfsetfillcolor{textcolor}%
\pgftext[x=0.344444in, y=1.816489in, left, base]{\color{textcolor}\rmfamily\fontsize{10.000000}{12.000000}\selectfont \(\displaystyle {\ensuremath{-}30}\)}%
\end{pgfscope}%
\begin{pgfscope}%
\pgfpathrectangle{\pgfqpoint{0.688581in}{0.565123in}}{\pgfqpoint{5.303240in}{2.585803in}}%
\pgfusepath{clip}%
\pgfsetbuttcap%
\pgfsetroundjoin%
\pgfsetlinewidth{0.803000pt}%
\definecolor{currentstroke}{rgb}{0.690196,0.690196,0.690196}%
\pgfsetstrokecolor{currentstroke}%
\pgfsetdash{{0.800000pt}{1.320000pt}}{0.000000pt}%
\pgfpathmoveto{\pgfqpoint{0.688581in}{2.155574in}}%
\pgfpathlineto{\pgfqpoint{5.991820in}{2.155574in}}%
\pgfusepath{stroke}%
\end{pgfscope}%
\begin{pgfscope}%
\pgfsetbuttcap%
\pgfsetroundjoin%
\definecolor{currentfill}{rgb}{0.000000,0.000000,0.000000}%
\pgfsetfillcolor{currentfill}%
\pgfsetlinewidth{0.803000pt}%
\definecolor{currentstroke}{rgb}{0.000000,0.000000,0.000000}%
\pgfsetstrokecolor{currentstroke}%
\pgfsetdash{}{0pt}%
\pgfsys@defobject{currentmarker}{\pgfqpoint{-0.048611in}{0.000000in}}{\pgfqpoint{-0.000000in}{0.000000in}}{%
\pgfpathmoveto{\pgfqpoint{-0.000000in}{0.000000in}}%
\pgfpathlineto{\pgfqpoint{-0.048611in}{0.000000in}}%
\pgfusepath{stroke,fill}%
}%
\begin{pgfscope}%
\pgfsys@transformshift{0.688581in}{2.155574in}%
\pgfsys@useobject{currentmarker}{}%
\end{pgfscope}%
\end{pgfscope}%
\begin{pgfscope}%
\definecolor{textcolor}{rgb}{0.000000,0.000000,0.000000}%
\pgfsetstrokecolor{textcolor}%
\pgfsetfillcolor{textcolor}%
\pgftext[x=0.344444in, y=2.107349in, left, base]{\color{textcolor}\rmfamily\fontsize{10.000000}{12.000000}\selectfont \(\displaystyle {\ensuremath{-}25}\)}%
\end{pgfscope}%
\begin{pgfscope}%
\pgfpathrectangle{\pgfqpoint{0.688581in}{0.565123in}}{\pgfqpoint{5.303240in}{2.585803in}}%
\pgfusepath{clip}%
\pgfsetbuttcap%
\pgfsetroundjoin%
\pgfsetlinewidth{0.803000pt}%
\definecolor{currentstroke}{rgb}{0.690196,0.690196,0.690196}%
\pgfsetstrokecolor{currentstroke}%
\pgfsetdash{{0.800000pt}{1.320000pt}}{0.000000pt}%
\pgfpathmoveto{\pgfqpoint{0.688581in}{2.446434in}}%
\pgfpathlineto{\pgfqpoint{5.991820in}{2.446434in}}%
\pgfusepath{stroke}%
\end{pgfscope}%
\begin{pgfscope}%
\pgfsetbuttcap%
\pgfsetroundjoin%
\definecolor{currentfill}{rgb}{0.000000,0.000000,0.000000}%
\pgfsetfillcolor{currentfill}%
\pgfsetlinewidth{0.803000pt}%
\definecolor{currentstroke}{rgb}{0.000000,0.000000,0.000000}%
\pgfsetstrokecolor{currentstroke}%
\pgfsetdash{}{0pt}%
\pgfsys@defobject{currentmarker}{\pgfqpoint{-0.048611in}{0.000000in}}{\pgfqpoint{-0.000000in}{0.000000in}}{%
\pgfpathmoveto{\pgfqpoint{-0.000000in}{0.000000in}}%
\pgfpathlineto{\pgfqpoint{-0.048611in}{0.000000in}}%
\pgfusepath{stroke,fill}%
}%
\begin{pgfscope}%
\pgfsys@transformshift{0.688581in}{2.446434in}%
\pgfsys@useobject{currentmarker}{}%
\end{pgfscope}%
\end{pgfscope}%
\begin{pgfscope}%
\definecolor{textcolor}{rgb}{0.000000,0.000000,0.000000}%
\pgfsetstrokecolor{textcolor}%
\pgfsetfillcolor{textcolor}%
\pgftext[x=0.344444in, y=2.398209in, left, base]{\color{textcolor}\rmfamily\fontsize{10.000000}{12.000000}\selectfont \(\displaystyle {\ensuremath{-}20}\)}%
\end{pgfscope}%
\begin{pgfscope}%
\pgfpathrectangle{\pgfqpoint{0.688581in}{0.565123in}}{\pgfqpoint{5.303240in}{2.585803in}}%
\pgfusepath{clip}%
\pgfsetbuttcap%
\pgfsetroundjoin%
\pgfsetlinewidth{0.803000pt}%
\definecolor{currentstroke}{rgb}{0.690196,0.690196,0.690196}%
\pgfsetstrokecolor{currentstroke}%
\pgfsetdash{{0.800000pt}{1.320000pt}}{0.000000pt}%
\pgfpathmoveto{\pgfqpoint{0.688581in}{2.737294in}}%
\pgfpathlineto{\pgfqpoint{5.991820in}{2.737294in}}%
\pgfusepath{stroke}%
\end{pgfscope}%
\begin{pgfscope}%
\pgfsetbuttcap%
\pgfsetroundjoin%
\definecolor{currentfill}{rgb}{0.000000,0.000000,0.000000}%
\pgfsetfillcolor{currentfill}%
\pgfsetlinewidth{0.803000pt}%
\definecolor{currentstroke}{rgb}{0.000000,0.000000,0.000000}%
\pgfsetstrokecolor{currentstroke}%
\pgfsetdash{}{0pt}%
\pgfsys@defobject{currentmarker}{\pgfqpoint{-0.048611in}{0.000000in}}{\pgfqpoint{-0.000000in}{0.000000in}}{%
\pgfpathmoveto{\pgfqpoint{-0.000000in}{0.000000in}}%
\pgfpathlineto{\pgfqpoint{-0.048611in}{0.000000in}}%
\pgfusepath{stroke,fill}%
}%
\begin{pgfscope}%
\pgfsys@transformshift{0.688581in}{2.737294in}%
\pgfsys@useobject{currentmarker}{}%
\end{pgfscope}%
\end{pgfscope}%
\begin{pgfscope}%
\definecolor{textcolor}{rgb}{0.000000,0.000000,0.000000}%
\pgfsetstrokecolor{textcolor}%
\pgfsetfillcolor{textcolor}%
\pgftext[x=0.344444in, y=2.689069in, left, base]{\color{textcolor}\rmfamily\fontsize{10.000000}{12.000000}\selectfont \(\displaystyle {\ensuremath{-}15}\)}%
\end{pgfscope}%
\begin{pgfscope}%
\pgfpathrectangle{\pgfqpoint{0.688581in}{0.565123in}}{\pgfqpoint{5.303240in}{2.585803in}}%
\pgfusepath{clip}%
\pgfsetbuttcap%
\pgfsetroundjoin%
\pgfsetlinewidth{0.803000pt}%
\definecolor{currentstroke}{rgb}{0.690196,0.690196,0.690196}%
\pgfsetstrokecolor{currentstroke}%
\pgfsetdash{{0.800000pt}{1.320000pt}}{0.000000pt}%
\pgfpathmoveto{\pgfqpoint{0.688581in}{3.028154in}}%
\pgfpathlineto{\pgfqpoint{5.991820in}{3.028154in}}%
\pgfusepath{stroke}%
\end{pgfscope}%
\begin{pgfscope}%
\pgfsetbuttcap%
\pgfsetroundjoin%
\definecolor{currentfill}{rgb}{0.000000,0.000000,0.000000}%
\pgfsetfillcolor{currentfill}%
\pgfsetlinewidth{0.803000pt}%
\definecolor{currentstroke}{rgb}{0.000000,0.000000,0.000000}%
\pgfsetstrokecolor{currentstroke}%
\pgfsetdash{}{0pt}%
\pgfsys@defobject{currentmarker}{\pgfqpoint{-0.048611in}{0.000000in}}{\pgfqpoint{-0.000000in}{0.000000in}}{%
\pgfpathmoveto{\pgfqpoint{-0.000000in}{0.000000in}}%
\pgfpathlineto{\pgfqpoint{-0.048611in}{0.000000in}}%
\pgfusepath{stroke,fill}%
}%
\begin{pgfscope}%
\pgfsys@transformshift{0.688581in}{3.028154in}%
\pgfsys@useobject{currentmarker}{}%
\end{pgfscope}%
\end{pgfscope}%
\begin{pgfscope}%
\definecolor{textcolor}{rgb}{0.000000,0.000000,0.000000}%
\pgfsetstrokecolor{textcolor}%
\pgfsetfillcolor{textcolor}%
\pgftext[x=0.344444in, y=2.979929in, left, base]{\color{textcolor}\rmfamily\fontsize{10.000000}{12.000000}\selectfont \(\displaystyle {\ensuremath{-}10}\)}%
\end{pgfscope}%
\begin{pgfscope}%
\definecolor{textcolor}{rgb}{0.000000,0.000000,0.000000}%
\pgfsetstrokecolor{textcolor}%
\pgfsetfillcolor{textcolor}%
\pgftext[x=0.288889in,y=1.858025in,,bottom,rotate=90.000000]{\color{textcolor}\rmfamily\fontsize{10.000000}{12.000000}\selectfont Transmission Loss (dB)}%
\end{pgfscope}%
\begin{pgfscope}%
\pgfpathrectangle{\pgfqpoint{0.688581in}{0.565123in}}{\pgfqpoint{5.303240in}{2.585803in}}%
\pgfusepath{clip}%
\pgfsetrectcap%
\pgfsetroundjoin%
\pgfsetlinewidth{2.007500pt}%
\definecolor{currentstroke}{rgb}{0.121569,0.466667,0.705882}%
\pgfsetstrokecolor{currentstroke}%
\pgfsetdash{}{0pt}%
\pgfpathmoveto{\pgfqpoint{0.678581in}{0.689278in}}%
\pgfpathlineto{\pgfqpoint{0.690063in}{0.716400in}}%
\pgfpathlineto{\pgfqpoint{0.704348in}{0.682660in}}%
\pgfpathlineto{\pgfqpoint{0.718632in}{0.911276in}}%
\pgfpathlineto{\pgfqpoint{0.732916in}{0.921165in}}%
\pgfpathlineto{\pgfqpoint{0.761485in}{0.914766in}}%
\pgfpathlineto{\pgfqpoint{0.775769in}{0.914766in}}%
\pgfpathlineto{\pgfqpoint{0.790053in}{0.904295in}}%
\pgfpathlineto{\pgfqpoint{0.804337in}{0.897314in}}%
\pgfpathlineto{\pgfqpoint{0.818622in}{0.904295in}}%
\pgfpathlineto{\pgfqpoint{0.832906in}{0.904295in}}%
\pgfpathlineto{\pgfqpoint{0.847190in}{0.907785in}}%
\pgfpathlineto{\pgfqpoint{0.861475in}{0.904295in}}%
\pgfpathlineto{\pgfqpoint{0.875759in}{0.897314in}}%
\pgfpathlineto{\pgfqpoint{0.890043in}{0.893824in}}%
\pgfpathlineto{\pgfqpoint{0.904327in}{0.911276in}}%
\pgfpathlineto{\pgfqpoint{0.932896in}{0.917675in}}%
\pgfpathlineto{\pgfqpoint{0.947180in}{0.931636in}}%
\pgfpathlineto{\pgfqpoint{0.961464in}{0.931636in}}%
\pgfpathlineto{\pgfqpoint{0.975749in}{0.935126in}}%
\pgfpathlineto{\pgfqpoint{0.990033in}{0.935126in}}%
\pgfpathlineto{\pgfqpoint{1.004317in}{0.917675in}}%
\pgfpathlineto{\pgfqpoint{1.018602in}{0.938617in}}%
\pgfpathlineto{\pgfqpoint{1.032886in}{0.921165in}}%
\pgfpathlineto{\pgfqpoint{1.047170in}{0.911276in}}%
\pgfpathlineto{\pgfqpoint{1.061454in}{0.904295in}}%
\pgfpathlineto{\pgfqpoint{1.075739in}{0.866483in}}%
\pgfpathlineto{\pgfqpoint{1.090023in}{0.907785in}}%
\pgfpathlineto{\pgfqpoint{1.104307in}{0.911276in}}%
\pgfpathlineto{\pgfqpoint{1.118591in}{0.921165in}}%
\pgfpathlineto{\pgfqpoint{1.132876in}{0.917675in}}%
\pgfpathlineto{\pgfqpoint{1.147160in}{0.979337in}}%
\pgfpathlineto{\pgfqpoint{1.161444in}{0.921165in}}%
\pgfpathlineto{\pgfqpoint{1.190013in}{0.900805in}}%
\pgfpathlineto{\pgfqpoint{1.204297in}{0.883935in}}%
\pgfpathlineto{\pgfqpoint{1.218581in}{0.696039in}}%
\pgfpathlineto{\pgfqpoint{1.232866in}{0.682660in}}%
\pgfpathlineto{\pgfqpoint{1.247150in}{0.846123in}}%
\pgfpathlineto{\pgfqpoint{1.261434in}{0.917675in}}%
\pgfpathlineto{\pgfqpoint{1.275719in}{0.866483in}}%
\pgfpathlineto{\pgfqpoint{1.290003in}{0.917675in}}%
\pgfpathlineto{\pgfqpoint{1.304287in}{0.924655in}}%
\pgfpathlineto{\pgfqpoint{1.318571in}{0.911276in}}%
\pgfpathlineto{\pgfqpoint{1.332856in}{0.921165in}}%
\pgfpathlineto{\pgfqpoint{1.347140in}{0.921165in}}%
\pgfpathlineto{\pgfqpoint{1.361424in}{0.941525in}}%
\pgfpathlineto{\pgfqpoint{1.375708in}{0.791441in}}%
\pgfpathlineto{\pgfqpoint{1.389993in}{0.778062in}}%
\pgfpathlineto{\pgfqpoint{1.404277in}{0.696039in}}%
\pgfpathlineto{\pgfqpoint{1.418561in}{0.682660in}}%
\pgfpathlineto{\pgfqpoint{1.432846in}{0.747231in}}%
\pgfpathlineto{\pgfqpoint{1.447130in}{0.873464in}}%
\pgfpathlineto{\pgfqpoint{1.461414in}{0.917675in}}%
\pgfpathlineto{\pgfqpoint{1.475698in}{0.914766in}}%
\pgfpathlineto{\pgfqpoint{1.489983in}{0.907785in}}%
\pgfpathlineto{\pgfqpoint{1.504267in}{0.822273in}}%
\pgfpathlineto{\pgfqpoint{1.518551in}{0.778062in}}%
\pgfpathlineto{\pgfqpoint{1.532835in}{0.682660in}}%
\pgfpathlineto{\pgfqpoint{1.547120in}{0.699530in}}%
\pgfpathlineto{\pgfqpoint{1.561404in}{0.730361in}}%
\pgfpathlineto{\pgfqpoint{1.575688in}{0.737341in}}%
\pgfpathlineto{\pgfqpoint{1.589973in}{0.710001in}}%
\pgfpathlineto{\pgfqpoint{1.604257in}{0.757702in}}%
\pgfpathlineto{\pgfqpoint{1.618541in}{0.723380in}}%
\pgfpathlineto{\pgfqpoint{1.632825in}{0.767591in}}%
\pgfpathlineto{\pgfqpoint{1.647110in}{0.692549in}}%
\pgfpathlineto{\pgfqpoint{1.661394in}{0.706510in}}%
\pgfpathlineto{\pgfqpoint{1.675678in}{0.723380in}}%
\pgfpathlineto{\pgfqpoint{1.689963in}{0.733851in}}%
\pgfpathlineto{\pgfqpoint{1.704247in}{0.785043in}}%
\pgfpathlineto{\pgfqpoint{1.718531in}{0.754211in}}%
\pgfpathlineto{\pgfqpoint{1.732815in}{0.730361in}}%
\pgfpathlineto{\pgfqpoint{1.747100in}{0.733851in}}%
\pgfpathlineto{\pgfqpoint{1.761384in}{0.801912in}}%
\pgfpathlineto{\pgfqpoint{1.775668in}{0.761192in}}%
\pgfpathlineto{\pgfqpoint{1.789952in}{0.761192in}}%
\pgfpathlineto{\pgfqpoint{1.804237in}{0.726871in}}%
\pgfpathlineto{\pgfqpoint{1.818521in}{0.740250in}}%
\pgfpathlineto{\pgfqpoint{1.832805in}{0.778062in}}%
\pgfpathlineto{\pgfqpoint{1.847090in}{0.764682in}}%
\pgfpathlineto{\pgfqpoint{1.861374in}{0.757702in}}%
\pgfpathlineto{\pgfqpoint{1.889942in}{0.805403in}}%
\pgfpathlineto{\pgfqpoint{1.904227in}{0.836234in}}%
\pgfpathlineto{\pgfqpoint{1.918511in}{0.791441in}}%
\pgfpathlineto{\pgfqpoint{1.932795in}{0.829253in}}%
\pgfpathlineto{\pgfqpoint{1.947079in}{0.860084in}}%
\pgfpathlineto{\pgfqpoint{1.961364in}{0.808893in}}%
\pgfpathlineto{\pgfqpoint{1.975648in}{0.825763in}}%
\pgfpathlineto{\pgfqpoint{1.989932in}{0.873464in}}%
\pgfpathlineto{\pgfqpoint{2.004217in}{0.904295in}}%
\pgfpathlineto{\pgfqpoint{2.018501in}{0.938617in}}%
\pgfpathlineto{\pgfqpoint{2.032785in}{0.962467in}}%
\pgfpathlineto{\pgfqpoint{2.075638in}{1.102080in}}%
\pgfpathlineto{\pgfqpoint{2.104207in}{1.218424in}}%
\pgfpathlineto{\pgfqpoint{2.132775in}{1.351056in}}%
\pgfpathlineto{\pgfqpoint{2.161344in}{1.491250in}}%
\pgfpathlineto{\pgfqpoint{2.189912in}{1.637844in}}%
\pgfpathlineto{\pgfqpoint{2.218481in}{1.794908in}}%
\pgfpathlineto{\pgfqpoint{2.247049in}{1.965352in}}%
\pgfpathlineto{\pgfqpoint{2.289902in}{2.241669in}}%
\pgfpathlineto{\pgfqpoint{2.347039in}{2.630840in}}%
\pgfpathlineto{\pgfqpoint{2.361323in}{2.719261in}}%
\pgfpathlineto{\pgfqpoint{2.375608in}{2.797793in}}%
\pgfpathlineto{\pgfqpoint{2.389892in}{2.862946in}}%
\pgfpathlineto{\pgfqpoint{2.404176in}{2.917046in}}%
\pgfpathlineto{\pgfqpoint{2.418461in}{2.961838in}}%
\pgfpathlineto{\pgfqpoint{2.432745in}{2.992088in}}%
\pgfpathlineto{\pgfqpoint{2.447029in}{3.015938in}}%
\pgfpathlineto{\pgfqpoint{2.461313in}{3.029899in}}%
\pgfpathlineto{\pgfqpoint{2.475598in}{3.033390in}}%
\pgfpathlineto{\pgfqpoint{2.489882in}{3.026409in}}%
\pgfpathlineto{\pgfqpoint{2.504166in}{3.013030in}}%
\pgfpathlineto{\pgfqpoint{2.518451in}{2.985689in}}%
\pgfpathlineto{\pgfqpoint{2.532735in}{2.954858in}}%
\pgfpathlineto{\pgfqpoint{2.547019in}{2.920536in}}%
\pgfpathlineto{\pgfqpoint{2.589872in}{2.811754in}}%
\pgfpathlineto{\pgfqpoint{2.604156in}{2.780923in}}%
\pgfpathlineto{\pgfqpoint{2.618440in}{2.753583in}}%
\pgfpathlineto{\pgfqpoint{2.647009in}{2.712280in}}%
\pgfpathlineto{\pgfqpoint{2.661293in}{2.698901in}}%
\pgfpathlineto{\pgfqpoint{2.675578in}{2.688430in}}%
\pgfpathlineto{\pgfqpoint{2.704146in}{2.682031in}}%
\pgfpathlineto{\pgfqpoint{2.718430in}{2.685521in}}%
\pgfpathlineto{\pgfqpoint{2.732715in}{2.691920in}}%
\pgfpathlineto{\pgfqpoint{2.761283in}{2.719261in}}%
\pgfpathlineto{\pgfqpoint{2.775567in}{2.739621in}}%
\pgfpathlineto{\pgfqpoint{2.804136in}{2.787904in}}%
\pgfpathlineto{\pgfqpoint{2.846989in}{2.869345in}}%
\pgfpathlineto{\pgfqpoint{2.861273in}{2.886796in}}%
\pgfpathlineto{\pgfqpoint{2.875557in}{2.896686in}}%
\pgfpathlineto{\pgfqpoint{2.889842in}{2.900176in}}%
\pgfpathlineto{\pgfqpoint{2.904126in}{2.889705in}}%
\pgfpathlineto{\pgfqpoint{2.918410in}{2.872835in}}%
\pgfpathlineto{\pgfqpoint{2.932695in}{2.848985in}}%
\pgfpathlineto{\pgfqpoint{2.946979in}{2.821644in}}%
\pgfpathlineto{\pgfqpoint{2.961263in}{2.797793in}}%
\pgfpathlineto{\pgfqpoint{2.975547in}{2.770452in}}%
\pgfpathlineto{\pgfqpoint{2.989832in}{2.746602in}}%
\pgfpathlineto{\pgfqpoint{3.004116in}{2.726242in}}%
\pgfpathlineto{\pgfqpoint{3.018400in}{2.702391in}}%
\pgfpathlineto{\pgfqpoint{3.032684in}{2.671560in}}%
\pgfpathlineto{\pgfqpoint{3.046969in}{2.634330in}}%
\pgfpathlineto{\pgfqpoint{3.061253in}{2.583139in}}%
\pgfpathlineto{\pgfqpoint{3.075537in}{2.517986in}}%
\pgfpathlineto{\pgfqpoint{3.089822in}{2.439454in}}%
\pgfpathlineto{\pgfqpoint{3.104106in}{2.351032in}}%
\pgfpathlineto{\pgfqpoint{3.189811in}{1.780947in}}%
\pgfpathlineto{\pgfqpoint{3.218380in}{1.607013in}}%
\pgfpathlineto{\pgfqpoint{3.246949in}{1.447040in}}%
\pgfpathlineto{\pgfqpoint{3.275517in}{1.293466in}}%
\pgfpathlineto{\pgfqpoint{3.318370in}{1.092191in}}%
\pgfpathlineto{\pgfqpoint{3.332654in}{1.037509in}}%
\pgfpathlineto{\pgfqpoint{3.346938in}{0.972356in}}%
\pgfpathlineto{\pgfqpoint{3.361223in}{0.928146in}}%
\pgfpathlineto{\pgfqpoint{3.389791in}{0.856594in}}%
\pgfpathlineto{\pgfqpoint{3.404076in}{0.836234in}}%
\pgfpathlineto{\pgfqpoint{3.418360in}{0.791441in}}%
\pgfpathlineto{\pgfqpoint{3.432644in}{0.754211in}}%
\pgfpathlineto{\pgfqpoint{3.446928in}{0.747231in}}%
\pgfpathlineto{\pgfqpoint{3.461213in}{0.730361in}}%
\pgfpathlineto{\pgfqpoint{3.475497in}{0.767591in}}%
\pgfpathlineto{\pgfqpoint{3.489781in}{0.716400in}}%
\pgfpathlineto{\pgfqpoint{3.504066in}{0.696039in}}%
\pgfpathlineto{\pgfqpoint{3.518350in}{0.682660in}}%
\pgfpathlineto{\pgfqpoint{3.532634in}{0.682660in}}%
\pgfpathlineto{\pgfqpoint{3.546918in}{0.754211in}}%
\pgfpathlineto{\pgfqpoint{3.561203in}{0.733851in}}%
\pgfpathlineto{\pgfqpoint{3.575487in}{0.682660in}}%
\pgfpathlineto{\pgfqpoint{3.589771in}{0.703020in}}%
\pgfpathlineto{\pgfqpoint{3.604055in}{0.699530in}}%
\pgfpathlineto{\pgfqpoint{3.618340in}{0.686150in}}%
\pgfpathlineto{\pgfqpoint{3.632624in}{0.733851in}}%
\pgfpathlineto{\pgfqpoint{3.646908in}{0.682660in}}%
\pgfpathlineto{\pgfqpoint{3.675477in}{0.682660in}}%
\pgfpathlineto{\pgfqpoint{3.689761in}{0.713491in}}%
\pgfpathlineto{\pgfqpoint{3.704045in}{0.730361in}}%
\pgfpathlineto{\pgfqpoint{3.718330in}{0.733851in}}%
\pgfpathlineto{\pgfqpoint{3.732614in}{0.713491in}}%
\pgfpathlineto{\pgfqpoint{3.746898in}{0.822273in}}%
\pgfpathlineto{\pgfqpoint{3.761182in}{0.805403in}}%
\pgfpathlineto{\pgfqpoint{3.775467in}{0.713491in}}%
\pgfpathlineto{\pgfqpoint{3.789751in}{0.771081in}}%
\pgfpathlineto{\pgfqpoint{3.804035in}{0.710001in}}%
\pgfpathlineto{\pgfqpoint{3.818320in}{0.682660in}}%
\pgfpathlineto{\pgfqpoint{3.832604in}{0.726871in}}%
\pgfpathlineto{\pgfqpoint{3.846888in}{0.740250in}}%
\pgfpathlineto{\pgfqpoint{3.861172in}{0.703020in}}%
\pgfpathlineto{\pgfqpoint{3.875457in}{0.703020in}}%
\pgfpathlineto{\pgfqpoint{3.904025in}{0.794932in}}%
\pgfpathlineto{\pgfqpoint{3.918310in}{0.754211in}}%
\pgfpathlineto{\pgfqpoint{3.932594in}{0.719890in}}%
\pgfpathlineto{\pgfqpoint{3.946878in}{0.682660in}}%
\pgfpathlineto{\pgfqpoint{3.961162in}{0.710001in}}%
\pgfpathlineto{\pgfqpoint{3.975447in}{0.682660in}}%
\pgfpathlineto{\pgfqpoint{4.046868in}{0.682660in}}%
\pgfpathlineto{\pgfqpoint{4.061152in}{0.689059in}}%
\pgfpathlineto{\pgfqpoint{4.075437in}{0.706510in}}%
\pgfpathlineto{\pgfqpoint{4.089721in}{0.682660in}}%
\pgfpathlineto{\pgfqpoint{4.104005in}{0.706510in}}%
\pgfpathlineto{\pgfqpoint{4.118289in}{0.699530in}}%
\pgfpathlineto{\pgfqpoint{4.132574in}{0.686150in}}%
\pgfpathlineto{\pgfqpoint{4.146858in}{0.699530in}}%
\pgfpathlineto{\pgfqpoint{4.161142in}{0.682660in}}%
\pgfpathlineto{\pgfqpoint{4.175426in}{0.754211in}}%
\pgfpathlineto{\pgfqpoint{4.189711in}{0.682660in}}%
\pgfpathlineto{\pgfqpoint{4.203995in}{0.682660in}}%
\pgfpathlineto{\pgfqpoint{4.218279in}{0.761192in}}%
\pgfpathlineto{\pgfqpoint{4.232564in}{0.682660in}}%
\pgfpathlineto{\pgfqpoint{4.246848in}{0.682660in}}%
\pgfpathlineto{\pgfqpoint{4.261132in}{0.686150in}}%
\pgfpathlineto{\pgfqpoint{4.289701in}{0.754211in}}%
\pgfpathlineto{\pgfqpoint{4.303985in}{0.778062in}}%
\pgfpathlineto{\pgfqpoint{4.318269in}{0.710001in}}%
\pgfpathlineto{\pgfqpoint{4.332554in}{0.740250in}}%
\pgfpathlineto{\pgfqpoint{4.346838in}{0.682660in}}%
\pgfpathlineto{\pgfqpoint{4.389691in}{0.682660in}}%
\pgfpathlineto{\pgfqpoint{4.403975in}{0.750721in}}%
\pgfpathlineto{\pgfqpoint{4.418259in}{0.682660in}}%
\pgfpathlineto{\pgfqpoint{4.446828in}{0.682660in}}%
\pgfpathlineto{\pgfqpoint{4.461112in}{0.689059in}}%
\pgfpathlineto{\pgfqpoint{4.475396in}{0.682660in}}%
\pgfpathlineto{\pgfqpoint{4.489681in}{0.689059in}}%
\pgfpathlineto{\pgfqpoint{4.503965in}{0.703020in}}%
\pgfpathlineto{\pgfqpoint{4.518249in}{0.682660in}}%
\pgfpathlineto{\pgfqpoint{4.532533in}{0.689059in}}%
\pgfpathlineto{\pgfqpoint{4.546818in}{0.682660in}}%
\pgfpathlineto{\pgfqpoint{4.561102in}{0.682660in}}%
\pgfpathlineto{\pgfqpoint{4.575386in}{0.686150in}}%
\pgfpathlineto{\pgfqpoint{4.589670in}{0.723380in}}%
\pgfpathlineto{\pgfqpoint{4.603955in}{0.682660in}}%
\pgfpathlineto{\pgfqpoint{4.618239in}{0.682660in}}%
\pgfpathlineto{\pgfqpoint{4.632523in}{0.713491in}}%
\pgfpathlineto{\pgfqpoint{4.646808in}{0.716400in}}%
\pgfpathlineto{\pgfqpoint{4.661092in}{0.682660in}}%
\pgfpathlineto{\pgfqpoint{4.675376in}{0.686150in}}%
\pgfpathlineto{\pgfqpoint{4.689660in}{0.686150in}}%
\pgfpathlineto{\pgfqpoint{4.703945in}{0.733851in}}%
\pgfpathlineto{\pgfqpoint{4.718229in}{0.737341in}}%
\pgfpathlineto{\pgfqpoint{4.732513in}{0.682660in}}%
\pgfpathlineto{\pgfqpoint{4.746798in}{0.682660in}}%
\pgfpathlineto{\pgfqpoint{4.761082in}{0.723380in}}%
\pgfpathlineto{\pgfqpoint{4.775366in}{0.805403in}}%
\pgfpathlineto{\pgfqpoint{4.789650in}{0.723380in}}%
\pgfpathlineto{\pgfqpoint{4.803935in}{0.682660in}}%
\pgfpathlineto{\pgfqpoint{4.818219in}{0.682660in}}%
\pgfpathlineto{\pgfqpoint{4.832503in}{0.703020in}}%
\pgfpathlineto{\pgfqpoint{4.846787in}{0.682660in}}%
\pgfpathlineto{\pgfqpoint{4.932493in}{0.682660in}}%
\pgfpathlineto{\pgfqpoint{4.946777in}{0.703020in}}%
\pgfpathlineto{\pgfqpoint{4.961062in}{0.682660in}}%
\pgfpathlineto{\pgfqpoint{5.003914in}{0.682660in}}%
\pgfpathlineto{\pgfqpoint{5.018199in}{0.740250in}}%
\pgfpathlineto{\pgfqpoint{5.032483in}{0.682660in}}%
\pgfpathlineto{\pgfqpoint{5.046767in}{0.682660in}}%
\pgfpathlineto{\pgfqpoint{5.061052in}{0.699530in}}%
\pgfpathlineto{\pgfqpoint{5.075336in}{0.750721in}}%
\pgfpathlineto{\pgfqpoint{5.089620in}{0.710001in}}%
\pgfpathlineto{\pgfqpoint{5.103904in}{0.710001in}}%
\pgfpathlineto{\pgfqpoint{5.118189in}{0.682660in}}%
\pgfpathlineto{\pgfqpoint{5.175326in}{0.682660in}}%
\pgfpathlineto{\pgfqpoint{5.189610in}{0.713491in}}%
\pgfpathlineto{\pgfqpoint{5.203894in}{0.723380in}}%
\pgfpathlineto{\pgfqpoint{5.218179in}{0.699530in}}%
\pgfpathlineto{\pgfqpoint{5.232463in}{0.682660in}}%
\pgfpathlineto{\pgfqpoint{5.246747in}{0.699530in}}%
\pgfpathlineto{\pgfqpoint{5.261031in}{0.682660in}}%
\pgfpathlineto{\pgfqpoint{5.332453in}{0.682660in}}%
\pgfpathlineto{\pgfqpoint{5.346737in}{0.692549in}}%
\pgfpathlineto{\pgfqpoint{5.361021in}{0.682660in}}%
\pgfpathlineto{\pgfqpoint{5.375306in}{0.682660in}}%
\pgfpathlineto{\pgfqpoint{5.389590in}{0.730361in}}%
\pgfpathlineto{\pgfqpoint{5.403874in}{0.754211in}}%
\pgfpathlineto{\pgfqpoint{5.418158in}{0.719890in}}%
\pgfpathlineto{\pgfqpoint{5.432443in}{0.730361in}}%
\pgfpathlineto{\pgfqpoint{5.446727in}{0.703020in}}%
\pgfpathlineto{\pgfqpoint{5.461011in}{0.726871in}}%
\pgfpathlineto{\pgfqpoint{5.475296in}{0.692549in}}%
\pgfpathlineto{\pgfqpoint{5.489580in}{0.689059in}}%
\pgfpathlineto{\pgfqpoint{5.503864in}{0.716400in}}%
\pgfpathlineto{\pgfqpoint{5.518148in}{0.682660in}}%
\pgfpathlineto{\pgfqpoint{5.532433in}{0.703020in}}%
\pgfpathlineto{\pgfqpoint{5.546717in}{0.726871in}}%
\pgfpathlineto{\pgfqpoint{5.561001in}{0.740250in}}%
\pgfpathlineto{\pgfqpoint{5.575285in}{0.726871in}}%
\pgfpathlineto{\pgfqpoint{5.589570in}{0.699530in}}%
\pgfpathlineto{\pgfqpoint{5.603854in}{0.737341in}}%
\pgfpathlineto{\pgfqpoint{5.618138in}{0.706510in}}%
\pgfpathlineto{\pgfqpoint{5.632423in}{0.713491in}}%
\pgfpathlineto{\pgfqpoint{5.646707in}{0.682660in}}%
\pgfpathlineto{\pgfqpoint{5.660991in}{0.682660in}}%
\pgfpathlineto{\pgfqpoint{5.675275in}{0.706510in}}%
\pgfpathlineto{\pgfqpoint{5.689560in}{0.761192in}}%
\pgfpathlineto{\pgfqpoint{5.703844in}{0.778062in}}%
\pgfpathlineto{\pgfqpoint{5.718128in}{0.706510in}}%
\pgfpathlineto{\pgfqpoint{5.732413in}{0.692549in}}%
\pgfpathlineto{\pgfqpoint{5.746697in}{0.682660in}}%
\pgfpathlineto{\pgfqpoint{5.760981in}{0.710001in}}%
\pgfpathlineto{\pgfqpoint{5.775265in}{0.757702in}}%
\pgfpathlineto{\pgfqpoint{5.789550in}{0.743740in}}%
\pgfpathlineto{\pgfqpoint{5.803834in}{0.747231in}}%
\pgfpathlineto{\pgfqpoint{5.818118in}{0.682660in}}%
\pgfpathlineto{\pgfqpoint{5.832402in}{0.757702in}}%
\pgfpathlineto{\pgfqpoint{5.860971in}{0.757702in}}%
\pgfpathlineto{\pgfqpoint{5.875255in}{0.723380in}}%
\pgfpathlineto{\pgfqpoint{5.889540in}{0.726871in}}%
\pgfpathlineto{\pgfqpoint{5.903824in}{0.750721in}}%
\pgfpathlineto{\pgfqpoint{5.918108in}{0.706510in}}%
\pgfpathlineto{\pgfqpoint{5.932392in}{0.747231in}}%
\pgfpathlineto{\pgfqpoint{5.946677in}{0.781552in}}%
\pgfpathlineto{\pgfqpoint{5.960961in}{0.682660in}}%
\pgfpathlineto{\pgfqpoint{5.975245in}{0.716400in}}%
\pgfpathlineto{\pgfqpoint{5.975245in}{0.716400in}}%
\pgfusepath{stroke}%
\end{pgfscope}%
\begin{pgfscope}%
\pgfpathrectangle{\pgfqpoint{0.688581in}{0.565123in}}{\pgfqpoint{5.303240in}{2.585803in}}%
\pgfusepath{clip}%
\pgfsetbuttcap%
\pgfsetroundjoin%
\pgfsetlinewidth{1.505625pt}%
\definecolor{currentstroke}{rgb}{1.000000,0.000000,0.000000}%
\pgfsetstrokecolor{currentstroke}%
\pgfsetdash{{5.550000pt}{2.400000pt}}{0.000000pt}%
\pgfpathmoveto{\pgfqpoint{2.544715in}{0.565123in}}%
\pgfpathlineto{\pgfqpoint{2.544715in}{3.150926in}}%
\pgfusepath{stroke}%
\end{pgfscope}%
\begin{pgfscope}%
\pgfsetrectcap%
\pgfsetmiterjoin%
\pgfsetlinewidth{0.803000pt}%
\definecolor{currentstroke}{rgb}{0.000000,0.000000,0.000000}%
\pgfsetstrokecolor{currentstroke}%
\pgfsetdash{}{0pt}%
\pgfpathmoveto{\pgfqpoint{0.688581in}{0.565123in}}%
\pgfpathlineto{\pgfqpoint{0.688581in}{3.150926in}}%
\pgfusepath{stroke}%
\end{pgfscope}%
\begin{pgfscope}%
\pgfsetrectcap%
\pgfsetmiterjoin%
\pgfsetlinewidth{0.803000pt}%
\definecolor{currentstroke}{rgb}{0.000000,0.000000,0.000000}%
\pgfsetstrokecolor{currentstroke}%
\pgfsetdash{}{0pt}%
\pgfpathmoveto{\pgfqpoint{5.991820in}{0.565123in}}%
\pgfpathlineto{\pgfqpoint{5.991820in}{3.150926in}}%
\pgfusepath{stroke}%
\end{pgfscope}%
\begin{pgfscope}%
\pgfsetrectcap%
\pgfsetmiterjoin%
\pgfsetlinewidth{0.803000pt}%
\definecolor{currentstroke}{rgb}{0.000000,0.000000,0.000000}%
\pgfsetstrokecolor{currentstroke}%
\pgfsetdash{}{0pt}%
\pgfpathmoveto{\pgfqpoint{0.688581in}{0.565123in}}%
\pgfpathlineto{\pgfqpoint{5.991820in}{0.565123in}}%
\pgfusepath{stroke}%
\end{pgfscope}%
\begin{pgfscope}%
\pgfsetrectcap%
\pgfsetmiterjoin%
\pgfsetlinewidth{0.803000pt}%
\definecolor{currentstroke}{rgb}{0.000000,0.000000,0.000000}%
\pgfsetstrokecolor{currentstroke}%
\pgfsetdash{}{0pt}%
\pgfpathmoveto{\pgfqpoint{0.688581in}{3.150926in}}%
\pgfpathlineto{\pgfqpoint{5.991820in}{3.150926in}}%
\pgfusepath{stroke}%
\end{pgfscope}%
\begin{pgfscope}%
\definecolor{textcolor}{rgb}{0.000000,0.000000,0.000000}%
\pgfsetstrokecolor{textcolor}%
\pgfsetfillcolor{textcolor}%
\pgftext[x=3.340201in,y=3.234260in,,base]{\color{textcolor}\rmfamily\fontsize{12.000000}{14.400000}\selectfont IF Filter Transmission Loss with 6 dB Bandwidth}%
\end{pgfscope}%
\begin{pgfscope}%
\pgfsetbuttcap%
\pgfsetmiterjoin%
\definecolor{currentfill}{rgb}{1.000000,1.000000,1.000000}%
\pgfsetfillcolor{currentfill}%
\pgfsetfillopacity{0.800000}%
\pgfsetlinewidth{1.003750pt}%
\definecolor{currentstroke}{rgb}{0.800000,0.800000,0.800000}%
\pgfsetstrokecolor{currentstroke}%
\pgfsetstrokeopacity{0.800000}%
\pgfsetdash{}{0pt}%
\pgfpathmoveto{\pgfqpoint{4.344442in}{2.652470in}}%
\pgfpathlineto{\pgfqpoint{5.894598in}{2.652470in}}%
\pgfpathquadraticcurveto{\pgfqpoint{5.922376in}{2.652470in}}{\pgfqpoint{5.922376in}{2.680247in}}%
\pgfpathlineto{\pgfqpoint{5.922376in}{3.053704in}}%
\pgfpathquadraticcurveto{\pgfqpoint{5.922376in}{3.081482in}}{\pgfqpoint{5.894598in}{3.081482in}}%
\pgfpathlineto{\pgfqpoint{4.344442in}{3.081482in}}%
\pgfpathquadraticcurveto{\pgfqpoint{4.316664in}{3.081482in}}{\pgfqpoint{4.316664in}{3.053704in}}%
\pgfpathlineto{\pgfqpoint{4.316664in}{2.680247in}}%
\pgfpathquadraticcurveto{\pgfqpoint{4.316664in}{2.652470in}}{\pgfqpoint{4.344442in}{2.652470in}}%
\pgfpathlineto{\pgfqpoint{4.344442in}{2.652470in}}%
\pgfpathclose%
\pgfusepath{stroke,fill}%
\end{pgfscope}%
\begin{pgfscope}%
\pgfsetrectcap%
\pgfsetroundjoin%
\pgfsetlinewidth{2.007500pt}%
\definecolor{currentstroke}{rgb}{0.121569,0.466667,0.705882}%
\pgfsetstrokecolor{currentstroke}%
\pgfsetdash{}{0pt}%
\pgfpathmoveto{\pgfqpoint{4.372220in}{2.977315in}}%
\pgfpathlineto{\pgfqpoint{4.511109in}{2.977315in}}%
\pgfpathlineto{\pgfqpoint{4.649998in}{2.977315in}}%
\pgfusepath{stroke}%
\end{pgfscope}%
\begin{pgfscope}%
\definecolor{textcolor}{rgb}{0.000000,0.000000,0.000000}%
\pgfsetstrokecolor{textcolor}%
\pgfsetfillcolor{textcolor}%
\pgftext[x=4.761109in,y=2.928704in,left,base]{\color{textcolor}\rmfamily\fontsize{10.000000}{12.000000}\selectfont Transmission Loss}%
\end{pgfscope}%
\begin{pgfscope}%
\pgfsetbuttcap%
\pgfsetroundjoin%
\pgfsetlinewidth{1.505625pt}%
\definecolor{currentstroke}{rgb}{1.000000,0.000000,0.000000}%
\pgfsetstrokecolor{currentstroke}%
\pgfsetdash{{5.550000pt}{2.400000pt}}{0.000000pt}%
\pgfpathmoveto{\pgfqpoint{4.372220in}{2.783642in}}%
\pgfpathlineto{\pgfqpoint{4.511109in}{2.783642in}}%
\pgfpathlineto{\pgfqpoint{4.649998in}{2.783642in}}%
\pgfusepath{stroke}%
\end{pgfscope}%
\begin{pgfscope}%
\definecolor{textcolor}{rgb}{0.000000,0.000000,0.000000}%
\pgfsetstrokecolor{textcolor}%
\pgfsetfillcolor{textcolor}%
\pgftext[x=4.761109in,y=2.735031in,left,base]{\color{textcolor}\rmfamily\fontsize{10.000000}{12.000000}\selectfont Center Frequency}%
\end{pgfscope}%
\end{pgfpicture}%
\makeatother%
\endgroup%

		\caption{Measured return loss and phase of a quartz resonatoras a function of frequency.}
		\label{fig:3514}
	\end{figure}
	
	
	\vfill
	
		
	% --- Bibliography ---
	\bibliographystyle{plain}
	\bibliography{referenzen}
	
	
\end{document}